\documentclass[11pt,a4paper]{article}
\usepackage{fontspec}
\setmainfont{Latin Modern Roman}
\usepackage[margin=23mm,headheight=15pt]{geometry}
\usepackage{amsmath,amssymb,amsthm,mathtools,mathrsfs,amscd}
\usepackage{longtable,booktabs,array,calc}
\usepackage{xcolor,fancyvrb,xurl}
\usepackage[unicode,hidelinks]{hyperref}
\usepackage{bookmark}
\newtheorem{theorem}{Theorem}[section]
\newtheorem{lemma}[theorem]{Lemma}
\newtheorem{proposition}[theorem]{Proposition}
\newtheorem{corollary}[theorem]{Corollary}
\newtheorem{definition}[theorem]{Definition}
\newtheorem{remark}[theorem]{Remark}
\newtheorem{claim}[theorem]{Claim}
\DefineVerbatimEnvironment{Highlighting}{Verbatim}{commandchars=\\\{\},fontsize=\small}
\newenvironment{Shaded}{}{}
\allowdisplaybreaks
\emergencystretch=3em
\hypersetup{pdftitle={Arithmetic Values, Endpoint Control and Exact Receiving Maps},pdfauthor={Split-Zero research programme}}
\title{Arithmetic Values, Endpoint Control\\and Exact Receiving Maps\\[0.6em]\large Complete proof continuation and original providers}
\author{Split-Zero research programme}
\date{19 September 2026}
\begin{document}
\maketitle
\noindent This reader gives the complete accepted calculations received and
extended on 19 September 2026, with their original provider proofs and
human-source citations. The accompanying three cumulative LaTeX successors
retain every byte of their preceding editions by reversible insertion.
This edition has no newly assigned deposited DOI.
\tableofcontents
\begingroup
\clearpage
\part{September 19: arithmetic values and their exact receiving maps}
\noindent The current arithmetic-baseline scalar is evaluated by the complete
proofs below. Earlier statements that this scalar was uncomputed record their
historical state and are superseded by the new value theorem and MRI1--3.
The original absolute directional allocations and projected signs retain the
exact attained metrics specified in MRI4--11. Every earlier source paragraph
remains byte-recoverable; no former edition is overwritten.
\providecommand{\tightlist}{\setlength{\itemsep}{0pt}\setlength{\parskip}{0pt}}
\providecommand{\C}{\mathbb C}\providecommand{\R}{\mathbb R}
\providecommand{\Z}{\mathbb Z}\providecommand{\Q}{\mathbb Q}
\providecommand{\N}{\mathbb N}
\providecommand{\im}{\operatorname{im}}
\providecommand{\diag}{\operatorname{diag}}
\providecommand{\Gr}{\operatorname{Gram}}
\providecommand{\Gram}{\operatorname{Gram}}
\providecommand{\Res}{\operatorname{Res}}
\providecommand{\cond}{\operatorname{cond}}
\providecommand{\ExtensionTok}[1]{#1}\providecommand{\NormalTok}[1]{#1}
\providecommand{\StringTok}[1]{#1}\providecommand{\BuiltInTok}[1]{#1}
\providecommand{\ControlFlowTok}[1]{#1}

\clearpage
\begingroup


\setcounter{equation}{0}
\renewcommand{\theequation}{SA.\arabic{equation}}
\section*{Signed arithmetic value: full received proof and finite repairs}
\addcontentsline{toc}{section}{Signed arithmetic value: full received proof and finite repairs}
% BEGIN PRESERVED INPUT 01_SIGNED_ARITHMETIC_INTEGRATED.tex SHA256 4707e57b9dbd2856a6735e9bfd7eb7ea75e083a359cc2e3dbe3bbed3487e6c14


\section{\texorpdfstring{Signed arithmetic baseline: evaluated return at
order
\(q\)}{Signed arithmetic baseline: evaluated return at order q}}\label{sa:signed-arithmetic-baseline-evaluated-return-at-order-q}

18 September 2026. New written research derivation on the original
fixed-packet domain. The proof and its imported inputs are identified
below; this is not an independent formal verification of the entire
programme.

\subsection{M1. The original assigned
quantity}\label{sa:m1.-the-original-assigned-quantity}

Retain the programme's stipulated quartet and its complete multiplicity:
\[
\rho=\tfrac12+\delta+i\gamma,\qquad 0<\delta<\tfrac12,\quad\gamma>2,
\qquad h(s)=\prod_{z\in\{\rho,\bar\rho,1-\rho,1-\bar\rho\}}(s-z)^{m_0}.
\] The original source and reference are \[
v_h=2\xi/h,\qquad w_h(y)=\frac{|v_h(1/2+iy)|^2}{2\pi},
\quad m_k=w_h^{*k},\quad
\sigma(y)=\frac{|\Gamma(1/4+iy/2)|^2}{2\pi}.
\] Their complete masses are retained; in particular
\(\int\sigma=\sqrt{2\pi}\). The admitted tensor orders and packet
dimension are exactly \[
k=4\ell+1\ge5,\qquad e=1+k(m_0-1),\qquad q=e(k+1)^2.
\] The full relation polynomial is \[
\chi_k(S)=\prod_{a,b_0=0}^k
\left[S-\frac k2-(2a-k)\delta-i(2b_0-k)\gamma\right]^e.
\] No numerical off-critical zero is supplied or asserted by this
calculation. The fixed original packet is the same one used in the
incoming arithmetic construction.

For its remainder map \(J_N=\operatorname{rem}_{\chi_k}\), write \[
G_N[\nu]=(J_NH_N[\nu]^{-1}J_N^*)^{-1},
\] \[
\mathcal B_k[\nu]=\log\det G_{q-1}[\nu]+\log\det G_q[\nu]
-\log\det G_{2q-1}[\nu]-\log\det G_{2q}[\nu].
\] The assigned scalar is \[
\boxed{\delta_k=\delta_{k,1}^{\rm ar}
=\mathcal B_k[m_k]-\mathcal B_k[\sigma].}\tag{M1}
\] The original full arithmetic unit remains in \(w_h\), in the
programme's inclusion, and in the fixed coefficient identifications. The
source-coordinate map is explicitly \(P(S)\mapsto P(k/2+iy)\), with
inverse \(y=(S-k/2)/i\); it carries the complete relation to
\(i^qQ(y)\), where \(Q(y)=i^{-q}\chi_k(k/2+iy)\). Monic phases have
modulus one. A fixed change of quotient coordinates contributes the same
squared determinant at all four endpoints, so its four-sign contribution
is exactly zero rather than an assumed isometry.

\subsection{M2. The evaluated
coefficient}\label{sa:m2.-the-evaluated-coefficient}

Define the fixed one-factor quantities \[
\alpha=\pi/2,\qquad M_h(\theta)=\int_{\mathbb R}e^{\theta y}w_h(y)dy,
\qquad \mu_h(\theta)=\frac{M_h'(\theta)}{M_h(\theta)},
\quad0\le\theta\le\alpha,
\] \[
\boxed{I_h=\int_0^{\pi/2}\mu_h(\theta)^2d\theta.}\tag{M2}
\] The RTM envelope gives finite first and second boundary-tilted
moments, and the preceding central calculation proves \(0<I_h<\infty\).
In particular this coefficient is a fixed integral of the original
one-factor source, not a limit of growing minimum matrices.
Equivalently, retaining the exact evenness factors, \[
I_h=\int_0^{\pi/2}
\left[
\frac{\int_0^\infty y\sinh(\theta y)w_h(y)dy}
{\int_0^\infty \cosh(\theta y)w_h(y)dy}
\right]^2d\theta.
\tag{M3}
\]

Retain the programme's already evaluated universal constant \[
\boxed{\mathcal C_\Gamma=2\mathcal L(2,\pi)-8\log2+4>0.}\tag{M4}
\] Its complete fixed endpoint integral is specified in M4 below. Put
\(d=4m_0-2\).

The evaluated signed return is \[
\boxed{\delta_k
=-(4m_0-2)\mathcal C_\Gamma q
+\frac{\log2}{\pi}I_h k^2+o_h(q).}\tag{M5}
\] More precisely, there is a fixed-source constant \(C_h\), on the
cofinal range established in the proof, such that \[
\boxed{
\left|\delta_k+(4m_0-2)\mathcal C_\Gamma q
-\frac{\log2}{\pi}I_h k^2\right|
\le C_hq\,\mathscr R(q),
}\tag{M6}
\] \[
\mathscr R(q)=
\frac{\log\log(q+16)}{\log(q+2)}
+q^{-1/200}\log^2(q+2)
\longrightarrow0.
\] The exponent in this explicit remainder is deliberately conservative.
It comes from the averaged original zeta tail, the full Gamma projection
bound, the Schur localization, and the quantified convex-secant
argument; it is not a fitted numerical error.

Consequently the requested coefficients are \[
\boxed{d_{\log}(h)=0,}\tag{M7}
\] \[
\boxed{
 d_0(h)=
 \begin{cases}
 -2\mathcal C_\Gamma+\dfrac{\log2}{\pi}I_h,&m_0=1,\\[5pt]
 -(4m_0-2)\mathcal C_\Gamma,&m_0\ge2.
 \end{cases}}
\tag{M8}
\] This follows from the actual dimension, not a redefinition of scale:
\[
\frac{k^2}{q}=
\frac{k^2}{[1+k(m_0-1)](k+1)^2}
\longrightarrow
\begin{cases}1,&m_0=1,\\0,&m_0\ge2.\end{cases}
\] Thus \(\delta_k=q\,d_0(h)+o_h(q)\). There is no additional term
between orders \(q\) and \(q\log k\), including no \(q\log\log k\) term.
For higher multiplicity, the separately evaluated central term in (M5)
is of order \(k^2\), but the displayed full-return remainder is only
claimed to be \(o(q)\); (M5) is not an assertion that every term of the
total return through order \(k^2\) has been evaluated there.

\subsection{M3. How the two previously unfinished values are
completed}\label{sa:m3.-how-the-two-previously-unfinished-values-are-completed}

Let \(\mathfrak M=M_h(\alpha)\), \(b=\mu_h(\alpha)\), and let
\(\theta(t)\) invert \(\mu_h\) on \([0,b]\). Define \[
D(t)=\log\mathfrak M-\log M_h(\theta(t))-(\alpha-\theta(t))t,
\] \[
W_k(y)=kD(|y|/k)1_{|y|\le kb}.
\] The previous two source calculations have exactly the same rate
function, with opposite notational sign: \(F_h^{\rm cen}(t)=-D(|t|)\).
They already calculate \(\int W_k=k^2I_h\). The present calculation
evaluates the determinant response to this rate, rather than relabeling
its Lebesgue integral as a determinant coefficient.

The literal endpoint decomposition is \[
\omega_k(y)=
\begin{cases}\sigma(y),&|y|\le kb,\\
\mathfrak M^{-k}m_k(y),&|y|>kb,
\end{cases}
\] \[
\delta_k^{\rm out}=\mathcal B_k[\omega_k]-\mathcal B_k[\sigma],\qquad
\delta_k^{\rm cen,end}=\mathcal B_k[m_k]-\mathcal B_k[\omega_k],
\] \[
\boxed{\delta_k=\delta_k^{\rm out}+\delta_k^{\rm cen,end}.}\tag{M9}
\] The mass factor is legitimate through the exact relation
\(G_N[A\nu]=AG_N[\nu]\): its logarithm contributes
\((q+q-q-q)\log A=0\). The outer density is fixed throughout the central
endpoint change.

\textbf{Outer value.} The boundary-tilted convolution is averaged
against its actual zeta factor, using a single-large-summand
decomposition, the proved uniform local Fourier estimate, and Ivić's
short-interval second-moment bound for \(E(t+U)-E(t)\). The resulting
Lebesgue averaged approximation is carried into the full source and
relation determinants by an explicit Meixner--Pollaczek projection bound
and the two-sided determinant inequality (O20). This yields \[
\delta_k^{\rm out}
=\mathcal B_k[(1+y^2)^{-d}\sigma]-\mathcal B_k[\sigma]+o(q).
\] The actual rational model is then evaluated: its source contribution
is \[
4dq\log q+(8d\log2-4d)q+o(q),
\] while its full relation contribution, with its original signs, is \[
-4dq\log q-2dq\mathcal L(2,\pi)+o(q).
\] Therefore \[
\boxed{\delta_k^{\rm out}=-d\mathcal C_\Gamma q+o_h(q).}\tag{M10}
\] The error is quantified in (O35). The relation calculation uses an
exact convex function of the original determinant; the quantitative
appendix proves the shrinking-secant estimate without differentiating a
free-energy error.

\textbf{Central value.} The established localization gives the original
central endpoint change as two bands of Gamma monic-norm losses, with
\(O_h(k\log^2q)\) error. An explicit Gaussian regularization of the
convex potentials \(\alpha|y|+tW_k(y)\) is returned through the same
Taylor-jet Schur quotients with \(o(k^2)\) error. For these regularized
potentials the full mass-\(n\) equilibrium is constructed explicitly and
the monic squared norm is proved uniformly to be \[
H_n(t)=\pi A_{n,t}e^{-\lambda_{n,t}}[1+O_h(n^{-1})],
\qquad q\le n\le2q.
\] The exact Robin variation is \[
\partial_t\lambda_{n,t}
=\frac1\pi\int_{-A_{n,t}}^{A_{n,t}}
\frac{\widetilde W_k(y)}{\sqrt{A_{n,t}^2-y^2}}dy
=\frac{k^2I_h}{2\pi n}+O_h(k^4/n^3).
\] Summing the two original bands retains the weights
\(1,2,\ldots,2,1\), so that \(\sum_{r=0}^q c_r/(q+r)=2\log2+O(1/q)\).
Consequently \[
\boxed{\delta_k^{\rm cen,end}
=\frac{\log2}{\pi}I_hk^2+o_h(k^2).}\tag{M11}
\] The stronger error is in (C32). Appendix A gives the exact matrix
jumps, Airy normalization, and the shrinking-strip Cauchy estimate
needed for uniformity. No fixed-potential theorem is presumed uniform on
this changing family.

Equations (M9)--(M11) prove (M5). Their explicit remainders prove (M6),
and the exact dimension gives (M7)--(M8).

\subsection{M4. The universal constant has no remaining growing
determinant}\label{sa:m4.-the-universal-constant-has-no-remaining-growing-determinant}

The EIQ convention is \[
V_{\alpha_E,\beta}(x)=\beta\sqrt x-\alpha_E\log x,
\] with \(\mathcal F\) equal to the supremum of logarithmic energy minus
this potential. At \((\alpha_E,\beta)=(2,\pi)\), take the unique
\(0<\kappa<1\), \(0<u<v\) satisfying \[
uK(\kappa)=2,\qquad vE(\kappa)=4,\qquad
\kappa^2=1-u^2/v^2.
\tag{M12}
\] Here \(K,E\) are the complete elliptic integrals in the original real
convention. Put \[
A_E=u^2,\quad B_E=v^2,\quad m_E=(A_E+B_E)/2,\quad
R_s=\sqrt{(A_E+s)(B_E+s)},\quad R_0=uv,
\] \[
C_E=(u+v)^2/4,\qquad \eta_E=(v-u)/(v+u),\qquad
\zeta_s=\frac{B_E-A_E}{2(m_E+s+R_s)},
\] \[
L_s=(m_E+s-R_s)\log C_E-(m_E-R_0)
-2R_s\log(1-\eta_E\zeta_s).
\] Then the same original logarithmic moment is \[
\boxed{\mathcal L(2,\pi)=
\frac14\int_0^\infty\frac{L_0-L_s}{\sqrt s\,R_s}ds.}\tag{M13}
\] The integral converges at both ends. This is precisely EIQ24--33, not
a new fitted constant. The source proves \(\mathcal C_\Gamma>0\)
independently of numerical quadrature. A high-precision non-interval
diagnostic gives \[
\mathcal C_\Gamma\approx0.654099894824426,\qquad
\log2/\pi\approx0.220635600152652.
\] The exact expressions (M2), (M4), and (M12)--(M13) govern the result;
these decimals are not used to prove it. Appendix I1--I4 gives an
explicit finite truncation and quadrature error for the fixed source
functional \(I_h\).

\subsection{M5. Return to the complete original
action}\label{sa:m5.-return-to-the-complete-original-action}

Retain the source's exact identity \[
J_k^{\rm action}=\mathcal B_k[\sigma]+\delta_k
+W_k^{\rm mono}-P_{k,-}-P_{k,+}.
\] The superscript on \(W_k^{\rm mono}\) distinguishes the original
monic-window scalar from the central insertion \(W_k(y)\). Substitution
of the evaluated correction gives \[
\boxed{\begin{aligned}
J_k^{\rm action}={}&\mathcal B_k[\sigma]
-(4m_0-2)\mathcal C_\Gamma q
+\frac{\log2}{\pi}I_hk^2\\
&+W_k^{\rm mono}-P_{k,-}-P_{k,+}
+\mathcal E_{h,k},
\qquad |\mathcal E_{h,k}|\le C_hq\mathscr R(q).
\end{aligned}}\tag{M14}
\] The Gamma baseline remains the exact baseline, not just its positive
leading \(q^2\) term. The correction is \(O_h(q)\), and therefore does
not cancel that \(q^2\) term. The independent proper-source metric,
directional allocations, and retained-class projected action are not
assigned values by this scalar calculation.

The assignment recorded as section 15.1 is thus evaluated through its
requested \(q\)-scale precision. No purity or Frobenius value has been
substituted for the arithmetic minima. The derivation does not assert
programme closure or the existence of an off-critical zero.

\newpage

\section{The outer arithmetic return: its value, not only a
bound}\label{sa:the-outer-arithmetic-return-its-value-not-only-a-bound}

This is a new derivation on the fixed original packet. The source inputs
are the RTM envelopes, the consolidated polynomial-prefactor comparison,
the original Gamma normalization, and the original source/relation
determinant identity. The external analytic input is the established
short-interval mean-square bound for the zeta mean-square error in Ivić
{[}I, (2.1){]}. All subsequent comparisons below are made on the
specified polynomial subspaces and attained quotients.

\subsection{O1. Objects and exact scalar
bookkeeping}\label{sa:o1.-objects-and-exact-scalar-bookkeeping}

Retain \[
\alpha=\pi/2,\quad m=m_0\ge1,\quad d=4m-2,\quad
p=2d-\tfrac12,\quad a=2d+\tfrac12=p+1,
\] \[
k=4\ell+1,\quad e=1+k(m-1),\quad q=e(k+1)^2,
\quad Q(y)=i^{-q}\chi_k(k/2+iy).
\] The degree of the real even monic polynomial \(Q\) is exactly \(q\),
with all original repetitions. Its roots have modulus at most
\(R=k\sqrt{\delta^2+\gamma^2}\).

Write \[
M(\theta)=\int e^{\theta y}w_h(y)\,dy,\quad
\mathfrak M=M(\alpha),\quad b=(\log M)'(\alpha),
\quad f(y)=\mathfrak M^{-1}e^{\alpha y}w_h(y).
\] Thus \(f\) is an auxiliary probability density of mean \(b\) and
variance \(V_b>0\), and the original convolution is recovered exactly by
\[
\boxed{\mathfrak M^{-k}m_k(y)=e^{-\alpha y}f^{*k}(y).}\tag{O1}
\] The original measure is not replaced by this auxiliary probability.
Equation (O1) records the complete inverse factor.

Let \[
\omega_k(y)=\begin{cases}\sigma(y),&|y|\le kb,\\
\mathfrak M^{-k}m_k(y),&|y|>kb,
\end{cases}
\qquad \tau_d(y)=(1+y^2)^{-d}\sigma(y).
\] The outer endpoint return is \[
\delta_k^{\rm out}=\mathcal B_k[\omega_k]-\mathcal B_k[\sigma].
\] We prove \[
\boxed{\delta_k^{\rm out}=-d\mathcal C_\Gamma q+o_h(q),
\qquad \mathcal C_\Gamma=2\mathcal L(2,\pi)-8\log2+4.}\tag{O2}
\] Here \(\mathcal L\) is the original EIQ logarithmic moment, with its
original potential and normalization. Its fixed integral is reproduced
in the main proof.

For a scalar \(A>0\), the identity on every affine remainder fibre gives
\(G_N[A\nu]=A G_N[\nu]\). Therefore \[
\mathcal B_k[A\nu]-\mathcal B_k[\nu]
=(q+q-q-q)\log A=0.\tag{O3}
\] This is the scalar cancellation used below. No scalar is removed from
an individual determinant.

\subsection{O2. The actual one-factor
tail}\label{sa:o2.-the-actual-one-factor-tail}

The retained arithmetic density is exactly \[
w_h(y)=\frac{(y^2+1/4)^2}{\sqrt\pi}
\frac{|\zeta(1/2+iy)|^2}
{[((y-\gamma)^2+\delta^2)((y+\gamma)^2+\delta^2)]^{2m}}
\sigma(y).
\] Consequently, as \(y\to+\infty\), \[
\boxed{f(y)=c_f y^{-a}|\zeta(1/2+iy)|^2(1+O_h(y^{-1})),
\qquad c_f=\frac{\sqrt{2/\pi}}{\mathfrak M}.}\tag{O4}
\] Indeed the rational factor is \(y^{-2d}(1+O_h(y^{-2}))\), and
\(e^{\alpha y}\sigma(y)=\sqrt2 y^{-1/2}(1+O(y^{-1}))\). The Gamma mass
and the factor \(\sqrt\pi\) are thus retained.

The original envelope yields \[
0\le f(y)\le C_h(1+|y|)^{-p},\qquad p\ge7/2.
\tag{O5}
\] Its negative tail has the additional factor \(e^{-2\alpha|y|}\). In
particular the variance and a \(9/4\)-moment are finite.

Put \(Z(y)=|\zeta(1/2+iy)|^2\) and \[
\mathcal L_\zeta(y)=\log y+c_\zeta,\qquad
c_\zeta=2\gamma_{\rm E}-\log(2\pi).
\] The classical mean-square formula gives, uniformly for
\(|z|\le T^{19/20}\), \[
\int_T^{2T}f(y-z)\,dy\le C_h T^{1-a}\log T.\tag{O6}
\] Only this upper bound is used in the convolution comparison; no
pointwise lower bound for \(Z\) is used. For example, the usual error
bound \(E(T)=O(T^{1/2}\log^2T)\), or any established \(o(T\log T)\)
bound, implies \(\int_{T/2}^{3T}Z\ll T\log T\). Equations (O4) and (O6)
then follow on the indicated shifted interval. The same estimate at a
finite number of comparable dyadic intervals will always be understood
literally.

\subsection{O3. Uniform total variation at the boundary
tilt}\label{sa:o3.-uniform-total-variation-at-the-boundary-tilt}

The incoming Fourier proof gives, for every real \(x\), \[
\left|f^{*n}(nb+x)-
\frac{e^{-x^2/(2nV_b)}}{\sqrt{2\pi nV_b}}\right|
\le C_h n^{-(1+1/4)/2}.\tag{O7}
\] Translation in Fourier inversion has modulus one, so the absolute
error at the mean has this same bound at every shifted point.

Let \(g_n\) be the Gaussian density in (O7), centered at \(nb\).
Integrate (O7) on \(|x|\le\sqrt n\,n^{1/24}\). This costs
\(O_h(n^{-1/12})\). On its complement, Chebyshev's inequality for each
of the two probability measures costs \(O_h(n^{-1/12})\). Therefore \[
\boxed{\|f^{*n}-g_n\|_{L^1(\mathbb R)}\le C_h n^{-1/12}.}\tag{O8}
\] This is a total-variation estimate for the actual tilted convolution;
all masses in (O1) remain unchanged.

\subsection{O4. A single large summand, in the averaged norm actually
needed}\label{sa:o4.-a-single-large-summand-in-the-averaged-norm-actually-needed}

Throughout this section take \[
q^{9/10}\le T\le128q.
\] For each fixed original multiplicity and packet, this gives \[
c_hT^{1/3}\le k\le C_hT^{5/9}.\tag{O9}
\] Set \(R_T=T^{19/20}\), \(f_s=f1_{|y|\le R_T}\), and \(f_l=f-f_s\).
These are subprobability densities, not normalized substitutes.

Expand \((f_s+f_l)^{*k}\). Its terms with at least two large summands
have total mass at most \[
C_h k^2R_T^{-2(p-1)}.\tag{O10}
\] Relative to \(kT^{1-a}\log T=kT^{-p}\log T\), this is \[
O_h\bigl(kT^pR_T^{-2(p-1)}/\log T\bigr)=O_h(T^{-2/3})
\] on the cofinal range; the worst exponent occurs at \(p=7/2\), and is
\(5/9+7/2-(19/10)(5/2)=-25/36\).

The all-small convolution has exponentially small mass in \([T,2T]\). To
check this directly, use \(\lambda=(\log T)/R_T\). Taylor's formula and
the finite second moment give \[
\int e^{\pm\lambda y}f_s(y)dy
\le\exp(C_h\lambda+C_h\lambda^2T).
\] Thus either tail beyond \(T/2\) is at most \[
\exp[-c_hT^{1/20}\log T],\tag{O11}
\] for all sufficiently large original \(k\). The inequalities
\(k/T\to0\) and \(k\log T/R_T\to0\), from (O9), justify the last
assertion.

The remaining term is exactly \(k f_l*f_s^{*(k-1)}\). Split its second
variable \(z\) into \(|z|\le R_T\), \(R_T<|z|\le T/2\), and \(|z|>T/2\).
On the middle region, (O5) and the variance of the original sum give the
integrated bound \[
C_h k^2T^{1-p}/R_T^2.\tag{O12}
\] For clarity, \(f_s^{*(k-1)}\le f^{*(k-1)}\), and on \(|z|>R_T\) one
has \(|z-(k-1)b|\ge R_T/2\). Its mass is therefore at most
\(C_h k/R_T^2\). The first density is at most \(C_hT^{-p}\) on the
relevant shifted window. Relative to the leading mass, (O12) is
\(O_h(kT/R_T^2)\), hence a negative fixed power of \(T\). The final
region is bounded by (O11), using the full mass of \(f_l\).

On \(|z|\le R_T\), the condition \(y-z>R_T\) holds for every
\(T\le y\le2T\). Replacing the truncated rest convolution by the full
one has total-variation cost at most \((k-1)\int_{|u|>R_T}f(u)du\). The
important bound for its convolution into this window is (O6), not a
pointwise bound charged over the whole interval. The resulting relative
error is at most \(C_h kR_T^{1-p}\).

Replace this full rest convolution by its Gaussian using (O8) and the
same shifted-window estimate (O6). Its relative error is
\(O_h(k^{-1/12})\). The Gaussian mass outside \(|z|\le R_T\) is
exponentially small. Finally (O4), uniformly on the retained shifted
window, allows replacement of \((y-z)^{-a}\) by \(y^{-a}\) at relative
integrated cost \(O_h(R_T/T)\).

Put \(n=k-1\), \(G=\sqrt{2nV_b}\), and \[
\mathcal Z_G(t)=\frac1{\sqrt\pi G}
\int_{\mathbb R}e^{-(u/G)^2}Z(t+u)\,du.
\] The preceding decomposition proves \[
\boxed{\begin{aligned}
\int_T^{2T}\left|f^{*k}(y)-kc_fy^{-a}
\mathcal Z_G(y-nb)\right|dy
\le C_h kT^{1-a}\log T\,[k^{-1/12}+T^{-1/20}].
\end{aligned}}\tag{O13}
\] Every discarded event in deriving this inequality has the displayed
quantitative bound. No pointwise regular-variation assertion about the
zeta factor has been made.

\subsection{O5. The zeta factor is averaged with an adequate
error}\label{sa:o5.-the-zeta-factor-is-averaged-with-an-adequate-error}

The external input {[}I, (2.1){]} is the established bound \[
\int_T^{2T}|E(t+U)-E(t)|^2dt
\ll_\varepsilon TU\log^3T+
T^{1/2+\varepsilon}U^2+T^{1+\varepsilon}U^{1/2},
\quad 1\ll U\le\tfrac12\sqrt T,
\tag{O14}
\] where \[
E(T)=\int_0^T Z(t)dt-T[\log(T/2\pi)+2\gamma_{\rm E}-1].
\] This bound concerns the actual \(E\), not a conjectural pointwise
short-interval estimate.

Integration by parts against the Gaussian, subtracting the constant
\(E(t)\) under the odd kernel, gives \[
\mathcal Z_G(t)-\mathcal L_\zeta(t)
=\frac2{\sqrt\pi G^3}\int u e^{-(u/G)^2}
[E(t+u)-E(t)]du+O(G/T)+O(T^{-100}).\tag{O15}
\] One can restrict \(|u|\le G\log T\) before this operation. The
omitted Gaussian tails are bounded by any chosen inverse power, using
polynomial bounds already implied by (O4)--(O5). On that range the
arguments are positive and comparable to \(T\).

For \(|u|\ge1\), apply (O14), using a finite dyadic cover after a
negative translation. Minkowski's inequality yields \[
\boxed{T^{-1/2}\|\mathcal Z_G-\mathcal L_\zeta\|_{L^2[T,2T]}
\ll_\varepsilon
G^{-1/2}\log^{3/2}T+T^{-1/4+\varepsilon}
+T^\varepsilon G^{-3/4}+G/T.}\tag{O16}
\] The part \(|u|<1\) needs no fourth-moment input: by positivity of
\(Z\), its error increment is bounded by a length-four increment of
\(E\), plus \(O(\log T)\). Equation (O14) at this fixed length gives a
contribution \(O(T^\varepsilon G^{-3})\), already smaller than the right
side of (O16).

In the present range, \(c_hT^{1/6}\le G\le C_hT^{5/18}\). Thus the use
of (O14), including \(G\log T<\sqrt T/2\), is valid on a cofinal range.
Fix \(\varepsilon=1/1000\). Apply Cauchy--Schwarz to (O16), account for
the shift \(nb=O_h(k)\), and combine with (O13). With, for example,
\(\eta_0=1/200\), one obtains \[
\boxed{\int_T^{2T}
\left|\frac{\mathfrak M^{-k}m_k(y)}
{\frac{k}{\sqrt\pi\mathfrak M}\mathcal L_\zeta(y)
(1+y^2)^{-d}\sigma(y)}-1\right|dy
\le C_hT^{1-\eta_0}.}\tag{O17}
\] The original evenness supplies the corresponding negative interval.
This is the averaged estimate subsequently used in full determinants. It
is not a pointwise replacement of the arithmetic density.

\subsection{O6. A uniform bound for the original Gamma projection
kernel}\label{sa:o6.-a-uniform-bound-for-the-original-gamma-projection-kernel}

Let \(K_N^\sigma(y,z)=\sum_{j=0}^N\varphi_j(y)\overline{\varphi_j(z)}\).
Then \[
\boxed{K_N^\sigma(y,y)\sigma(y)\le20\log(N+2),
\quad y\in\mathbb R,\ N\ge0.}\tag{O18}
\] Here is a derivation preserving the normalization. In the orthonormal
Meixner--Pollaczek convention of Koelink--Van der Jeugt {[}K,
Proposition 2.1{]}, take \(\lambda=1/4\), angle \(\pi/2\), and
\(x=y/2\). Their weight is
\(w_\lambda(x)=2^{2\lambda}|\Gamma(\lambda+ix)|^2/(2\pi)\), and the
exact map is \(\varphi_j(y)=2^{\lambda-1/2}p_j^{(\lambda)}(y/2)\). Hence
the original diagonal density is one half of their diagonal density.

For \(0<t<1\), put \(R_t=4t/(1-t)^2\). Their Poisson kernel on the
diagonal is \[
\frac{(1+t)^{2ix}(1-t)^{-2\lambda-2ix}}{\Gamma(2\lambda)}
{}_2F_1(\lambda+ix,\lambda+ix;2\lambda;-R_t).
\] Euler's integral (DLMF 15.6.1) is valid because
\(\Re(2\lambda)>\Re(\lambda+ix)>0\). Its Gamma denominator is exactly
\(|\Gamma(\lambda+ix)|^2\), which cancels the displayed weight. Every
remaining imaginary power has modulus one. Thus the weighted kernel is
at most \[
\frac{2^{2\lambda}}{2\pi}(1-t)^{-2\lambda}
\int_0^1u^{\lambda-1}(1-u)^{\lambda-1}(1+R_tu)^{-\lambda}du.
\] For \(t\ge1/2\), split at \(1/R_t\) and \(1/2\). At \(\lambda=1/4\),
the integral is at most \(R_t^{-\lambda}(16+2\log R_t)\). The original
half-weight factor and
\((1-t)^{-2\lambda}R_t^{-\lambda}=(4t)^{-\lambda}\) therefore bound the
original Poisson density by \((16+2\log R_t)t^{-1/4}/(4\pi)\).

All diagonal summands are nonnegative. Take \(t=(N+1)/(N+2)\), multiply
this bound by \(t^{-N}\le e\), and use \(R_t\le4(N+2)^2\). This proves
the enlarged explicit bound (O18).

For every fixed subspace \(S\subset\mathcal P_{2q}\), its
Gamma-orthogonal projector satisfies \[
K_S(y,y)\le K_{2q}^\sigma(y,y).\tag{O19}
\] This includes each complete original relation space, rather than only
monomials.

\subsection{O7. A determinant sandwich on the entire
subspace}\label{sa:o7.-a-determinant-sandwich-on-the-entire-subspace}

Let \(S\) be any finite polynomial subspace, in a Gamma-orthonormal
frame, and let \(r>0\) have the required integrable moments and
logarithms. Expansion of the Gram determinant gives \[
\det\operatorname{Gram}_S[r\sigma]
=\frac1{s!}\int |\det[\phi_i(y_j)]|^2\prod_jr(y_j)\,d\sigma(y_j).
\] At \(r=1\) the right-hand measure has total mass one, and its
one-point density is \(K_S(y,y)d\sigma(y)\). Jensen's inequality, and
separately \(\log\det A\le\operatorname{Tr}(A-I)\), give \[
\boxed{\int\log r\,K_Sd\sigma
\le\log\det\operatorname{Gram}_S[r\sigma]
\le\int(r-1)K_Sd\sigma.}\tag{O20}
\] All determinant cross terms remain in the squared determinant. This
is why a Lebesgue averaged tail estimate can be used without assuming an
arithmetic-kernel limit.

\subsection{O8. Patch a sublinear interval through a proved Schur
comparison}\label{sa:o8.-patch-a-sublinear-interval-through-a-proved-schur-comparison}

Set \(K=q^{9/10}\), \[
A_k=\frac{k\mathcal L_\zeta(q)}{\sqrt\pi\mathfrak M}>0,
\qquad
\widehat\omega_k(y)=\begin{cases}A_k\tau_d(y),&|y|\le K,\\
\omega_k(y),&|y|>K.
\end{cases}\tag{O21}
\] This operation has an \(o(q)\) signed determinant cost; it is not a
deletion of the inner interval.

Indeed the consolidated whole-line polynomial bounds give, on degree at
most \(2q\), \[
q^{-C_h}\|P\|_\sigma^2\le
\|P\|_{\omega_k}^2,\ \|P\|_{\widehat\omega_k}^2,
\ A_k\|P\|_{\tau_d}^2
\le q^{C_h}\|P\|_\sigma^2.\tag{O22}
\] Polynomial lower weights are transferred by Jensen and the original
Gamma multiplication bound. The scalar \(A_k\) is also between fixed
powers of \(q\).

Use the first \(r\) Taylor coefficients at zero, restricted to each
actual source or relation subspace. Its image has dimension at most
\(r\); its full kernel consists of the elements vanishing to order
\(r\). The generating function of the original Gamma polynomials and
Cauchy's estimate give \[
\int_{-K}^{K}|P|^2d\sigma
\le4^{1-r}e^2(2q+1)^2\sqrt{2q+2}
[2(2q+2)]^{2K}\|P\|_\sigma^2
\tag{O23}
\] on that kernel. Choose \(r=\lceil C_h(K+1)\log(q+2)\rceil\), with the
constant enlarged so that the right-hand fraction, after (O22), is at
most \(q^{-20}\). This has \(r=o(q)\).

Factor the determinant into its restriction to this kernel and its
attained quotient on the actual Taylor image. The kernel's relative
log-determinant change is \(O(q^{-18})\). The quotient, of dimension at
most \(r\), costs \(O_h(r\log q)\) by (O22). This proves at every
original endpoint and on every original relation space \[
|\log\det H_S[\widehat\omega_k]-\log\det H_S[\omega_k]|
=O_h(q^{9/10}\log^2q)=o(q).\tag{O24}
\] The fixed coordinate determinant in this Schur factorization cancels
between the two metrics. The Taylor observation is only an auxiliary
comparison; it does not replace the programme's observation or its
relation ideal.

\subsection{O9. Evaluate the patched density in the full
determinants}\label{sa:o9.-evaluate-the-patched-density-in-the-full-determinants}

Put \[
r_k(y)=\widehat\omega_k(y)/(A_k\tau_d(y)).
\] It is exactly one on \([-K,K]\). On each dyadic outer interval below
\(128q\), equation (O17) says that it differs in \(L^1\) from \[
a_q(y)=\mathcal L_\zeta(|y|)/\mathcal L_\zeta(q)
\] by at most \(C_hT^{1-\eta_0}\). Here \(a_q\) stays between two fixed
positive constants. The polynomial lower and upper bounds on \(r_k\)
imply on this region \[
|\log r_k-\log a_q|\le C_h\log q\,|r_k-a_q|.\tag{O25}
\] This follows by separating \(|r_k-a_q|\le a_q/2\) from its
complement; no pointwise zeta lower bound is used.

Equations (O18)--(O19) show that the errors in both sides of the
determinant sandwich (O20) are \(O_h(q^{1-\eta_0}\log^2q)=o(q)\). For
the deterministic factor \(a_q\), on \(q/\log^4q\le|y|\le128q\) it
differs from one by \(O(\log\log q/\log q)\), so use the exact kernel
mass \(\dim S\le2q+1\). On the remaining interval
\(K\le|y|\le q/\log^4q\), use its length and (O18). Both costs are
\(o(q)\).

The whole remaining tail is retained. From (O22) and the pointwise
envelopes, \(r_k(y)\le q^{C_h}(1+y^2)^d\) and
\(|\log r_k(y)|\le C_h[\log q+\log(1+y^2)]\). Repeated Gamma
multiplication, with an extra power \(y^q\), bounds the integral of
these functions against \(K_{2q}^\sigma d\sigma\) on \(|y|>128q\) by
\(e^{-c_hq}\) times a fixed power of \(q\). For example \[
\int_{|y|>128q}y^{2d}|P|^2d\sigma
\le(128q)^{-2q}[2(3q+d+1)]^{2(q+d)}\|P\|_\sigma^2
\] for \(\deg P\le2q\), and then sum an orthonormal frame. Thus \[
\boxed{\int|r_k-1|K_Sd\sigma=o_h(q),\qquad
\int|\log r_k|K_Sd\sigma=o_h(q)}\tag{O26}
\] uniformly for all the subspaces needed here.

The reference for our desired comparison is \(\tau_d\), not \(\sigma\).
Remove this fixed denominator by the exact multiplication map \[
P\longmapsto(y+i)^dP.
\] It is an isometry from \(\mathcal P_{N-d}\) in \(\sigma\) onto
\((y+i)^d\mathcal P_{N-d}\) in \(\tau_d\), because
\(|y+i|^{2d}=(1+y^2)^d\). With the factor \(r_k\), the same map
transports \(r_k\sigma\) to \(r_k\tau_d\). On a relation space use
instead \(QP\mapsto Q(y+i)^dP\). The omitted fixed quotient has
dimension at most \(d\), and its two attained metrics have logarithmic
relative cost \(O_h(d\log q)\) by (O22). All these are fixed subspaces
and exact minimum quotients.

Apply (O20) and (O26) after this multiplication. We obtain the complete
source and relation comparisons, each with error \(o_h(q)\). Summing
their original signs, using (O3) for \(A_k\), and returning the patch by
(O24), proves \[
\boxed{\delta_k^{\rm out}
=\mathcal B_k[\tau_d]-\mathcal B_k[\sigma]+o_h(q).}\tag{O27}
\] The zeta smoothing has therefore been carried through the original
minima and their full cross blocks. It is not only a density estimate.

\subsection{O10. Evaluate the fixed rational model: the source
part}\label{sa:o10.-evaluate-the-fixed-rational-model-the-source-part}

In \(\mathcal P_N\), use the monic frame \[
1,y,\ldots,y^{d-1},(y+i)^d,y(y+i)^d,\ldots,y^{N-d}(y+i)^d.
\] Its determinant is one. The last block has exactly the Gamma Gram on
\(\mathcal P_{N-d}\). The first \(d\) divided derivatives at \(-i\)
identify its attained quotient; their matrix on the first \(d\)
monomials has determinant one. Its Gram eigenvalues are bounded above by
fixed trial lifts and below by a negative fixed power of \(N+1\). The
latter follows from the polynomial lower comparison
\(\tau_d\succeq C_d(N+1)^{-2d}\sigma\) and the Gamma jet covariance
estimate at the fixed point \(-i\). Therefore \[
\log\det H_N[\tau_d]
=\sum_{j=0}^{N-d}\log\gamma_j+O_d(\log(N+2)).\tag{O28}
\] This includes the fixed jet quotient rather than setting its
determinant to one.

It follows, at the four original source signs, that \[
\boxed{\sum_N\epsilon_N\log\frac{\det H_N[\tau_d]}{\det H_N[\sigma]}
=4dq\log q+(8d\log2-4d)q+o(q).}\tag{O29}
\] To check the coefficients directly, the relative determinant at \(N\)
is \(-\sum_{j=N-d+1}^{N}\log\gamma_j+O(\log N)\), and
\(\log\gamma_j=2j\log j-2j+O(\log(j+2))\). A fixed shift in \(j\) costs
only \(O_d(\log q)\).

\subsection{O11. Evaluate the fixed rational model: the full relation
part}\label{sa:o11.-evaluate-the-fixed-rational-model-the-full-relation-part}

For a relation rank \(n=q\) or \(q+1\), define \[
D_{k,n}(t)=\det\operatorname{Gram}
\{Q,yQ,\ldots,y^{n-1}Q\}_{(1+y^2)^{qt}\sigma},
\] \[
F_{k,n}(t)=q^{-2}
\left[\log D_{k,n}(t)-\log D_{k,n}(0)-2qnt\log q\right].\tag{O30}
\] Andreief's identity on these exact columns proves convexity of
\(F_{k,n}\): its second derivative is a variance of the sum of the
logarithms of \(1+y_j^2\), multiplied by a positive scalar. This is the
whole relation Gram, not a product of independent trial norms.

We verify its leading limit at every fixed \(|t|\le1/8\): \[
\boxed{F_{k,n}(t)\longrightarrow
\frac12\{\mathcal F(2+2t,\pi)-\mathcal F(2,\pi)\}.}\tag{O31}
\] Here \(\mathcal F\) is exactly the original GEL energy, the negative
of the EIQ minimum energy.

The complete root comparison needed for (O31) is as follows. Set
\(Y=2^{-16}q\). Outside \([-Y,Y]\), opposite-root pairing gives \[
\left|\log\frac{|Q(y)|^2(1+y^2)^{qt}}
{|y|^{2q(1+t)}}\right|
\le \frac{qR^2}{Y^2-R^2}+q|t|\log(1+Y^{-2})=O_h(1).
\tag{O32}
\] Inside, use \((y^2+R^2)^q(1+y^2)^{qt}\). For \(t\ge0\) its maximum is
at \(Y\). For \(t<0\), its only possible interior stationary point is a
minimum, as differentiation gives numerator \((1+t)y^2+1+tR^2\); hence
its maximum is at zero or at \(Y\). The zero-endpoint bound \(R^{2q}\)
is negligible relative to \(q^{2q(1+t)}\), uniformly for \(|t|\le1/8\),
because \(R^2/q^{7/4}\to0\). The \(Y\)-endpoint bound is at most
\([(5/4)2^{-28}]^q q^{2q(1+t)}\). The same complete shifted-Legendre
estimate as in the original root comparison, now on \([q,2q]\), shows
that both inner integrals are exponentially small relative to the full
reference-power norm, for every polynomial of degree at most \(q\). Thus
the full quadratic forms are comparable with log width \(O_h(1)\); their
rank-\(n\) log determinants differ by \(O_h(q)\). Their leading
\(q^2\)-scaled limits agree.

For the reference power norm, use the literal parity map \(x=(y/q)^2\).
The two determinant blocks have sizes
\(\lceil n/2\rceil,\lfloor n/2\rfloor\), and starting powers
\(q(1+t),q(1+t)+1\), respectively. Each block has exactly the GEL3
Jacobian and source factor. GEL10 gives the energy at \((2+2t,\pi)\),
since the block size divided by \(q\) tends to \(1/2\). The two
\(q^2\)-coefficients add to \(1/2\). The change of the exact scaling
power is precisely \(2qnt\log q\), the term subtracted in (O30). This
proves (O31). The quantitative appendix Q1--Q3 strengthens it to a
uniform error of \(O_h(\log(q+2)/q)\) on the same fixed parameter
interval. Its proof retains the original Gamma ensemble and uses an
explicitly accounted circle regularization only to bound each finite
configuration.

The EIQ variational calculation proves
\(\partial_\alpha\mathcal F(\alpha,\beta)=\mathcal L(\alpha,\beta)\),
continuously. Thus the right side of (O31) is differentiable at zero,
with derivative \(\mathcal L(2,\pi)\). Convexity now determines the
shrinking secant from \(-d/q\) to zero. To be explicit, enclose that
secant between secants on two fixed small intervals to the left and
right of zero, apply (O31) there, and then shrink the fixed intervals.
This gives \[
\boxed{\log\frac{D_{k,n}(-d/q)}{D_{k,n}(0)}
=-2dn\log q-dq\mathcal L(2,\pi)+O_h(\sqrt{q\log(q+2)}),
\qquad n=q,q+1.}\tag{O33}
\] For the stated rate use the uniform free-energy remainder and a
convex secant of width \(\sqrt{\log(q+2)/q}\), as proved in Q3. No error
term from a \(q^2\)-asymptotic has been differentiated.

The low rank-one relation is retained. Its relative norm is at most one
and, by Jensen and Gamma multiplication applied to \(Q\), at least
\([1+4(q+1)^2]^{-d}\). Its log contribution is therefore
\(O_d(\log q)\). The other low relation is empty. Applying the actual
signs in the source/relation identity gives a total relation
contribution to the baseline of \[
-4dq\log q-2dq\mathcal L(2,\pi)+o_h(q).\tag{O34}
\] Combining this with (O29) proves \[
\mathcal B_k[\tau_d]-\mathcal B_k[\sigma]
=-d[2\mathcal L(2,\pi)-8\log2+4]q+o_h(q).
\] Together with (O27), this proves the value (O2).

The errors in the density and determinant transfer were quantified
before taking their limits. The sublinear patch costs
\(O_h(q^{9/10}\log^2q)\), the averaged zeta error costs
\(O_h(q^{199/200}\log^2q)\), and the residual factor
\(\mathcal L_\zeta(y)/\mathcal L_\zeta(q)\) costs
\(O_h(q\log\log(q+16)/\log(q+2))\). The remaining model error is
\(O_h(\sqrt{q\log(q+2)})\) by Q3. Enlarging the fixed-source constant
absorbs the smaller terms and gives \[
\boxed{\left|\delta_k^{\rm out}+d\mathcal C_\Gamma q\right|
\le C_h\left[
q\frac{\log\log(q+16)}{\log(q+2)}
+q^{199/200}\log^2(q+2)\right]}
\tag{O35}
\] on the cofinal range specified in the proof. Both terms divided by
\(q\) tend to zero. In particular no \(q\log k\), \(q\log\log k\), or
other contribution exceeding order \(q\) remains hidden in this outer
return.

\newpage

\section{Quantitative free-energy remainder used in the outer
value}\label{sa:quantitative-free-energy-remainder-used-in-the-outer-value}

This appendix supplies a uniform remainder for the convex-secant step.
It uses the original EIQ equilibrium and proves the finite-size upper
bound without differentiating an asymptotic error.

\subsection{Q1. A finite-configuration bound from the actual
equilibrium}\label{sa:q1.-a-finite-configuration-bound-from-the-actual-equilibrium}

Let \(V(x)=\beta\sqrt x-\alpha_E\log x\) on \((0,\infty)\), with
\((\alpha_E,\beta)\) in a fixed compact neighborhood of
\([7/4,9/4]\times\{\pi\}\). Let \(\rho\) be its original EIQ probability
density, \(U(x)=\int\log|x-y|\rho(y)dy\), and \(\lambda\) its Euler
constant. Thus \[
2U(x)-V(x)\le-\lambda\quad(x>0),\qquad
\mathcal F(\alpha_E,\beta)=-\lambda-\iint\log|x-y|\rho(x)\rho(y)dxdy.
\tag{Q1}
\] The EIQ endpoint and density formulas show that the supports lie in
one fixed compact subinterval of \((0,\infty)\), their lengths stay
positive, and \(\|\rho\|_\infty\) is uniformly bounded. Their
logarithmic energies and entropies relative to the reference probability
\[
d\omega(x)=\tfrac12x^{-1/2}e^{-\sqrt x}dx
\] are uniformly finite. These assertions follow directly from the
positive density formula with its square-root endpoints and its smooth
parameter dependence.

For a real configuration \(x_1,\ldots,x_n>0\), replace each point by
uniform mass on the complex circle \(x_j+\varepsilon e^{i\theta}\), and
let \(\nu_\varepsilon\) be their average. This is an auxiliary embedding
of the finite configuration into \(\mathbb C=\mathbb R^2\), not a change
in the original ensemble integration. The circle mean identity gives \[
\frac1{2\pi}\int_0^{2\pi}\log|z-\varepsilon e^{i\theta}|d\theta
=\log\max\{|z|,\varepsilon\}.
\tag{Q2}
\] It follows by the harmonic mean-value property outside the circle,
its value at zero inside, and continuity at the circle. Applying (Q2)
twice gives an off-diagonal mean at least \(\log|x_i-x_j|\); the
diagonal self-energy is exactly \(\log\varepsilon\).

The logarithmic energy of a compactly supported real signed measure of
total mass zero on \(\mathbb R^2\) is nonpositive in the \(+\log\)
convention: \[
-\iint\log|z-w|\,d\eta(z)d\eta(w)
=\frac12\int_0^\infty\frac1t\iint e^{-t|z-w|^2}\,d\eta(z)d\eta(w)dt\ge0.
\tag{Q3}
\] The equality follows from the same scalar integral used in EIQ20;
total mass zero cancels its constant. The Gaussian Fourier transform on
\(\mathbb R^2\) makes every inner integral nonnegative. All logarithmic
integrals for the circle measures and the bounded EIQ density are
finite, so the identity follows by truncation and dominated convergence.
This proves the particular two-dimensional comparison used here rather
than assuming a new equilibrium problem in the plane.

Apply (Q3) to \(\nu_\varepsilon-\rho(x)dx\). By (Q2), \[
0\le\int U(z)d\nu_\varepsilon(z)-\frac1n\sum_jU(x_j)
\le2\|\rho\|_\infty\varepsilon.
\] Consequently \[
\frac2{n^2}\sum_{i<j}\log|x_i-x_j|+\frac{\log\varepsilon}{n}
\le\frac2n\sum_jU(x_j)-\iint\log|x-y|\rho(x)\rho(y)dxdy
+4\|\rho\|_\infty\varepsilon.
\] Subtract \(n^{-1}\sum_jV(x_j)\), use (Q1), and take
\(\varepsilon=1/n\). We obtain the uniform pointwise inequality \[
\boxed{2\sum_{i<j}\log|x_i-x_j|-n\sum_jV(x_j)
\le n^2\mathcal F(\alpha_E,\beta)+n\log n+C n.}
\tag{Q4}
\] Repeated configuration points cause no difficulty: their original
left side is \(-\infty\). No original integration region is removed in
this bound.

\subsection{Q2. Apply it to the original Gamma
ensemble}\label{sa:q2.-apply-it-to-the-original-gamma-ensemble}

The exact GEL coordinate map \(y^2=q^2x\) has reference factor \[
B_q(x)=2q e^{\sqrt x}\sigma(q\sqrt x),
\] and ensemble \[
J_{n,q}(n\alpha_E)=\frac1{n!}\int\prod_{i<j}(x_i-x_j)^2
\prod_jx_j^{n\alpha_E}B_q(x_j)d\omega(x_j).
\tag{Q5}
\] All Gamma masses and coordinate powers remain those in GEL3.

The original GEL5 bound is \[
B_q(x)\le Cq^{3/2}\exp[-(\pi q/2-3/2)\sqrt x].
\] Thus (Q4), integrated against the probability \(\omega^n\), proves
the upper estimate with potential coefficient
\(\beta_n=\pi q/(2n)-3/(2n)\). Replacing this coefficient by
\(\pi q/(2n)\) costs \(O(n)\) in the exponent, because the EIQ
derivative \(\partial_\beta\mathcal F=-\int\sqrt x\rho\) is uniformly
bounded here. The factor \((Cq^{3/2})^n/n!\) costs at most
\(O(n\log(q+n))\).

For the lower estimate, restrict to the actual equilibrium support for
\((\alpha_E,\pi q/(2n))\), use GEL6, and change the integration measure
to \(\rho^n\). Jensen's inequality gives the equilibrium energy minus
the entropy and the missing diagonal term. The uniform bounds in Q1 show
that these costs are \(O(n)\), apart from the explicit Gamma prefactor
and \(n!\), whose costs are \(O(n\log(q+n))\).

Therefore, uniformly when \(q/n\) remains in a fixed neighborhood of two
and \(\alpha_E\) remains in the stated compact interval, \[
\boxed{\left|\log J_{n,q}(n\alpha_E)
-n^2\mathcal F\!\left(\alpha_E,\frac{\pi q}{2n}\right)\right|
\le C n\log(q+n).}
\tag{Q6}
\] This gives an explicit order for the remainder in the already
established free-energy limit, on exactly the parameters required by the
rational-model comparison.

\subsection{Q3. Quantify the shrinking secant on the full original
relation
Gram}\label{sa:q3.-quantify-the-shrinking-secant-on-the-full-original-relation-gram}

Use the convex function \(F_{k,n}(t)\) in (O30), for relation rank
\(n=q\) or \(q+1\). The full paired-root comparison (O32) has
logarithmic quadratic-form width \(O_h(1)\). Its rank-\(n\) determinant
cost is \(O_h(q)\). Apply (Q6) to the two literal parity blocks. Their
sizes are \(q/2+O(1)\), and their potential parameters differ from
\((2+2t,\pi)\) by \(O(1/q)\). EIQ parameter smoothness and the exact
scaling powers in (O30) therefore give, uniformly for \(|t|\le1/8\), \[
\boxed{\left|F_{k,n}(t)-F_*(t)\right|\le C_h\frac{\log(q+2)}q,
\qquad
F_*(t)=\tfrac12[\mathcal F(2+2t,\pi)-\mathcal F(2,\pi)].}
\tag{Q7}
\] The derivative \(F_*'(0)=\mathcal L(2,\pi)\), and \(F_*''\) is
bounded on this fixed interval by the original smooth endpoint formulas.

Set \(\varepsilon_q=\sqrt{\log(q+2)/q}\). On its cofinal range
\(d/q<\varepsilon_q/2\), convexity places the secant on \([-d/q,0]\)
between secants on fixed-width adjacent intervals contained in
\([-2\varepsilon_q,2\varepsilon_q]\). The error from (Q7), divided by
their length, is \(O_h(\log(q+2)/(q\varepsilon_q))\). The difference of
the corresponding \(F_*\) secants from \(F_*'(0)\) is
\(O_h(\varepsilon_q)\). Thus \[
\frac{F_{k,n}(0)-F_{k,n}(-d/q)}{d/q}
=\mathcal L(2,\pi)+O_h\!\left(\sqrt{\frac{\log(q+2)}q}\right).
\] Return the exact powers from (O30): \[
\boxed{\log\frac{D_{k,n}(-d/q)}{D_{k,n}(0)}
=-2dn\log q-dq\mathcal L(2,\pi)
+O_h(\sqrt{q\log(q+2)}).}
\tag{Q8}
\] This estimates a secant of the complete attained relation
determinant. It does not differentiate an error term and does not assign
independently favorable values at its endpoints.

\newpage

\section{The central endpoint contribution: evaluation of its
coefficient}\label{sa:the-central-endpoint-contribution-evaluation-of-its-coefficient}

This part evaluates the central Gamma quantity left by the consolidated
localization. It does not replace the arithmetic source-band kernel by
an assumed constant. Instead it proves a uniform monic-norm estimate for
an explicitly smoothed convex family and evaluates its Robin variation.
The smoothing is returned through fixed Taylor-jet quotients with an
error smaller than \(k^2\).

\subsection{C1. The quantity inherited from the exact
localization}\label{sa:c1.-the-quantity-inherited-from-the-exact-localization}

Retain \[
M(\theta)=\int e^{\theta y}w_h(y)dy,\quad L=\log M,
\quad \mu=L',\quad V=L'',\quad \alpha=\pi/2,
\quad b=\mu(\alpha)>0.
\] The inverse of \(\mu\) is denoted by \(\theta(t)\). Put \[
D(t)=L(\alpha)-L(\theta(t))-(\alpha-\theta(t))t,
\quad0\le t\le b,
\] \[
W_k(y)=kD(|y|/k)1_{|y|\le kb},\qquad
I_h=\int_0^\alpha\mu(\theta)^2d\theta
=\frac1{k^2}\int_{\mathbb R}W_k(y)dy>0.
\tag{C1}
\] The equality of the two expressions for \(I_h\) was derived in the
incoming central calculation. The whole function \(W_k\) is nonnegative,
Lipschitz with constant \(\alpha\), and bounded above by
\(\Lambda_k=k[L(\alpha)-L(0)]\).

For a positive weight \(w\) and a nonnegative insertion \(W\), write \[
D_N[w,W]=\log\det H_N[w]-\log\det H_N[e^{-W}w],
\] \[
\Phi_{q,k}[w,W]=D_{2q-1}[w,W]+D_{2q}[w,W]
-D_{q-1}[w,W]-D_q[w,W].\tag{C2}
\] The established exact-source localization gives \[
\delta_k^{\rm cen,end}
=\Phi_{q,k}[\sigma,W_k]+O_h(k\log^2q).
\tag{C3}
\] Its original relation determinants and all four source signs are part
of that result. We now prove \[
\boxed{\Phi_{q,k}[\sigma,W_k]
=\frac{\log2}{\pi}I_h k^2+o_h(k^2).}\tag{C4}
\] This supplies a value for the central contribution on the simple
stratum, not only an \(O(q)\) enclosure.

\subsection{C2. An explicit analytic regularization with a controlled
determinant
cost}\label{sa:c2.-an-explicit-analytic-regularization-with-a-controlled-determinant-cost}

Let \[
a=a_k=\frac{k}{[\log(q+2)]^4},\qquad
G_a(u)=\frac{e^{-u^2/a^2}}{\sqrt\pi a}.
\] We work on the cofinal range \(a\ge1\). Define the actual comparison
potentials \[
Q_{k,t}(y)=\alpha|y|+tW_k(y),\qquad
\widetilde Q_{k,t}=G_a*Q_{k,t},\qquad0\le t\le1,
\] \[
\widetilde W_k=G_a*W_k,
\qquad \widetilde Q_{k,t}=\widetilde Q_{k,0}+t\widetilde W_k.
\tag{C5}
\] The convolution here is an auxiliary comparison of real potentials,
not a replacement of the original arithmetic class. Its return into (C2)
is proved next.

The exact scalar properties are \[
0\le\widetilde Q_{k,0}(y)-\alpha|y|\le\alpha a/\sqrt\pi,
\qquad
\|\widetilde W_k-W_k\|_\infty\le\alpha a/\sqrt\pi,
\tag{C6}
\] \[
\int\widetilde W_k=\int W_k=k^2I_h,
\qquad
\int y^2\widetilde W_k(y)dy=O_h(k^4).
\tag{C7}
\] The first inequalities follow respectively from convexity of absolute
value and from the exact Lipschitz constant of \(W_k\), since
\(\int|u|G_a(u)du=a/\sqrt\pi\). Fubini proves the area identity. The
second moment is the original compactly supported moment plus
\((a^2/2)k^2I_h\). These are moments of the auxiliary rate function, not
nonexistent fourth moments of the boundary-tilted arithmetic density.

We claim the following determinant comparison: \[
\boxed{\Phi_{q,k}[\sigma,W_k]
=\Phi_{q,k}[e^{-\widetilde Q_{k,0}},\widetilde W_k]
+O_h\!\left(k\log q\,[a+\log q]\right).}\tag{C8}
\] Its error is \(o(k^2)\), by the specified choice of \(a\).

Here are the fixed-fibre details. Put
\(g=e^{-\widetilde Q_{k,0}}/\sigma\). The original Gamma bounds and (C6)
give \[
c_he^{-C_ha}\le g(y)\le C_h(1+y^2)^{1/4}.
\] On every polynomial of degree at most \(2q\), both before and after
multiplication of the weight by \(e^{-tW_k}\), these imply \[
c_he^{-C_ha}\|P\|_{e^{-tW_k}\sigma}^2
\le\|P\|_{g e^{-tW_k}\sigma}^2
\le C_h(q+1)^{1/2}\|P\|_{e^{-tW_k}\sigma}^2.
\tag{C9}
\] For the upper bound use Jensen and the consolidated full-tail
second-moment bound \[
\int y^2|P|^2e^{-tW_k}d\sigma\le272(2q+1)^2\|P\|_{e^{-tW_k}\sigma}^2.
\] Thus the logarithmic width in (C9) is \(O_h(a+\log q)\), at both
endpoints of the nonlinear insertion.

Let \(K'=kb+10a\sqrt{\log q}\). Outside this interval the Gaussian tail
gives \(\widetilde W_k\le\Lambda_k q^{-100}\). Restrict the first \(r\)
Taylor coefficients at zero to each source space, with \[
r=\left\lceil C_h\{(K'+1)\log(q+2)+\Lambda_k+a+\log(q+2)\}\right\rceil
=O_h(k\log q).
\tag{C10}
\] The constant can be chosen so that the Gamma central restriction
fraction on its full kernel is at most \(e^{-3\Lambda_k-4C_ha}q^{-50}\),
by the explicit generating-function estimate (O23). The actual
Taylor-image quotient has dimension at most \(r\).

First compare the loss \(D_N[g\sigma,W_k]\) with \(D_N[\sigma,W_k]\). On
the Taylor kernel, both insertion losses are \(O(q^{-20})\): the
unsuppressed central norm fraction is bounded by the preceding estimate
times \(e^{C_ha}\) and a fixed power of \(q\). On the two attained
Taylor-image quotients, (C9) bounds the difference of the two endpoint
losses by \(r\,O_h(a+\log q)\). This is a comparison of two complete
Schur losses, not a bound on two unrelated trial metrics.

Second compare the insertions \(W_k\) and \(\widetilde W_k\) on the same
background \(g\sigma\). Their pointwise relative factor is between
\(e^{-C_ha}\) and \(e^{C_ha}\). On the Taylor kernel its deviation from
one costs at most \(e^{\Lambda_k+2C_ha}\) times the chosen central
restriction fraction, plus \(O(\Lambda_kq^{-100})\) outside. Hence its
kernel log-determinant cost is again negligible. Its attained quotient
costs at most \(rC_ha\). The fixed frame determinant cancels in each
comparison. Sum the four original endpoint signs. This proves (C8), with
every preceding source direction included in the Schur minimum.

\subsection{C3. The regularized potentials have an exact positive
curvature
measure}\label{sa:c3.-the-regularized-potentials-have-an-exact-positive-curvature-measure}

Define the even finite positive measure \[
d\kappa_{k,t}(u)=2\alpha(1-t)\delta_0(du)
+\frac{t}{kV(\theta(|u|/k))}1_{|u|<kb}\,du.
\tag{C11}
\] Its mass is exactly \(2\alpha\): on the positive half-line the
substitution \(s=\mu(\theta)\) gives \(ds=V(\theta)d\theta\). Its first
moment is zero. Moreover \[
\boxed{Q_{k,t}(y)=\frac12\int|y-u|\,d\kappa_{k,t}(u).}\tag{C12}
\] To verify the constant as well as the derivative, both sides equal
\(\alpha|y|\) for \(|y|\ge kb\). On \(0<y<kb\), both derivatives are
\((1-t)\alpha+t\theta(y/k)\). Evenness gives the other half-line.

Let \(\kappa^a_{k,t}=G_a*\kappa_{k,t}\). Then \[
\widetilde Q_{k,t}''(y)=\kappa^a_{k,t}(y)>0,
\qquad \int\kappa^a_{k,t}=2\alpha,
\] \[
\int u^2\kappa^a_{k,t}(u)du=2tk^2I_h+\alpha a^2=O_h(k^2).
\tag{C13}
\] The original unsmoothed curvature has compact support \([-kb,kb]\).
These moment assertions therefore do not require any additional tilted
arithmetic moments.

The analytic continuation is explicit: \[
\widetilde Q_{k,t}(z)=\frac12\int
\left[(z-u)\operatorname{erf}\!\left(\frac{z-u}{a}\right)
+\frac a{\sqrt\pi}e^{-(z-u)^2/a^2}\right]d\kappa_{k,t}(u).
\tag{C14}
\] It is entire, real on the real axis, even, and strictly convex there.
Its derivative tends to \(\pm\alpha\) at the two real infinities.

\subsection{C4. Construct the equilibrium at the actual monic
degree}\label{sa:c4.-construct-the-equilibrium-at-the-actual-monic-degree}

Fix an integer \(n\) with \(q\le n\le2q\), and a parameter
\(t\in[0,1]\). There is exactly one \(A=A_{n,t}>0\) satisfying \[
\boxed{n=\frac1{2\pi}\int_{-A}^{A}
\sqrt{A^2-u^2}\,\kappa^a_{k,t}(u)du.}\tag{C15}
\] The right side is continuous, strictly increasing from zero to
infinity; its derivative is the positive integral with numerator \(A\).
This proves existence, uniqueness, and smooth dependence on \(t\).

The same equation is \[
\int_{-A}^A\frac{\widetilde Q_{k,t}'(s)}{\sqrt{A^2-s^2}}ds=0,
\qquad
\int_{-A}^A\frac{s\widetilde Q_{k,t}'(s)}{\sqrt{A^2-s^2}}ds=2\pi n.
\tag{C16}
\] The first identity is parity. Integration by parts gives the second
directly from (C15).

Write \(Q_t=\widetilde Q_{k,t}\) temporarily and define \[
h_{n,t}(z)=\frac1\pi\int_{-A}^A
\frac{Q_t'(s)-Q_t'(z)}{s-z}\frac{ds}{\sqrt{A^2-s^2}},
\] \[
\rho_{n,t}(y)=\frac{\sqrt{A^2-y^2}}{2\pi}h_{n,t}(y),
\qquad |y|<A.
\tag{C17}
\] The divided difference is removable at \(s=z\). Thus \(h_{n,t}\) is
analytic, and strict convexity makes it positive on the real axis. The
density vanishes as a square root at each endpoint.

For a direct verification of its mass and Euler equation, let
\(R_A(z)=\sqrt{z^2-A^2}\), with \(R_A(z)/z\to1\). The function \[
G(z)=\frac{R_A(z)}{2\pi}\int_{-A}^A
\frac{Q_t'(s)}{(z-s)\sqrt{A^2-s^2}}ds
\] has \(G(z)=n/z+O(z^{-2})\) by (C16), and jump
\(G_+-G_-=-2\pi i\rho_{n,t}\). Hence it is the Cauchy transform of this
density. Its real boundary value is \(Q_t'/2\), proving \[
\boxed{Q_t(y)-2\int\log|y-s|\rho_{n,t}(s)ds
=\lambda_{n,t},\qquad |y|\le A.}\tag{C18}
\] Outside, the derivative of the left side is \(R_A(y)h_{n,t}(y)\),
negative to the left and positive to the right. The full exterior
inequality follows. Positivity of the logarithmic energy of a zero-mass
difference, as proved in EIQ20--23, then proves that this is the unique
mass-\(n\) equilibrium, not just a candidate supported on one interval.

Its radius has the uniform estimate \[
\boxed{A_{n,t}=\frac{\pi n}{\alpha}+O_h(k^2/n)=2n+O_h(k^2/n).}\tag{C19}
\] Indeed, if \(m_2=\int u^2\kappa^a_{k,t}(u)du\), then (C15) and
\(A-\sqrt{A^2-u^2}\le u^2/A\) on \(|u|\le A\), together with
\(A\le u^2/A\) outside, give \[
0\le\frac{\alpha A}{\pi}-n\le\frac{m_2}{2\pi A}.
\] Use (C13). In particular \(kb/A\to0\) and \(a/A\to0\), uniformly
throughout the stated degree and parameter ranges.

\subsection{C5. Uniformity of the monic-norm estimate for this moving
family}\label{sa:c5.-uniformity-of-the-monic-norm-estimate-for-this-moving-family}

Let \(H_n(t)\) be the squared norm of the monic degree-\(n\) polynomial
in the weight \(e^{-Q_t(y)}dy\). We prove \[
\boxed{H_n(t)=\pi A_{n,t}e^{-\lambda_{n,t}}
\left[1+O_h(n^{-1})\right],
\qquad q\le n\le2q,\quad0\le t\le1.}\tag{C20}
\] The uniformity in this statement is essential. We verify it below for
the explicit family (C14); a theorem for one fixed analytic potential is
not simply applied to a changing potential.

\subsubsection{C5.1. Positivity and analytic control in the narrow
central
strip}\label{sa:c5.1.-positivity-and-analytic-control-in-the-narrow-central-strip}

For \(|y|,|u|<A\), put \[
\mathcal G_A(y,u)=
\log\left|\frac{A^2-yu+\sqrt{A^2-y^2}\sqrt{A^2-u^2}}
{A(y-u)}\right|.
\] Integration of the elementary arcsine resolvent in (C17) gives \[
\rho_{n,t}(y)=\frac1{2\pi^2}\int_{-A}^A
\mathcal G_A(y,u)\kappa^a_{k,t}(u)du.\tag{C21}
\] For example, set \(y=A\cos\theta\), \(u=A\cos\phi\); the kernel is
\(\log|\sin((\theta+\phi)/2)/\sin((\theta-\phi)/2)|\). Its sine
expansion is \(2\sum_{j\ge1}\sin(j\theta)\sin(j\phi)/j\), and
integration against \(A\sin\theta\) gives \(\pi\sqrt{A^2-u^2}\),
confirming the mass in (C15).

There are constants \(c,C>0\), independent of \(n,k,t\) on the cofinal
fixed-packet range, such that \[
\rho_{n,t}(y)\ge c\quad(|y|\le A/3),
\qquad
|\rho_{n,t}'(z)|\le C/a+C/A
\quad(|\Re z|\le A/3,\ |\Im z|\le a/10).
\tag{C22}
\] For the real lower bound, pair in (C17) the intervals \([A/2,3A/4]\)
and their negatives. There \(Q_t'\) is exponentially close to \(\alpha\)
and \(-\alpha\), respectively, uniformly by (C14), while
\(|Q_t'(y)|\le\alpha\). Their positive divided differences supply
\(h_{n,t}(y)\ge c/A\), hence the first assertion.

For the analytic derivative bound, decompose the kernel in (C21) into
\(-\log|y-u|\) and the logarithm of
\((A^2-yu+\sqrt{A^2-y^2}\sqrt{A^2-u^2})/A\). The latter has derivative
bounded by \(C/A\) in the indicated central region, uniformly for
\(|u|\le A\). The missing tails of \(\kappa^a\) have exponentially small
mass and are analytic there. The logarithmic convolution with the
Gaussian has an entire real-analytic continuation \(J_a\), whose
derivative is \[
J_a'(z)=\int_0^\infty e^{-a^2v^2/4}\sin(zv)\,dv.
\] This follows either by differentiating the Gaussian logarithmic
convolution or by solving \(J_a''+2zJ_a'/a^2=2/a^2\), \(J_a'(0)=0\). On
\(|\Im z|\le a/10\), its absolute value is at most \(C/a\). Integrating
against the unsmoothed curvature measure of exact mass \(2\alpha\)
proves the second assertion in (C22).

Choose a fixed sufficiently small \(\varepsilon_0>0\). Equation (C22)
then proves \[
\Re\rho_{n,t}(x+iv)\ge c/2
\quad(|x|\le A/3,\ |v|\le\varepsilon_0a).\tag{C23}
\] Thus the central strip has an explicitly growing physical height
proportional to \(a\), not a fixed scaled width that the varying family
might fail to possess.

\subsubsection{C5.2. Uniform soft edges and the rest of the
lens}\label{sa:c5.2.-uniform-soft-edges-and-the-rest-of-the-lens}

After scaling \(y=Az\), the equilibrium probability density is
\((A/n)\rho_{n,t}(Az)\). Near either scaled endpoint \(\pm1\), its
analytic square-root coefficient converges, with three derivatives on
fixed slightly larger disks, to the coefficient for \(\pi|z|\). The
error is \[
O_h((k^2+a^2)/n^2)+O_h(e^{-c_hn^2/a^2}).\tag{C24}
\] Here is the estimate behind this assertion. On \(\Re y\ge A/2\),
\(|\Im y|\le A/20\), (C14) gives \(Q_t'(y)=\alpha+O(e^{-c_h A^2/a^2})\),
with the corresponding differentiated estimates; the negative side is
identical. In the integral (C17), split the real variable at
\(|s|=A/2\). On the outer pieces the divided differences of this error
are exponentially small, including the removable point. On the middle
piece the denominator stays a fixed distance from the endpoint disks.
The difference from \(\alpha\operatorname{sgn}s\) is odd, so its zeroth
moment cancels. Its first absolute moment is \(O_h(k^2+a^2)\), by
(C12)--(C14). Taylor expansion of the remaining denominator therefore
gives an error \(O_h((k^2+a^2)/A^3)\) in the physical \(h\), and the
same estimate for fixed scaled derivatives. Multiplication by \(A^2/n\),
and (C19), prove (C24).

The limiting scaled density is \(\pi^{-1}\operatorname{arcosh}(1/|z|)\)
on \((-1,1)\). Its analytic edge coefficient in
\((2\pi)^{-1}\sqrt{1-z^2}\,\widehat h(z)\) has value
\(\widehat h(\pm1)=2\). Thus fixed endpoint disks have uniformly nonzero
edge coefficients and uniformly bounded analytic derivatives. The same
integral argument on fixed scaled compact sets away from zero proves
uniform analytic convergence there. In particular the positivity needed
for the rest of the lens persists in a fixed sufficiently small scaled
neighborhood.

Choose the upper and lower lens boundaries at physical height
\(\varepsilon_0a\) over the middle third. In the overlap with the region
away from zero, let them rise along fixed-slope segments and enter the
fixed scaled endpoint disks transversally. This uses a bounded number of
line segments and smooth arcs, with uniformly controlled angles. By
(C23), and the uniform analytic convergence away from zero, the
off-diagonal lens jumps after the equilibrium transformation are
\(O(e^{-c_ha})\) outside the endpoint disks. The full exterior Euler
inequality gives exponentially small real jumps outside those disks. At
large real distance this is integrable uniformly: \(Q_t'(y)\ge\alpha/2\)
and the equilibrium resolvent is at most \(n/(y-A)\).

\subsubsection{C5.3. The Riemann--Hilbert calculation and its error
coefficient}\label{sa:c5.3.-the-riemannhilbert-calculation-and-its-error-coefficient}

Use the exact orthogonal-polynomial matrix problem of
Fokas--Its--Kitaev: its jump is
\(\left(\begin{smallmatrix}1&e^{-Q_t}\\0&1\end{smallmatrix}\right)\),
its first entry is the monic polynomial, and the coefficient of
\(z^{-n-1}\) in its Cauchy-transform entry is \(-H_n(t)/(2\pi i)\).
These statements follow directly from the Cauchy expansion and the
defining orthogonality. They are also the normalization used in {[}D,
Theorem 31{]}.

Scale by the actual \(A\), and let \(g\) be the logarithmic transform of
the equilibrium probability. Its Euler constant in the scaled convention
is \[
\ell_{n,t}=-\lambda_{n,t}/n-2\log A.\tag{C25}
\] The equilibrium transformation and lens factorization produce the
constant interval jump
\(\left(\begin{smallmatrix}0&1\\-1&0\end{smallmatrix}\right)\). Its
exact solution is \[
N(z)=\frac12\begin{pmatrix}
\beta+\beta^{-1}&(\beta-\beta^{-1})/i\\
-(\beta-\beta^{-1})/i&\beta+\beta^{-1}
\end{pmatrix},\qquad
\beta(z)=\left(\frac{z-1}{z+1}\right)^{1/4},
\] with \(\beta\to1\) at infinity. In particular
\(N_{12}(z)=i/(2z)+O(z^{-2})\).

In each fixed endpoint disk the analytic phase has a simple soft edge:
its \(3/2\)-power coefficient is bounded above and away from zero by
(C24). The Airy fundamental matrix, with its exact constant Stokes
jumps, gives the local solution. Multiplying its leading analytic factor
by \(N\) matches the exterior solution. The Airy asymptotic, uniform on
the fixed disk boundary, gives the matching error \(I+O_h(1/n)\). The
coefficient bound is uniform because the phase and its inverse conformal
coordinate have the bounds just proved. This is precisely the standard
Airy construction in {[}D, Section 3{]}; the new moving-family bounds
needed to use it are (C22)--(C24).

For completeness the small-norm step remains valid despite the narrow
middle of the lens. On the scaled contour its Cauchy boundary operator
has norm at most \(C_hA/a\). A direct coarse proof splits it into the
finitely many segments, arcs, and rays. Self terms have the ordinary
line or circle Cauchy bound. Intersecting pieces have fixed transverse
angles and the usual Carleman-kernel bound (obtainable by the weight
\(s^{-1/2}\)). The other cross kernels have separation at least a fixed
multiple of \(a/A\), giving the displayed bound. No uniform separation
independent of \(k\) is assumed.

Write the final jump as \(I+w\). Its \(L^1,L^2,L^\infty\) norms are
bounded by \[
C_h/n+C_he^{-c_ha}.
\] Therefore \[
\|C_- (\,\cdot\,w)\|
\le\frac{C_hA}{a}\left(n^{-1}+e^{-c_ha}\right)=o(1).
\tag{C26}
\] The singular-integral equation is solved by its convergent Neumann
series. Its solution \(\nu\) obeys
\(\|\nu-I\|_2=O_h(a^{-1})+O_h((n/a)e^{-c_ha})\). The coefficient of
\(z^{-1}\) in the final error matrix is \(-\frac1{2\pi i}\int\nu w\),
hence is \(O_h(1/n)\): use the \(L^1\) bound for \(w\) and
Cauchy--Schwarz for \((\nu-I)w\). The term \(e^{-c_ha}\) is smaller than
every inverse power of \(q\), because \(a=k/\log^4(q+2)\).

Undoing the exact transformations now gives \[
H_n(t)=\pi A^{2n+1}e^{n\ell_{n,t}}[1+O_h(1/n)]
=\pi A e^{-\lambda_{n,t}}[1+O_h(1/n)],
\] by (C25). This proves (C20), including its normalization and the
required uniformity in \(n,k,t\). Appendix A1--A3 writes the exact
matrix transformations, all four Airy jumps, the analytic matching
factor, and the finite-piece Cauchy-operator bound explicitly. No
fixed-potential remainder is differentiated or presumed uniform without
these estimates.

\subsection{C6. The Robin variation has an exact arithmetic-source
integral}\label{sa:c6.-the-robin-variation-has-an-exact-arithmetic-source-integral}

Differentiate the Euler equality (C18) at a fixed interior point, then
integrate against the arcsine probability on the current interval
\([-A,A]\). The derivative of the equilibrium density has total mass
zero; its endpoint contributions vanish because the density has
square-root decay. The arcsine logarithmic potential is exactly
\(\log(A/2)\) on this whole interval. Therefore its pairing with that
zero-mass derivative vanishes, giving the exact identity \[
\boxed{\partial_t\lambda_{n,t}
=\frac1\pi\int_{-A_{n,t}}^{A_{n,t}}
\frac{\widetilde W_k(y)}{\sqrt{A_{n,t}^2-y^2}}dy.}\tag{C27}
\] The parameter derivatives exist from (C15)--(C17). This
differentiates an exact finite equilibrium equation, not an asymptotic
error.

By (C7), (C19), and the Gaussian tail outside \([-kb,kb]\), \[
\partial_t\lambda_{n,t}
=\frac{k^2I_h}{2\pi n}+O_h(k^4/n^3)+O_h(q^{-20}),
\tag{C28}
\] uniformly in the same range. To verify the error, on \(|y|\le A/2\)
expand \((A^2-y^2)^{-1/2}=A^{-1}+O(y^2/A^3)\) and use the second moment
in (C7). Beyond this interval the Gaussian convolution of the compact
rate is exponentially small, including the integrable arcsine endpoint
singularity. Finally \(A^{-1}=(2n)^{-1}+O_h(k^2/n^3)\) by (C19).

Also \(\log(A_{n,1}/A_{n,0})=O_h(k^2/n^2)\). Integrate (C28) over \(t\),
and use (C20) only at its two endpoints. It follows that the actual
monic norm loss is \[
\boxed{\log\frac{H_n(0)}{H_n(1)}
=\frac{k^2I_h}{2\pi n}
+O_h\left(\frac{k^4}{n^3}+\frac{k^2}{n^2}+\frac1n\right).}\tag{C29}
\] This is the missing value of the central high-degree norm response,
with the coupling to all preceding polynomials already included in
\(H_n(t)\).

\subsection{C7. Sum the original bands, with the endpoint weights
intact}\label{sa:c7.-sum-the-original-bands-with-the-endpoint-weights-intact}

A source determinant is the product of its successive monic norms. Thus
the four signs in (C2) give exactly \[
\Phi_{q,k}[e^{-\widetilde Q_{k,0}},\widetilde W_k]
=\sum_{r=0}^q c_r\log\frac{H_{q+r}(0)}{H_{q+r}(1)},
\quad c_0=c_q=1,\quad c_r=2\ (0<r<q).
\tag{C30}
\] No determinant has been replaced by an arithmetic mean. Substitute
(C29): \[
\Phi_{q,k}[e^{-\widetilde Q_{k,0}},\widetilde W_k]
=\frac{k^2I_h}{2\pi}\sum_{r=0}^q\frac{c_r}{q+r}
+O_h(1+k^2/q+k^4/q^2).
\] The displayed sum is the exact trapezoidal sum for
\(2\int_1^2x^{-1}dx\), with error \(O(q^{-1})\). Hence \[
\Phi_{q,k}[e^{-\widetilde Q_{k,0}},\widetilde W_k]
=\frac{\log2}{\pi}I_h k^2+O_h(1+k^2/q+k^4/q^2).
\tag{C31}
\] Return the smoothing by (C8), and then the original arithmetic
central endpoint by (C3). The completed central result is \[
\boxed{\delta_k^{\rm cen,end}
=\frac{\log2}{\pi}I_h k^2+
O_h\left(\frac{k^2}{\log^3(q+2)}+k\log^2(q+2)+1+\frac{k^4}{q^2}\right).}\tag{C32}
\] The error is \(o_h(k^2)\), and in particular \(o_h(q)\), on every
fixed multiplicity stratum. The positive coefficient in (C32), including
its factor \(\log2/\pi\), is now evaluated. It is not the unweighted
negative area coefficient from the preceding source-only calculation.

\newpage

\section{Explicit matrix construction for the uniform central monic
norm}\label{sa:explicit-matrix-construction-for-the-uniform-central-monic-norm}

This appendix expands the matrix step in C5. The object is the already
constructed even analytic potential \(Q_t=\widetilde Q_{k,t}\), at an
actual degree \(q\le n\le2q\). Its support radius is \(A=A_{n,t}\), and
its physical Robin constant is \(\lambda=\lambda_{n,t}\).

\subsection{A1. Exact scaling and equilibrium
transformation}\label{sa:a1.-exact-scaling-and-equilibrium-transformation}

Set \(V_n(z)=Q_t(Az)/n\). The scaled equilibrium probability is \[
\varrho(x)=\frac A n\rho_{n,t}(Ax),\qquad -1<x<1,
\] and \[
g(z)=\int_{-1}^1\log(z-s)\varrho(s)ds,\qquad
\ell=-\lambda/n-2\log A.
\] Its boundary equation is \(g_++g_--V_n=\ell\). Write
\(\widehat h(z)=A^2h_{n,t}(Az)/n\), so that
\(\varrho(x)=(2\pi)^{-1}\sqrt{1-x^2}\widehat h(x)\).

The exact orthogonal-polynomial matrix for \(e^{-nV_n(x)}dx\) has jump
\(Y_+=Y_-\left(\begin{smallmatrix}1&e^{-nV_n}\\0&1\end{smallmatrix}\right)\),
and \[
Y(z)z^{-n\sigma_3}=I+Y_1/z+O(z^{-2}),\qquad
(Y_1)_{12}=-\frac{A^{-2n-1}H_n(t)}{2\pi i}.
\tag{A1}
\] The second equality follows from expanding the Cauchy transform: all
lower moments vanish by the defining monic orthogonality, and its first
nonzero moment is the squared monic norm.

Define \[
T(z)=e^{-n\ell\sigma_3/2}Y(z)e^{-n(g(z)-\ell/2)\sigma_3}.
\tag{A2}
\] Then \(T\to I\). Near the right edge put \[
\phi(z)=\frac12\int_1^z\widehat h(s)\sqrt{s^2-1}\,ds.
\tag{A3}
\] On \(z>1\), this is one half of \(V_n-2g+\ell\), and is positive. On
the upper side of \((-1,1)\), \[
\phi_+(x)=-i\pi\int_x^1\varrho(s)ds,
\qquad g_+-g_-=-2\phi_+.
\] Thus the interval jump of \(T\) is \[
\begin{pmatrix}e^{2n\phi_+}&1\\0&e^{-2n\phi_+}\end{pmatrix}
=
\begin{pmatrix}1&0\\e^{-2n\phi_+}&1\end{pmatrix}
J_0
\begin{pmatrix}1&0\\e^{2n\phi_+}&1\end{pmatrix},
\quad J_0=\begin{pmatrix}0&1\\-1&0\end{pmatrix}.
\tag{A4}
\] Inside the upper lens multiply \(T\) on the right by
\(\left(\begin{smallmatrix}1&0\\-e^{2n\phi}&1\end{smallmatrix}\right)\);
inside the lower lens multiply by the same matrix with the minus sign
replaced by plus. The resulting matrix \(S\) has the constant jump
\(J_0\) on the interval, lower-triangular jumps with entry
\(e^{2n\phi}\) on the two lens boundaries, and upper-triangular entry
\(e^{-2n\phi}\) on the exterior real rays. The phases on the left are
their even-potential continuations; integer \(n\) removes the full-mass
\(2\pi i\) branch increment.

The explicit global solution of the remaining constant interval jump is
\[
N(z)=\frac12\begin{pmatrix}
\beta+\beta^{-1}&(\beta-\beta^{-1})/i\\
-(\beta-\beta^{-1})/i&\beta+\beta^{-1}
\end{pmatrix},\qquad
\beta(z)=\left(\frac{z-1}{z+1}\right)^{1/4},\quad\beta(\infty)=1.
\tag{A5}
\] In particular \((N_1)_{12}=i/2\). All of these matrices act on the
original two-by-two polynomial/\allowbreak Cauchy-transform problem; they do not
alter the programme's arithmetic quotient.

\subsection{A2. Airy matrix and all four constant
jumps}\label{sa:a2.-airy-matrix-and-all-four-constant-jumps}

Let \(\omega=e^{2\pi i/3}\), and define entire column vectors \[
F(\zeta)=\binom{\mathrm{Ai}(\zeta)}{\mathrm{Ai}'(\zeta)},\qquad
G(\zeta)=-\omega^2\binom{\mathrm{Ai}(\omega^2\zeta)}{\omega^2\mathrm{Ai}'(\omega^2\zeta)}.
\] With \(c_A=\sqrt{2\pi}e^{-\pi i/4}\), define the four sectors \[
\Psi(\zeta)=c_A\begin{cases}
[F,G],&0<\arg\zeta<2\pi/3,\\
[F-G,G],&2\pi/3<\arg\zeta<\pi,\\
[G,G-F],&-\pi<\arg\zeta<-2\pi/3,\\
[F,G-F],&-2\pi/3<\arg\zeta<0.
\end{cases}
\tag{A6}
\] Orient the positive ray outwards and the other three rays inwards.
Direct column multiplication gives the jumps \[
\begin{array}{c|c}
\arg\zeta=0&\begin{pmatrix}1&1\\0&1\end{pmatrix}\\[3pt]
\arg\zeta=\pm2\pi/3&\begin{pmatrix}1&0\\1&1\end{pmatrix}\\[3pt]
\arg\zeta=\pi&J_0.
\end{array}
\tag{A7}
\] The exact Airy connection identity
\(\mathrm{Ai}(\zeta)+\omega\mathrm{Ai}(\omega\zeta)+\omega^2\mathrm{Ai}(\omega^2\zeta)=0\)
and its derivative verify that all sector columns are the appropriate
rotated Airy solutions. Their Wronskian gives \(\det\Psi=1\). These are
DLMF 9.2.8--9.2.12 in the displayed convention.

Their sectorial asymptotic is \[
\Psi(\zeta)=\zeta^{-\sigma_3/4}C_A
[I+O(\zeta^{-3/2})]e^{-(2/3)\zeta^{3/2}\sigma_3},
\qquad
C_A=\frac{e^{-\pi i/4}}{\sqrt2}
\begin{pmatrix}1&i\\-1&i\end{pmatrix}.
\tag{A8}
\] For example, in the upper right sector the two columns have leading
vectors \((1,-1)^{\mathsf T}\) and \((i,i)^{\mathsf T}\), after the
common factors are extracted. This verifies the constant in (A8), rather
than leaving it as an unspecified normalization. The other sectors
follow from the just-written connection identities. DLMF 9.7 supplies
the uniform Airy remainder in the closed subsectors used for the disk
boundaries.

Near \(z=1\), set \[
f(z)=\left(\frac32\phi(z)\right)^{2/3}.
\] It is analytic with a simple zero at one. Since
\(\widehat h(1)\to2\), \[
f'(1)=\left(\frac{\widehat h(1)}{\sqrt2}\right)^{2/3}
\] is bounded above and away from zero. The three-derivative bounds in
(C24), on a slightly larger fixed disk, give a uniformly invertible
conformal coordinate after the disk radius is fixed sufficiently small.
The lens arcs within this disk can be chosen as the inverse images of
the two Airy rays; their intersection angles with the disk boundary are
uniformly nonzero.

The exact local matrix is \[
P(z)=E(z)\Psi(n^{2/3}f(z))e^{n\phi(z)\sigma_3},
\qquad
E(z)=N(z)C_A^{-1}f(z)^{\sigma_3/4}n^{\sigma_3/6}.
\tag{A9}
\] It has precisely the jumps in (A4): conjugation of the two triangular
jumps in (A7) produces \(e^{-2n\phi}\) and \(e^{2n\phi}\), while
\(\phi_-=-\phi_+\) leaves the interval jump equal to \(J_0\). The factor
\(E\) is analytic across the interval because \[
C_AJ_0C_A^{-1}=\operatorname{diag}(-i,i).
\] Its apparent endpoint singularity is removable: the fourth-root
singularity of \(N\) is canceled by that of \(f^{\sigma_3/4}\). Equation
(A8) gives, uniformly on the fixed boundary circle, \[
P(z)N(z)^{-1}=I+O_h(n^{-1}).\tag{A10}
\] The left-edge construction is obtained by the exact reflection
\(z\mapsto-z\) and conjugation by \(\sigma_3\). Reflection reverses
orientations; \(\sigma_3\) inverts all three constant jumps in (A7), and
\(N(-z)=\sigma_3N(z)\sigma_3\). Thus the same matching estimate holds at
the left edge, with no new metric or omitted endpoint.

\subsection{A3. The narrow contour has an adequate finite Cauchy
bound}\label{sa:a3.-the-narrow-contour-has-an-adequate-finite-cauchy-bound}

Outside the two endpoint disks put \(R=SN^{-1}\); inside put
\(R=SP^{-1}\). The central portions of the lens have scaled distance
\(\varepsilon_0a/A\) from the real line. The remaining portions are a
fixed finite number of segments and smooth arcs chosen in the uniformly
analytic regions, as described in C5. All intersection angles are
bounded away from zero.

The boundary Cauchy operator on this contour satisfies \[
\|C_-\|_{2\to2}\le C_h A/a.\tag{A11}
\] One explicit coarse justification is to split the contour into its
finitely many pieces. On each line or circle piece the usual Cauchy
bound follows from the Fourier Hilbert transform or its circular
version. Adjacent pieces meeting at a fixed angle have cross kernel
bounded by \(C/(s+t)\); the weighted Schur test with weight \(s^{-1/2}\)
gives the finite Carleman bound \(C\pi\). Nonadjacent bounded pieces,
separated by at least \(ca/A\), have Hilbert--Schmidt norm at most their
length product's square root divided by \(2\pi ca/A\). Cross terms with
an unbounded exterior ray have a square-integrable \(1/(1+s)\) tail; two
opposite rays again use the Carleman bound. Summing this finite block
matrix gives (A11). Thus a uniform lower separation independent of \(k\)
has not been assumed.

The two circle jumps are \(I+O(1/n)\) by (A10). Every remaining jump is
\(I+O(e^{-c_ha})\), with integrable tails, by the real positivity and
the analytic-strip estimate (C23). Therefore, for \(w=J_R-I\), \[
\|w\|_1+\|w\|_2+\|w\|_\infty
\le C_h/n+C_he^{-c_ha}.
\] Since \(A\asymp n\) and \(a\to\infty\), the operator
\(C_-(\,\cdot\,w)\) has norm less than one half on a cofinal range. Its
Neumann solution \(\nu\) has \[
\|\nu-I\|_2\le
2\|C_-\|\|w\|_2=O_h(a^{-1})+O_h((n/a)e^{-c_ha}).
\] The exact Cauchy solution then has coefficient \[
R_1=-\frac1{2\pi i}\int\nu(s)w(s)ds=O_h(n^{-1}),\tag{A12}
\] by the \(L^1\) bound for \(w\) and Cauchy--Schwarz for the remaining
product. This establishes the error coefficient actually needed for the
monic norm.

Finally, \((N_1)_{12}=i/2\), and undoing (A2) shows
\((Y_1)_{12}=e^{n\ell}[(N_1)_{12}+(R_1)_{12}]\). Equation (A1) yields \[
H_n(t)=\pi A^{2n+1}e^{n\ell}[1+O_h(n^{-1})]
=\pi A e^{-\lambda}[1+O_h(n^{-1})].
\tag{A13}
\] This proves the exact constant and the growing-family uniformity used
in (C20).

\newpage

\section{Finite evaluation of the remaining fixed source
functional}\label{sa:finite-evaluation-of-the-remaining-fixed-source-functional}

The coefficient \(I_h\) in the result is a fixed one-factor integral,
not a growing determinant or its unevaluated asymptotic limit. This
appendix gives a finite truncation and quadrature error for that
integral using exactly the original RTM envelope. It does not require a
numerical off-critical packet to be invented.

Retain \[
e^{\theta y}w_h(y)\le C(1+|y|)^{-p},\qquad
|\theta|\le\alpha=\pi/2,\qquad p=8m_0-9/2>3.
\] Let \[
m_*:=\int_{-1}^{1}w_h(y)dy>0,\qquad
C_1=\frac{2C}{p-2},\qquad C_2=\frac{2C}{p-3},\qquad B_\mu=C_1/m_*.
\tag{I1}
\] These constants retain the original source mass. For \(R\ge1\),
define finite integrals \[
M_R(\theta)=\int_{-R}^{R}e^{\theta y}w_h(y)dy,\qquad
\mu_R(\theta)=\frac{\int_{-R}^{R}y e^{\theta y}w_h(y)dy}{M_R(\theta)}.
\] Evenness gives \(M_R(\theta)\ge m_*\) on \([0,\alpha]\). The original
envelope gives \[
|M_h(\theta)-M_R(\theta)|\le T_0(R)
:=\frac{2C}{p-1}(1+R)^{1-p},
\] \[
|M_h'(\theta)-M_R'(\theta)|\le T_1(R)
:=\frac{2C}{p-2}(1+R)^{2-p}.
\] The full and truncated first moments are at most \(C_1\) in absolute
value. Comparing the two exact quotients therefore gives \[
\boxed{|\mu_h(\theta)-\mu_R(\theta)|\le
\varepsilon_R:=\frac{T_1(R)}{m_*}+\frac{C_1T_0(R)}{m_*^2}.}
\tag{I2}
\] Both means have absolute value at most \(B_\mu\). Their squared
difference is consequently at most \(2B_\mu\varepsilon_R\).

The derivative of the truncated mean is its actual positive tilted
variance. The second-moment envelope gives \[
0\le\mu_R'(\theta)\le C_2/m_*,\qquad
\left|\frac{d}{d\theta}\mu_R(\theta)^2\right|\le2B_\mu C_2/m_*.
\] For \(N\ge1\), take the literal midpoint values \[
\theta_j=\frac{\alpha(j+1/2)}{N},\qquad
I_{h;R,N}=\frac\alpha N\sum_{j=0}^{N-1}\mu_R(\theta_j)^2.
\tag{I3}
\] The midpoint error for a Lipschitz function of constant \(L\) on
\([0,\alpha]\) is at most \(L\alpha^2/(4N)\), as follows by integrating
the distance to each midpoint. Hence \[
\boxed{
|I_h-I_{h;R,N}|
\le2\alpha B_\mu\varepsilon_R
+\frac{\alpha^2B_\mu C_2}{2Nm_*}.
}
\tag{I4}
\] This error tends to zero at the explicit rate
\(O_h(R^{2-p}+N^{-1})\). The finite integrals in (I3) use the original,
explicitly given xi/zeta density on a bounded interval; any numerical
errors in those finite integrals must be added when performing a
certified numerical evaluation. The formula does not replace them by
artificial packet data.

Equations (M3), (M13), and (I4) specify the coefficient by fixed real
integrals with convergent error control. There is no remaining limit of
growing arithmetic minimum matrices in \(d_0(h)\). No numerical sign is
asserted for the simple-stratum coefficient without evaluating its
actual fixed-source functional.

\newpage

\section{Sources, calculations, and verification
record}\label{sa:sources-calculations-and-verification-record}

\subsection{Original programme inputs}\label{sa:original-programme-inputs}

The original fixed-packet definitions, source and relation minima, and
the complete receiver are retained from the supplied programme.
\texttt{SOURCES.json} records exact physical line ranges and SHA-256
hashes for every source extract. The two preceding central calculations
and their consolidation are preserved unchanged in \texttt{inputs/},
including the exact preceding ZIP.

\textbf{{[}P1{]}}
\emph{06\_Foundations\_and\_Original\_Six\_Readers(1).tex}, physical
lines 111601--111947: SXR original arithmetic source, all monic minima,
four endpoint signs, and complete relation/quotient exchange.

\textbf{{[}P2{]}} The same collection, physical lines 112323--112945:
RTM original tilted source and full convolution mass, and GTM original
whole-line comparison and Schur transfers.

\textbf{{[}P3{]}} \emph{01\_Current\_Cumulative\_and\_Proofs(1).tex},
physical lines 4350--5254: EIQ explicit equilibrium, its variational
derivatives and endpoint integral; GEL exact Gamma ensemble, source
scaling, and the constant \(\mathcal C_\Gamma\).

\textbf{{[}P4{]}} \emph{Untitled 1775.md}, physical lines 9383--9400 and
6437--6456: the assigned signed-coefficient precision and the complete
original action receiver. Its previous successful finite approximations
were not treated as values for their unevaluated centers.

\textbf{{[}P5{]}} \emph{CONSOLIDATED\_LOCALIZATION.md}: the original
central/outer endpoint split, full nonlinear central transfer,
Taylor-jet restriction bounds, relation suppression, and the
\(O_h(k^2)\) central bound. Its two remaining value questions are
evaluated in the present outer and central proofs.

\subsection{External primary
references}\label{sa:external-primary-references}

\textbf{{[}I{]}} Aleksandar Ivić. \emph{On the mean square of the
Riemann zeta-function in short intervals}. arXiv:0803.0132v3, 15 January
2009. \href{https://arxiv.org/html/0803.0132}{Primary text}. The actual
input is the established second-moment estimate for \(E(t+U)-E(t)\)
stated after (2.1), with its original range. The paper's conjectural
fourth-moment and pointwise estimates are not used.

\textbf{{[}K{]}} H. T. Koelink and J. Van der Jeugt. \emph{Bilinear
generating functions for orthogonal polynomials}. arXiv:q-alg/9704016.
\href{https://arxiv.org/html/q-alg/9704016}{Primary text}. The input is
the orthonormal Meixner--Pollaczek measure and Proposition 2.1, equation
(2.9). The change to the programme's \(y\) coordinate and Gamma mass is
written in O6.

\textbf{{[}D{]}} P. Deift, T. Kriecherbauer, K. T.-R. McLaughlin, S.
Venakides, and X. Zhou. \emph{Uniform asymptotics for orthogonal
polynomials}. Proceedings of the International Congress of
Mathematicians 1998, volume III, pp.~491--501.
\href{https://ems.press/content/book-chapter-files/27262?nt=1}{Publisher
text}. The exact polynomial/Cauchy-transform matrix formulation and the
equilibrium, lens, and Airy method are used. The moving-family strip,
edge, and contour estimates are derived separately in C5 and Appendix A;
uniformity is not inferred from a fixed-potential statement.

\textbf{{[}N{]}} NIST Digital Library of Mathematical Functions.
\href{https://dlmf.nist.gov/5.11}{§5.11} for Gamma asymptotics;
\href{https://dlmf.nist.gov/15.6.E1}{§15.6.1} for the hypergeometric
integral; \href{https://dlmf.nist.gov/9.2}{§9.2} and
\href{https://dlmf.nist.gov/9.7}{§9.7} for Airy connection formulas,
Wronskians, and sectorial asymptotics. All substitutions and matrix
constants used here are displayed in the proof.

\subsection{What was executed}\label{sa:what-was-executed}

The self-contained \texttt{checks/verify\_exact.py} uses explicit
exceptions rather than Python assertions. Its 42 exact checks and five
deliberately false-formula controls produce identical output in ordinary
Python and under \texttt{python\ -O}. They cover finite coefficient
arithmetic, original endpoint weights, full Schur determinants and
nonzero cross blocks, scalar mass cancellation, curvature/area algebra,
the Airy sector jumps, the exact norm scaling, and the cancellation of
the outer logarithmic terms.

The controls reject dropping the relation return, discarding the Schur
cross block, assigning the unweighted source area as the determinant
coefficient, losing the coordinate Jacobian, and deleting an Airy
connection factor. These are exact auxiliary diagnostics, not formal
verification of every analytic estimate.

The central numerical diagnostic uses an explicitly artificial quadratic
rate obtained from a Gaussian tilted-mass function, not an off-critical
xi packet. It evaluates the four original Gamma determinant endpoints.
With \(v=0.1\), its ratios \(\Phi/(k^2I)\) are:

\begin{longtable}[]{@{}rrr@{}}
\toprule\noalign{}
\(k\) & \(q=(k+1)^2\) & Computed ratio \\
\midrule\noalign{}
\endhead
\bottomrule\noalign{}
\endlastfoot
5 & 36 & 0.219020284196899 \\
9 & 100 & 0.220053135775213 \\
17 & 324 & 0.220462759517792 \\
25 & 676 & 0.220549730730594 \\
41 & 1764 & 0.220597259596785 \\
65 & 4356 & 0.220614166598438 \\
\end{longtable}

The derived limiting value is \(\log 2/\pi\approx0.220635600152652\).
These non-interval numerical results are a diagnostic of the scale and
factor, not the proof of that limit.

Three separate high-precision quadratures of the original EIQ endpoint
integral, with increasing precision and different tail splits, give \[
\mathcal C_\Gamma\approx0.6540998948244260328147439735.
\] The calculation uses a cancellation-free form of the original
integrand and an explicit reciprocal tail coordinate. The stored
decimals are not an interval certificate. The analytic expression and
the original proved positivity, rather than these decimals, govern the
theorem.

No numerical off-critical zero, unspecified original period, or
numerical value for the unspecified original functional \(I_h\) was
invented. No independent peer review, Lean certification, or full-corpus
re-verification is claimed. The current result is the written derivation
and its stated supporting checks, with the original broader programme
questions left separate.


% BEGIN SAI RECEIVING CALCULATION 2026-09-19
\section{Receiving calculation: exact central moments and finite Robin error}
\label{sa:sec:central-finite-repair}

This calculation retains the source order \(s=1\), the actual complete
packet and all the definitions of M1--M14.  In particular
\[
 q=[1+k(m_0-1)](k+1)^2,\qquad
 \chi_k(S)=\prod_{u,v=0}^k
 [S-k/2-(2u-k)\delta-i(2v-k)\gamma]^{1+k(m_0-1)}.
\]
The coordinate is still \(S=k/2+iy\).  The polynomial
\(Q(y)=i^{-q}\chi_k(k/2+iy)\), every repeated root, and both
integration half-lines occur in the source/relation determinants.
None of the calculations below assigns an allocation to a proper
source, an observation kernel, or a projected arithmetic class.

\subsection{The phase correction and its exact effect}

In Appendix A1 the displayed boundary convention is
\(g_++g_--V_n=\ell\).  With that convention the real exterior
identity is
\begin{equation}\tag{SAI1}
 \phi(x)=\frac12\{V_n(x)-2g(x)+\ell\},\qquad x>1.
\end{equation}
The minus sign before \(\ell\) in the received prose was incorrect.
Here is the verification including its constant.  By C17 and A3,
\[
 V_n'(z)-2g'(z)=\widehat h(z)\sqrt{z^2-1}=2\phi'(z).
\]
Both \(2\phi\) and \(V_n-2g+\ell\) vanish at \(x=1\) by the
Euler equation.  Integration on \(x>1\) proves SAI1.  In the exact
transformation
\[
 T=e^{-n\ell\sigma_3/2}Y e^{-n(g-\ell/2)\sigma_3},
\]
the upper-right real exterior jump is consequently
\(e^{-n(V_n-2g+\ell)}=e^{-2n\phi}\).  Thus A4--A13 use the
correct jump and need no sign change.  In particular
\[
 (Y_1)_{12}=e^{n\ell}\{i/2+(R_1)_{12}\},\qquad
 H_n=\pi A^{2n+1}e^{n\ell}\{1+2(R_1)_{12}/i\}.
\]
Since \(n\ell=-\lambda-2n\log A\), the original norm factor is
\(\pi A e^{-\lambda}\).  This calculation follows the
polynomial/Cauchy-transform formulation of Fokas, Its and Kitaev
as explicitly stated in Deift, Kriecherbauer, McLaughlin, Venakides
and Zhou, \emph{Uniform asymptotics for orthogonal polynomials},
ICM 1998, Theorem 31 and the subsequent transformation
(\url{https://ems.press/content/book-chapter-files/27262?nt=1}).
The Airy identities are those cited in A2 from
\url{https://dlmf.nist.gov/9.2} and
\url{https://dlmf.nist.gov/9.7}.  Their moving-family estimates
remain C22--C26 and A11--A12, rather than a fixed-potential
uniformity assumption.

\subsection{All even moments of the original central rate}

Write \(\alpha=\pi/2\), \(L=\log M_h\), \(\mu=L'\),
\(b=\mu(\alpha)\), and let \(\theta\) be the inverse of \(\mu\)
on \([0,b]\).  The variance bounds ACW4--ACW5 imply
\(0<v_*\le\mu'\le V_*\), so this inverse exists through its
endpoints.  Define, without changing the source measure,
\[
 D(t)=L(\alpha)-L(\theta(t))-(\alpha-\theta(t))t,\quad
 W_k(y)=kD(|y|/k)\mathbf1_{\{|y|\le kb\}}.
\]
Differentiation gives \(D'(t)=\theta(t)-\alpha\) and \(D(b)=0\);
hence \(D(t)=\int_t^b(\alpha-\theta(u))\,du\).
For every integer \(j\ge0\) there is the exact identity
\begin{equation}\tag{SAI2}
 \int_{\mathbb R}y^{2j}W_k(y)\,dy
 =
 \frac{k^{2j+2}}{(j+1)(2j+1)}
 \int_0^\alpha\mu(\theta)^{2j+2}\,d\theta.
\end{equation}
Indeed, evenness and the literal substitution \(y=kt\) give
\[
 2k^{2j+2}\int_0^b t^{2j}
       \int_t^b(\alpha-\theta(u))\,du\,dt
 =
 \frac{2k^{2j+2}}{2j+1}
       \int_0^b u^{2j+1}(\alpha-\theta(u))\,du.
\]
The substitution \(u=\mu(\theta)\) has Jacobian \(\mu'(\theta)\).
Integration by parts using
\((\mu^{2j+2})'=(2j+2)\mu^{2j+1}\mu'\)
has zero boundary terms: \(\mu(0)=0\) and
\(\alpha-\theta=0\) at \(\theta=\alpha\).  This proves SAI2.
Every integral is finite because \(0\le\mu\le b\); it makes
no assertion about higher moments of the boundary-tilted density.

In particular set
\begin{equation}\tag{SAI3}
 Z=k^2 I_h,\qquad
 J_{4,h}=\int_0^\alpha\mu(\theta)^4\,d\theta,\qquad
 M_2=\frac{k^4}{6}J_{4,h}+\frac{a^2}{2}Z,
 \quad a=\frac{k}{\log^4(q+2)}.
\end{equation}
For \(\widetilde W_k=G_a*W_k\), Gaussian mass, mean and variance
give exactly
\[
 \int\widetilde W_k=Z,\qquad
 \int y^2\widetilde W_k(y)\,dy=M_2.
\]
More generally,
\[
 \int y^{2j}\widetilde W_k(y)\,dy
 =
 \sum_{r=0}^j
  {2j\choose2r}\frac{(2r)!}{4^r r!}a^{2r}
  \int y^{2j-2r}W_k(y)\,dy.
\]
This last formula follows by expanding \((y+u)^{2j}\), using
the vanishing odd Gaussian moments and
\(\int u^{2r}G_a(u)\,du=(2r)!a^{2r}/(4^r r!)\).
Thus every rate moment used in a controlled expansion has an
explicit fixed-source integral and its exact return under smoothing.

\subsection{A finite radius and Robin-variation estimate}

For \(q\le n\le2q\) and \(0\le t\le1\), let \(A=A_{n,t}\)
be the unique solution C15 for the actual Gaussian-convolved
curvature \(\kappa^a_{k,t}\).  Its total mass is \(2\alpha\) and
its second moment is exactly \(2tZ+\alpha a^2\).  Put
\[
 A_*=\frac{\pi n}{\alpha}=2n,\qquad
 B_2=2Z+\alpha a^2.
 \]
The radius bounds, with no omitted tail of the curvature, are
\begin{equation}\tag{SAI4}
 A_*\le A_{n,t}\le A_*+\frac{B_2}{2\alpha A_*}.
\end{equation}
To prove them, write C15 as
\[
 \frac{\alpha A}{\pi}-n
 =
 \frac1{2\pi}\left[
  \int_{|u|\le A}(A-\sqrt{A^2-u^2})\kappa^a_{k,t}(u)\,du
  +A\int_{|u|>A}\kappa^a_{k,t}(u)\,du
 \right].
 \]
Every term is nonnegative.  Inside the interval,
\(A-\sqrt{A^2-u^2}\le u^2/A\); outside,
\(A\le u^2/A\).  The entire right side is at most
\(B_2/(2\pi A)\).  This proves the lower bound, and then
the upper bound after \(A\ge A_*\) is inserted.

Keep the explicit finite guards \(a\ge1\) and \(K=kb\le A_*/4\).
They hold on a cofinal range for every fixed packet and fixed
complete multiplicity.  Define
\begin{equation}\tag{SAI5}
 \begin{split}
 E_{\lambda}(n)={}&
 \frac{M_2+ZB_2/(2\alpha)}{\pi A_*^3}\\
 &+Z\left(\frac1{\sqrt\pi a}+\frac1{\pi A_*}\right)
          \exp\!\left(-\frac{A_*^2}{16a^2}\right).
 \end{split}
\end{equation}
Then the exact Robin variation C27 obeys
\begin{equation}\tag{SAI6}
 \left|\partial_t\lambda_{n,t}
                 -\frac{Z}{\pi A_*}\right|
 \le E_{\lambda}(n).
\end{equation}
Here is the complete estimate on both regions.  On \(|y|\le A/2\),
\[
 0\le\frac1{\sqrt{A^2-y^2}}-\frac1A\le\frac{y^2}{A^3}.
\]
For \(0\le z\le1/4\), the scalar inequality behind this line is
\((1-z)^{-1/2}\le1+z\): after squaring it becomes
\(z(1-z-z^2)\ge0\).  Integrating yields at most
\(M_2/(\pi A^3)\).

For \(|y|\ge A/2\), since \(W_k\) is nonnegative, supported in
\([-K,K]\), and has mass \(Z\),
\[
 \widetilde W_k(y)
 \le \frac Z{\sqrt\pi a}
       \exp\!\left(-\frac{(|y|-K)^2}{a^2}\right)
 \le \frac Z{\sqrt\pi a}e^{-A^2/(16a^2)}.
\]
The first remaining contribution integrates this bound against
the arcsine measure on \(A/2<|y|<A\), whose total mass is at
most one.  The second, subtracting the unweighted tail mass
in \(Z/(\pi A)\), is at most
\(Z e^{-A^2/(16a^2)}/(\pi A)\).
For this last statement use Tonelli and the Gaussian tail
\(\operatorname{erfc}(x)\le e^{-x^2}\), \(x\ge0\).  A direct
proof is to differentiate their difference: its derivative is
\(2e^{-x^2}(1/\sqrt\pi-x)\); the difference starts at zero,
first increases, then decreases to zero, and stays nonnegative.
These estimates account for the entire line, including
\(|y|>A\).  Finally SAI4 gives
\[
 \left|\frac{Z}{\pi A}-\frac{Z}{\pi A_*}\right|
 \le\frac{ZB_2}{2\pi\alpha A_*^3}.
\]
Together they prove SAI5--SAI6.

\subsection{A finite error for the actual monic norm}

The contour of A1--A3 is the one belonging to the potential C14
at this exact \(n,k,t\).  Let \(w=J_R-I\) be its error jump.
For the same oriented contour and arclength measure, set
\[
 B=\|C_-\|_{L^2\to L^2},\quad
 u_j=\|w\|_{L^j}\ (j=1,2),\quad
 u_\infty=\|w\|_\infty,\quad
 r_{n,t}=\frac1\pi
  \left(u_1+\frac{B u_2^2}{1-Bu_\infty}\right).
\]
Impose the explicit guards \(Bu_\infty<1\) and \(r_{n,t}<1\).
They hold on the cofinal range constructed in C22--C26:
\(B\le C_h A/a\) and \(u_1+u_2+u_\infty\le
C_h(n^{-1}+e^{-c_h a})\).
The singular-integral equation
\(\nu=I+C_-(\nu w)\) gives
\[
 \|\nu-I\|_2\le\frac{B u_2}{1-Bu_\infty}.
\]
In its Cauchy solution
\(R_1=-(2\pi i)^{-1}\int\nu w\), the constant and remaining
parts are bounded respectively by \(u_1\) and
\(B u_2^2/(1-Bu_\infty)\).  SAI1 therefore proves
\begin{equation}\tag{SAI7}
 \left|\log\frac{H_n(t)}
                  {\pi A_{n,t}e^{-\lambda_{n,t}}}\right|
 \le-\log(1-r_{n,t}).
\end{equation}
The quantity inside the logarithm is positive because it is
the norm of its actual nonzero monic minimizer.  Thus no
complex logarithm branch is selected to obtain this inequality.
The derived bounds give \(r_{n,t}=O_h(n^{-1})\), uniformly;
SAI7 also records the finite Neumann denominator, rather than
discarding it when the limiting order is taken.

Write \(f(x)=-\log(1-x)\), and put
\begin{equation}\tag{SAI8}
 E_H(n)=E_{\lambda}(n)
 +\frac{B_2}{2\alpha A_*^2}
 +f(r_{n,0})+f(r_{n,1}).
\end{equation}
The two endpoint radii lie in the interval SAI4, so their
logarithmic ratio has absolute value at most
\(B_2/(2\alpha A_*^2)\).  Integrating SAI6 in \(t\), and
subtracting the two estimates SAI7, proves
\begin{equation}\tag{SAI9}
 \left|\log\frac{H_n(0)}{H_n(1)}
                 -\frac{k^2 I_h}{2\pi n}\right|
 \le E_H(n).
\end{equation}
This is the same actual monic minimum as in C29, with an
explicit decomposition of its equilibrium, radius and contour
errors.  The lower polynomial space and its Schur inverse
have already been included in that minimum.

\subsection{Exact weighted-band summation and complete-action receiver}

Keep \(c_0=c_q=1\) and \(c_r=2\) for \(0<r<q\), exactly as
in SXR10.  The two overlapping original source bands give
\[
 S_q=\sum_{r=0}^q\frac{c_r}{q+r}
   =2(H_{2q}^{\rm harm}-H_q^{\rm harm})+\frac1{2q}.
\]
There is the sharpened finite inequality
\begin{equation}\tag{SAI10}
 0<S_q-2\log2<\frac1{8q^2}.
\end{equation}
Indeed integration of exponentials gives
\[
 S_q-2\log2
 =\int_0^\infty(e^{-qt}-e^{-2qt})
         \left(\coth(t/2)-\frac2t\right)dt.
\]
For \(u>0\), \(0<\coth u-1/u<u/3\).  The first inequality
follows from \((u\cosh u-\sinh u)'=u\sinh u>0\).
For the second, expand
\((u^2+3)\sinh u-3u\cosh u\):
the coefficient of \(u^{2j+1}\) is
\(4j(j-1)/(2j+1)!\), positive for \(j\ge2\), and the earlier
coefficients vanish.  Integrating the resulting upper bound
\((e^{-qt}-e^{-2qt})t/6\) proves SAI10.

Consequently, for the regularized four-endpoint expression
\(\widetilde\Phi\) in C30, SAI9 proves the finite return
\begin{equation}\tag{SAI11}
 \left|\widetilde\Phi-\frac{\log2}{\pi}k^2I_h\right|
 \le\sum_{r=0}^q c_r E_H(q+r)
       +\frac{k^2I_h}{16\pi q^2}.
\end{equation}
The smoothing and arithmetic localization are the attained
Taylor-image comparisons C8 and L22--L35; their full kernel
restriction and quotient factors remain present.  Adding
their proved errors to SAI11 gives exactly C32, with a
sharper harmonic-sum contribution.  Thus the central endpoint
value, rather than only its source-density area, is
\[
 \delta_k^{\rm cen,end}
 =\frac{\log2}{\pi}I_h k^2
 +O_h\!\left(\frac{k^2}{\log^3(q+2)}
       +k\log^2(q+2)+1+\frac{k^4}{q^2}\right).
\]
The outer proof O1--O35 uses the actual convolution on
\(|y|>kb\), both half-lines, the full relation polynomial,
and the original four signs.  Its bound is added before
any scale comparison.  Therefore the accepted receiver is
\begin{equation}\tag{SAI12}
 \begin{split}
 J_k^{\rm action}={}&\mathcal B_k[\sigma]
 -(4m_0-2)\mathcal C_\Gamma q
 +\frac{\log2}{\pi}I_h k^2\\
 &+W_k^{\rm mono}-P_{k,-}-P_{k,+}+\mathcal E_{h,k},
 \end{split}
\end{equation}
with the explicit sum of the central bound above and O35
as a bound for \(|\mathcal E_{h,k}|\).
Here \(\mathcal B_k[\sigma]\) is its full existing value.
It is not replaced by its leading positive term; neither
penalty nor the monic-window scalar has been dropped.
This new scalar evaluation changes the previously uncomputed
\(s=1\) arithmetic-baseline entry.  It does not compute the
independent absolute allocations or same-class projected signs.

% END SAI RECEIVING CALCULATION 2026-09-19

% END PRESERVED INPUT 01_SIGNED_ARITHMETIC_INTEGRATED.tex
\endgroup

\clearpage
\begingroup
\def\PP{\mathcal P}
\def\MM{\mathfrak M}
\def\BB{\mathcal B}

\setcounter{equation}{0}
\renewcommand{\theequation}{OR.\arabic{equation}}
\section*{Outer arithmetic value: independent complete reconstruction}
\addcontentsline{toc}{section}{Outer arithmetic value: independent complete reconstruction}
% BEGIN PRESERVED INPUT 02_OUTER_VALUE_RECONSTRUCTION.tex SHA256 3efd41217f4ec85da5f1d862df8997d220a11441da3a2ad871ef6bf2ab6e38fa



This receiving proof concerns the outer part of the arithmetic density.
The original source, packet, relation ideal, four endpoint signs and
Gamma mass are retained.  The original EIQ equilibrium and RTM
one-factor bounds are the proved programme providers specified in the
accompanying acceptance record; no arithmetic allocation or projected
sign is inferred from a full-quotient determinant.

\section{Objects and exact receiver}
Fix the original stipulated quartet
\(\rho=1/2+\delta+i\gamma\), \(0<\delta<1/2\), \(\gamma>2\),
with original multiplicity \(m\geq1\).  Define
\[
 h(s)=\prod_{z\in\{\rho,\bar\rho,1-\rho,1-\bar\rho\}}(s-z)^m,\qquad
 w_h(y)=\frac{|(2\xi/h)(1/2+iy)|^2}{2\pi},\qquad m_k=w_h^{*k}.
\]
Keep \(k=4\ell+1\), \(e=1+k(m-1)\), \(q=e(k+1)^2\), \(c=k/2\), and
\[
 \chi_k(S)=
 \prod_{a,b_0=0}^k
 [S-c-(2a-k)\delta-i(2b_0-k)\gamma]^e,\qquad
 Q(y)=i^{-q}\chi_k(c+iy).
\]
Thus \(Q\) is real, even and monic of degree \(q\); all its roots,
including every repetition, have modulus at most
\(R=k(\delta^2+\gamma^2)^{1/2}\).
Write
\[
\begin{gathered}
 \sigma(y)=\frac{|\Gamma(1/4+iy/2)|^2}{2\pi},\quad
 \int_{\mathbb R}\sigma=\sqrt{2\pi},\quad
 \alpha=\pi/2,\\
 M(\theta)=\int e^{\theta y}w_h(y)\,dy,\quad
 \MM=M(\alpha),\quad b=(\log M)'(\alpha)>0 .
\end{gathered}
\]
For a positive density \(\nu\), \(H_N[\nu]\) is the Gram of
\(1,y,\ldots,y^N\), \(J_N\) is remainder modulo the full \(Q\), and
\[
\begin{gathered}
 G_N[\nu]=(J_NH_N[\nu]^{-1}J_N^*)^{-1},\\
 \BB[\nu]=\log\det G_{q-1}[\nu]+\log\det G_q[\nu]
 -\log\det G_{2q-1}[\nu]-\log\det G_{2q}[\nu].
  \end{gathered}\tag{OR1}
\]
The monic frame consisting of the first \(q\) monomials followed by
\(Q,yQ,\ldots,y^{N-q}Q\) has determinant one.  Its Schur complement
proves, with its cross blocks retained,
\[
 \det G_N[\nu]=\frac{\det H_N[\nu]}
 {\det\operatorname{Gram}_{\nu}(Q,yQ,\ldots,y^{N-q}Q)}. \tag{OR2}
\]
The polynomial substitution \(P(S)\mapsto P(c+iy)\) is invertible, has
coefficient determinant of modulus one, and carries \(\chi_kP\) to
\(i^q QP(c+iy)\); this proves that (OR1)--(OR2) compute the original
quotient.  A positive scalar \(A\) multiplies every \(G_N\) by \(A\).
Each quotient has dimension \(q\), hence \(\BB[A\nu]=\BB[\nu]\).

Set
\[
 d=4m-2,\quad p=2d-\tfrac12,\quad a=p+1,\qquad
 \tau_d=(1+y^2)^{-d}\sigma,\qquad
 \omega_k(y)=
 \begin{cases}\sigma(y),&|y|\leq kb,\\
 \MM^{-k}m_k(y),&|y|>kb.\end{cases}                    \tag{OR3}
\]
The notation \(a\) in (OR3) is a tail exponent, distinct from the
integer indexing a root in \(\chi_k\).

\begin{theorem}
On the original fixed-packet cofinal range, with all original finite
guards retained,
\[
 \BB[\omega_k]-\BB[\sigma]
 =-d\mathcal C_\Gamma q+\mathcal E_k,\qquad
 \mathcal C_\Gamma=2\mathcal L(2,\pi)-8\log2+4,          \tag{OR4}
\]
where \(\mathcal L(\alpha_E,\beta)=\int\log x\,\rho_{\alpha_E,\beta}(x)\,dx\)
is the original EIQ logarithmic moment and
\[
 |\mathcal E_k|\leq C_h\left[
 \frac{q\log\log(q+16)}{\log(q+2)}
 +q^{199/200}\log^2(q+2)\right]=o_h(q).                \tag{OR5}
\]
In particular the outer term has no nonzero contribution of magnitude
\(q\log k\) or \(q\log\log k\).  This theorem is about the outer density
in (OR3); it does not delete the central density in \(m_k\).
\end{theorem}

\section{The growing convolution, including both tails}
The exact identity following from the definition of the convolution is
\[
 f(y)=\MM^{-1}e^{\alpha y}w_h(y),\qquad
 \MM^{-k}m_k(y)=e^{-\alpha y}f^{*k}(y).                \tag{OR6}
\]
The auxiliary \(f\) has mass one, mean \(b\) and variance \(V>0\).
The original density is recovered with the entire factor in (OR6).
From the completed-zeta definition one obtains directly
\[
 w_h(y)=\frac{(y^2+1/4)^2}{\sqrt\pi}\,
 \frac{|\zeta(1/2+iy)|^2}
 {[(y-\gamma)^2+\delta^2]^{2m}[(y+\gamma)^2+\delta^2]^{2m}}\sigma(y).
\]
Stirling's formula \cite{or:DLMFgamma}, on this fixed vertical line, gives
\[
 f(y)=c_fy^{-a}Z(y)(1+O_h(y^{-1})),\quad
 c_f=\frac{\sqrt{2/\pi}}{\MM},\quad
 Z(y)=|\zeta(1/2+iy)|^2,\quad y\longrightarrow+\infty. \tag{OR7}
\]
The RTM envelope gives \(f(y)\leq C_h(1+|y|)^{-p}\), \(p\geq7/2\);
on the negative half-line it has an additional
\(e^{-2\alpha|y|}\) factor.  Thus its \(9/4\)-moment is finite.
The mean-square zeta formula gives, whenever \(|z|\leq T^{19/20}\),
\[
 \int_T^{2T}f(y-z)\,dy\leq C_h T^{1-a}\log T.          \tag{OR8}
\]
Indeed \(y-z\) lies in a fixed enlargement of \([T,2T]\), the power
factor in (OR7) is \(O(T^{-a})\), and the zeta integral over that
enlargement is \(O(T\log T)\).

The proved Fourier local-limit estimate in the supplied localization
provider is uniform in the evaluation point:
\[
 \sup_x\left|f^{*n}(nb+x)
 -(2\pi nV)^{-1/2}e^{-x^2/(2nV)}\right|
 \leq C_h n^{-5/8}.                                  \tag{OR9}
\]
Its Fourier proof bounds an \(L^1\) Fourier difference; the factor
introduced by translation has modulus one.  Integrating (OR9) on
\(|x|\leq n^{1/2+1/24}\) costs \(O(n^{-1/12})\).
Chebyshev's inequality for both probability densities bounds the two
complementary masses by \(O(n^{-1/12})\).  If \(g_n\) is the Gaussian
centered at \(nb\), this proves
\[
 \|f^{*n}-g_n\|_1\leq C_h n^{-1/12}.                  \tag{OR10}
\]

For \(q^{9/10}\leq T\leq128q\), the exact definition of \(q\) implies
\[
 c_hT^{1/3}\leq k\leq C_hT^{5/9}.                    \tag{OR11}
\]
Put \(A=T^{19/20}\), \(f_s=f1_{\{|y|\leq A\}}\) and \(f_l=f-f_s\).
Neither subprobability density is normalized.  In the binomial
convolution expansion, terms with at least two large summands have
total mass at most \(C_hk^2A^{2-2p}\).  Dividing by
\(L_T=kT^{-p}\log T\) gives
\[
 C_h kT^pA^{2-2p}/\log T
 \leq C_hT^{-25/36}/\log T.                           \tag{OR12}
\]
The exponent is largest at \(p=7/2\); larger original multiplicity
only improves it.

For \(\lambda=(\log T)/A\), Taylor's formula on the truncated support
gives
\[
 \int e^{\pm\lambda y}f_s(y)\,dy
 \leq\exp(C_h\lambda+C_h\lambda^2T).
\]
Consequently the mass of \(f_s^{*j}\), \(j\leq k\), outside
\([-T/2,T/2]\) is at most
\(2\exp[-c_hT^{1/20}\log T]\), eventually.  Here
\(k/T\to0\) and \(k\lambda\to0\) by (OR11), so the two positive
exponent terms are strictly smaller than the Chernoff exponent.
This bounds the all-small convolution.

The remaining convolution is \(k f_l*f_s^{*(k-1)}\).
On \(A<|z|\leq T/2\), \(f_l(y-z)\leq C_hT^{-p}\), whereas the
rest density has mass at most \(C_h k/A^2\) by its variance and
\(A\geq2(k-1)b\).  Its window integral is at most
\(C_hk^2T^{1-p}/A^2\); divided by \(L_T\), it is
\(O_h(kT/A^2)\), a fixed negative power of \(T\).
On \(|z|>T/2\) use the preceding Chernoff bound and the full mass
of \(f_l\).  On \(|z|\leq A\), one has \(y-z>A\).
Replacing the small rest convolution by the full one has total
variation at most \(C_hkA^{1-p}\).  The window estimate (OR8) bounds
its contribution by \(L_TO_h(kA^{1-p})\).
Replacing that full rest convolution by \(g_{k-1}\), on the same
region, costs \(L_TO_h(k^{-1/12})\) by (OR8) and (OR10).
The complementary Gaussian region is exponentially small, including
the part near \(z=y\), since its center is \(O_h(k)\) and its
variance is \(O_h(k)\).
Finally (OR7) and \(|z|\leq A\) permit replacing
\((y-z)^{-a}\) by \(y^{-a}\), with relative integrated error
\(O_h(A/T)\).  Thus, with \(n=k-1\) and \(G=\sqrt{2nV}\),
\[
\begin{gathered}
 \int_T^{2T}\left|f^{*k}(y)-kc_fy^{-a}\mathcal Z_G(y-nb)\right|dy
 \leq C_h L_T(k^{-1/12}+T^{-1/20}),\\
 \mathcal Z_G(t)=\frac1{\sqrt\pi G}\int e^{-(u/G)^2}Z(t+u)\,du .
  \end{gathered}\tag{OR13}
\]
This estimate uses an averaged zeta norm; it asserts no pointwise
regular variation across zeta zeros.

\section{The arithmetic smoothing error}
Let \(c_\zeta=2\gamma_{\rm E}-\log(2\pi)\) and
\(\mathcal L_\zeta(t)=\log t+c_\zeta\).  The established theorem
of Ivi\'c \cite[Section 2, (2.1) and its stated \(E\)-analogue]{or:Ivic}
gives, for \(1\ll U\leq\sqrt T/2\),
\[
 \int_T^{2T}|E(t+U)-E(t)|^2dt
 \ll_\varepsilon TU\log^3 T+
 T^{1/2+\varepsilon}U^2+T^{1+\varepsilon}U^{1/2},       \tag{OR14}
\]
where \(E(t)=\int_0^tZ(v)\,dv-t(\log(t/(2\pi))+2\gamma_{\rm E}-1)\).
The cited result is a second-moment theorem, without a pointwise or
fourth-moment conjecture.
After truncation at \(|u|\leq G\log T\), integration by parts and
subtraction of \(E(t)\) under the odd kernel give
\[
 \mathcal Z_G(t)-\mathcal L_\zeta(t)
 =\frac2{\sqrt\pi G^3}\int
 u e^{-(u/G)^2}(E(t+u)-E(t))\,du+O(G/T)+O(T^{-100}).
 \tag{OR15}
\]
Polynomial bounds from (OR7) justify the Gaussian tails.  Minkowski's
inequality applied to the square roots of the three terms in
(OR14) yields respectively
\[
 G^{-1/2}\log^{3/2}T,\quad
 T^{-1/4+\varepsilon},\quad T^\varepsilon G^{-3/4}
\]
after division by \(\sqrt T\).  The Gaussian integrals giving these
powers are \(G^{-3}\int_0^\infty u^{3/2}e^{-(u/G)^2}du\),
\(G^{-3}\int_0^\infty u^2e^{-(u/G)^2}du\), and
\(G^{-3}\int_0^\infty u^{5/4}e^{-(u/G)^2}du\).
Negative \(u\) is reduced by translation and a bounded dyadic cover.
For \(|u|\leq1\), positivity of \(Z\) bounds the error increment by
the length-four increment plus \(O(\log T)\); (OR14) at length four
gives the smaller term \(O(T^\varepsilon G^{-3})\).
Since
\(c_hT^{1/6}\leq G\leq C_hT^{5/18}\), all uses of (OR14) lie in
its stated range eventually.  Taking \(\varepsilon=1/1000\),
Cauchy--Schwarz, (OR11)--(OR15), and the exact (OR6) imply
\[
 \int_T^{2T}\left|
 \frac{\MM^{-k}m_k(y)}
 { [k/(\sqrt\pi\MM)]\mathcal L_\zeta(y)(1+y^2)^{-d}\sigma(y)}
 -1\right|dy\leq C_hT^{199/200}.                      \tag{OR16}
\]
The shift \(nb=O_h(k)\) changes the logarithmic main term by
\(O_h(k/T)\).  Division in (OR16) is legitimate because its
power-log denominator on this window is comparable to
\(kT^{-a}\log T\); the exponential factors cancel exactly.
The original evenness of \(m_k,\sigma,\tau_d\) supplies the negative
half-line.

\section{Full polynomial metrics and their attained quotients}
For the Gamma-orthonormal polynomials \(\varphi_j\), define
\(K_N^\sigma(y,y)=\sum_{j=0}^N|\varphi_j(y)|^2\).
The normalization of Koelink--Van der Jeugt
\cite[Proposition 2.1, (2.9)]{or:KVJ} is transported by
\[
 x=y/2,\quad\lambda=1/4,\quad\phi=\pi/2,\qquad
 \varphi_j(y)=2^{\lambda-1/2}p_j^{(\lambda)}(y/2;\pi/2).
\]
The weighted diagonal density in \(y\) is one half its \(x\) version.
Euler's integral \cite[(15.6.1)]{or:DLMFhyper}, with its
\(\Gamma(2\lambda)\) conversion from the regularized
hypergeometric function, cancels the weight's
\(|\Gamma(\lambda+ix)|^2\).  Its absolute integrand is
\[
 u^{\lambda-1}(1-u)^{\lambda-1}(1+R_tu)^{-\lambda},
 \qquad R_t=4t/(1-t)^2 .
 \]
For \(t\geq1/2\), splitting at \(1/R_t\) and \(1/2\) bounds this
integral by \(R_t^{-1/4}(16+2\log R_t)\).
Taking \(t=(N+1)/(N+2)\), using positivity of each diagonal
summand and \(t^{-N}\leq e\), proves
\[
 K_N^\sigma(y,y)\sigma(y)\leq20\log(N+2).              \tag{OR17}
\]
For any subspace \(S\subset\PP_{2q}\), orthogonal projection gives
\(K_S\leq K_{2q}^\sigma\) and \(\int K_S\,d\sigma=\dim S\).

In a Gamma-orthonormal frame of \(S\), expansion of the squared
determinant proves
\[
 \det H_S[r\sigma]
 =\frac1{s!}\int|\det[\varphi_i(y_j)]|^2
              \prod_{j=1}^s r(y_j)\,d\sigma(y_j).
 \]
The unweighted right-hand measure has mass one and one-point
density \(K_S\,d\sigma\).  Jensen's inequality and
\(\log\lambda\leq\lambda-1\), applied to the Gram eigenvalues, give
\[
 \int\log r\,K_S\,d\sigma
 \leq\log\frac{\det H_S[r\sigma]}{\det H_S[\sigma]}
 \leq\int(r-1)K_S\,d\sigma.                          \tag{OR18}
\]
This is a statement about the full Gram determinant, including all
cross blocks.

Let \(K=q^{9/10}\) and
\[
 A_k=\frac{k\mathcal L_\zeta(q)}{\sqrt\pi\MM},\qquad
 \widehat\omega_k=
 \begin{cases}A_k\tau_d,&|y|\leq K,\\
 \omega_k,&|y|>K,\end{cases}\qquad
 r_k=\widehat\omega_k/(A_k\tau_d).                    \tag{OR19}
\]
The proved localization bounds, obtained by placing the
\((k-3)\)-fold tilted convolution in its central Gaussian interval
and retaining the three-factor RTM lower tail, give
\[
 c_h k^{-1/2}(1+y^2)^{-M_*}
 \leq\omega_k(y)/\sigma(y)\leq C_hk^{p+1}.            \tag{OR20}
\]
They start at \(|y|=kb\); inside, the quotient is exactly one.
Here \(M_*=21+4m\).  In detail, the interval
\([(k-3)b-\sqrt{(k-3)V},(k-3)b]\) has tilted probability bounded
below by a positive fixed constant, by (OR9).  For \(y\geq kb\)
the remaining positive argument is at most \(2y\), eventually.
Inserting the original three-factor estimate
\(w_h^{*3}(u)\geq c_he^{-\alpha|u|}(1+|u|)^{-42-8m}\)
in that part of the exact convolution proves (OR20); evenness
handles the negative side.  This derivation retains \(\MM^{k-3}\),
which becomes the fixed factor \(\MM^{-3}\) after (OR6).

On \(\PP_{2q}\), Gamma multiplication has norm at most \(2(2q+1)\).
Jensen's inequality applied to the convex function
\((1+x)^{-M_*}\), in the probability
\(|P|^2d\sigma/\|P\|_\sigma^2\), and (OR20), therefore prove
\[
 q^{-C_h}\|P\|_\sigma^2\leq
 \|P\|_{\omega_k}^2,\ \|P\|_{\widehat\omega_k}^2,\
 A_k\|P\|_{\tau_d}^2
 \leq q^{C_h}\|P\|_\sigma^2 .                        \tag{OR21}
\]
The constants here can be enlarged to absorb fixed finite source
constants; the assertion applies on the explicitly cofinal original
sequence.  The same inequality holds for every restriction, and
for a quotient by taking the minimum in each unchanged affine
fibre.

For completeness the small central replacement in (OR19) has its
own finite metric comparison.  The exact generating function
\[
 \sum_{j=0}^\infty\frac{b_j(z)}{j!}t^j
 =(1+t^2)^{-1/4}\exp(z\arctan t),\qquad
 \gamma_j=\|b_j\|_\sigma^2=\sqrt{2\pi}(2j)!/4^j
\]
follows by summing the original three-term recurrence.
Cauchy's coefficient estimate at \(|t|=(2q+1)/(2q+2)\)
and then at \(|z|=2K\) shows, for any \(P\) vanishing to order \(r\),
\[
 \int_{-K}^K|P|^2d\sigma
 \leq A_{q,K,r}\|P\|_\sigma^2,\quad
 A_{q,K,r}=4^{1-r}e^2(2q+1)^2\sqrt{2q+2}
                    [2(2q+2)]^{2K}.                \tag{OR22}
\]
Choose the actual integer \(r\) so that
\(A_{q,K,r}q^{2C_h}\leq q^{-20}\); one may take
\(r=\lceil C'_h(K+1)\log(q+2)\rceil\).
On any actual source or relation space use its own restricted
Taylor map to \(\mathbb C^r\), its full kernel and its attained
image.  The kernel metric changes by relative \(O(q^{-20})\);
its log determinant changes by \(O(q^{-19})\).
The attained image has dimension at most \(r\), so (OR21) costs
at most \(C_h r\log q\).  Factoring the original Gram by this
exact short sequence proves
\[
 \left|\log\frac{\det H_S[\widehat\omega_k]}
                        {\det H_S[\omega_k]}\right|
 \leq C_h q^{9/10}\log^2(q+2)                       \tag{OR23}
\]
for each of the original source and relation spaces.
No quotient frame is discarded; its fixed determinant cancels in
the ratio.

On each dyadic interval \(K\leq |y|\leq128q\), (OR16) bounds
\(r_k-a_q\) in \(L^1\), where
\(a_q(y)=\mathcal L_\zeta(|y|)/\mathcal L_\zeta(q)\).
This \(a_q\) is bounded above and below by fixed positive constants.
The pointwise polynomial bounds in (OR20) imply
\[
 |\log r_k-\log a_q|\leq C_h\log(q+2)|r_k-a_q|.
 \tag{OR24}
\]
To verify this inequality, when \(|r_k-a_q|\leq a_q/2\) use the
bounded derivative of \(\log\) on that interval.  On the complement
the right-hand difference is bounded below by a fixed positive
constant and the logarithms are \(O_h(\log q)\).
Thus (OR17) bounds both error integrals in (OR18) by
\(C_hq^{199/200}\log^2(q+2)\).
For \(q/\log^4q\leq|y|\leq128q\), the deterministic \(a_q-1\)
and its logarithm are \(O(\log\log q/\log q)\); use the exact
kernel mass at most \(2q+1\).  On
\(K\leq|y|\leq q/\log^4q\), use (OR17) and the interval length.
The full remaining tails are bounded by repeated Gamma
multiplication: for \(\deg P\leq2q\),
\[
 \int_{|y|>128q} y^{2d}|P|^2d\sigma
 \leq (128q)^{-2q}[2(3q+d+1)]^{2(q+d)}\|P\|_\sigma^2 .
 \tag{OR25}
\]
Together with (OR20), this gives an exponentially small tail for
both \(r_k-1\) and \(\log r_k\), after summing a Gamma-orthonormal
frame.  Both integration regions are retained.
Writing
\[
 E_q=C_h\left[q^{199/200}\log^2(q+2)
       +q\log\log(q+16)/\log(q+2)\right],
\]
we have therefore proved, uniformly in the required subspaces,
\[
 \int\bigl(|r_k-1|+|\log r_k|\bigr)K_S\,d\sigma\leq E_q.
 \tag{OR26}
\]

The desired reference is \(\tau_d\), and its exact transport is
\(P\mapsto(y+i)^dP\).  It is an isometry from
\(\PP_{N-d}\) in \(\sigma\) onto
\((y+i)^d\PP_{N-d}\) in \(\tau_d\), since
\(|y+i|^{2d}=(1+y^2)^d\).
For the factor \(r_k\) the same identity holds.  On a relation
space the map is \(QP\mapsto Q(y+i)^dP\), preserving the complete
factor \(Q\).  Its omitted quotient is the actual image of
the first \(d\) divided derivatives of \(P\) at \(-i\), hence has
dimension at most \(d\).  Its two metrics have logarithmic
relative cost \(O_h(d\log q)\) by (OR21).
Apply (OR18),(OR26) on the exact image, then include this quotient.
Sum the original signs using (OR2), undo (OR23), and cancel \(A_k\)
only after each quotient's dimension has been retained.  This proves
\[
 \BB[\omega_k]-\BB[\sigma]
 =\BB[\tau_d]-\BB[\sigma]+O_h(E_q).                  \tag{OR27}
\]

\section{The rational model and its finite free-energy error}
Let \(\rho_{\alpha_E,\beta}\) be the original EIQ density for
\(V(x)=\beta\sqrt x-\alpha_E\log x\).  Its Euler inequality is
\(2U-V\leq-\lambda\), and
\(\mathcal F=-\lambda-\iint\log|x-y|\rho(x)\rho(y)\,dxdy\).
The EIQ endpoint formulas give smooth dependence on both positive
parameters and support in a fixed positive compact interval on
the parameter set used below.  The density is bounded there,
has square-root edges, finite logarithmic energy and finite
entropy relative to
\(d\omega(x)=\tfrac12x^{-1/2}e^{-\sqrt x}dx\).

Here is an independent finite-configuration estimate underlying the
shrinking secant.  Replace each configuration point \(x_j\) by
uniform mass on the complex circle \(x_j+\epsilon e^{i\theta}\);
denote the average of these measures by \(\nu_\epsilon\).
The circle mean identity is
\[
 \frac1{2\pi}\int_0^{2\pi}\log|z-\epsilon e^{i\theta}|\,d\theta
 =\log\max(|z|,\epsilon).
 \tag{OR28}
\]
It follows from harmonic mean value outside and inside the circle,
with continuity at its boundary.  For compactly supported
zero-mass signed measures with finite logarithmic energy,
\[
 -\iint\log|z-w|\,d\eta(z)d\eta(w)
 =\frac12\int_0^\infty
 \iint e^{-t|z-w|^2}\,d\eta(z)d\eta(w)\frac{dt}{t}\geq0.
 \tag{OR29}
\]
The scalar logarithm identity proves the equality after truncation;
the zero mass cancels the constant.  The Fourier transform of
the planar Gaussian proves the nonnegative sign.  Absolute
logarithmic integrability of the present circle measures and
bounded density justifies passage to the limit.
Use (OR29) with \(\eta=\nu_\epsilon-\rho\,dx\).
The diagonal circle energy is \(\log\epsilon\), off-diagonal
circle means are at least \(\log|x_i-x_j|\), and (OR28) bounds
the increase of the cross energy by \(2\|\rho\|_\infty\epsilon\).
The Euler inequality now yields
\[
 2\sum_{i<j}\log|x_i-x_j|-n\sum_j V(x_j)
 \leq n^2\mathcal F+n\log n+4n\|\rho\|_\infty
 \quad(\epsilon=1/n).                               \tag{OR30}
\]

The literal parity coordinate is \(y^2=q^2x\).  With
\(B_q(x)=2qe^{\sqrt x}\sigma(q\sqrt x)\), its exact ensemble is
\[
 J_{n,q}(n\alpha_E)=\frac1{n!}\int
 \prod_{i<j}(x_i-x_j)^2
 \prod_j x_j^{n\alpha_E}B_q(x_j)\,d\omega(x_j).
 \tag{OR31}
\]
The original Gamma bound gives
\(B_q(x)\leq Cq^{3/2}
e^{-(\pi q/2-3/2)\sqrt x}\).
Apply (OR30) with
\(\beta_n=\pi q/(2n)-3/(2n)\), and integrate against the
probability \(\omega^n\).
The explicit prefactor and \(n!\) cost \(O(n\log(q+n))\).
Changing \(\beta_n\) to \(\pi q/(2n)\) costs \(O(n)\), since
\(\partial_\beta\mathcal F=-\int\sqrt x\,\rho\,dx\) is uniformly
bounded here.  For the lower bound restrict to the actual
equilibrium support, use the original Gamma lower bound there,
and change measure to \(\rho^n\).  Jensen's inequality gives
the same \(n^2\mathcal F\), with \(O(n)\) missing-diagonal
energy and entropy, plus \(O(n\log(q+n))\) prefactors.  Hence
\[
 |\log J_{n,q}(n\alpha_E)
   -n^2\mathcal F(\alpha_E,\pi q/(2n))|
 \leq Cn\log(q+n),                                  \tag{OR32}
\]
uniformly when \(q/n\) is near \(2\), and \(\alpha_E\) is near
\([7/4,9/4]\).

For the original relation rank \(r=q,q+1\), define
\[
\begin{gathered}
 D_{k,r}(t)=\det\operatorname{Gram}_{(1+y^2)^{qt}\sigma}
                                     (Q,yQ,\ldots,y^{r-1}Q),\\
 F_{k,r}(t)=q^{-2}\{\log[D_{k,r}(t)/D_{k,r}(0)]
                                             -2qrt\log q\}.\end{gathered} \tag{OR33}
\]
The squared-Vandermonde integral proves convexity by
differentiating twice: the second derivative of its logarithm
is a nonnegative variance.  The paired-root comparison on
\(|y|\geq2^{-16}q\), and the complete inner-interval estimate
in the accompanying root-comparison proof, give logarithmic
quadratic-form width \(O_h(1)\) uniformly for \(|t|\leq1/8\).
Its determinant cost is \(O_h(q)\), because the rank is \(r\).
For the reference power \(|y|^{2q(1+t)}\sigma\), its even and
odd blocks have sizes
\[
 n_0=\lceil r/2\rceil,\quad n_1=\lfloor r/2\rfloor,
 \qquad a_0=q(1+t),\quad a_1=q(1+t)+1.
\]
Their parameters satisfy
\[
 (a_i/n_i,\pi q/(2n_i))=(2+2t,\pi)+O(1/q)
 \quad(i=0,1).
 \tag{OR34}
\]
The exact scaling factor of each block is
\(q^{2a_i n_i+2n_i(n_i-1)}\).
Its change is exactly \(2qrt\log q\) after the two blocks
are added.  Applying (OR32) and EIQ smoothness therefore gives
\[
 \sup_{|t|\leq1/8}|F_{k,r}(t)-F_*(t)|
 \leq C_h\frac{\log(q+2)}q,\qquad
 F_*(t)=\tfrac12[\mathcal F(2+2t,\pi)-\mathcal F(2,\pi)] .
 \tag{OR35}
\]
Here \(F_*'(0)=\mathcal L(2,\pi)\) and \(F_*''\) is bounded.
Take \(\epsilon=\sqrt{\log(q+2)/q}\) with
\(d/q<\epsilon/2\) and \(2\epsilon<1/8\).
The convex secant over \([-d/q,0]\) is bounded below and
above by adjacent secants of length \(\epsilon\).
Equation (OR35) changes these latter slopes by at most
\(C_h\log(q+2)/(q\epsilon)\); bounded \(F_*''\) changes their
centers from \(F_*'(0)\) by \(O_h(\epsilon)\).
Returning the exact powers in (OR33) proves
\[
 \log\frac{D_{k,r}(-d/q)}{D_{k,r}(0)}
 =-2dr\log q-dq\mathcal L(2,\pi)
   +O_h(\sqrt{q\log(q+2)}),\qquad r=q,q+1.            \tag{OR36}
\]
No derivative of an asymptotic remainder was taken.

For the source determinant, use the monic frame
\[
 1,\ldots,y^{d-1},(y+i)^d,\ldots,y^{N-d}(y+i)^d .
\]
Its determinant is one and the last block is exactly the Gamma
Gram on \(\PP_{N-d}\).
The first \(d\) divided derivatives at \(-i\) give the quotient
and have determinant one on the first \(d\) monomials.
Their minimum Gram is bounded above by the fixed monomial
trial lifts.  It is bounded below by a negative fixed power
of \(N+2\): use
\(\tau_d\succeq[1+4(N+1)^2]^{-d}\sigma\) on \(\PP_N\),
then the above generating function on fixed circles about
\(-i\) bounds every Gamma jet covariance by a fixed power
of \(N+2\).  There are only \(d\) rows, so the reciprocal
covariance lower bound is again polynomial.
Consequently, retaining this quotient determinant,
\[
 \log\frac{\det H_N[\tau_d]}{\det H_N[\sigma]}
 =-\sum_{j=N-d+1}^N\log\gamma_j+O_d(\log(N+2)).       \tag{OR37}
\]
The exact \(\gamma_j=\sqrt{2\pi}(2j)!/4^j\) gives
\(\log\gamma_j=2j\log j-2j+O(\log(j+2))\).
Insert the four source endpoints and their original signs:
\[
 \sum_{N=q-1,q,2q-1,2q}\epsilon_N
 \log\frac{\det H_N[\tau_d]}{\det H_N[\sigma]}
 =4dq\log q+(8d\log2-4d)q+O_d(\log q).              \tag{OR38}
\]
For the low relation, rank zero contributes \(0\); rank one
has norm ratio between \([1+4(q+1)^2]^{-d}\) and \(1\), by
the same Jensen estimate applied to the full \(Q\).
It contributes \(O_d(\log q)\).
Using (OR36) at the two high relation ranks \(q,q+1\),
and the minus sign of the relation determinant in (OR2),
the net relation contribution to the baseline is
\[
 -4dq\log q-2dq\mathcal L(2,\pi)
 +O_h(\sqrt{q\log(q+2)}).                            \tag{OR39}
\]
Adding (OR38) and (OR39) proves
\[
 \BB[\tau_d]-\BB[\sigma]
 =-d[2\mathcal L(2,\pi)-8\log2+4]q
 +O_h(\sqrt{q\log(q+2)}).                            \tag{OR40}
\]
Together with (OR27), this proves (OR4)--(OR5).

\section{Integration scope}
All constants denoted \(C_h,c_h\) are fixed-source constants with
existence proved by the indicated estimates.  They have not been
evaluated numerically for an unspecified off-critical quartet.
The cofinal assertion retains the original packet cutoffs and uses,
in addition, the following literal finite inequalities wherever the
corresponding step is applied:
\[
 \begin{gathered}
 q\geq2^{16},\quad Y=2^{-16}q,\quad R\leq Y/4,\quad
 R^2/q^{7/4}\leq(5/4)2^{-28},\\
 K=q^{9/10}\geq kb,\quad
 T^{19/20}\geq2(k-1)b,\quad
 2T^{19/20}<T,\quad G\log T<\sqrt T/2,\\
 \epsilon_k\leq1/2,\quad
 \epsilon_{k-3}\leq e^{-1/2}/2,\quad k\geq V/b^2,\quad
 r\leq2q+1,\\
 \epsilon=\sqrt{\log(q+2)/q},\quad
 d/q<\epsilon/2,\quad 2\epsilon<1/8 ,
 \qquad q^{9/10}\leq T\leq128q .
 \end{gathered}                                      \tag{OR41}
\]
Here \(\epsilon_n\) is the original relative local-limit error
from the localization provider and \(r\) is the selected Taylor
rank in (OR22).  The Chernoff step also takes its explicit eventual
range \(C_hk/T+C_hk\log T/T^{19/20}\leq1/4\), enlarging the
fixed-source constants if needed.  The root-comparison appendix
retains its displayed exponentially small \(\varepsilon_q\leq1/2\).
Every guard holds eventually on each original fixed-multiplicity
sequence: \(q\asymp_h k^2\) for \(m=1\), \(q\asymp_h k^3\) for
\(m>1\), and \(R=O_h(k)\).  Thus these are finite-domain
specifications, not new assumptions replacing any missing estimate.

The fixed multiplicity, every root and repetition, both real
half-lines, the source-to-remainder map and its complete relation
Gram were used throughout.  The scalar \(A_k\) cancelled only at
the full quotient receiver (OR1).  The estimated complete
quotient determinant does not determine the seven absolute
allocations or projected arithmetic signs.  Those quantities
retain their original restricted metrics and same-class pairings.
The exact metric maps established here are (OR2), the Taylor
short sequence used in (OR23), and multiplication by \((y+i)^d\)
used in (OR27).



% END PRESERVED INPUT 02_OUTER_VALUE_RECONSTRUCTION.tex
\endgroup

\clearpage
\begingroup


\setcounter{equation}{0}
\renewcommand{\theequation}{RC.\arabic{equation}}
\section*{Finite comparison with the complete original root polynomial}
\addcontentsline{toc}{section}{Finite comparison with the complete original root polynomial}
% BEGIN PRESERVED INPUT 03_root_comparison.tex SHA256 6de05482c50493626cf1a9901e207488c78b73e7f706df1b30bcfe72bab5555f






\paragraph{Scope and sources.}
This calculation supplies the full polynomial quadratic-form comparison
used between O32 and the reference-power reduction in O31 of
\path{Signed_Arithmetic_Value_20260918/proof/OUTER_VALUE.md}.
It uses the original Gamma mass and lower bound RTM1--2 in that packet's
\path{sources/RTM_GTM_original_providers.tex}.
The Legendre normalization and recurrence are the classical ones recorded
in NIST DLMF, Table 18.3.1 and Section 18.9; their use and the resulting
finite extrapolation inequality are detailed below.
This note establishes the root-comparison step. It makes no assertion
that the separate reference-ensemble free-energy calculation has thereby
been proved.

\paragraph{Original objects and finite range.}
Retain the original fixed multiplicity $m\geq1$, actual quartet parameters
$0<\delta<1/2$, $\gamma>2$, and original orders $k\equiv1\pmod4$.
Write
\[
 r_h=\sqrt{\delta^2+\gamma^2},\qquad
 e=1+k(m-1),\qquad q=e(k+1)^2,\qquad R=kr_h,\qquad
 c=2^{-16},\qquad Y=cq.
\]
The original real even monic polynomial is
$Q(y)=i^{-q}\chi_k(k/2+iy)$, of exact degree $q$, with every
repetition retained and every root of modulus at most $R$.
In particular there are exactly $q/2$ opposite-root pairs.
Set
\[
 \sigma(y)=\frac{|\Gamma(1/4+iy/2)|^2}{2\pi},\quad
 c_\Gamma=\sqrt2 e^{-7/6},\quad M_\Gamma=\sqrt{2\pi},\quad
 B=\frac54\,2^{-28}.
\]
The exact mass and original lower bound are
\[
 \int_{\mathbb R}\sigma(y)\,dy=M_\Gamma,\qquad
 \sigma(y)\geq c_\Gamma e^{-\pi|y|/2}(1+|y|)^{-1/2}.
 \tag{RC1}
\]
For the explicit estimates below use the finite guard
\[
 q\geq2^{16},\qquad R\leq Y/4,\qquad
 \frac{R^2}{q^{7/4}}\leq B.                         \tag{RC2}
\]
This guard is attained on every original fixed packet; it is not an
extra asymptotic assumption. Indeed $q\geq(k+1)^2\geq k^2$, so
\[
 \frac RY\leq\frac{r_h}{ck},\qquad
 \frac{R^2}{q^{7/4}}\leq r_h^2 k^{-3/2}.
\]
Thus (RC2) follows from the completely explicit sufficient condition
\[
 k\geq
 \max\left\{256,\ 2^{18}r_h,\
                 (r_h^2/B)^{2/3}\right\},
 \qquad k\equiv1\pmod4.                           \tag{RC3}
\]
Any earlier finite guards imposed on the original packet remain in force.
The enlargement here only makes the cofinal range explicit.

\paragraph{Both full forms.}
For every real $t\in[-1/8,1/8]$ and every complex polynomial
$P\in\mathcal P_q$, define
\[
 A_t(P)=\int_{\mathbb R}|P(y)|^2|Q(y)|^2
                    (1+y^2)^{qt}\sigma(y)\,dy,\qquad
 B_t(P)=\int_{\mathbb R}|P(y)|^2|y|^{2q(1+t)}
                    \sigma(y)\,dy.               \tag{RC4}
\]
All integrals converge, and both forms are positive definite.
The compared linear map is the identity
$\mathcal P_q\longrightarrow\mathcal P_q$ in the original $y$ coordinate;
in the relation space it is the exact map
$P\mapsto QP$, with every coefficient and root of $Q$ retained.
The nonintegral reference power occurs only in the comparison metric.
No relation polynomial is replaced in the original form.

\begin{lemma}[Exterior comparison]
Define
\[
 H_{q,h}=\frac{qR^2}{Y^2-R^2}
                 +\frac q8\log(1+Y^{-2}).         \tag{RC5}
\]
For $|y|\geq Y$ and $|t|\leq1/8$ the quotient of the original weight
in $A_t$ by the reference weight in $B_t$ belongs to
$[e^{-H_{q,h}},e^{H_{q,h}}]$.
Moreover
\[
 H_{q,h}\leq
 c^{-2}\left(\frac{16}{15}r_h^2+\frac1{8q}\right).
 \tag{RC6}
\]
\end{lemma}
\begin{proof}
Write $Q(y)=\prod_{j=1}^{q/2}(y^2-z_j^2)$, where
$|z_j|\leq R$. For $|w|<1$ the absolutely convergent power series for
$\log(1-w)$ gives
$|\log|1-w||\leq |w|/(1-|w|)$.
Consequently
\[
 \left|\log\frac{|Q(y)|^2}{|y|^{2q}}\right|
 \leq 2\sum_{j=1}^{q/2}
       \frac{|z_j|^2}{y^2-|z_j|^2}
 \leq\frac{qR^2}{y^2-R^2}.
\]
The logarithm contributed by the other factor is exactly
$qt\log(1+y^{-2})$.
This proves (RC5) on each exterior half-line.
Since $R/Y\leq1/4$, $R^2/q\leq r_h^2$, and
$\log(1+Y^{-2})\leq Y^{-2}$, (RC6) follows.
\end{proof}

\begin{lemma}[Uniform weight bound on the retained inner interval]
For $|y|\leq Y$ and $|t|\leq1/8$,
\[
 |Q(y)|^2(1+y^2)^{qt}
       \leq B^q q^{2q(1+t)},\qquad
 |y|^{2q(1+t)}
       \leq B^q q^{2q(1+t)}.                     \tag{RC7}
\]
\end{lemma}
\begin{proof}
Opposite-root pairing gives $|Q(y)|^2\leq(y^2+R^2)^q$.
For $u=y^2$, the derivative of
$\log[(u+R^2)(1+u)^t]$ has the sign of
$(1+t)u+1+tR^2$. Its slope is strictly positive.
Thus its maximum on $[0,Y^2]$ occurs at an endpoint, for every allowed
$t$; any interior stationary point is a minimum.
At $u=0$, division by $q^{2(1+t)}$ gives at most
$R^2/q^{7/4}\leq B$.
At $u=Y^2$, the quotient equals
\[
 \frac{Y^2+R^2}{q^2}
              \left(\frac{1+Y^2}{q^2}\right)^t .
\]
Here $(1+Y^2)/q^2=q^{-2}+c^2<1$ by (RC2), so this expression is largest
at $t=-1/8$. Since $1+Y^2\geq Y^2$ and $R^2/Y^2\leq1/16$, its value is
at most $(17/16)c^{7/4}\leq(5/4)2^{-28}=B$.
Raising the two endpoint bounds to the $q$th power proves the first
claim. The reference bound follows from
$c^{2(1+t)}\leq c^{7/4}<B$.
\end{proof}

\begin{lemma}[Complete shifted-Legendre extrapolation]
Every $P\in\mathcal P_q$ and every $|y|\leq Y$ satisfy
\[
 |P(y)|^2\leq\frac{(q+1)^2}{q}\,100^q
                    \int_q^{2q}|P(u)|^2\,du.     \tag{RC8}
\]
\end{lemma}
\begin{proof}
Use the classical Legendre polynomials
$L_j(x)=(2^j j!)^{-1}(d/dx)^j(x^2-1)^j$.
Integration by parts $j$ times proves orthogonality to every polynomial
of degree below $j$, because $(x^2-1)^j$ has zeros of order $j$ at both
endpoints.
Their leading coefficient is $(2j)!/[2^j(j!)^2]$.
Applying the same integration by parts with $L_j$ itself gives
\[
 \int_{-1}^1L_j(x)^2\,dx
 =\frac{(2j)!}{2^{2j}(j!)^2}
                 \int_{-1}^1(1-x^2)^j\,dx
 =\frac2{2j+1};
\]
the middle integral is
$2^{2j+1}(j!)^2/(2j+1)!$, obtained by a beta integral or
successive integration by parts.
These are the normalizations in DLMF Table 18.3.1
\cite{rc:DLMF183}.
Their recurrence
$(j+1)L_{j+1}=(2j+1)xL_j-jL_{j-1}$
(DLMF Section 18.9 \cite{rc:DLMF189}) also follows directly by matching
the leading coefficient and using the just-proved orthogonality:
$xL_j$ has only $L_{j+1}$ and $L_{j-1}$ components by parity and
orthogonality, and their coefficients are
$(j+1)/(2j+1)$ and $j/(2j+1)$.
For $|x|\leq4$ it proves $|L_j(x)|\leq10^j$ by induction:
$L_0=1$, $|L_1|\leq4$, and the next value is at most
$8\,10^j+10^{j-1}\leq10^{j+1}$.
The functions
\[
 \ell_j(u)=\sqrt{\frac{2j+1}{q}}\,L_j(2u/q-3),
 \qquad 0\leq j\leq q,
\]
are therefore a complete orthonormal basis for $\mathcal P_q$ on
$[q,2q]$. If $|y|\leq Y$, then $|2y/q-3|<4$.
Expanding $P$ in that entire basis and applying Cauchy--Schwarz gives
\[
 |P(y)|^2\leq
 \left(\int_q^{2q}|P(u)|^2du\right)
 \frac1q\sum_{j=0}^q(2j+1)100^j.
\]
Finally the sum is at most $100^q(q+1)^2$.
All polynomial directions, including the endpoint degree $q$, were
included in this calculation.
\end{proof}

\paragraph{The two inner integrals, with their original mass.}
For nonzero $P$ the original Gamma lower bound on $[q,2q]$ gives
\[
 B_t(P)\geq
 \frac{c_\Gamma e^{-\pi q}}{\sqrt{1+2q}}\,
 q^{2q(1+t)}\int_q^{2q}|P(u)|^2\,du.              \tag{RC9}
\]
Here $2q(1+t)>0$, so its power has its minimum at $u=q$.
Combining (RC7)--(RC9) with the entire mass $M_\Gamma$ proves
\[
 \begin{split}
 \int_{-Y}^Y|P(y)|^2|Q(y)|^2(1+y^2)^{qt}\sigma(y)dy
       &\leq\varepsilon_q B_t(P),\\
 \int_{-Y}^Y|P(y)|^2|y|^{2q(1+t)}\sigma(y)dy
       &\leq\varepsilon_q B_t(P),
 \end{split}                                    \tag{RC10}
\]
where the same completely explicit number works for both:
\[
 \varepsilon_q=
 \frac{\sqrt{2\pi}}{c_\Gamma}
 \frac{(q+1)^2\sqrt{1+2q}}{q}
          \left(100e^\pi B\right)^q.             \tag{RC11}
\]
In particular $100e^\pi B<1/16$ and
$\sqrt{2\pi}/c_\Gamma<6$, so, for $q\geq1$,
$\varepsilon_q\leq42q^2 16^{-q}$.
This is less than $1/2$ throughout (RC2).
Thus both inner intervals have been quantitatively retained.

\begin{theorem}[Finite forms and full determinant return]
Throughout (RC2), uniformly in $|t|\leq1/8$ and
$P\in\mathcal P_q$,
\[
 e^{-H_{q,h}}(1-\varepsilon_q)B_t(P)
 \leq A_t(P)\leq(e^{H_{q,h}}+\varepsilon_q)B_t(P).
 \tag{RC12}
\]
For every fixed linearly independent frame $P_1,\ldots,P_n$
in $\mathcal P_q$, the two full Gram matrices satisfy
\[
 \left|\log
 \frac{\det\operatorname{Gram}_{A_t}(P_1,\ldots,P_n)}
      {\det\operatorname{Gram}_{B_t}(P_1,\ldots,P_n)}
 \right|
 \leq n\{H_{q,h}-\log(1-\varepsilon_q)\}.
 \tag{RC13}
\]
In particular this is $O_h(q)$ for the original frames
$1,y,\ldots,y^{n-1}$ with $n=q$ and $n=q+1$.
\end{theorem}
\begin{proof}
On the union of the two exterior half-lines use the pointwise bounds
(RC5). For the lower bound, discard only the nonnegative original inner
integral and retain the estimate
$B_t(\{|y|\geq Y\})\geq(1-\varepsilon_q)B_t(P)$.
For the upper bound, add the first inequality of (RC10) to the exterior
upper bound. This proves (RC12), without changing either full norm.

Since (RC12) holds for every linear combination of the frame, it is a
Loewner comparison of the complete Gram matrices, including all cross
blocks. Conjugate by the positive inverse square root of the reference
Gram. Every eigenvalue of the resulting positive matrix lies in the
interval from (RC12); their product is the displayed determinant ratio.
The lower endpoint has logarithm
$-H_{q,h}+\log(1-\varepsilon_q)$.
The upper logarithm is
$H_{q,h}+\log(1+\varepsilon_q e^{-H_{q,h}})$, which is at most
$H_{q,h}-\log(1-\varepsilon_q)$.
This proves (RC13).
The bound (RC6) and exponentially decreasing (RC11) give the claimed
uniform $O_h(q)$, with its finite constant and domain displayed.
\end{proof}

\paragraph{Exact consequence for O30--O31.}
Let $D^{\rm ref}_{k,n}(t)$ denote the full Gram determinant on
$1,y,\ldots,y^{n-1}$ for the second weight in (RC4), and put
\[
 F^{\rm ref}_{k,n}(t)=q^{-2}
  \{\log D^{\rm ref}_{k,n}(t)-\log D^{\rm ref}_{k,n}(0)
                                   -2qnt\log q\}.
\]
The determinant denoted $D_{k,n}(t)$ in O30 is exactly the
$A_t$ determinant, because its columns are $QP_j$ in the original
weighted Gamma norm.
Applying (RC13) twice, at $t$ and at $0$, proves
\[
 |F_{k,n}(t)-F^{\rm ref}_{k,n}(t)|
 \leq\frac{2n}{q^2}
          \{H_{q,h}-\log(1-\varepsilon_q)\}
 =O_h(q^{-1}).                                  \tag{RC14}
\]
This uniform finite bound proves precisely the root-comparison
reduction required in O31, including its $q^2$ scaling.
The zero-endpoint sentence in O32 is valid on a cofinal range, but the
finite guard $R^2/q^{7/4}\leq B$ must be stated to use the displayed
uniform endpoint constant. The $Y$-endpoint constant
$[(5/4)2^{-28}]^q$ is valid, as proved above.



% END PRESERVED INPUT 03_root_comparison.tex
\endgroup

\clearpage
\begingroup
\def\im{\operatorname{im}}
\def\diag{\operatorname{diag}}
\def\Gr{\operatorname{Gram}}
\def\C{\mathbb C}

\setcounter{equation}{0}
\renewcommand{\theequation}{PEQ\arabic{equation}}
\section*{Primitive-endpoint control and exact CPQ integration}
\addcontentsline{toc}{section}{Primitive-endpoint control and exact CPQ integration}
% BEGIN PRESERVED INPUT 04_PRIMITIVE_ENDPOINT_INTEGRATION.tex SHA256 276f46d56c1c672a83b5d25018c6e6fe09537a31e20dad417f06607ef0738289


\begin{abstract}
The original PRT comparison between primitive endpoints has three quotient directions. Its preceding finite estimate had logarithmic width $O(q\log a)$. The original full-remainder AKS estimate and the supplied LRC conductor bounds imply a native Loewner growth bound of width $O_{\rm actual}(q)$ before any fixed flag restriction or attained quotient. The three original PRT returns are therefore $O(q)$, as is the difference of the original common losses. In particular the primitive-endpoint difference has zero $q\log a$ coefficient, with an $O(q)$ remainder, on every retained rank stratum. This resolves that specific logarithmic ambiguity. It does not evaluate either absolute directional allocation, the independent proper-source standard return, or the retained-class projected signs; the requested directional completion target remains open.
\end{abstract}

The programme origin and its existing attribution are retained
\cite{peq:human:globalization2026,peq:human:splitzero2026}. This additive
integration retains the complete supplied calculation below and appends
the exact provider correspondence and a sharper finite endpoint receiver.
The original supplied bytes remain in the accompanying input directory.

\section{Original domain and scope}
Fix the stipulated simple quartet, its actual full unit, the original five-orbit period and its logarithm branch. Retain
\begin{equation}
\begin{gathered}
 a=k+4\equiv1\pmod4,\quad k\ge9,\quad q=(a+1)^2,\quad Q=(a-7)^2,\\
 \Delta=16a-48,\quad m=\Delta-v,\quad s\in\{1,a\},\quad
 c=a/2,\quad c'=a/2-4,\\
 E_A(z)=\sum_{u,w=0}^8a_{uw}e^{\beta_{8,u,w}z},\quad
 \mu_t=E_A^{(t)}(0),\quad v=\operatorname{ord}_0E_A\le80,\\
 \mu_v\ne0,\quad \alpha_A=v!/\mu_v,\qquad
 \beta_{a,u,w}=c+(2u-a)\delta+i(2w-a)\gamma,\\
 0<\delta<1/2,\qquad\gamma>2.
\end{gathered}\label{peq:peq:data}
\end{equation}
Here $m$ is a coefficient dimension, not a zero multiplicity. The lower grid has degree $a-8$ and source order $s$, not order $a-8$. The independent proper source keeps degree $k$, its own centre and orders $1,k$. All original five-orbit and PRT coefficient conditions remain. Neither a new large-root separation guard nor a numerical period choice is introduced by the bound below. The constants depend on the fixed actual conductor; joint degree/period uniformity is not asserted.

Use exactly
\begin{equation}
\begin{gathered}
 dm_s(y)=\frac{(2\pi)^{s/2}2^{s/2}}{4\pi\Gamma(s/2)}
 |\Gamma(s/4+iy/2)|^2dy,\quad M_s=(2\pi)^{s/2},\\
 h_{n,s}=M_sn!(s/2)_n,\quad
 \phi_{n,s,c}(S)=p_{n,s}((S-c)/i)/\sqrt{h_{n,s}},\\
 p_{n+1,s}=yp_{n,s}-n(n-1+s/2)p_{n-1,s},\quad p_0=1,\quad p_1=y.
\end{gathered}\label{peq:peq:gamma}
\end{equation}
These classical orthogonal-polynomial formulae retain the exact
Meixner--Pollaczek dictionary
$p_{n,s}(y)=n!P_n^{(s/4)}(y/2;\pi/2)$
\cite[eqs.~(18.19.8)--(18.19.9), (18.23.7)]{peq:human:dlmf18}.
In particular the change $y=2x$ keeps the full source mass.
Write $D_N=D_{a,N-v}^{(s)}$ for the original divided-native metric at $N\ge q-1$, and $\mathcal D_N=\diag(D_N,4)$. This abbreviation retains the cutoff shift $v$.

The new evaluated logarithmic conclusion is
\begin{equation}
 \boxed{\mathcal L_0-\mathcal L_1=O_{\rm actual}(q),\qquad
 \frac{\mathcal L_0-\mathcal L_1}{q\log a}\longrightarrow0.}
 \label{peq:peq:main}
\end{equation}
No existence of an asymptotic expansion for either absolute loss is assumed. The same conclusion applies to the PRT-related complementary differences specified below. Every result is restricted to the simple-quartet metric lane.

\section{A full-numerator estimate in the original coordinates}
Set $R=a\sqrt{\delta^2+\gamma^2}$ and $b=s/2$. The exact AKS3 envelope at degree $a$ is
\begin{equation}
\begin{gathered}
 t_q=\pi/2-1/q,\quad A_q=2\log3/t_q,\quad
 S_q=\sum_{j=0}^{q-1}A_q^j,\\
 B_{s,q}=[(9/25)\sin(1/q)]^{-b}
              \exp(2R(\log3+t_q/2)),\\
 Z_{a,s}=1+B_{s,q}S_q^2\sum_{n=q}^{2q}
                \frac{n!}{(b)_n}(25/16)^n.
\end{gathered}\label{peq:peq:Z}
\end{equation}
These are scalar envelopes, not boundary subspaces or coefficient frames. The original $q\ge36$ requirement is satisfied in \eqref{peq:peq:data}.

Let $\mathsf E_N[n,\lambda]=\phi_{n,s,c}(\beta_\lambda)$ be the literal evaluation matrix on the full original grid, and $K_N=\mathsf E_N^*\mathsf E_N$. Let $v_j$ be the remainder-coefficient column of $\phi_{q+j,s,c}$ in the low original orthonormal frame, and $E_0=\mathsf E_{q-1}$. The added evaluation row is $v_j^{\mathsf T}E_0$, so
\begin{equation}
 K_N=E_0^*\overline{\left(I+\sum_{j=0}^{N-q}v_jv_j^*\right)}E_0.
 \label{peq:peq:numgram}
\end{equation}
The ordinary transpose and the subsequent conjugation are required. The full original AKS18--24 calculation gives
\[
 I\preceq I+\sum_{j=0}^{N-q}v_jv_j^*\preceq Z_{a,s}I.
\]
Consequently
\begin{equation}
 \boxed{K_{q-1}\preceq K_N\preceq Z_{a,s}K_{q-1},
                  \quad q-1\le N\le2q.}\label{peq:peq:num-bound}
\end{equation}
This imports the full remainder operator bound, not an observation-rank estimate.

For its normalization, AKS applies the literal remainder map to the original generating identity
\[
 \sum_{n\ge0}p_{n,s}(y)z^n/n!
       =(1+z^2)^{-s/4}e^{y\arctan z}.
\]
The divided-difference integral at the complete root list proves
$|[\omega_1,\ldots,\omega_{j+1}]e^{v\cdot}|\le |v|^je^{R|v|}/j!$.
The whole Newton polynomial, expanded with every coefficient, satisfies
\[
 \left\|\prod_{i=1}^j(y-\omega_i)\right\|_s
 \le\sqrt{M_s}(\cos t)^{-s/4}(2/t)^j j!e^{Rt/2}.
\]
Cauchy extraction on $|z|=4/5$ and $t=t_q$, followed by the exact normalization $h_{n,s}$, gives
\[
 \|v_{n-q}\|^2\le
 B_{s,q}S_q^2\frac{n!}{(s/2)_n}(25/16)^n.
\]
Adding the complete positive rank-one matrices through $n=2q$ yields \eqref{peq:peq:Z}. The original mass cancels only between this squared source norm and its matching $h_{n,s}$.

The explicit envelope has size
\begin{equation}
 \log Z_{a,s}=O_h(q),\qquad
 \limsup_{a\to\infty}\frac{\log Z_{a,s}}q
 \le2\log\frac{25\log3}{4\pi}.
 \label{peq:peq:Z-size}
\end{equation}
Indeed $A_q\to4\log3/\pi>1$, hence $\log S_q=q\log A_q+O(1)$. For $s\ge1$,
\[
 \frac{n!}{(s/2)_n}\le\frac{n!}{(1/2)_n}
 =\frac{4^n}{\binom{2n}{n}}\le2n+1.
\]
The sum in \eqref{peq:peq:Z} is at most $(q+1)(4q+1)(25/16)^{2q}$. Finally $\log B_{s,q}=O_h(a\log q)=o(q)$ in both original orders. This proves \eqref{peq:peq:Z-size}, a bound on the scalar envelope, not an individual directional eigenvalue.

Use the original numerator matrices $C_a,Z_a$ of NCT/NPK/IGI. Their attained form is
\begin{equation}
\begin{split}
 \widehat D_N=Z_a^*\bigl[ K_N-K_NC_a^{\mathsf T}
 (\overline C_aK_NC_a^{\mathsf T})^{-1}
                  \overline C_aK_N\bigr]Z_a.
\end{split}\label{peq:peq:numerator-min}
\end{equation}
Equivalently its quadratic form is the minimum of
$(Z_ax+C_a^{\mathsf T}\eta)^*K_N(Z_ax+C_a^{\mathsf T}\eta)$ over the entire lower coefficient vector $\eta$. Applying \eqref{peq:peq:num-bound} to every such vector and then minimizing proves
\begin{equation}
 \widehat D_{q-1}\preceq\widehat D_N
                  \preceq Z_{a,s}\widehat D_{q-1}.
 \label{peq:peq:numerator-order}
\end{equation}
In particular neither of the two cross factors or the old lower inverse in \eqref{peq:peq:numerator-min} has been removed.

\section{Transport the entire bound through the original conductor}
At $D=N-v$ the original conductor is the square map in the original native frames,
\[
 A_{D,s}:\mathcal P_N/\mathcal P_{v-1}\longrightarrow\mathcal P_D,
 \qquad[P]\longmapsto\mathcal T_AP.
\]
Its exact WCF/LRC matrix is
\begin{equation}
\begin{gathered}
 A_{D,s}=\diag(\rho_i)^{-1}T_D(g)\diag(\rho_{i+v}),\qquad
 \rho_n^2=n!/(s/2)_n,\\
 g(t)=t^{-v}e^{-4w(t)}E_A(w(t)),\quad w(t)=-i\arctan t,\qquad
 g(0)=(-i)^v\mu_v/v!.
\end{gathered}\label{peq:peq:conductor}
\end{equation}
The nilpotent Toeplitz matrix uses its literal cutoff $D$. No interior zero, pole direction, full unit value or branch factor is replaced. Both source centres are the original ones in \eqref{peq:peq:data}.

Here are the complete finite LRC6--9 constants used for the transport. Define
\begin{equation}
\begin{gathered}
 A_{\rm abs}=\sum_{u,w}|a_{uw}|,\quad
 F_0=4+4\sqrt{\delta^2+\gamma^2},\quad
 B_0=2^{v+8}A_{\rm abs}e^{2\pi\gamma},\\
 B_2=B_0\{F_0(F_0+1)+4vF_0+4v(v+1)\},\qquad
 H=B_0(1+\sqrt3)+2\sqrt{10}B_2,\\
 H_s=\begin{cases}\max(1,H\sqrt{2v+1}),&s=1,\\
                 \max(1,H),&s=a,\end{cases}
 \qquad
 \kappa_s=\begin{cases}1+\sqrt{8/3},&s=1,\\2,&s=a.\end{cases}
\end{gathered}\label{peq:peq:H}
\end{equation}
Write $d=D+1$ and $\mathsf H_d=\sum_{n=1}^d1/n$. For $1\le r\le bd$ and $b=1$ or $4$, put
\begin{equation}
\begin{gathered}
 F_{D,s,r,b}=r\log H_s+4(bd-r)\log\kappa_s,\qquad F_{D,s,0,b}=0,\\
 J_{D,s}=d\log^+\frac1{|g(0)|}
       +\tfrac12(s/2-1)_+v\mathsf H_d,\\
 I_{D,s,r,b}=bJ_{D,s}+F_{D,s,bd-r,b},\qquad I_{D,s,0,b}=0.
\end{gathered}\label{peq:peq:FI}
\end{equation}
Here $b$ is a copy count only in \eqref{peq:peq:FI}; it is not the $s/2$ abbreviation in \eqref{peq:peq:Z}. LRC proves the exact original exterior bounds
\[
 \log\|\wedge^r(A_{D,s}^{\mathsf T})^{\oplus b}\|\le F_{D,s,r,b},
 \qquad
 \log\|\wedge^r(A_{D,s}^{-\mathsf T})^{\oplus b}\|\le I_{D,s,r,b}.
\]
Only rank one and one copy are needed below. In the admitted cutoff range $d\ge2$, the explicit rank-one inverse bound is nondecreasing in $D$.

For the exact map on a complex-linear functional, multiplying its germ by $E_A$ sends its native value column $f$ to $A_{D,s}^{\mathsf T}f$, with zero first $v$ coordinates in the upper polynomial source. The original lower exponential columns are carried by this same map to $C_a^{\mathsf T}$-numerator columns. Thus the map bijects the full affine lower fibres for $D_N$ and $\widehat D_N$. On every representative,
\[
 e^{-2I_{D,s,1,1}}\|f\|^2\le
      \|A_{D,s}^{\mathsf T}f\|^2
       \le e^{2F_{D,s,1,1}}\|f\|^2.
\]
Taking the two attained minima gives
\begin{equation}
 \boxed{e^{-2I_{N-v,s,1,1}}D_N
 \preceq\widehat D_N\preceq e^{2F_{N-v,s,1,1}}D_N.}
 \label{peq:peq:native-transfer}
\end{equation}
The transpose is the original ordinary dual transpose. Its singular values equal those of $A_{D,s}$; no bilinear residue pairing is substituted for a Hermitian norm. The actual coefficient numerator is $Z_ax$, so the equality of value coordinates is the NCT/NPK square, not an independently assigned scale. In particular $\alpha_A$ remains in \eqref{peq:peq:conductor} through $g(0)$.

\begin{theorem}[Native growth before any flag compression]
Set
\begin{equation}
 D_-=q-v-1,\qquad D_+=2q-v,\qquad
 \boxed{\mathfrak E_{a,s}
 =\log Z_{a,s}+2I_{D_+,s,1,1}+2F_{D_-,s,1,1}.}
 \label{peq:peq:E}
\end{equation}
Then for every original $q-1\le N\le2q$,
\begin{equation}
 \boxed{D_{q-1}\preceq D_N\preceq
 e^{\mathfrak E_{a,s}}D_{q-1},\qquad
 \mathfrak E_{a,s}=O_{\rm actual}(q).}
 \label{peq:peq:native-bound}
\end{equation}
The same scalar bounds hold on the original four-copy metric, and on every fixed restriction or attained quotient of either metric, in its original value frame.
\end{theorem}
\begin{proof}
Native row extension adds a positive rank-one form after the entire old lower minimization; hence $D_N\succeq D_{q-1}$. The upper estimate follows from the complete chain
\[
 D_N\preceq e^{2I_{D_+,s,1,1}}\widehat D_N
 \preceq e^{2I_{D_+,s,1,1}}Z_{a,s}\widehat D_{q-1}
 \preceq e^{\mathfrak E_{a,s}}D_{q-1}.
\]
No two unrelated minimum sections have been identified in this chain.
For the size, \eqref{peq:peq:Z-size} controls the first term. All constant conductor data in \eqref{peq:peq:H} are fixed. The terms linear in $D$ in \eqref{peq:peq:FI} are $O_{\rm actual}(q)$, and its harmonic term is $O_{\rm actual}(av\log q)=o(q)$ since $v$ is fixed. This proves the asserted order for both original source orders.

Taking four copies preserves the scalar Loewner inequality without multiplying its logarithm by four. Restriction follows by substitution in a fixed frame. For a fixed quotient $Y/X$, apply the inequalities to every representative of the same affine fibre and minimize. The resulting quotient Gram is exactly
\[
 X_Q^*\mathcal D_NX_Q-X_Q^*\mathcal D_NX
            (X^*\mathcal D_NX)^{-1}X^*\mathcal D_NX_Q,
\]
with both cross factors and the entire denominator retained. The scalar inequality and monotonicity therefore pass to that form, irrespective of the condition number of the chosen coefficient frame.
\end{proof}

On the PZU domain $v=0$, $D_-=q-1,D_+=2q$, and $g(0)=E_A(0)=1/\alpha_A$. These are exact simplifications; $\alpha_A$ is not set to one. Its possibly large fixed-period magnitude remains inside $J_{D,s}$. At an admitted period with $v>0$, no low coordinate or derivative condition has been removed. The bound is not uniform in varying periods.

\section{Evaluate the primitive-endpoint difference above the quadratic scale}
Retain the PRT/CPQ coefficient subspaces
\begin{equation}
\begin{gathered}
 \mathscr B=\im B_{\rm cof},\quad \mathbb J=\ker\mathbb B,\quad
 W=\mathscr B+\mathbb J,\\
 S_e=\im[B_{\rm cof},F_e],\quad K_e=S_e\cap W,\quad
 T_e=S_e+W\qquad(e=0,1).
\end{gathered}\label{peq:peq:spaces}
\end{equation}
The defining primitive endpoint $e$ is distinct from the later variable norm cutoff $N$. The original PRT6--10 matrix $L$ has rank exactly three and gives the literal identity $S_1=\mathscr B+F_0\ker L\subset S_0$. With the original three lifts $F_0Y$ it proves the exact sequence
\begin{equation}
 0\longrightarrow K_0/K_1\longrightarrow S_0/S_1
            \longrightarrow T_0/T_1\longrightarrow0.
 \label{peq:peq:sequence}
\end{equation}
Write $r_\partial=\dim(T_0/T_1)\in\{0,1,2,3\}$. Then the three ranks are respectively $3-r_\partial,3,r_\partial$. This retains the exact rank of the original three-column image, which may depend on the packet, period and degree; no generic rank is assigned.

Give these spaces their full PRT11--12 attained Grams
\[
 G_A(N)=\Gr_{\mathcal D_N}(S_0/S_1),\quad
 G_B(N)=\Gr_{\mathcal D_N}(K_0/K_1),\quad
 G_C(N)=\Gr_{\mathcal D_N}(T_0/T_1).
\]
The original frame factors in the Gram factorizations are fixed in $N$. They have not been normalized to one.

\begin{lemma}[Fixed-rank quotient return]
For any of the fixed restrictions or quotients in the preceding theorem, let $G_X(N)$ be its attained form of rank $r_X$. With the unchanged return
\[
 \mathcal R f=-f(q-1)-f(q)+f(2q-1)+f(2q),
\]
one has the finite inequality
\begin{equation}
 \boxed{0\le\mathcal R\log\det G_X\le2r_X\mathfrak E_{a,s}.}
 \label{peq:peq:quotreturn}
\end{equation}
\end{lemma}
\begin{proof}
Put $f_X(N)=\log\det G_X(N)-\log\det G_X(q-1)$. The full scalar Loewner bounds, positive congruence and multiplication of the $r_X$ eigenvalues give
$0\le f_X(N)\le r_X\mathfrak E_{a,s}$. Minimum monotonicity makes $f_X$ increasing. Thus
\[
 \mathcal R\log\det G_X
       =[f_X(2q-1)-f_X(q-1)]+[f_X(2q)-f_X(q)]
\]
is nonnegative and at most $2r_X\mathfrak E_{a,s}$.
For $r_X=0$, the return is exactly zero by the original empty determinant convention.
\end{proof}

Set $U_X=\mathcal R\log\det G_X$ for $X=A,B,C$. Substitution of their actual ranks yields
\begin{equation}
\boxed{\begin{gathered}
 0\le U_A\le6\mathfrak E_{a,s},\qquad
 0\le U_B\le2(3-r_\partial)\mathfrak E_{a,s},\qquad
 0\le U_C\le2r_\partial\mathfrak E_{a,s}.
\end{gathered}}\label{peq:peq:threebounds}
\end{equation}
These are new $O(q)$ bounds for the existing three attained matrices. They do not replace those matrices by three arbitrary rank-one directions.

PRT14--17 already proves the exact identities
\begin{equation}
\begin{aligned}
 V_{{\rm conn},0}-V_{{\rm conn},1}&=U_A,&
 V_{h,0}-V_{h,1}&=U_B,\\
 V_{s,0}-V_{s,1}&=U_A-U_B,&
 V_{t,0}-V_{t,1}&=U_C,\\
 V_{{\rm rem},1}-V_{{\rm rem},0}&=U_C,&
 \mathcal L_0-\mathcal L_1&=U_C-U_A+U_B.
\end{aligned}\label{peq:peq:PRT}
\end{equation}
Their fixed frame constants cancel only after the four signs are taken. Applying \eqref{peq:peq:threebounds} to these exact identities gives
\begin{theorem}[No primitive-choice logarithmic coefficient]
On the original domain \eqref{peq:peq:data},
\begin{equation}
\boxed{\begin{gathered}
 |\mathcal L_0-\mathcal L_1|\le6\mathfrak E_{a,s}=O_{\rm actual}(q),\\
 0\le V_{{\rm rem},1}-V_{{\rm rem},0}
                   \le2r_\partial\mathfrak E_{a,s},\\
 -2(3-r_\partial)\mathfrak E_{a,s}
 \le V_{s,0}-V_{s,1}\le6\mathfrak E_{a,s}.
\end{gathered}}\label{peq:peq:result}
\end{equation}
Every difference in \eqref{peq:peq:PRT} divided by $q\log a$ tends to zero. The $aq$ and every other scale strictly larger than $q$ are unchanged by the primitive endpoint. This assertion does not assume absence of intermediate scales, a particular arithmetic subsequence, or convergence of either absolute directional return.
\end{theorem}
\begin{proof}
The two positive terms $U_B+U_C$ in the loss difference have total rank three by \eqref{peq:peq:sequence}, and hence sum at most $6\mathfrak E_{a,s}$. The negative term $-U_A$ is at least $-6\mathfrak E_{a,s}$. This proves the first bound, rather than charging three unrelated rank-three errors. The other inequalities follow from their indicated differences. Since $\mathfrak E_{a,s}=O(q)$ and $\log a\to\infty$, all the normalized conclusions follow. They are uniform over the four possible ranks at each fixed period. At $r_\partial=0$, $U_C=0$ exactly; at $r_\partial=3$, $U_B=0$ exactly.
\end{proof}

The old PRT18--20 finite envelope $\log A^{\rm end}_{a,s}$ remains valid. In every previous directed interval replace it by
\begin{equation}
 \mathfrak E^{\rm best}_{a,s}
 =\min\{\log A^{\rm end}_{a,s},\mathfrak E_{a,s}\}.
 \label{peq:peq:intersection}
\end{equation}
Thus the new interval is an intersection, never a weakening of an earlier finite one. The order $O(q)$ in the new result follows from the explicit new term.

\section{Correlated receiving identities and the uncompleted terminal target}
Let $Z_e=V_b+V_{{\rm rem},e}+V_{s,e}+V_A$. Here $Z_e$ is precisely the original complementary return sum, not the scalar $Z_{a,s}$ in \eqref{peq:peq:Z}. PRT gives, with the same original maps,
\begin{equation}
\begin{gathered}
 Z_1-Z_0=\mathcal L_0-\mathcal L_1=U_C-U_A+U_B,\\
 \mathcal P^{\rm rest}_0-\mathcal P^{\rm rest}_1=U_C-U_A+U_B,\\
 4S_K=\mathcal L_0+Z_0=\mathcal L_1+Z_1,\qquad
 S_K+S_R=\mathcal T.
\end{gathered}\label{peq:peq:coupling}
\end{equation}
The $V_b,V_A,V_I$ terms do not depend on the primitive endpoint, since their coefficient subspaces and frames are unchanged. The two labelled occurrences in a signed flag expression remain labelled even when the underlying subspaces coincide. The independent degree-$k$ proper-source form has not been replaced by a target-order form, and its standard-window value has not been evaluated by \eqref{peq:peq:coupling}. The known prefix correction remains $-V_A^{\rm prefix}$.

The new estimate resolves the former possibility that switching between the two primitive constructions could change a $q\log a$ coefficient. It permits either construction to be used for coefficients above the $q$ scale, with the exact PRT map and the $O(q)$ error recorded. It does not determine those coefficients. In particular no rank-proportional allocation is implied by \eqref{peq:peq:quotreturn}.

The prescribed signed profile is still the clipped sum over the actual angular data, with the original eight labels and zeros. CRV's $O(q)$ radial error and UIT's $O(q)$ attained/trial error are not replaced or claimed as new. The present native Loewner ratio is $e^{O(q)}$, not a small relative error. It therefore cannot be used to infer a retained-class pairing at the arithmetic $k\delta F_{33}$ scale.

For the coordinator's exact retained-class data, keep
\[
 U=K_{13}/F_{13},\quad V=K_{23}/F_{23},\quad
 P_K=\bar UV,\quad P_B=(1-\bar U)(1-V),\quad
 P_\times=\bar U+V-2\bar UV.
\]
Its already evaluated full current implies
\[
 \frac{\Phi_X}{2F_{33}}
 =k[\delta\Re P_X-(-1)^r\gamma\Im P_X]
       +(-1)^r\varepsilon_{N,\lambda}\Im P_X.
\]
The same positive $\varepsilon_{N,\lambda}$ occurs in all three terms, and each absolute replacement error still includes $|\Im P_X|$. No upper bound for that amplification is supplied here. Both leakage blocks remain in
$(\Lambda AI)^*Q_N+H_NA_{KB}$; neither is identified with their sum.

Finally $J=\mathcal B+\mathcal W-\mathcal P$ retains its complete original cancellation and its positive $C_Bq^2$ baseline. Since \eqref{peq:peq:coupling} changes no allocation sum, the $O(q)$ primitive comparison creates no new negative term in that action. The closing criterion is neither proved nor disproved by this partial result. The outstanding terminal targets are the absolute original angular clipped sums, the independent proper-source standard return, and the individual retained-class projected signs at their stipulated scale.

% Additive integration and sharper finite endpoint receiver.
% PEQ intake, 19 September 2026. All original objects are retained.
\section{Exact provider correspondence and finite endpoint refinement}
\label{peq:peq:intake}

The incoming calculation is received into the frozen LRC transport
proof, LT1--28 and XIM1--23, and the original PRT1--20/CPQ
construction. Their source locations and exact byte pins accompany this
manuscript in \texttt{SOURCE\_PINS.json}. The domain is the fixed actual
simple quartet, the original five-orbit coefficient construction and
logarithm branch, $k=4\ell+1\ge9$, $a=k+4$, and $s=1,a$.
Every original coefficient condition needed to define $B_{\rm cof}$,
$\mathbb B$ and $F_e$ is retained. The estimates here add no root
separation condition. If these results are used alongside a theorem
requiring a larger degree or a quantitative root guard, that theorem's
guard remains in force.

There are two distinct uses of the letter $a$ and two distinct uses of
$Q$ in the old PRT notation. The exact substitution into this manuscript is
\[
 \begin{array}{c|c}
 \text{PRT notation}&\text{notation here}\\ \hline
 j=k+4&a\\
 Q=(j+1)^2&q=(a+1)^2\\
 q=(k+1)^2&(a-3)^2\\
 a=\dim(T_0/T_1)&r_\partial\\
 D=N-v&D=N-v\\
 c'_j=j/2-4&c'=a/2-4
 \end{array}                                                     \tag{PEI1}
\]
In particular PRT's proper-source size $(k+1)^2$ is not this
manuscript's lower-grid size $Q=(a-7)^2$. The coefficient spaces and
linear maps are preserved by the identity map on their original
coefficients; PEI1 merely specifies the substitutions in their formulae.
The proper-source orders remain $1,k$, and the target orders remain $1,a$.

\subsection{The full affine-fibre square at every original cutoff}

Here is the specific typed map used in the native inequality.
Let $\mathscr N_v$ be the subspace of numerator coefficient columns
$\nu\in\mathbb C^q$ for which
$\sum_\beta\nu_\beta\beta^r=0$ for $0\le r<v$.
The injective map $C_a^{\mathsf T}:\mathbb C^Q\to\mathscr N_v$
is defined by the coefficient identity
\[
 \sum_\beta(C_a^{\mathsf T}\eta)_\beta e^{\beta z}
   =E_A(z)\sum_{\beta'}\eta_{\beta'}e^{\beta' z}.
                                                               \tag{PEI2}
\]
The roots on the right are the full degree-$a-8$ grid. Their addition
to every degree-eight conductor root gives the original degree-$a$
root, so PEI2 retains every convolution coefficient. If its left side
vanishes, the lower exponential sum vanishes wherever $E_A\ne0$ and
then identically. Its Vandermonde is nonsingular, so $\eta=0$.
The first $v$ moments vanish because $E_A$ vanishes to order $v$.

Let $Z_a:\mathbb C^m\to\mathscr N_v$ be the original complement,
so $[C_a^{\mathsf T},Z_a]$ is the fixed numerator frame. For its
column $b$ define
$F_b(z)=\sum_\beta (Z_a)_{\beta b}e^{\beta z}/E_A(z)$,
an analytic germ at zero. At $D=N-v$ put
\[
 \begin{split}
 (H_D)_{n,\beta'}&=\phi_{n,s,c'}(\beta'),\\
 (J_D)_{n,b}&=\Lambda_{F_b}(\phi_{n,s,c'}),\qquad
      \Lambda_F(S'^r)=F^{(r)}(0),\qquad 0\le n\le D,\\
 B_D&=A_{D,s}^{\mathsf T}:\mathbb C^{D+1}\longrightarrow
                         \mathbb C^{D+1}.
 \end{split}                                                     \tag{PEI3}
\]
The source and target of $B_D$ are respectively the complex-linear
functional value coordinates in the original lower polynomial frame
and the upper frame with its first $v$ zero coordinates removed.
It is an isomorphism since the diagonal of $A_{D,s}$ is
$g(0)\rho_{i+v}/\rho_i\ne0$.
Leibniz's rule gives
$\Lambda_{F_b}(\mathcal T_AP)=\Lambda_{E_AF_b}(P)$.
The same equation holds for each lower exponential functional.
Applying it to every upper basis polynomial proves the literal square
\[
 \mathsf E_N[Z_a,C_a^{\mathsf T}]
   =\begin{bmatrix}
       0_{v\times(m+Q)}\\
       B_D[J_D,H_D]
     \end{bmatrix}.                                             \tag{PEI4}
\]
This is XIM12, including its ordinary transpose, both centres and the
cutoff shift. Every lower coefficient $\eta\in\mathbb C^Q$ on one side
is carried to that identical coefficient on the other side. Thus
\[
 \begin{split}
 x^*D_Nx&=\min_{\eta\in\mathbb C^Q}
                          \|J_Dx+H_D\eta\|_2^2,\\
 x^*\widehat D_Nx&=\min_{\eta\in\mathbb C^Q}
                          \|B_D(J_Dx+H_D\eta)\|_2^2\\
 &=\min_{\eta\in\mathbb C^Q}
                  (Z_ax+C_a^{\mathsf T}\eta)^*K_N
                         (Z_ax+C_a^{\mathsf T}\eta).
 \end{split}                                                     \tag{PEI5}
\]
There is no requirement that the minimizing $\eta$ agree between
these lines. The second equality follows before minimization from
PEI4 and the zero first rows.

For clarity, if $K>0$, $H$ has full column rank and $J$ is the retained
lift matrix, expansion about $\eta_* =-(H^*KH)^{-1}H^*KJx$ gives
\[
 \begin{split}
 (Jx+H\eta)^*K(Jx+H\eta)
  ={}&x^*\{J^*KJ-J^*KH(H^*KH)^{-1}H^*KJ\}x\\
    &+(\eta-\eta_*)^*H^*KH(\eta-\eta_*).
 \end{split}                                                     \tag{PEI6}
\]
This proves the minimum and retains both cross blocks and the complete
inverse. It is the complex Hermitian Schur calculation; the
corresponding classical block identity is recorded by
Boyd and Vandenberghe, Appendices A.5.5 and C.4.1
\cite{peq:human:boyd-vandenberghe}. Here PEI6 proves the complex form
directly. Since $[H_D,J_D]$ has full column rank, these attained
forms are positive definite.

For any $\eta$, increasing $N$ appends nonnegative squared coordinates
to $\|J_Dx+H_D\eta\|^2$. The old rows of $J_D,H_D$ are unchanged
because their germs, centres and source order are fixed. Taking the
minimum over the same full coefficient space therefore proves
$D_{N+1}\succeq D_N$. No assertion about the equality of successive
minimum sections is needed.

\subsection{A sharper finite bound using each high endpoint}

For integers $q\le N\le2q$, define the finite original AKS envelope
\[
 Z_{a,s}(N)=1+B_{s,q}S_q^2
             \sum_{n=q}^{N}\frac{n!}{(s/2)_n}(25/16)^n,
 \qquad Z_{a,s}(q-1)=1.                                         \tag{PEI7}
\]
All constants on the right are those of \eqref{peq:peq:Z}. The full
positive remainder sum in \eqref{peq:peq:numgram}, truncated at its actual
last row, proves
$K_{q-1}\preceq K_N\preceq Z_{a,s}(N)K_{q-1}$.
Equation PEI5 then proves the same scalar comparison of the attained
numerator forms. For $q\le N\le2q$ put
\[
 E_{a,s}(N)=\log Z_{a,s}(N)
           +2I_{N-v,s,1,1}+2F_{q-v-1,s,1,1},
 \qquad E_{a,s}(q-1)=0.                                        \tag{PEI8}
\]
The value zero at $q-1$ is the exact identity of the endpoint with
itself, not the positive expression applied outside its stated range.
Applying the two rank-one bounds for $B_D$ to every representative
in PEI5 gives the explicit chain
\[
 \begin{split}
 D_N&\preceq e^{2I_{N-v,s,1,1}}\widehat D_N\\
    &\preceq e^{2I_{N-v,s,1,1}}Z_{a,s}(N)\widehat D_{q-1}\\
    &\preceq e^{E_{a,s}(N)}D_{q-1}.
 \end{split}                                                     \tag{PEI9}
\]
It holds alongside $D_{q-1}\preceq D_N$.
At positive rank one the inverse allowance, with $d=D+1\ge2$, is
\[
 I_{D,s,1,1}
 =d\log^+(1/|g(0)|)
  +\tfrac12(s/2-1)_+v\mathsf H_d
  +(d-1)\log H_s+4\log\kappa_s.
                                                               \tag{PEI10}
\]
Every term in this expression is nondecreasing in $D$.
Consequently $E_{a,s}(N)\le\mathfrak E_{a,s}$ for all the displayed
cutoffs. This also checks explicitly the monotonicity of the inverse
allowance used in the incoming manuscript.

Let $A^{\rm end}_{a,s}$ be the exact old PRT18 scalar, with the
substitution PEI1. It retains the original reconstruction norm, mass,
Vandermonde factors and conductor. Set
\[
 e_{a,s}(N)=\min\{\log A^{\rm end}_{a,s},E_{a,s}(N)\},\qquad
 B^{\rm end}_{a,s}=e_{a,s}(2q-1)+e_{a,s}(2q).
                                                               \tag{PEI11}
\]
Both original scalar inequalities hold on the same native form, so
their intersection gives
\[
 D_{q-1}\preceq D_N\preceq e^{e_{a,s}(N)}D_{q-1},
 \qquad
 0\le B^{\rm end}_{a,s}\le2\mathfrak E^{\rm best}_{a,s}.
                                                               \tag{PEI12}
\]
Taking four copies is a direct sum of these inequalities and keeps
the same scalar exponent.

Let $X\subset Y$ be any fixed pair of the original coefficient
spaces and let $G_{Y/X}(N)$ be its original attained metric in a
fixed value frame, of rank $t=\dim(Y/X)$. Applying PEI12 to every
representative of each fixed affine fibre and taking its attained
minimum preserves both inequalities. Congruence by the inverse
positive square root of $G_{Y/X}(q-1)$ shows that each of its $t$
generalized eigenvalues lies in $[1,e^{e_{a,s}(N)}]$. Multiplication,
followed by logarithm, proves
\[
 0\le f_{Y/X}(N):=
       \log\frac{\det G_{Y/X}(N)}{\det G_{Y/X}(q-1)}
       \le t\,e_{a,s}(N).
                                                               \tag{PEI13}
\]
The minima are increasing in $N$, by the same pointwise argument.
The original signed endpoint functional therefore satisfies
\[
 \begin{split}
 0\le\mathcal R\log\det G_{Y/X}
   &=[f_{Y/X}(2q-1)-f_{Y/X}(q-1)]
          +[f_{Y/X}(2q)-f_{Y/X}(q)]\\
   &\le t\,B^{\rm end}_{a,s}.
 \end{split}                                                     \tag{PEI14}
\]
The rank-zero determinant remains one and its return zero.
This is a finite refinement of \eqref{peq:peq:quotreturn}; no new
asymptotic assumption enters it.

\subsection{All primitive CPQ receivers with the correlated ranks}

The maps in \eqref{peq:peq:sequence} send a class to the same coefficient
class. Their exactness follows explicitly as follows. For the first
map, $K_0\cap S_1=K_1$ proves injectivity. If $x\in S_0$ lies in
$T_1=S_1+W$, write $x=s+w$. Then $w=x-s\in K_0$ and its class
maps to $[x]$. Every element of $T_0=S_0+W$ has the class of its
$S_0$ summand modulo $T_1$, proving surjectivity. PRT6--8 uses the
rank-three original relation map, so $\dim(S_0/S_1)=3$.
This proves that the ranks in PEI14 are exactly $3$, $3-r_\partial$
and $r_\partial$.

For a fixed inclusion $X\subset Y$, let $P_X,P_Y$ be the original
frames and let $V$ be the quotient lifts. The actual invertible
coefficient matrix $R$ is defined by $[P_X,V]=P_YR$.
Taking the determinant of PEI6 in this full frame gives
\[
 \det G_{Y/X}(N)=|\det R|^2
   \frac{\det(P_Y^*\mathcal D_NP_Y)}
        {\det(P_X^*\mathcal D_NP_X)}.                            \tag{PEI15}
\]
The factor $|\det R|^2$ is retained. It disappears only under
$\mathcal R$, since it is independent of $N$ and $\mathcal R1=0$.
Applying PEI15 to $S_e/\mathscr B$, $K_e/\mathscr B$,
$S_e/K_e$, $T_e/W$ and $\mathcal U/T_e$ proves exactly the
six original signed identities \eqref{peq:peq:PRT}, including their
orders of subtraction. Thus the complete sharper receiving intervals are
\[
 \boxed{\begin{gathered}
 0\le U_A\le3B^{\rm end}_{a,s},\\
 0\le U_B\le(3-r_\partial)B^{\rm end}_{a,s},\qquad
 0\le U_C\le r_\partial B^{\rm end}_{a,s},\\
 0\le V_{{\rm conn},0}-V_{{\rm conn},1}\le3B^{\rm end}_{a,s},\\
 0\le V_{h,0}-V_{h,1}\le(3-r_\partial)B^{\rm end}_{a,s},\\
 -(3-r_\partial)B^{\rm end}_{a,s}
       \le V_{s,0}-V_{s,1}\le3B^{\rm end}_{a,s},\\
 0\le V_{t,0}-V_{t,1}
      =V_{{\rm rem},1}-V_{{\rm rem},0}
      \le r_\partial B^{\rm end}_{a,s},\\
 |\mathcal L_0-\mathcal L_1|
   =|Z_1-Z_0|
   =|\mathcal P^{\rm rest}_0-\mathcal P^{\rm rest}_1|
   \le3B^{\rm end}_{a,s}.
 \end{gathered}}                                                 \tag{PEI16}
\]
For the last line, $U_B+U_C$ lies in $[0,3B^{\rm end}_{a,s}]$,
and $U_A$ lies in that same interval, so their difference has
absolute value at most $3B^{\rm end}_{a,s}$. This uses the exact
sequence rather than three unrelated ranks. At $r_\partial=0$,
$U_C=0$ exactly; at $r_\partial=3$, $U_B=0$ exactly.

Since $B^{\rm end}_{a,s}\le2\mathfrak E^{\rm best}_{a,s}=O(q)$,
the incoming $6\mathfrak E^{\rm best}_{a,s}$ bound and all its
above-$q$ conclusions follow. The complete identity
$4S_K=\mathcal L_0+Z_0=\mathcal L_1+Z_1$ is unchanged.
The prefix term remains $-V_A^{\rm prefix}$. No independent
proper-source return has been replaced by a target-order return.

For an absolute rank-$t$ restriction or attained quotient, PEI13
contains the factor $t$. At $t=O(a)$ its allowed logarithmic
variation remains $O(aq)$. Accordingly PEI16 supplies no value
for an absolute common loss, $S_K$, $S_R$, an individual complementary
return, the independent proper-source standard-window return, or an
individual retained-class projected sign. Those quantities remain
the original quantities with the original coefficient and metric maps.
The current result advances the precision of their primitive-choice
comparison and its correlated receiving identities.

\subsection{Explicit asymptotic coefficient of the proved envelope}

The finite formula also gives a useful size record without assigning
an allocation coefficient. Put
$G=\log^+(1/|g(0)|)$ and
$z_* =2\log(25\log3/(4\pi))$.
For each of the two order families, $H_s,\kappa_s$ are the fixed
constants of \eqref{peq:peq:H}. From \eqref{peq:peq:Z-size},
PEI10, and $d_+=2q-v+1$, $d_-=q-v$, one obtains
\[
 \limsup_{a\to\infty}\frac{\mathfrak E_{a,s}}q
 \le z_*+4G+4\log H_s+8\log\kappa_s.                            \tag{PEI17}
\]
Indeed $d_+/q\to2$, $(d_+-1)/q\to2$ and
$(d_--1)/q\to1$; the harmonic term divided by $q$ is bounded
by $O(av\log q/q)\to0$. The remaining rank-one constants divided
by $q$ vanish. These limits give respectively $4G$,
$4\log H_s$, and $8\log\kappa_s$ in the displayed bound.
At $s=a$, $H_s=\max\{1,H\}$ and $\kappa_s=2$;
at $s=1$, their literal values in \eqref{peq:peq:H} remain.
The right side is a bound for the proved metric envelope and is
not an evaluated coefficient of any absolute determinant.


\section*{Source and verification record}
The finite source amplification is read from the supplied complete AKS3,15--24 and transported to literal IGI6--10 coordinates. The conductor formulas are read from LRC6--10 with its exact WCF transpose map, and NCT/NPK supplies its original affine-fibre meaning. The three-dimensional coefficient maps, original rank strata and all receiving identities are read directly from PRT1--20. The kernel-current input is the supplied EC37--41. These are written programme proofs, not external claims of numerical period certification. The integration checks the complete AKS envelope, the LRC/XIM affine-fibre map and the PRT coefficient sequence against their pinned provider sources. No Lean execution, actual-period computation or public edition is claimed by this local source. The native affine-fibre and nested-flag calculations are proved explicitly in PEI2--16. No numerical diagnostic is used as proof of an unresolved directional coefficient.



% END PRESERVED INPUT 04_PRIMITIVE_ENDPOINT_INTEGRATION.tex
\endgroup

\clearpage
\begingroup
\def\Gram{\operatorname{Gram}}
\def\Res{\operatorname{Res}}
\def\im{\operatorname{im}}
\def\cond{\operatorname{cond}}

\setcounter{equation}{0}
\renewcommand{\theequation}{ET.\arabic{equation}}
\section*{Original endpoint retractions and reflected-pair period calculation}
\addcontentsline{toc}{section}{Original endpoint retractions and reflected-pair period calculation}
% BEGIN PRESERVED INPUT 05_ENDPOINT_PERIOD_INTEGRATION.tex SHA256 c7f9b7db3b10e23ed4d4331c9e96d7d59cbd939e243e7165aee6c04034c7390f



\section{Objects, source conventions, and the receiving maps}
Keep a complete stipulated quartet
\[
 \rho=\tfrac12+\delta+i\gamma,\quad
 \mathcal R=\{\rho,\bar\rho,1-\rho,1-\bar\rho\},\quad
 0<\delta<\tfrac12,\quad\gamma>2,
\]
each of order \(m\). Set
\[
 h(s)=\prod_{\alpha\in\mathcal R}(s-\alpha)^m,\quad
 v_h=2\xi/h,\quad U_h=j_hv_h,\quad
 E_h=\mathbb C[s]/(h).
\]
For \(k\equiv1\pmod4,\ k\ge5\), retain
\[
 e=1+k(m-1),\quad c=k/2,\quad q=e(k+1)^2,\qquad
 \chi(S)=\prod_{a,b=0}^k
 \bigl(S-c-(2a-k)\delta-i(2b-k)\gamma\bigr)^e.
\]
The spaces and operators are \(E_k=\mathbb C[S]/(\chi)\),
\(A=M_S\), \(A_\Sigma=\sum_iM_{s_i}\), and
\[
 \eta_k:E_k\longrightarrow E_h^{\otimes k},\qquad
 [P]\longmapsto U_h^{\otimes k}P(A_\Sigma)1.                 \tag{ET1}
\]
These are the original maps OPG1--5, not altered unit values
\cite[OPG1--5]{et:Current}. The exact source is
\[
 w_h(y)=|v_h(1/2+iy)|^2/(2\pi),\qquad m_k=w_h^{*k},\qquad
 \sigma(y)=|\Gamma(1/4+iy/2)|^2/(2\pi).
                                                               \tag{ET2}
\]
All norms below keep their stated measures. In particular, a coordinate
congruence used to calculate eigenvalues does not change either metric.

\section{The complete finite endpoint retraction}
Write \(\zeta(\alpha+z)=z^m z_\alpha(z)\), \(z_\alpha(0)\ne0\).
The finite singular germ attached to a local class \(u\) is exactly
\[
 \mathfrak b_\alpha u
 =x^{-\alpha}\Res_{z=0}
       \frac{e^{zL}u(z)}{2z^mz_\alpha(z)}\,dz
 =x^{-\alpha}\sum_{j=0}^{m-1}
 [z^{m-1-j}]\frac{u(z)}{2z_\alpha(z)}\frac{L^j}{j!},
 \quad L=\log(1/x).                                    \tag{ET3}
\]
The equality follows by multiplying the convergent Taylor germs and
selecting the coefficient of \(z^{-1}\). Thus no unknown lower local
coefficient has been suppressed. This verifies the finite-germ formula
E1 of the incoming note \cite[E1]{et:Incoming}. Identifying it with a
particular infinite-dimensional analytic connecting map additionally
uses that map's specified source extension; the construction here works
directly with the displayed residue and does not change its domain.

The linear reversal \(J_m(z^r)=L^{m-1-r}/(m-1-r)!\) satisfies
\(J_mM_z=\partial_LJ_m\), including the top nilpotent power. Multiplication
by the entire unit commutes with \(M_z\). Hence ET3 intertwines local
multiplication by \(\alpha+z\) with \(\alpha+\partial_L\).
For an ordered tuple \(\boldsymbol\alpha\), put
\(\beta=\sum_i\alpha_i\), \(n=k(m-1)\). Diagonal multiplication is
\[
 \mathfrak d_{\boldsymbol\alpha}
 \left(\bigotimes_i x_i^{-\alpha_i}L_i^{j_i}/j_i!\right)
 =x^{-\beta}\frac{(\sum_i j_i)!}{\prod_i j_i!}
                  \frac{L^{\sum_i j_i}}{(\sum_i j_i)!}.       \tag{ET4}
\]
The product rule proves exact intertwining with
\(\sum_i(\alpha_i+\partial_{L_i})\).
In particular the factor \(n!/((m-1)!)^k\) on the longest polynomial
does not disappear.

Let \(A_*(s)=s(s-1)\pi^{-s/2}\Gamma(s/2)/2\) and
\(h_\alpha(s)=h(s)/(s-\alpha)^m\).
The identity \(\xi=A_*\zeta\) gives, in the original local algebra,
\[
 \frac{U_h(\alpha+z)}{2z_\alpha(z)}
       =\frac{A_*(\alpha+z)}{h_\alpha(\alpha+z)}
                  \pmod {z^m}.                            \tag{ET5}
\]
All these units are invertible: the stipulated roots are nontrivial,
the Gamma function has no zeros and no pole at these points
\cite[\S5.2(i)]{et:DLMF}, and
the full orders have already been divided out.
Define the actual polynomial
\[
 C_{\boldsymbol\alpha}(L)=
 \prod_{i=1}^k\left(
  \sum_{j=0}^{m-1}[z^{m-1-j}]
  \frac{A_*(\alpha_i+z)}{h_{\alpha_i}(\alpha_i+z)}
                  \frac{L^j}{j!}\right)
 =\sum_{j=0}^n c_jL^j/j!.
                                                               \tag{ET6}
\]
Its top coefficient is
\[
 c_n=\frac{n!}{((m-1)!)^k}
                 \prod_i\frac{U_h(\alpha_i)}{2z_{\alpha_i}(0)}
 \ne0.                                                        \tag{ET7}
\]
With \(D_{\boldsymbol\alpha}=\mathfrak d_{\boldsymbol\alpha}
\bigotimes_i\mathfrak b_{\alpha_i}\),
\[
 D_{\boldsymbol\alpha}\eta_kP
       =x^{-\beta}P(\beta+\partial_L)C_{\boldsymbol\alpha}(L).
                                                               \tag{ET8}
\]
For an arbitrary output \(x^{-\beta}\sum b_jL^j/j!\), set
\(B(T)=\sum_{r=0}^n b_{n-r}T^r\) and
\(H(T)=\sum_{r=0}^n c_{n-r}T^r\). Reversal intertwines \(\partial_L\)
with multiplication by \(T\), so ET8 is precisely
\[
 B(T)=P(\beta+T)H(T)\pmod{T^{n+1}}.
\]
The complete inverse is finite:
\[
 P(\beta+T)=H(T)^{-1}B(T),\qquad
 H(T)^{-1}=c_n^{-1}\sum_{j=0}^n
                  \left(-\frac{H(T)-c_n}{c_n}\right)^j
                   \pmod{T^{n+1}}.                         \tag{ET9}
\]

Choose and record one ordered tuple at every original centre
\(\beta_{ab}=c+(2a-k)\delta+i(2b-k)\gamma\).
Every such tuple exists because the two signs in a quartet root are
independent. Restrict to that tuple, apply ET9, and use the ordinary
Chinese remainder inverse for the pairwise coprime polynomials
\((S-\beta_{ab})^e\). This defines the exact map
\[
 R_\partial:E_h^{\otimes k}\longrightarrow E_k,\qquad
 R_\partial\eta_k=I,\qquad R_\partial A_\Sigma=AR_\partial.
                                                               \tag{ET10}
\]
These two identities follow component by component from ET9 and from
the derivative intertwining. They also prove injectivity of ET1
without a dimension argument. All unselected tuples remain in the
domain and belong to the specified kernel. For any earlier
equivariant retraction \(R\) of this same inclusion, the exact
comparison is
\[
 R_\partial-R=(D\eta_k)^{-1}D(I-\eta_kR),                 \tag{ET11}
\]
where \(D\) is the direct sum of the selected maps, with the
Chinese remainder target understood. Indeed its right side is
\(R_\partial-R_\partial\eta_kR=R_\partial-R\).
This is an explicit map through the old complementary module;
neither complement is asserted orthogonal.

The complete Jordan list of that kernel can be recovered directly.
Put \(d=m-1\) and give the auxiliary polynomial basis \(L^j\),
\(0\le j\le d\), squared lengths \(1/\binom dj\). This auxiliary
inner product is used only to prove a finite representation identity;
it is not substituted for ET2. The operators
\[
 N=\partial_L,\quad H=2L\partial_L-d,\quad
 R=L(d-L\partial_L)
\]
satisfy \(N^*=R\), \(H^*=H\),
\([H,N]=-2N\), \([H,R]=2R\), \([R,N]=H\).
Their sums act on the ordered tensor. If \(Rv=0\), \(Hv=\ell v\),
then
\[
 RN^jv=j(\ell-j+1)N^{j-1}v,\qquad
 \|N^jv\|^2=j(\ell-j+1)\|N^{j-1}v\|^2.
\]
These identities follow by induction from the commutators and
adjointness. Taking a highest nonzero weight in any invariant
orthogonal complement therefore constructs a string of length
\(\ell+1\); its orthogonal complement is again invariant because
both adjoints belong to the operator family. Induction on the finite
dimension decomposes the complete tensor into these strings.
If \(p_r=[t^r](1+\cdots+t^d)^k\), the weight multiplicities and their
symmetry give exactly \(p_r-p_{r-1}\) strings of length \(kd-2r+1\),
for \(0\le r\le\lfloor kd/2\rfloor\), with \(p_{-1}=0\).
In particular the longest string is unique. Diagonal multiplication
ET4 intertwines all three displayed operators with those of degree
\(kd\), by the product rule. It is onto by its nonzero top value,
and kills each shorter string: a highest weight vector of smaller
weight maps to a vector killed by \(R\) in that same smaller weight,
where the degree-\(kd\) target has no such vector.

At \(\beta_{ab}\), there are exactly
\(M_{ab}=\binom ka\binom kb\) ordered tuples, by choosing the two
sign sets independently. ET10 removes one longest string from their
direct sum. Its kernel consequently has \(M_{ab}-1\) strings of
length \(e\) and
\[
 M_{ab}(p_r-p_{r-1})\quad\hbox{strings of length }e-2r
 \quad(1\le r\le\lfloor(e-1)/2\rfloor).              \tag{ET11a}
\]
Multiplication by the full local units commutes with the original
nilpotent action, so it changes the explicit splitting but preserves
this complete list.

\section{Actual period coordinates and single-class survival}
This section uses the original simple lane \(m=1\). Put
\[
 x=\delta+i\gamma,\quad
 h_c(w)=w^4+\beta w^2+\eta,\quad
 \beta=2(\gamma^2-\delta^2),\quad\eta=(\delta^2+\gamma^2)^2,
\]
\[
 e_x(w)=\frac{h_c(w)}{(w-x)h_c'(x)}
 =\frac{w^3+xw^2+(x^2+\beta)w+x(x^2+\beta)}{h_c'(x)}.
                                                               \tag{ET12}
\]
Use its column \(\mathbf v_x\) in the unchanged frame
\(1,w,w^2,w^3\). The recovered provider gives
\[
 P(u)=e^{\mathfrak a/u}G_u\mathcal R(u^{-1}),\quad
 G_u=\operatorname{diag}(u^{a/5}g_a)_{a=1}^4,\quad
 g_a=5^{a/5-1}e^{\pi ia/5}\Gamma(a/5),
                                                               \tag{ET13}
\]
\[
 \mathfrak a=1/160+\beta/24+\eta/2,\quad
 \det\mathcal R(z)=1,\quad
 \mathcal R(0)=
 \begin{pmatrix}1&0&-\beta/3&0\\0&1&0&-2\beta/3\\
 0&0&1&0\\0&0&0&1\end{pmatrix}=:J_\infty .
                                                               \tag{ET14}
\]
The exact coefficients and their convergent bound are PCL2--6
\cite[PCL1--6, RPC8]{et:Current}:
\[
 R_{ar,n}=
 \sum_{\substack{p,h\ge0\\2p+4h=5n+r+1-a}}
 \frac{\beta^p\eta^h}{3^pp!h!}(-1)^{p+h-n}
 5^{p+h-n}(a/5)_{p+h-n},\quad
 C_1=\sum_{n\ge1}\|R_n\|_1<\infty ,
                                                               \tag{ET15}
\]
\[
 \|\mathcal R(z)-J_\infty\|_1\le C_1|z|\quad(|z|\le1).
\]
The original logarithm branch remains in \(G_u\). Let
\(\zeta_5=e^{2\pi i/5}\), \(D=\operatorname{diag}(\zeta_5^a)\) and
\[
 \mathcal P_k=\frac15\sum_{j=0}^4(D^{-j})^{\otimes k}.
                                                               \tag{ET16}
\]
OPG5 and its Fourier conjugacy show that this is the actual
observation projector in the stated Fourier coordinates, not a new
observation \cite[OPG5, PQO5]{et:Current}. It is an orthogonal projector
for the Euclidean metric in those coordinates.
Writing \(b=P(u)\mathbf v_x\), the original extremal class has output
\[
 U_h(\rho)^k\mathcal P_kb^{\otimes k}.                  \tag{ET17}
\]
Since \(P(u)\) is invertible, \(b\ne0\). For \(5\mid k\), choose
any nonzero entry \(b_a\). The all-\(a\) ordered tensor entry is
nonzero and survives ET16. Thus ET17 is nonzero on the original
cofinal sequence \(k=20n+5\). If two entries \(b_a,b_b\) are nonzero,
choose \(j\in\{0,\ldots,4\}\) with \(ka+j(b-a)=0\bmod5\).
For \(k\ge4\), the entry with \(j\) labels \(b\) and \(k-j\) labels
\(a\) exists and proves survival.

With \(p_a=|b_a|^2/\|b\|^2\), exact character orthogonality gives
\[
 f_k(b):=\frac{\|\mathcal P_kb^{\otimes k}\|^2}{\|b\|^{2k}}
  =\frac15\sum_{j=0}^4
       \left(\sum_a p_a\zeta_5^{ja}\right)^k.          \tag{ET18}
\]
If at least two entries are nonzero,
\(r_b=\max_{1\le j\le4}|\sum_ap_a\zeta_5^{ja}|<1\),
since distinct fifth roots cannot all have the same argument.
Consequently \(|f_k(b)-1/5|\le4r_b^k/5\).
If exactly one entry is nonzero, \(f_k(b)\) is one for \(5\mid k\)
and zero otherwise. These statements verify E15--18 in their
actual period metric, with the common unit cancelling only in the ratio.

\section{A new reflected-pair observation calculation}
Retain \(b_+=P(u)\mathbf v_x\), \(b_-=P(u)\mathbf v_{-x}\) and
\[
 a=\|b_+\|^2,\quad b_0=\|b_-\|^2,\quad
 z_j=b_+^*D^jb_-,\quad
 a_j=b_+^*D^jb_+,\quad b_j=b_-^*D^jb_- .
\]
The exact output Gram on the original reflected pair is
\[
 \mathcal O_k=
 |U_h(\rho)|^{2k}\frac15
 \begin{pmatrix}
 \sum_{j=0}^4a_j^k&\sum_{j=0}^4z_j^k\\
 \sum_{j=0}^4\overline{z_j}^{\,k}&\sum_{j=0}^4b_j^k
 \end{pmatrix}.                                      \tag{ET19}
\]
The lower left uses the summed conjugates; replacing \(j\) by
\(-j\) proves Hermitian symmetry directly. To prove ET19, insert
ET16 between each pair of tensor vectors and use
\(\langle v^{\otimes k},w^{\otimes k}\rangle
=\langle v,w\rangle^k\).
The original reflection has the same unit at \(\rho\) and \(1-\rho\)
because both \(h\) and \(\xi\) satisfy the same reflection; this is
why the displayed unit is common. No cross term is discarded.

Here is an explicit nonalignment guard derived from the actual PCL
coefficients. Put
\[
 A_0=x/h_c'(x),\quad B_0=1/h_c'(x),\quad
 c_5=\min_j|1+\zeta_5^j|=(\sqrt5-1)/2,
\]
\[
 V=\|\mathbf v_x\|_1=\|\mathbf v_{-x}\|_1,\qquad
 U_*=\max\left\{1,\frac{16C_1V}
                    {c_5\min\{|A_0|,|B_0|\}}\right\}. \tag{ET20}
\]
The last two entries of \(J_\infty\mathbf v_x\) are
\((A_0,B_0)\), and those of \(J_\infty\mathbf v_{-x}\) are
\((A_0,-B_0)\), because \(h_c'(-x)=-h_c'(x)\).
For \(|u|\ge U_*\), every perturbation of those four entries is at
most \(d=c_5\min(|A_0|,|B_0|)/16\).
For each \(j\), the two-row minor between these two columns after
multiplying the second by \(D^j\) has unperturbed modulus
\(|A_0B_0||1+\zeta_5^j|\).
Its perturbation is at most
\[
 2d(|A_0|+|B_0|)+2d^2
 \le |A_0B_0|(c_5/4+c_5^2/128)
 <\tfrac12c_5|A_0B_0|.
                                                               \tag{ET21}
\]
Thus all five minors are nonzero. Multiplication by the original
nonzero diagonal \(e^{\mathfrak a/u}G_u\) does not change their
nonvanishing. In particular \(b_+\) and \(D^jb_-\) are independent
for every \(j\). This also guarantees two nonzero entries in both
columns.

At a fixed admitted period satisfying ET20 define the actual numbers
\[
 r_+=\max_{1\le j\le4}|a_j|/a<1,\quad
 r_-=\max_{1\le j\le4}|b_j|/b_0<1,\quad
 r_\times=\max_{0\le j\le4}|z_j|/\sqrt{ab_0}<1.
                                                               \tag{ET22}
\]
Strictness in the third inequality follows from ET21 and the
equality case of Cauchy--Schwarz. These are finite sums of the four
original coordinates; no unknown asymptotic Gram replaces them.
On the finite guards \(4r_+^k\le1\) and \(4r_-^k\le1\),
ET19 gives the lower determinant bound
\[
 \frac{\det\mathcal O_k}
 {|U_h(\rho)|^{4k}(ab_0)^k}
 \ge\frac{(1-4r_+^k)(1-4r_-^k)}{25}
                 -r_\times^{2k}.                    \tag{ET23}
\]
For example the right side is positive whenever
\[
 4r_+^k\le\tfrac14,\qquad4r_-^k\le\tfrac14,\qquad
 25r_\times^{2k}<\tfrac9{16}.                         \tag{ET24}
\]
Taking logarithms supplies a finite threshold directly from ET22;
if a radius is zero its condition is already satisfied.
This proves eventual injectivity of the original observation on the
entire reflected pair at every fixed admitted period in ET20.
It strengthens single-class nonvanishing without asserting a
characteristic-changing comparison.

There is a quantitative full-pair version of the \(1/5\) calculation.
The unprojected period Gram is
\[
 \mathcal H_k=|U_h(\rho)|^{2k}
       \begin{pmatrix}a^k&z_0^k\\\bar z_0^{\,k}&b_0^k\end{pmatrix}.
\]
Define
\[
 \epsilon_k=
 \frac{\frac45\max(r_+^k,r_-^k)+\frac45r_\times^k}
 {1-r_\times^k}.
                                                               \tag{ET25}
\]
Then
\[
 (1/5-\epsilon_k)\mathcal H_k\preceq
       \mathcal O_k\preceq(1/5+\epsilon_k)\mathcal H_k.
                                                               \tag{ET26}
\]
Indeed congruence by
\(|U_h(\rho)|^{-k}\operatorname{diag}(a^{-k/2},b_0^{-k/2})\)
leaves the unprojected Gram with diagonal one and off-diagonal
modulus at most \(r_\times^k\).
The difference \(\mathcal O_k-\mathcal H_k/5\) has diagonal entries
at most \(4\max(r_+^k,r_-^k)/5\) in modulus, and off-diagonal
entry at most \(4r_\times^k/5\), since its \(j=0\) summand cancels
exactly. Its operator norm is bounded by their sum. Dividing by
the lower eigenvalue \(1-r_\times^k\) proves ET26, after reversing
the congruence. Every unit and coordinate factor is therefore
present in the asserted original inequality.

\section{Reconstructing the first-cutoff arithmetic parity coefficient}
This proof recovers S2's parity coefficient from the source envelopes,
without using an unavailable separate parity calculation.
Set \(\lambda=k\rho=c+\mu\), \(\mu=k(\delta+i\gamma)\),
\(R=|\mu|\), and keep the terminal class
\[
 v=\left[\frac{\chi(S)}
  {(S-\lambda)\chi_\lambda(\lambda)}\right],\qquad
 \chi_\lambda=\chi/(S-\lambda)^e .
\]
At \(N=q-1\) the quotient representative is unique. Define the
complete polynomial
\[
 B_k(y)=\frac{\chi(c+iy)}{(iy-\mu)(iy+\mu)},\qquad n=q-2,\qquad
 M_j=\int_{\mathbb R}y^j|B_k(y)|^2m_k(y)\,dy .
                                                               \tag{ET27}
\]
Its leading coefficient has modulus one; its remaining \(n\) roots,
with their full remaining multiplicities, have modulus at most
\(R=k\sqrt{\delta^2+\gamma^2}\).
The polynomial is even. The literal representative is
\((iy+\mu)B_k(y)/\chi_\lambda(\lambda)\).
Evenness of the source kills the even/odd cross term exactly.
Consequently the even-energy fraction is
\[
 \theta_{q-1}=\frac{R^2M_0}{M_2+R^2M_0}.              \tag{ET28}
\]
This derivation is valid for every complete multiplicity; it never
uses \(\chi'(\lambda)\).

For completeness the source estimates used here are exactly GTM1
and RTM2 \cite[GTM1, RTM2]{et:Envelopes}:
\[
 a_k(1+y^2)^{-M}\sigma(y)\le m_k(y)\le A_k\sigma(y),
 \quad M=21+4m,\quad
 \log(A_k/a_k)=O_h(k+\log k),
                                                               \tag{ET29}
\]
\[
 c_\Gamma e^{-\alpha|y|}(1+|y|)^{-1/2}
 \le\sigma(y)\le C_\Gamma e^{-\alpha|y|}(1+|y|)^{-1/2},
\quad\alpha=\pi/2,
\]
\[
 c_\Gamma=\sqrt2e^{-7/6},\qquad
 C_\Gamma=\sqrt{2\sqrt5}e^{2/3}.
\]
In particular the full source mass remains in \(a_k,A_k\).
Their definitions are
\[
 a_k=\frac{c_h\vartheta_h^{k-3}e^{-\alpha(k-3)}(k-2)^{-42-8m}}
 {C_\Gamma 2^{21+4m}},\quad
 A_k=\frac{C_h^{\rm tilt}}{c_\Gamma}k^{8m-7/2}
                    \mathfrak M^{k-1},
\]
\[
 \vartheta_h=\int_{-1}^1w_h(y)\,dy,\qquad
 \mathfrak M=\int_{\mathbb R}e^{\alpha y}w_h(y)\,dy.
\]

Here are explicit finite bounds establishing concentration, including
the unbounded tails. Set \(x_0=2n/\alpha\), and impose only the
eventual guard \(R<x_0\). Restriction to \([x_0,x_0+1]\) gives
\[
 M_0\ge L_k:=
 \frac{a_kc_\Gamma e^{-\alpha(x_0+1)}(x_0-R)^{2n}}
 {(1+(x_0+1)^2)^M\sqrt{x_0+2}}.                      \tag{ET30}
\]
For \(j=0,2\), the upper bound on either half-line is
\[
 |y|^j|B_k(y)|^2m_k(y)
 \le A_kC_\Gamma |y|^j(|y|+R)^{2n}e^{-\alpha|y|}.
                                                               \tag{ET31}
\]
Put \(s_*=2/\alpha=4/\pi\) and fix \(0<\varepsilon<s_*\).
For \(t=|y|+R\), the region
\(\bigl||y|/q-s_*\bigr|>\varepsilon\) lies in
\[
 t<q(s_*-\varepsilon)+R
 \quad\hbox{or}\quad t>q(s_*+\varepsilon)+R.
\]
For \(\nu_j=2n+j+1\), let \(T_j\) denote the auxiliary integral
with density \(\alpha^{\nu_j}t^{\nu_j-1}e^{-\alpha t}/(\nu_j-1)!\)
on \(t>0\); this auxiliary notation does not rescale ET2.
Integrating ET31 yields the exact bound
\[
 \frac{\int_{\{||y|/q-s_*|>\varepsilon\}}
              |y|^j|B_k(y)|^2m_k(y)\,dy}{q^jM_0}
 \le
 \frac{2A_kC_\Gamma e^{\alpha R}(2n+j)!}
 {q^j\alpha^{2n+j+1}L_k}\,
 \Pr\{T_j\notin[q(s_*-\varepsilon)+R,
                         q(s_*+\varepsilon)+R]\}.       \tag{ET32}
\]
For \(v>1\) and \(0<v<1\), respectively, elementary exponential
Markov estimates give
\[
 \Pr\{T_j\ge\nu_jv/\alpha\}\le
 e^{-\nu_j(v-1-\log v)},\qquad
 \Pr\{T_j\le\nu_jv/\alpha\}\le
 e^{-\nu_j(v-1-\log v)}.                             \tag{ET33}
\]
To prove them integrate \(e^{s\alpha t}\) against the displayed
density, obtaining \((1-s)^{-\nu_j}\), and choose \(s=1-1/v\).
This is positive for the upper tail and negative for the lower tail.
All integrals follow by repeated integration by parts.

The logarithm of the prefactor in ET32 is \(O_h(k+\log q)\).
Indeed \(\log(2n+j)!=(2n+j)\log(2n+j)-(2n+j)+O(\log n)\),
obtained by bounding the sum of \(\log r\) between its two adjacent
integrals; \(R=O_h(k)\), \(R/n\to0\), and
\(-2n\log(1-R/x_0)=O_h(k)\).
The remaining factors in ET30 are polynomial in \(q\), and ET29
gives the stated source ratio. Each tail in ET33 has exponent
\(-c_\varepsilon q+o(q)\), since its \(v\) tends to
\(1\pm\alpha\varepsilon/2\). Thus the right side of ET32 tends to
zero for \(j=0\) and \(j=2\). On the complementary region
\((s_*-\varepsilon)^2\le y^2/q^2\le(s_*+\varepsilon)^2\).
Letting \(\varepsilon\downarrow0\) proves
\[
 \frac{M_2}{q^2M_0}\longrightarrow\frac{16}{\pi^2},\qquad
 \theta_{q-1}=
 \frac{\pi^2(\delta^2+\gamma^2)}{16}
                  \frac{k^2}{q^2}(1+o_h(1)).         \tag{ET34}
\]
This proof used the unchanged complete root polynomial, all roots in
the degree count, both half-lines, and the actual source envelopes.

\section{Exact arithmetic pair metrics and their period comparison}
Reflection \(\Sigma P(S)=P(k-S)\) is an isometry of the original
quotient minimum because ET2 is even and \(\chi(k-S)=\chi(S)\).
For \(P_N\), its minimum representative of \(v\), put
\(E_N=\|P_N\|^2\) and
\(\theta_N=\|(P_N+\Sigma P_N)/2\|^2/E_N\).
The reflected pair is independent, and its exact Gram is
\[
 G_N^{\rm pair}=E_N
  \begin{pmatrix}1&h_N\\h_N&1\end{pmatrix},
 \qquad h_N=2\theta_N-1,\quad0<\theta_N<1.             \tag{ET35}
\]
The decomposition into even and odd parts proves the off-diagonal
entry and positivity. On this pair \(A-cI=\operatorname{diag}(\mu,-\mu)\).
The spectral coordinate projector onto the first line has norm
\[
 \|E_+\|_{G_N^{\rm pair}}=(1-h_N^2)^{-1/2}.           \tag{ET36}
\]
For a fixed first coordinate \(z_1\), minimization over \(z_2\) in
the denominator gives \((1-h_N^2)|z_1|^2\), proving the formula.
For real \(\tau\), \(T=\exp(\tau(A-cI))\) has squared singular
values with product one and sum
\[
 t_N(\tau)=\frac{2\cosh(2k\delta\tau)
                     -2h_N^2\cos(2k\gamma\tau)}{1-h_N^2}.
                                                               \tag{ET37}
\]
This follows by multiplying \(G^{-1}T^*GT\); therefore its two
eigenvalues are \((t_N\pm\sqrt{t_N^2-4})/2\).
Using ET34, for each fixed \(\tau\ne0\),
\[
 \|T\|_{G_{q-1}^{\rm pair}}
 =\frac{2q}{\pi k\sqrt{\delta^2+\gamma^2}}
              e^{|\tau|k\delta}(1+o_{h,\tau}(1)).     \tag{ET38}
\]
The original uncentred action has the additional factor
\(e^{\tau k/2}\), retained exactly.

For comparison with the proposed normal-action metric in S9, any
positive metric making this two-eigenvalue action normal is diagonal
in the displayed eigenbasis: \(H=\operatorname{diag}(a_*,b_*)\),
\(a_*,b_*>0\). Distinct eigenvectors of a normal operator are
orthogonal, as follows by applying the adjoint eigenvalue equation
to their inner product; conversely a diagonal metric makes this
diagonal action normal. The squared generalized trace divided by
the generalized determinant is
\[
 \frac{(a_*+b_*)^2}{a_*b_*(1-h_N^2)}.
\]
Its minimum is at \(a_*=b_*\), since
\((a_*-b_*)^2\ge0\). The positive generalized eigenvalue ratio
increases with that expression. Thus, exactly,
\[
\begin{gathered}\inf_{H:\,A\ {\rm normal\ in}\ H}
 \cond((G_N^{\rm pair})^{-1}H)
   =\frac{1+|h_N|}{1-|h_N|},\\
 \inf_{H:\,A\ {\rm normal\ in}\ H}
 \cond((G_{q-1}^{\rm pair})^{-1}H)
 =\frac{16q^2}{\pi^2(\delta^2+\gamma^2)k^2}
                                             (1+o_h(1)).\end{gathered} \tag{ET38a}
\]
This evaluates the metric cost of that specified construction; the
calculation of the original arithmetic metric still uses ET35.

In the simple period lane the unprojected period metric on the pair
is \(\mathcal H_k\) from ET25. Its relative generalized eigenvalue
condition number \(\kappa\) obeys the exact equation
\[
 \kappa+\kappa^{-1}+2=
 \frac{(a^k+b_0^k-2h_N\Re z_0^k)^2}
 {(1-h_N^2)((ab_0)^k-|z_0|^{2k})}.                  \tag{ET39}
\]
Indeed the right side is the squared generalized trace divided by
the generalized determinant. Both scalar factors
\(E_N\) and \(|U_h(\rho)|^{2k}\) cancel only in this condition number.
At every fixed admitted period, independence of \(b_+,b_-\) gives
\(|z_0|/\sqrt{ab_0}<1\). Combining ET34 with ET39 gives
\[
 \kappa_{q-1,k}=
 \frac{4q^2}{\pi^2(\delta^2+\gamma^2)k^2}
 \left[(a/b_0)^k+(b_0/a)^k+2\right](1+o_{h,u}(1)).
                                                               \tag{ET40}
\]
For example, the cross term in the trace divided by
\(a^k+b_0^k\) is at most
\((|z_0|/\sqrt{ab_0})^k\); the determinant correction is its
square, and the positive solution of ET39 divided by its right
side tends to one. These estimates prove ET40 also when the period
norms are unbalanced. Hence S14 is recovered from explicit source
envelopes and the original period provider, rather than imported
as an unsupported asymptotic.

\section{The exact next receiving calculation}
Let \(W=P(u)^{\otimes k}\eta_kE_k\) be the original cyclic image
in Fourier tensor coordinates and \(L=\mathcal P_k|_W\).
The period-metric quotient norm of the original output \(Lv\) is
the minimum of \(\|w\|^2\) over \(w\in W\) with \(Lw=Lv\).
As \(\mathcal P_k\) is orthogonal, it is at least \(\|Lv\|^2\)
and at most \(\|v\|^2\); this proves precisely the bounds in E18.
ET19--26 calculate \(L^*L\) on the reflected pair, including its
nonzero cross term. They do not replace the full cyclic fibre by
that two-dimensional domain.

In any fixed original basis of \(E_k\), with its full period Gram
\(H\) and the original observation matrix \(\Lambda\), the attained
quotient matrix in a full basis of the actual image is
\[
 Q_{\rm per}=(\Lambda H^{-1}\Lambda^*)^{-1}.           \tag{ET41}
\]
To prove this, set
\(S=H^{-1}\Lambda^*(\Lambda H^{-1}\Lambda^*)^{-1}\).
Then \(\Lambda S=I\) and \(S^*HI=0\) for every kernel column \(I\),
so \(Sv\) is the unique minimum lift and \(S^*HS=Q_{\rm per}\).
The arithmetic receiver has the identical exact formula with the
original full arithmetic \(G_N\). This specifies the two exact
maps whose actual inverse Grams must enter the seven allocations;
ET40 is the proven pair restriction of their metric comparison.
No individual arithmetic projected sign follows merely from ET18
or ET26, and none is assigned here.

The incoming E19--21 additionally refer to a finite-field exponential
family and a lower-boundary exact sequence. Deligne's weight theorem
is a theorem about its specified finite-field complexes and Frobenius
\cite[Thm.\ 3.3.1]{et:Deligne}. ET3--41 do not change coefficient fields
or assign Frobenius to the analytic dilation. The exact identity
\(\pi C\partial^{\otimes k}=\pi(C\partial^{\otimes k}-\iota b)\)
uses \(\pi\iota=0\); it neither constructs the comparison \(C\) nor
evaluates it. The present accepted receiving results are the explicit
complex maps and metrics above. The supplied family's stationary-phase
and boundary-weight providers have not been replaced by that identity.



% END PRESERVED INPUT 05_ENDPOINT_PERIOD_INTEGRATION.tex
\endgroup

\clearpage
\begingroup
\def\C{\mathbb C}
\def\R{\mathbb R}
\def\PP{\mathcal P}
\def\Bcal{\mathcal B}
\def\Rcal{\mathcal R}
\def\Scal{\mathcal S}
\def\norm#1{\left\lVert#1\right\rVert}
\def\ip#1#2{\left\langle#1,#2\right\rangle}
\def\Tr{\operatorname{Tr}}
\def\diag{\operatorname{diag}}
\def\im{\operatorname{im}}
\def\rem{\operatorname{rem}}
\def\rad{\operatorname{rad}}

\setcounter{equation}{0}
\renewcommand{\theequation}{CPT.\arabic{equation}}
\section*{Central path transfer through the attained quotient}
\addcontentsline{toc}{section}{Central path transfer through the attained quotient}
% BEGIN PRESERVED INPUT 18_CENTRAL_PATH_TRANSFER.tex SHA256 cd1aec3fb1899d572d622e11abcd901403936bf8d07789155f201b6fcb05d200



\begin{abstract}
The full two-parameter arithmetic source interpolation has a uniformly
small central derivative on every specified complete annihilator
polynomial space.  Its explicit derivative bound is the product of the
original logarithmic density contrast and the central localization
fraction.  Integrating this derivative proves exponential relative
metric bounds and an explicit bound for the actual minimum sections.
For the complete relation ideal this gives a finite central-leg error
and a finite rectangle error in the original four-endpoint signed
determinant.  The remaining source-band interaction is expressed by its
exact two-variable reproducing-kernel integral.  Every statement retains
the full primary orders, original physical coordinate, original
polynomial remainder map, complete real integration domain and both
endpoint multiplicities.
\end{abstract}

\section{Original source, analytic providers and fixed notation}

Fix the original complete hypothetical quartet
\[
 \rho=\tfrac12+\delta+i\gamma,\qquad
 0<\delta<\tfrac12,\quad\gamma>2,\quad m_0\in\mathbb Z_{\ge1}.
\]
The divisor, its entire Taylor unit, and its arithmetic density are
\begin{equation}\tag{CPT1}
\begin{gathered}
 h(z)=\prod_{\lambda\in\{\rho,\bar\rho,1-\rho,1-\bar\rho\}}
             (z-\lambda)^{m_0},\qquad
 v_h(z)=2\xi(z)/h(z),\\
 w_h(y)=|v_h(1/2+iy)|^2/(2\pi),\qquad m_k=w_h^{*k},\\
 k\equiv1\pmod4,\quad k\ge5,\quad
 e=1+k(m_0-1),\quad q=e(k+1)^2,\quad c=k/2,\\
 \chi(S)=\prod_{a,b=0}^k
  [S-c-(2a-k)\delta-i(2b-k)\gamma]^e,\qquad
 E=\C[S]/(\chi).
\end{gathered}
\end{equation}
The symbol \(\xi\) and the zero hypothesis have their original
Riemann-xi meaning.  This note computes consequences on that original
counterfactual source; it supplies no numerical off-critical zero.
Every limit below fixes the complete packet.

The Gamma measure, full tilted mass and central endpoint are
\begin{equation}\tag{CPT2}
\begin{gathered}
 \sigma(y)=\frac{|\Gamma(1/4+iy/2)|^2}{2\pi},\quad
 d\sigma=\sigma(y)\,dy,\quad c_\sigma=\int_{\R}d\sigma=\sqrt{2\pi},
 \quad \alpha=\pi/2,\\
 M_h(\theta)=\int_{\R}e^{\theta y}w_h(y)\,dy,\quad
 \ell_h(\theta)=\log M_h(\theta),\quad
 M=M_h(\alpha),\quad M_0=M_h(0),\\
 \mu=\ell_h'(\alpha)>0,\quad T=k\mu,\quad
 R=k\sqrt{\delta^2+\gamma^2},\quad
 C=[-T,T],\quad O=\R\setminus C.
\end{gathered}
\end{equation}
The notation \(M\) here always denotes the full tilted mass; set
\(\beta_0=21+4m_0\) for the polynomial exponent, to avoid confusing
these two different quantities.  In particular \(B_0=2\beta_0=42+8m_0\)
and \(p=8m_0-9/2>3\).

The following original analytic estimates are the GTM1 and RTM2
providers, including their constants \cite{cpt:providers}:
\begin{equation}\tag{CPT3}
\begin{gathered}
 a_k(1+y^2)^{-\beta_0}\sigma(y)\le m_k(y)\le A_k\sigma(y),\\
 a_k=
 \frac{c_h\vartheta_h^{\,k-3}e^{-\alpha(k-3)}(k-2)^{-B_0}}
 {C_\Gamma2^{B_0/2}},\qquad
 A_k=\frac{C_h^{\rm tilt}}{c_\Gamma}k^{p+1}M^{k-1},
 \qquad \vartheta_h=\int_{-1}^1w_h(y)\,dy,\\
 c_\Gamma e^{-\alpha|y|}(1+|y|)^{-1/2}
 \le\sigma(y)\le
 C_\Gamma e^{-\alpha|y|}(1+|y|)^{-1/2},\\
 c_\Gamma=\sqrt2e^{-7/6},\qquad
 C_\Gamma=\sqrt{2\sqrt5}e^{2/3}.
\end{gathered}
\end{equation}
Here \(c_h,C_h^{\rm tilt}>0\) are the original threefold convolution
and boundary-tilt constants.  Their original source bounds are
\[
 w_h^{*3}(y)\ge c_h e^{-\alpha|y|}(1+|y|)^{-B_0},
 \qquad e^{\alpha y}w_h(y)\le
 C_h^{\rm tilt}(1+|y|)^{-p}.
\]
These statements, with their source derivations, are preserved in the
provider file; this continuation does not select new density constants.
No positive lower estimate for the one-factor zeta numerator is used.

Define the exact logarithmic contrast and the entire parameter square:
\begin{equation}\tag{CPT4}
\begin{gathered}
 H(y)=\log\frac{m_k(y)}{M^k\sigma(y)},\qquad
 H_C=1_CH,\quad H_O=1_OH,\\
 d\nu_{s,t}(y)=\sigma(y)e^{sH_C(y)+tH_O(y)}\,dy,
 \qquad (s,t)\in[0,1]^2.
\end{gathered}
\end{equation}
Thus \(\nu_{0,0}=\sigma\), \(\nu_{1,1}=M^{-k}m_k(y)dy\), and
\(\nu_{0,1}\) retains the exact arithmetic exterior while its centre is
the specified Gamma density.  The factor \(M^k\) has not disappeared:
its exact cancellation in the particular signed determinant is proved
in Section~\ref{cpt:sec:determinants}.

All polynomial inner products are conjugate linear in their first
argument:
\[
 \ip{P}{Q}_{s,t}
 =\int_{\R}\overline{P(c+iy)}Q(c+iy)\,d\nu_{s,t}(y).
\]
The coefficient map
\begin{equation}\tag{CPT5}
 U_{c,D}:\C[S]_{\le D}\longrightarrow\C[y]_{\le D},
 \qquad U_{c,D}P(y)=P(c+iy)
\end{equation}
is an exact complex-linear isomorphism with inverse
\(F(y)\mapsto F((S-c)/i)\).  In increasing monomial frames its
diagonal is \(1,i,\ldots,i^D\), so
\(\det U_{c,D}=i^{D(D+1)/2}\).  This coordinate map, with modulus-one
determinant, will be displayed whenever the real-coordinate proof is
returned to the original coefficient quotient.

\section{Uniform bounds on the entire interpolation square}

Put
\begin{equation}\tag{CPT6}
 b_*=\max\{1,A_k/M^k\},\quad a_*=\min\{1,a_k/M^k\},
 \qquad
 a_D=\frac{a_*}{[1+4(D+1)^2]^{\beta_0}}.
\end{equation}
For every \((s,t)\in[0,1]^2\),
\begin{equation}\tag{CPT7}
 a_*(1+y^2)^{-\beta_0}\,d\sigma
 \le d\nu_{s,t}\le b_*\,d\sigma,\qquad
 a_D\norm{F}_\sigma^2\le\norm{F}_{s,t}^2
 \le b_*\norm{F}_\sigma^2\quad(\deg F\le D).
\end{equation}

To prove the pointwise bounds, let
\(r(y)=m_k(y)/(M^k\sigma(y))>0\).
At each \(y\), precisely one exponent in (CPT4) acts, and its value
is some \(u\in[0,1]\).  For any positive \(r\),
\[
 \min\{1,r\}\le r^u\le\max\{1,r\}.
\]
Applying (CPT3) gives the first two bounds in (CPT7).
In particular this proves a pointwise bound before any polynomial norm
is considered.

For the polynomial lower estimate, the monic Gamma polynomials
\(b_j(y)\) obey
\begin{equation}\tag{CPT8}
 yb_j=b_{j+1}+j(j-\tfrac12)b_{j-1},\qquad
 \gamma_j=\norm{b_j}_\sigma^2
 =\sqrt{2\pi}\frac{(2j)!}{4^j}.
\end{equation}
These are GTM7 \cite{cpt:providers}; the recurrence is also the
Meixner--Pollaczek recurrence in the convention
\(b_j(y)=j!P_j^{(1/4)}(y/2;\pi/2)\)
\cite[(18.22.8)]{cpt:dlmfMP}.  Indeed substituting
\(\lambda=1/4,\varphi=\pi/2,x=y/2\) into that cited recurrence and
multiplying by \(j!\) gives exactly (CPT8).
For the orthonormal polynomials \(e_j=b_j/\sqrt{\gamma_j}\),
\[
 ye_j=\sqrt{(j+1)(j+\tfrac12)}\,e_{j+1}
       +\sqrt{j(j-\tfrac12)}\,e_{j-1}.
\]
The upward and downward shifts on
\(\operatorname{span}(e_0,\ldots,e_D)\) each have operator norm at
most \(D+1\), including the output coefficient at degree \(D+1\).
The triangle inequality therefore proves
\[
 \norm{yF}_\sigma\le2(D+1)\norm{F}_\sigma.
\]
For a nonzero \(F\) divide the positive measure
\(|F|^2d\sigma\) by its actual mass \(\norm{F}_\sigma^2\), solely
to apply the convex scalar inequality
\(\mathbb E(1+X)^{-\beta_0}\ge(1+\mathbb EX)^{-\beta_0}\)
to \(X=y^2\).  This inequality follows from the supporting tangent
line to the function \((1+x)^{-\beta_0}\), whose second derivative
is positive on \(x\ge0\).  Multiplying the mass back gives
\[
 \int |F|^2(1+y^2)^{-\beta_0}d\sigma
 \ge[1+4(D+1)^2]^{-\beta_0}\norm{F}_\sigma^2.
\]
Combining this with the pointwise lower bound proves (CPT7).
No source measure or source polynomial is replaced in that argument.

The pointwise envelopes also imply
\[
 |H(y)|\le
 |\log(a_k/M^k)|+|\log(A_k/M^k)|
       +\beta_0\log(1+y^2).
\]
Thus every finite polynomial moment and every finite parameter
derivative of it is dominated on the closed square by a polynomial
times a power of \(\log(1+y^2)\) times the exponentially decreasing
Gamma density.  Differentiation under the integral is valid to all
finite orders, with one-sided derivatives at the square's boundary.
Every Gram is positive definite, since its density is positive
everywhere and a nonzero polynomial has only finitely many real
zeros.  Its inverse and every fixed-fibre minimum section therefore
depend smoothly on the square.

\section{Finite central localization with every root retained}
\label{cpt:sec:localization}

Let \(g\mid\chi\) be an actual monic divisor, with its selected full
primary orders, and write
\begin{equation}\tag{CPT9}
 J=\deg g,\quad d_0=q-J,\quad
 Q_g(y)=i^{-d_0}(\chi/g)(c+iy),\quad
 L=D-d_0,\quad
 V_{g,D}=Q_g\,\C[y]_{\le L}.
\end{equation}
Assume the explicit finite guards
\[
 0\le J\le q,\qquad d_0\le D\le2q+1,\qquad
 R<q,\qquad T\le q/2.
\]
Each root of \(Q_g\), counted with its remaining multiplicity, has
modulus at most \(R\).  Define the finite positive bound and its cap
\begin{equation}\tag{CPT10}
\begin{split}
 \mathfrak e_{k,J,D}
 ={}&\frac{b_*}{a_D}\frac{c_\sigma}{c_\Gamma}
 \frac{(L+1)^2\sqrt{1+2q}}q\,e^{\pi q}100^L
 \left(\frac{T+R}{q-R}\right)^{2(q-J)},\\
 &\varepsilon_{k,J,D}=\min\{1,\mathfrak e_{k,J,D}\}.
\end{split}
\end{equation}
Then the entire parameter square satisfies
\begin{equation}\tag{CPT11}
 \int_C |F(y)|^2d\nu_{s,t}(y)
 \le\varepsilon_{k,J,D}\norm{F}_{s,t}^2,
 \qquad F\in V_{g,D}.
\end{equation}

Here is the full polynomial estimate.  On \([q,2q]\) the
orthonormal Legendre basis of polynomials of degree at most \(L\)
is
\[
 \varphi_j(y)=q^{-1/2}\sqrt{2j+1}\,
 P_j(2y/q-3),\qquad 0\le j\le L.
\]
The Legendre norm \(2/(2j+1)\) and recurrence used here are
respectively \cite[Table 18.3.1]{cpt:dlmfLegendreNorm} and
\cite[Table 18.9.1]{cpt:dlmfLegendreRec}.  The latter gives
\[
 (j+1)P_{j+1}(x)=(2j+1)xP_j(x)-jP_{j-1}(x).
\]
For \(|x|\le4\), \(P_0=1,P_1=x\) and induction prove
\(|P_j(x)|\le10^j\): the inductive upper bound is
\[
 \frac{2j+1}{j+1}\,4\,10^j+
 \frac{j}{j+1}\,10^{j-1}
 \le8\,10^j+10^{j-1}<10^{j+1}.
\]
If \(|y|\le T\le q/2\), the argument \(2y/q-3\) belongs to
\([-4,-2]\).  Expand the full polynomial
\(P=\sum_{j=0}^L d_j\varphi_j\), without omitting any coefficient.
Orthogonality and Cauchy--Schwarz give
\begin{equation}\tag{CPT12}
 |P(y)|^2\le
 \left(\sum_{j=0}^L|d_j|^2\right)
 \left(\frac1q\sum_{j=0}^L(2j+1)100^j\right)
 \le\frac{(L+1)^2 100^L}{q}
                       \int_q^{2q}|P(x)|^2dx .
\end{equation}
For \(F=Q_gP\), all its roots give
\[
 |Q_g(y)|\le(T+R)^{d_0}\quad(y\in C),\qquad
 |Q_g(x)|\ge(q-R)^{d_0}\quad(q\le x\le2q).
\]
The upper bound in (CPT7), (CPT12) and the unchanged total Gamma
mass yield
\[
 \int_C|F|^2d\nu_{s,t}
 \le b_*c_\sigma(T+R)^{2d_0}
       \frac{(L+1)^2 100^L}{q}\int_q^{2q}|P(x)|^2dx.
\]
The polynomial lower bound in (CPT7) is applied on the \emph{whole}
real line before its Gamma integral is restricted to \([q,2q]\).
The exact Gamma envelope in (CPT3) gives
\[
 \norm{F}_{s,t}^2
 \ge a_D\frac{c_\Gamma e^{-\pi q}}{\sqrt{1+2q}}
           (q-R)^{2d_0}\int_q^{2q}|P(x)|^2dx .
\]
For \(P\ne0\), division gives \(\mathfrak e_{k,J,D}\).
The trivial central-fraction bound is one.  The zero polynomial
satisfies (CPT11) directly.  This proves (CPT11) with both full
real integration regions present in the norm.

For fixed complete packet and \(D\le2q+1\),
\(\log(b_*/a_D)=O_h(k+\log q)\).  Since \(T+R\) is \(k\) times
a fixed positive constant and \(R/q\to0\), direct logarithms in
(CPT10) give, for its uncapped positive bound,
\begin{equation}\tag{CPT13}
 \log\mathfrak e_{k,J,2q+1}
 =-2(q-J)\log(q/k)+O_h(q+J+k+\log q).
\end{equation}
For every fixed \(0<\vartheta<1\) and \(A>0\), this implies
\[
 q^A\mathfrak e_{k,J,2q+1}\longrightarrow0
 \quad\hbox{uniformly for }0\le J\le(1-\vartheta)q.
\]
The finite formula (CPT10), rather than this asymptotic statement,
is used in all bounds below.  At \(J=q\) the root factor is absent;
the capped bound remains valid but supplies no such decay.

\section{An explicit logarithmic central contrast}

The original central-source calculation ACW11--13
\cite{cpt:centralProvider} proves, on its explicit local-density guard
\(\epsilon_k^{\rm loc}\le1/2\), the following formula.  Let
\(\theta(u)\) invert \(\ell_h'\) on \([-\alpha,\alpha]\), and put
\[
 V(\theta)=\ell_h''(\theta),\qquad
 0<v_*\le V(\theta)\le V_*,
 \qquad c_{\sigma}^{\rm env}(y)
 =\sigma(y)e^{\alpha|y|}\sqrt{1+|y|}.
\]
Then, for \(|u|\le\mu\),
\begin{equation}\tag{CPT14}
\begin{split}
 H(ku)={}&kF_h(u)+\tfrac12\log(|u|+k^{-1})
 -\tfrac12\log(2\pi V(\theta(u)))
 -\log c_{\sigma}^{\rm env}(ku)
 +\log(1+\epsilon_{k,\theta(u)}),\\
 F_h(u)={}&
 \ell_h(\theta(u))-u\theta(u)-\ell_h(\alpha)
                           +\alpha|u|,\qquad
 |\epsilon_{k,\theta}|\le\epsilon_k^{\rm loc}.
\end{split}
\end{equation}
For completeness, all constants and the local-density proof yielding
this provider are spelled out in the appendix.  In particular this
formula retains the actual error, not a selected saddle coefficient.

The function \(F_h\) is even, and on \(0<u<\mu\),
\[
 F_h'(u)=\alpha-\theta(u)>0,\qquad F_h(\mu)=0,\qquad
 F_h(0)=\log M_0-\log M.
\]
Consequently \(-(\log M-\log M_0)\le F_h(u)\le0\).
The exact argument \(|u|+1/k\) lies between \(1/k\) and \(\mu+1\).
Set
\begin{equation}\tag{CPT15}
\begin{split}
 L_k={}&k(\log M-\log M_0)
 +\tfrac12\max\{\log k,|\log(\mu+1)|\}\\
 &+\tfrac12\max\{|\log(2\pi v_*)|,|\log(2\pi V_*)|\}\\
 &+\max\{|\log c_\Gamma|,|\log C_\Gamma|\}
 +2\epsilon_k^{\rm loc}.
\end{split}
\end{equation}
Since \(|\log(1+x)|\le2|x|\) for \(|x|\le1/2\),
(CPT14) proves the finite estimate
\begin{equation}\tag{CPT16}
 |H_C(y)|\le L_k,\qquad L_k=O_h(k+\log k).
\end{equation}
All terms in (CPT15) are nonnegative.  In particular
\(\log M-\log M_0\ge0\) by the original evenness of \(w_h\).

There is also a uniform upper bound on the positive part:
\begin{equation}\tag{CPT17}
 H(y)\le K_h\quad(y\in C),\qquad
 K_h=\left[
 \tfrac12\log(\mu+1)-\tfrac12\log(2\pi v_*)
 -\log c_\Gamma+\log2\right]_+.
\end{equation}
Indeed the \(kF_h\) term is nonpositive, and each other term of
(CPT14) has the corresponding upper bound in (CPT17).

\section{Central derivative, transported metrics and attained sections}
\label{cpt:sec:sections}

Fix \(g,D\) as in Section~\ref{cpt:sec:localization}, and abbreviate
\begin{equation}\tag{CPT18}
 \Lambda_{k,J,D}=L_k\varepsilon_{k,J,D}.
\end{equation}
On \(V_{g,D}\) the actual derivative form is
\[
 \dot h^C_{s,t}(F,G)
 =\partial_s\ip{F}{G}_{s,t}
 =\int_C H(y)\overline{F(y)}G(y)d\nu_{s,t}(y).
\]
The bound (CPT16), Cauchy--Schwarz on \(C\), and (CPT11) prove
\begin{equation}\tag{CPT19}
 |\dot h^C_{s,t}(F,G)|
 \le\Lambda_{k,J,D}\norm{F}_{s,t}\norm{G}_{s,t}.
\end{equation}
In any fixed coefficient frame, write \(A_{s,t}\) for the
positive Gram and \(A^C_{s,t}=\partial_s A_{s,t}\).
Thus (CPT19) is the proved operator inequality
\[
 \norm{A_{s,t}^{-1/2}A^C_{s,t}A_{s,t}^{-1/2}}_{\rm op}
 \le\Lambda_{k,J,D}.
\]
For a fixed nonzero \(F\), its squared norm satisfies
\(|\partial_s\log\norm{F}_{s,t}^2|\le\Lambda_{k,J,D}\).
Integrating from \(s_0\) to \(s_1\) gives
\begin{equation}\tag{CPT20}
 e^{-\Lambda_{k,J,D}|s_1-s_0|}A_{s_0,t}
 \preceq A_{s_1,t}
 \preceq e^{\Lambda_{k,J,D}|s_1-s_0|}A_{s_0,t}.
\end{equation}
This bound holds for every \(t\in[0,1]\).
It requires no smallness assumption on
\((e^{L_k}-1)\varepsilon_{k,J,D}\).

Here is the exact transfer through every fixed restriction and
attained quotient of this polynomial space.  A fixed injective
matrix \(I\) gives the restriction \(I^*A_{s,t}I\); congruence
preserves (CPT20).  A fixed onto matrix
\(J:V_{g,D}\twoheadrightarrow Y\), with a fixed value frame in \(Y\),
has minimum metric and section
\begin{equation}\tag{CPT21}
 G^Y_{s,t}=(J A_{s,t}^{-1}J^*)^{-1},\qquad
 R^Y_{s,t}=A_{s,t}^{-1}J^*G^Y_{s,t}.
\end{equation}
Direct multiplication gives \(JR^Y_{s,t}=I_Y\).
For \(z\in\ker J\),
\[
 z^*A_{s,t}R^Y_{s,t}v=z^*J^*G^Y_{s,t}v=0.
\]
Every representative of \(v\) is \(R^Y_{s,t}v+z\); orthogonality
proves that its squared norm is
\(v^*G^Y_{s,t}v+z^*A_{s,t}z\).
This proves both formulas in (CPT21) and attainment of the unique
minimum.  Taking the infimum over the \emph{same} fibre in each
side of (CPT20) gives the same bounds for \(G^Y\).
Repeating these two operations proves the bounds for any stated
fixed sequence of restrictions and attained quotients inside this
ambient polynomial space.  For a rank-\(r\) resulting metric \(G\),
\begin{equation}\tag{CPT22}
 \left|\log\frac{\det G_{s_1,t}}{\det G_{s_0,t}}\right|
 \le r\Lambda_{k,J,D}|s_1-s_0|.
\end{equation}
The rank is retained, including rank zero with determinant one.

There is a direct estimate for the actual minimizing polynomials.
Let \(P_s=R^Y_{s,t}v\), for \(v\in Y\) fixed and \(t\) fixed.
Because \(JP_s=v\), differentiation gives \(P_s'\in\ker J\).
For every \emph{fixed} \(z\in\ker J\), differentiate its exact
orthogonality to obtain
\[
 \ip{z}{P_s'}_{s,t}=-\dot h^C_{s,t}(z,P_s).
\]
This identity holds for every vector of the fixed kernel at each
parameter, so it may then be evaluated at \(z=P_s'\).
Using (CPT19) and dividing when \(P_s'\ne0\) proves
\begin{equation}\tag{CPT23}
 \norm{P_s'}_{s,t}\le
 \Lambda_{k,J,D}\norm{P_s}_{s,t}.
\end{equation}
The minimum property and (CPT20) give
\(\norm{P_s}_{s,t}\le e^{s\Lambda_{k,J,D}/2}
\norm{P_0}_{0,t}\).
The same metric comparison gives
\(\norm{P_s'}_{0,t}\le e^{s\Lambda_{k,J,D}/2}
\norm{P_s'}_{s,t}\).
Integrating the resulting bound
\(\norm{P_s'}_{0,t}\le
\Lambda_{k,J,D}e^{s\Lambda_{k,J,D}}\norm{P_0}_{0,t}\)
in the fixed coefficient space proves
\begin{equation}\tag{CPT24}
 \boxed{\norm{P_1-P_0}_{0,t}
 \le(e^{\Lambda_{k,J,D}}-1)\norm{P_0}_{0,t}.}
\end{equation}
At \(\Lambda=0\) this formula follows by the zero derivative.
For arbitrary endpoints the same proof, with the oriented segment
reparametrized, gives \(e^{\Lambda|s_1-s_0|}-1\) and the starting
norm.  Thus the original coefficients and phases of the entire
minimum section, rather than only its determinant, are controlled.

\section{The exact map to the original primary submodule}

Let \(A:E\to E\) denote multiplication by the original \(S\).
The polynomial-domain division identity gives
\begin{equation}\tag{CPT25}
 E[g]:=\ker g(A)=(\chi/g)\C[S]/(\chi),\qquad
 \dim_\C E[g]=J.
\end{equation}
Indeed \(g[P]_\chi=0\) says that \(\chi\mid gP\).
Writing \(\chi=g(\chi/g)\) and cancelling the nonzero polynomial
\(g\) in \(\C[S]\) is equivalent to \((\chi/g)\mid P\).
Multiplication by \(\chi/g\) identifies
\(\C[S]/(g)\) bijectively with this submodule, proving the
dimension assertion and retaining every primary multiplicity.

For \(N\ge q-1\), put
\[
 W_{g,N}=(\chi/g)\C[S]_{\le N-q+J},\qquad
 J_{g,N}:W_{g,N}\longrightarrow E[g],\quad
 J_{g,N}P=[P]_\chi.
\]
For \(J=0\) the value space is zero and the formula has that
literal meaning.  The onto property follows from the
degree-\(<J\) representatives in \(\C[S]/(g)\).
Its kernel is exactly \(\chi\C[S]_{\le N-q}\), with negative
degree interpreted as the zero space.  Moreover every
degree-\(\le N\) representative in the original remainder fibre
of a class in \(E[g]\) already lies in \(W_{g,N}\), by the
divisibility just proved.  Therefore minimizing in \(W_{g,N}\)
is precisely the original minimum for that class.

In real coordinates \(U_{c,N}W_{g,N}=V_{g,N}\), because the
constant factor \(i^{q-J}\) changes no subspace.  This gives the
commuting diagram in explicit map form:
\begin{equation}\tag{CPT26}
\begin{gathered}
 W_{g,N}\ \xrightarrow{\ U_{c,N}\ }\ V_{g,N},\\
 J_{g,N}
 =\left(P\longmapsto[P]_\chi\right)\big|_{W_{g,N}},
 \qquad
 J_{g,N}\circ U_{c,N}^{-1}:V_{g,N}\twoheadrightarrow E[g].
\end{gathered}
\end{equation}
If \(G_N[\nu_{s,t}]\) is the original full quotient Gram on \(E\)
and \(I_g\) is a fixed frame inclusion of \(E[g]\) into \(E\),
the minimum metric in (CPT21) is exactly
\begin{equation}\tag{CPT27}
 G^{g}_{N;s,t}=I_g^*G_N[\nu_{s,t}]I_g.
\end{equation}
Equations (CPT20)--(CPT24) consequently apply to the original
restriction, its minimum sections and its fixed subquotients,
with \(D=N\), or with \(D=2q+1\) as one common ambient degree.
In particular, at all four original cutoff degrees the same
explicit large-degree constant is valid.

An opposite terminal eigenpair uses the actual divisor
\(g=(S-\lambda)(S-(k-\lambda))\) and \(J=2\), including when
the full primary length \(e\) exceeds one.  The remaining
\(e-1\) powers at these roots stay in \(\chi/g\).
All terminal eigenlines together use \(g=\rad\chi\) and
\(J=(k+1)^2=q/e\).  For \(m_0>1\) this whole terminal space
has the decay (CPT13).  For \(m_0=1\), this last divisor is
\(\chi\), hence \(J=q\); it is the full quotient.

The exact determinant relation records the complementary metric
that is still present when \(J<q\).  Complete \(I_g\) by any fixed
frame \(I_\perp\) to a frame of \(E\), and in that frame write
\[
 G_N=\begin{pmatrix} A_g&B_g\\B_g^*&C_g\end{pmatrix},
 \qquad A_g=I_g^*G_NI_g.
\]
The identity induced on the quotient \(E/E[g]\) has attained
metric
\begin{equation}\tag{CPT28}
 Q_g^{\rm comp}=C_g-B_g^*A_g^{-1}B_g,\qquad
 \det G_N=\det A_g\,\det Q_g^{\rm comp}.
\end{equation}
To verify it, a value \(z\) in this quotient is represented by
\((x,z)\).  Completing the square in its norm gives the unique
minimum at \(x=-A_g^{-1}B_gz\) and the displayed \(Q_g^{\rm comp}\).
The block triangular elimination has determinant one, proving
the determinant formula.  Thus the map \(E[g]\hookrightarrow
E\twoheadrightarrow E/E[g]\) and its precise complementary
metric remain in the calculation.  A bound on \(A_g\) alone
has not evaluated this second factor.

\section{Original Schur identity and its signed central receiver}
\label{cpt:sec:determinants}

For any source \(\nu_{s,t}\) define the original degree-\(N\)
moment matrix in the \(S\)-monomial frame by
\[
 (H_N)_{ij}=\int_\R
       \overline{(c+iy)^i}(c+iy)^j\,d\nu_{s,t}(y),
 \qquad 0\le i,j\le N.
\]
Let \(\omega_n(s,t)\) be its squared monic source minimum of
degree \(n\), and let \(\eta_r(s,t)\) be the squared monic
relation minimum
\[
 \eta_r(s,t)=
 \min_{\substack{u\ {\rm monic}\\\deg u=r}}
 \int_\R|\chi(c+iy)u(c+iy)|^2\,d\nu_{s,t}(y).
\]
Every preceding source polynomial, respectively preceding
relation \(\chi\C[S]_{<r}\), is included in these minima.
Gram--Schmidt in monic triangular frames gives
\[
 \det H_N=\prod_{n=0}^N\omega_n,\qquad
 \det L_r=\prod_{j=0}^{r-1}\eta_j,
 \qquad
 L_r=\operatorname{Gram}(\chi,\chi S,\ldots,\chi S^{r-1}),
 \quad \det L_0=1.
\]
These are identities of the actual source moments.

The concatenated coefficient frame
\[
 1,S,\ldots,S^{q-1},\chi,\chi S,\ldots,\chi S^{N-q}
\]
is monic triangular, with determinant one.
Its lower Gram block is \(L_{N-q+1}\).  Completing the square
in its relation coefficients yields precisely the remainder
minimum metric
\[
 G_N=(J_NH_N^{-1}J_N^*)^{-1},\qquad
 J_NP=[P]_\chi.
\]
Taking the block determinant proves
\begin{equation}\tag{CPT29}
 \det G_N=
 \frac{\prod_{n=0}^N\omega_n}
      {\prod_{j=0}^{N-q}\eta_j}\quad(N\ge q-1).
\end{equation}
For \(N=q-1\) the denominator is one.
Under (CPT5) a monic degree-\(n\) polynomial becomes one of
leading coefficient \(i^n\); multiplying it by \(i^{-n}\)
returns monicity and leaves its norm unchanged.  The relation
\(\chi u\) has the corresponding factor \(i^{q+r}\).
Hence no determinant modulus or monic norm acquires a phase
factor in (CPT29).

Retain the four signs, the full relation band and its weights:
\begin{equation}\tag{CPT30}
\begin{gathered}
 \Bcal(s,t)=\log\det G_{q-1}+\log\det G_q
             -\log\det G_{2q-1}-\log\det G_{2q},\\
 c_0=c_q=1,\qquad c_r=2\quad(1\le r\le q-1),\\
 \Rcal(s,t)=\sum_{r=0}^q c_r\log\eta_r(s,t),\qquad
 \Scal(s,t)=\sum_{r=0}^q c_r\log\omega_{q+r}(s,t).
\end{gathered}
\end{equation}
Substitute (CPT29) at each of the four cutoffs.
The source prefix sums give coefficient \(-1\) at \(q,2q\)
and \(-2\) at \(q+1,\ldots,2q-1\).
The relation prefix sums give coefficient \(+1\) at \(0,q\)
and \(+2\) at \(1,\ldots,q-1\).  Thus
\begin{equation}\tag{CPT31}
 \boxed{\Bcal(s,t)=\Rcal(s,t)-\Scal(s,t),\qquad
 \sum_{r=0}^q c_r=2q.}
\end{equation}
This independently proves the signs and endpoint weights of
ACW20/SXR10 \cite{cpt:centralProvider,cpt:sourceRelation}.

If a source measure is multiplied by \(A>0\), all these minima
and every quotient metric are multiplied by \(A\).
The quotient determinant at each cutoff is multiplied by
\(A^q\), whose four signed logarithms cancel.
Equivalently each ratio \(\eta_r/\omega_{q+r}\) is unchanged.
Therefore the original arithmetic correction is exactly
\begin{equation}\tag{CPT32}
 \delta_k^\sigma
 :=\Bcal[m_k(y)dy]-\Bcal[\sigma(y)dy]
 =\Bcal(1,1)-\Bcal(0,0).
\end{equation}
This is the proved scalar cancellation of the original \(M^k\).

Use \(J=0,D=2q+1\), \(g=1\) in (CPT9), and abbreviate
\(\Lambda_0=\Lambda_{k,0,2q+1}\).
Every minimizing polynomial \(\chi u_r\), \(0\le r\le q\),
and each of its allowed coefficient variations belongs to
this complete relation space.
Differentiating \(\eta_r\) makes the coefficient-derivative
terms vanish by its defining orthogonality to
\(\chi\C[S]_{<r}\).  The remaining derivative is exactly
\[
 \partial_s\eta_r
 =\int_C H(y)|\chi(c+iy)u_{r;s,t}(c+iy)|^2
                                             d\nu_{s,t}(y).
\]
Equations (CPT11) and (CPT16) consequently prove
\begin{equation}\tag{CPT33}
 |\partial_s\log\eta_r(s,t)|\le\Lambda_0,\qquad
 |\partial_s\Rcal(s,t)|\le2q\Lambda_0.
\end{equation}
The common degree \(2q+1\) includes the output face needed when
one original linear factor multiplies a degree-\(2q\) relation.
For the monic relation band alone, \(D=2q\) also suffices and
gives its correspondingly sharper constant.

Splitting (CPT32) at the actual endpoint \((0,1)\), and using
(CPT31) on its central leg, gives the exact finite receiver
\begin{equation}\tag{CPT34}
\begin{split}
 \boxed{\delta_k^\sigma
 =\Bcal(0,1)-\Bcal(0,0)
 -[\Scal(1,1)-\Scal(0,1)]+E_{\rm rel}},\\
 E_{\rm rel}=\Rcal(1,1)-\Rcal(0,1),\qquad
 |E_{\rm rel}|\le2q\Lambda_0.
\end{split}
\end{equation}
For any function \(F\) on the square, define its oriented
rectangle difference
\(\Delta_\square F=F(1,1)-F(0,1)-F(1,0)+F(0,0)\).
Both central relation legs obey (CPT33), so
\begin{equation}\tag{CPT35}
 \boxed{\Delta_\square\Bcal+\Delta_\square\Scal
       =\Delta_\square\Rcal,\qquad
 |\Delta_\square\Rcal|\le4q\Lambda_0.}
\end{equation}
No claim of a small mixed derivative follows from this integrated
rectangle estimate; its derivation controls the two actual legs.

The sign estimate (CPT17) propagates through every monic
source minimum because
\(d\nu_{1,t}\le e^{K_h}d\nu_{0,t}\).
Thus
\(\Scal(1,t)-\Scal(0,t)\le2qK_h\).
Combining this with (CPT31),(CPT33) proves
\begin{equation}\tag{CPT36}
 \Bcal(1,t)-\Bcal(0,t)\ge-2qK_h-2q\Lambda_0
 \quad(0\le t\le1).
\end{equation}
At fixed packet (CPT13) implies \(q^A\Lambda_0\to0\) for every
fixed \(A>0\).  In particular the lower logarithmic-scale bound
\[
 \liminf_{k\to\infty}
 \frac{\Bcal(1,t)-\Bcal(0,t)}{q\log k}\ge0
\]
holds uniformly for \(t\in[0,1]\).
It leaves the source-band value and the exterior leg in
(CPT34) explicitly present.

\section{Exact mixed source-band kernel and complete cross terms}

Write \(\Delta_N(s,t)=\det H_N(s,t)\).
The monic source-product identity above gives
\begin{equation}\tag{CPT37}
 \Scal=
 \log\Delta_{2q-1}+\log\Delta_{2q}
          -\log\Delta_{q-1}-\log\Delta_q.
\end{equation}
Let \(f_N(y)=(1,c+iy,\ldots,(c+iy)^N)^T\).
With the conjugate-linear-first Gram convention,
\[
 H_N=\int\overline{f_N(y)}f_N(y)^T\,d\nu_{s,t}(y).
\]
Define its reproducing kernel by the original coefficient formula
\begin{equation}\tag{CPT38}
 K_N^{s,t}(y,z)
   =f_N(y)^T H_N(s,t)^{-1}\overline{f_N(z)}.
\end{equation}
If \(F(y)=f_N(y)^Ta\), matrix multiplication gives
\[
 \int K_N^{s,t}(y,z)F(z)d\nu_{s,t}(z)=F(y).
\]
Also \(K_N^{s,t}(z,y)=\overline{K_N^{s,t}(y,z)}\).
This verifies the kernel and its conjugations directly in the
original complex coefficient frame.

At each real point \(H_CH_O=0\).  Hence
\(\partial_s\partial_tH_N=0\), while
\[
 H_{N,C}=\int_C H(y)\overline{f_N(y)}f_N(y)^T\,d\nu_{s,t}(y),
 \quad
 H_{N,O}=\int_O H(z)\overline{f_N(z)}f_N(z)^T\,d\nu_{s,t}(z).
\]
Differentiating \(H_N^{-1}H_N=I\) gives
\(\partial_tH_N^{-1}=-H_N^{-1}H_{N,O}H_N^{-1}\).
The determinant derivative follows by differentiating the
finite determinant polynomial and using its adjugate:
\(\partial_s\log\det H_N=\Tr(H_N^{-1}H_{N,C})\).
Consequently
\begin{equation}\tag{CPT39}
\begin{split}
 \partial_s\partial_t\log\Delta_N
 &=-\Tr(H_N^{-1}H_{N,C}H_N^{-1}H_{N,O})\\
 &=-\int_C\int_O H(y)H(z)|K_N^{s,t}(y,z)|^2
                          d\nu_{s,t}(y)d\nu_{s,t}(z).
\end{split}
\end{equation}
For the second equality, expand the two integrals and use the
trace of the rank-one products.  The resulting scalar product
is \(K_N(y,z)K_N(z,y)=|K_N(y,z)|^2\).
Every integral is absolutely convergent by the moment bounds
proved after (CPT8); the finite inverse Gram introduces only
finite coefficients on the compact parameter square.

Combining all four terms of (CPT37) yields
\begin{equation}\tag{CPT40}
\begin{split}
 \partial_s\partial_t\Scal(s,t)
 =-\int_C\int_O H(y)H(z)\big[
 &|K_{2q-1}^{s,t}(y,z)|^2+|K_{2q}^{s,t}(y,z)|^2\\
 &-|K_{q-1}^{s,t}(y,z)|^2-|K_q^{s,t}(y,z)|^2
 \big]\,d\nu_{s,t}(y)d\nu_{s,t}(z).
\end{split}
\end{equation}
This is an exact expression for the source-band interaction,
with its original two consecutive windows.  Integrating over
the parameter square gives \(\Delta_\square\Scal\).
Equations (CPT35),(CPT40) therefore return the full signed
rectangle to this explicit mixed kernel integral with the
proved remainder \(4q\Lambda_0\).
Disjoint scalar supports made \(\partial_s\partial_tH_N=0\);
they do not make the two noncommuting compressed matrices
\(H_{N,C},H_{N,O}\) have zero product.

The relation determinant has the same exact formula.
If \(B_N\) is the full coefficient matrix of
\(\chi,\chi S,\ldots,\chi S^{N-q}\), its Gram and derivatives
are \(B_N^*H_NB_N,B_N^*H_{N,C}B_N,B_N^*H_{N,O}B_N\).
The relation version of (CPT39) therefore contains its complete
cross matrix.  Subtracting it from the source version in
(CPT29) gives the exact quotient mixed derivative.  At
\(N=q-1\) the relation matrix is empty, its determinant is one,
and its derivative is zero.  This proves the relation-to-source
and source-to-quotient maps without omitting the zero-cutoff
endpoint.

\section{Propagation into the earlier central statements}

Equations (CPT10)--(CPT11) reproduce the root-sensitive finite
annihilator localization C3--C7 from the supplied complete
proofs \cite{cpt:complete}, now uniformly on the entire square.
For the replacement \(\nu_{1,t}\to\nu_{0,t}\), the earlier
endpoint estimate multiplied that localization by
\(e^{L_k}-1\) and required a final smallness guard.
The proved derivative argument replaces its use by
\[
 \Lambda=L_k\varepsilon_{k,J,D},\quad
 e^{-\Lambda}G_{0,t}\preceq G_{1,t}\preceq e^\Lambda G_{0,t},
 \quad
 \norm{P_1-P_0}_{0,t}\le(e^\Lambda-1)\norm{P_0}_{0,t}.
\]
These are valid finite inequalities for every admitted
\(k\) satisfying the displayed density, degree and root guards.
The old endpoint inequalities remain preserved statements
on their own domains.

The previous central-leg relation radius in C17 is replaced,
for this exact interpolation, by \(2qL_k\varepsilon_{k,0,2q+1}\).
The uniform one-sided conclusion C21--C22 now has the radius
in (CPT36) for every exterior parameter.
The additional rectangle identity is (CPT35), with the
literal mixed kernel (CPT40).  These substitutions propagate
new constants and new path information; they do not evaluate
an uncomputed source-band integral.

For each localized original submodule and its actual fixed
subquotients, (CPT27) proves the quantitative correspondence.
For the complementary quotient, (CPT28) records its precise
Schur metric.  For the full signed determinant, (CPT34) and
(CPT40) retain the actual source-band and exterior quantities.
Thus no seven absolute allocation, proper-source
standard-window value, individual projected sign, or
Riemann-hypothesis conclusion is supplied by a substituted
rank or a two-class comparison in this note.

\par\bigskip\noindent\textbf{Appendix}\par
\section{The retained local-density constants and proof}

This appendix states and verifies the particular central-density
provider used in (CPT14)--(CPT16), with the exact original constants
of ACW4--11 \cite{cpt:centralProvider}.  It uses the two original
source envelopes recorded below (CPT3), rather than any new
analytic hypothesis.
Put \(\eta=1/4\) and \(C_h^{\rm tilt}=C_{\rm tilt}\).
For \(|\theta|\le\alpha\), define the auxiliary density
\[
 p_\theta(y)=e^{\theta y}w_h(y)/M_h(\theta),\qquad
 \mu(\theta)=\ell_h'(\theta).
\]
Its unit mass is used only in this Fourier computation.  The
exact convolution identity is
\begin{equation}\tag{CPT41}
 m_k(x)=M_h(\theta)^k e^{-\theta x}p_\theta^{*k}(x).
\end{equation}
It follows by pulling the convolution back to its \(k\)
original variables, whose sum is \(x\) with Jacobian one.
Thus the original arithmetic mass is explicitly restored.

The envelope and \(M_h(\theta)\ge M_0\) give
\[
 V_*=\frac{2C_{\rm tilt}}{M_0}\operatorname B(3,p-3),
 \qquad
 H_*=
 2^{1+\eta}\left[
 \frac{2C_{\rm tilt}}{M_0}\operatorname B(3+\eta,p-3-\eta)
                      +\mu^{2+\eta}\right].
\]
Integrating \(y^2(1+|y|)^{-p}\), and using
\(|x-y|^{2+\eta}\le2^{1+\eta}(|x|^{2+\eta}+|y|^{2+\eta})\),
proves \(V(\theta)\le V_*\) and
\(\int|y-\mu(\theta)|^{2+\eta}p_\theta(y)dy\le H_*\).
The beta integrals converge, since \(p\ge7/2>3+\eta\).

On \([-1,1]\), (CPT41) at \(k=3\) and the original lower
convolution envelope give
\[
 p_\theta^{*3}(x)\ge
 c_h e^{-2\alpha}2^{-B_0}/M^3.
\]
Define
\[
 a_{\rm mix}=\min\{1/4,c_he^{-2\alpha}2^{-B_0}/M^3\},\quad
 \beta_{\rm mix}=2a_{\rm mix},\quad v_*=\beta_{\rm mix}/9.
\]
The triple convolution is the sum of \(\beta_{\rm mix}\)
times the uniform density on \([-1,1]\), and a nonnegative
residual of mass \(1-\beta_{\rm mix}\).
The variance of a mixture is at least the mixture-weighted
within-component variances, as follows by expanding each
centred square about the full mean.
It is therefore at least \(\beta_{\rm mix}/3\).
The variance of the three independent summands is
\(3V(\theta)\), proving \(V(\theta)\ge v_*>0\).
Continuity and differentiation of \(M_h\) through order two
follow from the same integrable envelope.
Consequently \(\mu(\theta)\) is strictly increasing and maps
\([-\alpha,\alpha]\) onto \([-\mu,\mu]\), establishing the
inverse used in (CPT14).

For the centred characteristic function
\(\phi_\theta(\zeta)=\int
 e^{i\zeta(y-\mu(\theta))}p_\theta(y)dy\),
Taylor's formula on \(|u|\le1\), and
\(2+|u|+u^2/2\le(7/2)|u|^{2+\eta}\) on \(|u|>1\), imply
\[
 |\phi_\theta(\zeta)-1+V(\theta)\zeta^2/2|
       \le4H_*|\zeta|^{2+\eta}.
\]
Put
\[
 \zeta_0=\min\{1,(2V_*)^{-1/2},
                  (v_*/(16H_*))^{1/\eta}\}.
\]
For \(|\zeta|\le\zeta_0\), the quadratic term is nonnegative
and the Taylor error is at most \(v_*\zeta^2/4\).  Hence
\(|\phi_\theta(\zeta)|\le e^{-v_*\zeta^2/4}\).
The triple-density mixture gives
\[
 |\phi_\theta(\zeta)|^3
 \le1-\beta_{\rm mix}
           +\beta_{\rm mix}|\sin\zeta/\zeta|.
\]
For \(|\zeta|\ge\zeta_0\) one has
\(|\sin\zeta/\zeta|\le1-\zeta_0^2/7\).
To check the constant, on \([\zeta_0,2]\) the positive
function \(\sin u/u\) decreases, and its Taylor upper bound
at \(0<\zeta_0\le1\) is
\(1-\zeta_0^2/6+\zeta_0^4/120
\le1-\zeta_0^2/7\).  Beyond \(2\), its absolute value is
at most \(1/2\), which is smaller.
Thus, with
\[
 r_*=(1-\beta_{\rm mix}\zeta_0^2/7)^{1/3}<1,\qquad
 C_2=\frac{4\pi(C_{\rm tilt}/M_0)^2}{2p-1},
\]
we have \(|\phi_\theta|\le r_*\) off the small interval.
The Fourier Plancherel identity \cite[(1.14.8)]{cpt:dlmfFourier}
in the convention
\(\widehat f(\zeta)=\int e^{i\zeta y}f(y)dy\) gives
\(\int|\phi_\theta|^2\le C_2\).
Relative to the DLMF convention
\(\mathcal F f=(2\pi)^{-1/2}\int e^{i\zeta y}f(y)dy\),
our transform is \(\sqrt{2\pi}\mathcal F f\).  Thus its
squared \(L^2\) norm is \(2\pi\int|f|^2\); integrating
the squared density envelope yields exactly \(C_2\).

Let
\[
 C_0=4H_*+V_*^2/8,\qquad a_k^{\rm loc}=v_*(k-1)/4,
\]
and define the finite error
\begin{equation}\tag{CPT42}
\begin{split}
 A_k^{\rm loc}=\frac1{2\pi}\bigg[
 &C_0k\,\Gamma((3+\eta)/2)
                    (a_k^{\rm loc})^{-(3+\eta)/2}
 +C_2r_*^{\,k-2}\\
 &+\frac{2e^{-kv_*\zeta_0^2/2}}{kv_*\zeta_0}
 \bigg],\qquad
 \epsilon_k^{\rm loc}=\sqrt{2\pi kV_*}\,A_k^{\rm loc}.
\end{split}
\end{equation}
On the small interval,
\[
 |\phi_\theta(\zeta)-e^{-V(\theta)\zeta^2/2}|
 \le C_0|\zeta|^{2+\eta}.
\]
This follows from the previous Taylor bound and the exponential
remainder \(V_*^2\zeta^4/8\), using
\(|\zeta|^4\le|\zeta|^{2+\eta}\) when \(|\zeta|\le1\).
Factor the difference of the \(k\)-th powers as a sum of \(k\)
terms.  Both factors have modulus at most
\(e^{-v_*\zeta^2/4}\), so its integral is bounded by
\[
 C_0k\int_\R|\zeta|^{2+\eta}
             e^{-a_k^{\rm loc}\zeta^2}\,d\zeta
 =C_0k\,\Gamma((3+\eta)/2)
                      (a_k^{\rm loc})^{-(3+\eta)/2}.
\]
The characteristic-function integral on the complementary
region is at most \(C_2r_*^{k-2}\), by its \(L^2\) bound.
The corresponding Gaussian tail is at most
\(2e^{-kv_*\zeta_0^2/2}/(kv_*\zeta_0)\), by replacing
the factor \(|\zeta|/\zeta_0\) with its lower bound one
before integrating the Gaussian derivative.
Fourier inversion therefore yields
\begin{equation}\tag{CPT43}
 \left|p_\theta^{*k}(k\mu(\theta))
           -(2\pi kV(\theta))^{-1/2}\right|
 \le A_k^{\rm loc}.
\end{equation}
All bounds are uniform for the closed interval
\(|\theta|\le\alpha\), including its two endpoints.
The use of inversion is justified by
\(\phi_\theta^k\in L^1\), already implied by the \(L^2\)
bound and \(|\phi_\theta|\le1\) for \(k\ge2\).
Multiplying (CPT43) by \(\sqrt{2\pi kV(\theta)}\) proves
the relative error in (CPT14), with
\(\epsilon_k^{\rm loc}=O_h(k^{-1/8})\).
In particular its stated finite guard is eventually satisfied.

Finally apply (CPT41) with \(\theta=\theta(u)\),
insert the relative-error expression, and divide by
\(M^k\sigma(ku)\).  Writing the Gamma density as
\[
 \sigma(ku)=c_{\sigma}^{\rm env}(ku)
                 e^{-\alpha k|u|}(1+k|u|)^{-1/2}
\]
gives (CPT14) by exact logarithms.
This completes the density-to-contrast calculation used by
the new path estimates.



% END PRESERVED INPUT 18_CENTRAL_PATH_TRANSFER.tex
\endgroup

\clearpage
\begingroup
\def\C{\mathbb C}
\def\R{\mathbb R}
\def\Pcal{\mathcal P}
\def\norm#1{\left\lVert#1\right\rVert}
\def\abs#1{\left|#1\right|}

\setcounter{equation}{0}
\renewcommand{\theequation}{JW.\arabic{equation}}
\section*{Joint endpoint estimates and localized two-window transport}
\addcontentsline{toc}{section}{Joint endpoint estimates and localized two-window transport}
% BEGIN PRESERVED INPUT 06_JOINT_LIMIT_AND_WINDOW_PROOF.tex SHA256 23bb18f459cd0a49237f6b66984053185c665c0a81da5b887c5fa2fdedcfcf64



\paragraph{Source and precise scope.}
This note verifies and supplies finite constants for the deductions labelled
J1--J10 and S15--S16 in the supplied \texttt{COMPLETE\_PROOFS.md}, dated
18 September 2026.  The source bytes have SHA-256
\path{e4e54261c91ae899a78956383482650f88933024f5f3cec6b41332ceeb5a4172}.
The endpoint block is at lines 799--966; the two-window block is at
lines 746--770.  This note retains every singular coefficient, the original
Lebesgue product measure, the full polynomial relation, and both boundary
windows.  The incoming endpoint decomposition, the actual tilt envelope,
and the old window coefficients are explicitly identified as input results.
The named earlier files \texttt{TWO\_WINDOW\_SAME\_CLASS\_DEFECT.md} and the
theta-endpoint note were not located in the bounded intake.  Their numerical
window coefficient and full theta existence theorem are therefore not
claimed as independently recovered here.  The finite deductions below are
proved from the actual formulas recorded in the supplied source; no missing
allocation or projected arithmetic sign is assumed.

\section{The entire remainder has an explicit uniform norm allowance}

Write the original quartet as
\[
\mathcal Z=\{1/2+\epsilon\delta+i\epsilon'\gamma:
                 \epsilon,\epsilon'\in\{-1,1\}\},
\qquad 0<\delta<1/2,\quad\gamma>2,
\]
with common multiplicity \(m\), \(d=m-1\), and retain the given inverse
primitive
\[
\psi(x)=s_h(x)+R_h(x),\qquad
s_h(x)=\sum_{\rho\in\mathcal Z}\sum_{j=0}^{d}
a_{\rho,j}x^{-\rho}\frac{\log(1/x)^j}{j!}.
\tag{JL1}
\]
Here \(R_h\) is the actual entire even remainder, not a freely chosen
regularization.  If
\(R_h(x)=\sum_{n\ge0}r_{2n}x^{2n}\), put
\[
B_h=\sum_{n\ge0}|r_{2n}|<\infty,\qquad
C_h^{\rm reg}=B_h^2+
2B_h\sum_{\rho,j}
\frac{|a_{\rho,j}|}{(1-\Re\rho)^{j+1}}.
\tag{JL2}
\]
The sum defining \(B_h\) converges absolutely because the power series
is entire.  One can alternatively use the possibly smaller
\(\sup_{0\le x\le1}|R_h(x)|\).  These constants refer to that same remainder
and leave all singular coefficients unchanged.

\begin{proposition}
For every \(L\ge0\),
\[
\left|\int_{e^{-L}}^1|\psi(x)|^2\,dx
      -\int_{e^{-L}}^1|s_h(x)|^2\,dx\right|
\le C_h^{\rm reg}.
\tag{JL3}
\]
\end{proposition}
\begin{proof}
The integrand difference is
\(2\Re(s_h\overline{R_h})+|R_h|^2\).  Since
\(\Re\rho<1\),
\[
\int_0^1 x^{-\Re\rho}\frac{\log(1/x)^j}{j!}\,dx
=\frac1{(1-\Re\rho)^{j+1}}.
\]
This identity follows from \(x=e^{-t}\) and \(j\) integrations by parts
in \(\int_0^\infty e^{-(1-\Re\rho)t}t^j\,dt\).
Take absolute values of each of the finitely many singular terms,
use \(|R_h|\le B_h\), and enlarge the integration interval to \((0,1)\).
This proves (JL3).  In particular the cross terms with the entire
remainder do not grow with the endpoint cutoff.
\end{proof}

\section{All factorials in the oscillating endpoint leading term}

Let
\(\rho_+=1/2+\delta+i\gamma\) and
\(\mathcal R=\{\rho_+,\overline{\rho_+}\}\).
The local endpoint coefficient convention in the supplied source is
\[
a_{\rho,j}=[z^{d-j}]
\frac{p_\delta(\rho+z)}{2z_\rho(z)},\qquad
p_\delta(s)=(s-1/2)^2+\gamma^2-\delta^2,\qquad
\zeta(\rho+z)=z^m z_\rho(z).
\tag{JL4}
\]
Thus
\[
a_{\rho_+,d}
=\frac{2i\delta\gamma}{2z_{\rho_+}(0)}
=\frac{m!\,i\delta\gamma}{\zeta^{(m)}(\rho_+)}
=:A_h\ne0,\qquad
a_{\overline{\rho_+},d}=\overline{A_h}.
\tag{JL5}
\]
This is the top Laurent coefficient in the recorded convention; no
phase has been made real by a coordinate change.

Put \(I_h(L)=\int_{e^{-L}}^1|\psi|^2\,dx\).
The change of variables \(x=e^{-t}\) gives the exact finite identity
\[
I_h^{\rm sing}(L)=
\sum_{\rho,\sigma\in\mathcal Z}\sum_{j,l=0}^d
\frac{a_{\rho,j}\overline{a_{\sigma,l}}}{j!\,l!}
E_{j+l}(\rho+\bar\sigma-1,L),\qquad
E_p(z,L)=\int_0^L e^{zt}t^p\,dt.
\tag{JL6}
\]
There is no division by \(z\) in this definition, so the terms with
\(z=0\), including the reflected-root pairs, remain present.
For \(z\ne0\), repeated integration by parts gives
\[
E_p(z,L)=e^{zL}\sum_{r=0}^p
\frac{(-1)^r p!L^{p-r}}{(p-r)!z^{r+1}}
-\frac{(-1)^p p!}{z^{p+1}}.
\tag{JL7}
\]

For every ordered right/right pair and \(p=j+l\) put
\[
\begin{split}
A_1={}&\sum_{\rho,\sigma\in\mathcal R}\sum_{j,l=0}^d
\frac{|a_{\rho,j}a_{\sigma,l}|}{j!\,l!}
\sum_{\substack{0\le r\le j+l\\(j,l,r)\ne(d,d,0)}}
\frac{(j+l)!}{(j+l-r)!\,|\rho+\bar\sigma-1|^{r+1}},\\
A_2={}&C_h^{\rm reg}+
\sum_{\rho,\sigma\in\mathcal R}\sum_{j,l=0}^d
\frac{|a_{\rho,j}a_{\sigma,l}|}{j!\,l!}
\frac{(j+l)!}{|\rho+\bar\sigma-1|^{j+l+1}},\\
A_3={}&
\sum_{(\rho,\sigma)\in\mathcal Z^2\setminus\mathcal R^2}
\sum_{j,l=0}^d
\frac{|a_{\rho,j}a_{\sigma,l}|}{j!\,l!\,(j+l+1)}.
\end{split}
\tag{JL8}
\]
These are exactly J2 with (JL2) as a fully displayed valid regular
allowance.  Define
\[
F_h(L)=\frac1{(d!)^2}
\left(\frac{|A_h|^2}{\delta}
+\Re\frac{A_h^2e^{2i\gamma L}}{\delta+i\gamma}\right),
\qquad
f_*=\frac{|A_h|^2}{(d!)^2}
\left(\frac1\delta-\frac1{\sqrt{\delta^2+\gamma^2}}\right)>0.
\tag{JL9}
\]

\begin{proposition}
For every \(L\ge1\),
\[
F_h(L)\ge f_*,
\qquad
\left|\frac{I_h(L)}{e^{2\delta L}L^{2d}}-F_h(L)\right|
\le \frac{A_1}{L}+(A_2+A_3L)e^{-2\delta L}.
\tag{JL10}
\]
\end{proposition}
\begin{proof}
For the two equal right/right roots, the \(j=l=d,r=0\) terms in
(JL7) sum to
\(|A_h|^2/(\delta(d!)^2)\) after dividing by the displayed scale.
The other two ordered right/right pairs sum to
\(\Re(A_h^2e^{2i\gamma L}/(\delta+i\gamma))/(d!)^2\).
These are precisely the terms in (JL9).
Every remaining upper endpoint term in (JL7), after that division,
contains \(L^{j+l-r-2d}\) with exponent at most \(-1\).
The sum of its absolute coefficients is \(A_1\).
The lower endpoint terms and (JL3) are bounded by
\(A_2e^{-2\delta L}\), because \(L^{-2d}\le1\).
For every pair outside \(\mathcal R^2\), \(\Re z\le0\), and therefore
\[
|E_{j+l}(z,L)|\le \frac{L^{j+l+1}}{j+l+1}.
\]
Its scaled contribution is at most the corresponding coefficient in
\(A_3Le^{-2\delta L}\).
The lower bound for \(F_h\) follows by taking the modulus of its
oscillating summand.  This proves all assertions, retaining all ordered
cross terms.
\end{proof}

\paragraph{Joint tensor/cutoff error.}
Put
\[
\eta_h(L)=
\frac{A_1/L+(A_2+A_3L)e^{-2\delta L}}{f_*}.
\]
If \(\eta_h(L)\le1/2\), the actual ratio
\(I_h(L)/(e^{2\delta L}L^{2d}F_h(L))\) lies in
\([1-\eta_h(L),1+\eta_h(L)]\).
For real \(|x|\le1/2\), integration of \(1/(1+t)\) from \(0\) to \(x\)
proves \(|\log(1+x)|\le2|x|\).  Consequently
\[
\left|\log
\frac{I_h(L)^k}{e^{2\delta kL}L^{2dk}F_h(L)^k}\right|
\le2k\eta_h(L),
\tag{JL11}
\]
with relative error at most \(e^{2k\eta_h(L)}-1\).
At \(L=k^2\), the exponent is \(O_h(k^{-1})\), so
\[
I_h(k^2)^k=e^{2\delta k^3}k^{4dk}F_h(k^2)^k
\bigl(1+O_h(k^{-1})\bigr).
\tag{JL12}
\]
When \(m=1\), the empty sum \(A_1\) is zero.  At
\(L=(\log k)/\delta\ge1\), the exponent in (JL11) is at most
\[
\frac{2}{f_*k}\left(A_2+\frac{A_3}{\delta}\log k\right).
\tag{JL13}
\]
Fubini's theorem gives exactly \(I_h(L)^k\) for the fixed original
product representative on \([e^{-L},1]^k\).
These estimates concern that specified representative.  They do not
evaluate a minimum over added source corrections.

\section{The literal diagonal map and the terminal factorial}

For \(b_\rho x^{-\rho}\), \(\Re\rho=1/2+\delta\), the unchanged product
measure gives
\[
N_{\rm prod}(k,L)=|b_\rho|^{2k}
\left(\frac{e^{2\delta L}-1}{2\delta}\right)^k.
\tag{JL14}
\]
The literal diagonal map substitutes all \(x_i=x\), hence its value is
\(b_\rho^kx^{-k\rho}\).  Its norms in the two expressly specified measures are
\[
\begin{split}
N_{\rm diag,dx}(k,L)&=|b_\rho|^{2k}
\frac{e^{(k-1+2k\delta)L}-1}{k-1+2k\delta},\\
N_{\rm diag,k}(k,L)&=|b_\rho|^{2k}
\frac{e^{2k\delta L}-1}{2k\delta}
\quad\text{in }x^{k-1}\,dx.
\end{split}
\tag{JL15}
\]
Integrating the respective powers of \(x\) proves these formulas.
Dividing (JL14) by the second formula gives the exact discrepancy
\[
\frac{N_{\rm prod}}{N_{\rm diag,k}}
=k(2\delta)^{1-k}
\frac{(1-e^{-2\delta L})^k}{1-e^{-2k\delta L}}.
\tag{JL16}
\]

The common terminal factorial can also be checked without dropping the
original unit.  In one local algebra \(\C[z]/(z^m)\), the actual terminal
vector is \(z^d\); its endpoint coefficient is
\(b_\rho=1/(2z_\rho(0))\) in convention (JL4).
In the extremal ordered tuple, put \(n=kd\) and
\(N_\Sigma=z_1+\cdots+z_k\).  The multinomial theorem and \(z_i^{d+1}=0\)
give the exact polynomial identity
\[
N_\Sigma^{kd}=\frac{(kd)!}{(d!)^k}\prod_{i=1}^k z_i^d.
\tag{JL17}
\]
Multiplying by the full \(U_h^{\otimes k}\) kills only its positive-degree
terms in this terminal vector, so its endpoint is
\[
\frac{(kd)!}{(d!)^k}\,U_h(\rho)^k b_\rho^k
\prod_{i=1}^k x_i^{-\rho}.
\tag{JL18}
\]
Both product and diagonal squared norms therefore acquire
\(\left|(kd)!/(d!)^k\right|^2|U_h(\rho)|^{2k}\).
This verifies the sentence following J8 for that exact terminal vector.
No factorial is absorbed into an unnamed choice of basis.

\section{Uniform propagation of moments to the closed tilt interval}

Retain the actual nonnegative one-packet density \(w_h\) and its stated
envelope
\[
e^{\alpha|y|}w_h(y)\le C_h^{\rm tilt}(1+|y|)^{-p},
\qquad \alpha=\pi/2,\quad p=8m-9/2>3.
\tag{JL19}
\]
Let \(\mu_n=\int_\R y^nw_h(y)\,dy\), without changing the source mass.
For \(j=0,1,2\),
\[
\sum_{r\ge0}\frac{\alpha^r}{r!}
\int_\R |y|^{r+j}w_h(y)\,dy
=\int_\R |y|^je^{\alpha|y|}w_h(y)\,dy<\infty.
\tag{JL20}
\]
Tonelli's theorem applies to this nonnegative majorant.  It proves
uniform absolute convergence on \(0\le\theta\le\alpha\), interchange
with the signed integrands, and continuity of the derivatives at
the endpoint:
\[
M_h^{(j)}(\theta)
=\int_\R y^je^{\theta y}w_h(y)\,dy
=\sum_{r\ge0}\mu_{r+j}\frac{\theta^r}{r!}.
\tag{JL21}
\]

For an integer cutoff \(L\ge0\), let \(n=L+1\) and
\(T_L=n/(2e\alpha)\).  For \(z\ge0\) put
\[
\mathcal Q_n(z)=e^{-z}\sum_{r=n}^{\infty}\frac{z^r}{r!}.
\]
For \(0<z\le n/(2e)<n\), multiplication of each term by
\((n/z)^{r-n}\ge1\) yields
\[
\mathcal Q_n(z)
\le (z/n)^n e^{n-z}
\le(ez/n)^n\le2^{-n}.
\tag{JL22}
\]
For \(z=0\) the sum is zero; globally it is at most one.
The absolute tail of (JL21) is bounded by
\[
\int_\R |y|^jw_h(y)e^{\alpha|y|}
\mathcal Q_n(\alpha|y|)\,dy.
\]
Use \(2^{-n}\) on \(|y|\le T_L\), and one elsewhere.
Since
\[
\int_0^\infty y^j(1+y)^{-p}\,dy
=\operatorname B(j+1,p-j-1),\qquad
\int_{T_L}^\infty y^j(1+y)^{-p}\,dy
\le\frac{(1+T_L)^{j-p+1}}{p-j-1},
\]
the resulting explicit bound is
\[
\sup_{0\le\theta\le\alpha}
\left|M_h^{(j)}(\theta)
-\sum_{r=0}^{L}\mu_{r+j}\frac{\theta^r}{r!}\right|
\le2C_h^{\rm tilt}\left[
2^{-(L+1)}\operatorname B(j+1,p-j-1)
+\frac{(1+T_L)^{j-p+1}}{p-j-1}\right].
\tag{JL23}
\]
The beta identity follows by \(u=y/(1+y)\) in its defining integral.
Thus no fourth tilted moment is used.  These formulas verify J9--J10
but do not by themselves verify the separately cited arithmetic
moment recurrence or evaluate a growing projection-kernel integral.

\section{The localized ambient space really contains the window projections}

Here restore the simple target order in S15: \(a=k+4\),
\(a\equiv1\bmod4\), \(q=(a+1)^2\), \(c=a/2\), and
\[
\chi(S)=\prod_{r,s=0}^{a}
\bigl[S-c-(2r-a)\delta-i(2s-a)\gamma\bigr].
\tag{JW1}
\]
The actual eigenvalue \(\lambda\) is one of these roots.
Set \(D=2q+1\),
\[
\mathcal W=\frac{\chi(S)}{S-\lambda}\Pcal_{D-q+1},
\quad
\mathcal R_N=\chi\Pcal_{N-q}
\quad (\Pcal_j=\{0\}\text{ for }j<0).
\tag{JW2}
\]
All norms are integrals on the original line \(S=c+iy\).
Let \(\omega\) be the actual arithmetic density and
\(\widetilde\omega\) its specified central replacement at order \(a\),
center \(c\), and interval \([-a\mu_h,a\mu_h]\).
The central-localization calculation C3--C13 gives on this \emph{entire}
fixed space the relative form bound
\[
|(\widetilde\omega-\omega)(f,f)|\le\tau\|f\|_\omega^2
\quad(f\in\mathcal W),\qquad 0\le\tau<1.
\tag{JW3}
\]
For clarity, the valid finite parameter is
\[
\tau=B_a\min\{1,\mathfrak e_{a,1,2q+1}(\omega)\},
\qquad
B_a=e^{H_*}-1,
\tag{JW4}
\]
with \(H_*\) and \(\mathfrak e\) the actual constants in C3 and C13.
The guards of those expressions, including \(R<q\) and \(a\mu_h\le q/2\),
are retained.  Self-adjointness and polarization of (JW3) yield
\[
|(\widetilde\omega-\omega)(f,g)|
\le\tau\|f\|_\omega\|g\|_\omega.
\tag{JW5}
\]

\paragraph{Exact fixed relation band.}
Let \(S_Nv\) be the original attained polynomial representative of \(v\)
at cutoff \(N\).  It is orthogonal to \(\mathcal R_N\).
Every representative of the chosen eigenclass belongs to \(\mathcal W\);
all its differences belong to \(\mathcal R_N\).
For \(n=q-1\) or \(n=q\), put
\[
F_n=S_{n+q}v-S_nv,\qquad
\mathcal B_n=\mathcal R_{n+q}\cap\mathcal R_n^{\perp_\omega}.
\tag{JW6}
\]
Then \(F_n\in\mathcal B_n\), and the \(\omega\)-orthogonal projector onto
this band is \(B_n=P_{\mathcal R_{n+q}}-P_{\mathcal R_n}\).
The analogous statement holds in the replaced metric.
The \emph{full} section-difference map from the cyclic module to
\(\mathcal B_n\) has rank \(q\), not merely on the single eigenline.
Indeed
\(\dim\mathcal R_{n+q}-\dim\mathcal R_n=q\).
If the section difference of any class is zero, its common representative
\(P\), of degree at most \(n\), is orthogonal to every high relation.
In particular it is orthogonal to \(\chi P\), whose degree is at most
\(n+q\).

To check the sign used here, pair each root in (JW1) with the root
having opposite real displacement and the same imaginary displacement.
On \(S=c+iy\), each pair contributes
\(-[(y-(2s-a)\gamma)^2+(2r-a)^2\delta^2]\).
There are \(q/2\) pairs, and \(q/2\) is even because \(a\equiv1\bmod4\).
Every \((2r-a)\delta\) is nonzero because \(a\) is odd.
Thus \(\chi(c+iy)>0\) for every real \(y\).  For \(P\ne0\),
\(\int\chi(c+iy)|P(c+iy)|^2d\omega(y)>0\), contradicting orthogonality.
The map is injective and hence fills the band.

\begin{lemma}[Explicit fixed-subspace projection bound]
For any fixed \(\mathcal R\subseteq\mathcal W\), let \(P,\widetilde P\)
be its orthogonal projectors in the two metrics.  Put
\[
d_\tau=\frac{\tau}{1-\tau},\qquad
g_\tau=\sqrt{\frac{1+\tau}{1-\tau}}.
\]
Then, as operators on \((\mathcal W,\omega)\),
\[
\|\widetilde P-P\|\le d_\tau,\qquad
\|\widetilde P\|\le g_\tau.
\tag{JW7}
\]
For the window band,
\[
\|\widetilde B_n-B_n\|\le2d_\tau,\qquad
\|\widetilde B_n\|,\ \|I-\widetilde B_n\|\le g_\tau.
\tag{JW8}
\]
\end{lemma}
\begin{proof}
For \(w\in\mathcal W\), put \(r=w-Pw\) and
\(h=\widetilde Pw-Pw\in\mathcal R\).  The original orthogonality gives
\(\langle r,h\rangle_\omega=0\); the new orthogonality gives
\(\langle r-h,h\rangle_{\widetilde\omega}=0\).
Consequently
\[
(1-\tau)\|h\|_\omega^2
\le \|h\|_{\widetilde\omega}^2
= (\widetilde\omega-\omega)(r,h)
\le\tau\|r\|_\omega\|h\|_\omega
\le\tau\|w\|_\omega\|h\|_\omega.
\]
This proves the difference bound.  For the second bound use
\(\|\widetilde Pw\|_{\widetilde\omega}\le
\|w\|_{\widetilde\omega}\) and (JW3).
Subtracting the two fixed nested relation projectors proves the first
part of (JW8).  Their difference and its complement are orthogonal
projectors in the new metric, proving its second part.
\end{proof}

\paragraph{Sections, multiplication, and both projected errors.}
The same argument in a fixed affine observation fibre proves
\[
\|\widetilde S_Nv-S_Nv\|_\omega
\le d_\tau\|S_Nv\|_\omega.
\tag{JW9}
\]
Put \(E_n=\|S_nv\|_\omega^2\).
Nested minima and orthogonality give
\[
\|S_{n+q}v\|_\omega\le\sqrt{E_n},\qquad
\|F_n\|_\omega^2=E_n-E_{n+q}\le E_n,\qquad
\|\widetilde F_n-F_n\|_\omega\le2d_\tau\sqrt{E_n}.
\tag{JW10}
\]
Multiplication by \(S-\lambda\) takes all these section differences
into \(\chi\Pcal_{D-q}\subseteq\mathcal W\), with degree at most \(D\).
Let \(M_a\) be a valid operator allowance for this multiplication on
the input face in the original metric.  The actual original
Gamma comparison C2 gives the following fully explicit choice:
\[
M_a=
\sqrt{\frac{b_a}{a_{a,\,2q+1}}}
\bigl(4q+2+|\lambda-c|\bigr).
\tag{JW11}
\]
In fact the incoming Gamma multiplication bound on degree \(2q\)
is \(\|yP\|_\sigma\le2(2q+1)\|P\|_\sigma\);
the triangle inequality for
\((S-\lambda)P=(iy-(\lambda-c))P\), the upper comparison for the output,
and the lower comparison for the input give (JW11).
Thus this bound uses only the displayed C2 provider.  It does not
silently replace \(b_a/a_{a,2q+1}\) by a fixed constant.
If the norms are written with the unscaled density \(m_a\) instead of
\(M_h(\alpha)^{-a}m_a\), both source/Gamma comparison constants acquire
the same factor \(M_h(\alpha)^a\).  Their ratio in (JW11) is unchanged;
all section energies in the formulas still carry their original mass.

Set \(w_n=(S-\lambda)F_n\) and
\(\widetilde w_n=(S-\lambda)\widetilde F_n\).  From (JW10)--(JW11),
\[
\|\widetilde w_n-w_n\|_\omega
\le2d_\tau M_a\sqrt{E_n}.
\tag{JW12}
\]
Equations (JW8) and (JW12), with the triangle inequality, now give
\[
\begin{split}
\|\widetilde B_n\widetilde w_n-B_nw_n\|_\omega
&\le2d_\tau(g_\tau+1)M_a\sqrt{E_n},\\
\|(I-\widetilde B_n)\widetilde w_n-(I-B_n)w_n\|_\omega
&\le2d_\tau(g_\tau+1)M_a\sqrt{E_n}.
\end{split}
\tag{JW13}
\]
Every vector and every projection in these inequalities is in
\(\mathcal W\).  No restricted pair-metric estimate is applied to
an arbitrary projector on the full cyclic quotient.

For a squared-norm version put
\(\varepsilon_\tau=2d_\tau(g_\tau+1)\).
For either projected vector \(V_n\) and its corresponding
\(\widetilde V_n\),
\[
\left|\|\widetilde V_n\|_{\widetilde\omega}^2
              -\|V_n\|_\omega^2\right|
\le
\left[\tau(1+\varepsilon_\tau)^2
      +2\varepsilon_\tau+\varepsilon_\tau^2\right]M_a^2E_n.
\tag{JW14}
\]
To see this, use \(\|V_n\|_\omega\le M_a\sqrt{E_n}\),
(JW13), and (JW3), then expand the difference of the two squared
norms.  This includes the cross terms.
Also, by comparing each of the two attained minima,
\[
\left|\|\widetilde F_n\|_{\widetilde\omega}^2
              -\|F_n\|_\omega^2\right|
\le\tau(E_n+E_{n+q})\le2\tau E_n.
\tag{JW15}
\]

\paragraph{Consequence for the inherited leading coefficients.}
The supplied original two-window result records
\(\|F_n\|_\omega^2/E_n\to1\) and the total coefficient
\(16/\pi^2\) after division by \(q^2\|F_n\|_\omega^2\);
it records that this coefficient is tangential at \(n=q-1\)
and normal at \(n=q\).  Equations (JW14)--(JW15) prove stability
of each of those exact existing coefficients under the specified
central replacement.  Indeed the displayed C2 constants give
\[
\log(b_a/a_{a,2q+1})=O_h(a+\log q),\qquad
\log\tau\le-2(q-1)\log(q/a)+O_h(q+a+\log q).
\tag{JW16}
\]
Since \(q=(a+1)^2\) and \(|\lambda-c|\le
a\sqrt{\delta^2+\gamma^2}\), it follows that
\(\tau M_a^2/q^2\to0\).  Hence the right side of (JW14),
divided by \(q^2E_n\), tends to zero for each window.
The denominator perturbation (JW15) also tends to zero relatively.
This proves the transfer without needing an unlocated sharp
\(O_h(q)\) multiplication theorem.

The bare rate \(O_h(\tau q\sqrt{E_n})\) stated in S16 should therefore
be replaced in the received source by the finite (JW13), unless its
separately named sharper multiplication provider is supplied.
Its asserted leading-coefficient stability is already justified by
the weaker explicit rate.  The numerical coefficient of the old
two-window theorem remains an inherited value: this proof does not
claim to have independently calculated that absent source.

\paragraph{Integration boundary.}
The results verified here neither calculate the seven absolute native
allocations nor the individual three-column projected arithmetic signs.
They provide a quantitative fixed-representative endpoint limit,
the exact diagonal norm discrepancy, uniform fixed-source moment input,
and complete transport of two specified boundary-window projections.
All original allocation and same-class response calculations remain in
their original metrics.


% END PRESERVED INPUT 06_JOINT_LIMIT_AND_WINDOW_PROOF.tex
\endgroup

\clearpage
\begingroup
\def\CC{\mathbb C}
\def\AA{\mathbb A}
\def\QQl{\overline{\mathbb Q}_\ell}
\def\End{\operatorname{End}}
\def\Sym{\operatorname{Sym}}
\def\im{\operatorname{im}}
\def\coker{\operatorname{coker}}
\def\rank{\operatorname{rank}}
\def\tr{\operatorname{tr}}
\def\diag{\operatorname{diag}}
\def\Res{\operatorname{Res}}

\setcounter{equation}{0}
\renewcommand{\theequation}{MII.\arabic{equation}}
\section*{Marked exponential image: exact connection, periods and metric receivers}
\addcontentsline{toc}{section}{Marked exponential image: exact connection, periods and metric receivers}
% BEGIN PRESERVED INPUT 19_MARKED_IMAGE_RECEIVING_PROOF.tex SHA256 bb3c6e2bd3f8cc86b3ce0b56e3072980edeff80710c3be8aec3b571a67cf018f



This note receives the complete nine-page \emph{Deligne marked image} calculation
\cite{mii:received}, including its finite-field endomorphism cohomology, its tensor
constituents, and its characteristic-zero marked frame. Every formula used below
is proved with the specified full polynomial and the original coefficient frame.
The additional receiving calculations are the exact second fundamental form of
the retained polynomial line, an explicit inverse to the regularized marked
frame, and the transport of every original restriction and attained quotient
through that frame. These are finite identities. Their statements contain no
unevaluated asymptotic hypothesis.

\section{Objects, coefficient models, and imported theorems}
Retain the stipulated simple quartet with
\[
 \rho=\tfrac12+\delta+i\gamma,\qquad 0<\delta<\tfrac12,\quad\gamma>2,
 \qquad k\equiv1\pmod4,
\]
and set
\begin{equation}\tag{MII1}\label{mii:mii:objects}
 \begin{gathered}
 c=k/2,\qquad q=(k+1)^2,\qquad
 \beta_{ab}=c+(2a-k)\delta+i(2b-k)\gamma,\\
 \chi(S)=\prod_{a,b=0}^k(S-\beta_{ab}),\qquad
 \Phi'(S)=\chi(S),\quad\Phi(0)=0,\qquad\lambda=k\rho,\\
 p_\lambda(S)=\frac{\chi(S)}{S-\lambda},\qquad
 b_\lambda=\frac1{\chi'(\lambda)},\qquad v_\lambda=b_\lambda p_\lambda.
 \end{gathered}
\end{equation}
All the roots are distinct: equality of two roots first equates their real
parts and then their imaginary parts. The numbers $\delta,\gamma$ are nonzero.
In particular $b_\lambda\ne0$. The arithmetic unit $v_h=2\xi/h$, its factor
$v_h(\rho)^k$ in the original inclusion, and the original convolution measure
remain unchanged. If that original inclusion specifies a scalar multiple
$a v_\lambda$, every occurrence of $b_\lambda$ in a frame determinant is replaced
by $a b_\lambda$; the present formulas use the cardinal polynomial in
\eqref{mii:mii:objects} and leave the inclusion itself in its original coordinates.

The marked complex is
\begin{equation}\tag{MII2}\label{mii:mii:complex}
 \left[\CC[u,t,S]\xrightarrow{\ u\partial_S+\chi(S)-t\ }
                 \CC[u,t,S]\,dS\right].
\end{equation}
The mark $t$, the period $u$, and the arithmetic cutoff $N$ are three separate
specified parameters. The maps connecting their uses will be given explicitly.

For the finite-field calculation choose a finitely generated coefficient ring
containing the finitely many coefficients of $\Phi$, invert $u$ and
$1,2,\ldots,q+1$, and choose a good closed fibre of characteristic $p>q+1$.
Use $\mathbb F_b$ for its finite ground field, $Q=b^n$ for an extension
cardinality, a nontrivial character $\psi:\mathbb F_p\to\QQl^*$, and $\ell\ne p$.
No reduction of the analytic function $\xi$ is used. The finite-field
polynomial and the complex polynomial arise from that coefficient model; an
identification of their cohomology fibres is not inserted.

We use the following results of Pierre Deligne with their exact scope
\cite[\S3.7, (3.7.2.3), Lemma (3.7.3); (3.4.1)(iii)]{mii:deligne}.
For a polynomial of degree $d$ prime to $p$ whose highest homogeneous part
defines a smooth projective hypersurface, exponential compact cohomology in
$n$ variables is concentrated in degree $n$, has dimension $(d-1)^n$, and is
pure of weight $n$. In the corresponding open coefficient family these
cohomology sheaves are lisse. A lisse pointwise-pure sheaf on a normal
finite-field scheme is geometrically semisimple. For the one-variable phases
here, the highest projective zero locus is empty in $\mathbb P^0$, and the
leading coefficient is invertible. Thus all mark values are included. These
are the cited external theorems; the endomorphism computations below are
proved here. Compact base change, K\"unneth and smooth Poincar\'e duality are
used with the conventions of \cite[\S1.4]{mii:deligne}.

\section{The entire endomorphism cohomology on the mark line}
Let
\begin{equation}\tag{MII3}\label{mii:mii:V}
 \mathcal V=R^1\pi_!\mathcal L_\psi((\Phi(x)-tx)/u),
 \qquad\pi:\AA_x^1\times\AA_t^1\longrightarrow\AA_t^1.
\end{equation}
It is lisse of rank $q$, pure of weight one, and the full compact direct image
is $\mathcal V[-1]$.

\begin{lemma}[The exact dual and Fourier kernels]\label{mii:mii:fourierlemma}
If $\mathcal V_-$ is defined by the negative of the complete phase, then
\begin{equation}\tag{MII4}\label{mii:mii:dual}
 \mathcal V_-\simeq\mathcal V^\vee(-1).
\end{equation}
For $\operatorname{pr}_a:\AA_a^1\times\AA_z^1\to\AA_a^1$ and the zero
inclusion $i_0$, one has
\begin{equation}\tag{MII5}\label{mii:mii:fourier}
 R\operatorname{pr}_{a!}\mathcal L_\psi(az)
       \simeq i_{0*}\QQl(-1)[-2].
\end{equation}
\end{lemma}
\begin{proof}
The degree $q+1$ of the phase is prime to $p$. A polar expression $g^p-g$
has highest pole order divisible by $p$, so the rank-one Artin--Schreier
character at infinity is nontrivial on wild inertia. Its wild-inertia
invariants vanish. Wild invariants are exact for $\ell$-adic coefficients
when $\ell\ne p$: one averages over the finite wild quotients, whose orders
are units. The tame-inertia spectral sequence consequently has zero input.
For the compactification $j:\AA^1\to\mathbb P^1$ this proves
$j_!\mathcal L\simeq Rj_*\mathcal L$. Poincar\'e duality pairs the two
opposite exponential degree-one groups perfectly into $\QQl(-1)$, proving
\eqref{mii:mii:dual}, including its twist.

For \eqref{mii:mii:fourier}, compact base change gives at $a=0$ the compact
complex of $\AA^1$, namely $\QQl(-1)[-2]$. If $a\ne0$, the cover
$y^p-y=az$ is an affine line with coordinate $y$. Its compact cohomology is
one Tate line in degree two, on which all deck translations act trivially.
The nontrivial character summand defining $\mathcal L_\psi(az)$ therefore
has zero compact cohomology. The relative complex is supported on the
closed zero point. Recollement identifies it with the direct image of its
zero-fibre complex, proving the stated relative identity.
\end{proof}

\begin{theorem}[Complete mark-line endomorphism complex]\label{mii:mii:irreducible}
Over the algebraic closure of the specified finite ground field,
\begin{equation}\tag{MII6}\label{mii:mii:endV}
 R\Gamma_c(\AA_t^1,\End\mathcal V)=\QQl(-1)[-2].
\end{equation}
The sheaf $\mathcal V$ is geometrically irreducible and
$H^1(\AA^1,\End\mathcal V)=0$.
\end{theorem}
\begin{proof}
Compute the compact complex of
$\mathcal L_\psi((\Phi(x)-\Phi(y)-t(x-y))/u)$ on the full ordered
space $\AA^3_{x,y,t}$. Integration first in $x,y$, the concentration in
degree one and \eqref{mii:mii:dual} give
\[
 R\Gamma_c(\AA_t^1,\End\mathcal V)(-1)[-2].
\]
Integration first in $t$ applies \eqref{mii:mii:fourier} to
$-(x-y)/u$. The result is supported exactly on $x=y$, with factor
$(-1)[-2]$. The remaining phase is identically zero on that diagonal, so
the same complex is
$R\Gamma_c(\AA_x^1,\QQl)(-1)[-2]$. Cancelling the identical twist and
shift proves \eqref{mii:mii:endV}. The trace pairing identifies
$(\End\mathcal V)^\vee$ with $\End\mathcal V$. Duality on the smooth
line now gives $H^0(\End\mathcal V)=\QQl$ and $H^1(\End\mathcal V)=0$.
Deligne's geometric semisimplicity applies to the pure lisse sheaf
$\mathcal V$. In a semisimple representation $\bigoplus_i W_i^{m_i}$
over the algebraically closed coefficient field, the dimension of the
endomorphism algebra is $\sum_i m_i^2$. Its value one forces exactly one
simple summand of multiplicity one. This proves irreducibility.
\end{proof}

The corresponding exact trace identity is
\begin{equation}\tag{MII7}\label{mii:mii:moment2}
 \sum_{t\in\mathbb F_Q}
 \left|\sum_{x\in\mathbb F_Q}
 \psi\operatorname{Tr}_{\mathbb F_Q/\mathbb F_p}
             \frac{\Phi(x)-tx}{u}\right|^2=Q^2.
\end{equation}
Indeed the $t$-sum of the expanded square is zero for $x\ne y$ and is
$Q$ for each of the $Q$ diagonal pairs. The complex identity
\eqref{mii:mii:endV} also determines the individual cohomology groups.

A morphism from $\mathcal V$ to any lisse target has a
monodromy-invariant kernel. Thus its rank is either zero or $q$. The
same statement holds after restricting to a nonempty open subcurve: its
fundamental group surjects onto that of the original smooth curve, because
every connected finite \`etale cover remains connected on restriction to
the dense open. This is a precise restriction on a proposed lisse
extension of a fibre map. The exact characteristic-zero coefficient
and metric receivers of actual fibre maps are constructed below; no
comparison of a complex fibre with an unspecified finite-field fibre
is required by them.

\section{The complete rank-four tensor-square calculation}
Retain the one-packet centred family and primitive constant
\begin{equation}\tag{MII8}\label{mii:mii:quintic}
 \begin{gathered}
 \phi_{\beta,\eta}(x)=x^5/5+\beta x^3/3+\eta x,\qquad
 A(\beta,\eta)=1/160+\beta/24+\eta/2,\\
 \beta=2(\gamma^2-\delta^2),\qquad
 \eta=(\delta^2+\gamma^2)^2
 \quad\hbox{at the original complex quartet}.
 \end{gathered}
\end{equation}
In characteristic $p>5$ with $u\ne0$ let
\[
 \mathcal U=R^1\pi_!\mathcal L_\psi(\phi_{\beta,\eta}(x)/u),\qquad
 \mathcal A=\mathcal L_\psi(A(\beta,\eta)/u),\qquad
 \mathcal V^{\rm orig}=\mathcal A\otimes\mathcal U.
\]
These rank-four sheaves on $\AA^2_{\beta,\eta}$ are lisse and pure of
weight one. The constant phase is retained in $\mathcal V^{\rm orig}$.

\begin{theorem}[The three tensor constituents]\label{mii:mii:tensor}
One has
\begin{equation}\tag{MII9}\label{mii:mii:endU}
 R\Gamma_c(\AA^2,\End\mathcal U)=\QQl(-2)[-4]
\end{equation}
and
\begin{equation}\tag{MII10}\label{mii:mii:endtensor}
 H_c^i(\AA^2,\End(\mathcal U^{\otimes2}))=
 \begin{cases}
 \QQl&i=2,\\
 \QQl(-1)^{\oplus3}&i=3,\\
 \QQl(-2)^{\oplus3}&i=4,\\
 0&\text{otherwise}.
 \end{cases}
\end{equation}
The following direct summands are geometrically simple, pairwise
nonisomorphic, and have ranks $10,1,5$, respectively:
\begin{equation}\tag{MII11}\label{mii:mii:decomposition}
 (\mathcal V^{\rm orig})^{\otimes2}
 =\Sym^2\mathcal V^{\rm orig}
 \oplus\mathcal A^2(-1)
 \oplus(\bigwedge^2\mathcal V^{\rm orig})_0.
\end{equation}
The last summand is the kernel of the alternating contraction to
$\mathcal A^2(-1)$.
\end{theorem}
\begin{proof}
For \eqref{mii:mii:endU}, use ordered $x,y$. Integration in $\eta$ of the
full opposite-phase product imposes $x=y$ by \eqref{mii:mii:fourier}.
The cubic and quintic differences then vanish. The remaining free
coordinates are $x,\beta$. The twist $(-1)$ and shift $[-2]$ from this
Fourier integration equal those from degree-one clean duality in $x,y$.
Cancelling them leaves exactly $R\Gamma_c(\AA^2,\QQl)$.

For the tensor square use ordered $x_1,x_2,y_1,y_2$. Integrating those
four fibre variables first gives
\[
 R\Gamma_c(\AA^2,\End(\mathcal U^{\otimes2}))(-2)[-4].
\]
Instead integrate $\eta,\beta$ first. Their Fourier constraints are
$x_1+x_2=y_1+y_2$ and $x_1^3+x_2^3=y_1^3+y_2^3$, with the identical
twist and shift $(-2)[-4]$. Write $x_1=x,x_2=s-x,y_1=y,y_2=s-y$.
The exact polynomial identities are
\begin{equation}\tag{MII12}\label{mii:mii:phasefactor}
 \begin{aligned}
 x^3+(s-x)^3-y^3-(s-y)^3
     &=3s(x-y)(x+y-s),\\
 x^5+(s-x)^5-y^5-(s-y)^5
     &=5s(x-y)(x+y-s)\bigl(x^2-sx+y^2-sy+s^2\bigr).
 \end{aligned}
\end{equation}
Their signs and all the opposite-pair solutions are retained. The
second polynomial vanishes on the exact reduced zero scheme of the
first. The linear map
$(s,x,y)\mapsto(s,x-y,s-x-y)$ has determinant $-2$, so this scheme
is $Z=\{z_1z_2z_3=0\}\subset\AA^3$. Consequently
\[
 R\Gamma_c(\AA^2,\End(\mathcal U^{\otimes2}))
       =R\Gamma_c(Z,\QQl).
\]
The complement is $(\mathbb G_m)^3$. Its compact cohomology in degrees
$3,4,5,6$ is respectively
$\QQl,\QQl(-1)^3,\QQl(-2)^3,\QQl(-3)$, by three K\"unneth products
of $H_c^1(\mathbb G_m)=\QQl$ and $H_c^2(\mathbb G_m)=\QQl(-1)$.
The open--closed exact sequence in $\AA^3$ shifts the first three
groups down by one to the groups of $Z$. The top map
$H_c^6((\mathbb G_m)^3)\to H_c^6(\AA^3)$ is an isomorphism of
fundamental classes. This proves \eqref{mii:mii:endtensor}, including its
Frobenius twists and its zero groups.

Duality on the smooth two-dimensional base gives dimensions one and
three for $H^0(\End\mathcal U)$ and
$H^0(\End(\mathcal U^{\otimes2}))$. To identify three actual
summands, the odd involution $\iota(x)=-x$ identifies the negative
phase with the positive phase. Clean degree-one duality gives
\[
 B(a,b)=\int a\smile\iota^*b,
 \qquad B:\mathcal U\otimes\mathcal U\longrightarrow\QQl(-1).
\]
It is nondegenerate. Pullback by $\iota$ preserves the top trace class
and $\iota^2=1$; exchanging the two degree-one factors changes the
sign. Therefore $B(b,a)=-B(a,b)$. In a symplectic basis
$e_1,f_1,e_2,f_2$, the invariant bivector
$e_1\wedge f_1+e_2\wedge f_2$ contracts to $2$, so it splits a
Tate line from $\bigwedge^2\mathcal U$. Its contraction kernel has
rank five and $\Sym^2\mathcal U$ has rank ten. Twisting both factors
by the original $\mathcal A$ gives \eqref{mii:mii:decomposition}.
Deligne's semisimplicity applies to the pure lisse tensor square.
Its three nonzero direct summands contribute at least three to its
endomorphism dimension, already proved to be exactly three. Hence
each summand is simple and no two are isomorphic. Their distinct
ranks give the same last conclusion directly.
\end{proof}

For every extension field the full trace sum $S(\beta,\eta)$ satisfies
\begin{equation}\tag{MII13}\label{mii:mii:moment4}
 \sum_{\beta,\eta\in\mathbb F_Q}|S(\beta,\eta)|^4
       =Q^2(3Q^2-3Q+1).
\end{equation}
The constant character $\mathcal A$ cancels in this fourth absolute
power. The two parameter sums give $Q^2$ times the number of points
of $Z$. At $s=0$ there are $Q^2$ pairs $(x,y)$; at each $s\ne0$
there are $2Q-1$. The exact number is
$Q^2+(Q-1)(2Q-1)=3Q^2-3Q+1$, proving the formula.

For completeness the two finite fibre maps mentioned in the received
continuation have elementary exact ranks. For
$D=\diag(\zeta,\zeta^2,\zeta^3,\zeta^4)$ with $\zeta^5=1$ primitive,
the actual cyclic average on the symmetric square is
\begin{equation}\tag{MII14}\label{mii:mii:cyclic}
 P_5=\tfrac15\sum_{j=0}^4\Sym^2(D^j),\qquad
 \im P_5=\langle x_1x_4,x_2x_3\rangle.
\end{equation}
The characters on $x_rx_s$ are $\zeta^{r+s}$, and precisely the
two listed pairs have exponent zero modulo five. Thus the rank is
two. For the original four distinct one-packet roots $\alpha_i$,
let $e_i$ be their cardinal idempotents and let $A e_i=\alpha_i e_i$.
The nine distinct sums $\alpha_i+\alpha_j$ form the $3$ by $3$
grid with increments $2\delta,2i\gamma$. Their sum algebra maps to
the symmetric square by the exact evaluation map
\begin{equation}\tag{MII15}\label{mii:mii:sumembedding}
 \begin{aligned}
 \iota_2:\CC[T]/\prod_{\nu}(T-\nu)&\longrightarrow\Sym^2\CC^4,\\
 f&\longmapsto f(A\otimes1+1\otimes A)(1\otimes1),\\
 \iota_2(e_\nu)&=
 \sum_{2\alpha_i=\nu}e_i\otimes e_i
 +2\sum_{i<j,\ \alpha_i+\alpha_j=\nu}e_i\odot e_j,
 \end{aligned}
\end{equation}
where $e_i\odot e_j=(e_i\otimes e_j+e_j\otimes e_i)/2$.
The nine displayed images are nonzero and have disjoint supports in
the symmetric tensor basis. This proves injectivity and rank nine,
including the multiplicity two of an off-diagonal symmetric pair.
Neither this image nor \eqref{mii:mii:cyclic} can be the image of a
proper nonzero lisse projector on the simple rank-ten summand of
\eqref{mii:mii:decomposition}. These exact finite maps remain available
as fibre maps; the proved assertion concerns their extension as
projectors of this particular finite-field family.

\section{The marked coefficient module and the complete derivative frame}
Return to the complex polynomial \eqref{mii:mii:objects}, put
$R=\CC[u,u^{-1},t]$, and write
\[
 \mathscr H=R[S]/(u\partial_S+\chi-t)(R[S]).
\]
This is a quotient of $R$-modules. Polynomial division in the
differential expression gives a unique representative of degree less
than $q$: subtract the leading coefficient times
$(u\partial_S+\chi-t)S^{d-q}$ from a polynomial of degree $d\ge q$,
and repeat. Termination follows from the strictly decreasing degree.
Uniqueness follows because every nonzero polynomial $f$ has
$\deg((u\partial_S+\chi-t)f)=q+\deg f\ge q$.
Thus $1,S,\ldots,S^{q-1}$ is a free $R$-basis.

The cochain commutator and its induced connection are
\begin{equation}\tag{MII16}\label{mii:mii:connection}
 [\partial_t-S/u,\ u\partial_S+\chi-t]=0,
 \qquad\nabla_t=\partial_t-A(t)/u.
\end{equation}
Here $A(t)$ is the companion matrix of $\chi-t$ on these
representatives: multiplying a polynomial of degree at most $q-1$
by $S$ requires only the relation
$(u\partial_S+\chi-t)1=\chi-t$. This proves the displayed matrix
connection. It does not define a ring multiplication on the
differential quotient. At $t=0$ the representative map identifies
its underlying vector space with $E=\CC[S]/\chi$ exactly.

Let $R_j$ denote the representative of $[S^j p_\lambda]$. The
polynomial relation $(S-\lambda)p_\lambda=\chi$ and reduction of
$(\chi-t)S^{j-1}$ give
\begin{equation}\tag{MII17}\label{mii:mii:recurrence}
 R_j-\lambda R_{j-1}
       =t S^{j-1}-u(j-1)S^{j-2},\qquad 1\le j\le q.
\end{equation}
The last term is zero at $j=1$. Its right side already has degree
less than $q$, so this is equality of the literal representatives.
Applying the cochain connection to the unreduced fixed polynomial
before reduction gives, for every $j\ge0$,
\begin{equation}\tag{MII18}\label{mii:mii:jets}
 \nabla_t^j[v_\lambda]=b_\lambda(-u)^{-j}R_j.
\end{equation}
The commutator in \eqref{mii:mii:connection} proves that differentiating
the reduced representative gives exactly the same class.

\begin{theorem}[Determinants and the entire supported cokernel]
In the increasing original $S$-coefficient basis, let
\[
 \mathcal O_\lambda(t)=[v_\lambda,\nabla_t v_\lambda,
                  \ldots,\nabla_t^{q-1}v_\lambda],
 \qquad\epsilon_q=(-1)^{(q-1)(q+2)/2}.
\]
Then
\begin{equation}\tag{MII19}\label{mii:mii:Odet}
 \det\mathcal O_\lambda(t)
   =\epsilon_q b_\lambda^q u^{-q(q-1)/2}t^{q-1},
 \qquad\coker\mathcal O_\lambda\simeq R/(t^{q-1}).
\end{equation}
The fibre matrix has rank $q-1$ at zero. The regularized original
frame
$\mathcal J_\lambda=[v_\lambda,\nabla_t^2v_\lambda,
\ldots,\nabla_t^qv_\lambda]_{t=0}$ is invertible with
\begin{equation}\tag{MII20}\label{mii:mii:Jdet}
 \det\mathcal J_\lambda
    =\epsilon_q(q-1)!b_\lambda^q u^{-q(q-1)/2}.
\end{equation}
\end{theorem}
\begin{proof}
First remove the factors $b_\lambda(-u)^{-j}$ from the columns
of $\mathcal O_\lambda$. Perform the unipotent column operation
$R_j\mapsto R_j-\lambda R_{j-1}$ simultaneously, or from right
to left. Equation \eqref{mii:mii:recurrence} gives the columns
$p_\lambda$ and $tS^{j-1}-u(j-1)S^{j-2}$, $1\le j\le q-1$.
The first is monic of degree $q-1$; the other columns have degrees
at most $q-2$. Moving the first column to the last position gives
the sign $(-1)^{q-1}$. The remaining upper-bidiagonal block has
diagonal entries $t$. Restoring the column factors gives sign
$(-1)^{q-1+q(q-1)/2}=\epsilon_q$ and proves the determinant.

The first column eliminates the top coefficient generator. On the
remaining generators $e_0,\ldots,e_{q-2}$ the presentation is exactly
\[
 te_0=0,\qquad te_j=u j e_{j-1}\quad(1\le j\le q-2).
\]
The explicit presentation isomorphism is
\begin{equation}\tag{MII21}\label{mii:mii:cokeriso}
 e_j\longmapsto
 \frac{j!}{(q-2)!}u^{-(q-2-j)}t^{q-2-j}\pmod{t^{q-1}},
 \qquad\text{inverse: }1\longmapsto e_{q-2}.
\end{equation}
Every relation maps to zero. Conversely the displayed relations,
using the invertible constants $uj$, express every generator in
terms of $e_{q-2}$ by precisely this formula; its sole remaining
relation is $t^{q-1}e_{q-2}=0$. This proves the isomorphism and
shows that the special fibre of the cokernel has dimension one.

For the regularized frame remove the column scalars with orders
$0,2,\ldots,q$. At zero $R_1=\lambda p_\lambda$; replace
$R_2$ by $R_2-\lambda^2p_\lambda$, and, from right to left, each
$R_j$ for $j\ge3$ by $R_j-\lambda R_{j-1}$. The columns become
$p_\lambda,-u,-2uS,\ldots,-(q-1)uS^{q-2}$. Their determinant
is $u^{q-1}(q-1)!$. The restored exponent of $-u$ is
$-[q(q+1)/2-1]$, giving exactly \eqref{mii:mii:Jdet}.
\end{proof}

The proof of this finite algebraic theorem also works for a
terminal primary class when
$\chi=(S-\lambda)^e\chi_\lambda$ and $\chi_\lambda(\lambda)\ne0$:
use $p_\lambda=\chi/(S-\lambda)$ and
$b_\lambda=1/\chi_\lambda(\lambda)$. This statement follows from
the same polynomial recurrence, which used no simple-root
division. It does not change the simple-quartet metric in this note.

\section{The explicit inverse and the retained-line leakage}
Put $v_j=(\nabla_t^j v_\lambda)|_{t=0}$. Formula
\eqref{mii:mii:recurrence} gives the following inverse reconstruction
in the original coefficient space:
\begin{equation}\tag{MII22}\label{mii:mii:inverse}
 \begin{aligned}
 1&=-\frac{u}{b_\lambda}v_2
            +\frac{\lambda^2}{b_\lambda u}v_0,\\
 S^r&=\frac{(-u)^{r+1}v_{r+2}
                    -\lambda(-u)^r v_{r+1}}{b_\lambda(r+1)}
                &&(1\le r\le q-2),\\
 S^{q-1}&=b_\lambda^{-1}v_0-
                         \sum_{r=0}^{q-2}p_rS^r,
 \qquad p_\lambda=S^{q-1}+\sum_{r=0}^{q-2}p_rS^r.
 \end{aligned}
\end{equation}
To verify the first identity, set $j=2,t=0$ in the recurrence to
obtain $R_2=\lambda^2p_\lambda-u$. For the middle identities use
$R_{r+2}-\lambda R_{r+1}=-u(r+1)S^r$ and then
$R_j=b_\lambda^{-1}(-u)^jv_j$. The last identity is the monic
definition of $p_\lambda$. These identities express each original
basis vector in the $q$ specified jet columns and are an explicit
inverse matrix for $\mathcal J_\lambda$; no basis existence argument
replaces the map.

The exact bundle map measuring the retained polynomial line is
also explicit. Let $L=R v_\lambda\subset\mathscr H$, and let
$\pi_L:\mathscr H\to\mathscr H/L$ be its quotient map. Since
$v_\lambda$ is monic up to the unit $b_\lambda$, $L$ is a direct
summand as an $R$-module. The second fundamental form of this
specified line is the $R$-linear map
\begin{equation}\tag{MII23}\label{mii:mii:secondfund}
 \begin{aligned}
 \mathfrak b_L:L&\longrightarrow(\mathscr H/L)\,dt,\\
 f v_\lambda&\longmapsto
       \pi_L(\nabla_t(fv_\lambda))\,dt
       =-\frac{b_\lambda t}{u}f\,\pi_L(1)\,dt.
 \end{aligned}
\end{equation}
The Leibniz term $f'v_\lambda$ maps to zero. The formula itself
is \eqref{mii:mii:recurrence} at $j=1$. The class $\pi_L(1)$ is
nonzero because $q\ge2$ and a degree-zero polynomial cannot be
an $R$-multiple of the nonconstant $p_\lambda$. In particular,
the second fundamental form vanishes to exactly first order
at the original mark, and
\begin{equation}\tag{MII24}\label{mii:mii:secondjet}
 \pi_L(v_1)=0,\qquad
 \pi_L(v_2)=-\frac{b_\lambda}{u}\pi_L(1)\ne0.
\end{equation}
This computes the first actual marked leakage from the original
arithmetic eigenline $A(0)v_\lambda=\lambda v_\lambda$.
Any connection-stable subbundle containing the specified
polynomial section contains every $v_j$ and hence, by
\eqref{mii:mii:inverse}, the entire fibre at zero.

There is also an exact formal bridge to a horizontal extension
of its initial vector. Write
$A(t)=A_0+tB$, $B=e_0e_{q-1}^{\mathsf T}$ in the original
increasing coefficient frame. The formal matrix
$H(t)=\sum_{n\ge0}H_nt^n$ is uniquely defined by
\begin{equation}\tag{MII25}\label{mii:mii:horizontal}
 H_0=I,\quad H_{-1}=0,\qquad
 (n+1)H_{n+1}=u^{-1}(A_0H_n+B H_{n-1}).
\end{equation}
It has invertible constant coefficient and satisfies
$\nabla_t H=0$. Thus $a(t)\mapsto H(t)a(t)$ is a
connection-preserving isomorphism from the trivial formal
connection to the completed marked connection. Its line through
$v_\lambda$ is the horizontal extension of that fibre vector.
The polynomial line's actual leakage remains
\eqref{mii:mii:secondfund}; these two sections are related by the
explicit matrix \eqref{mii:mii:horizontal}, without identifying them.

\section{Period determinant, Gamma constants, and scalar equation}
The full root polynomial satisfies $\chi(c+w)=\chi(c-w)$.
Its degree $q$ is even. Write
\[
 \Phi(c+w)=C_\Phi+\phi(w),\quad C_\Phi=\Phi(c),\quad
 \phi(0)=0,\quad\phi\text{ odd}.
\]
Choose the original branch of $u^{1/(q+1)}$ and retain the ordered
decay rays with angles $(\arg u+\pi+2\pi j)/(q+1)$, $0\le j\le q$.
The cycles are the original differences $\Gamma_j=-\ell_0+\ell_j$,
$1\le j\le q$. Moving their common finite starting point to
$w=0$ changes no integral: the same finite connecting path cancels
between the two rays, and the joining arcs at infinity decay in
the same sectors. Define the period matrix by
\[
 \Pi_{ja}(u,t)=\int_{\Gamma_j}w^{a-1}
    \exp\bigl((C_\Phi-ct+\phi(w)-tw)/u\bigr)\,dw,
 \qquad1\le j,a\le q.
\]
Let
\begin{equation}\tag{MII26}\label{mii:mii:Fg}
 \zeta=e^{2\pi i/(q+1)},\quad F_{ja}=\zeta^{ja}-1,\qquad
 g_a=(q+1)^{a/(q+1)-1}e^{\pi ia/(q+1)}
             \Gamma\!\left(\frac a{q+1}\right).
\end{equation}

\begin{theorem}[Complete original period determinant]\label{mii:mii:periodtheorem}
With the just specified branches, orientations and column order,
\begin{equation}\tag{MII27}\label{mii:mii:perioddet}
 \det\Pi(u,t)=
 e^{q(C_\Phi-ct)/u}(\det F)
                 \left(\prod_{a=1}^qg_a\right)u^{q/2}.
\end{equation}
\end{theorem}
\begin{proof}
Separate the constant $C_\Phi-ct$ in the exponent. Vary a lower
odd coefficient of the centred phase, with inserted odd power
$w^p$, $p\le q-1$. For the column $w^j$, $0\le j\le q-1$,
perform ordinary division
$w^{p+j}=(\chi(c+w)-t)Q(w)+r(w)$. If the quotient is nonzero
its degree is at most $p+j-q\le j-1$. In the differential
quotient, the inserted column is represented by
$r-uQ'$. This last term has degree at most $j-2$ and needs no
further reduction. The ordinary remainder map is multiplication
by $w^p$ in the quotient by the even polynomial $\chi(c+w)-t$;
it reverses parity and so has zero trace. The correction
$-uQ'$ is strictly lower triangular and also has zero trace.
Therefore the coefficient derivative of the period matrix has
the form $u^{-1}\Pi M_p$ with $\tr M_p=0$. Multilinearity, or
the adjugate formula, gives zero derivative of its determinant
without presuming its invertibility. This proves independence
of every lower odd coefficient, including the marked coefficient
$-t$ in the centred phase.

All these differentiations are dominated along the fixed rays
on compact coefficient sets by
$\exp[-R^{q+1}/((q+1)|u|)+O(R^{q-1})]$ times a polynomial.
At the monomial phase, use
$w=u^{1/(q+1)}e^{\pi i/(q+1)}\zeta^j r$ with $r>0$ on the
$j$th ray, followed by $s=r^{q+1}/(q+1)$. The literal integral
is $\zeta^{ja}u^{a/(q+1)}g_a$. Taking the specified differences
gives $F\diag(u^{a/(q+1)}g_a)$. Finally
$\sum_{a=1}^q a/(q+1)=q/2$, and restoring the separated
constant proves \eqref{mii:mii:perioddet}.
\end{proof}

Let $\mathbf y_\lambda(t)$ be the vector of integrals of the
original amplitude $v_\lambda(S)$. Translation of coefficient
vectors from $S$ to $w$ is the exact matrix
\[
 T_{rj}=\binom jr c^{j-r}\quad(0\le r\le j\le q-1),\qquad
 T_{rj}=0\quad(r>j),\qquad\det T=1.
\]
Differentiating integrals gives the cochain connection in
\eqref{mii:mii:connection}. Hence the derivative matrix is literally
$\Pi T\mathcal O_\lambda$, and
\begin{equation}\tag{MII28}\label{mii:mii:wronskian}
 \det[\mathbf y_\lambda,\mathbf y'_\lambda,\ldots,
                       \mathbf y^{(q-1)}_\lambda]
 =\epsilon_q b_\lambda^q u^{-q(q-1)/2}t^{q-1}\det\Pi(u,t).
\end{equation}
The regularized determinant at the original mark is
\begin{equation}\tag{MII29}\label{mii:mii:regperiod}
 \left|\det[\mathbf y_\lambda,\mathbf y''_\lambda,\ldots,
                    \mathbf y_\lambda^{(q)}]_{t=0}\right|^2
 =(2\pi)^q((q-1)!)^2|b_\lambda|^{2q}
     |u|^{-q(q-2)}e^{2q\Re(C_\Phi/u)}.
\end{equation}
To check every constant, summing characters gives
$F^*F=(q+1)(I+\mathbf1\mathbf1^*)$, whence
$|\det F|^2=(q+1)^{q+1}$. Gauss's multiplication formula
\cite[(5.5.6)--(5.5.7)]{mii:dlmf} gives
$\prod_{a=1}^q\Gamma(a/(q+1))=(2\pi)^{q/2}/\sqrt{q+1}$.
Consequently $|\det F\prod g_a|^2=(2\pi)^q$.
Combining this with \eqref{mii:mii:Jdet} and
\eqref{mii:mii:perioddet} proves \eqref{mii:mii:regperiod}.

For any one original cycle let $Y$ be its period with amplitude
$v_\lambda$, let $F_0$ be its period with amplitude one, and
put $D=-u\partial_t$. Integration by parts in $S$ gives
$(D-\lambda)Y=b_\lambda tF_0$ and
$Y=b_\lambda p_\lambda(D)F_0$. For $t\ne0$,
$tDt^{-1}=D+u/t$, so
\begin{equation}\tag{MII30}\label{mii:mii:scalar}
 [p_\lambda(D+u/t)(D-\lambda)-t]Y=0.
\end{equation}
The operator factors have this order. The leading two coefficients
are $D^q+[\chi_{q-1}+(q-1)u/t]D^{q-1}$, where
$\chi_{q-1}=-qc$. Abel's determinant equation therefore gives
$W'/W=(q-1)/t-qc/u$, as in \eqref{mii:mii:wronskian}.
Every component period is holomorphic at zero. Its apparent
singularity is represented by the explicit cyclic cokernel
\eqref{mii:mii:cokeriso} and the regularized frame, without discarding
its multiplicity $q-1$.

\section{Exact receiving maps for the original arithmetic metrics}
Fix $u\ne0$ in the original branch. Let $G_N$ be the unchanged
positive Hermitian metric on the original $q$-dimensional space
$E=\CC[S]/\chi$, using its actual masses, cutoffs, root factors,
and full arithmetic unit. Set $J=\mathcal J_\lambda$ with its
explicit inverse \eqref{mii:mii:inverse}. The diagram of vector
spaces and metrics is the literal pullback
\begin{equation}\tag{MII31}\label{mii:mii:gram}
 J:(\CC^q,J^*G_NJ)\longrightarrow(E,G_N),\qquad
 \det(J^*G_NJ)=((q-1)!)^2|b_\lambda|^{2q}
                       |u|^{-q(q-1)}\det G_N.
\end{equation}
The map $J$ sends standard vectors to the original specified
jet columns; it introduces no source normalization. The determinant
identity follows by multiplying determinants and using
\eqref{mii:mii:Jdet}. For the actual positive extremal root the
normalizing constant itself is completely specified:
\begin{equation}\tag{MII32}\label{mii:mii:rootproduct}
 |b_\lambda|^{-1}
  =\prod_{\substack{0\le i,j\le k\\(i,j)\ne(0,0)}}
                   2\sqrt{\delta^2 i^2+\gamma^2j^2}.
\end{equation}
This is the product of all differences
$\lambda-\beta_{ab}=2(k-a)\delta+2i(k-b)\gamma$.

The map also transports every specified arithmetic subquotient,
not only its total determinant. Let $I_X:\CC^r\hookrightarrow E$
be any actual injective column map and $Q_X:E\twoheadrightarrow
\CC^s$ any actual surjection. Define their jet-coordinate maps
$\widetilde I_X=J^{-1}I_X$ and $\widetilde Q_X=Q_XJ$.
Then
\begin{equation}\tag{MII33}\label{mii:mii:restrictions}
 \begin{aligned}
 \widetilde I_X^*(J^*G_NJ)\widetilde I_X&=I_X^*G_NI_X,\\
 \widetilde Q_X(J^*G_NJ)^{-1}\widetilde Q_X^*
                         &=Q_XG_N^{-1}Q_X^*,\\
 \left[\widetilde Q_X(J^*G_NJ)^{-1}\widetilde Q_X^*\right]^{-1}
                         &=(Q_XG_N^{-1}Q_X^*)^{-1}.
 \end{aligned}
\end{equation}
The first equality cancels $J,J^{-1}$ directly. For the second
use $(J^*G_NJ)^{-1}=J^{-1}G_N^{-1}J^{-*}$; the third follows
because surjectivity and positivity make both matrices invertible.
Equivalently the change $x=Jz$ bijects the exact affine fibres
$\{x:Q_Xx=y\}$ and $\{z:\widetilde Q_Xz=y\}$ and preserves their
quadratic values. Their attained minima agree. Restricting first
and taking a further quotient applies these same equalities to
$I_X^*G_NI_X$ and the unchanged next quotient map. Thus every
finite composition of the programme's actual restrictions and
attained quotients has this exact receiver. No rank-proportional
allocation has been substituted for any of these matrices.

Retain the original four-endpoint signs by defining
\[
 \mathcal R_{\rm ar}f=f(q-1)+f(q)-f(2q-1)-f(2q).
\]
The multiplier in \eqref{mii:mii:gram} is independent of $N$;
$1+1-1-1=0$. Therefore
\begin{equation}\tag{MII34}\label{mii:mii:return}
 \mathcal R_{\rm ar}\log\det(J^*G_NJ)
    =\mathcal R_{\rm ar}\log\det G_N
    =\mathcal B_k^{\rm ar}.
\end{equation}
For a restriction or an attained quotient transported by
\eqref{mii:mii:restrictions}, even its full matrix is identical in its
original domain or codomain, so its signed return agrees term by
term. This supplies the explicit morphisms relating the geometric
marked frame to the arithmetic receivers. It does not evaluate
the still-present matrices $I_X^*G_NI_X$ or $Q_XG_N^{-1}Q_X^*$.
In particular the seven original allocations and the three-column
current responses retain those exact matrices. Their calculation
continues from these equalities with the original inverse and
complete action retained.

\section{Acceptance scope and provenance}
The admitted conclusions are the full formulas MII6--MII13 for
the specified good finite-field families, MII16--MII30 for the
original complex marked family, and the exact original-metric
receivers MII31--MII34. MII22--MII25 and MII33 explicitly extend
the received calculation. They retain the original full root
polynomial, mark, period branch, and all coefficient factors.
The degree-two ranks in MII14--MII15 were proved here independently.

The received text also quotes a large-$k$ five-orbit observation
rank $q-(8k-16)$. Its precise admission domain and its earlier
proof are not an input to any theorem in this receiving note;
the formula is not newly certified here. The general obstruction
to an intermediate-rank lisse morphism is proved above without it.
No assertion here identifies an arithmetic Frobenius, makes a
cross-characteristic fibre isomorphism, or evaluates an absolute
native angular allocation. In particular none asserts that the
broader tau-base or $\mathbb F_1$ programme has failed.

The entire received PDF and transcript were read. The local full
D032 source was consulted specifically at (3.4.1)(iii) and
(3.7.2.3) with its following lissity argument. The primary Numdam
record and the DLMF multiplication formula were checked on
19 September 2026. The source hashes, exact scope, and independent
algebraic checks are recorded in the accompanying machine-readable
acceptance file. Finite rational fixtures are used only to detect
algebraic errors; they do not represent actual zeros of $\xi$.



% END PRESERVED INPUT 19_MARKED_IMAGE_RECEIVING_PROOF.tex
\endgroup

\clearpage
\begingroup


\setcounter{equation}{0}
\renewcommand{\theequation}{PEQR.\arabic{equation}}
\section*{Primitive-choice receiving identities}
\addcontentsline{toc}{section}{Primitive-choice receiving identities}
% BEGIN PRESERVED INPUT 07_RECEIVER_UPDATES_INTEGRATED.tex SHA256 2fe76f4de4a834f095ae7264afc3fd101ee308bfd412e36071efe93db4cf88df
% Exact additive PEQ/PEI receiver; no predecessor is overwritten.
\paragraph{Primitive-choice quadratic control in the original target order.}
Use the unchanged PRT/CPQ coefficient maps at target degree $a=k+4$,
target size $q=(a+1)^2$, lower-grid size $(a-7)^2$, cutoff $N-v$,
and orders $s=1,a$. Retain the actual rank
$r_\partial=\dim(T_0/T_1)\in\{0,1,2,3\}$.
With the literal constants of PEQ and PEI7--11 set
\[
 B^{\rm end}_{a,s}=
 \min\{\log A^{\rm end}_{a,s},E_{a,s}(2q-1)\}
 +\min\{\log A^{\rm end}_{a,s},E_{a,s}(2q)\}.
\]
Then
\[
 \begin{gathered}
 0\le U_A\le3B^{\rm end}_{a,s},\quad
 0\le U_B\le(3-r_\partial)B^{\rm end}_{a,s},\quad
 0\le U_C\le r_\partial B^{\rm end}_{a,s},\\
 \mathcal L_0-\mathcal L_1=Z_1-Z_0
     =\mathcal P^{\rm rest}_0-\mathcal P^{\rm rest}_1
     =U_C-U_A+U_B,\\
 |\mathcal L_0-\mathcal L_1|
 \le3B^{\rm end}_{a,s}\le6\mathfrak E^{\rm best}_{a,s}
       =O_{\rm actual}(q).
 \end{gathered}
\]
PEI16 gives all directed component intervals. The primitive-choice
difference divided by $q\log a$ tends to zero. The exact allocation
identities remain
$4S_K=\mathcal L_0+Z_0=\mathcal L_1+Z_1$ and
$S_R=\mathcal T-S_K$. The prefix contribution remains
$-V_A^{\rm prefix}$. An absolute rank-$t$ quotient obeys the bound
with its explicit factor $t$; when $t=O(a)$ the bound permits
$O(aq)$. This update therefore supplies no absolute allocation
coefficient. The independent degree-$k$ proper source and each
response factor in the projected-sign calculation retain their
original forms and unresolved values.

% END PRESERVED INPUT 07_RECEIVER_UPDATES_INTEGRATED.tex
\endgroup

\clearpage
\begingroup


\setcounter{equation}{0}
\renewcommand{\theequation}{MRI.\arabic{equation}}
\section*{Backward propagation into the complete action}
\addcontentsline{toc}{section}{Backward propagation into the complete action}
% BEGIN PRESERVED INPUT 08_RECEIVER_UPDATES.tex SHA256 4f2c98d651691872d95aa6cf0e73fd8c6d43a7c018d084fb741bc963c2b4376d
% Current additive receiver. Original scalar, orders and errors are retained.
\section{Backward propagation of the September 19 calculations}
\label{mri:sec:multi-intake-receiver}

This section updates the live status of the earlier SXR10--12, GTM20,
DRU11--12 and primitive-choice receivers. Their earlier proofs and
versions remain in the cumulative source. All claims here refer to the
original stipulated quartet, its complete root polynomial, the physical
coordinate and the attained minima specified in the receiving proofs.
In particular, the dimension is not substituted for the tensor order.

\subsection{The arithmetic scalar in the complete action}
For the original degree $a\equiv1\pmod4$, fixed multiplicity $m_0\ge1$,
and dimension $q_a=[1+a(m_0-1)](a+1)^2$, the new scalar theorem gives
\[
 \delta_{a,1}^{\rm ar}
 =-(4m_0-2)\mathcal C_\Gamma q_a
       +\frac{\log2}{\pi}I_h a^2+E_{h,a},
 \tag{MRI1}
\]
where $I_h=\int_0^{\pi/2}\mu_h(\theta)^2\,d\theta$ uses the original
one-factor tilted mean and
$\mathcal C_\Gamma=2\mathcal L(2,\pi)-8\log2+4$.
The full definitions, convergent evaluation, original finite guards,
and proofs of these constants occur in M1--M8, I1--I4, OR1--41 and
SAI1--12 in this edition. The quantitative error is
\[
 |E_{h,a}|\le C_hq_a
 \left\{\frac{\log\log(q_a+16)}{\log(q_a+2)}
       +q_a^{-1/200}\log(q_a+2)^2\right\}=o_h(q_a).
 \tag{MRI2}
\]
Here $C_h$ is the fixed-source constant from that proof, not a
numerically evaluated uniform constant over all quartets.

The arithmetic receiver is the exact existing equality
\[
 J_a=\mathcal B_a^{(s)}+\delta_{a,s}^{\rm ar}
               +\mathcal W_a^{\rm mono}-\mathcal P_{a,-}-\mathcal P_{a,+}.
 \tag{MRI3}
\]
The source-order map, including both original masses, is
\[
 \delta_{a,a}^{\rm ar}-\delta_{a,1}^{\rm ar}
     =\mathcal B_a^{(1)}-\mathcal B_a^{(a)}.
 \tag{MRI4}
\]
Indeed both sides follow by subtracting the two expressions MRI3 for
the same action and unchanged remaining terms. Substitution of MRI1
therefore evaluates its arithmetic contribution at order one, and
transports that evaluation to the other order through the displayed
actual baseline difference. In particular
\[
 \begin{split}
 J_a={}&\mathcal B_a^{(1)}-(4m_0-2)\mathcal C_\Gamma q_a
       +\frac{\log2}{\pi}I_h a^2\\
      &+\mathcal W_a^{\rm mono}-\mathcal P_{a,-}-\mathcal P_{a,+}+E_{h,a}.
 \end{split}\tag{MRI5}
\]
This proves the exact backward receiver, retaining both penalties and
the full monic window.

In the simple-quartet domain of DRU12, $q_a=(a+1)^2$. Define the
already controlled analytic remainder by its actual value
\[
 \mathfrak R_{h,a}:=J_a-\delta_{a,1}^{\rm ar}
                          -C_Bq_a^2-8q_a\log a.
 \tag{MRI6}
\]
DRU12 proves $\mathfrak R_{h,a}=O_h(q_a)$. Equations MRI1 and MRI6
give the more informative exact decomposition
\[
 \begin{split}
 J_a={}&C_Bq_a^2+8q_a\log a-2\mathcal C_\Gamma q_a
        +\frac{\log2}{\pi}I_h a^2+\mathfrak R_{h,a}+E_{h,a}.
 \end{split}\tag{MRI7}
\]
Thus the identified arithmetic order-$q_a$ term has been propagated,
while the pre-existing analytic $O_h(q_a)$ remainder retains its
own value. In particular the proof does not assign a limiting
coefficient to $\mathfrak R_{h,a}/q_a$. The leading $C_Bq_a^2$ and
$8q_a\log a$ terms remain those of the proved complete action.
Dividing MRI1 by its literal dimension proves
\[
 \lim_{a\to\infty}\frac{\delta_{a,1}^{\rm ar}}{q_a}
 =\begin{cases}
 -2\mathcal C_\Gamma+(\log2/\pi)I_h,&m_0=1,\\
 -(4m_0-2)\mathcal C_\Gamma,&m_0\ge2.
 \end{cases}\tag{MRI8}
\]
The limits use $a^2/(a+1)^2\to1$ and
$a^2/([1+a(m_0-1)](a+1)^2)\to0$, respectively.

\subsection{The original allocation maps after the scalar update}
Let $G_M$ be the unchanged source form, $E$ the original lift matrix,
$\mathbf U_L$ its complete relation matrix, and $Z_k$ the literal
kernel inclusion. The same attained native form is
\[
 H_{Y,N}=E^*G_ME-E^*G_M\mathbf U_L
    (\mathbf U_L^*G_M\mathbf U_L)^{-1}\mathbf U_L^*G_ME,
 \qquad H_{K,N}=Z_k^*H_{Y,N}Z_k.
 \tag{MRI9}
\]
These are the exact receiving maps for the unresolved individual
allocation values. To verify the quotient map, write every lift as
$Ey+\mathbf U_Lz$ and complete the square in its original $G_M$
norm. The unique minimizing coefficient is
$z=-(\mathbf U_L^*G_M\mathbf U_L)^{-1}\mathbf U_L^*G_MEy$;
substitution gives MRI9. Restricting to $y=Z_kx$ gives its second
formula without discarding any cross block. This calculation also
proves why the relation determinant alone cannot determine the
kernel restriction: the latter depends on the explicitly retained
mixed Gram $E^*G_M\mathbf U_L$ and inclusion $Z_k$.

At every original endpoint, restriction followed by the attained
residual quotient gives Schur determinant factorization. Applying
the unchanged four-endpoint return to that factorization preserves
\[
 S_K+S_R=\mathcal T,\qquad
 4S_K=\mathcal L_e+V_b+V_{\rm rem}+V_s+V_A.
 \tag{MRI10}
\]
The new scalar MRI1 does not alter these maps. The current task is
to evaluate the seven displayed quantities on MRI9 and its actual
complementary quotient maps, through every term larger than $q_a$.
The PEI bounds below provide sharper correlated comparisons between
the primitive choices on these same spaces. The rank of an absolute
restriction or quotient remains in its determinant allowance.

\subsection{The exact reflected-pair receiving map}
Write $L:E\to W$ for the literal full cyclic observation and
$G_N$ for the original positive form. Choose coordinates on the
actual image $W_0=\operatorname{im}L$, and denote the resulting
surjective map by $L_0:E\to W_0$. Its attained form and section are
\[
 Q_{W_0,N}=(L_0G_N^{-1}L_0^*)^{-1},\qquad
 s_{W_0,N}=G_N^{-1}L_0^*Q_{W_0,N}.
 \tag{MRI11}
\]
Indeed $L_0s_{W_0,N}=I$ and, for $z\in\ker L_0$,
$z^*G_Ns_{W_0,N}=0$. Every other lift is this section plus such a
$z$, so Pythagoras proves the stated minimum and its form.
If $C$ embeds the two reflected observed classes into $W_0$, their
attained quotient metric is exactly $C^*Q_{W_0,N}C$. ET19--26
evaluate their Fourier observation Gram and prove eventual
injectivity at the displayed fixed original periods. MRI11 is
the exact morphism that carries those classes into the original
arithmetic quotient; its full inverse and all remaining image
directions are retained. The projected arithmetic responses
$U=K_{13}/F_{13}$ and $V=K_{23}/F_{23}$ still require evaluation
in their original three-column receiver, including the common
error multiplied by $\operatorname{Im}P_X$.

\paragraph{Edition corrections and mathematical scope.}
The Airy exterior phase in the incoming text is corrected to
$(V_n-2g+\ell)/2$ by SAI1. The two-window perturbation is the
explicit JW13 finite bound with its source-ratio norm; the earlier
bare allowance is superseded by this proved formula. All original
source texts are preserved, and each revised formula is linked to
its proof. The newly computed scalar resolves the earlier signed
baseline coefficient entry. The original proper-source standard
window, individual attained allocations and retained-class signs
remain active calculations with their displayed receiving maps.

% END PRESERVED INPUT 08_RECEIVER_UPDATES.tex
\endgroup

\clearpage
\begingroup


\setcounter{equation}{0}
\renewcommand{\theequation}{EPORIG.\arabic{equation}}
\section*{Residue prices, constrained composition and nonperiodic schedules}
\addcontentsline{toc}{section}{Residue prices, constrained composition and nonperiodic schedules}
% BEGIN PRESERVED INPUT 09_body.tex SHA256 973a8231ab4eedd589b640ae6a4b63450b57568ccf09c2d1246b4bec69869f73
\section{Residue prices, constrained composition, and nonperiodic optimal schedules}\label{eporig:residue-prices-constrained-composition-and-nonperiodic-optimal-schedules}

The Clankers. 17 September 2026.

\subsection{1. Source, definitions, and evidence boundary}\label{eporig:source-definitions-and-evidence-boundary}

This continuation starts from the integrated Erdős Problem 817 workbench at commit \path|dbd0a93f2cf935103e88e5c2b2b71fe85ae7537b|. The corrected 29-note edition, its \(L>3Q\) separation condition, its current cone-dual domain, and all earlier source identities remain upstream. The present addition changes none of them. Its cumulative delivery begins a new text-first series from that baseline rather than nesting the historical archives.

The original \(k=4\) mathematical construction remains credited to the anonymous/deleted Reddit contributor. The finite upper-bound Lean extension remains separately credited to \texttt{sneed-and-feed}. The present results have ordinary proofs and the specified exact computational certificates. They have not undergone new Lean elaboration or independent human mathematical review. No literature-priority claim is made.

For a finite set \(A\) of distinct positive integers, retain

\[
H_q(A)=\left\{\sum_{a\in A}c_a a:0\le c_a<q\right\},
\qquad S(A)=\sum_{a\in A}a.
\]

Each numerical value is counted once. The original coefficient words and their representation multiplicities remain available separately. The coefficient arity \(q\) and forbidden progression length \(k\) are different parameters.

A \textbf{canonical modular-\(k\) block} is \((b,A)\) with \(b>S(A)\) and with \(H_2(A)\bmod b\) containing no ordered \(k\)-term progression with nonzero modular step. Repetition later in such a modular tuple is allowed. This convention checks steps of small additive order as well as full-order steps. A sequence of these blocks gives

\[
P_0=1,\quad P_{j+1}=b_jP_j,\quad
A_w=\bigcup_{j<m}P_jA_j,\quad N(w)=\sum_{j<m}|A_j|.
\]

The actual generators are distinct, because every earlier generator is below the next position. Lowest-digit reduction proves that every binary prefix image is \(k\)-AP-free. Thus all of the applications below retain actual positive integer constructions.

The first new result gives sharp ternary residue prices at ranks one through four, with no upper bound on radix or generator height. The second determines the entire composition-constrained spectrum of the original radix-ten dictionary, including an interval where the optimal infinite rate is attained but no periodic schedule attains it. A general finite-horizon linear-programming theorem gives a uniform error for all prescribed composition profiles.

\subsection{2. Sharp residue prices with unbounded arithmetic parameters}\label{eporig:sharp-residue-prices-with-unbounded-arithmetic-parameters}

\textbf{Theorem 1.} Let \(k\in\{5,6\}\) and let \((b,A)\) be a canonical modular-\(k\) block. If \(n=|A|\le4\), then

\[
\boxed{|H_3(A)\bmod b|\ge L_n,\qquad
(L_1,L_2,L_3,L_4)=(3,5,13,23).}
\tag{2.1}
\]

All four constants are attained. The proof uses Kneser's classical addition theorem and a complete finite list of canonical obstructions. Both ingredients are stated below.

\subsubsection{2.1 The required binary counts}\label{eporig:the-required-binary-counts}

For distinct positive \(a_1<\cdots<a_n\), adjoining \(a_j\) produces at least \(j\) new values above \(S_{j-1}\): the values \(a_j+S_{j-1}\) and \(a_j+S_{j-1}-a_i\) for \(i<j\). Hence

\[
|H_2(A)|\ge1+\frac{n(n+1)}2.
\tag{2.2}
\]

At rank three, equality seven is possible only when the largest weight is the sum of the other two. Indeed, a binary collision is a signed relation on disjoint supports; positivity and distinctness leave only that possibility.

At rank four, equality eleven in (2.2) forces

\[
A=a\{1,2,3,4\}.
\tag{2.3}
\]

Here is the complete argument. Equality forces the first three weights to have exactly seven values, so write them as \(a<b<c=a+b\). Their image is

\[
D_3=\{0,a,b,c,a+c,b+c,2c\}.
\]

For a new \(d>c\), an intersection of \(D_3\) with \(d+D_3\) uses a lower representative in \(\{0,a,b\}\). There are at most three intersection values. Equality eleven requires all three. Their ordered upper representatives are precisely \(a+c,b+c,2c\). Consequently \(d=a+c\), \(d+a=b+c\), and \(d+b=2c\), which give \(b=2a\), \(c=3a\), \(d=4a\).

The binary image in (2.3) contains \(a[0,10]\). It therefore contains both five- and six-term nonconstant progressions. For the blocks of Theorem 1, the rank-four binary count is consequently at least twelve. The binary counts needed below are therefore four, seven, and twelve at ranks two, three, and four.

\subsubsection{2.2 Kneser's theorem and its complete case reduction}\label{eporig:knesers-theorem-and-its-complete-case-reduction}

We use the following classical statement (DeVos, 2013): for finite nonempty subsets \(U,V\) of an abelian group, with \(H=\operatorname{Stab}(U+V)\),

\[
|U+V|\ge|U+H|+|V+H|-|H|.
\tag{2.4}
\]

Apply it in \(\mathbb Z/b\mathbb Z\) to \(D=H_2(A)\). Put

\[
v=|D+D|=|H_3(A)\bmod b|,\quad h=|H|,\quad z=|D+H|.
\]

Then \(z\) is a multiple of \(h\), \(|D|\le z\le v\), and

\[
v\ge2z-h.
\tag{2.5}
\]

If \(h=1\), the binary lower counts immediately give the rank-three bound thirteen and the rank-four bound twenty-three.

If \(D+H\) has only one coset, it equals \(H\) because \(0\in D\). Let \(g=b/h\). Every actual generator is a multiple of \(g\), and the maps

\[
\iota:\mathbb Z/h\mathbb Z\overset\sim\longrightarrow H,
\quad[x]\longmapsto[gx],\qquad
A'=\{a/g:a\in A\}
\tag{2.6}
\]

preserve the original modular progression equations. The inverse divides the displayed representatives by \(g\). The positive scale \(g\), original weights, and original modulus are retained. In particular \(S(A')<h\), and \(A'\) is again a canonical modular-\(k\) block of the same rank. This is a specified subgroup isomorphism, not an assumption that every nontrivial stabilizer contains the original source.

For rank two, \(D\) has four elements. If \(v\le4\), then \(D+D=D\), making \(D\) a subgroup of order four. Its full cyclic orbit gives a forbidden ordered progression.

For rank three, assume \(v\le12\). The complete possibilities with \(h\ge2\) in (2.5) are

\[
(h,z)=(4,8),\qquad (h,h)\quad(7\le h\le12).
\tag{2.7}
\]

In the first case, seven source points in two cosets of size four force a full coset of order four. It gives an ordered six-term progression, and hence a five-term prefix. The remaining cases reduce by (2.6) to rank-three blocks with modulus at most twelve.

For rank four, assume \(v\le22\). The complete possibilities are

\[
(h,z)\in\{(2,12),(3,12),(4,12),(6,12),(7,14)\},
\quad (h,h)\ (12\le h\le22).
\tag{2.8}
\]

In the first four cases, \(|D|\ge12=z\) forces a full nontrivial coset of order at most six. In the \((7,14)\) case, one seven-element coset contains at least six of the source points. Removing at most one point from a cyclic seven-element orbit leaves a six-term consecutive progression in that orbit. Each case is forbidden for both \(k=5\) and \(k=6\). The one-coset cases again reduce by (2.6).

The lists (2.7) and (2.8) follow simply by enumerating the integer inequalities \(z\ge7\) or \(z\ge12\), \(h\mid z\), and \(2z-h\le12\) or \(22\). The proof data list every pair, including all one-coset cases.

\subsubsection{2.3 The finite obstruction lemma}\label{eporig:the-finite-obstruction-lemma}

The remaining finite statement is:

\begin{itemize}
\tightlist
\item
  every positive distinct triple with sum below its modulus \(b\le12\) has a modular six-term progression in its binary image;
\item
  every positive distinct quadruple with sum below \(b\le22\) has such a progression.
\end{itemize}

There are exactly 41 triples and 448 quadruples in these complete domains. The certificate additionally includes the one canonical pair at modulus at most four, for 490 source blocks in total.

For every source block the text certificate records its actual weights, modulus, start, nonzero step, six residue values, and six original binary subset masks. The following terminating enumeration specifies both the domain and witness test:

\begin{enumerate}
\def\labelenumi{\arabic{enumi}.}
\tightlist
\item
  For each stated modulus, enumerate all increasing positive tuples with sum below that modulus.
\item
  Form all binary subset sums, keeping their masks.
\item
  For every two distinct values \(x,y\), set \(d=y-x\pmod b\) and test all six values \(x+id\pmod b\).
\item
  Retain the first successful tuple and its masks.
\end{enumerate}

All 490 sources have a recorded witness. The separate auditor generates the increasing tuples recursively and checks every original mask. It also checks modular freeness in an additional complete domain by rotating residue bit sets, rather than using the producer's pair enumeration. This is finite proof data for the stated lemma; it is not a search extrapolation over larger moduli.

Equations (2.2)--(2.8), with this finite lemma, prove Theorem 1. At rank one the generator's residue has order at least three: order two would itself supply a forbidden alternating modular tuple. Thus its three ternary residues are distinct.

The sharp blocks are

\[
(3,\{1\}),\quad(5,\{1,2\}),\quad
(13,\{1,3,4\}),\quad(23,\{1,3,4,7\}).
\tag{2.9}
\]

Their binary modular freeness and their residue sections are checked directly. Their fifth-arity images are full integer intervals, whose repeated-image growth radii are their radices. This proves sharpness as both a residue count and a homogeneous image-growth bound.

\subsection{3. An all-rank generator-profile inequality}\label{eporig:an-all-rank-generator-profile-inequality}

For a level \((b,A)\), choose one actual ternary value \(d_r\in H_3(A)\) for each occurring residue \(r\). The map

\[
\{d_r\}\times Y\longrightarrow H_3(A)+bY,
\qquad(d_r,y)\longmapsto d_r+by
\tag{3.1}
\]

is injective: reduction modulo \(b\) recovers the residue and hence the chosen actual value; subtraction and division then recover \(y\). No requirement that all chosen ternary values be below \(b\) is used.

Iterating (3.1) through any changing canonical modular-\(k\) schedule proves

\[
|H_3(A_w)|\ge\prod_{j<m}|H_3(A_j)\bmod b_j|.
\tag{3.2}
\]

For arbitrary rank, the usual one-shift argument gives

\[
|H_2(A)|\ge\gamma_k^{|A|},\qquad
\gamma_k=\frac{k-1}{k-2}.
\tag{3.3}
\]

For clarity, before adjoining a generator \(a\), decompose the existing binary image into maximal \(a\)-chains. Each has length at most \(k-2\): adjoining \(a\) extends its end, so a longer chain would yield a forbidden \(k\)-AP. The union with its translate has size at least \(\gamma_k\) times the preceding size. Iteration proves (3.3). The canonical binary values remain distinct modulo the radix, so (3.3) also lower-bounds the number of ternary residues. This elementary general-bound mechanism belongs to the existing workbench and surrounding literature.

Let \(m_j\) count levels having exactly \(j\) generators, and let \(N_{\ge5}\) count all generators belonging to larger levels. Combining Theorem 1 and (3.2)--(3.3) gives the new explicit rank profile:

\[
\boxed{
|H_5(A_w)|\ge|H_3(A_w)|
\ge3^{m_1}5^{m_2}13^{m_3}23^{m_4}
\gamma_k^{N_{\ge5}},\qquad k=5,6.
}
\tag{3.4}
\]

It applies at every finite horizon and at unbounded changing heights and radices. It makes no assertion that the original generator list has a unique block presentation.

For the family in which every level is individually modular-\(k\)-free and has rank at most \(r\), the exact optimum fifth-image rate per generator is

\[
\boxed{3,\quad\sqrt5,\quad\sqrt5,\quad23^{1/4}
\qquad(r=1,2,3,4).}
\tag{3.5}
\]

The lower bound is (3.4); repeated sharp blocks (2.9) attain it. With exactly \(r\) generators on every level, the optimum is instead \(L_r^{1/r}\).

This domain is an explicitly specified subclass of general carry-safe schedules. For example the earlier singleton schedule with radices \(2,4\) has a first level whose binary image is the full group of order two and is individually modular-unsafe. Its two-level aggregation is the actual rank-two block \(\{1,2\}\) in radix eight. The aggregation keeps both generators and product radix eight. Thus the broader singleton optimum \(\sqrt8\) and (3.5)'s singleton optimum three concern different domains connected by a rank-changing aggregation map.

Set \(c=\log(23)/4\), \(N=N(w)\), and \(\theta=N_{\ge5}/N\). All four low-rank prices per generator are at least \(c\). Therefore

\[
\boxed{
\frac{\log |H_5(A_w)|}{N}
\ge(1-\theta)c+\theta\log\gamma_k.
}
\tag{3.6}
\]

Along any sequence of such words whose logarithmic image cost is at most \(\log u+o(1)\), with \(u<23^{1/4}\),

\[
\liminf\theta\ge
\frac{c-\log u}{c-\log\gamma_k}.
\tag{3.7}
\]

For the existing comparison rates \(u=97^{1/6}\) at \(k=5\) and \(u=1651^{1/10}\) at \(k=6\), the right sides lie respectively in

\[
\left(\frac{431}{10000},\frac{432}{10000}\right),
\qquad
\left(\frac{766}{10000},\frac{767}{10000}\right).
\tag{3.8}
\]

The certificate checks these intervals by integer powers after clearing rational logarithmic coefficients. Consequently at least about \(4.31\%\) or \(7.66\%\) of the generators must lie in blocks of rank at least five in any individually modular-safe presentation approaching those respective costs. This is a necessary profile constraint, not a new numerical bound on the unrestricted \(\lambda_k\).

\subsection{4. Composition-constrained capacity for any finite dictionary}\label{eporig:composition-constrained-capacity-for-any-finite-dictionary}

Let \(\mathcal D=\{(b_i,A_i):1\le i\le d\}\) be any finite canonical dictionary with nonempty positive distinct generator sets. Put \(n_i=|A_i|\). Modular admissibility is needed for arithmetic applications, but the count statements in this section hold for every such dictionary.

For a word \(w\) retain

\[
F_q(w)=|H_q(A_w)|,
\qquad \mathbf n(w)=(\text{number of each letter in }w).
\]

The numerical cut map \((x,y)\mapsto x+P_uy\) is onto. Its first value lies in \([0,(q-1)(P_u-1)]\), so each fibre has at most \(C=q-1\) elements. Hence

\[
\boxed{C^{-1}F_q(u)F_q(v)\le F_q(uv)\le F_q(u)F_q(v).}
\tag{4.1}
\]

This is the integrated finite-period theorem's original counting mechanism, retained with its exact cut constant.

For a profile \(p=(p_i)\) in the probability simplex, define \(\mathcal C_q(p)\) as the infimum of

\[
\liminf_{m\to\infty}\frac{\log F_q(w_0\cdots w_{m-1})}{m}
\]

over infinite words whose empirical letter frequencies converge to \(p\). The mean generator reward is

\[
\nu(p)=\sum_i n_i p_i>0.
\]

Thus the corresponding per-generator rate is \(\exp(\mathcal C_q(p)/\nu(p))\). Frequency in this section counts original levels; it is not silently identified with generator frequency when the \(n_i\) differ.

For each \(m\ge1\) solve the finite linear program

\[
E_m(p)=\min\left\{
\frac1m\sum_{w\in\mathcal D^m}\lambda_w\log F_q(w):
\lambda_w\ge0,\ \sum_w\lambda_w=1,
\sum_w\lambda_w\mathbf n(w)=mp
\right\}.
\tag{4.2}
\]

The program is feasible because the pure words supply the simplex vertices. An optimum uses at most \(d\) words: the \(d\) count constraints include the probability constraint, since each word has length \(m\).

\textbf{Theorem 2.} For every profile,

\[
\boxed{
E_m(p)-\frac{\log(q-1)}m
\le\mathcal C_q(p)\le E_m(p).
}
\tag{4.3}
\]

The convergence is uniform on the complete simplex. In particular \(\mathcal C_q\) is continuous and convex. Every profile has an infinite schedule attaining its value with a full logarithmic cost limit. Exact periodic attainment is a separate property.

\textbf{Proof.} Partition a long word into length-\(m\) chunks and a final shorter piece. Applying the lower inequality in (4.1) gives the sum of chunk log counts minus at most one \(\log C\) per chunk. Along a subsequence realizing the lower limit, extract a convergent empirical distribution of the finitely many chunks. Its mean letter profile is \(p\). Its mean chunk log count is at least \(mE_m(p)\), proving the lower bound. The final piece has bounded length and reward, so it does not change the limit.

For the upper bound, take an optimal mixture in (4.2). Concatenate its finitely many words in frequencies tending to that mixture. The upper inequality in (4.1) bounds the log count by the sum of their actual log counts. The resulting profile is \(p\), proving the upper bound.

The finite LP value is a continuous convex piecewise-affine function of \(p\). The uniform bound in (4.3) therefore gives continuity and convexity of the limit.

For actual attainment, choose increasing horizons and rational approximations to the corresponding optimal word mixtures, with profile and mean-cost errors tending to zero. Each rational approximation gives a finite period \(V_j\). Concatenate many copies of \(V_j\) before moving to \(V_{j+1}\). Choose the repetitions at stage \(j\) so the accumulated length is at least \(j|V_{j+1}|\). All partial-period frequency errors then tend to zero. The upper multiplicative bound gives a limiting upper cost at most \(\mathcal C_q(p)\); the already-proved lower bound gives the reverse lower limit. The full cost limit exists and equals \(\mathcal C_q(p)\). Every period and transition word is an actual original dictionary word. \(\square\)

Dividing (4.3) by \(\nu(p)\) gives a uniform per-generator logarithmic error at most \(\log(q-1)/(m\min_i n_i)\).

The finite dual is also explicit. A vector of logarithmic prices \(y_i\) satisfying

\[
\sum_i y_i n_i(w)\le\log F_q(w)
\quad\text{for every length-}m\text{ word}
\tag{4.4}
\]

gives the all-infinite lower certificate

\[
\mathcal C_q(p)\ge y\cdot p-\frac{\log C}{m}.
\tag{4.5}
\]

A primal and dual pair with equal values certifies the finite optimum. At rational profiles, all basis coefficients and all coefficients expressing \(y_i\) as linear combinations of integer logarithms are rational. Every inequality is checked by integer powers. These prices are auxiliary count budgets; they change no original radix, generator, or representation mass.

\subsection{5. Exact minima at every composition in the radix-ten source}\label{eporig:exact-minima-at-every-composition-in-the-radix-ten-source}

Retain the original dictionary

\[
\mathsf A=(10,\{1,3\}),\qquad
\mathsf B=(10,\{2,4\}).
\tag{5.1}
\]

Its binary residues are \(\{0,1,3,4\}\) and \(\{0,2,4,6\}\). Neither contains an order-two pair differing by five; a step of order five or ten would require at least five distinct values. Thus both are modular-five-free. Every schedule has actual five-admissible binary prefixes and two generators per level.

On the seven integrated fifth-image coordinates

\[
\{0\},\{0,1\},\{0,2\},\{0,1,2\},
\{0,3\},\{0,1,3\},\{0,1,2,3\},
\]

the unchanged arithmetic actions are

\[
\begin{split}
M_Af={}&(3f_0+7f_1,\ 2f_0+8f_1,\ f_0+9f_1,\ f_0+9f_1,\\
&10f_1,\ 10f_1,\ 10f_1),
\end{split}
\]

\[
\begin{split}
M_Bf={}&(2f_1+3f_3,\ 4f_1+6f_3,\ f_1+4f_3,\ 3f_1+7f_3,\\
&4f_1+6f_3,\ 3f_1+7f_3,\ 2f_1+8f_3).
\end{split}
\tag{5.2}
\]

The initial tail vector is \((1,2,2,3,2,3,4)^{\mathsf T}\). These matrices and the following rank-one identity are inherited from the integrated local-flow source and are rechecked independently here by interpolation on actual finite numerical tails:

\[
M_AM_B=uv,\qquad
u=(17,18,19,19,20,20,20)^{\mathsf T},\quad
v=(0,2,0,3,0,0,0).
\tag{5.3}
\]

For \(a,b\ge1\), define

\[
T(a,b)=vM_B^{b-1}M_A^{a-1}u
=\frac{8\,10^{a+b}+4\,10^b-3}{9}.
\tag{5.4}
\]

Both vector sequences satisfy the recurrence with characteristic polynomial \((X-1)(X-10)\). The four original values at \(a,b\in\{1,2\}\) prove (5.4) at every pair of lengths.

For a mixed cyclic word with successive runs \(A^{a_i}B^{b_i}\), its nonzero spectral value is

\[
\rho(M_w)=\prod_i T(a_{i+1},b_i),\qquad a_{s+1}=a_1.
\tag{5.5}
\]

This follows by multiplying the original rank-one factors. Cyclic rotation preserves the nonzero eigenvalues, not the original finite image count. The latter retains the bound

\[
\rho(M_w)\le F_5(w)\le4\rho(M_w).
\tag{5.6}
\]

For completeness, the weighted norm \(\|x\|=\max_R|x_R|/|R|\) gives \(\|M_w\|=F_5(w)\): translated copies supply the upper bound and the singleton coordinate supplies equality. Repetition and (4.1) prove (5.6).

Let \(F(s,t)\) be the minimum actual fifth-image count among words having exactly \(s\) letters \(A\) and \(t\) letters \(B\).

\textbf{Theorem 3.} For \(s,t\ge0\),

\[
F(0,0)=1,\quad
F(s,0)=\frac{16\,10^s-7}{9},\quad
F(0,t)=\frac{4\,10^t-1}{3}.
\tag{5.7}
\]

For \(s\ge t>0\), write \(s=at+r\), \(0\le r<t\). Then

\[
\boxed{F(s,t)=T(a,1)^{t-r}T(a+1,1)^r.}
\tag{5.8}
\]

For \(0<s<t\),

\[
\boxed{F(s,t)=93^{s-1}T(1,t-s+1).}
\tag{5.9}
\]

\subsubsection{5.1 The lower optimization}\label{eporig:the-lower-optimization}

Put \(G(a,b)=10^{a+b}/T(a,b)\). Its exact bounds are

\[
\frac{15}{14}<G(a,b)<\frac98.
\tag{5.10}
\]

Aligning large \(a\) with large \(b\) does not decrease the product of the gains. To check this directly, put \(x=10^{-a}\le x'=10^{-a'}\) and \(y=4-3\,10^{-b}\ge y'=4-3\,10^{-b'}\). The relevant numerator difference between the two gain products is

\[
(8+xy')(8+x'y)-(8+xy)(8+x'y')
=8(x'-x)(y-y')\ge0.
\tag{5.11}
\]

If the number of crossed run pairs is below \(\min(s,t)\), some \(a\ge2\) and some \(b\ge2\) exist. Re-pair them using (5.11), and split that pair into two positive pairs. The gain product increases strictly, since

\[
(15/14)^2>9/8.
\]

Every such re-pairing is realized by a cyclic run word with the same original letter counts. Hence an optimal spectral product has exactly \(\min(s,t)\) run pairs.

When \(s\ge t\), all \(b\)-lengths are one. Put \(K_a=T(a,1)=(80\,10^a+37)/9\). The identity

\[
K_aK_{a+2}-K_{a+1}^2=2960\,10^a>0
\tag{5.12}
\]

proves strict logarithmic convexity. Transferring one unit from a run at least two longer to a shorter run decreases the cost. Thus the lengths are \(a\) and \(a+1\) with the multiplicities in (5.8).

When \(s<t\), all \(a\)-lengths are one. Put \(J_b=T(1,b)=(28\,10^b-1)/3\). Here

\[
J_bJ_{b+2}-J_{b+1}^2=-252\,10^b<0.
\tag{5.13}
\]

Strict logarithmic concavity moves all excess \(B\)-length to one run and gives (5.9) as the minimum spectral product. Inequality (5.6) transfers both lower bounds to the original counts.

\subsubsection{5.2 Actual finite attainers}\label{eporig:actual-finite-attainers}

The required endpoint equalities are supplied by original numerical images:

\[
H_5(BA^a)=2[0,K_a-1],\qquad
H_5(B^bA)=2[0,J_b-1].
\tag{5.14}
\]

For the first identity, the upper \(A\)-levels have interval image \([0,16(10^a-1)/9]\). The low \(B\)-image is \(2[0,12]\). A single division of the whole numerical set by two gives

\[
[0,12]+5[0,16(10^a-1)/9]
=[0,K_a-1].
\]

For the second, the lower \(B\)-levels, after that same whole-image division, give \([0,12(10^b-1)/9]\). Adding the upper \(A\)-level fills consecutive intervals and gives \([0,J_b-1]\). The map \(x\mapsto x/2\) has inverse \(x\mapsto2x\) on these specified images. No individual source coordinate is dropped.

The divided digit maxima in (5.14) are below the actual macro radices \(10^{a+1}\) and \(10^{b+1}\). Thus concatenating the macros gives exact multiplication of their counts. The attaining words are

\[
(BA^a)^{t-r}(BA^{a+1})^r\quad(s\ge t),
\tag{5.15}
\]

\[
B^{t-s+1}A(BA)^{s-1}\quad(s<t).
\tag{5.16}
\]

Their original generator sets and radices are the original concatenations. Equations (5.14)--(5.16) establish equality in the spectral lower bounds and complete the proof of Theorem 3.

\subsection{6. The entire frequency spectrum and an unattained periodic regime}\label{eporig:the-entire-frequency-spectrum-and-an-unattained-periodic-regime}

Let \(p\) be the limiting fraction of \(B\)-levels. The mean generator reward is two. Let \(\mathcal C(p)\) denote the optimal logarithmic fifth-image cost per level.

\textbf{Theorem 4.} At \(p=0,1\), \(\mathcal C(p)=\log10\). For \(0<p\le1/2\), write

\[
a=\left\lfloor\frac{1-p}{p}\right\rfloor,\qquad
u=\frac{1-p}{p}-a.
\]

Then

\[
\boxed{\mathcal C(p)=p\bigl((1-\nu)\log K_a+\nu\log K_{a+1}\bigr).}
\tag{6.1}
\]

For \(1/2\le p\le1\),

\[
\boxed{\mathcal C(p)=(1-p)\log93+(2p-1)\log10.}
\tag{6.2}
\]

The per-generator image rate is \(\exp(\mathcal C(p)/2)\).

\textbf{Proof.} Theorem 3 bounds every finite prefix below by the formula at its actual composition. Passing to the prescribed limiting frequency yields (6.1)--(6.2); at the junctions the expressions agree. The attainable schedules below supply the reverse bounds.

For \(0<p\le1/2\), concatenate \(BA^a\) and \(BA^{a+1}\) with macro frequencies \(1-\nu\) and \(\nu\). The number of \(B\)-levels per macro is one and the mean macro length is \(a+1+\nu=1/p\). The exact product in (5.14) proves (6.1). Rational macro frequencies can be periodic. For irrational frequencies, the difference sequence of \(\lfloor j\nu\rfloor\) supplies a specific nonperiodic choice with bounded discrepancy.

For \(1/2<p<1\), at stage \(j\) use

\[
(BA)^{\lfloor(1-p)j\rfloor}
B^{\lfloor(2p-1)j\rfloor}.
\tag{6.3}
\]

Each stage has length \(j+O(1)\), while accumulated length is \(j^2/2+O(j)\). The empirical frequencies converge to \(p\) at every prefix, including prefixes inside a stage. The stage count is at most

\[
93^{\lfloor(1-p)j\rfloor}
\frac{4\,10^{\lfloor(2p-1)j\rfloor}-1}{3}.
\]

Submultiplicativity makes its total logarithmic overhead \(O(j)\), which vanishes per accumulated length. This gives (6.2). Pure schedules attain the endpoints. \(\square\)

The function is continuous and convex. On \([1/(a+2),1/(a+1)]\) it is affine with slope

\[
s_a=(a+2)\log K_a-(a+1)\log K_{a+1}.
\tag{6.4}
\]

Identity (5.12) gives \(s_a>s_{a+1}\), and \(s_a\to-\log(9/8)\). The final affine segment has slope \(\log(100/93)\). The unconstrained minimum is at \(p=1/3\), recovering the earlier rate \(893^{1/6}\) rather than improving it.

For a fixed rational \(p\in(1/2,1)\), no periodic schedule attains (6.2). Indeed any period with \(s\) letters \(A\) and \(t>s\) letters \(B\) has spectral radius at least \(F(s,t)\). With \(b=t-s+1\ge2\),

\[
\boxed{
\frac{F(s,t)}{93^s10^{t-s}}
=\frac{280}{279}-\frac{1}{279\,10^{b-1}}>1.
}
\tag{6.5}
\]

Repeating the period retains this strict logarithmic excess per period. Yet (6.3) attains the constrained optimum. Irrational profiles of course also exclude a periodic frequency, but (6.5) gives the stronger obstruction at every rational profile in this open interval.

The excess has a local explanation. For a crossed run pair \((a,b)\), relative to the linear branch (6.2), its logarithmic defect is

\[
\epsilon(a,b)=a\log(100/93)-\log G(a,b)\ge0.
\tag{6.6}
\]

Equality holds only at \((1,1)\). For \(a=1,b\ge2\) it lies between \(\log(311/310)\) and \(\log(280/279)\). For \(a\ge2\), it is bounded below by a positive constant times \(a-1\), since \(\log G(a,b)<\log(9/8)<2\log(100/93)\). It is also bounded above by a constant times \(a-1\) in that range. Consequently optimal schedules on the open linear branch have a vanishing density of merged \(A\)-runs and of nontrivial \(B\)-runs. Their excess \(B\) mass remains positive, so the lengths of the latter runs must grow. Formula (6.3) realizes that precise regime.

An auxiliary price \(\xi\) per \(B\)-level makes the optimal value \(\min_p(\mathcal C(p)+\xi p)\). Equations (6.4) give infinitely many transition prices \(-s_a\), accumulating at \(\log(9/8)\), and periodic phases \(BA^a\) with unbounded lengths. The parameter is a specified composition price; it does not change the arithmetic radices or add generators. This observation explains why a parameter-independent bound for exact optimal periods would discard real information, even though the finite-horizon approximation in Theorem 2 is uniform.

\subsection{7. Original receiving maps and a supported empty position}\label{eporig:original-receiving-maps-and-a-supported-empty-position}

The numerical cut inequality (4.1) comes from the original onto map

\[
\pi_{u,v}:H_q(A_u)\times H_q(A_v)\to H_q(A_{uv}),
\quad(x,y)\mapsto x+P_uy.
\tag{7.1}
\]

If \(q_u,q_v\) are the original coefficient-word to local-value quotients, the relation sequence is

\[
0\to\ker(q_u\otimes q_v)
\to\ker\bigl(\pi_{u,v}(q_u\otimes q_v)\bigr)
\xrightarrow{q_u\otimes q_v}\ker\pi_{u,v}\to0.
\tag{7.2}
\]

Surjectivity at the right follows by choosing an actual coefficient-word representative of each local value. The original local boundary space remains separate from the extra cut kernel.

For positive original pair masses \(w(x,y)=\mu_u(x)\mu_v(y)\), put

\[
W_z=\sum_{x+P_uy=z}w(x,y).
\]

The minimum-norm section in the original quotient metric is

\[
s(e_z)=\sum_{x+P_uy=z}\frac{w(x,y)}{W_z}e_{(x,y)},
\qquad G_{\rm receiving}=\operatorname{diag}_z(1/W_z).
\tag{7.3}
\]

Neither the cut cardinality constant nor a dual composition price replaces this metric.

A simple original example distinguishes removing generators from erasing an active radix position. In the word \(ABA\), retain only the two \(A\)-levels. Keeping the actual positions gives generators

\[
\{1,3,100,300\},\qquad |H_5|=289.
\]

Erasing the middle radix gives instead

\[
\{1,3,10,30\},\qquad |H_5|=177.
\]

On the distinct local values \([0,16]^2\), the first evaluation \((x,y)\mapsto x+100y\) is injective. The second \((x,y)\mapsto x+10y\) is onto \([0,176]\). The additional receiving kernel of the free linearized map has dimension \(289-177=112\). All original pair values remain before taking that quotient.

The full block word after deleting the middle generators still has an active empty position with radix ten. Its source has one empty choice and a present numerical zero. Absorbing its radix into the preceding level keeps the first numerical set exactly. Simply omitting its cost produces the second set and the explicitly calculated new kernel.

For comparison, if \(u\) is obtained from a word \(w\) by erasing \(d\) original letters of total generator reward \(E\), repeated use of (4.1) gives the fully costed inequality

\[
-2d\log(q-1)\le\log F_q(w)-\log F_q(u)
\le E\log q+d\log(q-1).
\tag{7.4}
\]

It follows by inserting one removed block between its two retained neighbours and bounding both original cuts. In a dictionary of individually modular-safe levels, the contracted word is again admissible; in a more general carry-controlled language that preservation requires its own witness. The two operations are never identified without their maps and costs.

The source framework for (7.2) is the original SplitZero internal quotient and extra receiving-kernel square, equations D6--D8 in the inspected support-diagram source. On a nonbottom label, a boundary maps to that label's coefficient zero. A count or residue observation can introduce an additional kernel, recorded above, rather than changing the old relation module or global absence. This contribution transfers those interfaces to explicit finite arithmetic sources; it supplies no new analytic theta estimate or Riemann-hypothesis conclusion.

\subsection{8. Verification, reproducibility, and remaining general question}\label{eporig:verification-reproducibility-and-remaining-general-question}

The producer retains all 490 finite canonical obstruction blocks and their original six-term masks. A separate scan covers all 36,310 canonical blocks of ranks at most four and radix at most forty. The auditor checks 31,083 admissible source/length cases and their actual finite Kneser stabilizers, the sharp residue sections and all 44 weighted residue fibres, and 84 original mixed-rank image joins. These bounded checks corroborate the general reduction; the all-height theorem rests on Section 2 and the complete exceptional certificate.

For composition, all 524,286 binary words of lengths one through eighteen are accounted for, with their literal multiplicities. The 189 fixed-composition classes agree with Theorem 3. The producer uses prefix rows; the auditor recovers the two matrices by interpolation on original numerical sets and propagates tail columns. It checks 126 original-generator images directly, all 961 attainers with \(0\le s,t\le30\), and the actual phase-separated prefixes in the stated finite domain.

There are 84 finite composition LP certificates, including unequal generator rewards and boundary profiles. Each stores its complete count-signature domain, rational primal weights, logarithmic dual prices, and exact cut-factor correction. The separate auditor checks 570 dual inequalities and 360 further original-generator images. Eight deliberately corrupted records are rejected. It imports neither the producer nor its helper library.

The command pair is

\begin{Shaded}
\begin{Highlighting}[]
\ExtensionTok{python}\NormalTok{ research/residue\_composition/certificates/verify\_residue\_composition.py}
\ExtensionTok{python}\NormalTok{ research/residue\_composition/certificates/audit\_residue\_composition.py}
\end{Highlighting}
\end{Shaded}

All check functions remain active with \texttt{python\ -O}. Normal and optimized receipts are compared byte-for-byte in the delivery record. The one compiled PDF is a reading copy of this full text, not a substitute for the source or certificate. The package contains no earlier ZIP, compressed proof payload, font file, or compiled Lean binary.

The unbounded problem remains the best arithmetic cost over all admissible block ranks. The rank-profile theorem imposes a uniform constraint on original individually modular-safe presentations approaching the known costs. It does not prove that enough generators must remain in the low-rank part to obtain a stronger unrestricted lower bound. The profile spectrum solves all prescribed mixtures in its original two-block family and supplies uniform constrained approximation for every finite dictionary; it does not assert a rank cutoff for the entire problem.

The continuation therefore contributes an all-height arithmetic lower certificate and a complete composition optimization, while keeping the remaining all-rank inequality visible. Existing results, benchmark rates, and mathematical credits retain their earlier scopes.

\subsection{References}\label{eporig:references}

DeVos, M. (2013). \emph{A short proof of Kneser's addition theorem for abelian groups} (arXiv:\allowbreak1303.3539v1). The primary abstract states the exact stabilizer theorem used in (2.4). The classical theorem is imported with attribution; no claim is made to reprove it or newly formalize it here. \url{https://arxiv.org/abs/1303.3539}

The Clankers. (2026, September 16). \emph{Exact local-pattern quotients and switching image capacity}. Integrated workbench source, \path|research/local_flow/notes/local-flow-and-switching.md|, at commit \path|dbd0a93f2cf935103e88e5c2b2b71fe85ae7537b|. Role: original seven matrices and crossed-run identity; these finite inputs are independently replayed in this contribution.

The Clankers. (2026, September 16). \emph{Uniform finite-period approximation and exact finite-word switching minima}. Integrated workbench source, \path|research/finite_period/notes/finite-period-control.md|, at the same commit. Role: cut-fibre estimate and preceding unconstrained approximation theorem. The composition LP and fixed-count spectrum above extend those results.

KokunoYumeto. (2026). \emph{Reconstruction with changes of support index} {[}TeX source{]}. Zeta Function Research Reader, \path|workbenches/splitzero-tandem/tex/support_diagrams.tex|, revision \path|58626a62cd5648e8ae7450fb54d4a4aab2330981|, blob \path|d3493f891291ee6e94dbf2c77649f7d85d240df2|, equations D6--D8. Role: original internal quotient and additional receiving-kernel interface. No new analytic statement from that workbench is assumed.

% END PRESERVED INPUT 09_body.tex
\endgroup

\clearpage
\begingroup
\def\Z{\mathbb Z}
\def\C{\mathbb C}
\def\R{\mathbb R}

\setcounter{equation}{0}
\renewcommand{\theequation}{EP.\arabic{equation}}
\section*{Exact supported-quotient and metric interface for EP817}
\addcontentsline{toc}{section}{Exact supported-quotient and metric interface for EP817}
% BEGIN PRESERVED INPUT 10_EXACT_SUPPORTED_METRIC_RETURN.tex SHA256 fabebe11a8f7981313a2b8c277051b9c63fa7f35be976ad9eef5aee6960e3b9d



This note integrates the supplied \emph{Residue prices, constrained composition,
and nonperiodic optimal schedules} of 17 September 2026
\cite{ep:continuation}.  Its full proof, original LaTeX, and exact certificates
are preserved without alteration in the accompanying source subtree.
The maps below implement the internal quotient and receiving-kernel construction
of the original Split-Zero support diagram, equations (D6)--(D8)
\cite{ep:support}.  All word multiplicities, actual integer values, moduli,
support labels, and Euclidean word metrics are retained.  In the
composition results the coefficient arity is an integer \(q\geq2\);
the applications at \(q=3,5\) therefore have the original cut constant \(q-1\).

\section{The actual source and its two cochain complexes}

Let \(A\) be a fixed nonempty finite set of distinct positive integers and
let \(b>\sum_{a\in A}a\) be the original radix.  Set
\[
 \Omega=\{0,1,2\}^{A},\qquad
 v(c)=\sum_{a\in A}c_a a,\qquad
 \bar v(c)=v(c)\bmod b .
 \tag{E1}
\]
For every \(S\subseteq\Omega\) use the actual images
\[
 X_S=v(S)\subset\Z,\quad R_S=\bar v(S)\subset\Z/b\Z,\quad
 V_S=\Z[S],\quad E_S=\Z[X_S],\quad \bar E_S=\Z[R_S].
 \tag{E2}
\]
Here square brackets denote free abelian groups on the indicated distinct
bases, including a basis vector \(e_0\) whenever the image contains zero.
The maps on these displayed bases are
\[
 q_S(e_c)=e_{v(c)},\qquad
 \pi_S(e_x)=e_{x\bmod b},\qquad p_S=\pi_Sq_S.
 \tag{E3}
\]
All three have their stated codomains by (E2), and both \(q_S,p_S\) are
surjective.  Their kernels are
\[
 B_S=\ker q_S,\qquad B'_S=\ker p_S,\qquad B_S\subseteq B'_S .
 \tag{E4}
\]
For an empty image its free group is zero and the same formulas apply.

Define two complexes concentrated in degrees \(-1,0\):
\[
 C_S=[\,B_S\overset{\iota_S}{\longrightarrow}V_S\,],\qquad
 D_S=[\,B'_S\overset{\iota'_S}{\longrightarrow}V_S\,].
 \tag{E5}
\]
The map \(f_S:C_S\to D_S\) is inclusion in degree \(-1\) and identity
in degree zero.  It is a cochain map because
\(\iota'_S(B_S\hookrightarrow B'_S)=\iota_S\).

\begin{theorem}[The exact receiving kernel in the original complexes]
The maps in (E3) induce the isomorphisms and exact sequence
\[
 \begin{gathered}
 H^0(C_S)\overset{\sim}{\longrightarrow}E_S,\quad[c]\longmapsto q_Sc,
 \qquad
 H^0(D_S)\overset{\sim}{\longrightarrow}\bar E_S,\quad[c]\longmapsto p_Sc,\\
 H^{-1}(C_S)=H^{-1}(D_S)=0,\qquad H^0(f_S)=\pi_S,\\
 0\longrightarrow B_S\longrightarrow B'_S
       \overset{q_S}{\longrightarrow}\ker\pi_S\longrightarrow0,\qquad
 B'_S/B_S\overset{\sim}{\longrightarrow}\ker H^0(f_S),
       \quad[c]\longmapsto q_Sc .
 \end{gathered}
 \tag{E6}
\]
\end{theorem}
\begin{proof}
Injectivity of the two differentials gives the degree-\(-1\) statements.
Their degree-zero cohomologies are \(V_S/B_S\) and \(V_S/B'_S\).
The maps \(q_S,p_S\) are onto with exactly those kernels, proving the two
displayed isomorphisms.  Their commutative square is (E3), so their induced
degree-zero map is precisely \(\pi_S\).

If \(c\in B'_S\), then \(\pi_Sq_Sc=p_Sc=0\), so the third map of the
sequence has the asserted codomain.  Its kernel is exactly \(B_S\).
To prove its surjectivity without an unspecified lift, order the original
coefficient words lexicographically using the increasing order on \(A\).
For each \(x\in X_S\), let \(c_x\) be the first word in \(S\) with value
\(x\).  An element \(h=\sum_xh_xe_x\in\ker\pi_S\) is the image under
\(q_S\) of \(\sum_xh_xe_{c_x}\); that lift belongs to \(B'_S\) because
its \(p_S\)-image is \(\pi_Sh=0\).  This proves the exact sequence and
the final isomorphism.
\end{proof}

For \(S\subseteq T\), use basis inclusions in each of \(V,E,\bar E\).
The squares with \(q\) and \(\pi\) commute on every basis vector.
Consequently \(B_S\to B_T\) and \(B'_S\to B'_T\) are the restrictions
of \(V_S\to V_T\), and the complexes and (E6) form actual diagrams.
Use the join-semilattice
\[
 L=\{\bot\}\sqcup\mathcal P(\Omega),\qquad S\vee T=S\cup T,
 \tag{E7}
\]
with a distinct bottom \(\bot\) and zero fibre there.  The label
\(\varnothing\in\mathcal P(\Omega)\) remains distinct from \(\bot\),
despite its zero amplitude group.  Over each displayed integral diagram
the split carrier consists of \((S,x)\), with addition by transport to
\(S\cup T\), supported scalar multiplication \(r^\bullet(S,x)=(S,rx)\),
and \(\tau(S,x)=(\bot,0)\).  These are the original operations
of \cite[(D2)]{ep:support}.  Commutativity of the above squares proves
that the reconstructed \(q,\pi,f\) are split-linear.

In particular the all-label zero-amplitude preimage of the reconstructed
\(\pi\) is
\[
 \coprod_{S\in L}\ker\pi_S .
 \tag{E8}
\]
The preimage of its global zero contains only the bottom fibre, since its
label map is identity.  The exact comparison is the inclusion of that
bottom fibre into (E8); (E6) specifies each nonbottom fibre of (E8).
A supported boundary \((S,c)\) goes to \((S,0)\) at the same label.
Also the residue-amplitude map
\(\bar E_S\to\Z/b\Z,\ \sum_rz_re_r\mapsto\sum_rz_rr\) is integral linear;
it sends the present vector \(e_0\) to zero.  These statements require no
map from \(\C\) or \(\mathbb Q\) to the finite group.

\section{Original masses and an evaluated determinant loss}

For the metric statements extend the preceding groups and maps to \(\C\).
Place the orthonormal metric on the original word basis of \(\C[S]\).
Define the positive masses only on their actual images:
\[
 \mu_S(x)=|\{c\in S:v(c)=x\}|,\qquad
 \bar\mu_S(r)=\sum_{\substack{x\in X_S\\x\bmod b=r}}\mu_S(x).
 \tag{E9}
\]
No zero-mass value is inserted into an inverse diagonal matrix.

\begin{theorem}[Exact integral-section loss]
The quotient metric on \(E_S\otimes\C\), the residue quotient metric,
and the unique minimum-norm right inverse of \(\pi_S\) are respectively
\[
 G_{E,S}=\operatorname{diag}_{x\in X_S}\frac1{\mu_S(x)},\quad
 G_{\bar E,S}=\operatorname{diag}_{r\in R_S}\frac1{\bar\mu_S(r)},\quad
 s_S(e_r)=\sum_{x\bmod b=r}\frac{\mu_S(x)}{\bar\mu_S(r)}e_x .
 \tag{E10}
\]
Let \(d_S(r)\) be the least actual integer value in the residue fibre
and let \(\sigma_S(e_r)=e_{d_S(r)}\).  Put \(K_S=\sigma_S-s_S\).
Then
\[
 \begin{gathered}
 \pi_SK_S=0,\qquad
 \sigma_S^*G_{E,S}\sigma_S
   =G_{\bar E,S}+K_S^*G_{E,S}K_S,\\
 K_S^*G_{E,S}K_S
   =\operatorname{diag}_{r\in R_S}
       \left(\frac1{\mu_S(d_S(r))}-\frac1{\bar\mu_S(r)}\right),\\
 \log\frac{\det(\sigma_S^*G_{E,S}\sigma_S)}{\det G_{\bar E,S}}
   =\sum_{r\in R_S}\log\frac{\bar\mu_S(r)}{\mu_S(d_S(r))}.
 \end{gathered}
 \tag{E11}
\]
Every determinant is in the displayed original residue basis.  Empty
determinants are one.
\end{theorem}
\begin{proof}
For a fixed \(x\), minimizing \(\sum_{v(c)=x}|z_c|^2\) subject to
\(\sum_{v(c)=x}z_c=u_x\) gives the unique solution
\(z_c=u_x/\mu_S(x)\) and cost \(|u_x|^2/\mu_S(x)\), by completing
the sum of squared deviations.  Different value fibres are disjoint.
This proves \(G_{E,S}\).

Within residue \(r\), write \(u_x=\mu_S(x)t/\bar\mu_S(r)+h_x\)
where \(\sum_{x\bmod b=r}h_x=0\).  Direct expansion gives
\[
 \sum_{x\bmod b=r}\frac{|u_x|^2}{\mu_S(x)}
 =\frac{|t|^2}{\bar\mu_S(r)}
  +\sum_{x\bmod b=r}\frac{|h_x|^2}{\mu_S(x)} .
 \tag{E12}
\]
The cross term is zero because its coefficient is
\(\sum_xh_x=0\).  The displayed cost is uniquely minimized by \(h=0\).
This proves the last two parts of (E10), as well as orthogonality of
\(\operatorname{im}s_S\) to \(\ker\pi_S\).  Both \(s_S\) and
\(\sigma_S\) are right inverses, so \(K_S\) has image in that kernel.
Substituting \(h=K_St\) in (E12) proves the first identity in (E11).
Columns for different residues have disjoint value support.
The squared norm of \(\sigma_Se_r\) is
\(1/\mu_S(d_S(r))\), while that of \(s_Se_r\) is
\(1/\bar\mu_S(r)\).  This proves the diagonal formula, and taking
the quotient of the diagonal determinants proves the final formula.
\end{proof}

Thus the squared generalized singular values of the actual integral
section relative to the attained quotient metric are
\(\bar\mu_S(r)/\mu_S(d_S(r))\).  Formula (E11) evaluates its full loss,
including every original representation multiplicity.

\section{Source enlargement has an exact metric boundary}

Write \(i_{ST}:E_S\otimes\C\to E_T\otimes\C\) and
\(\bar i_{ST}:\bar E_S\otimes\C\to\bar E_T\otimes\C\) for the original
basis inclusions when \(S\subseteq T\).  Extend \(\mu_S(x)\) by zero
for \(x\in X_T\setminus X_S\).  The minimum-norm section defect is the
map with its stated source and target
\[
 K_{ST}=i_{ST}s_S-s_T\bar i_{ST}:
    \bar E_S\otimes\C\longrightarrow\ker\pi_T.
 \tag{E13}
\]
\begin{proposition}[Full source-change norm and its composition law]
For every \(r\in R_S\),
\[
 \begin{split}
 K_{ST}e_r
   &=\sum_{\substack{x\in X_T\\x\bmod b=r}}
      \left(\frac{\mu_S(x)}{\bar\mu_S(r)}
           -\frac{\mu_T(x)}{\bar\mu_T(r)}\right)e_x,\\
 \|K_{ST}e_r\|_{G_{E,T}}^2
   &=\frac{1}{\bar\mu_S(r)^2}
        \sum_{x\bmod b=r}\frac{\mu_S(x)^2}{\mu_T(x)}
        -\frac1{\bar\mu_T(r)},\\
 0\leq\|K_{ST}e_r\|_{G_{E,T}}^2
   &\leq\frac1{\bar\mu_S(r)}-\frac1{\bar\mu_T(r)}.
 \end{split}
 \tag{E14}
\]
The columns for distinct residues are orthogonal in \(G_{E,T}\).  For
\(S\subseteq T\subseteq U\) the exact typed composition identity is
\[
 K_{SU}=i_{TU}K_{ST}+K_{TU}\bar i_{ST}.
 \tag{E15}
\]
\end{proposition}
\begin{proof}
The two right-inverse identities and the commutative square give
\(\pi_TK_{ST}=\bar i_{ST}-\bar i_{ST}=0\).
Insertion of (E10) gives the first formula.  In the \(G_{E,T}\) norm,
expansion of its square has cross term
\(-2/\bar\mu_T(r)\) and last term \(1/\bar\mu_T(r)\), yielding the
second formula.  Nonnegativity follows either from this norm or
Cauchy--Schwarz.  Since \(S\subseteq T\) gives
\(0\leq\mu_S(x)\leq\mu_T(x)\), one has
\(\sum_x\mu_S(x)^2/\mu_T(x)\leq\bar\mu_S(r)\), giving the upper bound.
The disjoint residue supports prove column orthogonality.  Finally,
\[
 i_{TU}(i_{ST}s_S-s_T\bar i_{ST})
 +(i_{TU}s_T-s_U\bar i_{TU})\bar i_{ST}
 =i_{SU}s_S-s_U\bar i_{SU},
 \]
which is (E15), with all domains as in (E13).
\end{proof}

The integral sections have the parallel defect
\[
 i_{ST}\sigma_Se_r-\sigma_T\bar i_{ST}e_r
     =e_{d_S(r)}-e_{d_T(r)}\in\ker\pi_T .
 \tag{E16}
\]
Its squared \(G_{E,T}\)-norm is zero if the representatives coincide,
and otherwise is \(1/\mu_T(d_S(r))+1/\mu_T(d_T(r))\), because those
are two different original basis vectors.  The same addition and
subtraction proves its version of (E15).

\section{A universal derivation of the seven arithmetic matrices}

The composition source retains
\[
 A=(10,\{1,3\}),\quad B=(10,\{2,4\}),\quad
 D_A=H_5(\{1,3\})=[0,16],\quad
 D_B=H_5(\{2,4\})=2[0,12].
 \tag{E17}
\]
For \(A\), the intervals \([3j,3j+4]\), \(0\leq j\leq4\), overlap
and fill \([0,16]\).  For \(B\), the original values are
\(2(i+2j)\) with \(0\leq i,j\leq4\); the intervals
\([2j,2j+4]\) overlap and fill \([0,12]\).  Thus (E17) is an equality
of original numerical sets.

Keep the seven source states in their stated order:
\[
 \mathcal R=((0),(0,1),(0,2),(0,1,2),(0,3),(0,1,3),(0,1,2,3)).
 \tag{E18}
\]
For any finite \(Y\subset\Z\), put \(f_R(Y)=|Y+R|\).  For
\(X=A,B\), \(R\in\mathcal R\), and \(j\in\{0,\ldots,9\}\), define
\[
 C_{X,R,j}
  =\left\{\frac{d+r-j}{10}:
        d\in D_X,\ r\in R,\ d+r\equiv j\pmod {10}\right\}.
 \tag{E19}
\]
If this set is nonempty, write \(m_{X,R,j}=\min C_{X,R,j}\) and
\(R_{X,R,j}=C_{X,R,j}-m_{X,R,j}\).
The change of presentation is the explicit bijection
\[
 Y+R_{X,R,j}\longrightarrow
 (D_X+10Y+R)\cap(j+10\Z),\qquad
 z\longmapsto j+10(m_{X,R,j}+z),
 \tag{E20}
\]
whose inverse is \(t\mapsto(t-j)/10-m_{X,R,j}\).  Each shift and
residue is retained in \path{UNIVERSAL_INTERFACE_CHECKS.json}.
Direct insertion of the finite sets in (E17)--(E19) gives the following
numbers of residues of each resulting state:
\[
 M_A=\begin{pmatrix}
3&7&0&0&0&0&0\\2&8&0&0&0&0&0\\1&9&0&0&0&0&0\\
1&9&0&0&0&0&0\\0&10&0&0&0&0&0\\0&10&0&0&0&0&0\\
0&10&0&0&0&0&0
\end{pmatrix},
\quad
 M_B=\begin{pmatrix}
0&2&0&3&0&0&0\\0&4&0&6&0&0&0\\0&1&0&4&0&0&0\\
0&3&0&7&0&0&0\\0&4&0&6&0&0&0\\0&3&0&7&0&0&0\\
0&2&0&8&0&0&0
\end{pmatrix}.
 \tag{E21}
\]
For completeness these entries can be obtained without a numerical tail:
enumerate each of the at most \(25\cdot4\) displayed pairs \((d,r)\),
place \((d+r-j)/10\) into its residue set in (E19), subtract its recorded
minimum with inverse (E20), and count the state.  The accompanying
independent verifier records all \(140\) such original fibres and
checks each entry of (E21).

The residues in (E20) are disjoint.  Summing their cardinalities proves,
for \emph{every} finite numerical tail \(Y\), the identities
\[
 f(D_X+10Y)=M_Xf(Y),\qquad
 f(\{0\})=(1,2,2,3,2,3,4)^{\mathsf T}.
 \tag{E22}
\]
The matrices thus describe the exact count observations of the retained
numerical sets.  They are induced through the original residue partition
and bijections (E20); no change to coefficient-word amplitudes is involved.
Multiplication of the displayed integer matrices gives
\[
 M_AM_B=(17,18,19,19,20,20,20)^{\mathsf T}(0,2,0,3,0,0,0).
 \tag{E23}
\]
This supplies the universal finite-source input to the supplied
composition proofs, alongside their exact finite and limiting attainers.

\section{The retained empty radix and its actual kernel}

Retaining the two \(A\)-positions of \(ABA\) gives
\(\{1,3,100,300\}\).  The distinct local values are
\([0,16]^2\), and the two evaluations at issue are
\[
 e_{x,y}\longmapsto e_{x+100y},\qquad
 e_{x,y}\longmapsto e_{x+10y}.
 \tag{E24}
\]
The first is a bijection onto a \(289\)-element basis: if its values
coincide, \(100(y-y')=x'-x\in[-16,16]\), forcing both differences
to vanish.  The second has image \([0,176]\), because the intervals
\([10y,10y+16]\), \(0\leq y\leq16\), overlap and have precisely these
endpoints.  It is therefore an onto map from a free rank-\(289\) group
to a free rank-\(177\) group, with free kernel rank \(112\).
An explicit kernel basis chooses for each value \(z\) one pair
\((x_z,y_z)\) evaluating to \(z\) and uses
\(e_{x,y}-e_{x_z,y_z}\) for every other pair in that fibre.
Within each fibre these vectors are independent and span its
coefficient-sum-zero group, proving the rank calculation integrally.

The first evaluation identifies its basis with the original retained
image, so composing its inverse with the second evaluation gives the
exact contraction map between the two value modules.  Formula (E6)
then gives its corresponding additional receiving kernel.  The mass of
each local value pair remains the product of the two original local
word multiplicities.  Applying (E10)--(E12) to the fibres of (E24)
supplies its original minimum-norm quotient, with no unit-mass
substitution.

\section{Accepted return and next arithmetic scope}

The preserved paper proves sharp ternary residue prices
\(3,5,13,23\) for ranks \(1,2,3,4\) of individually modular-five- or
modular-six-free canonical blocks, at unrestricted height and radix.
Its classical addition-theorem input is Martin Kneser's stabilizer
theorem, in the stated form attributed to Matt DeVos's primary proof
\cite{ep:devos}.  The remaining \(490\) exceptional source blocks have
explicit original subset-mask witnesses, replayed in this intake.
The preserved composition paper also gives every fixed-composition
minimum and the full limiting frequency spectrum in the actual
radix-ten dictionary.  All supplied credit to the original
anonymous/deleted contributor and to \texttt{sneed-and-feed}'s distinct
Lean extension remains in the original paper.

The present note proves the exact finite cochain realization (E6),
the integral-section determinant loss (E11), the source-change
boundary and composition law (E14)--(E15), and the universal original
matrix return (E19)--(E23).  These are usable instances of the retained
Split-Zero supported quotient, with their full metrics specified.
The larger-rank modular arithmetic price remains an uncomputed part
of this EP817 continuation.  The original theta source and its signed
analytic allocations receive the algebraic construction (E6)--(E16);
no matrix identifying those analytic sources with the finite residue
source has been supplied or proved in this intake.  Therefore no
analytic allocation coefficient or statement about zeta zeros is
entered in the accepted-result index from this package.



% END PRESERVED INPUT 10_EXACT_SUPPORTED_METRIC_RETURN.tex
\endgroup

\clearpage
\begingroup


\setcounter{equation}{0}
\renewcommand{\theequation}{EIQ.\arabic{equation}}
\section*{Original complete equilibrium provider}
\addcontentsline{toc}{section}{Original complete equilibrium provider}
This is the preserved original EIQ/GEL provider used in the arithmetic value calculation; its exact input hash is recorded in the insertion manifest.
% BEGIN PRESERVED INPUT 11_EIQ_GEL_COMPLETE.tex SHA256 b5eee56d9bac3572fe5eb8582323d6e393a76977e19e32d5c49a8f6563cab838
% END RX INLINE 2

\paragraph{Exact dictionary for the equilibrium provider.}
The elliptic modulus $k$ used locally in EIQ is the $\kappa$ of GEL,
and EIQ's local $m=(u^2+v^2)/2$ is an endpoint midpoint. They are
distinctly indexed coordinates from the original packet
$k_{\rm packet}=4l+1$ and multiplicity $m_{\rm packet}$.
GEL3 gives their exact parameter map
$\alpha=a/n$, $\beta=\pi q/(2n)$, tending to $(2,\pi)$ on
the original sequences. GEL's constant $C_\Gamma=\Gamma(1/4)^2/(2\pi)$
in its source bound retains that definition; the leading coefficient
is the separately defined $\mathcal C_\Gamma$ in GEL18.

% BEGIN RX INLINE 3; SHA256 a0a0346a19d0e035a186590dbd69881cdf7a30e8fef602d47bba0932f175e731
% Independent proof. No external one-cut or equilibrium theorem is used.
% May be included as a section in the cumulative paper.
\section{Exact equilibrium for the square-root logarithmic potential}
\label{eiq:sec:independent-sqrt-log-equilibrium}

Throughout this section $\alpha>0$ and $\beta>0$.  On $(0,\infty)$ set
\begin{equation}\tag{EIQ1}
 V_{\alpha,\beta}(x)=\beta\sqrt{x}-\alpha\log x,
 \qquad
 \mathcal E_{\alpha,\beta}(\mu)
 =\int V_{\alpha,\beta}(x)\,d\mu(x)
       -\iint\log|x-y|\,d\mu(x)d\mu(y).
\end{equation}
The measures considered initially are probability measures for which
$\int V_{\alpha,\beta}\,d\mu$ is finite.  Because this potential is bounded
below and tends to $+\infty$ at both ends, this implies
$\int(\sqrt{x}+|\log x|)\,d\mu<\infty$.
The positive part of $\log|x-y|$ is then integrable, by
$\log^+|x-y|\leq\log(1+x)+\log(1+y)$.
Thus the energy is well defined with value in $\mathbb R\cup\{+\infty\}$.
All comparisons below also apply to the usual extended energy defined by
integrating the kernel
$\tfrac12V_{\alpha,\beta}(x)+\tfrac12V_{\alpha,\beta}(y)-\log|x-y|$:
the kernel is bounded below, and finite energy for this kernel implies the
same moment conditions.  To see the latter assertion without cancellation,
put $g(x)=\tfrac12V_{\alpha,\beta}(x)-\log(1+x)$.
The kernel is at least $g(x)+g(y)$; $g$ is bounded below and controls positive
constant multiples of both $\sqrt{x}$ near infinity and $|\log x|$ near zero.

\subsection{Existence and uniqueness of the two endpoints}
For $0<k<1$ define the complete elliptic integrals by their actual real
integrals
\begin{equation}\tag{EIQ2}
 K(k)=\int_0^{\pi/2}\frac{d\theta}
            {\sqrt{1-k^2\sin^2\theta}},\qquad
 E(k)=\int_0^{\pi/2}\sqrt{1-k^2\sin^2\theta}\,d\theta.
\end{equation}
There is exactly one $k\in(0,1)$ satisfying
\begin{equation}\tag{EIQ3}
 \frac{\sqrt{1-k^2}K(k)}{E(k)}=\frac{\alpha}{\alpha+2}.
\end{equation}
For this $k$ put
\begin{equation}\tag{EIQ4}
 r=\sqrt{1-k^2},\qquad
 v=\frac{(\alpha+2)\pi}{\beta E(k)},\qquad u=rv,
 \qquad a=u^2,\quad b=v^2.
\end{equation}
In particular $0<a<b<\infty$, and the endpoint equations, with all
constants retained, are
\begin{equation}\tag{EIQ5}
 uK(k)=\frac{\alpha\pi}{\beta},\qquad
 vE(k)=\frac{(\alpha+2)\pi}{\beta},\qquad
 k^2=1-\frac{u^2}{v^2}.
\end{equation}

Here is a proof of the assertion of uniqueness. Differentiation of the
integrals in (EIQ2), and integration of the derivative of
$\sin\theta\cos\theta/\sqrt{1-k^2\sin^2\theta}$ between its two zero
endpoint values, give
\[
 E'(k)=\frac{E(k)-K(k)}{k},\qquad
 K'(k)=\frac{E(k)}{k(1-k^2)}-\frac{K(k)}{k}.
\]
For completeness the second identity is equivalent, after differentiation
under the integral sign, to
\[
 \int_0^{\pi/2}
 \left\{\frac{k^2\sin^2\theta}{(1-k^2\sin^2\theta)^{3/2}}
 -\frac{1-k^2\sin^2\theta}{(1-k^2)\sqrt{1-k^2\sin^2\theta}}
 +\frac1{\sqrt{1-k^2\sin^2\theta}}\right\}d\theta=0;
\]
the expression in braces is
$-k^2(1-k^2)^{-1}\,d(\sin\theta\cos\theta/
\sqrt{1-k^2\sin^2\theta})/d\theta$.
With $t=E(k)/K(k)$, the integral definitions show strictly
$r^2<t<1$.  Consequently, if $f(k)=rK(k)/E(k)$, then
\begin{equation}\tag{EIQ6}
 \frac{kf'(k)}{f(k)}
 =\frac{t}{r^2}+\frac1t-1-\frac1{r^2}
 =\frac{(t-1)(t-r^2)}{r^2t}<0.
\end{equation}
At $k=0$ the continuous limiting value is $f(0)=1$.
As $k\uparrow1$, the integrand of $rK(k)$ tends to zero at every
$\theta<\pi/2$ and is bounded by $1$, while $E(k)\to1$ by dominated
convergence.  Hence $f(k)\to0$.  This proves (EIQ3), including existence.
It also proves smooth dependence of $k,u,v,a,b$ on $(\alpha,\beta)$ in
$(0,\infty)^2$, by the nonzero derivative in (EIQ6).

Let $m=(a+b)/2$, $R_s=\sqrt{(a+s)(b+s)}$ for $s\geq0$, and
$R_0=uv$.  Substitution
$t=a\cos^2\theta+b\sin^2\theta$ gives, directly from (EIQ5),
\begin{equation}\tag{EIQ7}
 \int_a^b\frac{V'_{\alpha,\beta}(t)}
                  {\sqrt{(b-t)(t-a)}}\,dt=0,
 \qquad
 \int_a^b\frac{tV'_{\alpha,\beta}(t)}
                  {\sqrt{(b-t)(t-a)}}\,dt=2\pi.
\end{equation}
Indeed the three integrals with numerators $t^{-1/2},t^{-1},t^{1/2}$
are respectively $2K(k)/v$, $\pi/(uv)$, and $2vE(k)$, and
$V'(t)=\beta/(2\sqrt t)-\alpha/t$.

\subsection{A positive density and its exact Cauchy transform}
The elementary Stieltjes integral, obtained by $s=xz^2$, is
\begin{equation}\tag{EIQ8}
 \int_0^\infty\frac{ds}{\sqrt s(x+s)}=\frac\pi{\sqrt x}
 \quad(x>0).
\end{equation}
Using this identity in (EIQ7) and applying Tonelli to its positive
constituents gives
\begin{equation}\tag{EIQ9}
 \frac\alpha{R_0}=
 \frac\beta{2\pi}\int_0^\infty\frac{ds}{\sqrt s R_s},
 \qquad
 \int_0^\infty\frac{1-s/R_s}{\sqrt s}\,ds=2vE(k).
\end{equation}
Both integrals converge; their integrands are $O(s^{-1/2})$ at zero
and $O(s^{-3/2})$ at infinity.  The second formula also follows by
averaging $\sqrt t=\pi^{-1}\int_0^\infty
s^{-1/2}t/(t+s)\,ds$ against the arcsine probability measure on $[a,b]$.

Define for every $x>0$
\begin{align}\tag{EIQ10}
 h(x)
 &=\frac\alpha{xR_0}
   -\frac\beta{2\pi}
       \int_0^\infty\frac{ds}{\sqrt s(x+s)R_s} \\
 &=\frac\beta{2\pi x}
       \int_0^\infty\frac{\sqrt s}{(x+s)R_s}\,ds>0.\notag
\end{align}
The second equality is exactly (EIQ9), with no change of the potential.
All integrals in (EIQ10) converge absolutely and locally uniformly in
$x>0$.  The candidate density on the original positive half-line is
\begin{equation}\tag{EIQ11}
 \rho_{\alpha,\beta}(x)=
 \begin{cases}
 \displaystyle
 \frac{\sqrt{(b-x)(x-a)}}{2\pi}\,h(x),&a<x<b,\\
 0,&x\notin(a,b).
 \end{cases}
\end{equation}
It is strictly positive on $(a,b)$ and vanishes continuously with the
square-root factor at each endpoint.  Positivity here follows from the
last integral in (EIQ10), not from a one-interval ansatz.

The following elementary transform computes its mass and resolvent.
Choose $R(z)=\sqrt{(z-a)(z-b)}$ on $\mathbb C\setminus[a,b]$ with
$R(z)/z\to1$ at infinity. Thus $R(x)>0$ for $x>b$,
$R(x)<0$ for $x<a$, and $R(x+i0)=i\sqrt{(b-x)(x-a)}$ on $(a,b)$.
Direct substitution $t=m+(b-a)\cos\theta/2$, followed by
$w=\tan(\theta/2)$, proves
\[
 \frac1\pi\int_a^b\frac{dt}{(z-t)\sqrt{(b-t)(t-a)}}=\frac1{R(z)}.
\]
For real $z>b$, the substitution reduces the integral to
$\pi^{-1}\int_0^\infty2\,dw/[(z-b)+(z-a)w^2]$;
the identity elsewhere follows by analytic continuation.
Polynomial division then gives
\begin{equation}\tag{EIQ12}
 \frac1\pi\int_a^b\frac{\sqrt{(b-t)(t-a)}}{z-t}\,dt
     =z-m-R(z),\qquad
 \frac1\pi\int_a^b\frac{\sqrt{(b-t)(t-a)}}{t+s}\,dt
     =m+s-R_s.
\end{equation}
For example the first identity follows from
$(b-t)(t-a)=-(t-z)^2+2(m-z)(t-z)-(z-a)(z-b)$,
the preceding arcsine transform, and the arcsine first moment $m$.
Partial fractions in these two identities yield
\begin{equation}\tag{EIQ13}
 H_s(z):=\frac1\pi\int_a^b
       \frac{\sqrt{(b-t)(t-a)}}{(t+s)(z-t)}\,dt
       =1-\frac{R_s+R(z)}{z+s}.
\end{equation}
The formula at $z=-s$ is interpreted by its removable limit.

The mass of (EIQ11), computed using (EIQ12), is
\begin{align}\tag{EIQ14}
 \int_a^b\rho_{\alpha,\beta}(t)\,dt
 &=\frac\alpha{2R_0}(m-R_0)
 -\frac\beta{4\pi}\int_0^\infty
       \frac{(m+s)/R_s-1}{\sqrt s}\,ds\\
 &=-\frac\alpha2+
   \frac\beta{4\pi}\int_0^\infty\frac{1-s/R_s}{\sqrt s}\,ds
 =-\frac\alpha2+\frac{\beta vE(k)}{2\pi}=1.\notag
\end{align}
In particular this is a probability density.

Define its Cauchy transform with the explicit sign convention
$G(z)=\int_a^b\rho_{\alpha,\beta}(t)/(z-t)\,dt$.
For $z$ outside both $[a,b]$ and $(-\infty,0]$, formulas (EIQ10) and
(EIQ13) give
\begin{align}\tag{EIQ15}
 G(z)
 &=\frac\alpha{2R_0}H_0(z)
       -\frac\beta{4\pi}
          \int_0^\infty\frac{H_s(z)}{\sqrt sR_s}\,ds\\
 &=-\frac\alpha{2z}+\frac\beta{4\sqrt z}
       -\frac{R(z)h(z)}2
 =\frac{V'_{\alpha,\beta}(z)-R(z)h(z)}2.\notag
\end{align}
Here the positive square root is continued with its negative-real-axis
cut, and $h(z)$ is the first integral formula of (EIQ10).  Every integral
exchange is absolutely convergent, uniformly on compact subsets of this
domain: near zero the auxiliary integrand is $O(s^{-1/2})$, and near
infinity the original integral defining $H_s$ is $O(s^{-1})$.
The constant terms on the second line cancel by (EIQ9), and (EIQ8)
supplies the term $\beta/(4\sqrt z)$.

\subsection{The Euler equality, its strict exterior inequality, and uniqueness}
Set
\begin{equation}\tag{EIQ16}
 U(x)=V_{\alpha,\beta}(x)
             -2\int_a^b\log|x-t|\rho_{\alpha,\beta}(t)\,dt
 \quad(x>0).
\end{equation}
This function is continuous on $(0,\infty)$.  The integral is continuous
because the density is bounded, compactly supported, and the integral
of $|\log|x-t||$ over an interval of length $\varepsilon$ tends to zero
uniformly in its centre.  In the interior of $[a,b]$, differentiation of
the logarithmic potential is the principal-value integral
$\operatorname{PV}\int\rho(t)/(x-t)\,dt$.
This follows by subtracting $\rho(x)$ on a symmetric interval about
$x$; the remaining numerator is $O(|t-x|)$ because the density is
continuously differentiable locally in the interior.  The real part of
(EIQ15) on the upper edge of the interval is precisely this principal
value.  Indeed subtracting $\rho(x)$ in
$\int\rho(t)/(x+i\varepsilon-t)\,dt$ gives the same limit, while the
constant term has real part equal to its elementary principal value.
Since $R(x+i0)$ is purely imaginary and $h(x)$ is real, (EIQ15) proves
$U'(x)=0$ on $(a,b)$.  There is therefore a real constant $\ell$ such that
\begin{equation}\tag{EIQ17}
 U(x)=\ell\quad(a\leq x\leq b).
\end{equation}
For positive $x$ outside this interval, (EIQ15) instead gives the exact
real derivative
\begin{equation}\tag{EIQ18}
 U'(x)=R(x)h(x)
 \begin{cases}<0,&0<x<a,\\>0,&x>b.\end{cases}
\end{equation}
Combining (EIQ17), continuity, and these signs proves
\begin{equation}\tag{EIQ19}
 U(x)>\ell\quad\hbox{for }x\in(0,a)\cup(b,\infty).
\end{equation}
This verifies the entire exterior inequality on the original domain,
including the interval next to zero.

Here is a complete positivity argument for the energy of the difference
of two measures.  If $\nu$ is a finite signed real measure of total mass
zero for which $|\log|x-y||$ is integrable against
$|\nu|\otimes|\nu|$, the scalar identity
\begin{equation}\tag{EIQ20}
 -\log d=\frac12\int_0^\infty
           \frac{e^{-td^2}-e^{-t}}t\,dt\quad(d>0)
\end{equation}
follows by differentiating in $d$ and evaluating at $d=1$.
One half of the integral over any finite subinterval of $(0,\infty)$
has the same sign as $-\log d$ and absolute value at most $|\log d|$.
Consequently dominated convergence and the zero total mass give
\begin{equation}\tag{EIQ21}
 -\iint\log|x-y|\,d\nu(x)d\nu(y)
 =\frac12\int_0^\infty\frac1t
       \iint e^{-t(x-y)^2}\,d\nu(x)d\nu(y)\,dt\geq0.
\end{equation}
The nonnegative sign is proved, rather than assumed, by the Gaussian
Fourier identity
\begin{equation}\tag{EIQ22}
 \iint e^{-t(x-y)^2}\,d\nu(x)d\nu(y)
 =\frac1{\sqrt{4\pi t}}\int_{\mathbb R}
          e^{-\xi^2/(4t)}|\widehat\nu(\xi)|^2\,d\xi,
 \qquad \widehat\nu(\xi)=\int e^{i\xi x}\,d\nu(x).
\end{equation}
The scalar Gaussian Fourier identity follows by completing the square
in the Gaussian integral (or by differentiating its Fourier transform
and its value at zero); its integration against the finite measures is
justified by absolute integrability.  Equality in (EIQ21) forces
$\widehat\nu=0$: the right side of (EIQ22) is strictly positive for every
$t>0$ if this continuous Fourier transform has a nonzero value.
Finally $\widehat\nu=0$ implies $\nu=0$ without an extra uniqueness
assumption. Convolve $\nu$ with a centred Gaussian approximate identity.
Its value at every point is zero by that same Fourier integral.
For every compactly supported continuous test function, its convolution
with these Gaussians converges uniformly to the test function; integration
against $\nu$ then proves that $\nu$ annihilates every such function.

Let $\mu$ be any probability measure of finite energy in (EIQ1), and
write $\mu_*(dx)=\rho_{\alpha,\beta}(x)\,dx$ and $\nu=\mu-\mu_*$.
The integrability needed for (EIQ21) holds.  The $\mu\otimes\mu$ term
has it by finite energy and the already established logarithmic moment;
the $\mu_*\otimes\mu_*$ term has it by bounded compact support; and the
cross term has it because the negative logarithmic singularity integrated
against the bounded density is uniformly bounded, while the positive
part is controlled by $\log(1+x)+\log(1+y)$.
Expanding the energy therefore yields exactly
\begin{equation}\tag{EIQ23}
 \mathcal E_{\alpha,\beta}(\mu)-
       \mathcal E_{\alpha,\beta}(\mu_*)
 =\int(U(x)-\ell)\,d\mu(x)
       -\iint\log|x-y|\,d\nu(x)d\nu(y)\geq0.
\end{equation}
The first term is nonnegative by (EIQ17)--(EIQ19), and the second is
strictly positive unless $\mu=\mu_*$. Infinite energy competitors cannot
improve this finite value. Thus (EIQ11) is the unique global minimizing
probability measure, with exactly the support $[a,b]$.

\subsection{An explicit endpoint integral for the logarithmic moment}
Retain $m,R_s,u,v$ as defined above, and define
\begin{equation}\tag{EIQ24}
 C=\frac{(u+v)^2}4,\qquad
 \eta=\frac{v-u}{v+u},\qquad
 \zeta_s=\frac{b-a}{2(m+s+R_s)}\quad(s\geq0).
\end{equation}
Thus $0<\zeta_s\leq\eta<1$ and $\zeta_0=\eta$.
For every $s\geq0$ put
\begin{equation}\tag{EIQ25}
 L_s=(m+s-R_s)\log C-(m-R_0)
                  -2R_s\log(1-\eta\zeta_s).
\end{equation}
Then the logarithmic moment of the equilibrium measure is exactly
\begin{equation}\tag{EIQ26}
 \mathcal L(\alpha,\beta):=\int_a^b\log x\rho_{\alpha,\beta}(x)\,dx
 =\frac\beta{4\pi}\int_0^\infty
                    \frac{L_0-L_s}{\sqrt s R_s}\,ds.
\end{equation}
This is a convergent one-dimensional endpoint integral of explicitly
given real elementary functions; its only endpoints are those specified
uniquely in (EIQ3)--(EIQ5).

Here is a direct derivation, including the constants in (EIQ25).
Set $h_0=(b-a)/2$ and substitute $t=m+h_0\cos\theta$ in arcsine
integrals.  The identities
\[
 t=C(1+2\eta\cos\theta+\eta^2),\qquad
 t+s=\frac{m+s+R_s}2
                (1+2\zeta_s\cos\theta+\zeta_s^2)
\]
give the absolutely convergent Fourier series
\[
 \log t=\log C+
       2\sum_{j\geq1}\frac{(-1)^{j+1}\eta^j}j\cos(j\theta),
 \qquad
 \frac1{t+s}=\frac1{R_s}
       \left(1+2\sum_{j\geq1}(-\zeta_s)^j\cos(j\theta)\right).
\]
Orthogonality on $[0,\pi]$ gives
\begin{align}\tag{EIQ27}
 \frac1\pi\int_a^b\frac{\log t}{\sqrt{(b-t)(t-a)}}\,dt
 &=\log C,\\
 \frac1\pi\int_a^b\frac{t\log t}{\sqrt{(b-t)(t-a)}}\,dt
 &=m\log C+(m-R_0),\notag\\
 \frac1\pi\int_a^b\frac{\log t}{(t+s)\sqrt{(b-t)(t-a)}}\,dt
 &=\frac{\log C+2\log(1-\eta\zeta_s)}{R_s}.\notag
\end{align}
The second identity uses $h_0\eta=m-R_0$.
Polynomial division
\[
 \frac{(b-t)(t-a)}{t+s}
       =-t+2m+s-\frac{R_s^2}{t+s}
\]
now proves, exactly,
\begin{equation}\tag{EIQ28}
 L_s=\frac1\pi\int_a^b
               \frac{\log t\sqrt{(b-t)(t-a)}}{t+s}\,dt.
\end{equation}
Using the positive integral expression for $h$ in (EIQ10), and the
identity $1/[t(t+s)]=(1/t-1/(t+s))/s$, yields (EIQ26).
The exchanges are absolutely convergent: $\log t$ is bounded on
$[a,b]$, $L_0-L_s=O(s)$ at zero by (EIQ28), and $L_s=O(s^{-1})$
at infinity.  Thus the integrand in (EIQ26) is $O(s^{1/2})$ at zero
and $O(s^{-3/2})$ at infinity.

\subsection{Parameter continuity and the variational derivative}
The density and its logarithmic moment depend smoothly on
$(\alpha,\beta)\in(0,\infty)^2$ in the following precise sense.
On a compact parameter neighbourhood, (EIQ3)--(EIQ6) put $a,b$ in a
fixed compact subset of $(0,\infty)$, with $b-a$ bounded below by a
positive number.  The integral for $h$ and each of its parameter and
$x$ derivatives converge uniformly for $x$ in this fixed compact set:
their absolute integrands are bounded by constants times $s^{1/2}$
near zero and $s^{-3/2}$ near infinity. Parameter derivatives of the
endpoints are bounded on a smaller compact neighbourhood.
Under the fixed-interval coordinate
$x=a+(b-a)y$, $0\leq y\leq1$, the measure is exactly
\begin{equation}\tag{EIQ29}
 \rho_{\alpha,\beta}(x)\,dx
 =\frac{(b-a)^2}{2\pi}\sqrt{y(1-y)}
                  h\bigl(a+(b-a)y\bigr)\,dy.
\end{equation}
All parameter derivatives of its coefficient are bounded by a constant
times $\sqrt{y(1-y)}$. The same is true after multiplication by
$\log(a+(b-a)y)$ or by $\sqrt{a+(b-a)y}$.
Differentiation under this integral proves smoothness, in particular
continuity, of both moments on the entire positive parameter quadrant.

Define $E_{\min}(\alpha,\beta)=\mathcal E_{\alpha,\beta}(\mu_{\alpha,\beta})$,
where the measure is the unique minimizer proved above.
Testing the minimizer for a parameter pair against the neighbouring
potential, and then reversing the two tests, gives for $h>0$
\[
 -\mathcal L(\alpha+h,\beta)
 \leq\frac{E_{\min}(\alpha+h,\beta)-E_{\min}(\alpha,\beta)}h
 \leq-\mathcal L(\alpha,\beta).
\]
The corresponding inequalities for $h<0$ and the proved continuity
therefore give
\begin{equation}\tag{EIQ30}
 \partial_\alpha E_{\min}(\alpha,\beta)=-\mathcal L(\alpha,\beta),\qquad
 \partial_\beta E_{\min}(\alpha,\beta)
        =\int\sqrt{x}\,d\mu_{\alpha,\beta}(x).
\end{equation}
No unproved differentiability of a minimizing selection is being used.
For the opposite convention
$F(\alpha,\beta)=\sup_\mu\{\iint\log|x-y|\,d\mu(x)d\mu(y)
-\int V_{\alpha,\beta}\,d\mu\}=-E_{\min}(\alpha,\beta)$,
the exact derivative is therefore
$\partial_\alpha F(\alpha,\beta)=\mathcal L(\alpha,\beta)$.

The second moment in (EIQ30) has an exact value.  Push the minimizing
measure forward by the explicit dilation $x\mapsto c^2x$, $c>0$.
Writing $M=\int\sqrt{x}\,d\mu_{\alpha,\beta}(x)$, its energy minus
the energy at $c=1$ is precisely
$\beta(c-1)M-2(\alpha+1)\log c$.
Its derivative at the minimizing value $c=1$ is zero. Hence
\begin{equation}\tag{EIQ31}
 \int\sqrt{x}\,d\mu_{\alpha,\beta}(x)
       =\frac{2(\alpha+1)}\beta.
\end{equation}
The same explicit dilation, or (EIQ3)--(EIQ5), proves that the map
$x\mapsto(\pi/\beta)^2x$ pushes $\mu_{\alpha,\pi}$ to
$\mu_{\alpha,\beta}$, including the density and its two endpoints.
Consequently
\begin{equation}\tag{EIQ32}
 \mathcal L(\alpha,\beta)
       =\mathcal L(\alpha,\pi)+2\log(\pi/\beta).
\end{equation}

For the particular original limiting potential
$V(x)=\pi\sqrt{x}-2\log x$, the substitutions are exactly
$\alpha=2$, $\beta=\pi$. Equations (EIQ5), (EIQ10), and (EIQ26) become
\begin{equation}\tag{EIQ33}
 uK(k)=2,\qquad vE(k)=4,\qquad
 h(x)=\frac1{2x}\int_0^\infty\frac{\sqrt s}{(x+s)R_s}\,ds,
 \qquad
 \mathcal L(2,\pi)=\frac14\int_0^\infty
                         \frac{L_0-L_s}{\sqrt sR_s}\,ds.
\end{equation}
In particular the preceding proof includes positivity, unit mass,
the full Euler inequality, the unique global minimizer, and the
parameter continuity needed at these actual parameters.

% END RX INLINE 3


% BEGIN RX INLINE 4; SHA256 2014e7fefdcd963b70b00fb89d597e24c0826582a2541dd1ff44d63051979fc1
% Original-source large-degree calculation, 14 September 2026.
% Input fragment. Equilibrium provider:
% independent/EXACT_SQRT_LOG_EQUILIBRIUM.tex (EIQ1--33).
\section{The exact Gamma ensemble and the leading growing return}
\label{eiq:sec:GEL}

\paragraph{GEL1. Source, exponents, and the literal coordinate map.}
The source throughout is
\[
 d\sigma(y)=\frac{|\Gamma(1/4+iy/2)|^2}{2\pi}\,dy,
 \qquad \sigma(\mathbb R)=\sqrt{2\pi},\qquad
 d\nu(t)=\frac{|\Gamma(1/4+i\sqrt t/2)|^2}{2\pi\sqrt t}\,dt.
 \tag{GEL1}
\]
Here \(\nu\) is the full pushforward under \(y\mapsto y^2\); both
half-axes are present in its density. For integers \(a,n\geq0\), retain
\[
 H_n^{(a)}=\det\left[\int_0^\infty t^{a+i+j}\,d\nu(t)\right]_{i,j<n}.
\]
The definition extends to real \(a>-1/2\); all its integrals converge.
For \(n\geq1\), expansion of the two Vandermonde determinants followed
by integration term by term gives the exact identity
\[
 H_n^{(a)}=\frac1{n!}\int_{(0,\infty)^n}
       \prod_{i<j}(t_i-t_j)^2\prod_{i=1}^n t_i^a\,d\nu(t_i).
 \tag{GEL2}
\]
Indeed expansion gives \(n!\) copies of the Gram determinant after one
of the two permutation indices is fixed. Absolute integrability follows
by expanding the squared polynomial and using the existing moments.

Fix a positive real coordinate scale \(q\), independently of \(a,n\),
and make the literal map \(t_i=q^2x_i\). Introduce the reference
probability measure
\[
 d\omega(x)=\tfrac12x^{-1/2}e^{-\sqrt x}\,dx,
 \qquad B_q(x)=2q e^{\sqrt x}\sigma(q\sqrt x).
\]
Its mass is one by \(x=u^2\). This reference measure does not replace
or rescale the source: the exact Radon--Nikodym identity is
\(d\nu(q^2x)=B_q(x)d\omega(x)\), including its whole mass. Consequently
\[
 \boxed{H_n^{(a)}=q^{2an+2n(n-1)}J_{n,q}(a),\quad
 J_{n,q}(a)=\frac1{n!}\int\prod_{i<j}(x_i-x_j)^2
                         \prod_i x_i^aB_q(x_i)\,d\omega(x_i).}
 \tag{GEL3}
\]
The powers in this identity come separately from the \(n\) factors
\(t_i^a\) and the \(n(n-1)/2\) squared differences. No determinant,
source mass, or Jacobian has been omitted.

\paragraph{GEL2. An elementary global Gamma bound with the exact exponential.}
Put \(\lambda=1/4\), \(b\geq0\), and
\(C_\Gamma=\Gamma(1/4)^2/(2\pi)\). Euler's finite limit gives
\[
 \frac{|\Gamma(\lambda+ib)|^2}{\Gamma(\lambda)^2}
 =\prod_{j=0}^\infty\left(1+\frac{b^2}{(j+\lambda)^2}\right)^{-1}.
\]
For clarity, the finite limit follows by applying the beta integral to
\(\int_0^1t^{z-1}(1-t)^Ndt=N!/[z(z+1)\cdots(z+N)]\), replacing
\(t\) by \(u/N\), and using dominated convergence for
\(\operatorname{Re}z>0\). Taking squared moduli and dividing by the
same limit at \(z=\lambda\) proves the displayed product. Its logarithm
converges absolutely because its summand is \(O((j+1)^{-2})\).

The decreasing function \(v\mapsto\log(1+b^2/v^2)\) gives, with
\(S=\sum_{j\geq0}\log(1+b^2/(j+\lambda)^2)\),
\[
 I\leq S\leq I+\log(1+b^2/\lambda^2),\qquad
 I=\int_\lambda^\infty\log(1+b^2/v^2)\,dv
  =\pi b-\lambda\log(1+b^2/\lambda^2)
                      -2b\arctan(\lambda/b).
\]
The integral is evaluated by differentiating
\(v\log(1+b^2/v^2)+2b\arctan(v/b)\); at \(b=0\) it is zero
by continuity. Since \(0\leq2b\arctan(\lambda/b)\leq2\lambda\),
substitution of \(b=|y|/2\) proves the global, finite inequalities
\[
 \boxed{C_\Gamma(1+4y^2)^{-3/4}e^{-\pi|y|/2}
 \leq\sigma(y)\leq
 C_\Gamma e^{1/2}(1+4y^2)^{1/4}e^{-\pi|y|/2}.}
 \tag{GEL4}
\]
Both inequalities include \(y=0\). In particular, for
\(C_0=\Gamma(1/4)^2/\pi\),
\[
 B_q(x)\leq C_0q e^{1/2}(1+2q)^{1/2}
            e^{-(\pi q/2-3/2)\sqrt x}\quad(x>0).
 \tag{GEL5}
\]
Here \((1+4q^2x)^{1/4}\leq(1+2q\sqrt x)^{1/2}
\leq(1+2q)^{1/2}(1+\sqrt x)^{1/2}
\leq(1+2q)^{1/2}e^{\sqrt x/2}\). On each fixed compact interval
\([A,B]\subset(0,\infty)\), the lower bound is
\[
 B_q(x)\geq C_0q(1+4q^2B)^{-3/4}e^{-\pi q\sqrt x/2}.
 \tag{GEL6}
\]
Thus the exact exponential, and the size and position of the remaining
source factor, have been proved directly from the original Gamma function.

\paragraph{GEL3. The energy object and its original parameter domain.}
For \(\alpha,\beta>0\) put
\[
 V_{\alpha,\beta}(x)=\beta\sqrt x-\alpha\log x,
 \quad K_{\alpha,\beta}(x,y)=\log|x-y|
           -\tfrac12V_{\alpha,\beta}(x)-\tfrac12V_{\alpha,\beta}(y),
\]
\[
 \mathcal F(\alpha,\beta)=\sup_{\mu\in\mathcal P((0,\infty))}
                   \iint K_{\alpha,\beta}(x,y)\,d\mu(x)d\mu(y).
 \tag{GEL7}
\]
The kernel is bounded above, so its integral has an unambiguous value
in \(\mathbb R\cup\{-\infty\}\) for every probability in this domain.
Whenever its value is finite it equals
\(\iint\log|x-y|\,d\mu^2-\int V_{\alpha,\beta}\,d\mu\).
Indeed \(K\leq r(x)+r(y)\), with
\(r=\log(1+x)-V_{\alpha,\beta}/2\) bounded above and tending to
\(-\infty\) at both endpoints; finite kernel integral therefore forces
the \(\sqrt x\) and \(|\log x|\) moments to be finite. The positive
logarithmic part is then integrable, and finite kernel integral forces
the negative part to be integrable as well. This defines the extended
energy without forming an indeterminate difference of infinite quantities.
The equilibrium calculation supplied with this section proves existence,
uniqueness, and the full exterior variational inequality for its density
\(\rho_{\alpha,\beta}\), with support \([u^2,v^2]\). Its exact endpoints
and positive density are
\[
 uK(\kappa)=\frac{\alpha\pi}{\beta},\quad
 vE(\kappa)=\frac{(\alpha+2)\pi}{\beta},\quad
 \kappa^2=1-u^2/v^2,
\]
\[
 \rho_{\alpha,\beta}(x)=
 \frac{\beta\sqrt{(v^2-x)(x-u^2)}}{4\pi^2x}
 \int_0^\infty
 \frac{\sqrt s\,ds}{(x+s)\sqrt{(u^2+s)(v^2+s)}}
       \quad(u^2<x<v^2).
 \tag{GEL8}
\]
The endpoint proof shows continuous parameter dependence, finite entropy
relative to \(\omega\), finite logarithmic energy, and
\[
 \partial_\alpha\mathcal F(\alpha,\beta)
       =\int\log x\,\rho_{\alpha,\beta}(x)\,dx.
 \tag{GEL9}
\]
In particular \(\mathcal F\) and this derivative are continuous on the
positive parameter quadrant. The following ensemble argument proves the
connection to GEL3; it does not presume a random-matrix asymptotic theorem.

\paragraph{GEL4. Complete free-energy limit from a compactified kernel.}
Let \(n\to\infty\), \(q/n\to\theta>0\), and let \(\alpha>0\) be fixed.
Then the actual integral in GEL3 satisfies
\[
 \boxed{\frac1{n^2}\log J_{n,q}(n\alpha)
       \longrightarrow\mathcal F(\alpha,\pi\theta/2).}
 \tag{GEL10}
\]
Here is a proof of both bounds. For the upper bound GEL5 removes the
factor \([C_0q e^{1/2}(1+2q)^{1/2}]^n/n!\), whose logarithm divided by
\(n^2\) tends to zero. The remaining exponent is
\[
 2\sum_{i<j}\log|x_i-x_j|-n\sum_iV_{\alpha,\beta_n}(x_i),
 \qquad \beta_n=\frac{\pi q}{2n}-\frac3{2n}.
\]
For \(n>1\) it is exactly
\(\sum_{i\ne j}K_{Q_n}(x_i,x_j)\), where
\[
 Q_n=\frac n{n-1}V_{\alpha,\beta_n},\qquad
 K_Q(x,y)=\log|x-y|-\tfrac12Q(x)-\tfrac12Q(y).
\]
Write \(\beta=\pi\theta/2\). Choose \(\varepsilon>0\) so that
\(\alpha-\varepsilon>0\) and \(\beta-3\varepsilon>0\).
For all sufficiently large \(n\), the two coefficients of \(Q_n\)
differ from \((\alpha,\beta)\) by at most \(\varepsilon\). It follows
pointwise that
\(Q_n\geq Q_-:=V_{\alpha-\varepsilon,\beta-3\varepsilon}\).
Indeed its difference has the form
\(b\sqrt x-a\log x\), where \(b\geq2\varepsilon\) and
\(0\leq a\leq2\varepsilon\); it is nonnegative for \(x<1\), and for
\(x\geq1\) use \(\log x\leq\sqrt x\).

Compactify \((0,\infty)\) by adjoining its two endpoints.
The function \(K_{Q_-}\), assigned value \(-\infty\) on the diagonal
and at either added endpoint, has continuous truncations
\(K_M=\max\{K_{Q_-},-M\}\). To verify the endpoint assertion uniformly
in the other variable, use
\[
 K_{Q_-}(x,y)\leq r(x)+r(y),\qquad
 r(x)=\log(1+x)-Q_-(x)/2.
\]
The function \(r\) is bounded above and tends to \(-\infty\) at both
endpoints. Continuity at a diagonal point in the interior follows from
\(\log|x-y|\to-\infty\) and local boundedness of \(Q_-\).
For the empirical probability \(L_n=n^{-1}\sum_i\delta_{x_i}\),
\[
 \frac1{n^2}\sum_{i\ne j}K_{Q_-}(x_i,x_j)
 \leq\iint K_M\,dL_n^2+\frac Mn
 \leq s_M+\frac Mn,
 \quad s_M=\max_{\mu}\iint K_M\,d\mu^2.
\]
Since \(\omega\) has mass one, integration preserves this upper bound.
Moreover \(s_M\downarrow\mathcal F(\alpha-\varepsilon,
\beta-3\varepsilon)\). For completeness, choose maximizing probabilities
\(\mu_M\). Compactness gives a convergent subsequence. For every fixed
\(L\), the inequality \(K_M\leq K_L\) when \(M\geq L\), followed by
continuity of the integral of \(K_L\), bounds the subsequential limit
of \(s_M\) by \(\iint K_Ld\mu^2\). Decrease \(L\)'s truncation to the
untruncated kernel. Monotone convergence after subtracting a common upper
bound gives \(\lim s_M\leq\iint K_{Q_-}d\mu^2\leq\mathcal F\).
The reverse inequality follows by testing every \(K_M\) on an arbitrary
probability of finite energy. Probabilities charging an added endpoint
have energy \(-\infty\), so they do not alter the supremum.
First take \(n\to\infty\), then \(M\to\infty\), and finally
\(\varepsilon\downarrow0\). Continuity proved in GEL9 yields the
upper bound in GEL10.

For the lower bound choose exactly
\(d\mu(x)=\rho_{\alpha,\beta}(x)dx\), supported on \([A,B]\), and use
GEL6. Write \(p=d\mu/d\omega\). Restricting the integral to this support
and changing measure to \(\mu^n\), Jensen's inequality gives
\begin{align*}
 \log J_{n,q}(n\alpha)\geq{}&-
 \log n!+n\log\{C_0q(1+4q^2B)^{-3/4}\}\\
 &+n(n-1)\iint\log|x-y|\,d\mu(x)d\mu(y)\\
 &-n^2\int V_{\alpha,\pi q/(2n)}\,d\mu
       -n\int\log p\,d\mu .
\end{align*}
All logarithms on the right are integrable: the equilibrium density is
bounded, has square-root decay at its finite positive endpoints, and
\(\omega\) has a positive smooth density there. The prefactor contributes
\(O(n\log q+n\log n)\), and the entropy contributes \(O(n)\).
Division by \(n^2\) therefore yields exactly the value of GEL7 at its
maximizer. This proves GEL10.

\paragraph{GEL5. Finite positive shifts by a proved convex squeeze.}
For each fixed \(n,q\), differentiation under the integral in GEL3 on
every compact subinterval of \(a>-1/2\) gives
\[
 \frac{d}{da}\log J_{n,q}(a)=
    \mathbb E_a\sum_i\log x_i,\qquad
 \frac{d^2}{da^2}\log J_{n,q}(a)=
    \operatorname{Var}_a\!\left(\sum_i\log x_i\right)\geq0.
 \tag{GEL11}
\]
The expectation is with respect to the entire positive integrand in
GEL3 divided by its integral. Extra logarithmic factors are integrable
at zero because \(a>-1/2\), and at infinity by GEL4. Hence the
functions \(F_{n,q}(z)=n^{-2}\log J_{n,q}(nz)\) are differentiable
and convex. Suppose now
\[
 q/n\longrightarrow2,\qquad a/n\longrightarrow2,\qquad
 l>0,\quad l/n\longrightarrow0.
\]
Set
\[
 \mathcal L=\int_{u^2}^{v^2}\log x\,\rho_{2,\pi}(x)\,dx,
 \qquad uK(\kappa)=2,\quad vE(\kappa)=4.
 \tag{GEL12}
\]
For fixed \(0<\eta<2\), both endpoints of the secant from \(a/n\) to
\((a+l)/n\) lie in \((2-\eta/2,2+\eta/2)\) for large \(n\).
Convexity bounds that secant from below by the secant from
\(2-\eta\) to \(2-\eta/2\), and from above by the secant from
\(2+\eta/2\) to \(2+\eta\). Apply GEL10 at these four fixed positive
arguments, then let \(\eta\downarrow0\) and use GEL9. This proves
\[
 \frac{\log J_{n,q}(a+l)-\log J_{n,q}(a)}{nl}
       \longrightarrow\mathcal L.
\]
Combining this with the exact source factors in GEL3 gives
\[
 \boxed{\log\frac{H_n^{(a+l)}}{H_n^{(a)}}
       =2nl\log q+nl\mathcal L+o(nl).}
 \tag{GEL13}
\]
This proof covers moving starting exponents \(a\); it never differentiates
the error in a free-energy asymptotic. The remainder assertion has the
precise scope \(o(nl)\) along the displayed sequences.

\paragraph{GEL6. Both original matrices and the low moment.}
Return to the actual packet
\[
 k=4l+1,\quad e=1+k(m-1),\quad q=e(k+1)^2=2n,
 \qquad m\geq1\text{ fixed},\qquad l\to\infty.
\]
Retain the exact parity identity from HCT35,
\(Z_a=H_n^{(a)}H_{n+1}^{(a)}(H_n^{(a+1)})^2\).
Use GEL13 at the four literal factors, with sizes \(n,n+1,n,n\)
and starting exponents \(q,q,q+1,q+1\), respectively. Every sequence
satisfies its stated domain. Addition gives
\[
 \boxed{\log\frac{Z_{q+l}}{Z_q}
   =2(2q+1)l\log q+(2q+1)l\mathcal L+o(ql).}
 \tag{GEL14}
\]
Thus the contribution from the two original full matrices has been
evaluated, with all four parity blocks still identifiable.

Write \(\mu_a=\int y^{2a}d\sigma(y)\). GEL4 gives, for integer \(a\geq1\),
\[
 \log\mu_a=2a\log a+O(a).
 \tag{GEL15}
\]
One elementary proof of the upper bound uses
\((1+4y^2)^{1/4}\leq1+2y\) for \(y\geq0\); its integrals are
\((2a)!/(\pi/2)^{2a+1}\) and
\(2(2a+1)!/(\pi/2)^{2a+2}\), multiplied by the exact constant
\(2C_\Gamma e^{1/2}\). The bound \(N!\leq N^N\) suffices.
For the lower bound restrict the positive-axis integral to \([a,a+1]\):
there \((1+4y^2)^{-3/4}\geq5^{-3/4}(a+1)^{-3/2}\),
\(y^{2a}\geq a^{2a}\), and
\(e^{-\pi y/2}\geq e^{-\pi(a+1)/2}\). These prove both sides of GEL15
with constants independent of \(a\).

The real-exponent function \(a\mapsto\log\mu_a\) is convex by the
same variance proof as GEL11. Since \(q\) is even and \(1\leq l\leq q\),
its secant on \([q,q+l]\) lies between the secants on \([q/2,q]\)
and \([2q,4q]\). GEL15 bounds both latter secants in absolute value
by \(O(\log q)\). Therefore
\[
 \left|\log\frac{\mu_{q+l}}{\mu_q}\right|=O(l\log q)=o(ql).
 \tag{GEL16}
\]
The original low moment is precisely
\(\mu_{q+l}/\mu_q=m_{2q+2l}/m_{2q}\), so no index is reinterpreted.

\paragraph{GEL7. The leading value of the unchanged signed centre.}
Keep the source product and its exact row intervals:
\[
 P_k=-\sum_{a=q}^{2q-1}\sum_{j=0}^{2l-1}\log(a+j+1/2)
     -\sum_{a=q+1}^{2q}\sum_{j=0}^{2l-1}\log(a+j+1/2).
\]
The two outer sums each have \(q\) terms. Since \(l/q\to0\), direct
Riemann sums, with \(\log\) Lipschitz on \([1,3]\), give
\[
 P_k=-4ql\log q-4ql\int_1^2\log x\,dx+o(ql)
     =-4ql\log q-4ql(2\log2-1)+o(ql).
 \tag{GEL17}
\]
Explicitly, omitting \((j+1/2)/q\) changes each logarithm by at most
\(2l/q\), hence changes the total by at most \(8l^2\); the ordinary
left and right Riemann sums together have error \(O(l)\).

The full finite product calculation WGP10 strengthens this leading
identity to the following enclosure for the very same \(P_k\):
\[
 \begin{split}
 P_k={}&-4lq\log q-4lq(2\log2-1)-4l^2\log2
              +\frac{8l^3+l}{6q}+\mathcal E_{q,l},\\
 -\left\{\frac l{6q}+\frac{8l^3+l}{12q^3}
                 +\frac{8l^4+2l^2}{3q^2}\right\}
 \leq{}&\mathcal E_{q,l}\leq\frac{l^2}{3q^2}.
 \end{split}
 \tag{GEL17a}
\]
Its proof uses the literal midpoint coordinates in WGP3, the signed
quadrature remainders WGP4--5, and the nonnegative cubic remainder
WGP6--9. The subsequent exact refinement WGP12--14 gives
\[
 \begin{split}
 P_k={}&-4lq\log q-4lq(2\log2-1)-4l^2\log2
 +\frac{8l^3+l}{6q}-\frac{l^4+l^2/4}{q^2}
          +\mathcal E^{(3)}_{q,l},\\
 -\frac l{6q}-\frac{8l^3+l}{12q^3}
 \leq{}&\mathcal E^{(3)}_{q,l}\leq
 \frac{l^2}{3q^2}+\frac{8l^4+2l^2}{6q^4}
                 +\frac{7(192l^5+80l^3-2l)}{1440q^3}.
 \end{split}
 \tag{GEL17b}
\]
These original endpoint contributions are retained when the finer
\(q\)-scale calculation is continued. In particular the term
\(-4l^2\log2\) has order \(q\) at \(m=1\), and the cubic term tends
to \(-1/256\) there; the leading limit does not remove either term.

The original quantity is exactly
\(W_k=P_k-\log(\mu_{q+l}/\mu_q)+\log(Z_{q+l}/Z_q)\).
Substitute GEL14, GEL16, and GEL17. The surviving
\(2l\log q\) is \(o(ql)\), so the leading growing-degree calculation is
\[
 \boxed{\frac{W_k}{ql}\longrightarrow
 \mathcal C_\Gamma:=2\mathcal L-4(2\log2-1),\qquad
 \mathcal L=\int_{u^2}^{v^2}\log x\,\rho_{2,\pi}(x)\,dx,\quad
 uK(\kappa)=2,\ vE(\kappa)=4.}
 \tag{GEL18}
\]
Every measure, exponent, matrix dimension, and source factor entering this
limit has been connected to its original object by GEL1--3 and GEL14.
The exact positive density and unique endpoint equations in GEL8 make
the constant an explicitly defined one-dimensional equilibrium integral.
This leading result has remainder \(o(ql)\). It supplies no claim of
an \(o(q)\) remainder, and hence does not erase the next finite-scale
calculation required by the original arithmetic receiver.

The fully proved original-polynomial recurrence estimate RWB20--27
also applies to this identical \(W_k\). It gives the finite bounds
\(lq/64\leq W_k\leq6lq\), and its narrower finite lower remainder
gives \(\liminf W_k/(lq)\geq4\log(27/(8\pi))\). Combining those
proved inequalities with the now proved existence of GEL18 yields
\[
 \boxed{4\log\frac{27}{8\pi}\leq\mathcal C_\Gamma\leq6,
              \qquad\mathcal C_\Gamma>0.}
 \tag{GEL18a}
\]
The strict sign is analytic: RWB22 proves
\(4\log(27/(8\pi))>104/365>0\) using rational integral bounds for
\(\pi\) and the elementary logarithm inequality. Numerical evaluation
of the equilibrium integral is unnecessary for this sign conclusion.

Finally LET20, whose error is \(O(k)=o(q)\) at fixed original
\(m,\delta,\gamma\), transfers the proved limit without changing its sign:
\[
 \frac{\mathfrak T_k}{ql}\longrightarrow\mathcal C_\Gamma,
 \qquad
 \frac{R_k^0-R_k^\sigma}{ql}\longrightarrow-\mathcal C_\Gamma.
 \tag{GEL19}
\]
The arithmetic identity remains
\(B_{\rm ar}=B_0+\delta_\sigma-\mathfrak T_k\), including its original
\(\delta_\sigma\) allowance. Thus GEL19 evaluates the Gamma summand
and its exact signed map into that identity; it makes no unsupported
growth assertion about either of the other two summands.

% END RX INLINE 4



% END PRESERVED INPUT 11_EIQ_GEL_COMPLETE.tex
\endgroup

\clearpage
\begingroup


\setcounter{equation}{0}
\renewcommand{\theequation}{RTM.\arabic{equation}}
\section*{Original complete tilted-mass and source-envelope providers}
\addcontentsline{toc}{section}{Original complete tilted-mass and source-envelope providers}
% BEGIN PRESERVED INPUT 12_RTM_GTM_original_providers.tex SHA256 788eac1809dcab880cbffa587d50a97db4392f449ba684d586daaffcb5c339f2
\section{The full tilted mass in the original arithmetic tail}
\label{rtm:sec:retained-tilted-mass-tail}

This calculation uses the original complete hypothetical actual quartet,
not a substituted arithmetic density.  Keep
\[
 \rho=\tfrac12+\delta+i\gamma,\quad0<\delta<\tfrac12,\quad\gamma>2,
 \qquad h(s)=\prod_{\lambda\in\{\rho,\bar\rho,1-\rho,1-\bar\rho\}}
 (s-\lambda)^m,\quad m\ge1,
\]
and retain the entire unit $v_h=2\xi/h$ and the measures
\begin{equation}\tag{RTM1}
 \begin{gathered}
 w_h(y)=|v_h(1/2+iy)|^2/(2\pi),\quad m_k=w_h^{*k},\quad
 \alpha=\pi/2,\\ f(y)=e^{\alpha y}w_h(y),\quad
 \mathfrak M=M_\alpha=\int_{\mathbb R}f(y)\,dy,\\
 d\sigma(y)=\sigma(y)\,dy
 =|\Gamma(1/4+iy/2)|^2\,dy/(2\pi),\quad
 \sigma(\mathbb R)=\sqrt{2\pi}.
 \end{gathered}
\end{equation}
The original convolution envelope at three factors, the boundary tilt
bound, and the two-sided Gamma bounds are exactly TW7 at $k=3$ and
BT9--17:
\begin{equation}\tag{RTM2}
 \begin{gathered}
 w_h^{*3}(y)\ge c_h e^{-\alpha|y|}(1+|y|)^{-B},\quad B=42+8m,\\
 0\le f(y)\le C_h^{\rm tilt}(1+|y|)^{-p},\quad p=8m-9/2\ge7/2,\\
 c_\Gamma e^{-\alpha|y|}(1+|y|)^{-1/2}
 \le\sigma(y)\le C_\Gamma e^{-\alpha|y|}(1+|y|)^{-1/2},\\
 c_\Gamma=\sqrt2e^{-7/6},\qquad
 C_\Gamma=\sqrt{2\sqrt5}e^{2/3}.
 \end{gathered}
\end{equation}
Here $c_h$ and $C_h^{\rm tilt}$ are the unchanged explicit constants
from those proved providers.  Their source is still
\[
 w_h(y)=\frac{(y^2+1/4)^2}{\sqrt\pi}
 \frac{|\zeta(1/2+iy)|^2}
 {\bigl[((y-\gamma)^2+\delta^2)((y+\gamma)^2+\delta^2)\bigr]^{2m}}
 \sigma(y).
\]
No positive lower bound for its one-factor zeta numerator is used.

\subsection{Exact moments of the original tilted convolution}

Define the following actual source integrals, preserving their masses:
\begin{equation}\tag{RTM3}
 \begin{gathered}
 J_h=\int |y|f(y)\,dy,
 \qquad \mu_h=\mathfrak M^{-1}\int y f(y)\,dy,\\
 v_h^{\rm tilt}=\mathfrak M^{-1}
       \int(y-\mu_h)^2f(y)\,dy,\\
 0<J_h\le\frac{2C_h^{\rm tilt}}{p-2},\qquad
 0<\mu_h\le J_h/\mathfrak M,\qquad
 0<v_h^{\rm tilt}\le\frac{2C_h^{\rm tilt}}{\mathfrak M(p-3)}.
 \end{gathered}
\end{equation}
The symbol $v_h^{\rm tilt}$ is the displayed scalar variance, not the
entire Taylor unit $v_h$.  Finiteness follows from BT17 at orders one
and two, which are both admitted even when $m=1$.  The mean is strictly
positive because $w_h$ is even and positive almost everywhere, and
\[
 \int y f(y)\,dy=\int_0^\infty
 y(e^{\alpha y}-e^{-\alpha y})w_h(y)\,dy>0.
\]
The variance is strictly positive because $f$ is not supported at one
point.  Expanding its centered square gives the upper bound by the
second uncentered moment in RTM3.

For $n=k-3\ge2$, let $F_n=f^{*n}$.  The original convolution map has
Jacobian one, and Tonelli, followed by the triangle inequality, gives
\begin{equation}\tag{RTM4}
 \begin{gathered}
 \int F_n(z)\,dz=\mathfrak M^n,\qquad
 \int |z|F_n(z)\,dz\le nJ_h\mathfrak M^{n-1},\\
 \int (z-n\mu_h)F_n(z)\,dz=0,\qquad
 \int (z-n\mu_h)^2F_n(z)\,dz
       =nv_h^{\rm tilt}\mathfrak M^n.
 \end{gathered}
\end{equation}
For the last identity, pull the integral back to the $n$ original
variables and expand $\bigl(\sum_j(y_j-\mu_h)\bigr)^2$.
Every cross term vanishes since $\int(y-\mu_h)f(y)dy=0$.
The $n$ diagonal terms are each $v_h^{\rm tilt}\mathfrak M^n$.
All exchanges are justified by the finite first and second absolute
moments in RTM3.  These are identities for the full, unaltered tilted
mass; no probability measure is substituted for an arithmetic source.

\subsection{A tail lower bound with the complete mass}

Put
\begin{equation}\tag{RTM5}
 H_k=\sqrt{2(k-3)v_h^{\rm tilt}},\qquad
 Y_k^{\rm cen}=\max\{2,(k-3)\mu_h+H_k\},\qquad
 Y_k^{\rm abs}=\max\{2,2(k-3)J_h/\mathfrak M\}.
\end{equation}
The centered-square identity in RTM4 proves
\[
 \int_{|z-n\mu_h|\le H_k}F_n(z)\,dz\ge\mathfrak M^n/2:
\]
on the complementary set its integrand multiplied by $H_k^2$ is at
most the integrand of the centered second moment.  The same argument
with the absolute first moment proves
\[
 \int_{|z|\le Y_k^{\rm abs}}F_n(z)\,dz\ge\mathfrak M^n/2.
\]
Inside every convolution product,
$f^{*j}(x)=e^{\alpha x}w_h^{*j}(x)$, exactly, because the original
coordinates sum to $x$.  Thus RTM2 gives
$f^{*3}(x)\ge c_h(1+x)^{-B}$ whenever $x\ge0$.
If $y\ge Y_k^{\rm cen}$ and $|z-n\mu_h|\le H_k$, then
\[
 0\le y-z\le y-n\mu_h+H_k.
\]
Restrict the positive convolution $f^{*3}*F_n$ to this entire
retained-mass set.  The preceding inequalities prove
\begin{equation}\tag{RTM6}
 \boxed{m_k(y)\ge\frac{c_h}{2}\mathfrak M^{k-3}e^{-\alpha|y|}
 \bigl(1+|y|-(k-3)\mu_h+H_k\bigr)^{-B}
 \quad(|y|\ge Y_k^{\rm cen}).}
\end{equation}
The negative half-line follows from the original evenness of $m_k$.
The denominator is positive on the stated domain; at its unrounded
threshold it equals $1+2H_k$.

Alternatively the absolute-moment retained set proves the simpler bound
\begin{equation}\tag{RTM7}
 \boxed{m_k(y)\ge\frac{c_h}{2}\mathfrak M^{k-3}e^{-\alpha|y|}
 (1+2|y|)^{-B}\quad(|y|\ge Y_k^{\rm abs}).}
\end{equation}
In this proof $0\le y-z\le y+Y_k^{\rm abs}\le2y$ on the
retained set when $y\ge Y_k^{\rm abs}$.  Both cutoffs are $O_h(k)$.
For example RTM3 gives the fully explicit upper cutoff
\[
 Y_k^{\rm abs}\le
 \max\left\{2,\frac{4(k-3)C_h^{\rm tilt}}
                         {\mathfrak M(p-2)}\right\}.
\]
The new feature is the exact factor $\mathfrak M^{k-3}$ instead of
a truncated mass $M_R^{k-3}$.  No loss of tilted mass is hidden in
the cutoff or a limiting argument.

\subsection{The original logarithmic shape on the growing source region}

Retain the exact moving degree
\[
 k=4l+1\ge5,\quad e=1+k(m-1),\quad q=e(k+1)^2,
 \qquad q\ge(k+1)^2.
\]
For $y\ge Y_k^{\rm cen}$, divide RTM6 by the upper Gamma bound in
RTM2.  Divide the existing BT21 upper convolution envelope by the
lower Gamma bound.  These operations give the exact two-sided result
\begin{equation}\tag{RTM8}
 \begin{aligned}
 \log\frac{m_k(y)}{\sigma(y)}\ge{}&(k-3)\log\mathfrak M
 +\log\frac{c_h}{2C_\Gamma}+\tfrac12\log(1+y)
 -B\log(1+y-(k-3)\mu_h+H_k),\\
 \log\frac{m_k(y)}{\sigma(y)}\le{}&(k-3)\log\mathfrak M
 +\log\frac{C_h^{\rm tilt}}{c_\Gamma}+2\log\mathfrak M
 +\log k+\tfrac12\log(1+y)-p\log(1+y/k).
 \end{aligned}
\end{equation}
There is no replacement of the Gamma density by its asymptotic in
these finite inequalities.

Fix $0<a<b<\infty$ and put $y=qx$ with $x\in[a,b]$.
Both cutoffs are eventually satisfied, uniformly in $x$, since
$k/q\to0$.  For sufficiently large $k$ the numbers
$1/y$, $k/y$, and $(1-(k-3)\mu_h+H_k)/y$ have absolute values at
most $1/2$.  The identity
$\log(1+z)=\int_0^z(1+t)^{-1}dt$ gives
$|\log(1+z)|\le2|z|$ there.  Hence RTM8 implies the following
two distinct polynomial-log bounds, with their full constants:
\begin{equation}\tag{RTM9}
 \begin{aligned}
 \log\frac{m_k(qx)}{\sigma(qx)}\ge{}&(k-3)\log\mathfrak M
 +(1/2-B)\log(qx)+\log\frac{c_h}{2C_\Gamma}
 -O_{h,a,b}(k/q),\\
 \log\frac{m_k(qx)}{\sigma(qx)}\le{}&(k-3)\log\mathfrak M
 +(1/2-p)\log(qx)+(p+1)\log k
 +\log\frac{C_h^{\rm tilt}}{c_\Gamma}+2\log\mathfrak M
 +O_{h,a,b}(k/q).
 \end{aligned}
\end{equation}
In particular, uniformly on this compact range,
\begin{equation}\tag{RTM10}
 \log\frac{m_k(qx)}{\sigma(qx)}
 =(k-3)\log\mathfrak M+O_{h,a,b}(\log q).
\end{equation}
RTM9 has different lower and upper powers.  It does not identify a
single polynomial-log shape with an $o(1)$ remainder for the actual
ratio, nor does it evaluate its signed determinant return.

The retained tail is connected to the original convolution by the
literal exact expression
\begin{equation}\tag{RTM11}
\begin{gathered}
 e^{\alpha y}m_k(y)=
 \int_{\mathbb R} f^{*3}(y-z)F_{k-3}(z)\,dz,\\
 \int F_{k-3}=\mathfrak M^{k-3},\quad
 \int(z-(k-3)\mu_h)^2F_{k-3}
 =(k-3)v_h^{\rm tilt}\mathfrak M^{k-3}.
\end{gathered}
\end{equation}
Every complementary contribution remains in this equality.  The
proved lower bounds use only a positive part of it, and make no
claim that the complementary part vanishes.

\paragraph{Exact remaining scope.}
These are estimates on the original density at the specified source
points.  Contributions from $|y|=O_h(k)$ to the original growing
polynomial Grams have not been removed.  In the simple-quartet case
their natural $k^2$ scale can be the same order as $q$.
Consequently RTM6--10 alone do not justify an order-$q$ signed
four-endpoint determinant calculation or a first-variation value.
The source-coordinate map remains $P(S)\mapsto P(k/2+iy)$ with
its inverse, the actual packet relation remains $\chi$, and all
proper-source kernels and Taylor-unit observations are unchanged.
\endgroup


\begingroup
\section{A global tilted-mass transfer with the central source retained}
\label{rtm:sec:global-tilted-mass-deficit}

Retain the original source and constants RTM1--3.  In particular
$\mathfrak M=M_\alpha$ is its full boundary-tilted mass, $J_h$ is its
actual first absolute tilted moment, $B=42+8m$, $M=B/2=21+4m$, and
$p=8m-9/2$.  Keep every full packet order and every original tensor
order $k=4l+1\ge5$, with $q=[1+k(m-1)](k+1)^2$.
The TW global bounds, on their unchanged measures, are
\begin{equation}\tag{GTM1}
 \begin{gathered}
 a_k(1+y^2)^{-M}\sigma(y)\le m_k(y)\le A_k\sigma(y),\\
 a_k=\frac{c_h\vartheta_h^{k-3}e^{-\alpha(k-3)}(k-2)^{-B}}
 {C_\Gamma2^{B/2}},\quad
 A_k=\frac{C_h^{\rm tilt}}{c_\Gamma}k^{p+1}\mathfrak M^{k-1},\quad
 \vartheta_h=\int_{-1}^1w_h(y)dy.
 \end{gathered}
\end{equation}
The following calculation retains the central source region, including
its finite polynomial cost.  It proves an $O_h(q\log q)$ bound,
including $m=1$; it does not evaluate an order-$q$ signed coefficient.

\subsection{An exact global lower measure}

Define finite positive scalars and an auxiliary integer by
\begin{equation}\tag{GTM2}
 \begin{gathered}
 Y=\max\{2,2(k-3)J_h/\mathfrak M\},\quad
 b_\dagger=\frac{c_h}{2C_\Gamma2^{3B/2}},\quad
 \lambda_k=b_\dagger\mathfrak M^{k-3},\\
 r=\max\left\{1,\left\lceil
       \frac{\log(\lambda_k/a_k)}{\log2}\right\rceil\right\},\quad
 Q(y)=(y+iY)^r(y+i)^M,\quad F(y)=y^rQ(y),\quad d=2r+M.
 \end{gathered}
\end{equation}
These auxiliary polynomials occur only in proved comparison maps on
the same source.  They do not replace the actual packet relation
$\chi$ or the observation on its quotient.  The ceiling argument is
exactly
\begin{equation}\tag{GTM3}
 \log(\lambda_k/a_k)
 =(k-3)\{\log(\mathfrak M/\vartheta_h)+\alpha\}
      +B\log((k-2)/2)-\log2.
\end{equation}
In particular $r=O_h(k)$ and $Y=O_h(k)$, including for a simple quartet.

RTM7 and the upper Gamma bound give, when $|y|\ge Y$,
\[
 m_k(y)\ge\frac{c_h}{2C_\Gamma}\mathfrak M^{k-3}
       (1+2|y|)^{-B}\sigma(y)
 \ge\lambda_k(1+y^2)^{-M}\sigma(y).
\]
The first inequality has only discarded the factor $(1+|y|)^{1/2}\ge1$.
For the second use
$(1+2|y|)^B\le2^B(1+|y|)^B
\le2^{3B/2}(1+y^2)^{B/2}$.
The multiplier $(y^2/(Y^2+y^2))^r$ is at most one on this region.
On $|y|\le Y$ it is at most $2^{-r}$, and GTM2 ensures
$\lambda_k2^{-r}\le a_k$.  The original global lower bound GTM1
therefore proves, on the entire real source,
\begin{equation}\tag{GTM4}
 \boxed{m_k(y)\ge\lambda_k
    \left|\frac{y^r}{(y+iY)^r(y+i)^M}\right|^2\sigma(y),
                      \qquad y\in\mathbb R.}
\end{equation}
The denominators have no real zeros.  In particular GTM4 has not
deleted the interval $|y|=O_h(k)$: its cost is recorded by $r,Y$.

\subsection{A finite determinant inequality on the full polynomial space}

On the same monomial frame $1,y,\ldots,y^N$, put
$H_{k,N}=\operatorname{Gram}_{m_k}(y^j)_{j=0}^N$ and
$H_N^\sigma=\operatorname{Gram}_\sigma(y^j)_{j=0}^N$.
They are positive definite.  Write $n=N+1$, $L=N+d$, and define
the positive $n$ by $n$ matrices
\begin{equation}\tag{GTM5}
 \begin{aligned}
 \mathcal A_{ij}&=\int y^{2r+i+j}|Q(y)|^{-2}d\sigma(y),\\
 \mathcal B_{ij}&=\int y^{2r+i+j}d\sigma(y),\qquad
 \mathcal C_{ij}=\int y^{2r+i+j}|Q(y)|^2d\sigma(y),
 \quad0\le i,j\le N.
 \end{aligned}
\end{equation}
The Gram of the lists $y^{r+j}/Q(y)$ and
$y^{r+j}\overline{Q(y)}$ is
$\left(\begin{smallmatrix}\mathcal A&\mathcal B\\
\mathcal B&\mathcal C\end{smallmatrix}\right)$.
Completing the square proves
$\mathcal A\succeq\mathcal B\mathcal C^{-1}\mathcal B$.
Together with GTM4 this gives
\begin{equation}\tag{GTM6}
 \det H_{k,N}\ge\lambda_k^n
             (\det\mathcal B)^2/\det\mathcal C.
\end{equation}
The order of the noncommuting matrices has not been changed.

The original Gamma monic polynomials satisfy the proved exact formulas
\begin{equation}\tag{GTM7}
 \gamma_j=\|b_j\|_\sigma^2
   =\sqrt{2\pi}(2j)!/4^j,\qquad
 yb_j=b_{j+1}+j(j-1/2)b_{j-1},\quad b_0=1.
\end{equation}
The missing term at $j=0$ is zero.  A monic change has determinant
one, so $\det H_N^\sigma=\prod_{j=0}^N\gamma_j$.
The exterior product of $y^r,\ldots,y^{N+r}$ has coefficient one
on $b_r\wedge\cdots\wedge b_{N+r}$.  Expanding in the orthogonal
exterior basis, whose squared coefficients are nonnegative, proves
\begin{equation}\tag{GTM8}
 \det\mathcal B\ge\prod_{j=r}^{N+r}\gamma_j.
\end{equation}

Let $\mathcal J$ take all divided derivatives at
$0,-iY,-i$ with multiplicities $r,r,M$ in that order.
Its kernel in degree at most $L$ is exactly $F\mathcal P_N$.
Its restriction $\mathcal V$ to degrees below $d$ has determinant
\begin{equation}\tag{GTM9}
 \det\mathcal V=(-iY)^{r^2}(-i)^{rM}\{i(Y-1)\}^{rM},
 \qquad |\det\mathcal V|=Y^{r^2}(Y-1)^{rM}\ge1.
\end{equation}
This is the divided-derivative confluent Vandermonde: replace each
root by distinct nearby points, divide the ordinary Vandermonde by
the three within-block factors, and expand their evaluation rows in
Taylor series.  The first nonzero coefficient gives the divided
derivatives and exactly the cross-root factors in GTM9.  In particular
the observation is onto, with every auxiliary jet retained.

In the Gamma orthonormal frame $b_j/\sqrt{\gamma_j}$ its covariance is
\begin{equation}\tag{GTM10}
 (\mathcal K_L)_{(z,a),(z',a')}
  =\sum_{j=0}^L\frac{b_j^{[a]}(z)
                    \overline{b_j^{[a']}(z')}}{\gamma_j}.
\end{equation}
The exact relation determinant is
\begin{equation}\tag{GTM11}
 \det\mathcal C=
  \left(\prod_{j=0}^L\gamma_j\right)
          \frac{\det\mathcal K_L}{|\det\mathcal V|^2}.
\end{equation}
Indeed the frame $1,\ldots,y^{d-1},F,\ldots,Fy^N$ has determinant
one.  Its lower Gram block is $\mathcal C$, and its Schur complement
is the minimum-observation Gram
$\mathcal V^*\mathcal K_L^{-1}\mathcal V$.
The latter follows directly by solving the constrained minimum:
for source Gram $H$ and onto map $J$, its section is
$H^{-1}J^*(JH^{-1}J^*)^{-1}$.  Subtracting it leaves a vector in
$\ker J$ orthogonal to the section.  Taking the determinant of this
block Gram proves GTM11, with its displayed factors.

\subsection{All finite auxiliary-jet and factorial costs}

Define
\begin{equation}\tag{GTM12}
 Z_{L,Y}=\frac{\exp(2)(L+1)(L+2)(2L+1)}{\sqrt{2\pi}}
                 [2(L+2)]^{Y+1}.
\end{equation}
Then $\det\mathcal K_L\le Z_{L,Y}^{d}$.
Here is the coefficient estimate.  GTM7 gives the generating function
\[
 \sum_{j\ge0}\frac{b_j(z)}{j!}t^j
 =(1+t^2)^{-1/4}\exp(z\arctan t),\qquad |t|<1.
\]
To derive it, multiply the recurrence by $t^j/j!$ and sum;
$(1+t^2)\partial_t G=(z-t/2)G$ with $G(0,z)=1$ determines the
displayed analytic solution and its coefficients.
On $|t|=(L+1)/(L+2)$,
$|\arctan t|\le\tfrac12\log(2(L+2))$,
$|1+t^2|\ge(L+2)^{-2}$, and $|t|^{-j}\le\exp(1)$ for $j\le L$.
Cauchy's coefficient bound therefore gives
\[
 |b_j(z)|/j!\le\exp(1)(L+2)^{1/2}[2(L+2)]^{(Y+1)/2}
                       \quad(|z|\le Y+1).
\]
Applying the derivative formula on circles of radius one about
$0,-iY,-i$ gives the same bound for every divided derivative used
in GTM10.  No factorial is removed.  Also
\[
 (j!)^2/\gamma_j
 =4^j/(\sqrt{2\pi}\binom{2j}{j})\le(2j+1)/\sqrt{2\pi},
\]
because the central coefficient is at least the average of all
$2j+1$ binomial coefficients.  Summing squared derivatives proves
that every diagonal entry of $\mathcal K_L$ is at most $Z_{L,Y}$.
Successive orthogonal projection in its Gram frame gives the
determinant bound as the product of at most these diagonal entries.

Put $g_0=\log\sqrt{2\pi}$ and retain the complete finite cost
\begin{equation}\tag{GTM13}
 \begin{aligned}
 \mathcal E_{N,r,Y}={}&2r^2\log(N+2r)
 +M\{g_0+2L\log L\}+2rg_0
 +4r^2\log\max(1,r)\ + (2r+M)\log Z_{L,Y},\\
 & L=N+2r+M.
 \end{aligned}
\end{equation}
Every term is finite and nonnegative on the present domain.
GTM6--12 give
\begin{equation}\tag{GTM14}
 \boxed{(N+1)\log\lambda_k-\mathcal E_{N,r,Y}
 \le\log\frac{\det H_{k,N}}{\det H_N^\sigma}
 \le(N+1)\log A_k.}
\end{equation}
For clarity, the factorial expression on the lower side before
estimation is exactly
\[
 T=2\sum_{j=r}^{N+r}\log\gamma_j
 -\sum_{j=0}^{N+2r+M}\log\gamma_j-\sum_{j=0}^N\log\gamma_j.
\]
Writing these sums as differences of prefix sums, without a condition
such as $r\le N$, gives
\[
 T=\sum_{j=1}^r\{\log\gamma_{N+j}-\log\gamma_{N+r+j}\}
 -\sum_{j=N+2r+1}^{N+2r+M}\log\gamma_j
 -2\sum_{j=0}^{r-1}\log\gamma_j.
\]
The exact ratio $\gamma_{j+1}/\gamma_j=(j+1)(j+1/2)$ bounds
the first negative sum by $2r^2\log(N+2r)$.
The elementary bound
$\log\gamma_j\le g_0+2j\log\max(1,j)$, obtained from
$(2j)!\le(2j)^{2j}$, bounds the other two negative sums by
the corresponding terms of GTM13.  Finally subtract
$d\log Z_{L,Y}$ and use $|\det\mathcal V|\ge1$.
This proves the lower side of GTM14 with all central-region,
auxiliary-jet and factorial costs retained.  The upper side is
the determinant of GTM1's full matrix upper bound.

\subsection{A finite deficit and the simple-quartet improvement}

Define on the original source
\begin{equation}\tag{GTM15}
 \Xi_{k,N}=(N+1)\log A_k
       -\log\frac{\det H_{k,N}}{\det H_N^\sigma}\ge0.
\end{equation}
The nonnegative sign follows from the full upper Gram inequality.
There is no loss of tilted mass in the following exact scalar:
\[
 \log(A_k/\lambda_k)
 =(p+1)\log k+\log\frac{C_h^{\rm tilt}}{c_\Gamma b_\dagger}
                  +2\log\mathfrak M.
\]
Thus GTM14 proves the finite bound
\begin{equation}\tag{GTM16}
 \boxed{0\le\Xi_{k,N}\le
 (N+1)\left\{(p+1)\log k
       +\log\frac{C_h^{\rm tilt}}{c_\Gamma b_\dagger}
       +2\log\mathfrak M\right\}+\mathcal E_{N,r,Y}.}
\end{equation}
Replacing the last two scalar logarithms in braces by the absolute
value of their sum gives a nonnegative, possibly looser, bound.
The displayed unaltered expression is also valid when
$\mathfrak M<1$.

At $N\in\{q-1,q,2q-1,2q\}$, GTM3 and RTM3 give
$r=O_h(k)$, $Y=O_h(k)$, $d=O_h(k)$ and $L=N+O_h(k)$.
GTM12 implies $\log Z_{L,Y}=O_h(k\log(N+k+2))$.
Every term of GTM13 is consequently bounded, with its actual degree,
by
\[
 \mathcal E_{N,r,Y}
 =O_h\bigl(N\log(N+k+2)+k^2\log(N+k+2)\bigr).
\]
Since $q\ge(k+1)^2$ and $q=O_m(k^3)$, GTM16 proves
\begin{equation}\tag{GTM17}
 \boxed{\Xi_{k,N}=O_h(q\log q)
       \quad(N=q-1,q,2q-1,2q),\qquad m\ge1.}
\end{equation}
In particular this includes $m=1$, $q=(k+1)^2$.
The term $O_h(k^2\log q)$ from the central bridge has been retained,
not declared $o(q)$ in that case.  This is precisely why GTM17 does
not assert an order-$q$ asymptotic.

\subsection{Exact transfer to the original quotient and its mixed maps}

Return to $S=c+iy$.  The coefficient map
$P(S)\mapsto P(c+iy)$ has diagonal $1,i,\ldots,i^N$ and
determinant of modulus one.  Applying the same congruence to both
source Grams leaves GTM14--17 and every relative determinant unchanged.
For the remainder map $J_N:\mathcal P_N\twoheadrightarrow E$
with its actual kernel $\chi\mathcal P_{N-q}$, put
\[
 G_N^\mu=(J_N(H_N^\mu)^{-1}J_N^*)^{-1},\qquad
 \Xi_{k,N}^{E}=q\log A_k-\log\frac{\det G_N^{\rm ar}}{\det G_N^\sigma}.
\]
Then
\begin{equation}\tag{GTM18}
 0\le\Xi_{k,N}^{E}\le\Xi_{k,N}.
\end{equation}
Here is the complete matrix argument.  Set
$D=A_k^{-1}(H_N^\sigma)^{-1/2}H_N^{\rm ar}(H_N^\sigma)^{-1/2}$,
so $0\prec D\preceq I$, and let
$V=(H_N^\sigma)^{-1/2}J_N^*(G_N^\sigma)^{1/2}$.
The identity $V^*V=I$ makes this an isometry.
In an orthonormal completion of its columns, the normalized quotient
Gram is the Schur complement
\[
 S=(V^*D^{-1}V)^{-1}
  =A_k^{-1}(G_N^\sigma)^{-1/2}G_N^{\rm ar}(G_N^\sigma)^{-1/2}.
\]
Writing $D$ in two blocks gives $\det D=\det D_{22}\det S$,
with $0\prec D_{22}\preceq I$ and $0\prec S\preceq I$.
Thus $-\log\det S$ is nonnegative and at most
$-\log\det D$, proving GTM18.  An empty complementary block
has determinant one, covering $N=q-1$.

On a fixed actual original mixed observation $O:E\twoheadrightarrow Y$
onto its image frame, keep the complete kernel inclusion $I_K$ and put
\[
 H_{K,N}^\mu=I_K^*G_N^\mu I_K,\quad
 Q_{Y,N}^\mu=(O(G_N^\mu)^{-1}O^*)^{-1},\quad a=\dim K,\quad b=\dim Y,
 \quad a+b=q.
\]
Define the kernel and image deficits using their own dimensions,
\[
 \Xi^K_{k,N}=a\log A_k-\log\frac{\det H_{K,N}^{\rm ar}}
                                      {\det H_{K,N}^\sigma},\qquad
 \Xi^Y_{k,N}=b\log A_k-\log\frac{\det Q_{Y,N}^{\rm ar}}
                                      {\det Q_{Y,N}^\sigma}.
\]
Then
\begin{equation}\tag{GTM19}
 \Xi^K_{k,N}\ge0,\quad\Xi^Y_{k,N}\ge0,\qquad
 \Xi^K_{k,N}+\Xi^Y_{k,N}=\Xi^E_{k,N}\le\Xi_{k,N}.
\end{equation}
To prove the identity, choose any fixed right inverse $S_0$ of $O$.
The fixed coordinate frame $[I_K,S_0]$ is invertible.  Its block Gram
has determinant $\det H_{K,N}^\mu\det Q_{Y,N}^\mu$ by the
same minimum-section Schur complement.  The squared determinant of
the coordinate frame cancels between the two original metrics.
Nonnegativity follows by restriction and by minimization from
$G_N^{\rm ar}\preceq A_kG_N^\sigma$.
This proof preserves the complete original kernel and image, including
empty ones, and never identifies either with the auxiliary jet kernel.
For the original QDC observed boundary, take $O=\Lambda$ and
$Y=\mathcal B$: then the observed quotient is literally
$Q_N^{\mathcal B}=(\Lambda G_N^{-1}\Lambda^*)^{-1}$ and the
kernel metric is $I_K^*G_NI_K$.  When the actual joint observation
$O$ is selected instead, its own fixed image and kernel are used.
All these statements concern an individual original $G_N$.
The proper polynomial-boundary difference metric
$G_N-G_{N'}$ is a different stated expression; subtracting two
source-comparison inequalities does not give a bound for its
cross-source relative determinant.  No such bound is asserted here.

For every $N\le2q$, $\Xi_{k,N}\le\Xi_{k,2q}$.
Indeed the literal coefficient inclusion into $\mathcal P_{2q}$,
transported by the Gamma square roots, is an isometry; its compression
of the normalized positive matrix $D_{2q}\preceq I$ is precisely
$D_N$.  Complete that isometry orthogonally.  The remaining Schur
complement is at most the identity, so
$\det D_{2q}\le\det D_N$ and the asserted deficit inequality follows.

Take the original endpoint signs $(+,+,-,-)$ at
$(q-1,q,2q-1,2q)$.  Each common scalar dimension times
$\log A_k$ cancels because the sign sum is zero.  The remaining
positive and negative sides each contain two nonnegative deficits
bounded by $\Xi_{k,2q}$.  Consequently, for the original signed
quotient correction and its separate kernel and observable corrections,
\begin{equation}\tag{GTM20}
 \boxed{\max\{|\delta_k^\sigma|,|\Delta F_k^K|,|\Delta F_k^Y|\}
 \le2\Xi_{k,2q}=O_h(q\log q),\qquad m\ge1.}
\end{equation}

% BEGIN CURRENT CF RECEIVER GTM20
\paragraph{Current nested transfer of the same full-source deficit.}
AFT3 proves the exact original primary Vandermonde congruence for
this source Gram. AFT6--17 then carries this same $\Xi_{k,2q}$
through the actual conductor, kernel, common and complementary
flag, with their full fixed frame determinant. The sum of the
four nonnegative deficits is precisely the primary deficit;
every subset of their signed returns has the single bound
$2\Xi_{k,2q}$. This is the transfer used in AFT27, without
discarding any proper kernel or mixed Schur block.
% END CURRENT CF RECEIVER GTM20

The finite upper bound for the right side is exactly twice GTM16
at $N=2q$, with the integer guard GTM2 and complete cost GTM13.
Every relative frame and dimension has been specified before
performing the signed cancellation.
For $O=\Lambda$, the third term of GTM20 is exactly the original
observed-boundary correction $|\Delta F_k^B|$; it is not the norm
of the proper polynomial boundary in the TCD construction.

The same improvement propagates through the existing complete scalar
return.  Retain its original identity
$S_k=\mathcal B_k^\sigma+\delta_k^\sigma-4q\log(D_hk)$,
its exact GQB22 Gamma-energy bound $E_B(q)$, and its complete
packet comparison $C_k^{\rm cmp}$.  The proved GQB7 and GQB22
give
$|\mathcal B_k^\sigma-C_Bq^2|
\le E_B(q)+(2q+2)C_k^{\rm cmp}$.
Inserting GTM20 in this unchanged signed identity proves
\begin{equation}\tag{GTM21}
 \boxed{\begin{gathered}
 |S_k-C_Bq^2|\le E_B(q)+(2q+2)C_k^{\rm cmp}
                   +2\Xi_{k,2q}+4q|\log(D_hk)|,\\
 S_k=C_Bq^2+O_h(q\log q),\qquad m\ge1.
 \end{gathered}}
\end{equation}
For the second line use the already proved
$E_B(q)=O(q\log q)$, $C_k^{\rm cmp}=O_h(e^{-1})=O_h(1)$,
the fixed positive $D_h$, and GTM17.  All finite packet cutoffs
required by those Gamma providers remain in the first line.

% END PRESERVED INPUT 12_RTM_GTM_original_providers.tex
\endgroup

\clearpage
\begingroup


\setcounter{equation}{0}
\renewcommand{\theequation}{SXR.\arabic{equation}}
\section*{Original signed source-relation receiver}
\addcontentsline{toc}{section}{Original signed source-relation receiver}
% BEGIN PRESERVED INPUT 13_SXR_COMPLETE.tex SHA256 ba2e763bcf453fc3e0c8ac324fe5c734dbfecf85065d7821da96ca7f468e8dbb
\begingroup
\section{The signed return through the actual adjacent relation and quotient}
\label{sxr:sec:sxr}

Retain the full original hypothetical packet
$\rho=1/2+\delta+i\gamma$, $0<\delta<1/2$, $\gamma>2$, and its
complete order $m$. Set
\begin{equation}\tag{SXR1}
\begin{gathered}
 g=2\xi,\quad h(s)=\prod_{z\in\{\rho,\bar\rho,1-\rho,1-\bar\rho\}}
                    (s-z)^m,\quad v_h=g/h,\\
 w_h(y)=|v_h(1/2+iy)|^2/(2\pi),\quad m_k=w_h^{*k},\quad
 k=4l+1\ge5,\quad e=1+k(m-1),\\
 q=e(k+1)^2,\quad c=k/2,\quad
 \chi(S)=\prod_{a,b=0}^k
  (S-c-(2a-k)\delta-i(2b-k)\gamma)^e.
\end{gathered}
\end{equation}
The full Taylor unit, source masses, period parameter and tensor
coordinates remain fixed. All limits fix this complete packet and
use the displayed admitted tensor orders. The original zero hypothesis
is the one in the total counterfactual; no numerical offcritical zero
is supplied or asserted by this calculation.

\subsection{The unchanged source path and its exact derivatives}

The original CA14a source interpolation is, with its parameter written
$x$ to distinguish it from the period $u$,
\begin{equation}\tag{SXR2}
 d\mu_x(y)=(1-x)d\sigma(y)+x m_k(y)dy,\quad
 d\sigma(y)=\frac{|\Gamma(1/4+iy/2)|^2}{2\pi}dy,
 \quad 0\le x\le1.
\end{equation}
In particular its mass is
$(1-x)\sqrt{2\pi}+x(\int w_h)^k$. There is no replacement by
a probability measure. On $\mathcal P_d=\mathbb C[S]_{\le d}$
use the basis $1,S,\ldots,S^d$ and
\[
 (H_d(x))_{ab}=\int\overline{(c+iy)^a}(c+iy)^b\,d\mu_x(y),
 \qquad 0\le a,b\le d.
\]
These positive matrices have $d+1$ rows. Embeddings into
$\mathcal P_{2q}$ are literal zero extensions. Write
\begin{equation}\tag{SXR3}
 M(x)=H_{2q}(x),\quad \dot M=H_{2q}^{\rm ar}-H_{2q}^{\sigma},
 \quad C(x)=M(x)^{-1}\dot M.
\end{equation}
The dot denotes this source derivative only. The operator $C$ is
self-adjoint in $M$, since $MC=\dot M=\dot M^*$. For any two
polynomials in the ambient space,
\[
 \langle P,CQ\rangle_x
   =\int\overline{P(c+iy)}Q(c+iy)
                     (m_k(y)-\sigma(y))dy.
\]
This identity proves its exact connection to the original density
change. It is not a freely chosen coefficient operator.

Let $p_d(x)$ be the original monic source orthogonal polynomial and
$u_r(x)$ the monic polynomial in the original relation norm. Thus
\begin{equation}\tag{SXR4}
 \omega_d(x)=\|p_d(x)\|_x^2,\quad
 \nu_r(x)=\|\chi u_r(x)\|_x^2,\quad
 T_r(x)=\frac{\nu_r(x)}{\omega_{q+r}(x)},\quad d=q+r.
\end{equation}
The polynomial $u_r$ is orthogonal to degrees below $r$ for the
literal measure $|\chi(c+iy)|^2d\mu_x(y)$. Its full $S$-coordinate
monicity is retained. Positive finite moment Grams give unique monic
minimizers, and solving their linear systems shows that their
coefficients are smooth rational functions of $x$. Denominators
are positive on the whole compact interval. Differentiating the
minimum norms gives exactly
\begin{equation}\tag{SXR5}
 \dot\omega_d=\langle p_d,Cp_d\rangle_x,\qquad
 \dot\nu_r=\langle\chi u_r,C\chi u_r\rangle_x.
\end{equation}
Indeed $\dot p_d$ has degree below $d$ and $\chi\dot u_r$ lies
in $\chi\mathcal P_{r-1}$. Both coefficient-derivative cross terms
vanish by their respective defining orthogonalities. This proves
SXR5 without differentiating an asymptotic remainder.

The source is even. In real $y$ coordinates its monic orthogonal
polynomial has simple real roots by the sign-change proof in TCN4.
The polynomial $\chi(c+iy)$ has nonreal roots, so $\chi u_r$ cannot
be that monic minimizer. Therefore $T_r(x)>1$ for every $x$ and
$0\le r\le q$. The phase change is exactly
$P(S)\mapsto i^{-d}P(c+iy)$ in degree $d$; no source norm changes.

\subsection{An orthogonal exchange in the actual quotient section}

Put $b_d(x)=[p_d(x)]_\chi$ in the original quotient
$E=\mathbb C[S]/(\chi)$ and keep
\[
 G_d=(J_dH_d^{-1}J_d^*)^{-1},\qquad R_d=H_d^{-1}J_d^*G_d.
\]
For $d=q+r$ define the following actual source vectors:
\begin{equation}\tag{SXR6}
\begin{gathered}
 f_r=\frac{p_d}{\sqrt{\omega_d}},\quad
 g_r=\frac{\chi u_r}{\sqrt{\nu_r}},\quad
 \tau_r=T_r^{-1},\quad
 z_r=\frac{b_d}{\sqrt{\omega_d(1-\tau_r)}},\\
 w_r=\frac{f_r-\sqrt{\tau_r}g_r}{\sqrt{1-\tau_r}}
           =R_d z_r.
\end{gathered}
\end{equation}
All square roots here are the positive square roots of the original
positive real norms. To prove the equalities, monicity and the
degree-$d$ orthogonality give
$\langle p_d,\chi u_r\rangle_x=\omega_d$, hence
$\langle g_r,f_r\rangle_x=\sqrt{\tau_r}$ with this exact positive
phase. The complete relation space in degree $d$ is
$\chi\mathcal P_r$. Both $p_d$ and $\chi u_r$ are orthogonal to
$\chi\mathcal P_{r-1}$, and the coefficient of the last orthogonal
relation vector in $p_d$ is $\omega_d/\nu_r$. Consequently
\begin{equation}\tag{SXR7}
 R_db_d=p_d-\frac{\omega_d}{\nu_r}\chi u_r,
 \qquad b_d^*G_db_d=\omega_d(1-T_r^{-1}).
\end{equation}
This is the actual full relation projection, not a deletion of its
highest row. It proves SXR6 and
$\|g_r\|_x=\|w_r\|_x=1$, $\langle g_r,w_r\rangle_x=0$.
The vector $g_r$ belongs to the original relation space and $w_r$
to the original minimum quotient section in the same degree.

For a unit vector $v$ write $P_v=vv^*M$ for its source-orthogonal
projector. SXR5 gives the exact identity
\begin{equation}\tag{SXR8}
 \frac{d}{dx}\log T_r
   =\operatorname{Tr}((P_{g_r}-P_{f_r})C)
   =\langle g_r,Cg_r\rangle_x-\langle w_r,Cw_r\rangle_x
      +\operatorname{Tr}((P_{w_r}-P_{f_r})C).
\end{equation}
In the orthonormal basis $(g_r,w_r)$, the last projector difference
is exactly
\[
 P_{w_r}-P_{f_r}=
 \begin{pmatrix}
 -\tau_r&-\sqrt{\tau_r(1-\tau_r)}\\
 -\sqrt{\tau_r(1-\tau_r)}&\tau_r
 \end{pmatrix}.
\]
Its eigenvalues are $\sqrt{\tau_r},-\sqrt{\tau_r}$; all remaining
ambient eigenvalues are zero. If
$\operatorname{osc}_x C=\lambda_{\max}(C)-\lambda_{\min}(C)$,
the Rayleigh quotients at the two unit eigenvectors therefore prove
\begin{equation}\tag{SXR9}
 |\operatorname{Tr}((P_{w_r}-P_{f_r})C)|
          \le T_r^{-1/2}\operatorname{osc}_x C.
\end{equation}
Every scalar part of $C$ cancels exactly because this projector
difference has trace zero. A bound by an uncentered source norm is
neither needed nor substituted for the oscillation.

\subsection{All four endpoint signs before the estimate}

Let $V_N(x)=\det G_N(x)$ and retain precisely
\begin{equation}\tag{SXR10}
\begin{gathered}
 \mathcal B_k(x)=\log V_{q-1}(x)+\log V_q(x)
                     -\log V_{2q-1}(x)-\log V_{2q}(x),\\
 \delta_k^\sigma=\mathcal B_k(1)-\mathcal B_k(0),\qquad
 c_0=c_q=1,\quad c_r=2\ (1\le r\le q-1).
\end{gathered}
\end{equation}
The monic source/relation change of basis has determinant one, so
$V_{q+r-1}=\prod_{d=0}^{q+r-1}\omega_d/\prod_{j=0}^{r-1}\nu_j$.
Taking consecutive ratios and then the four displayed endpoint signs
gives the exact finite equality
$\mathcal B_k(x)=\sum_{r=0}^q c_r\log T_r(x)$.
In particular $\sum c_r=2q$, with coefficient one at both endpoints.

Define the original relation-versus-quotient receiving integral
\begin{equation}\tag{SXR11}
 \mathcal Q_k=\int_0^1\sum_{r=0}^q c_r
  \bigl(\langle g_r,Cg_r\rangle_x
                  -\langle w_r,Cw_r\rangle_x\bigr)dx.
\end{equation}
This uses the original source path, actual minimum sections, full
relation rows and positive monic phases of SXR6. It is an explicitly
constructed scalar of that source. SXR8 gives
\begin{equation}\tag{SXR12}
 \delta_k^\sigma=\mathcal Q_k+\mathcal E_k,\qquad
 \mathcal E_k=\int_0^1\sum_{r=0}^q c_r
          \operatorname{Tr}((P_{w_r}-P_{f_r})C)dx.
\end{equation}
There is no assertion that either summand has a necessary sign.

Let $b_-,b_+$ be the extreme eigenvalues of the unchanged relative
Gram $(H_{2q}^\sigma)^{-1}H_{2q}^{\rm ar}$. Diagonalizing its
self-adjoint representative in the original Gamma metric gives
the exact eigenvalues of $C(x)$:
$(b_j-1)/(1-x+xb_j)$. This expression increases in $b_j$,
with derivative $(1-x+xb_j)^{-2}>0$. Direct integration gives
\begin{equation}\tag{SXR13}
 \int_0^1\operatorname{osc}_x C\,dx
       =\log(b_+/b_-).
\end{equation}
The formula at $b_j=1$ follows by continuity. It preserves even a
large common source mass until the ratio cancels it exactly.

\subsection{A finite certificate and an exponentially small error}

Keep the original TW constants, without changing their values:
\[
 \frac{a_k}{D_d}\|P\|_\sigma^2\le\|P\|_{\rm ar}^2
                      \le A_k\|P\|_\sigma^2,
 \quad \deg P\le d,\quad
 L_{k,d}=\log(A_kD_d/a_k).
\]
Their full definitions are TCN10 and the unchanged TW/BT providers
in the source closure; $L_{k,2q}=O_h(k+\log q)$. On SXR2 they
give lower and upper factors
$l_d(x)=1-x+x a_k/D_d$, $u(x)=1-x+xA_k$.
The ratio $u(x)/l_d(x)$ increases from $1$ to
$A_kD_d/a_k$: its derivative has numerator $A_k-a_k/D_d\ge0$.
The last inequality follows already by applying TW to a nonzero
polynomial. Taking the actual monic minima in numerator and
denominator therefore proves, simultaneously for every $x$,
\begin{equation}\tag{SXR14}
 T_r(x)\ge e^{-L_{k,2q}}T_r^\sigma,
 \qquad 0\le r\le q,\qquad
 \log(b_+/b_-)\le L_{k,2q}.
\end{equation}
Here $T_r^\sigma$ uses the same complete $\chi$ and its monic
relation norm in the original Gamma source, not a replacement
arithmetic packet.

It is a finite explicit determinant quantity. With
$L_r^\sigma=\operatorname{Gram}_\sigma(\chi S^j)_{0\le j<r}$,
$\det L_0^\sigma=1$, define
\begin{equation}\tag{SXR15}
 \Theta_k=e^{-L_{k,2q}}
   \min_{0\le r\le q}
   \frac{\det L_{r+1}^\sigma}{\det L_r^\sigma}
       \frac{4^{q+r}}{\sqrt{2\pi}(2q+2r)!}.
\end{equation}
Every determinant here is in its original monic frame and retains
the full Gamma mass. For example its moments are obtained exactly
by differentiating
$\int e^{ity}d\sigma(y)=\sqrt{2\pi}(\cosh t)^{-1/2}$ at zero.
This is the same Gamma Fourier identity used in the retained source
provider, with its stated mass. No arithmetic zero is needed to
compute these reference matrices at fixed packet parameters.

Combining SXR9, SXR12--15 and the strict inequality $T_r(x)>1$
proves the finite bound
\begin{equation}\tag{SXR16}
 \boxed{|\delta_k^\sigma-\mathcal Q_k|
 \le 2q\min\{1,\Theta_k^{-1/2}\}\log(b_+/b_-)
 \le 2q L_{k,2q}\min\{1,\Theta_k^{-1/2}\}.}
\end{equation}
All $q+1$ adjacent exchanges and both original endpoint weights
were included before this estimate was taken.

For completeness the new TCN/TCX estimates evaluate the size of
this finite bound. They give
\begin{equation}\tag{SXR17}
 \min_{0\le r\le q}\log T_r^\sigma
        =2q a_1+O_h(k+\log q),\qquad a_1=\psi(1)>0.
\end{equation}
Here is the required uniform argument rather than an inference
from a four-endpoint assertion. For
$0\le r\le\lfloor q/64\rfloor$, TCX2--9's original-root interval
and monic arcsine bound apply to $\sigma$ with comparison factor
one. They give a lower rate $c_*q$, $c_*>1/3>2a_1$, with
error $O_h(k+\log q)$. For the remaining integers, TCN13 and
TCN17--26 give $\log T_r^\sigma=2q\psi(r/q)+O_h(k+\log q)$
uniformly on the fixed compact ratio interval $[1/128,1]$,
retaining the original packet cutoff. Since $\psi$ decreases,
this is at least $2qa_1-O_h(k+\log q)$. At $r=q$ the same formula
gives a matching upper bound for the minimum. The original cutoff
holds eventually because $q\ge(k+1)^2$ and
$R=k\sqrt{\delta^2+\gamma^2}$. This proves SXR17 and
$\log\Theta_k=2qa_1+O_h(k+\log q)$.
Consequently the exact error in SXR12 satisfies
\begin{equation}\tag{SXR18}
 \boxed{|\mathcal E_k|
       \le\exp\{-qa_1+O_h(k+\log q)\}=o(q).}
\end{equation}
This is a bound on an actual signed correction term, not a claimed
evaluation of the still-present $\mathcal Q_k$.

\subsection{Return to the original joint kernel and observable image}

At fixed original nonzero period and branch, keep the complete
observation $O:E\to\mathbb C^b$ onto its fixed image frame, with
its full Taylor units, reduced mixed map and ordered periods. It
does not vary with the source interpolation. Let $I_K$ be the fixed
kernel frame. At each actual degree $d=q+r$ set
\begin{equation}\tag{SXR19}
\begin{gathered}
 Q_d=(OG_d^{-1}O^*)^{-1},\quad
 P_{K,d}=I_K(I_K^*G_dI_K)^{-1}I_K^*G_d,\\
 P_{Y,d}=G_d^{-1}O^*Q_dO,\quad
 w_{K,r}=R_dP_{K,d}z_r,\quad w_{Y,r}=R_dP_{Y,d}z_r.
\end{gathered}
\end{equation}
Empty kernels or images use their unique empty matrices. These
are the original source-Gram minimum sections, at their written
degrees; no boundary metric is substituted. Multiplying the
formulas gives $OP_{K,d}=0$, $OP_{Y,d}=O$ and orthogonality of
their ranges in $G_d$. Hence $P_K+P_Y=I$ and
$w_r=w_{K,r}+w_{Y,r}$, with the two summands orthogonal in $M(x)$.
The literal receiving integrand is therefore
\begin{equation}\tag{SXR20}
\begin{aligned}
 \langle g_r,Cg_r\rangle_x-\langle w_r,Cw_r\rangle_x
 ={}&\langle g_r,Cg_r\rangle_x
   -\langle w_{K,r},Cw_{K,r}\rangle_x
   -\langle w_{Y,r},Cw_{Y,r}\rangle_x\\
 &-2\Re\langle w_{K,r},Cw_{Y,r}\rangle_x.
\end{aligned}
\end{equation}
Orthogonality for $M$ does not remove the last pairing with $C$.
It is retained, with its original sign, inside $\mathcal Q_k$.
The quotient coefficient is exactly $z_r$ in SXR6, and its observed
value is exactly $Oz_r$. There is no assertion that it is the
distinguished primary or that its observable fraction is uniformly
nonzero. On the primary, the earlier nonvanishing subsequence and
the undivided JQC identity retain their proved domains.

The relation $\chi u_r$ has its original ordered theta primitive
at $W_*$ or a containing source. Its proper-source class and full
$h^2$-jet keep the TCD3--5 map and image restrictions. Applying the
linear source maps coefficientwise retains every coefficient face,
outer source label and nonempty leg mask. A zero remains at its
present support and is not the external $\tau$. Thus the controlled
exchange has been returned to the original mixed construction with
its exact morphisms; it has not erased that construction's kernel.

Equations SXR12, SXR16 and SXR20 leave a precisely specified original
mixed scalar with an exponentially bounded exchange error. The sign
and order-$q$ value of this scalar have not been computed here.
Neither the new error estimate nor the separate positive exterior
energy proves RH or certifies an offcritical zero.
\endgroup
% END PRESERVED INPUT 13_SXR_COMPLETE.tex
\endgroup

\clearpage
\begingroup


\setcounter{equation}{0}
\renewcommand{\theequation}{OPG.\arabic{equation}}
\section*{Original polynomial quotient and tensor maps}
\addcontentsline{toc}{section}{Original polynomial quotient and tensor maps}
% BEGIN PRESERVED INPUT 14_OPG_COMPLETE.tex SHA256 5c192e0fc7d46d95df2908ba6d92d1a325b02223ae1c247ac81112bf86ce52d1
\section{The original period vectors and their Gamma generating map}
\label{opg:sec:OPG}

The observation in this calculation is the literal BRU1--3 matrix. We
retain its entire amplitude unit, its primary multiplicities, its
period contours, its cyclic projector and its kernel. The only source
orders used are the already present orders $s\in\{1,k\}$. The letter
$s$ in a source-order subscript is an integer parameter; the coordinate
on the one-factor polynomial algebra will be denoted $z_1$ below to
keep it distinct. The generating variable is $z$.

\subsection{The original finite algebras and the observation}
Retain
\[
\begin{gathered}
 \rho=\tfrac12+\delta+i\gamma,\quad 0<\delta<\tfrac12,\quad\gamma>2,
 \quad\mathcal R=\{\rho,\bar\rho,1-\rho,1-\bar\rho\},\\
 h(z_1)=\prod_{\alpha\in\mathcal R}(z_1-\alpha)^m,
 \quad n=4m,\quad v_h=g/h,\quad g=2\xi,\quad m\geq1,\\
 E_h=\mathbb C[z_1]/(h),\qquad A_h=M_{z_1}:E_h\to E_h,
 \qquad U=M_{j_hv_h}:E_h\to E_h.
\end{gathered}\tag{OPG1}
\]
Here $m$ is the original complete zero order at each of the four
points. Consequently $v_h(\alpha)\ne0$ for every $\alpha\in\mathcal R$.
The entire-function remainder $j_hv_h$ records its derivatives of
orders $0,\ldots,m-1$ at every one of these points. In each primary
algebra it has nonzero constant term, and its inverse is the finite
geometric series for that constant plus a nilpotent. Thus $U$ is the
same invertible unit as in BRU, with none of its values prescribed.

For the original positive integer $k$, set
\[
\begin{gathered}
 c=k/2,\quad \ell_k=1+k(m-1),\quad q=(k+1)^2\ell_k,\\
 \chi(S)=\prod_{a,b=0}^k
   \bigl(S-c-(2a-k)\delta-i(2b-k)\gamma\bigr)^{\ell_k},\\
 E=\mathbb C[S]/(\chi),\quad A=M_S:E\to E,\quad
 Y=(A-cI)/i,\\
 A_\Sigma=\sum_{j=1}^k
      I_{E_h}^{\otimes(j-1)}\otimes A_h\otimes I_{E_h}^{\otimes(k-j)},
 \qquad
 \iota:E\hookrightarrow E_h^{\otimes k},\quad
 [P]\longmapsto[P(z_1+\cdots+z_k)].
\end{gathered}\tag{OPG2}
\]
For clarity, the stated map is injective with exactly this defining
polynomial. On an ordered primary tensor the eigenvalue of
$A_\Sigma$ is the sum of the ordered roots. Its nilpotent part is
the sum of the $k$ independent nilpotents of orders $m$.
Its power $k(m-1)$ has nonzero coefficient
$(k(m-1))!/((m-1)!)^k$ on their top monomial, and its next power
is zero. The distinct sums are exactly the $(k+1)^2$ displayed
roots of $\chi$, since the signs of the real and imaginary parts
can be selected independently in every factor. Distinct pairs
$(a,b)$ give distinct sums by $\delta,\gamma>0$. Their common
nilpotency index is $\ell_k$. These facts prove that the minimal
polynomial of $A_\Sigma$ is $\chi$, and therefore prove the
kernel of polynomial evaluation and the asserted injection.
In particular
\[
 \iota A=A_\Sigma\iota,\qquad
 \iota Y=\frac{A_\Sigma-cI}{i}\iota,\qquad
 \iota[1]=[1]^{\otimes k}.
 \tag{OPG3}
\]

Keep the original $u\ne0$, its logarithm branch, and $t=0$ in
the period construction. With $\Phi_h'=h$, $\Phi_h(0)=0$,
let $\ell_j$ be the ray of angle
$(\arg u+\pi+2\pi j)/(n+1)$ and let
$\Gamma_j=-\ell_0+\ell_j$, $1\leq j\leq n$, with the original
orientations. The original period map is
\[
 \Pi:E_h\to\mathbb C^n,\qquad
 (\Pi[P])_j=\int_{\Gamma_j}e^{\Phi_h(z_1)/u}
                    \operatorname{rem}_h P(z_1)\,dz_1.
 \tag{OPG4}
\]
The representative on the right has degree less than $n$. These
are exactly the original convergent polynomial period integrals;
their existence and full determinant are proved in PDE1--13 and
IPM1--7 in the retained source closure. This definition does not
assign a period to an arbitrary representative of a quotient class.
Let $C$ be the original matrix with rows corresponding to
$\Gamma_j\mapsto\Gamma_{j+1}-\Gamma_1$ for $j<n$ and
$\Gamma_n\mapsto-\Gamma_1$. The reduced endpoint basis gives
$C^{n+1}=I$. Hence the literal finite average below is idempotent:
\[
\begin{gathered}
 P_k=\frac1{n+1}\sum_{a=0}^{n}(C^{-a})^{\otimes k},\qquad
 L=P_k\Pi^{\otimes k}U^{\otimes k}\iota,\\
 B=\operatorname{im}L,\quad\jmath:B\hookrightarrow(\mathbb C^n)^{\otimes k},
 \qquad\Lambda:E\twoheadrightarrow B,\quad\jmath\Lambda=L,
 \qquad K=\ker\Lambda.
\end{gathered}\tag{OPG5}
\]
The basis inclusion $\jmath$, the inclusion $I:K\hookrightarrow E$,
and a coefficient section $S_*:B\to E$ with $\Lambda S_*=I_B$
are the fixed BRU choices. The full actual image is used at every
rank. Empty boundary matrices use the original empty conventions.

\subsection{An exact generating identity before and after observation}
Put $b_s=s/2$, $M_s=(2\pi)^{s/2}$, and retain the monic Gamma
polynomials and their full norms from GMB1:
\[
\begin{gathered}
 p_0(y)=1,\quad p_1(y)=y,\quad
 p_{d+1}(y)=yp_d(y)-d(b_s+d-1)p_{d-1}(y),\\
 h_d^{(s)}=M_s d!(b_s)_d,\quad
 r_d=p_d(Y)[1],\quad e_d=r_d/\sqrt{h_d^{(s)}},\quad
 \lambda_d=\Lambda e_d\in B.
\end{gathered}\tag{OPG6}
\]
The square root is the positive real square root. In particular
$e_d=[u_d]$ in the original notation, including the phase $i^{-d}$
in its leading coefficient. Define the holomorphic germs of
$\log(1+z^2)$ and $\arctan z$ by their values zero at $z=0$.
For $|z|<1$, the exact algebra-valued identity is
\[
 \sum_{d=0}^{\infty}r_d\frac{z^d}{d!}
 = (1+z^2)^{-s/4}\exp((\arctan z)Y)[1].
 \tag{OPG7}
\]
Indeed the right side is holomorphic in that disc, is $[1]$ at
zero, and satisfies
$(1+z^2)F'(z)=(Y-b_s zI)F(z)$.
Comparing its Taylor coefficients gives exactly the recurrence
in OPG6: the contribution of $z^2F'$ to coefficient $z^d/d!$
is $d(d-1)r_{d-1}$, and the contribution of $b_szF$ is
$b_sd r_{d-1}$. Thus its derivatives at zero are the $r_d$,
which proves the convergent identity and not merely a formal
coefficient equality.

All factors of $A_\Sigma$ act on different tensor coordinates
and therefore commute. Their matrix exponential power series
are absolutely convergent in finite dimension. The multinomial
expansion of every power consequently proves
\[
 \exp\left(w\frac{A_\Sigma-cI}{i}\right)[1]^{\otimes k}
 =\left(\exp\left(w\frac{A_h-\tfrac12I}{i}\right)[1]\right)^{\otimes k}.
 \tag{OPG8}
\]
The centre on the left is exactly $c=k/2$; each of the $k$
centres on the right is exactly $1/2$. Applying OPG3--5 to OPG7,
and then using OPG8, gives the original observed generating map
\[
\boxed{
 \jmath\sum_{d=0}^{\infty}\sqrt{h_d^{(s)}}\lambda_d\frac{z^d}{d!}
 = (1+z^2)^{-s/4}P_k
 \left[\Pi U\exp\left((\arctan z)
                \frac{A_h-\tfrac12I}{i}\right)[1]\right]^{\otimes k}.
}\tag{OPG9}
\]
The scalar outside $P_k$ has exponent $-s/4$. It has not acquired
an extra factor of $k$: OPG7 concerns one original Gamma source
on the sum coordinate, while OPG8 is the exact factorization of
that coordinate action in the already existing tensor algebra.
For $s=1$ and $s=k$ this gives the two stated source versions of
the same original observation. No additional order is selected.

\subsection{All primary coefficients and the original period integrals}
Here is a finite formula for every vector in OPG9. For
$\alpha\in\mathcal R$ write
\[
\begin{gathered}
 h_\alpha(z_1)=h(z_1)/(z_1-\alpha)^m,\qquad
 t_\alpha(z_1)=\sum_{j=0}^{m-1}
       \frac{(h_\alpha^{-1})^{(j)}(\alpha)}{j!}(z_1-\alpha)^j,\\
 \varepsilon_\alpha=[h_\alpha t_\alpha]\in E_h,
 \qquad E_{\alpha,a}(z_1)=
   \operatorname{rem}_h\bigl(h_\alpha(z_1)t_\alpha(z_1)
                                  (z_1-\alpha)^a\bigr),\quad0\leq a<m,\\
 \mathbf b_{\alpha,a}=\Pi[E_{\alpha,a}],\qquad
 (\mathbf b_{\alpha,a})_j=
   \int_{\Gamma_j}e^{\Phi_h(z_1)/u}E_{\alpha,a}(z_1)\,dz_1.
\end{gathered}\tag{OPG10}
\]
The Taylor inverse exists because $h_\alpha(\alpha)\ne0$.
Reduction modulo $(z_1-\alpha)^m$ gives
$\varepsilon_\alpha=1$ on its primary factor and zero on every
other primary factor. These observations prove
$\varepsilon_\alpha\varepsilon_\beta=0$ for $\alpha\ne\beta$,
$\varepsilon_\alpha^2=\varepsilon_\alpha$, and
$\sum_\alpha\varepsilon_\alpha=1$: their differences are divisible
by all four coprime factors and hence by $h$. They also prove that
$[E_{\alpha,a}]=\varepsilon_\alpha(z_1-\alpha)^a$ is a basis of
$E_h$. Every representative in OPG10 has degree less than $n$,
so each integral is precisely of the original polynomial type.

For $w\in\mathbb C$ put
$\mathbf v(w)=\Pi U\exp(w(A_h-\tfrac12I)/i)[1]$.
The primary nilpotent $N_\alpha=(z_1-\alpha)\varepsilon_\alpha$
satisfies $N_\alpha^m=0$. Taylor multiplication in that algebra
therefore proves, with all amplitude derivatives retained,
\[
\boxed{
 \mathbf v(w)=\sum_{\alpha\in\mathcal R}
 e^{w(\alpha-1/2)/i}\sum_{a=0}^{m-1}\mathbf b_{\alpha,a}
       \sum_{\nu=0}^{a}
       \frac{v_h^{(\nu)}(\alpha)}{\nu!}
       \frac{(w/i)^{a-\nu}}{(a-\nu)!}.
}\tag{OPG11}
\]
For a direct verification before integration, both
$U\exp(w(A_h-\tfrac12I)/i)[1]$ and
$j_h(v_h(z_1)e^{w(z_1-1/2)/i})$ have the same first $m$ Taylor
coefficients at every $\alpha$. Their difference is therefore
zero in $E_h$. The representative to which $\Pi$ is applied is
exactly the finite polynomial in the basis $E_{\alpha,a}$ whose
coefficients are displayed in OPG11. In particular this proof
does not replace the integrand in OPG10 by the unreduced entire
function $v_h(z_1)e^{w(z_1-1/2)/i}$. The precise map from that
entire function to the period vector is its finite Taylor
remainder $j_h$, followed by the degree-less-than-$n$ representative,
followed by the original polynomial period integral.

The cyclic projector is also fully evaluated as a finite operation:
\[
 \jmath\sum_{d\geq0}\sqrt{h_d^{(s)}}\lambda_d\frac{z^d}{d!}
 =\frac{(1+z^2)^{-s/4}}{n+1}
       \sum_{a=0}^{n}\bigl(C^{-a}\mathbf v(\arctan z)\bigr)^{\otimes k}.
 \tag{OPG12}
\]
This is obtained by applying each tensor power in OPG5 to a pure
tensor. It retains every term of the original average, its factor
$1/(n+1)$, and its inverse powers of $C$.

There is also a finite coefficient formula requiring no series
extraction. For an ordered $k$-tuple
$\boldsymbol\alpha\in\mathcal R^k$ set
$\omega_{\boldsymbol\alpha}=(\sum_i\alpha_i-c)/i$.
For $0\leq a_i<m$ and $0\leq\nu_i\leq a_i$ put
$D(\mathbf a,\boldsymbol\nu)=\sum_i(a_i-\nu_i)$ and
\[
 T_d(\boldsymbol\alpha,\mathbf a)=
 \sum_{\substack{0\leq\nu_i\leq a_i\\1\leq i\leq k}}
 \frac{i^{-D(\mathbf a,\boldsymbol\nu)}
        p_d^{(D(\mathbf a,\boldsymbol\nu))}(\omega_{\boldsymbol\alpha})}
        {\prod_{i=1}^k(a_i-\nu_i)!}
 \prod_{i=1}^k\frac{v_h^{(\nu_i)}(\alpha_i)}{\nu_i!}.
 \tag{OPG13}
\]
As usual the polynomial derivative is zero when its order exceeds
$d$. Multivariate Taylor expansion in the independent nilpotents
$(z_i-\alpha_i)$ proves the exact formula
\[
\boxed{
 \jmath\lambda_d=
 \frac1{\sqrt{h_d^{(s)}}}\frac1{n+1}
 \sum_{a=0}^{n}\ \sum_{\boldsymbol\alpha\in\mathcal R^k}
 \sum_{\substack{0\leq a_i<m\\1\leq i\leq k}}
 T_d(\boldsymbol\alpha,\mathbf a)
 \bigotimes_{i=1}^k C^{-a}\mathbf b_{\alpha_i,a_i}.
}\tag{OPG14}
\]
Indeed on that primary tensor, expand each $v_h(z_i)$ first.
For the remaining multi-degree $a_i-\nu_i$, Taylor expansion of
$p_d((\sum_i z_i-c)/i)$ supplies its derivative of the total
order $D$, the factor $i^{-D}$, and the factorials
$\prod_i(a_i-\nu_i)!$. These are exactly the coefficients in
OPG13. Apply the $k$ polynomial period maps and the literal cyclic
average to get OPG14. All sums are finite. Ordered tuples have
not been replaced by just their eigenvalue sum, so both the
amplitude data and every independent primary nilpotent survive.

\subsection{The actual observation increment without a coordinate inverse}
Retain the GMB maps
$C_s=[e_0,\ldots,e_{q-1}]:\mathbb C^q\to E$,
$v_j=C_s^{-1}e_{q+j}$,
$K_j=I_q+\sum_{a<j}v_av_a^*$, and $R_s=\Lambda C_s$.
The low-degree leading coefficients prove $C_s$ invertible as in
GMB2. On each column, applying the maps gives
\[
 R_s v_j=\lambda_{q+j},\qquad
 R_sK_jR_s^*=\sum_{d=0}^{q+j-1}\lambda_d\lambda_d^*.
 \tag{OPG15}
\]
For the second equality the identity term in $K_j$ gives exactly
the first $q$ columns of $R_s$, and each remaining rank-one term
gives the next displayed column. The low vectors $\lambda_d$,
$0\leq d<q$, span $B$ because $C_s$ is invertible and $\Lambda$
is onto. Thus, in the fixed original basis of $B$, the matrix
\[
 \mathcal M_D=\sum_{d=0}^{D-1}\lambda_d\lambda_d^*\quad(D\geq q)
 \tag{OPG16}
\]
is positive definite whenever $B$ has positive dimension:
for a nonzero covector its quadratic form is the sum of its
squared values on spanning vectors. Consequently the exact
original boundary increment of GMB8 is
\[
 \boxed{\displaystyle
 \xi_j=\lambda_{q+j}^*\mathcal M_{q+j}^{-1}\lambda_{q+j},\qquad
 1+\xi_j=\frac{\det\mathcal M_{q+j+1}}{\det\mathcal M_{q+j}}.}
 \tag{OPG17}
\]
The determinant equality follows by conjugating the rank-one
update by $\mathcal M_{q+j}^{-1/2}$; its sole possible nonzero
added eigenvalue is the displayed quadratic form. If $B=0$,
the vector is empty, $\xi_j=0$ and both determinants equal one.
Thus OPG11--14 compute the boundary increment from the original
period vectors, with no $C_s^{-1}$ in OPG17. This cancellation
is an identity of the proved maps, rather than a change of the
observation or the source measure.

At the four original endpoints the boundary contribution is
therefore the exact finite expression
\[
 F_B^{(s)}=
 \log\frac{\det\mathcal M_{2q}\det\mathcal M_{2q+1}}
                 {\det\mathcal M_q\det\mathcal M_{q+1}}
 =\log(1+\xi_0)+2\sum_{j=1}^{q-1}\log(1+\xi_j)
                         +\log(1+\xi_q).
 \tag{OPG18}
\]
Expanding the last sum of consecutive log determinant differences
leaves coefficient $-1$ at $\mathcal M_q,\mathcal M_{q+1}$ and
$+1$ at $\mathcal M_{2q},\mathcal M_{2q+1}$, proving every sign.
GMB13 proves that the other original factor is
$(1+t_j)/(1+\xi_j)$, with $t_j=v_j^*K_j^{-1}v_j$ and
$0\leq\xi_j\leq t_j$. Thus OPG18 is the actual BRU boundary
contribution, and the full original kernel contribution remains
$\mathcal W_q(\log((1+t_j)/(1+\xi_j)))$.

\subsection{A finite recurrence retaining the exact action defect}
The period formula also has a finite recurrence on the fixed
original coefficient spaces. Define the uniquely typed projection
$\kappa:E\to K$ by
$I\kappa=I_E-S_*\Lambda$. It exists because the right side has
image in $\ker\Lambda=\operatorname{im}I$, and $I$ is an
inclusion. It satisfies $\kappa I=I_K$ and $\kappa S_*=0$.
In particular $E=I(K)\oplus S_*(B)$ as coefficient spaces, without
an assertion that this splitting is preserved by $Y$.
Set
\[
\begin{gathered}
 \mathsf Y_B=\Lambda YS_*:B\to B,\quad
 \mathsf D=\Lambda YI:K\to B,\\
 \mathsf C=\kappa YS_*:B\to K,\quad
 \mathsf Y_K=\kappa YI:K\to K,\quad
 \zeta_d=\kappa e_d\in K.
\end{gathered}\tag{OPG19}
\]
All four maps are computable from OPG1--5 and the fixed BRU
section. For $d\geq0$ put
\[
 a_d=\frac1{\sqrt{(d+1)(b_s+d)}},\qquad
 b_d=\sqrt{\frac{d(b_s+d-1)}{(d+1)(b_s+d)}}\quad(d\geq1),
 \qquad b_0=0.
 \tag{OPG20}
\]
Set $\lambda_{-1}=0$ and $\zeta_{-1}=0$. Then the exact coupled
recurrence is
\[
\boxed{\begin{aligned}
 \lambda_{d+1}&=a_d(\mathsf Y_B\lambda_d+\mathsf D\zeta_d)
                                          -b_d\lambda_{d-1},\\
 \zeta_{d+1}&=a_d(\mathsf C\lambda_d+\mathsf Y_K\zeta_d)
                                          -b_d\zeta_{d-1},\\
 \lambda_0&=\Lambda[1]/\sqrt{M_s},\qquad
 \zeta_0=\kappa[1]/\sqrt{M_s}.
\end{aligned}}\tag{OPG21}
\]
To prove this, the norm ratio from OPG6 gives
$\sqrt{h_d^{(s)}/h_{d+1}^{(s)}}=a_d$ and
$d(b_s+d-1)\sqrt{h_{d-1}^{(s)}/h_{d+1}^{(s)}}=b_d$.
Hence $e_{d+1}=a_dYe_d-b_de_{d-1}$, including $d=0$ by its
initial values. Substituting $e_d=S_*\lambda_d+I\zeta_d$ and
applying $\Lambda$ and $\kappa$ proves both lines of OPG21.
It retains the kernel coordinate at every degree and supplies
the period vectors through degree $2q$, exactly the range used
in OPG17--18.

The failure of a putative boundary action has the precise map
\[
 \Lambda Y-\mathsf Y_B\Lambda=\mathsf D\kappa:E\to B,\qquad
 \jmath\mathsf D=
 \frac1iP_k\Pi^{\otimes k}U^{\otimes k}
                   (A_\Sigma-cI)\iota I:K\to(\mathbb C^n)^{\otimes k}.
 \tag{OPG22}
\]
The first equality follows from $I_E=S_*\Lambda+I\kappa$;
the second follows from OPG3--5. These equations preserve the
cyclic average and the full unit. They also prove exactly that
$Y(K)\subseteq K$ holds if and only if $\mathsf D=0$.
No such vanishing is used anywhere in OPG7--21. Thus every
possible boundary action defect is retained inside a proved
coupled action and inside the explicit original period formula.

For completeness the multiplication action does descend to a
canonical quotient that remembers every observed iterate. Define
\[
 \mathcal K_\infty=\bigcap_{a=0}^{q-1}\ker(\Lambda Y^a),\qquad
 \mathcal O:E\to B^{\oplus q},\quad
 x\longmapsto(\Lambda x,\Lambda Yx,\ldots,\Lambda Y^{q-1}x).
 \tag{OPG23}
\]
The polynomial $\chi(c+iT)$ annihilates $Y$ by OPG2, has degree
$q$ and leading coefficient $i^q\ne0$. It therefore expresses
$Y^q$ as a linear combination of $I,Y,\ldots,Y^{q-1}$, with
the exact coefficients of this original polynomial. If
$x\in\mathcal K_\infty$, each component of $\mathcal O(Yx)$
vanishes by this relation. Thus $\mathcal K_\infty$ is
$Y$-invariant and is precisely $\bigcap_{a\geq0}\ker(\Lambda Y^a)$.
Moreover any $Y$-invariant subspace contained in $K$ lies in
this intersection. Consequently the first isomorphism map
\[
\begin{gathered}
 E/\mathcal K_\infty\ \xrightarrow{\ \overline{\mathcal O}\ }
 \operatorname{im}\mathcal O,\quad[x]\longmapsto\mathcal O(x),\\
 \chi_j=[T^j]\chi(c+iT),\quad\chi_q=i^q,\\
 \mathsf{Shift}(b_0,\ldots,b_{q-1})=
 \left(b_1,\ldots,b_{q-1},-i^{-q}\sum_{j=0}^{q-1}\chi_j b_j\right).
\end{gathered}
 \tag{OPG24}
\]
is a bijection intertwining multiplication by $Y$ with the shift
of the $q$ components, whose last component is the stated exact
linear combination. It projects to the original observation by
the first component, with retained kernel $K/\mathcal K_\infty$.
This proves an action-preserving enlargement of the observed
data on the original algebra. It neither discards the original
kernel in OPG21 nor presumes a vanished defect in OPG22.

The original period isomorphism also locates that defect in the
cyclic projection. Its inverse exists by the retained period
determinant theorem used in OPG4. On the unchanged ambient tensor
space put
\[
 \mathcal V=\Pi^{\otimes k}U^{\otimes k}\iota,
 \qquad
 \mathcal T=\sum_{j=1}^k
 \left(\Pi\frac{A_h-\tfrac12I}{i}\Pi^{-1}\right)_j.
 \tag{OPG25}
\]
The subscript $j$ means that the displayed factor acts in position
$j$ and identities act in all other positions. Multiplication by
$j_hv_h$ commutes with $A_h$. Equations OPG3 and OPG25 therefore
give $\mathcal VY=\mathcal T\mathcal V$ exactly. Since
$L=P_k\mathcal V$ and $P_k^2=P_k$, splitting the ambient identity
as $P_k+(I-P_k)$ now proves
\[
\begin{gathered}
 LY=P_k\mathcal T L+P_k\mathcal T(I-P_k)\mathcal V,\\
 \jmath\mathsf D=P_k\mathcal T(I-P_k)\mathcal V I.
\end{gathered}\tag{OPG26}
\]
For the second equality, $P_k\mathcal V I=LI=0$, so the
first term of the first equality vanishes on $I(K)$. This is
the same defect as OPG22 on its original domain. The two ambient
terms in the first line sum to a vector in $\jmath(B)$; neither
individual term is asserted to preserve that image. OPG19--21
give the full recurrence within $K\oplus B$ without such an
assertion. Thus the preserved unit and period isomorphism carry
the multiplication action exactly, and OPG26 identifies precisely
the retained contribution of the original cyclic projection.

Equations OPG9, OPG11--14, OPG17 and OPG21 now give the same
original boundary increments by analytic generating functions,
finite primary coefficients, and finite coupled recurrences.
Their derivations establish the exact maps connecting these
calculations. In particular the new input into the original
mixed-control problem is the fully specified sequence
$\lambda_0,\ldots,\lambda_{2q}$ with its actual unit and periods,
rather than a freely selected observation or a source-order
replacement. The formulas themselves do not assign a growing-
parameter estimate to the individual weighted sums; they supply
their unchanged mathematical input for that calculation.

% END RX INLINE 16

% END PRESERVED INPUT 14_OPG_COMPLETE.tex
\endgroup

\clearpage
\begingroup


\setcounter{equation}{0}
\renewcommand{\theequation}{PCL.\arabic{equation}}
\section*{Original period-coefficient and pivot restriction provider}
\addcontentsline{toc}{section}{Original period-coefficient and pivot restriction provider}
% BEGIN PRESERVED INPUT 15_PCL_COMPLETE.tex SHA256 231acb7cd4902b6d3477f853d1d8dea814c71b94cff2a24e7fdee5481652ef54
\section{The literal original conductor limit and an invertible pivot restriction}
\label{pcl:sec:PCL}
Retain all original hypotheses and maps of PQO, RPC, OCF and CEP.
In particular $0<\delta<1/2$, $\gamma>2$, the original unit values
are $(a,\bar a,\bar a,a)$ with $a\ne0$, and the same logarithm
branch is used throughout. Write $\beta=2(\gamma^2-\delta^2)$,
$\eta=(\delta^2+\gamma^2)^2$, and
$\mathfrak a=1/160+\beta/24+\eta/2$. No source order is changed.

\subsection{An entire coefficient matrix in the inverse original period}
Let $P_0$ be RPC3's centered Fourier period matrix with its factor
$e^{\mathfrak a/u}$ removed by the stated exact identity, and retain
$G_u=\operatorname{diag}(u^{a/5}g_a)_{a=1}^4$ with RPC8's full
Gamma constants. Introduce the coefficient matrix $\mathcal R$ by
the following identity, rather than replacing the original period:
\[
 P= e^{\mathfrak a/u}P_0
   = e^{\mathfrak a/u}G_u\mathcal R(u^{-1}).
 \tag{PCL1}
\]
The row index below is $1\le a\le4$, and the column index is
$0\le r\le3$. The complete series is
\[
 \begin{split}
 \mathcal R_{a r}(z)&=\sum_{n=0}^{\infty}R_{a r,n}z^n,\\
 R_{a r,n}&=
 \sum_{\substack{p,h\ge0\\2p+4h=5n+r+1-a}}
 \frac{\beta^p\eta^h}{3^p p!h!}
      (-1)^{L}5^L(a/5)_L,
 \qquad L=p+h-n.
 \end{split}\tag{PCL2}
\]
Every sum at a fixed $n$ is finite. For each contributing pair,
$r+3p+h+1=a+5L$ and $L\ge0$. To prove the latter assertion,
the left side is positive and $1\le a\le4$, so a negative multiple
of five cannot be their difference. The exponent $n$ is nonnegative:
its numerator $2p+4h+a-r-1$ is an integer multiple of five and is
at least $-3$. Thus a negative value of $n$ is impossible.

For a direct derivation, divide every term of the absolutely
convergent original series PQO9 by $u^{a/5}g_a$. The quotient of
its Gamma factors is exactly
\[
 \frac{5^{(a+5L)/5-1}e^{\pi i(a+5L)/5}
                  \Gamma(a/5+L)}{g_a}
        =(-1)^L5^L(a/5)_L.
 \tag{PCL3}
\]
Its remaining power of $u$ is $u^{L-p-h}=u^{-n}$.
These identities prove PCL2 including its signs, factorials and
all unit-independent period coefficients. Absolute convergence of
PQO9 permits grouping these terms. For each prescribed complex
$z\ne0$, put $u=1/z$ on any of its original branches. The resulting
sum of absolute values is finite. Consequently PCL2 converges
absolutely for every complex $z$, and defines an entire matrix.
This also proves that this matrix is independent of the branch;
all branch factors remain in $G_u$ in PCL1.

The complete constant matrix in PCL2 is
\[
 J_\infty=
 \begin{pmatrix}
 1&0&-\beta/3&0\\
 0&1&0&-2\beta/3\\
 0&0&1&0\\
 0&0&0&1
 \end{pmatrix},\qquad
 J_\infty^{-1}=
 \begin{pmatrix}
 1&0&\beta/3&0\\
 0&1&0&2\beta/3\\
 0&0&1&0\\
 0&0&0&1
 \end{pmatrix}.
 \tag{PCL4}
\]
Indeed $n=0$ forces $2p+4h=r+1-a\le3$, so either $p=h=0$
or $p=1,h=0,r+1-a=2$. PCL3 gives respectively the diagonal
entries and the two negative entries displayed. In particular,
the original row-scaled period limit includes both cubic-phase
corrections. It is not a diagonal coefficient map.

RPC's determinant proof holds on the complete coefficient family
and gives $\det P_0=\det G_u$. Alternatively divide PQO11 by
the product of the four unchanged Gamma ray factors. Therefore
\[
       \det\mathcal R(z)=1,\qquad
       \mathcal R(z)^{-1}=\operatorname{adj}\mathcal R(z)
                       \quad(z\in\mathbb C).
 \tag{PCL5}
\]
The equality first holds for $z\ne0$ by PCL1 and extends to zero
by continuity. Both the matrix and its inverse are entire. For
any fixed $R>0$, the finite, completely specified constant
\[
 C_R=\sum_{n\ge1}\|R_n\|_1 R^{n-1}<\infty
 \tag{PCL6}
\]
gives $\|\mathcal R(z)-J_\infty\|_1\le C_R|z|$ on $|z|\le R$.
Here $\|\cdot\|_1$ is the induced column-sum matrix norm and
$R_n=(R_{a r,n})$. Convergence follows from the entire power
series in finite dimension, also at any larger radius. Thus
this is a bound in the actual coefficients, with no uncomputed
nonzero minor used as a premise.

\subsection{All eighty-one actual conductor coefficients and their limit}
Let $V$ be the original centered evaluation Vandermonde in PQO12,
$U=\operatorname{diag}(a,\bar a,\bar a,a)$, and
$D=\operatorname{diag}(\zeta,\zeta^2,\zeta^3,\zeta^4)$.
The complete period matrix from primary to Fourier coordinates is
\[
 \mathsf H(u)=e^{\mathfrak a/u}G_u\mathcal R(u^{-1})V^{-1}U.
 \tag{PCL7}
\]
For $1\le j\le4$ define the entire matrix and biform
\[
 \begin{split}
 L_j(z)&=U^{-1}V\mathcal R(z)^{-1}D^{-j}
                              \mathcal R(z)V^{-1}U,\\
 t(Z,W)&=(ZW,Z,W,1)^{\mathsf T},\\
 A_z(Z,W)&=\prod_{j=1}^4 Q_0(L_j(z)t(Z,W))
                  =\sum_{r,s=0}^8 a_{rs}(z)Z^rW^s.
 \end{split}\tag{PCL8}
\]
These are precisely all original OCF4 coefficients at $z=u^{-1}$.
To prove this, compute
\[
 \mathsf H(u)^{-1}D^{-j}\mathsf H(u)=L_j(u^{-1}):
 \tag{PCL9}
\]
the scalar $e^{\mathfrak a/u}$ cancels with its inverse, and the
two diagonal matrices $G_u,D^{-j}$ commute. OCF's definition
$Q_j(x)=Q_0(\mathsf H^{-1}D^{-j}x)$ followed by $x=\mathsf Ht$
then proves PCL8. The homogeneous version uses
$t=(z_0w_0,z_0w_1,z_1w_0,z_1w_1)$ and retains the bidegree
$(8,8)$. Thus no unit or period entry has been omitted by the
affine notation.

For an explicit finite formula for the individual coefficients,
write $L_{j,b e}$ for a matrix entry, and associate the four pairs
$(\rho_e,\sigma_e)=(1,1),(1,0),(0,1),(0,0)$ to its columns.
Put
\[
 \begin{split}
 c_{j,rs}(z)&=
 \sum_{\substack{1\le e,f\le4\\
       \rho_e+\rho_f=r,\ \sigma_e+\sigma_f=s}}
 (L_{j,1e}L_{j,4f}-L_{j,2e}L_{j,3f}),\\
 a_{rs}(z)&=
 \sum_{\substack{0\le r_j,s_j\le2\\
             \sum_jr_j=r,\ \sum_js_j=s}}
                       \prod_{j=1}^4c_{j,r_js_j}(z).
 \end{split}\tag{PCL10}
\]
Expanding the two products defining $Q_0$ proves the first line;
expanding its four original orbit factors proves the second.
Both lines are finite and retain all ordered-pair multiplicities.
Equations PCL2, PCL5 and PCL10 are a convergent formula for every
actual coefficient, with no freely assigned amplitude.

Their exact leading values are obtained by the explicit matrices
\[
 L_j(0)=U^{-1}VJ_\infty^{-1}D^{-j}J_\infty V^{-1}U,
 \qquad a_{rs}(u^{-1})=a_{rs}(0)+O_{\rm actual}(|u|^{-1}).
 \tag{PCL11}
\]
The equality for $L_j(0)$ follows from PCL4, and the bound follows
from the entire coefficient series. More explicitly, if
$a_{rs}(z)=\sum_{n\ge0}a_{rs,n}z^n$, the finite constant
$\sum_{r,s}\sum_{n\ge1}|a_{rs,n}|R^{n-1}$ bounds the sum of
all eighty-one errors divided by $|z|$ on $|z|\le R$.
The coefficient functions are single-valued entire functions of
$u^{-1}$ even though the original period itself retains its
logarithm branch. This assertion follows from the exact conjugation
PCL9, not from a change of the original branch.

At any specified $k$, the complete central convolution section is
therefore the explicitly evaluated entire matrix
\[
 C_{\rm in}(z)_{(a,b),(i,j)}=a_{i+4-a,j+4-b}(z),
 \qquad 0\le a,b,i,j\le k-8,
 \tag{PCL12}
\]
with all out-of-range coefficients zero. PCL10--12 give its full
constant matrix and all successive coefficients. They do not
discard off-diagonal coefficients or replace this matrix by its
central coefficient $a_{44}$. In particular the calculation of
the limit has not asserted a sign or diagonal-dominance inequality
for this complete original matrix.

\subsection{An exact positivity slice for the literal positive period ray}
For $u>0$ on any of its original logarithm lifts, PQO21 gives
$P_0=E_{\rm br}R$ with $R$ real invertible and $E_{\rm br}$ diagonal
unitary commuting with $D$. Its exact centered quadratic coefficient
matrix is $S$ of PQO13; all its entries are purely imaginary.
Consequently, for real $y\in\mathbb R^4$,
\[
 \overline{Q_j(E_{\rm br}y)}=-Q_{5-j}(E_{\rm br}y)
              \quad(1\le j\le4).
 \tag{PCL13}
\]
Indeed the original quadratic matrix in these coordinates is
$e^{-2\mathfrak a/u}R^{-\mathsf T}SR^{-1}$, which is purely
imaginary. Complex conjugation changes its sign and changes
$D^{-j}$ into $D^j$; this is precisely the displayed formula.

Let $\mathsf P_{\rm conj}$ be the permutation matrix exchanging
primary indices $1,2$ and $3,4$. For real $Z$ and $|W|=1$, direct substitution
gives $\mathsf P_{\rm conj}\overline{t(Z,W)}=W^{-1}t(Z,W)$.
The original roots and unit satisfy
$\bar V=\mathsf P_{\rm conj}V$ and
$\mathsf P_{\rm conj}\bar U=U\mathsf P_{\rm conj}$. Thus
\[
 \overline{V^{-1}Ut(Z,W)}=W^{-1}V^{-1}Ut(Z,W).
 \tag{PCL14}
\]
Choose either square root $w$ of $W$. Then
$p=w^{-1}V^{-1}Ut(Z,W)$ is real by PCL14, and
$\mathsf Ht=wE_{\rm br}y$ for the real vector
$y=e^{\mathfrak a/u}Rp$. Homogeneity and PCL13 prove the exact
nonnegative slice identity
\[
 \boxed{\quad
 W^{-4}A_{u^{-1}}(Z,W)
   =|Q_1(E_{\rm br}y)Q_2(E_{\rm br}y)|^2\ge0
       \quad(Z\in\mathbb R,\ |W|=1).\quad}
 \tag{PCL15}
\]
The choice of sign of $w$ has no effect because the quadratic
forms have even degree. In particular, comparing the Laurent
coefficients on the unit circle, and then the polynomial
coefficients in real $Z$, gives the exact symmetry
\[
             a_{rs}(u^{-1})=\overline{a_{r,8-s}(u^{-1})}.
 \tag{PCL16}
\]
PCL11 extends this identity and PCL15 to the coefficient limit
at $u^{-1}=0$ by continuity. The slice in PCL15 has one real
coordinate and one unit-circle coordinate; it has not changed
either coordinate into a second unit circle.

\subsection{A complete invertible restriction supplied by the original pivot}
At each actual $u$ in the proved five-orbit locus, use precisely
OCF5's lexicographically largest nonzero coefficient $(r_*,s_*)$,
$a_*=a_{r_*s_*}\ne0$, and its original pivot rectangle
$\mathcal P_k=\{(r_*+i,s_*+j):0\le i,j\le k-8\}$.
Let $\mathcal E_k$ be its complement and $S_k$ the original
insertion of its coordinates. No new coefficient has been selected.
The full original maps satisfy
\[
 C N_k^{\mathsf T}=0,\qquad
 S_k^{\mathsf T}N_k^{\mathsf T}=I_\Delta,\qquad
 \operatorname{im}N_k^{\mathsf T}=V_A.
 \tag{PCL17}
\]
The first is the transpose of $N_kC^{\mathsf T}=0$ in OCF7;
the second is the transpose of $N_kS_k=I$. Dimension $\Delta$
then proves the last equality, since $C$ has full row rank.
It follows that restriction to the actual strip is the exact
isomorphism
\[
 S_k^{\mathsf T}:V_A\longrightarrow\mathbb C^{\mathcal E_k},
        \qquad (S_k^{\mathsf T}|_{V_A})^{-1}=N_k^{\mathsf T}.
 \tag{PCL18}
\]
Equivalently its square pivot block has nonzero diagonal $a_*$
and is triangular in the elimination order; OCF's complete
division supplies its inverse. This is an alternative restriction
to the central one in PCL12, with the exact changed set of roots.

Let
$P_{\mathcal P,k}(y)=\prod_{(a,b)\in\mathcal P_k}(y-\omega_{ab})$
and $P_{\mathcal E,k}(y)=\prod_{(a,b)\in\mathcal E_k}(y-\omega_{ab})$,
where $\omega_{ab}=(2b-k)\gamma-i(2a-k)\delta$. Thus
$P_k=P_{\mathcal P,k}P_{\mathcal E,k}$ exactly. Let
$V_{\mathcal E}$ evaluate $1,y,\ldots,y^{\Delta-1}$ at those
actual outer roots and set
$D_{\mathcal P,\mathcal E}=\operatorname{diag}
(P_{\mathcal P,k}(\omega_{ab}))_{\mathcal E_k}$. These two matrices
are invertible by distinctness of the original roots. The full
polynomial-target restriction and its exact inverse are
\[
 \begin{split}
 \rho_{\mathcal E}
 &=V_{\mathcal E}^{-1}D_{\mathcal P,\mathcal E}^{-1}
                         S_k^{\mathsf T}|_{V_A},\\
 \rho_{\mathcal E}^{-1}
 &=N_k^{\mathsf T}D_{\mathcal P,\mathcal E}V_{\mathcal E}.
 \end{split}\tag{PCL19}
\]
Multiplying these matrices and using PCL17 proves both inverse
identities, with all inner products of root differences retained.
In particular the original residual becomes the following
specified onto map, with its exact complete kernel:
\[
 \Xi_{\mathcal E}=\Xi_k^{\rm fix}
       N_k^{\mathsf T}D_{\mathcal P,\mathcal E}V_{\mathcal E},
 \qquad\ker\Xi_{\mathcal E}=\rho_{\mathcal E}(K).
 \tag{PCL20}
\]
This follows by composition with the proved isomorphism, not by
deleting any common-kernel vector.

Finally the original metric and its entire source graph are carried
by this isomorphism as follows. Write
$J_{\mathcal E}=N_k^{\mathsf T}D_{\mathcal P,\mathcal E}V_{\mathcal E}$.
At every original endpoint the actual conductor metric in these
outer coordinates is exactly
\[
 H_{\mathcal E,N}=J_{\mathcal E}^*G_NJ_{\mathcal E},\qquad
 Q_{\Xi,N}^{\rm fix}
     =(\Xi_{\mathcal E}H_{\mathcal E,N}^{-1}
                      \Xi_{\mathcal E}^*)^{-1}.
 \tag{PCL21}
\]
The first formula is the pullback along $\rho_{\mathcal E}^{-1}$;
the minimum-section proof of GIP6 gives the second. On the full
GIP graph its target map is
$\rho_{\mathcal E}E_N^{\rm graph}$, with the same complete
$t$ ideal columns, and composing with $\Xi_{\mathcal E}$ gives
exactly $B_N^{\rm graph}$. Its residual kernel is therefore all
$t+r$ columns in GIP4, including the low-$v$ coordinates.
The coefficient-space bound for $N_k$ does not remove the displayed
root-value and Gamma matrices from PCL21. These formulas retain
the full metric transport needed for the next comparison.

\subsection{The exact inverse-period reflection of the central section}
Let $E_{\rm par}=\operatorname{diag}(1,-1,1,-1)$ in the centered
coefficient order, and let $\mathsf P_{\rm opp}$ exchange primary
indices $1,4$ and $2,3$. Equation PCL2 implies
\[
 \mathcal R(-z)=E_{\rm par}\mathcal R(z)E_{\rm par},\qquad
 V E_{\rm par}=\mathsf P_{\rm opp}V,\qquad
 U\mathsf P_{\rm opp}=\mathsf P_{\rm opp}U.
 \tag{PCL22}
\]
For the first equality a contributing coefficient satisfies
$5n+r+1-a=2p+4h$, so $n\equiv a-r-1\pmod2$. The signs from
the two parity matrices on entry $(a,r)$ are
$(-1)^{a-1+r}=(-1)^n$. Summing the entire series proves the
identity. The second equality evaluates each monomial at the
opposite original root. The third uses the full original unit
values $(a,\bar a,\bar a,a)$, which agree at opposite roots.

The two diagonal matrices $E_{\rm par}$ and $D$ commute. Substitution
into PCL8 and use of PCL22 therefore give the exact maps
\[
 L_j(-z)=\mathsf P_{\rm opp}L_j(z)\mathsf P_{\rm opp},\qquad
 A_{-z}(Z,W)=Z^8W^8 A_z(Z^{-1},W^{-1}).
 \tag{PCL23}
\]
For the second identity, $Q_0(\mathsf P_{\rm opp}t)=Q_0(t)$
by its two explicit products. Also
$\mathsf P_{\rm opp}t(Z,W)=ZWt(Z^{-1},W^{-1})$ when $ZW\ne0$.
Each of the four quadratic factors therefore contributes $(ZW)^2$.
This proves the identity on that open set; both sides are
polynomials after multiplication by $Z^8W^8$, so it holds
identically including both coordinate axes.

Comparing every coefficient in PCL23 proves
\[
 a_{rs}(-z)=a_{8-r,8-s}(z),\qquad
 C_{\rm in}(-z)=C_{\rm in}(z)^{\mathsf T},\qquad
 \det C_{\rm in}(-z)=\det C_{\rm in}(z).
 \tag{PCL24}
\]
Indeed the $(a,b),(i,j)$ entry on the left of the middle identity
is $a_{4+a-i,4+b-j}(z)$, precisely the transposed entry of
PCL12. Out-of-range coefficients are zero on both sides.
The determinant equality follows by the permutation expansion
of a determinant and its transpose. Thus the exact central
determinant is an even entire function of $u^{-1}$, with every
coefficient supplied by PCL2 and PCL10--12. This retains its
actual constant matrix and does not assign that determinant
a value by deleting any of its entries.

% END EXTENDED PCL
% BEGIN ADDITIVE MIXED PROFILE INTEGRATION RMR_PERIOD_RECEIVER.tex
% EXACT INPUT SHA256 f752149d660b38b978aef1e03f4e17d07e020ab041ce3b5a2d979d69e3392159
\label{pcl:MPBUILD:RMR_PERIOD_RECEIVER:start}
\subsection{The evaluated zero-order domain in the original period}
The complete PZ, PZX and PZU calculations below evaluate the previously
retained value $A_{u^{-1}}(1,1)$ of this same conductor. In their exact
notation put $r=a/\bar a$,
$T=(\delta+i\gamma)^2/(\delta-i\gamma)^2$,
$A_e=(2T-1)/(3(T-1))$ and $B_e=(2-T)/(3(T-1))$.
For the four actual factors
$d_j(\tau)=Q_0(L_j(\tau)\boldsymbol1)$, the first nonzero
inverse-period degree is
\[
 \ell_A=0\quad(r\notin\{1,A_e/B_e\}),\qquad
 \ell_A=2\quad(r\in\{1,A_e/B_e\}),\qquad
 b_j=[\tau^{\ell_A}]d_j(\tau)\ne0.
 \tag{RMR1}
\]
The complete second-order matrices and polynomial nonvanishing
certificate at both exceptional phases are PZX1--6; no phase is
excluded. Each $d_j$ is even entire by PCL22. Hence, with the actual
finite-valued convergent sum $C_j=\sum_{n\ge1}|[\tau^{\ell_A+2n}]d_j|$, define
$R_A=\max\{R_*,1,\max_j(2C_j/|b_j|)^{1/2}\}$.
For $|u|\ge R_A$ summing that Taylor tail gives
\[
 \frac{\prod_j|b_j|}{16|u|^{4\ell_A}}
 \le |E_{A,u}(0)|\le
 \frac{81\prod_j|b_j|}{16|u|^{4\ell_A}},\qquad
 v=0,\qquad \alpha_A=1/E_{A,u}(0).
 \tag{RMR2}
\]
Indeed each $u^{\ell_A}d_j(u^{-1})$ lies within $|b_j|/2$ of
$b_j$, and multiplication of the four complete factors proves
both inequalities. The nonzero entire product has only finitely
many zeros in $|\tau|\le1/R_*$; therefore $v=0$ also on the
complement of the explicitly defined finite set PZU5 in the original
domain $|u|\ge R_*$. At its exceptional periods the
complete PZ12 derivative test and $v\le32$ remain in force.

On this proved $v=0$ domain, substitution into the original maps gives
\[
 \begin{gathered}
 \mathbb C[S]_{<v}=0,\quad t_{0,n}=0,\quad
 m_n=16n-48,\quad t_n^{\rm EFP}=8n-16,\quad\\
 j_{0,n}=8n-32,\quad \kappa_e=d_{K,e}=d_{T,e}=0.
 \end{gathered}
 \tag{RMR3}
\]
The last three equalities follow from BTK25 and the exact RRC14
sequence. Thus the full low-cycle and connecting-cycle coefficient
spaces coincide; their original quotient maps have identical fibres
and identical attained metrics. All RRC16 tail corrections vanish.
The ROG low-coordinate shear is the identity. The conductor,
$\alpha_A$, every coefficient and the complete CPQ mixed forms
remain present: PZU9--11 gives the exact size of this scalar and
its unchanged division recursion, including the exceptional phases.
In particular these equalities determine the low-order correction;
the connecting rank remains the exact CPR18 intersection rank.

\label{pcl:MPBUILD:RMR_PERIOD_RECEIVER:end}
% END ADDITIVE MIXED PROFILE INTEGRATION RMR_PERIOD_RECEIVER.tex


\clearpage
% BEGIN EXTENDED ROG
% Complete original relation operator; no source ideal or low component omitted.

% END PRESERVED INPUT 15_PCL_COMPLETE.tex
\endgroup

\clearpage
\begingroup


\setcounter{equation}{0}
\renewcommand{\theequation}{RPC.\arabic{equation}}
\section*{Original period determinant and its constants}
\addcontentsline{toc}{section}{Original period determinant and its constants}
% BEGIN PRESERVED INPUT 16_RPC_COMPLETE.tex SHA256 b1ed1c284ac1759bd396382b4dd2dde8e7a5946353583cc3d123ddc2a85662ba
\section{The original period quadric as a residue commutator}
\label{rpc:sec:actual-period-residue-commutator}

Retain $0<\delta<1/2$, $\gamma>2$, the ordered simple quartet
$\alpha_1=1/2+\delta+i\gamma$, $\alpha_2=1/2+\delta-i\gamma$,
$\alpha_3=1/2-\delta+i\gamma$, $\alpha_4=1/2-\delta-i\gamma$,
and $h(s)=\prod_a(s-\alpha_a)$. Put
\[
 w=s-\tfrac12,\quad r_a=\alpha_a-\tfrac12,\quad
 \beta=2(\gamma^2-\delta^2),\quad \eta=(\delta^2+\gamma^2)^2,
 \quad H(w)=w^4+\beta w^2+\eta.
 \tag{RPC1}
\]
Thus $r_4=-r_1$, $r_3=-r_2$, $r_2=\overline{r_1}$ and
$r_1^2-r_2^2=4i\delta\gamma$. No root or original coordinate is
discarded by this translation. Its coefficient isomorphism is
$J_c:\mathbb C[w]_{<4}\to E_h$, $J_cp=[p(s-1/2)]$; its inverse
is substitution $s=w+1/2$ in the unique original representative.
The primitive with its original constant is
\[
 \Phi_h(s)=A+\phi(w),\quad
 \phi(w)=\frac{w^5}{5}+\frac{\beta w^3}{3}+\eta w,\quad
 A=\Phi_h(1/2)=\frac1{160}+\frac\beta{24}+\frac\eta2.
 \tag{RPC2}
\]
The last equality is forced by $\Phi_h(0)=0$ and follows by evaluating
the displayed odd polynomial at $w=-1/2$.

Let $u\ne0$ carry the specified original branch of its logarithm.
Let $\Pi$ be the original period map on the rays and orientations
of PDE2, and let $F_{ja}=\zeta^{ja}-1$, $\zeta=e^{2\pi i/5}$,
$1\leq j,a\leq4$. Its inverse has entries
$(F^{-1})_{aj}=\zeta^{-ja}/5$: multiplying the two matrices and using
$\sum_{j=1}^4\zeta^{jb}=4$ for $5\mid b$ and $-1$ otherwise proves
this entry by entry. The literal centered coefficient-to-Fourier map is
\[
 P=F^{-1}\Pi J_c,\qquad P_0=e^{-A/u}P.
 \tag{RPC3}
\]
These are invertible by the complete original period determinant
PDE3--13. The contours for $P_0$ may be taken to be the rays based at
$w=0$ with the original five angles. Indeed translating the old rays
gives rays based at $-1/2$. Join corresponding finite endpoints and
close each contour at radius $R$. The integrand is entire. On the
bounded-length outer joining curves its exponent is
$-R^5/(5|u|)+O(R^4)$, uniformly along the joining curves, so the
outer integrals tend to zero. The common finite joining path occurs
with opposite coefficients in $-\ell_0+\ell_j$ and cancels exactly.
Cauchy's integral theorem therefore proves this contour identity,
including its signs and the factor $e^{A/u}$.

\subsection{A flat alternating pairing with all constants}

On $\mathbb C[w]/(H)$ define the bilinear, rather than sesquilinear,
pairing
\[
 \mathcal J(p,q)=\sum_{H(r)=0}\frac{p(r)q(-r)}{H'(r)}.
 \tag{RPC4}
\]
The roots are simple at the original packet. The equivalent polynomial
formula is the coefficient of $w^3$ in the remainder of $p(w)q(-w)$
modulo $H$. To prove it, use Lagrange interpolation: the coefficient
of $w^3$ in the idempotent at $r$ is $1/H'(r)$. This polynomial
formula extends the pairing to every complex value of $\beta,\eta$,
including repeated roots, without a division by a repeated-root value.
In the increasing frame $1,w,w^2,w^3$, multiplication by $w$ and
this pairing are the exact matrices
\[
 M=\begin{pmatrix}0&0&0&-\eta\\1&0&0&0\\0&1&0&-\beta\\0&0&1&0\end{pmatrix},
 \qquad
 J=\begin{pmatrix}0&0&0&-1\\0&0&1&0\\0&-1&0&\beta\\1&0&-\beta&0\end{pmatrix}.
 \tag{RPC5}
\]
For the last matrix, the coefficient functional is zero through degree
two, one in degree three, zero in degree four and six, and $-\beta$
in degree five. Inserting the sign $(-1)^j$ from the second argument
gives every displayed entry. Thus $J^{\mathsf T}=-J$, $\det J=1$,
and direct multiplication gives $M^{\mathsf T}J+JM=0$.

Allow the two original odd-phase coefficients to vary complexly,
solely to calculate the existing period map. The ray domination
$e^{-r^5/(5|u|)+O(r^3+r)}$ on compact parameter sets justifies all
coefficient derivatives and integrations by parts below. Ordinary
division of $w^{j+3}$ by $H$, followed by differential reduction,
gives
\[
 \partial_\eta P_0=\frac1uP_0M,\qquad
 \partial_\beta P_0=\frac1{3u}P_0(M^3-uN),\qquad
 N=\begin{pmatrix}0&0&1&0\\0&0&0&2\\0&0&0&0\\0&0&0&0\end{pmatrix}.
 \tag{RPC6}
\]
For clarity, the quotient derivatives for columns $j=0,1,2,3$ are
$0,0,1,2w$, respectively; hence this is the full correction matrix.
The two outer boundary terms vanish, and the two values at the common
origin cancel. No differential remainder was replaced by an arithmetic
remainder in RPC6.

There is also a direct nonvanishing proof on this entire two-parameter
family. The displayed matrices satisfy $\operatorname{Tr}M=0$,
$\operatorname{Tr}M^3=0$ and $\operatorname{Tr}N=0$. The first and
third follow from their diagonals; the second follows by multiplying
RPC5 three times, or from the vanishing third Newton sum of the even
monic polynomial $H$. Multilinearity of the determinant in its columns
then gives $\partial_\eta\det P_0=\partial_\beta\det P_0=0$ from
RPC6, without first assuming an inverse exists. At the origin the
literal Gamma ray integrals give the diagonal matrix $G_u$ defined in
RPC8, whose determinant is nonzero. Hence $P_0$ is invertible at every
point of $\mathbb C^2$, justifying the inverses throughout this family.

The matrix identity $M^{\mathsf T}J+JM=0$ also gives
$(M^3)^{\mathsf T}J+JM^3=0$. Direct multiplication of RPC5--6 gives
\[
 \partial_\beta J+\frac13(N^{\mathsf T}J+JN)=0.
 \tag{RPC7}
\]
In particular $\partial_\beta J$ has entries $(2,3)=1$, $(3,2)=-1$
when indices begin at zero, while $N^{\mathsf T}J+JN$ has the
opposite threefold entries. Differentiate the inverse matrices in
$P_0^{-\mathsf T}JP_0^{-1}$. RPC6--7 show that both its $\beta$ and
its $\eta$ derivatives vanish. This matrix is consequently constant
on the connected coefficient space $\mathbb C^2$.

Define the unchanged positive-Gamma ray constants
\[
 g_a=5^{a/5-1}e^{\pi ia/5}\Gamma(a/5),\qquad
 G=\operatorname{diag}(g_1,g_2,g_3,g_4),\quad
 G_u=\operatorname{diag}(u^{a/5}g_a)_{a=1}^4.
 \tag{RPC8}
\]
At $\beta=\eta=0$, literal substitution in each original ray integral
gives $P_0=G_u$. Every $g_a$ is nonzero, since its Gamma factor is
the integral of a strictly positive function. We have therefore proved
the complete flat-pairing identity
\[
 \begin{split}
 \Omega_0&=P_0^{-\mathsf T}JP_0^{-1}
       =G_u^{-\mathsf T}J\big|_{\beta=0}G_u^{-1},\\
 (\Omega_0)_{14}&=-\frac1{u g_1g_4},\qquad
 (\Omega_0)_{23}=\frac1{u g_2g_3},\qquad
 (\Omega_0)_{ba}=-(\Omega_0)_{ab},\\
 \Omega&=P^{-\mathsf T}JP^{-1}=e^{-2A/u}\Omega_0.
 \end{split}
 \tag{RPC9}
\]
All entries not listed are zero; the indices in RPC9 begin at one.
This formulation retains the Gamma constants directly and requires
no unevaluated normalization of a period or source mass. For
$D=\operatorname{diag}(\zeta,\zeta^2,\zeta^3,\zeta^4)$, the two
paired sums of weights are five. It follows entry by entry that
\[
 D^{\mathsf T}\Omega D=\Omega,
 \qquad D^{\mathsf T}\Omega_0D=\Omega_0.
 \tag{RPC10}
\]

\subsection{The full arithmetic unit and the exact quadratic form}

Retain the original $v_h=2\xi/h$ and its complete simple-root unit
values. Reflection and conjugation of this original function give
\[
 v_h(\alpha_1)=v_h(\alpha_4)=a,\qquad
 v_h(\alpha_2)=v_h(\alpha_3)=\overline a,\qquad a\ne0.
 \tag{RPC11}
\]
Indeed both numerator and denominator are invariant under $s\mapsto1-s$
and conjugate under $s\mapsto\overline s$, first away from the roots
and then by the removable values. The complete zero orders give
nonzero values at these simple roots. Write
\[
 c_3=\frac{a^{-2}-\overline a^{-2}}{4i\delta\gamma},\qquad
 c_1=\frac{a^{-2}+\overline a^{-2}}2-(\delta^2-\gamma^2)c_3,
 \quad \mathfrak r(w)=c_1w+c_3w^3.
 \tag{RPC12}
\]
Both coefficients are real, and $(c_1,c_3)\ne(0,0)$. Directly,
$c_1+c_3r_1^2=a^{-2}$ and $c_1+c_3r_2^2=\overline a^{-2}$;
the same equalities hold at the respective opposite roots. Thus
$\mathfrak r(M)=M\,v_h(M+1/2)^{-2}$ in the original finite algebra, with the
entire function evaluated by its complete original arithmetic unit.
RPC12 is an exact interpolation of this specific multiplier.

For $Q_0(t)=t_1t_4-t_2t_3$, let $J_{\rm prim}$ be the original
Lagrange coordinate isomorphism and set
\[
 B=F^{-1}\Pi UJ_{\rm prim},\qquad Q(x)=Q_0(B^{-1}x),\qquad
 \mathcal H=P\mathfrak r(M)P^{-1}.
 \tag{RPC13}
\]
If $p=P^{-1}x$, the primary components of $B^{-1}x$ are exactly
$p(r_1)/a,p(r_2)/\overline a,p(r_3)/\overline a,p(r_4)/a$.
Since $H'(r)=4r(r^2+\gamma^2-\delta^2)$, its values give
$r/H'(r)=1/(8i\delta\gamma)$ on roots one and four, and
$r/H'(r)=-1/(8i\delta\gamma)$ on roots two and three. Inserting
these four values into RPC4 proves the exact identity
\[
 \boxed{\quad Q(x)=4i\delta\gamma\,
       p^{\mathsf T}\mathfrak r(M)^{\mathsf T}Jp
       =4i\delta\gamma\,x^{\mathsf T}\mathcal H^{\mathsf T}\Omega x.\quad}
 \tag{RPC14}
\]
There is no conjugate transpose in this quadratic identity. Because
$\mathfrak r$ is odd, $\mathfrak r(M)^{\mathsf T}J+J\mathfrak r(M)=0$; hence the matrices in
RPC14 are symmetric. This also proves
$\mathcal H^{\mathsf T}\Omega+\Omega\mathcal H=0$.
The four eigenvalues of $\mathfrak r(M)$ are exactly
$r_1/a^2,\overline{r_1/a^2},-\overline{r_1/a^2},-r_1/a^2$.
They are nonzero, although they need not be distinct. Both $B$ and
$Q_0$ are invertible as linear and quadratic matrices, respectively,
so $Q$ always has rank four, including the possible eigenvalue collisions.

By RPC10, substituting $Dx$ into RPC14 transports $\mathcal H$ to
$D^{-1}\mathcal H D$. The resulting symmetric quadratic matrices
are equal if and only if the corresponding $\mathcal H$ matrices
are equal, because $\Omega$ is invertible. Consequently
\[
 \boxed{\quad Q(Dx)=Q(x)\quad\Longleftrightarrow\quad
 [\mathcal H,D]=0\quad\Longleftrightarrow\quad
 \mathcal H_{ab}=0\text{ for every }a\ne b.\quad}
 \tag{RPC15}
\]
This is an exact test in the actual original period matrix and unit.
The final equivalence uses the four distinct diagonal entries of $D$.
It makes no assumption that their off-diagonal tests are independent.

For completeness, a rank-four quadratic semi-invariant of cyclic
weight $j\ne0$ is impossible: its coefficient in row $j$ could pair
only with the missing weight zero, so that row is zero. Thus any
semi-invariant rank-four quadratic must have weight zero. Since five
is prime, the projective orbit of $Q$ under $D$ has size one or five,
and size one is exactly RPC15. Every failure of any displayed
off-diagonal test proves size five.

\subsection{An explicit original-parameter radius forcing orbit five}

Here we prove that RPC15 fails for all sufficiently large $|u|$,
with a completely specified finite radius for the given original
quartet and unit. The argument evaluates the original periods; it
does not choose or insert a replacement amplitude.

Set $b=\beta u^{-2/5}$, $c=\eta u^{-4/5}$, and let $T(b,c)$
be the Fourier period matrix for
$z^5/5+bz^3/3+cz$ at $u=1$, in columns $1,z,z^2,z^3$ and on
the five rays of angles $(\pi+2\pi j)/5$. Substitution
$w=u^{1/5}z$, with exactly the original branch, gives
\[
 P_0=T(b,c)S_u,\quad S_u=\operatorname{diag}(u^{a/5})_{a=1}^4,
 \quad S_uMS_u^{-1}=u^{1/5}M(b,c),\quad T(0,0)=G.
 \tag{RPC16}
\]
Every column includes its original differential factor. The matrix
$M(b,c)$ is RPC5 with $\beta,\eta$ replaced by $b,c$; in particular
$N_0:=M(0,0)$ has its three subdiagonal entries equal to one.

The following positive constants are explicit finite real integrals
or finite expressions in the original Gamma constants:
\[
 \begin{split}
 E_j&=\frac85\int_0^\infty r^j(r^3/3+r)
                    e^{-r^5/5+r^3/3+r}\,dr\quad(0\le j\le3),\\
 C_T&=2\left(\sum_{j=0}^3E_j^2\right)^{1/2},\qquad
 g_+=\max_a|g_a|,\quad g_-^{-1}=\max_a|g_a|^{-1},\\
 \kappa&=\max_{1\le a<4}|g_{a+1}/g_a|,\qquad
 K_0=GN_0G^{-1},\\
 C_K&=4C_Tg_-^{-1}+2g_+g_-^{-1}
                       +2g_+g_-^{-2}C_T,\\
 C_{K,3}&=C_K(3\kappa^2+3\kappa C_K+C_K^2),\qquad
 \epsilon_0=\min\{1,(2C_Tg_-^{-1})^{-1}\}.
 \end{split}
 \tag{RPC17}
\]
Here $g_-^{-2}$ means $(\max_a|g_a|^{-1})^2$. The negative fifth
degree exponent proves convergence of every integral. All norms in
the following estimates are the Euclidean operator norm, with the
Frobenius norm used explicitly as an upper bound where stated.

Write $\epsilon=|b|+|c|$. When $|b|,|c|\le1$, on each ray
\[
 |e^{bz^3/3+cz}-1|
 \le \epsilon(r^3/3+r)e^{r^3/3+r}.
\]
Each contour has two rays, and each entry of $F^{-1}$ has absolute
value $1/5$ with four summands. The entry in column $j$ of
$T-G$ is therefore at most $E_j\epsilon$. Summing the squares
over all four rows proves
\[
 \|T-G\|\le\|T-G\|_{\rm F}\le C_T\epsilon.
 \tag{RPC18}
\]
For $\epsilon\le\epsilon_0$, the convergent inverse series for
$I+G^{-1}(T-G)$ gives
$\|T^{-1}\|\le2g_-^{-1}$ and
$\|T^{-1}-G^{-1}\|\le2g_-^{-2}C_T\epsilon$.
Also $\|M(b,c)-N_0\|\le\epsilon$, $\|M(b,c)\|\le2$,
and $\|N_0\|=1$. Subtract the products
$K:=T M(b,c)T^{-1}$ and $K_0$ in three terms, changing the left
factor, then the middle factor, then the inverse factor. The preceding
inequalities give, with the complete RPC17 constants,
\[
 \|K-K_0\|\le C_K\epsilon,\quad
 \|K\|\le\kappa+C_K\epsilon,\quad
 \|K^3-K_0^3\|\le C_{K,3}\epsilon.
 \tag{RPC19}
\]
For the last assertion use the exact noncommutative identity
$K^3-K_0^3=(K-K_0)K^2+K_0(K-K_0)K+K_0^2(K-K_0)$.
Its norm is at most $C_K\epsilon$ times
$(\kappa+C_K)^2+\kappa(\kappa+C_K)+\kappa^2$, as claimed.

The original multiplication and full unit in RPC13 now give exactly
\[
 \mathcal H=c_1u^{1/5}K+c_3u^{3/5}K^3,
 \qquad (K_0^3)_{41}=g_4/g_1,\quad (K_0)_{21}=g_2/g_1.
 \tag{RPC20}
\]
The scalar $e^{A/u}$ cancels from this conjugation, while it remains
in the quadratic form through RPC9. Put $B_*=\beta+\eta>0$.
For $|u|\ge1$, $\epsilon\le B_*|u|^{-2/5}$.

If $c_3\ne0$, define the completely specified radius
\[
 \begin{split}
 T_*&=\max\left\{1,\frac{B_*}{\epsilon_0},
 \frac{2\bigl(|c_3|C_{K,3}B_*+|c_1|(\kappa+C_K)\bigr)}
 {|c_3g_4/g_1|}\right\},\qquad R_*=T_*^{5/2}.
 \end{split}
 \tag{RPC21}
\]
For every $u$ on its original branch with $|u|\ge R_*$,
RPC19--20 imply
\[
 \left|u^{-3/5}\mathcal H_{41}-c_3g_4/g_1\right|
 \le |u|^{-2/5}\bigl(|c_3|C_{K,3}B_*
                           +|c_1|(\kappa+C_K)\bigr)
 \le\tfrac12|c_3g_4/g_1|.
 \tag{RPC22}
\]
In particular $\mathcal H_{41}\ne0$. In the quadratic matrix
RPC14 this is the nonzero $(1,1)$ entry
$4i\delta\gamma\,\mathcal H_{41}\Omega_{41}$, whose cyclic
weight is two.

If $c_3=0$, then $c_1\ne0$. In this case define
\[
 T_* =\max\left\{1,\frac{B_*}{\epsilon_0},
                   \frac{2C_KB_*}{|g_2/g_1|}\right\},
 \qquad R_*=T_*^{5/2}.
 \tag{RPC23}
\]
For every $|u|\ge R_*$, the same estimates give
\[
 \left|u^{-1/5}\mathcal H_{21}-c_1g_2/g_1\right|
       \le |c_1|C_KB_*|u|^{-2/5}
       \le\tfrac12|c_1g_2/g_1|.
 \tag{RPC24}
\]
Thus $\mathcal H_{21}\ne0$. The $(1,3)$ entry of the quadratic
matrix is then $4i\delta\gamma\mathcal H_{21}\Omega_{23}$,
which has cyclic weight four; symmetry retains its $(3,1)$ partner.
Both cases prove the unconditional conclusion for the actual original
parameters and the explicit radius belonging to those parameters:
\[
 \boxed{\quad |u|\ge R_*\quad\Longrightarrow\quad
       \#\{[Q(D^jx)]:0\le j<5\}=5.\quad}
 \tag{RPC25}
\]
The bracket denotes the projective class of the quadratic polynomial.
RPC21 or RPC23 supplies its radius for every nonzero original unit;
neither case assumes a condition on an uncomputed period coefficient.

Finally, the original periods are holomorphic on the connected
logarithm cover of $\mathbb C^*$, and on each branch domain: locally the contours may
be held in common decay sectors and differentiated under the dominated
integral, and contour deformation identifies these local definitions.
Their determinant never vanishes. Hence all entries of $\mathcal H$
are holomorphic there. RPC22 or RPC24 proves that its respective
decisive entry is not identically zero on the logarithm cover and,
by the identity theorem on that connected cover, on any open branch domain. The zeros
of that one entry are isolated, by its convergent local power series.
The exact exceptional set
\[
 \mathcal E=\{u:\mathcal H_{ab}(u)=0\text{ for all }a\ne b\}
 \tag{RPC26}
\]
is a subset of those isolated zeros and satisfies $|u|<R_*$.
Thus it has no accumulation inside the branch domain; accumulation
at $u=0$ or a branch boundary has not been excluded. RPC15 is the
full calculation specifying this set. No assertion that it is empty
at arbitrary finite complex $u$ is made or used.

% END RX INLINE 27

% END PRESERVED INPUT 16_RPC_COMPLETE.tex
\endgroup

\clearpage
\begingroup

\def\Fq{\mathbb F_q}\def\Ql{\mathbb Q_\ell}\def\Spec{\operatorname{Spec}}
\setcounter{equation}{0}
\renewcommand{\theequation}{D33.\arabic{equation}}
\section*{Deligne: the general mixed-weight direct-image theorem}
\addcontentsline{toc}{section}{Deligne: the general mixed-weight direct-image theorem}
The following complete supplied section is from P. Deligne, \emph{La conjecture de Weil. II}, section 3.3, with proof. The full source is retained in volume 16 of the programme source collection. See \cite{et:Deligne}.
% BEGIN PRESERVED INPUT 17_DELIGNE_3_3_COMPLETE.tex SHA256 686e81d152b2ff710bc3073013c242220fa01d3b91c27cb2a7f9b9c052580074
\subsection*{(3.3) The general case}

\begin{theorem}[3.3.1]
Let $f:X\to Y$ be a morphism of schemes of finite type over $\mathbb Z$, and let $\mathcal F$ be a sheaf on $X$. If $\mathcal F$ is mixed of weights $\le n$, then, for each $i$, $R^if_{!}\mathcal F$ is mixed of weights $\le n+i$.
\end{theorem}

The proof uses successive devissages. Exact sequences reduce the assertion in the coefficient sheaf; an open-closed decomposition reduces it in $X$; base change to an open part and its complement reduces it in $Y$; Leray spectral sequences give transitivity. Universal homeomorphisms do not affect etale cohomology. Finally, morphisms of relative dimension zero preserve pointwise purity and have no higher direct images. These reductions reduce to the case where $f$ has relative dimension one, $\mathcal F$ is lisse and pointwise pure, and $X$ is the complement in a smooth projective relative curve of a divisor finite etale over $Y$, with tame ramification along the divisor.

On each geometric fiber, Theorem (3.2.3) gives purity for $R^if_{!}j_{*}\mathcal F$. The local-monodromy theorem (1.8.4) gives the weights of the graded pieces of the local monodromy filtration on $i^*j_{*}\mathcal F$. The finite part is handled by the relative-dimension-zero case, and the exact sequence
\[
0\to j_{!}\mathcal F\to j_{*}\mathcal F\to i_{*}i^*j_{*}\mathcal F\to0
\]
finishes the proof.

\paragraph{(3.3.2).}
Lower bounds for weights are obtained from the $p$-adic estimates of SGA 7 XXI, 5. A sheaf is called integral if all Frobenius eigenvalues are algebraic integers. The weights of a mixed integral sheaf are necessarily non-negative, since the norm of a pure algebraic integer of weight $n$ has absolute value $q^{dn/2}$.

\begin{corollary}[3.3.3]
Let $f:X\to Y$ be of finite type over $\mathbb Z$, and let $\mathcal F$ be an integral sheaf on $X$. If $\mathcal F$ is mixed of weights $\le n$, then $R^if_{!}\mathcal F$ is mixed of weights between $0$ and $n+i$. If $i$ is larger than the maximum dimension $d$ of the fibers, the weights lie between $2(i-d)$ and $n+i$.
\end{corollary}

For $Y=\operatorname{Spec}\Fq$ this becomes:

\begin{corollary}[3.3.4]
Let $X_0$ be of finite type over $\Fq$, and let $\mathcal F_0$ be mixed of weights $\le n$ on $X_0$. Then $H^i_c(X,\mathcal F)$ is mixed of weights $\le n+i$. If $\mathcal F_0$ is integral, the weights are non-negative; if moreover $i>d=\dim X$, they are $\ge2(i-d)$.
\end{corollary}

\begin{corollary}[3.3.5]
If $X_0$ is smooth of finite type over $\Fq$ and $\mathcal F_0$ is lisse mixed of weights $\ge n$, then $H^i(X,\mathcal F)$ is mixed of weights $\ge n+i$.
\end{corollary}
This is Poincare duality applied to the dual sheaf. Hence:

\begin{corollary}[3.3.6]
If $X_0$ is smooth of finite type over $\Fq$ and $\mathcal F_0$ is lisse pure of weight $n$, then the image of
\[
H^i_c(X,\mathcal F)\longrightarrow H^i(X,\mathcal F)
\]
is pure of weight $n+i$.
\end{corollary}

\paragraph{(3.3.7)--(3.3.9).}
If an algebraic number $\alpha$ is pure of weight $n$ relative to $q$, normalize $p$-adic valuations by $v(q)=1$ and attach to $\alpha$ the pair $(v(\alpha),v(q^n\alpha^{-1}))$. For integral eigenvalues both coordinates are non-negative and sum to $n$. For the constant sheaf, the preceding bounds localize these pairs in the corresponding triangles in the $(r,s)$-plane. In particular, for $X_0$ proper and smooth one obtains again the main theorem of I, with projective replaced by proper:
\[
\det(1-Ft,H^i(X,\Ql))\in\mathbb Z[t]
\]
is independent of $\ell\ne p$, and all complex roots have absolute value $q^{i/2}$.

\paragraph{(3.3.10)--(3.3.11).}
The theorem remains valid with ``mixed'' replaced by ``$\iota$-mixed'' and real-valued weights. The only weights in $R^if_{!}\mathcal F$ are those $\le\beta+i$ congruent modulo $\mathbb Z$ to a weight already occurring in $\mathcal F$. The smoothness assumption is used only to invoke Poincare duality, and may therefore be replaced by the corresponding rational-homology-manifold condition
\[
Ra^!\Ql=\Ql(n)[2n]
\]
for the structural morphism $a:X_0\to\Spec\Fq$.


% END PRESERVED INPUT 17_DELIGNE_3_3_COMPLETE.tex
\endgroup

\clearpage
\begingroup


\setcounter{equation}{0}
\renewcommand{\theequation}{LOC.\arabic{equation}}
\section*{Complete attained-image localization proof}
\addcontentsline{toc}{section}{Complete attained-image localization proof}
% BEGIN PRESERVED INPUT provider_CONSOLIDATED_LOCALIZATION.tex SHA256 58c088b7934ef8af5cab6c615c6dcbcd9cd1b69a5119bea293ed7b51707e64ad
\section{Signed arithmetic baseline: consolidation and localization}\label{loc:signed-arithmetic-baseline-consolidation-and-localization}

18 September 2026. A written mathematical continuation on the original fixed-packet domain. The two incoming continuations are retained verbatim in \texttt{inputs/}. This note identifies their overlap exactly, preserves the stronger statements, and proves additional endpoint and degree-band estimates. It does not claim that the full signed arithmetic coefficient has been evaluated.

\subsection{Result reached}\label{loc:result-reached}

Write the original signed baseline as

\[
\delta_k=\mathcal B_k[m_k]-\mathcal B_k[\sigma],\qquad
\mathcal B_k[\nu]=\log\det G_{q-1}[\nu]+\log\det G_q[\nu]
-\log\det G_{2q-1}[\nu]-\log\det G_{2q}[\nu].
\]

Let \(C_k=[-kb,kb]\), with \(b\) the original boundary-tilted mean. Define the outer-only density, including its exact central value, by

\[
\omega_k(y)=
\begin{cases}\sigma(y),&y\in C_k,\\
\mathfrak M^{-k}m_k(y),&y\notin C_k.
\end{cases}
\]

The new decomposition, with all four original signs, is

\[
\boxed{\delta_k=\delta_k^{\rm out}+\widehat\Phi_{q,k}+\mathfrak e_k,
\qquad \delta_k^{\rm out}=\mathcal B_k[\omega_k]-\mathcal B_k[\sigma],}
\tag{R1}
\]

\[
\boxed{\widehat\Phi_{q,k}\ge0,\quad
|\mathfrak e_k|\le E_k=O_h(k(\log q)^2)=o_h(q).}
\tag{R2}
\]

Here \(\widehat\Phi_{q,k}\) is given in Section 8 by four explicit Gamma determinants of size

\[
r=O_h(k\log q)=o(q).
\]

Its entries use only the fixed one-factor tilted mass \(M(\theta)\), the original Gamma polynomials, and finite integrals on \(C_k\). The actual convolution density outside \(C_k\) remains solely in \(\delta_k^{\rm out}\). No original arithmetic kernel is substituted by rank or by an asserted isometry.

There is a further bound on the central endpoint contribution

\[
\delta_k^{\rm cen,end}:=\mathcal B_k[m_k]-\mathcal B_k[\omega_k]:
\]

\[
\boxed{\delta_k^{\rm cen,end}=O_h(k^2),\qquad
\liminf\frac{\delta_k^{\rm cen,end}}{k^2}\ge0.}
\tag{R3}
\]

Consequently its coefficient at the \(q\log k\) scale is zero, with an \(O_h(q)\) remainder, on every fixed multiplicity stratum. On each fixed stratum \(m_0\ge2\), this entire central endpoint contribution is \(o_h(q)\). On the simple stratum \(m_0=1\), its order-\(q\) value is reduced to \(\widehat\Phi_{q,k}\) with an \(o(q)\) error, but that limiting value is not calculated in this note.

The outer term in (R1) is still unevaluated at the requested scales. In particular this note does not assign the full baseline's \(q\log k\) coefficient or its order-\(q\) coefficient.

\begin{center}\rule{0.5\linewidth}{0.5pt}\end{center}

\subsection{1. Original source, exact coordinates, and consolidation}\label{loc:original-source-exact-coordinates-and-consolidation}

Retain the original stipulated quartet, its complete multiplicity, its full arithmetic amplitude, and its original tensor orders:

\[
\rho=\tfrac12+\delta+i\gamma,\quad 0<\delta<\tfrac12,\quad\gamma>2,
\qquad h(s)=\prod_{z\in\{\rho,\bar\rho,1-\rho,1-\bar\rho\}}(s-z)^{m_0},
\]

\[
v_h=2\xi/h,\qquad w_h(y)=|v_h(1/2+iy)|^2/(2\pi),\qquad m_k=w_h^{*k},
\]

\[
k=4l+1\ge5,\quad e=1+k(m_0-1),\quad q=e(k+1)^2,\quad c=k/2,
\]

\[
\chi_k(S)=\prod_{a,b_0=0}^k
\bigl(S-c-(2a-k)\delta-i(2b_0-k)\gamma\bigr)^e.
\]

This retains the programme's stipulated packet; it does not supply or assert an actual numerical off-critical zero.

The reference source is

\[
\sigma(y)=\frac{|\Gamma(1/4+iy/2)|^2}{2\pi},\quad
\sigma(\mathbb R)=\sqrt{2\pi},\quad\alpha=\pi/2.
\]

The original change of polynomial coordinates is

\[
\mathcal U_N:P(S)\longmapsto P(c+iy),\qquad
Q(y)=i^{-q}\chi_k(c+iy).
\tag{C1}
\]

Its inverse is substitution \(y=(S-c)/i\); its coefficient determinant has modulus one. Moreover

\[
\mathcal U_N(\chi_kP)=i^q Q\,\mathcal U_{N-q}P.
\]

Thus it carries the complete relation ideal, every remainder class and the degree cutoffs to their displayed real-coordinate versions. The phase has modulus one. We use the corresponding source norm, not an independent norm on the transported coefficients. For monic norms, the additional phase \(i^{-\deg P}\) restores real-coordinate monicity exactly. The original full-unit inclusion is transported by this coefficient isomorphism; its arithmetic coefficients are not set to one.

Let \(H_N[\nu]\) be the Gram of \(1,y,\ldots,y^N\), \(J_N\) the remainder map modulo \(Q\), and

\[
G_N[\nu]=(J_NH_N[\nu]^{-1}J_N^*)^{-1}.
\]

The original monic source/relation frame gives

\[
\det G_N[\nu]=\frac{\det H_N[\nu]}{\det H_N^{\rm rel}[\nu]},
\tag{C2}
\]

where the relation frame is \(Q,yQ,\ldots,y^{N-q}Q\). Empty determinants are one. Its frame determinant is one; both cross blocks remain in the Schur complement. Equivalently,

\[
\mathcal B_k[\nu]=\sum_{j=0}^q c_j
\log\frac{\nu_j[\nu]}{\gamma_{q+j}[\nu]},\quad
c_0=c_q=1,\quad c_j=2\ (0<j<q),\quad \sum c_j=2q.
\tag{C3}
\]

Here \(\nu_j\) and \(\gamma_{q+j}\) are the attained monic relation and source norms, respectively. We use (C2) and (C3), rather than redefining the original baseline.

\subsubsection{Exact dictionary between the incoming calculations}\label{loc:exact-dictionary-between-the-incoming-calculations}

Use

\[
M(\theta)=\int e^{\theta y}w_h(y)\,dy,\quad L=\log M,
\quad \mu=L',\quad V=L'',\quad
\mathfrak M=M(\alpha),\quad b=\mu(\alpha),
\]

and let \(\theta(t)\) invert \(\mu\) on \([-\alpha,\alpha]\). For \(0\le t\le b\), define

\[
D(t)=L(\alpha)-L(\theta(t))-(\alpha-\theta(t))t.
\tag{C4}
\]

The first continuation's \(\ell_h,a_h,v_h^{\rm tilt},\mu_h\) are respectively the second continuation's \(L,\mu,V,b\). Their rate functions satisfy the literal equality

\[
\boxed{F_h^{\rm cen}(t)=-D(|t|).}
\tag{C5}
\]

The integral coefficients are also equal, without a limiting identification:

\[
\boxed{\mathscr C_h=\mathfrak I_h^{\rm cen}
=\int_0^\alpha\mu(\theta)^2\,d\theta=:I_h,}
\]

\[
\boxed{\mathscr B_h=\mathfrak J_h^{\rm cen}
=\int_0^\alpha V(\theta)
\log\frac{\mu(\theta)}{4\pi V(\theta)}\,d\theta=:J_h.}
\tag{C6}
\]

The second integral converges at zero because \(\mu(\theta)\) is comparable to \(\theta\) there. Both incoming notes therefore calculate the same central density profile and the same two coefficients. The first retains a stronger explicit remainder and a source-to-quotient comparison; the second supplies a complete standalone proof and the sharper relation-kernel weight count \(2q\) instead of the valid larger count \(2q+1\). All are retained below, once.

\subsubsection{Consolidated source estimate}\label{loc:consolidated-source-estimate}

The local estimate in the incoming ACW10--11 is

\[
p_\theta^{*k}(k\mu(\theta))
=\frac{1+\varepsilon_{k,\theta}}{\sqrt{2\pi kV(\theta)}},
\quad p_\theta(y)=\frac{e^{\theta y}w_h(y)}{M(\theta)},
\quad |\varepsilon_{k,\theta}|\le\epsilon_k=O_h(k^{-1/8}).
\tag{C7}
\]

The exact inverse accounting is

\[
m_k(x)=M(\theta)^k e^{-\theta x}p_\theta^{*k}(x).
\tag{C8}
\]

This follows from the original addition map \(\sum y_j=x\). It retains the mass \(M(\theta)^k\). At the auxiliary Hilbert-space level the isometry is
\(F\mapsto M(\theta)^{1/2}e^{-\theta y/2}F\) from \(L^2(w_hdy)\) to \(L^2(p_\theta dy)\); it is not asserted to preserve polynomial degree.

Put \(c_\sigma(y)=\sigma(y)e^{\alpha|y|}\sqrt{1+|y|}\). On \(|t|\le b\),

\[
\ell_k(kt):=\log\frac{m_k(kt)}{\mathfrak M^k\sigma(kt)}
=-kD(|t|)+\tfrac12\log(|t|+k^{-1})
-\tfrac12\log(2\pi V(\theta(t)))-\log c_\sigma(kt)
+\log(1+\varepsilon_{k,\theta(t)}).
\tag{C9}
\]

The full origin term is present. With \(L_k^{\rm loc}=-\log(1-\epsilon_k)\) on \(\epsilon_k\le1/2\), the consolidated finite bound is

\[
\left|\int_{-kb}^{kb}\ell_k(y)dy+k^2I_h-kJ_h\right|
\le 2kbL_k^{\rm loc}+(kb+1)\log(kb+1)-kb\log(kb)
+2C_\Gamma^0+4\log^+(kb),
\tag{C10}
\]

\[
C_\Gamma^0=\max\{|\log(c_\Gamma/\sqrt2)|,
|\log(C_\Gamma/\sqrt2)|\}.
\]

In particular the error in (C10) is \(O_h(k^{7/8}+\log k)\), which implies the second note's \(o(k)\) remainder. To recover (C10) from the standalone note, use the complex Stirling bound
\(|R_\Gamma(z)|\le1/(6|z|)\) for \(\Re z>0\) in the principal term at \(z=1/4+iy/2\). Direct substitution gives

\[
|\log c_\sigma(y)-\tfrac12\log2|\le2/y\quad(y\ge1).
\]

The original Gamma envelopes control \([0,1]\). The exact integral
\(k\int_0^b\log(1+(kt)^{-1})dt
=(kb+1)\log(kb+1)-kb\log(kb)\)
and (C9) prove the bound. The rate-area identity is
\(2\int_0^bD(t)dt=I_h\), proved by Fubini or integration by parts using \(\mu'=V\).

The source-coordinate approximation in the first note is retained as well. Define \(\bar m_k\) from (C7)--(C8) by omitting \(1+\varepsilon\) on \(C_k\), and by setting \(\bar m_k=m_k\) outside \(C_k\). The identity on every original affine remainder fibre gives

\[
(1-\epsilon_k)G_N[\bar m_k]\preceq G_N[m_k]
\preceq(1+\epsilon_k)G_N[\bar m_k],
\]

\[
|\mathcal B_k[m_k]-\mathcal B_k[\bar m_k]|
\le2q\log\frac{1+\epsilon_k}{1-\epsilon_k}=o(q).
\tag{C11}
\]

(C10) is not assigned as a determinant coefficient. The additional work below controls an actual endpoint determinant response.

\begin{center}\rule{0.5\linewidth}{0.5pt}\end{center}

\subsection{2. Retained source constants and local-error formula}\label{loc:retained-source-constants-and-local-error-formula}

The only arithmetic envelope inputs used are RTM1--3 and the original three-factor lower bound:

\[
\begin{gathered}
e^{\alpha|y|}w_h(y)\le C(1+|y|)^{-p},\quad
p=8m_0-9/2,\\
w_h^{*3}(y)\ge c_h e^{-\alpha|y|}(1+|y|)^{-B},
\quad B=42+8m_0,\quad M_*=B/2,
\end{gathered}\tag{L1}
\]

\[
c_\Gamma e^{-\alpha|y|}(1+|y|)^{-1/2}\le\sigma(y)
\le C_\Gamma e^{-\alpha|y|}(1+|y|)^{-1/2},
\quad c_\Gamma=\sqrt2e^{-7/6},\quad C_\Gamma=\sqrt{2\sqrt5}e^{2/3}.
\]

The global upper convolution envelope is

\[
m_k(y)\le\frac C{c_\Gamma}k^{p+1}\mathfrak M^{k-1}\sigma(y).
\tag{L2}
\]

For clarity, the finite local-error constants from the standalone note are reproduced. Let \(\eta=1/4\), \(M_0=M(0)\), and

\[
V_* =\frac{2C}{M_0}{\rm B}(3,p-3),\quad
H_* =2^{1+\eta}\left[\frac{2C}{M_0}{\rm B}(3+\eta,p-3-\eta)+b^{2+\eta}\right],
\]

\[
a_* =\min\left\{\tfrac14,\frac{c_he^{-2\alpha}2^{-B}}{\mathfrak M^3}\right\},
\quad\beta=2a_*,\quad v_*=\beta/9,
\]

\[
\begin{gathered}
\xi_0=\min\left\{1,(2V_*)^{-1/2},
(v_*/(16H_*))^{1/\eta}\right\},\quad
r_*=(1-\beta\xi_0^2/7)^{1/3},\\
C_2=\frac{4\pi(C/M_0)^2}{2p-1},\quad C_0=4H_*+V_*^2/8.
\end{gathered}
\]

For an integer \(n\ge2\), put \(a_n=v_*(n-1)/4\) and

\[
A_n^{\rm loc}=\frac1{2\pi}\left[
C_0n\,\Gamma((3+\eta)/2)a_n^{-(3+\eta)/2}
+C_2r_*^{n-2}+
\frac{2e^{-nv_*\xi_0^2/2}}{nv_*\xi_0}\right],
\quad\epsilon_n=\sqrt{2\pi nV_*}\,A_n^{\rm loc}.
\tag{L3}
\]

These retain \(v_*\le V(\theta)\le V_*\). They are local-limit constants, not the arithmetic upper-envelope symbol in (L2).

The Fourier proof of (C7) establishes slightly more than its displayed mean value. The same absolute error holds at every shifted evaluation point:

\[
\boxed{\left|p_\theta^{*n}(n\mu(\theta)+x)
-\frac{e^{-x^2/(2nV(\theta))}}{\sqrt{2\pi nV(\theta)}}\right|
\le A_n^{\rm loc},\quad x\in\mathbb R.}
\tag{L4}
\]

Indeed, Fourier inversion inserts the factor \(e^{-i\xi x}\), of modulus one, into the very same integrable difference estimated in ACW6--10. No additional moment or large-deviation hypothesis is needed.

\subsection{3. Closing the transition gap and extracting the full central rate}\label{loc:closing-the-transition-gap-and-extracting-the-full-central-rate}

Define

\[
W_k(y)=\begin{cases}kD(|y|/k),&|y|\le kb,\\0,&|y|>kb,
\end{cases}\quad
D_0=D(0)=L(\alpha)-L(0)>0,\quad\Lambda_k=kD_0.
\tag{L5}
\]

Then \(0\le W_k\le\Lambda_k\). Define the exact remaining factor

\[
\psi_k(y)=\frac{m_k(y)}{\mathfrak M^k\sigma(y)}e^{W_k(y)}.
\quad\text{Thus}\quad
\boxed{\mathfrak M^{-k}m_k=\psi_k e^{-W_k}\sigma.}
\tag{L6}
\]

This factorization changes neither the source nor the packet. Its significance is the new whole-line polynomial bound proved next.

Set \(n=k-3\), \(V_b=V(\alpha)\), and
\(a_G=e^{-1/2}/(2\sqrt{2\pi})\). Work on the explicit cofinal range

\[
\epsilon_k\le1/2,\qquad
\epsilon_{k-3}\le e^{-1/2}/2,\qquad
k\ge V_b/b^2.
\tag{L7}
\]

By (L4), the tilted convolution places mass at least \(a_G\) in

\[
I_n=[nb-\sqrt{nV_b},nb].
\]

On this interval the Gaussian density is at least
\(e^{-1/2}/\sqrt{2\pi nV_b}\), and the error is at most half this value. Integrating over the length \(\sqrt{nV_b}\) proves the mass bound.

For every \(y\ge kb\) and \(z\in I_n\),

\[
0<y-z\le y+\sqrt{nV_b}\le2y.
\]

Apply the original three-factor lower bound within the exact convolution and retain the full tilted mass. This gives

\[
\boxed{m_k(y)\ge c_h a_G\mathfrak M^{k-3}e^{-\alpha y}(1+2y)^{-B},
\qquad y\ge kb.}
\tag{L8}
\]

Evenness gives the negative half-line. In contrast to the old positive-width guard beyond the boundary-tilted mean, this bound starts at \(|y|=kb\) itself. In particular there is no untreated interval between the saddle window and the retained tail.

Divide (L8) by the upper Gamma envelope and use
\(1+2y\le\sqrt5\sqrt{1+y^2}\). Define

\[
c_{\rm tail}=\frac{c_h a_G}{\mathfrak M^3C_\Gamma 5^{M_*}},
\quad
c_{\rm core}=\frac1{2\sqrt{2\pi V_*}C_\Gamma},
\quad
C_{\rm core}=\frac{3\sqrt{b+1}}{2\sqrt{2\pi v_*}c_\Gamma}.
\]

From (C9), on the core,
\(c_{\rm core}k^{-1/2}\le\psi_k\le C_{\rm core}\). Set

\[
c_- =\min\{1,c_{\rm tail},c_{\rm core}\},\qquad
c_+=\max\{1,C_{\rm core},C/(c_\Gamma\mathfrak M)\}.
\]

Then the exact whole-line result is

\[
\boxed{c_-k^{-1/2}(1+y^2)^{-M_*}\le\psi_k(y)
\le c_+k^{p+1}.}
\tag{L9}
\]

The factor \(\omega_k/\sigma\) satisfies the same bounds: it is one on \(C_k\), and equals \(\psi_k\) off \(C_k\). This statement retains the actual zeta-dependent density on the whole outer line. The original exponential-in-\(k\) variation is now in the explicit factor \(e^{-W_k}\), not hidden in a comparison constant.

\subsection{4. Polynomial comparison also holds after inserting the rate}\label{loc:polynomial-comparison-also-holds-after-inserting-the-rate}

Write \(d=2q\) and \(d\tau_t=e^{-tW_k}\sigma(y)dy\), \(0\le t\le1\). Impose the explicit eventual guard

\[
\Lambda_k\le d\log16.
\tag{L10}
\]

For every polynomial \(P\) of degree at most \(d\),

\[
\boxed{\int y^2|P|^2d\tau_t\le272(d+1)^2\|P\|_{\tau_t}^2.}
\tag{L11}
\]

Here is a proof retaining the whole tail. The original Gamma recurrence bounds multiplication by \(y\) on degree \(j\) by \(2(j+1)\). Hence

\[
\|y^{d+1}P\|_\sigma^2\le[4(d+1)]^{2d+2}\|P\|_\sigma^2.
\]

On \(|y|>16(d+1)\), this bounds the second-moment tail by

\[
16(d+1)^2\,16^{-d}\|P\|_\sigma^2
\le16(d+1)^2e^{\Lambda_k}16^{-d}\|P\|_{\tau_t}^2.
\]

The inner part is at most \(256(d+1)^2\|P\|_{\tau_t}^2\). Guard (L10) proves (L11). This argument does not infer a tail estimate merely by dropping the core suppression.

Jensen's inequality for \(x\mapsto(1+x)^{-M_*}\), applied to the measure \(|P|^2d\tau_t/\|P\|_{\tau_t}^2\), now proves

\[
L_k\|P\|_{\tau_t}^2\le\|P\|_{\psi_k\tau_t}^2
\le U_k\|P\|_{\tau_t}^2,
\]

\[
L_k\|P\|_{\tau_t}^2\le\|P\|_{(\omega_k/\sigma)\tau_t}^2
\le U_k\|P\|_{\tau_t}^2,
\tag{L12}
\]

where

\[
L_k=c_-k^{-1/2}[1+272(d+1)^2]^{-M_*},\quad
U_k=c_+k^{p+1},\quad\kappa_k=U_k/L_k\ge1.
\tag{L13}
\]

Thus \(\log\kappa_k=O_h(\log q)\). The same scalar inequalities hold on every fixed restriction and every attained quotient of these polynomial spaces: apply them to every representative in its specified affine fibre before taking the minimum. This is the comparison used in the next section; no metric equality is asserted.

\subsection{5. A fixed Taylor-jet kernel sees exponentially little central mass}\label{loc:a-fixed-taylor-jet-kernel-sees-exponentially-little-central-mass}

Let \(b_j\) be the original Gamma monic polynomials, and

\[
\gamma_j=\sqrt{2\pi}\,j!(1/2)_j
=\sqrt{2\pi}(2j)!/4^j,\quad \varphi_j=b_j/\sqrt{\gamma_j}.
\]

Their exact generating function, equivalently derived by summing their recurrence, is

\[
\sum_{j\ge0}\frac{b_j(z)}{j!}t^j
=(1+t^2)^{-1/4}\exp(z\arctan t).
\tag{L14}
\]

This is the original-coordinate Meixner--Pollaczek generating function; it is also DLMF 18.23.7 after \(\lambda=1/4,x=z/2,\varphi=\pi/2\).

Cauchy's estimate on \(|t|=(d+1)/(d+2)\), together with
\(|\arctan t|\le\frac12\log(2(d+2))\), gives, for \(|z|\le Z\),

\[
\sum_{j=0}^d|\varphi_j(z)|^2
\le\frac{e^2(d+1)^2\sqrt{d+2}}{\sqrt{2\pi}}
[2(d+2)]^Z.
\tag{L15}
\]

The factorial estimate used here is
\((j!)^2/\gamma_j\le(2j+1)/\sqrt{2\pi}\), and
\(\sum_{j=0}^d(2j+1)=(d+1)^2\). This preserves the original Gamma mass.

Let \(K=kb\), and let \(\mathcal J_rP\) be the first \(r\) Taylor coefficients at zero. Its kernel is the literal subspace

\[
Z_{N,r}=y^r\mathcal P_{N-r}\subset\mathcal P_N.
\tag{L16}
\]

For \(P\in Z_{N,r}\), \(N\le d\), apply Cauchy's coefficient formula on \(|z|=2K\) and sum the geometric tail on \(|y|\le K\). Equation (L15) and the full mass \(\sqrt{2\pi}\) give

\[
\boxed{\int_{C_k}|P|^2d\sigma\le a_{k,r}\|P\|_\sigma^2,\qquad
 a_{k,r}=4^{1-r}\mathcal A_{q,k},}
\tag{L17}
\]

\[
\mathcal A_{q,k}=e^2(d+1)^2\sqrt{d+2}[2(d+2)]^{2K}.
\]

There is also an operator version. Let \(B_N\) restrict a polynomial in the Gamma orthonormal frame to \(L^2(C_k,\sigma dy)\), and let \(\widetilde B_{N,r}\) first truncate its Taylor series below degree \(r\), then restrict. The same calculation proves

\[
\|B_N-\widetilde B_{N,r}\|\le\sqrt{a_{k,r}},\qquad
\operatorname{rank}\widetilde B_{N,r}\le r.
\tag{L18}
\]

Choose the explicit integer

\[
\boxed{r=\left\lceil 1+
\frac{\log\mathcal A_{q,k}+2\Lambda_k+2\log\kappa_k
+10\log q+\log256}{\log4}\right\rceil.}
\tag{L19}
\]

Use its explicit eventual guard \(r\le q\). Then

\[
a:=a_{k,r}\le\frac{e^{-2\Lambda_k}}{256\kappa_k^2q^{10}},
\quad r=O_h(k\log q)=o(q).
\tag{L20}
\]

This Taylor map is an auxiliary estimate and computational compression. It is not the programme's period observation and is not identified with its kernel. Its exact relation to each determinant is the Schur factorization used next.

\subsection{6. Transfer the nonlinear central loss through the actual metrics}\label{loc:transfer-the-nonlinear-central-loss-through-the-actual-metrics}

For \(g=1\) or \(g=\psi_k\), define the full source loss

\[
D_N(g)=\log\det H_N[g\sigma]-\log\det H_N[ge^{-W_k}\sigma]\ge0.
\tag{L21}
\]

Decompose \(\mathcal P_N\) by the exact sequence

\[
0\longrightarrow Z_{N,r}\longrightarrow\mathcal P_N
\xrightarrow{\mathcal J_r}\mathbb C^r\longrightarrow0.
\]

The quotient metric is the attained one,
\(G_{N,r}^{\rm jet}[\nu]=(\mathcal J_rH_N[\nu]^{-1}\mathcal J_r^*)^{-1}\).
Its Schur complement retains both cross factors and the full high-monomial Gram inverse. The determinant factors into the restriction determinant and this quotient determinant, in a fixed monic frame.

On the restriction \(Z_{N,r}\), insertion of \(e^{-W_k}\) loses at most \(a\) of the Gamma norm, and at most \(\kappa_k a\) of the \(\psi_k\sigma\) norm, by (L12) and (L17). Put \(f(x)=-\log(1-x)\) for \(0\le x<1\). The two restriction losses lie between zero and \((N+1-r)f(\kappa_k a)\).

On the jet quotient, (L12) bounds the relative determinant at both \(t=0\) and \(t=1\) in the same interval \([r\log L_k,r\log U_k]\). Subtracting those endpoints, before taking absolute values, gives

\[
\boxed{|D_N(\psi_k)-D_N(1)|
\le r\log\kappa_k+(N+1-r)f(\kappa_k a).}
\tag{L22}
\]

This is a comparison of the nonlinear, finite-amplitude central loss. It is not a first variation at the Gamma source.

Define

\[
\Phi(g)=D_{2q-1}(g)+D_{2q}(g)-D_{q-1}(g)-D_q(g).
\]

Adding the four estimates proves

\[
|\Phi(\psi_k)-\Phi(1)|
\le4r\log\kappa_k+(6q+2)f(\kappa_k a).
\tag{L23}
\]

\subsubsection{Remove the central prefactor, without changing the outer density}\label{loc:remove-the-central-prefactor-without-changing-the-outer-density}

The two densities \(\psi_k\sigma\) and \(\omega_k\) agree off \(C_k\). On \(C_k\) their difference has absolute value at most \(2U_k\sigma\). On \(Z_{N,r}\), (L17) and (L12) therefore bound the relative change by \(2\kappa_k a\). On the jet quotient the relative spectrum lies in \([\kappa_k^{-1},\kappa_k]\), giving absolute log determinant at most \(r\log\kappa_k\). Thus their full source determinant difference, at each endpoint, is bounded by

\[
r\log\kappa_k+(N+1-r)f(2\kappa_k a).
\tag{L24}
\]

The original relation part is included in Section 7. Equations (L22)--(L24) are why a polynomial bound on the remaining factor suffices: it is charged on \(r=O(k\log q)\) quotient coordinates, not on all \(q\) coordinates independently.

\subsection{7. Retain the original relation determinant and its sign}\label{loc:retain-the-original-relation-determinant-and-its-sign}

Set

\[
R=k\sqrt{\delta^2+\gamma^2},\quad Y=2^{-16}q,\quad
B_\chi=\frac{qR^2}{Y^2-R^2},
\]

\[
\eta_0=\frac{\sqrt{2\pi}}{c_\Gamma}
\frac{(2q+1)(4q+1)\sqrt{1+2q}}q16^{-q},\qquad
\eta_\Gamma=\frac{e^{B_\chi}\eta_0}{1-\eta_0}.
\tag{L25}
\]

Use the original cofinal guards \(q\ge100,R\le Y/4,K\le Y\). For every polynomial \(P\) of degree at most \(q\), the retained full-primary estimate is

\[
\int_{C_k}|QP|^2d\sigma\le\eta_\Gamma\|QP\|_\sigma^2.
\tag{L26}
\]

It follows by pairing the original opposite roots outside \([-Y,Y]\) and by the full shifted-Legendre expansion on \([q,2q]\) inside. The inner polynomial bound is

\[
|P(y)|^2\le\frac{(2q+1)(4q+1)100^{2q}}q
\int_q^{2q}|P(t)|^2dt,
\]

while the reference power norm is at least
\(c_\Gamma q^{2q}e^{-\pi q}(1+2q)^{-1/2}\int_q^{2q}|P|^2\).
The inner multiplier is at most \([(9/4)2^{-32}q^2]^q\), and
\((9/4)e^\pi10^4 2^{-32}<1/16\). These bounds give (L26) with both inner contributions and the full mass retained. This is also ACW27--29 in the frozen standalone note.

By (L12), the central fraction in the \(\psi_k\sigma\) relation norm is at most \(\kappa_k\eta_\Gamma\). Hence insertion of \(e^{-W_k}\) changes each attained monic relation norm by a factor in
\([1-\kappa_k\eta_\Gamma,1]\). Its full weighted loss is

\[
0\le\Psi^{\rm rel}(\psi_k)
:=\sum_{j=0}^qc_j\log\frac{\nu_j[\psi_k\sigma]}
{\nu_j[\psi_ke^{-W_k}\sigma]}
\le2q f(\kappa_k\eta_\Gamma).
\tag{L27}
\]

The exact source/relation identity gives the signed equation

\[
\mathcal B_k[\psi_ke^{-W_k}\sigma]-\mathcal B_k[\psi_k\sigma]
=\Phi(\psi_k)-\Psi^{\rm rel}(\psi_k).
\tag{L28}
\]

Thus the relation contribution has not been assigned zero or the wrong sign.

For the central prefactor comparison between \(\psi_k\sigma\) and \(\omega_k\), the analogous relative relation-Gram error is at most \(2\kappa_k\eta_\Gamma\). The sum of the four relation dimensions is \(2q+2\). Its absolute signed log-determinant contribution is bounded by

\[
(2q+2)f(2\kappa_k\eta_\Gamma).
\tag{L29}
\]

Use the explicit eventual guard \(2\kappa_k\eta_\Gamma\le1/2\). Both (L27) and (L29) are exponentially small: \(\eta_\Gamma=\exp[-q\log16+O_h(\log q+1)]\), and \(\kappa_k\) is polynomial in \(q,k\).

\subsection{8. Four explicit Gamma matrices, with a finite total error}\label{loc:four-explicit-gamma-matrices-with-a-finite-total-error}

The Gamma loss \(D_N(1)\) can be evaluated with no source interpolation. Define the central moment matrix

\[
M_{W,k}[i,j]=\int_{-K}^{K}
(1-e^{-W_k(y)})y^{i+j}\sigma(y)dy,
\qquad0\le i,j<r.
\tag{L30}
\]

Let

\[
T_{r,N}[j,n]=[y^j]\varphi_n(y),\qquad
E_{r,N}=T_{r,N}T_{r,N}^*,\quad 0\le n\le N.
\tag{L31}
\]

All entries in \(T\) come from (L14) or the exact Gamma recurrence. All entries in \(M_{W,k}\) come from the fixed function (C4) on a finite interval. Define

\[
\widehat D_N=-\log\det(I_r-M_{W,k}^{1/2}E_{r,N}M_{W,k}^{1/2}),
\]

\[
\boxed{\widehat\Phi_{q,k}
=\widehat D_{2q-1}+\widehat D_{2q}-\widehat D_{q-1}-\widehat D_q.}
\tag{L32}
\]

The square root can be avoided in a determinant computation using
\(\det(I_r-M_{W,k}E_{r,N})\), by Sylvester's determinant identity. Its positive self-adjoint representative is the displayed one; a non-Hermitian coordinate matrix is not silently treated as Hermitian.

To prove the approximation, let \(C_N\) be the Gamma compression of multiplication by \((1-e^{-W_k})1_{C_k}\). Then
\(D_N(1)=-\log\det(I-C_N)\), and \(I-C_N\succeq e^{-\Lambda_k}I\).
By (L18), its Taylor approximation \(\widetilde C_N=T^*M_{W,k}T\) satisfies

\[
\|C_N-\widetilde C_N\|\le2\sqrt a+a.
\]

Put \(z_k=e^{\Lambda_k}(2\sqrt a+a)\). Equation (L20) gives \(z_k<1/2\), so the approximate matrix is positive and

\[
|D_N(1)-\widehat D_N|\le(N+1)f(z_k).
\tag{L33}
\]

For example, conjugate the difference by \((I-C_N)^{-1/2}\); its norm is at most \(z_k\), so every eigenvalue of the relative matrix lies in \([1-z_k,1+z_k]\). This proves (L33) without assuming a small inverse norm independent of the rate depth.

Both \(D_N(1)\) and \(\widehat D_N\) increase with \(N\). For the latter, \(E_{r,N}\) increases in the positive order. For the former, write \(I-C_N\) in successive Gamma coordinates; its new Schur complement is at most one, or use compression interlacing. Hence \(\Phi(1)\ge0\) and \(\widehat\Phi_{q,k}\ge0\).

Combining (L23), (L24), (L27)--(L29), and (L33) yields (R1) with the fully explicit radius

\[
\boxed{\begin{aligned}
E_k={}&8r\log\kappa_k+(6q+2)
\{f(\kappa_k a)+f(2\kappa_k a)+f(z_k)\}\\
&+(2q+2)f(2\kappa_k\eta_\Gamma)
+2qf(\kappa_k\eta_\Gamma).
\end{aligned}}
\tag{L34}
\]

Every quantity on this right-hand side is given above. It contains no original growing Gram inverse. The prescribed rank gives

\[
E_k=O_h(k(\log q)^2)=o_h(k^2)=o_h(q).
\tag{L35}
\]

The pointwise relative approximation in (C11) is not being charged again; (L34) is a direct estimate on the actual source. In particular it is stronger at this stage than multiplying the earlier local-limit error by an assumed constant kernel density.

\subsection{\texorpdfstring{9. The central degree-band response is only \(O_h(k^2)\)}{9. The central degree-band response is only O\_h(k\^{}2)}}\label{loc:the-central-degree-band-response-is-only-o_hk2}

The rank \(r=O(k\log q)\) establishes the transfer, but a separate exact recurrence argument improves the size of the response itself. This is essential for the scale conclusion in (R3).

The Gamma orthonormal recurrence is

\[
y\varphi_n=a_{n+1}\varphi_{n+1}+a_n\varphi_{n-1},\qquad
a_n=\sqrt{n(n-1/2)}.
\tag{L36}
\]

Consider either original degree band \(n=n_0,\ldots,n_0+q-1\), where \(n_0=q\) or \(q+1\). Relative to the two initial functions \(\varphi_{n_0},\varphi_{n_0-1}\), write

\[
\varphi_n(y)=A_n(y)\varphi_{n_0}(y)+B_n(y)\varphi_{n_0-1}(y).
\tag{L37}
\]

These coefficients are obtained by the actual transfer matrices

\[
\begin{pmatrix}\varphi_{n+1}\\\varphi_n\end{pmatrix}
=\begin{pmatrix}z/a_{n+1}&-a_n/a_{n+1}\\1&0\end{pmatrix}
\begin{pmatrix}\varphi_n\\\varphi_{n-1}\end{pmatrix}.
\]

At \(z=0\) each transfer has norm one because \(a_n<a_{n+1}\). Its norm at \(z\) is at most \(1+|z|/a_{n+1}\). There are fewer than \(q\) factors, and \(a_{n+1}\ge q\), so

\[
|A_n(z)|,|B_n(z)|\le e^{|z|}.
\tag{L38}
\]

This removes the common high-degree oscillation by an exact two-function representation, not by an asymptotic kernel substitution.

Truncate \(A_n,B_n\) below Taylor degree \(s\), applying Cauchy's formula on \(|z|=2K\). The restriction map from this Gamma-orthonormal degree band to \(L^2(C_k,\sigma dy)\) then has a rank-at-most \(2s\) approximation, with operator error at most

\[
\nu_s=2^{2-s}\sqrt q\,e^{2K}.
\tag{L39}
\]

Indeed each coefficient tail is at most \(2^{1-s}e^{2K}\); the \(q\) coefficient columns cost \(\sqrt q\), and the two initial functions each have full Gamma norm one. The approximation range is spanned by
\(y^j\varphi_{n_0},y^j\varphi_{n_0-1}\), \(j<s\). No orthogonality of these range functions is assumed. These are auxiliary output functions for the restriction operator, not replacement representatives at the original degree cutoff.

\subsubsection{Retain the complete lower-degree Schur inverse}\label{loc:retain-the-complete-lower-degree-schur-inverse}

Let \(V_L,V_H\) be the low-degree and band restriction maps multiplied by \(\sqrt{1-e^{-W_k}}\). The full suppressed source Gram is \(I-V^*V\). Its attained band Gram is exactly

\[
\boxed{S=I-V_H^*\mathcal A_LV_H,\qquad
\mathcal A_L=I+V_L(I-V_L^*V_L)^{-1}V_L^*.}
\tag{L40}
\]

Here the complete preceding source Gram inverse remains. It satisfies
\(\|\mathcal A_L\|\le e^{\Lambda_k}\), and the minimum definition gives
\(S\succeq e^{-\Lambda_k}I\). Thus both possible amplifications of a band approximation must be paid.

Choose

\[
\boxed{s=\left\lceil 2+
\frac{2K+2\Lambda_k+(11/2)\log q+\log8}{\log2}\right\rceil.}
\tag{L41}
\]

Then \(\nu_s\le e^{-2\Lambda_k}q^{-5}/8\). Replacing only \(V_H\) by its rank-\(2s\) approximation changes the matrix in (L40) by at most
\(e^{\Lambda_k}\nu_s(2+\nu_s)\). Relative to \(S\), its norm is at most

\[
u_k=e^{2\Lambda_k}\nu_s(2+\nu_s)\le3q^{-5}/8<1/2.
\]

The approximate band Gram is positive, at least
\(e^{-\Lambda_k}(1-u_k)I\), and differs from the identity in at most \(2s\) directions. Therefore

\[
0\le-\log\det S\le2s\{\Lambda_k+f(u_k)\}+qf(u_k).
\tag{L42}
\]

The lower bound follows from \(S\preceq I\). For the upper bound, bound the at most \(2s\) approximate eigenvalue losses by \(\Lambda_k+f(u_k)\), then use the full \(q\)-dimensional relative log-determinant bound. This argument explicitly preserves the lower-space coupling in (L40).

Each of the two differences making up \(\Phi(1)\) is the loss (L42) on its original band. Hence

\[
\boxed{0\le\Phi(1)\le4s(\Lambda_k+1)+2qf(u_k)=O_h(k^2).}
\tag{L43}
\]

Equations (L33)--(L35) prove (R3). More quantitatively,

\[
\boxed{0\le\liminf\frac{\delta_k^{\rm cen,end}}{k^2}
\le\limsup\frac{\delta_k^{\rm cen,end}}{k^2}
\le\frac{8D_0(b+D_0)}{\log2}.}
\tag{L44}
\]

This is a signed asymptotic interval, not a computed limiting coefficient. It follows directly from \(s/k\to(2b+2D_0)/\log2\), \(\Lambda_k/k=D_0\), and \(E_k/k^2\to0\).

\subsubsection{The scale conclusions on the original multiplicity strata}\label{loc:the-scale-conclusions-on-the-original-multiplicity-strata}

For every fixed \(m_0\), \(q\ge(k+1)^2\), so

\[
\boxed{\frac{\delta_k^{\rm cen,end}}{q\log k}\longrightarrow0,
\qquad \delta_k^{\rm cen,end}=O_h(q).}
\tag{L45}
\]

On the simple stratum, \(q=(k+1)^2\), (L44) is an explicit nonnegative limiting interval at the \(q\) scale. On every fixed stratum \(m_0\ge2\),
\(q\asymp_{m_0}k^3\), and

\[
\boxed{\delta_k^{\rm cen,end}=o_h(q),\qquad
\delta_k=\delta_k^{\rm out}+o_h(q).}
\tag{L46}
\]

Unlike the incoming unweighted area comparison, (L45)--(L46) are proved for the actual endpoint quotient-determinant contribution with its full relation determinant retained.

\subsection{10. Endpoint accounting, the earlier spatial split, and the receiver}\label{loc:endpoint-accounting-the-earlier-spatial-split-and-the-receiver}

The outer-only density is explicitly \(\omega_k\), and

\[
\delta_k^{\rm out}+\delta_k^{\rm cen,end}
=\mathcal B_k[m_k]-\mathcal B_k[\sigma].
\tag{L47}
\]

The mass cancellation used here is exact: \(G_N[A\nu]=A G_N[\nu]\), so four signs cancel \(q\log A\). It is equivalently the scalar coefficient map \(P\mapsto\sqrt A P\), with the same map on remainder values, between source measures \(\nu\) and \(A^{-1}\nu\).

An isolated spatial central integral along the earlier diagonal interpolation is not identified with this central endpoint difference. Their exact relation can be given. Let

\[
f_C=1_{C_k}\ell_k,\quad f_O=1_{\mathbb R\setminus C_k}\ell_k,
\quad F(a,b_0)=\mathcal B_k[\sigma e^{af_C+b_0f_O}].
\]

The earlier central spatial term is \(\int_0^1F_a(t,t)dt\), whereas the endpoint central term is \(\int_0^1F_a(t,1)dt\). Consequently

\[
\boxed{\delta_k^{\rm cen,end}-\delta_k^{\rm cen,diag}
=\int_0^1\int_t^1F_{ab_0}(t,b_0)\,db_0dt.}
\tag{L48}
\]

The mixed derivative is determined by the original source and relation kernels. For a single full source determinant, differentiation twice gives

\[
\partial_a\partial_{b_0}\log\det H_N
=\operatorname{Tr}(H_N^{-1}H_{N,ab_0})
-\operatorname{Tr}(H_N^{-1}H_{N,b_0}H_N^{-1}H_{N,a}).
\]

The first term is zero since \(f_Cf_O=0\). The second is the negative of
\[\iint f_C(y)f_O(z)|K_N(y,z)|^2d\nu(y)d\nu(z).\]
The relation expression uses its full relation kernel; (C2) and the four signs give \(F_{ab_0}\). This explicitly records the path-interchange contribution, rather than presuming that the earlier two spatial pieces are individually path independent.

The complete action receives the new result as

\[
\boxed{J_k^{\rm action}=\mathcal B_k[\sigma]
+\delta_k^{\rm out}+\widehat\Phi_{q,k}
+\mathfrak e_k+W_k^{\rm mono}-P_{k,-}-P_{k,+}.}
\tag{L49}
\]

Here \(W_k^{\rm mono}\) denotes the programme's original monic-window term; it is distinguished from the explicitly defined rate function \(W_k(y)\). Every complementary term and both penalty terms remain. No favourable sign has been assigned to \(\delta_k^{\rm out}\), and the known positive canonical baseline is not removed.

\subsection{11. What remains to be evaluated}\label{loc:what-remains-to-be-evaluated}

There are now two precisely separated value calculations.

The outer correction is exactly \(\mathcal B_k[\omega_k]-\mathcal B_k[\sigma]\), with the actual convolution outside \([-kb,kb]\). The original unequal upper and lower tail powers have not been converted into a sharp arithmetic tail coefficient. Its \(q\log k\) and \(q\) terms remain to be evaluated.

On the simple stratum, the remaining central order-\(q\) value is that of the explicit matrices (L30)--(L32). This note proves an \(o(q)\) transfer and a signed \(O(q)\) interval, not a limit formula for those matrices. On higher fixed multiplicity strata that matrix contribution is \(o(q)\), by (L43), and is no longer needed for the order-\(q\) coefficient.

In particular, the density-area value \(-I_h\) is not inserted as the baseline coefficient. No unproved constant-density limit for the source band has been used. No current geometric purity theorem is applied without a constructed metric comparison. The new estimates are source and determinant calculations and do not constitute programme closure.

\subsection{Source record and verification scope}\label{loc:source-record-and-verification-scope}

Imported arithmetic inputs are the original RTM envelopes, the full even source and packet, SXR's source/relation identity, and the original Gamma recurrence and mass. Their upstream analytic proofs have not received a new exhaustive audit. The incoming central local-limit proof is retained in full, with its constants reproduced and the one additional shifted-point consequence proved above.

The generating-function convention and the complex Stirling bound were checked against the NIST DLMF, Sections 18.23 and 5.11. The degree-band transfer, whole-line prefactor comparison, Taylor-jet localization, nonlinear endpoint transfer, finite determinant error, and the \(O(k^2)\) band estimate are derived in this note.

The executable checks are auxiliary exact matrix and polynomial diagnostics. They do not evaluate an actual off-critical xi packet, prove every analytic inequality by computation, or supply an independent review or formal verification of this derivation. Original inputs are frozen with hashes; their earlier check counts are not counted as new analytic results.

\subsubsection{Executed auxiliary diagnostics}\label{loc:executed-auxiliary-diagnostics}

The new checker passed 88 exact checks and rejected five deliberately false formulas. Its ordinary and \texttt{python\ -O} outputs are byte-identical. The original checker was replayed unchanged: all 39 checks and four false-formula controls agree with the original JSON receipt and with its optimized run. These statements report the actual executions recorded in \texttt{checks/EXECUTION\_RECEIPT.json}; they do not upgrade the written analytic derivation to an independent or formal verification.

% END PRESERVED INPUT provider_CONSOLIDATED_LOCALIZATION.tex
\endgroup

\clearpage
\begingroup


\setcounter{equation}{0}
\renewcommand{\theequation}{ACW.\arabic{equation}}
\section*{Original central source and tilted local-limit proof}
\addcontentsline{toc}{section}{Original central source and tilted local-limit proof}
% BEGIN PRESERVED INPUT provider_CENTRAL_SOURCE_WINDOW.tex SHA256 e16d985242426fc2055720106e047a3447e0d7b13019be6b7f7e694d95330c80
\section{Signed arithmetic baseline: central-source continuation}\label{acw:signed-arithmetic-baseline-central-source-continuation}

18 September 2026. Written research derivation; not independently reviewed or formally verified.

\subsection{Completion status}\label{acw:completion-status}

The assigned target is the finer value of

\[
\delta_{k,1}^{\mathrm{ar}}=\mathcal B_k^{\mathrm{ar}}-\mathcal B_k^\sigma.
\]

\textbf{Its requested \(q\log k\)- and \(q\)-scale coefficients are not computed in this note.} The last integration in \emph{Untitled 1775}, lines 13324--13328, also leaves that coefficient uncomputed. This note proves a new source-side central-window expansion and writes its exact, still unevaluated, return into the assigned determinant. Its central coefficient must not be recorded as the coefficient of \(\delta_{k,1}^{\mathrm{ar}}\).

The precision target in \emph{Untitled 1775}, lines 9389--9397, is a value for the signed correction, not just a convergent enclosure around an unevaluated centre. That target is preserved here. No numerical hypothetical xi packet is supplied. All constants below retain their dependence on the fixed original packet.

\subsection{1. Original objects and the imported source inequalities}\label{acw:original-objects-and-the-imported-source-inequalities}

Fix the programme's stipulated full quartet

\[
\rho=\tfrac12+\delta+i\gamma,\qquad
0<\delta<\tfrac12,\quad\gamma>2,
\]

\[
h(s)=\prod_{z\in\{\rho,\bar\rho,1-\rho,1-\bar\rho\}}(s-z)^{m_0},
\qquad m_0\ge1,
\]

and retain

\[
v_h=2\xi/h,\quad
w_h(y)=\frac{|v_h(1/2+iy)|^2}{2\pi},\quad m_k=w_h^{*k},
\]

\[
k\equiv1\pmod4,\quad k\ge5,\quad e=1+k(m_0-1),\quad
q=e(k+1)^2,\quad c=k/2,
\]

\[
\chi(S)=\prod_{a,b=0}^k
\bigl(S-c-(2a-k)\delta-i(2b-k)\gamma\bigr)^e.
\]

The reference is exactly

\[
\sigma(y)=\frac{|\Gamma(1/4+iy/2)|^2}{2\pi},\quad
\int_{\mathbb R}\sigma(y)\,dy=\sqrt{2\pi},\quad \alpha=\pi/2.
\]

The RTM1--3 inputs used below, with their original constants, are

\[
e^{\alpha y}w_h(y)\le C(1+|y|)^{-p},\quad
p=8m_0-9/2\ge7/2,\quad C=C_h^{\mathrm{tilt}},
\tag{ACW1}
\]

\[
w_h^{*3}(y)\ge c_h e^{-\alpha|y|}(1+|y|)^{-B},
\quad B=42+8m_0,
\]

\[
c_\Gamma\frac{e^{-\alpha|y|}}{\sqrt{1+|y|}}
\le\sigma(y)\le
C_\Gamma\frac{e^{-\alpha|y|}}{\sqrt{1+|y|}},
\quad
c_\Gamma=\sqrt2e^{-7/6},\quad
C_\Gamma=\sqrt{2\sqrt5}e^{2/3}.
\]

The original density is even and positive almost everywhere. The constants and those properties are retained inputs, not re-audited upstream claims. The proof uses the three-fold lower bound only on the fixed interval \([-1,1]\). This continuation does not require its large-\(y\) lower power to equal the large-\(y\) upper power.

Exact source pin: \texttt{06\_Foundations\_and\_Original\_Six\_Readers(1).tex}, lines 112323--112602, RTM1--11. The full original source has not been replaced by a model density.

\subsection{2. Tilted mass, mean, and variance on the closed tilt interval}\label{acw:tilted-mass-mean-and-variance-on-the-closed-tilt-interval}

Define the actual one-factor integrals

\[
M(\theta)=\int_{\mathbb R}e^{\theta y}w_h(y)\,dy,
\quad L(\theta)=\log M(\theta),\quad |\theta|\le\alpha,
\]

\[
\mu(\theta)=L'(\theta),\quad V(\theta)=L''(\theta),\quad
\mathfrak M=M(\alpha),\quad b_h=\mu(\alpha).
\tag{ACW2}
\]

Evenness and ACW1 imply

\[
e^{\theta y}w_h(y)\le C(1+|y|)^{-p}
\quad(|\theta|\le\alpha).
\]

Take \(\eta=1/4\). Since \(p>3+\eta\), dominated differentiation through order two is valid on the closed interval, with finite uniformly bounded absolute moments of order \(2+\eta\). Thus \(\mu\) is continuously differentiable, \(\mu'=V\), and

\[
\mu(\theta)=\frac{\int y e^{\theta y}w_h(y)\,dy}{M(\theta)},
\]

\[
V(\theta)=\frac{\int (y-\mu(\theta))^2e^{\theta y}w_h(y)\,dy}{M(\theta)}>0.
\]

Strict positivity follows because the density is not concentrated at a point. Furthermore \(M\) is even, \(\mu\) odd, and \(V\) even. Consequently \(b_h>0\), and \(\mu\) maps \([-\alpha,\alpha]\) bijectively onto \([-b_h,b_h]\). Write its inverse as \(\theta(t)\).

No source probability normalization is made in ACW2. For the Fourier proof only, introduce

\[
p_\theta(y)=\frac{e^{\theta y}w_h(y)}{M(\theta)}.
\]

The exact map back to the original convolution is

\[
\boxed{m_k(x)=M(\theta)^k e^{-\theta x}p_\theta^{*k}(x).}
\tag{ACW3}
\]

Indeed, multiplication of the \(k\) densities gives the factor
\(M(\theta)^{-k}e^{\theta\sum y_j}\); the original addition map has \(\sum y_j=x\) and Jacobian one. Every original mass is therefore present in ACW3.

\subsection{3. Explicit uniform Fourier bounds, including the boundary tilts}\label{acw:explicit-uniform-fourier-bounds-including-the-boundary-tilts}

This section proves the local density estimate rather than invoking an unspecified uniform local limit theorem.

Put \(M_0=M(0)\), and write \(\mathrm B\) for Euler's beta function. Define

\[
V_* = \frac{2C}{M_0}\,\mathrm B(3,p-3),
\]

\[
H_* = 2^{1+\eta}\left[
\frac{2C}{M_0}\,\mathrm B(3+\eta,p-3-\eta)
+b_h^{2+\eta}\right].
\tag{ACW4}
\]

These follow by integrating the original envelope:

\[
V(\theta)\le\int y^2p_\theta(y)\,dy\le V_*,
\qquad
\int|y-\mu(\theta)|^{2+\eta}p_\theta(y)\,dy\le H_*.
\]

Here \(M(\theta)\ge M_0\), and
\(|x-y|^{2+\eta}\le2^{1+\eta}(|x|^{2+\eta}+|y|^{2+\eta})\).

On \([-1,1]\), the exact tilted-convolution identity gives

\[
p_\theta^{*3}(x)\ge
\frac{c_h e^{-2\alpha}2^{-B}}{\mathfrak M^3}.
\]

Set

\[
a_* = \min\left\{\frac14,
\frac{c_h e^{-2\alpha}2^{-B}}{\mathfrak M^3}\right\},
\quad \beta=2a_*,\quad v_* = \beta/9.
\tag{ACW5}
\]

The probability density \(p_\theta^{*3}\) is the sum of \(\beta\) times the uniform probability density on \([-1,1]\) and a nonnegative residual of mass \(1-\beta\). Variance of this mixture is at least \(\beta/3\). Its variance is also exactly \(3V(\theta)\), by expansion of the centred square of the sum. Hence

\[
0<v_*\le V(\theta)\le V_*.
\]

Let

\[
\phi_\theta(\xi)=\int e^{i\xi(y-\mu(\theta))}p_\theta(y)\,dy.
\]

Taylor's formula and ACW4 give

\[
\left|\phi_\theta(\xi)-1+\tfrac12V(\theta)\xi^2\right|
\le4H_*|\xi|^{2+\eta}.
\tag{ACW6}
\]

For completeness, the scalar remainder is at most \(|u|^3/6\) for \(|u|\le1\), and at most \(2+|u|+u^2/2\le(7/2)|u|^{2+\eta}\) for \(|u|>1\); the stated factor four bounds both regions.

Define

\[
\xi_0=\min\left\{1,(2V_*)^{-1/2},
\left(\frac{v_*}{16H_*}\right)^{1/\eta}\right\}.
\]

For \(|\xi|\le\xi_0\), the nonnegative quadratic principal term and ACW6 imply

\[
|\phi_\theta(\xi)|\le e^{-v_*\xi^2/4}.
\tag{ACW7}
\]

The fixed-interval mixture also gives

\[
|\phi_\theta(\xi)|^3\le1-\beta+\beta\left|\frac{\sin\xi}{\xi}\right|.
\]

For \(|\xi|\ge\xi_0\),
\(|\sin\xi/\xi|\le1-\xi_0^2/7\). To verify this bound, use monotonicity of \(\sin x/x\) on \([0,2]\), its Taylor bound at \(0<\xi_0\le1\), and \(|\sin x/x|\le1/2\) beyond two. Consequently

\[
|\phi_\theta(\xi)|\le r_*
:=\left(1-\frac{\beta\xi_0^2}{7}\right)^{1/3}<1
\quad(|\xi|\ge\xi_0).
\tag{ACW8}
\]

The envelope and Plancherel give the uniform bound

\[
\int_{\mathbb R}|\phi_\theta(\xi)|^2\,d\xi
\le C_2:=\frac{4\pi(C/M_0)^2}{2p-1}.
\tag{ACW9}
\]

In particular Fourier inversion of the \(k\)-fold density is valid for every \(k\ge2\).

\subsubsection{A fully specified error}\label{acw:a-fully-specified-error}

Put

\[
C_0=4H_*+V_*^2/8,\quad a_k=v_*(k-1)/4,
\]

\[
\begin{aligned}
A_k^{\mathrm{loc}}=\frac1{2\pi}\bigg[
&C_0 k\,\Gamma\!\left(\frac{3+\eta}{2}\right)
 a_k^{-(3+\eta)/2}
+C_2 r_*^{k-2}\\
&+\frac{2e^{-kv_*\xi_0^2/2}}{kv_*\xi_0}
\bigg],\qquad
\epsilon_k=\sqrt{2\pi kV_*}\,A_k^{\mathrm{loc}}.
\end{aligned}
\tag{ACW10}
\]

These local symbols are not the programme's arithmetic upper-envelope constant \(A_k\).

For \(|\xi|\le\xi_0\),

\[
|\phi_\theta(\xi)-e^{-V(\theta)\xi^2/2}|
\le C_0|\xi|^{2+\eta}.
\]

Telescoping the difference of the two \(k\)-th powers, and using ACW7, bounds its integral by the first term of ACW10. ACW8--9 bound the remaining characteristic-function integral by \(C_2r_*^{k-2}\). Integration of the remaining Gaussian tail gives the last term. Thus

\[
\left|
p_\theta^{*k}(k\mu(\theta))
-\frac1{\sqrt{2\pi kV(\theta)}}
\right|\le A_k^{\mathrm{loc}}.
\]

Equivalently,

\[
\boxed{
p_\theta^{*k}(k\mu(\theta))
=\frac{1+\varepsilon_{k,\theta}}
{\sqrt{2\pi kV(\theta)}},
\qquad
\sup_{|\theta|\le\alpha}|\varepsilon_{k,\theta}|\le\epsilon_k
=O_h(k^{-1/8}).}
\tag{ACW11}
\]

The exponent is \(-\eta/2=-1/8\). All constants are independent of \(\theta\), so the two endpoint tilts are included. For logarithms, use the eventual explicitly testable range \(\epsilon_k\le1/2\). Such a range exists by ACW10.

\subsection{4. The full central density profile}\label{acw:the-full-central-density-profile}

Define the source contrast, with the full mass explicitly recorded,

\[
\ell_k(y)=\log\frac{m_k(y)}{\mathfrak M^k\sigma(y)}.
\]

Define on \([-b_h,b_h]\)

\[
F_h^{\mathrm{cen}}(t)=
L(\theta(t))-t\theta(t)-L(\alpha)+\alpha|t|.
\tag{ACW12}
\]

The function is even. On \(0<t<b_h\), direct differentiation using \(L'(\theta(t))=t\) gives

\[
(F_h^{\mathrm{cen}})'(t)=\alpha-\theta(t),
\qquad
(F_h^{\mathrm{cen}})''(t)=-\frac1{V(\theta(t))},
\qquad F_h^{\mathrm{cen}}(b_h)=0.
\]

In particular \(F_h^{\mathrm{cen}}(t)<0\) for \(|t|<b_h\). Its cusp at zero comes from the explicitly retained term \(\alpha|t|\); the original convolution itself remains smooth.

Write

\[
c(y)=\sigma(y)e^{\alpha|y|}\sqrt{1+|y|},
\qquad c_\Gamma\le c(y)\le C_\Gamma.
\]

Substitute \(\theta=\theta(t)\) in the exact identity ACW3 and in ACW11. There is then the exact expression

\[
\boxed{
\begin{aligned}
\ell_k(kt)={}&kF_h^{\mathrm{cen}}(t)
+\tfrac12\log(|t|+k^{-1})
-\tfrac12\log(2\pi V(\theta(t)))\\
&-\log c(kt)+\log(1+\varepsilon_{k,\theta(t)}).
\end{aligned}}
\tag{ACW13}
\]

The last scalar is the actual relative local-density error, not a freely selected endpoint value. It satisfies
\(|\log(1+\varepsilon_{k,\theta(t)})|\le2\epsilon_k\) on the stated logarithmic range. This proves the uniform central profile including \(t=0\) and \(t=\pm b_h\), without a saddle point beyond the original boundary of exponential integrability.

\subsection{5. The integrated central source contrast}\label{acw:the-integrated-central-source-contrast}

The scalar now evaluated is explicitly

\[
C_{h,k}^{\mathrm{cen}}:=
\int_{-kb_h}^{kb_h}\ell_k(y)\,dy.
\]

It is a Lebesgue-integrated log-density statistic. It is not, by definition or by any established metric identification, the programme's signed quotient determinant.

Define the fixed original-source functional

\[
\mathfrak I_h^{\mathrm{cen}}
=\int_0^\alpha\mu(\theta)^2\,d\theta>0.
\tag{ACW14}
\]

It has finite explicit bounds

\[
\frac{v_*^2\alpha^3}{3}\le\mathfrak I_h^{\mathrm{cen}}
\le\frac{V_*^2\alpha^3}{3},
\]

because \(\mu(0)=0\) and \(v_*\le\mu'\le V_*\). In particular, this coefficient has a determined nonzero sign.

The area of the rate function is evaluated exactly:

\[
\boxed{
\int_{-b_h}^{b_h}F_h^{\mathrm{cen}}(t)\,dt
=-\mathfrak I_h^{\mathrm{cen}}.}
\tag{ACW15}
\]

Indeed,

\[
F_h^{\mathrm{cen}}(t)
=-\int_t^{b_h}(\alpha-\theta(u))\,du
\quad(0\le t\le b_h).
\]

Integrating this formula gives

\[
\int_{-b_h}^{b_h}F_h^{\mathrm{cen}}
=-2\int_0^{b_h}u(\alpha-\theta(u))\,du.
\]

The substitution \(u=\mu(\theta)\) followed by integration by parts gives ACW15. The boundary term is zero at \(\theta=0\), because \(\mu(0)=0\), and at \(\theta=\alpha\), because of \(\alpha-\theta\).

\subsubsection{Finite remainder}\label{acw:finite-remainder}

Set

\[
C_h^{(0)}=
\tfrac12\max\{|\log(2\pi v_*)|,|\log(2\pi V_*)|\}
+\max\{|\log c_\Gamma|,|\log C_\Gamma|\}.
\]

Integrating ACW13 and using ACW15 gives the entirely explicit bound

\[
\boxed{
\left|C_{h,k}^{\mathrm{cen}}+k^2\mathfrak I_h^{\mathrm{cen}}\right|
\le k\left[
1+b_h\log(1+b_h)+2b_h C_h^{(0)}+4b_h\epsilon_k
\right].}
\tag{ACW16}
\]

The elementary bound used here is

\[
\int_0^{b_h}|\log(t+k^{-1})|\,dt
\le1+b_h\log(1+b_h),\quad k\ge1.
\]

The negative part is bounded by \(\int_0^1-\log t\,dt=1\); the positive part is at most \(b_h\log(1+b_h)\). No interval around zero has been removed.

Consequently

\[
\boxed{C_{h,k}^{\mathrm{cen}}=-k^2\mathfrak I_h^{\mathrm{cen}}+O_h(k).}
\tag{ACW17}
\]

\subsubsection{Its next coefficient}\label{acw:its-next-coefficient}

The precise Gamma asymptotic also gives

\[
\sigma(y)e^{\alpha|y|}|y|^{1/2}\longrightarrow\sqrt2.
\]

This is the substitution \(z=1/4+iy/2\) in DLMF 5.11.9 (or the modulus of Stirling's formula 5.11.1). Therefore \(c(y)\to\sqrt2\). Define

\[
\mathfrak J_h^{\mathrm{cen}}
=\int_0^{b_h}
\log\frac{t}{4\pi V(\theta(t))}\,dt
=\int_0^\alpha
V(\theta)\log\frac{\mu(\theta)}{4\pi V(\theta)}\,d\theta.
\tag{ACW18}
\]

The logarithmic endpoint is integrable, since \(\mu(\theta)\) is bounded above and below by positive constants times \(\theta\). In ACW13, the variance term is bounded, the Gamma term is uniformly bounded and converges away from zero, and the local-density error tends uniformly to zero. Also \(\log(t+1/k)\to\log t\) in \(L^1(0,b_h)\). Dominated convergence consequently proves

\[
\boxed{
C_{h,k}^{\mathrm{cen}}
=-k^2\mathfrak I_h^{\mathrm{cen}}
+k\mathfrak J_h^{\mathrm{cen}}+o_h(k).}
\tag{ACW19}
\]

The fixed integrals in ACW14 and ACW18 contain only the original one-factor tilted mass and its first two derivatives. They do not contain an unknown growing quotient Gram or an assigned projected allocation.

On the simple-quartet family \(q=(k+1)^2\),

\[
C_{h,k}^{\mathrm{cen}}/q\longrightarrow-\mathfrak I_h^{\mathrm{cen}}.
\]

For every fixed \(m_0\ge2\), the same scalar remains of order \(k^2\), while \(q\asymp k^3\); its ratio to \(q\) tends to zero. These statements concern ACW19, not the determinant target.

\subsection{6. Exact return to the original signed determinant}\label{acw:exact-return-to-the-original-signed-determinant}

The source/relation identity in SXR10 gives, for any one of the original positive source metrics,

\[
\mathcal B_k^\mu
=\sum_{r=0}^q c_r
\log\frac{\nu_r^\mu}{\omega_{q+r}^\mu},
\qquad c_0=c_q=1,\quad c_r=2\ (0<r<q).
\tag{ACW20}
\]

Here \(\nu_r^\mu\) is the full monic minimum in \(|\chi|^2d\mu\), and \(\omega_{q+r}^\mu\) is the full monic source minimum. Every preceding relation is retained.

Multiplying a source by a scalar \(D>0\) multiplies its quotient Gram by \(D\) and its quotient determinant by \(D^q\). Thus the four signs give the exact cancellation

\[
\mathcal B_k^{D\mu}=\mathcal B_k^\mu,
\qquad (+1)+(+1)+(-1)+(-1)=0.
\]

Equivalently, both monic norms in each ratio in ACW20 are multiplied by \(D\). Hence

\[
\mathcal B_k^{m_k/\mathfrak M^k}=\mathcal B_k^{m_k}.
\tag{ACW21}
\]

This is the precise account of the scalar in \(\ell_k\); it does not alter the recorded original mass or the arithmetic class.

For an exact evaluation path, introduce on the same polynomial coefficient spaces

\[
d\nu_u(y)=\sigma(y)e^{u\ell_k(y)}\,dy,
\qquad0\le u\le1.
\tag{ACW22}
\]

Its endpoints are \(\sigma\) and \(m_k/\mathfrak M^k\). This is a new path used to compute the same endpoint difference; it is not identified with the original affine source path at intermediate parameter values.

The original TW/GTM global source envelopes imply, at each fixed \(k\),

\[
|\ell_k(y)|\le C_{h,k}+M\log(1+y^2),
\qquad M=21+4m_0.
\]

All polynomial moments and their parameter derivatives are consequently finite, locally uniformly in \(u\). The moment matrices are positive definite. Their monic minimizers depend smoothly on \(u\). Differentiating their attained norms, the coefficient-derivative cross terms vanish in the full preceding polynomial or relation space, exactly as in SXR5.

Let \(p_{q+r,u}\) and \(u_{r,u}\) be these actual monic minimizers along ACW22. Define

\[
\begin{aligned}
\mathcal K_{k,u}(y)=\sum_{r=0}^q c_r\bigg[
&\frac{|\chi(c+iy)u_{r,u}(c+iy)|^2}{\nu_{r,u}}\\
&-\frac{|p_{q+r,u}(c+iy)|^2}{\omega_{q+r,u}}
\bigg].
\end{aligned}
\tag{ACW23}
\]

Each of its two nonnegative sums has integral \(2q\) in \(d\nu_u\). In particular \(\int\mathcal K_{k,u}\,d\nu_u=0\). Differentiating ACW20 and integrating from zero to one gives the exact receiver

\[
\boxed{
\delta_{k,1}^{\mathrm{ar}}
=\int_0^1\int_{\mathbb R}
\ell_k(y)\mathcal K_{k,u}(y)\,d\nu_u(y)\,du.}
\tag{ACW24}
\]

Both endpoints, the full relation ideal, and the monic coordinate frames agree with the assigned target. The scalar \(\mathfrak M^k\) cancels by ACW21, before estimates. No eigenclass, observation kernel, proper-source metric, or quotient denominator has been replaced.

\subsection{7. The actual relation kernel has exponentially small central mass}\label{acw:the-actual-relation-kernel-has-exponentially-small-central-mass}

The central density calculation can also be returned to one of the two actual kernels in ACW24. This step does not replace the remaining source kernel by a reference kernel.

Retain GTM1's original constants

\[
M=21+4m_0,\quad
\vartheta_h=\int_{-1}^1 w_h(y)\,dy,
\]

\[
a_k^{\mathrm{TW}}=
\frac{c_h\vartheta_h^{k-3}e^{-\alpha(k-3)}(k-2)^{-B}}
 {C_\Gamma2^{B/2}},\qquad
A_k^{\mathrm{TW}}=\frac{C}{c_\Gamma}k^{p+1}\mathfrak M^{k-1}.
\]

Thus

\[
a_k^{\mathrm{TW}}(1+y^2)^{-M}\sigma(y)
\le m_k(y)\le A_k^{\mathrm{TW}}\sigma(y).
\]

Define

\[
D_n=[1+4(n+M)^2]^M,
\quad
L_k^- = D_{2q}^{-1}\min\{1,a_k^{\mathrm{TW}}/\mathfrak M^k\},
\quad
L_k^+ = \max\{1,A_k^{\mathrm{TW}}/\mathfrak M^k\}.
\tag{ACW25}
\]

On the entire geometric path ACW22,

\[
d\nu_u\le L_k^+\,d\sigma,
\qquad
\|P\|_{\nu_u}^2\ge L_k^-\|P\|_\sigma^2
\quad(\deg P\le2q).
\tag{ACW26}
\]

Here is a direct verification of the lower norm inequality. The Gamma recurrence in GTM7 bounds multiplication by \(y\) on polynomials of degree at most \(n\), including its output of degree \(n+1\), by \(2(n+1)\). Hence

\[
\|(y+i)^M P\|_\sigma^2\le D_n\|P\|_\sigma^2.
\]

Cauchy--Schwarz, applied to the factors \(|P|(1+y^2)^{-M/2}\) and \(|P|(1+y^2)^{M/2}\), proves

\[
\int|P|^2(1+y^2)^{-M}\,d\sigma\ge D_n^{-1}\|P\|_\sigma^2.
\]

For \(0\le u\le1\), raising the positive lower density ratio to the power \(u\) bounds it below by
\(\min\{1,a_k^{\mathrm{TW}}/\mathfrak M^k\}(1+y^2)^{-M}\).
This proves ACW26. In particular
\(\log(L_k^+/L_k^-)=O_h(k+\log q)\).

Put

\[
\widehat\chi(y)=i^{-q}\chi(c+iy),\quad
R=k\sqrt{\delta^2+\gamma^2},\quad
Y=2^{-16}q,\quad
B_\chi=\frac{qR^2}{Y^2-R^2},
\]

\[
\eta_0=
\frac{\sqrt{2\pi}}{c_\Gamma}
\frac{(2q+1)(4q+1)\sqrt{1+2q}}q\,16^{-q}.
\tag{ACW27}
\]

Use the eventual finite domain

\[
q\ge100,\quad R\le Y/4,\quad kb_h\le Y.
\]

It contains all sufficiently large original tensor orders at every fixed packet and fixed complete multiplicity. In particular \(\eta_0<1/2\) there: the bound \(\eta_0\le540q^2 16^{-q}\) holds, its value at \(q=100\) is below one half, and its consecutive ratio decreases thereafter.

For every polynomial \(P\) of degree at most \(q\), the original Gamma central estimate is

\[
\int_{|y|\le Y}|\widehat\chi(y)P(y)|^2\,d\sigma
\le
\frac{e^{B_\chi}\eta_0}{1-\eta_0}
\int_{\mathbb R}|\widehat\chi(y)P(y)|^2\,d\sigma.
\tag{ACW28}
\]

For completeness, pair every opposite root of \(\widehat\chi\), retaining repetitions. Outside \([-Y,Y]\), its squared modulus is at least
\(e^{-B_\chi}|y|^{2q}\). Inside that interval it is at most
\([(9/4)2^{-32}q^2]^q\).
Expand \(P\) in the complete shifted Legendre basis on \([q,2q]\). Its argument on \([-Y,Y]\) has modulus below four. The elementary recurrence bounds the degree-\(j\) Legendre polynomial there by \(10^j\). Cauchy--Schwarz therefore gives the enlarged valid bound

\[
|P(y)|^2\le
\frac{(2q+1)(4q+1)100^{2q}}q
\int_q^{2q}|P(t)|^2\,dt.
\]

The reference integral
\(\int|y|^{2q}|P|^2\,d\sigma\) is at least

\[
\frac{c_\Gamma q^{2q}e^{-\pi q}}{\sqrt{1+2q}}
\int_q^{2q}|P(t)|^2\,dt.
\]

Together with the full reference mass \(\sqrt{2\pi}\) and
\((9/4)e^\pi10^4\,2^{-32}<1/16\), this bounds both the central root-weighted and central reference-power integrals by \(\eta_0\) times the complete reference-power integral. The outside lower bound then proves ACW28. This is the same coefficient-space central comparison used in the supplied EC calculation; all its intervals and factors have been retained.

Since \(\deg(\widehat\chi P)\le2q\), ACW26 returns this estimate to the actual entire path:

\[
\boxed{
\int_{|y|\le kb_h}|\widehat\chi P|^2\,d\nu_u
\le\eta_k^{\mathrm{rel}}\int_{\mathbb R}|\widehat\chi P|^2\,d\nu_u,
\quad
\eta_k^{\mathrm{rel}}=
\frac{L_k^+}{L_k^-}\frac{e^{B_\chi}\eta_0}{1-\eta_0}.}
\tag{ACW29}
\]

The exact polynomial coordinate change \(P(S)\mapsto P(c+iy)\) changes no degree or absolute norm; thus ACW29 applies to every actual monic relation minimizer in ACW23.

Write \(\mathcal R_{k,u}\) for ACW23's nonnegative relation sum and \(\mathcal P_{k,u}\) for its nonnegative source sum. Their weights sum to \(2q\), so

\[
\int_{|y|\le kb_h}\mathcal R_{k,u}\,d\nu_u
\le2q\eta_k^{\mathrm{rel}}.
\]

ACW13 bounds the absolute contrast on this whole interval by

\[
T_k=k[L(\alpha)-L(0)]
+\tfrac12\log k+\tfrac12\log(1+b_h)+C_h^{(0)}+2\epsilon_k.
\]

Consequently the central part of the original signed correction has the exact directed enclosure

\[
\boxed{
\left|\delta_{k,1}^{\mathrm{ar,cen}}
+\int_0^1\int_{|y|\le kb_h}
\ell_k(y)\mathcal P_{k,u}(y)\,d\nu_u(y)\,du\right|
\le2qT_k\eta_k^{\mathrm{rel}}.}
\tag{ACW30}
\]

The right side is

\[
\exp[-q\log16+O_h(k+\log q)].
\]

Thus the actual relation part of this central return is exponentially controlled. The remaining central term is the original degree band of source minima, not an arbitrary projection or an assigned scalar density. Its weighting has not yet been evaluated at order \(q\).

\subsection{8. Exact stopping point}\label{acw:exact-stopping-point}

ACW13--19 evaluate the source contrast on the full interval \([-kb_h,kb_h]\). They do not evaluate its weighting by \(\mathcal K_{k,u}\,d\nu_u\) in ACW24. The explicit Fourier error in ACW11 cannot be multiplied by an assumed constant kernel density.

In particular, the two terms still required in the assigned calculation are the values, at the requested scales, of

\[
\int_0^1\int_{|y|\le kb_h}
\ell_k(y)\mathcal K_{k,u}(y)\,d\nu_u(y)\,du
\]

and

\[
\int_0^1\int_{|y|>kb_h}
\ell_k(y)\mathcal K_{k,u}(y)\,d\nu_u(y)\,du.
\]

ACW30 has bounded the first expression's relation term exponentially. Its remaining source-band kernel still needs an evaluation in the central window. The second still needs more than RTM9's two different tail powers. The existing trace-mass cancellations and GTM20's bound \(O_h(q\log q)\) do not evaluate either term with an \(o(q)\) residual.

No value of \(d_{\log}(h)\) or \(d_0(h)\) is asserted here. In particular, \(-\mathfrak I_h^{\mathrm{cen}}\) is not asserted to equal either of them.

\subsection{9. D032 check and the comparison still not established}\label{acw:d032-check-and-the-comparison-still-not-established}

The supplied full D032 collection was consulted at the relevant weight theorem and tensor-amplification argument. Theorem 3.3.1 on the original printed page 204 is stated for schemes of finite type over \(\mathbb Z\), not only for a single finite field. The original paper's printed page 140 gives the Kunneth/product argument that halves an existing weight error. These passages were checked against the original IAS PDF of Pierre Deligne, \emph{La conjecture de Weil. II}, Publications mathematiques de l'IHES 52 (1980), 137--252.

The current attempt has not constructed a comparison from the actual signed kernel expression ACW24 to a D032 cohomological object that returns its required value or metric estimate. This is an uncompleted comparison in this attempt, not a proof that no such mathematical correspondence exists. D032's weight bounds are not substituted for ACW24's unknown signed kernel correlation.

\subsection{Source and verification record}\label{acw:source-and-verification-record}

Imported programme inputs: RTM1--3, RTM8--11, GTM1 and GTM20, and SXR1--12. Exact source ranges and file hashes are recorded in \texttt{sources.json}. The current status and precision target were checked in \emph{Untitled 1775}. The source history in D032 was retained; no historical translation was silently promoted to its latest maintained revision.

Classical Gamma asymptotic: NIST DLMF 5.11.9 / 5.11.1, with the exact substitution written in Section 5. Fourier inversion and Plancherel are used after the explicit integrability bounds ACW8--9.

The accompanying verifier checks exact finite source/relation, monic-frame, endpoint-sign and mass-cancellation identities, together with the central area algebra on an explicitly artificial Gaussian source. These are auxiliary checks. They are not numerical values for the hypothetical xi packet, not an interval certificate for it, not an independent verification of all imported source estimates, and not a formal verification of this analytic derivation.

\textbf{Final status:} central-source profile and its integrated \(k^2\) and \(k\) coefficients derived, with an exponential bound for the actual central relation-kernel term; the original signed arithmetic-baseline coefficients remain uncomputed. The user's requested full completion has not been achieved.

% END PRESERVED INPUT provider_CENTRAL_SOURCE_WINDOW.tex
\endgroup

\clearpage
\begingroup


\setcounter{equation}{0}
\renewcommand{\theequation}{LRC.\arabic{equation}}
\section*{Complete original conductor transport and affine fibre proof}
\addcontentsline{toc}{section}{Complete original conductor transport and affine fibre proof}
% BEGIN PRESERVED INPUT provider_LRC_TRANSPORT_COMPLETE.tex SHA256 eadce8964ae128a4c1f5aab8361f46f6ae0b876d3aef2e6892093bf008953ca2
% Additive, independent reconstruction of incoming LRC6--10.
% Original WCF/NCT/TAC files are unchanged. Human bibliography keys retained.
\begingroup
\section{Uniform all-rank transport for the original conductor}
\label{lrc:sec:LRC-transport-audit}

Fix the original simple quartet, period, branch and conductor coefficients.
The programme origin retains its original attribution
\cite{human:globalization2026,human:splitzero2026}. Put
\[
\begin{gathered}
 a=4\ell+1\ge9,\qquad 0<\delta<\tfrac12,\qquad \gamma>2,
 \qquad q=(a+1)^2,\qquad Q=(a-7)^2,\\
 \beta_{n,u,w}=\tfrac n2+(2u-n)\delta+i(2w-n)\gamma,
 \qquad E_A(z)=\sum_{u,w=0}^{8}a_{uw}e^{\beta_{8,u,w}z},\\
 \mu_j=E_A^{(j)}(0),\qquad
 v=\operatorname{ord}_0E_A\le80,\quad \mu_v\ne0,
 \quad A_{\rm abs}=\sum_{u,w=0}^{8}|a_{uw}|,\\
 s\in\{1,a\},\quad b=s/2,\quad c_0=a/2,\quad c_1=a/2-4,
 \quad D\in\mathbb Z_{\ge0},\quad N=D+v,\quad d=D+1.
\end{gathered}                                                   \tag{LT1}
\]
The actual coefficients, $v,\delta,\gamma$ are fixed as $a$ varies.
The inner product is conjugate-linear in its first argument. Retain
the complete original measure and all its factors:
\[
\begin{gathered}
 dm_s(y)=\frac{(2\pi)^{s/2}2^{s/2}}{4\pi\Gamma(s/2)}
       |\Gamma(s/4+iy/2)|^2\,dy,\qquad M_s=(2\pi)^{s/2},\\
 \|P\|_{s,c}^2=\int_{\mathbb R}|P(c+iy)|^2\,dm_s(y),\qquad
 p_{0,s}=1,\quad p_{1,s}=y,\\
 p_{n+1,s}=yp_{n,s}-n(n-1+s/2)p_{n-1,s},\qquad
 h_{n,s}=M_sn!(s/2)_n,\\
 \phi_{n,s,c}(S)=p_{n,s}((S-c)/i)/\sqrt{h_{n,s}},\qquad
 [S^n]\phi_{n,s,c}=i^{-n}/\sqrt{h_{n,s}},\qquad
 \rho_n=\sqrt{n!/(s/2)_n}.
\end{gathered}                                                   \tag{LT2}
\]
These are the original orthonormal frames. To check the classical
dictionary with its mass, take
$p_{n,s}(y)=n!P_n^{(s/4)}(y/2;\pi/2)$ in the Meixner--Pollaczek
orthogonality formulas, DLMF (18.19.8)--(18.19.9)
\cite{human:dlmf18}. Their hypotheses hold because $s/4>0$.
With $y=2x$, the squared norm is exactly
\[
 \frac{M_s2^{s/2}}{4\pi\Gamma(s/2)}\,
 2(n!)^2\frac{2\pi\Gamma(n+s/2)}{2^{s/2}n!}
 =M_sn!(s/2)_n .
\]
The generating identity at these same parameters is DLMF (18.23.7)
\cite{human:dlmf18}:
\[
 \sum_{n\ge0}\frac{p_{n,s}(y)}{n!}t^n
 =(1+t^2)^{-s/4}e^{y\arctan t},\qquad |t|<1.          \tag{LT3}
\]
It also follows directly by multiplying the recurrence in LT2 by
$t^n/n!$: the series satisfies
$(1+t^2)\partial_t F=(y-(s/2)t)F$, $F(y,0)=1$.
The analytic factors and $\arctan$ take their branches fixed at zero.

For $\mathcal P_j=\mathbb C[S]_{\le j}$ and
$\mathcal P_{-1}=0$, the literal conductor is
\[
 (\mathcal T_AP)(S')=\sum_{u,w=0}^{8}a_{uw}P(S'+\beta_{8,u,w}),
 \qquad
 A_{D,s}:(\mathcal P_N/\mathcal P_{v-1},\|\cdot\|_{s,c_0,\rm quot})
       \longrightarrow(\mathbb C[S']_{\le D},\|\cdot\|_{s,c_1}).
                                                               \tag{LT4}
\]
Indeed $\mathcal T_AS^n=\sum_{j=v}^n\binom nj\mu_jS'^{n-j}$.
The leading coefficient for $n\ge v$ is $\binom nv\mu_v\ne0$,
and all polynomials of degree below $v$ vanish. Elimination by degree
proves surjectivity and that the full kernel is $\mathcal P_{v-1}$.
The quotient is represented isometrically by
$\phi_{v,s,c_0},\ldots,\phi_{D+v,s,c_0}$.
Set
\[
 w(t)=-i\arctan t=\tfrac12\log\frac{1-it}{1+it},\qquad
 g(t)=t^{-v}e^{-4w(t)}E_A(w(t)),\qquad
 g(0)=(-i)^v\mu_v/v! .                                  \tag{LT5}
\]
The zero-order definition of $v$ removes the singularity at zero.
Apply $\mathcal T_A$ to LT3 at $y=(S-c_0)/i$.
Since $c_0-c_1=4$, its image is $t^vg(t)$ times the same generating
function at $y=(S'-c_1)/i$. Comparing coefficients, including
$n!/\sqrt{h_{n,s}}=\rho_n/\sqrt{M_s}$, gives the exact matrix
\[
 A_{D,s}=\bigl[\operatorname{diag}(\rho_i)_{i=0}^D\bigr]^{-1}
             T_D(g)\operatorname{diag}(\rho_{i+v})_{i=0}^D,
 \quad T_D(g)_{ij}=\begin{cases}g_{j-i},&i\le j,\\0,&i>j.\end{cases}
                                                               \tag{LT6}
\]
Both copies of $M_s$ remain in the original norms; they cancel only
in this matrix connecting two frames with the same original order.
No scalar $v!/\mu_v$ is inserted.

\subsection{The finite analytic constants in LRC6}

Define, exactly as in the incoming calculation,
\[
\begin{gathered}
 g_\flat(t)=(1+t^2)^4g(t),\qquad
 F_0=4+4\sqrt{\delta^2+\gamma^2},\\
 B_0^{\flat}=2^{v+8}A_{\rm abs}e^{2\pi\gamma},\qquad
 B_2^{\flat}=B_0^{\flat}
 \{F_0(F_0+1)+4vF_0+4v(v+1)\}.
\end{gathered}                                                   \tag{LT7}
\]
The superscript on $B_0^{\flat},B_2^{\flat}$ only prevents collision
with WCF14's different zero-free-disk constant $B_0$; their numerical
values are the incoming LRC6 constants, without alteration.
For $\lambda_{uw}=(\beta_{8,u,w}-4)/2$, put
$\alpha_{uw}=4+\lambda_{uw}$ and $\beta^{\flat}_{uw}=4-\lambda_{uw}$.
Then
\[
\begin{gathered}
 t^vg_\flat(t)=H_\flat(t):=
 \sum_{u,w}a_{uw}(1-it)^{\alpha_{uw}}(1+it)^{\beta^{\flat}_{uw}},
 \\ |\Re\lambda_{uw}|\le4\delta<2,
 \quad |\Im\lambda_{uw}|\le4\gamma,
 \quad |\alpha_{uw}|,|\beta^{\flat}_{uw}|\le F_0.
\end{gathered}            \tag{LT8}
\]
For $|t|<1$, $z=(1-it)/(1+it)$ satisfies
$\Re z=(1-|t|^2)/|1+it|^2>0$. Each $1\pm it$ has positive
real part, and the branches fixed at zero obey
$\log z=\log(1-it)-\log(1+it)$ throughout the disk.
For nonnegative integers $j,k$ with $j+k\le2$, the real exponents
in $(1-it)^{\alpha_{uw}-j}(1+it)^{\beta^{\flat}_{uw}-k}$ are
positive: the smallest possible one is $2-4\delta>0$.
Their sum is $8-j-k$, while the imaginary contribution to its
modulus is $\exp[-\Im\lambda_{uw}\arg z]\le e^{2\pi\gamma}$.
Consequently
\[
 |(1-it)^{\alpha_{uw}-j}(1+it)^{\beta^{\flat}_{uw}-k}|
 \le 2^{8-j-k}e^{2\pi\gamma}.                              \tag{LT9}
\]
For $f=(1-it)^\alpha(1+it)^\beta$ its first derivative has
coefficients $-i\alpha,i\beta$, and its second derivative has
coefficients $-\alpha(\alpha-1),2\alpha\beta,-\beta(\beta-1)$.
LT9 therefore proves
\[
\begin{split}
 |H_\flat|&\le2^8A_{\rm abs}e^{2\pi\gamma},\\
 |H_\flat'|&\le2^7(2F_0)A_{\rm abs}e^{2\pi\gamma},\\
 |H_\flat''|&\le2^6\{2F_0(F_0+1)+2F_0^2\}
                      A_{\rm abs}e^{2\pi\gamma}
 \le2^8F_0(F_0+1)A_{\rm abs}e^{2\pi\gamma}.
\end{split}                                                     \tag{LT10}
\]
On $1/2\le|t|<1$, use $g_\flat=t^{-v}H_\flat$ and
\[
 g_\flat''=t^{-v}H_\flat''-2vt^{-v-1}H_\flat'
                         +v(v+1)t^{-v-2}H_\flat .
\]
It follows that $|g_\flat|\le B_0^{\flat}$ and
$|g_\flat''|\le B_2^{\flat}$ on this annulus. Both functions
are holomorphic on the disk, since $H_\flat$ has a zero of order
$v$ at zero. Applying the maximum-modulus principle on each circle
of radius $r\in[1/2,1)$ proves
\[
 \boxed{|g_\flat(t)|\le B_0^{\flat},\qquad
 |g_\flat''(t)|\le B_2^{\flat}\quad(|t|<1).}              \tag{LRC6}
\]
If $g_\flat(t)=\sum_{n\ge0}b_nt^n$, Cauchy's integral formula
on $|t|=r<1$, followed by $r\uparrow1$, gives
\[
 |b_0|,|b_1|\le B_0^{\flat},\qquad
 |b_n|\le B_2^{\flat}/(n(n-1))\quad(n\ge2).               \tag{LT11}
\]
This argument has retained every interior zero of $g$.

\subsection{Exact matrix factorization and every exterior rank}

Let $S_D$ be the upper shift, with entries $1$ in positions $(i,i+1)$,
and set
\[
 R=\operatorname{diag}(\rho_i)_{i=0}^D,\qquad
 W=R^{-1}S_DR,\qquad V_v=\operatorname{diag}(\rho_{i+v}/\rho_i)_{i=0}^D,
 \qquad U_\flat=R^{-1}T_D(g_\flat)R .                    \tag{LT12}
\]
For any two series analytic at zero, $T_D(fh)=T_D(f)T_D(h)$:
the $(i,j)$ entry is exactly the finite convolution for coefficient
$j-i$, and all entries below the diagonal vanish. Since $S_D^{D+1}=0$,
all power-series evaluations at $S_D$ and at $W$ are finite.
The identity $g=g_\flat(1+t^2)^{-4}$ and LT6 thus give
\[
 \boxed{A_{D,s}=U_\flat(I+W^2)^{-4}V_v.}                 \tag{LRC7}
\]
This fixes the position of $V_v$ on the right. It is never commuted
through the other two matrices.

The nonzero images of different coordinate vectors under $W^n$ are
orthogonal, so its norm is the maximum of $\rho_{i+n}/\rho_i$.
At $s=a$, each factor $(i+j+1)/(i+j+s/2)$ is at most one, hence
$\|W^n\|\le1$ and $\|V_v\|\le1$. At $s=1$,
\[
 \frac{\rho_{i+n}^2}{\rho_i^2}
 =\prod_{j=0}^{n-1}\frac{i+j+1}{i+j+1/2}
 \le\frac{n!}{(1/2)_n}\le2n+1.                         \tag{LT13}
\]
The first inequality follows because every factor decreases with $i$.
For the second, induction starts at $n=0$; multiplying $2n+1$ by
$(n+1)/(n+1/2)$ gives $2n+2\le2n+3$.
In particular $\|V_v\|\le\sqrt{2v+1}$.
Now $U_\flat=\sum_{n=0}^D b_nW^n$. For $n\ge2$,
\[
 \frac{\sqrt{2n+1}}{n(n-1)}\le\sqrt{10}\,n^{-3/2},\qquad
 \sum_{n=2}^\infty n^{-3/2}
 \le\int_1^\infty x^{-3/2}\,dx=2.
\]
LT11--13 prove the original finite constant
\[
 \|U_\flat\|\le H:=B_0^{\flat}(1+\sqrt3)+2\sqrt{10}B_2^{\flat}.
                                                               \tag{LT14}
\]
Put
\[
 H_s=\begin{cases}\max\{1,H\sqrt{2v+1}\},&s=1,\\
                  \max\{1,H\},&s=a,\end{cases}
 \qquad
 \kappa_s=\begin{cases}1+\sqrt{8/3},&s=1,\\2,&s=a.\end{cases}
                                                               \tag{LT15}
\]
Writing $Q_D=I+W^2$, triangularity gives $\det Q_D=1$.
For $s=1$, the decreasing-factor argument gives the sharper
$\|W^2\|\le\rho_2/\rho_0=\sqrt{8/3}$, including $D<2$ when
$W^2=0$. For $s=a$ it gives $\|W^2\|\le1$.
Thus $\|Q_D\|\le\kappa_s$ in both cases.

For any invertible complex $n\times n$ matrix $T$ with singular
values $\sigma_1\ge\cdots\ge\sigma_n>0$, diagonalizing its
positive squared norm, or taking its singular-value decomposition,
gives
\[
 \|\textstyle\bigwedge^rT\|=\prod_{j=1}^r\sigma_j,
 \qquad
 \|\textstyle\bigwedge^rT^{-1}\|
 =\frac{\|\bigwedge^{n-r}T\|}{|\det T|}
 \quad(0\le r\le n).                                  \tag{LT16}
\]
The second equality multiplies the largest inverse singular values;
its denominator is the full determinant, not a truncated product.
Apply it to $Q_D^{\oplus c}$, where $c\in\{1,4\}$, whose dimension
is $cd$ and determinant one. Exterior functoriality and
submultiplicativity give
\[
 \|\bigwedge^r\bigl((Q_D^{-4})^{\oplus c}\bigr)\|
 \le\|\bigwedge^r\bigl((Q_D^{-1})^{\oplus c}\bigr)\|^4
 =\|\bigwedge^{cd-r}\bigl(Q_D^{\oplus c}\bigr)\|^4
 \le\kappa_s^{4(cd-r)}.                                 \tag{LT17}
\]
This is valid at $r=cd$, where the right side equals one.
For a complex matrix $A=U\Sigma V^*$,
$A^{\mathsf T}=\overline V\Sigma U^{\mathsf T}$ is another
singular-value decomposition; ordinary transposition therefore
preserves every exterior norm. Apply LRC7 to $c$ orthogonal copies,
use LT14 for the first factor and LT13 for the last, and then LT17.
For $1\le r\le cd$ this proves exactly
\[
 \boxed{\log\|\bigwedge^r\bigl((A_{D,s}^{\mathsf T})^{\oplus c}\bigr)\|
 \le F_{D,s,r,c}:=r\log H_s+4(cd-r)\log\kappa_s.}         \tag{LRC8}
\]
Set $F_{D,s,0,c}=0$, the exact logarithm of the empty exterior norm.

LT6 also proves the full complex determinant identity
\[
 \det A_{D,s}=\bigl((-i)^v\mu_v/v!\bigr)^d
                         \prod_{i=0}^{D}\frac{\rho_{i+v}}{\rho_i}.
                                                               \tag{LT18}
\]
The phase is retained in this equality. At $s=1$ the product of
positive ratios is at least one. At $s=a$,
\[
 \log\prod_{i=0}^{D}\frac{\rho_i}{\rho_{i+v}}
 =\tfrac12\sum_{i=0}^{D}\sum_{j=0}^{v-1}
       \log\left(1+\frac{s/2-1}{i+j+1}\right)
 \le\tfrac12(s/2-1)v\sum_{j=1}^d\frac1j .                \tag{LT19}
\]
Here the elementary inequality $\log(1+x)\le x$ applies to each
nonnegative $x$; the sums are empty if $v=0$.
Define $\log^+x=\max\{0,\log x\}$, $H_d^{\rm harm}=\sum_{j=1}^d1/j$,
and retain WCF's exact finite allowance
\[
 J_{D,s}=d\log^+(1/|g(0)|)
       +\tfrac12(s/2-1)_+vH_d^{\rm harm}.                \tag{LT20}
\]
LT18--19 give $|\det A_{D,s}|^{-1}\le e^{J_{D,s}}$.
Apply LT16 to the whole matrix
$B=(A_{D,s}^{\mathsf T})^{\oplus c}$ and LRC8 to its complementary
rank. For $1\le r\le cd$,
\[
 \boxed{\log\|\bigwedge^r\bigl((A_{D,s}^{-\mathsf T})^{\oplus c}\bigr)\|
 \le I_{D,s,r,c}:=cJ_{D,s}+F_{D,s,cd-r,c}.}              \tag{LRC9}
\]
Set $I_{D,s,0,c}=0$. This does not assert that the displayed
positive-rank formula extends to rank zero. At full rank it gives
$I_{D,s,cd,c}=cJ_{D,s}$; at $D=0$ the original matrix is the scalar
$g(0)\rho_v$. Thus the smallest cutoff, $v=0$, empty exterior powers
and full exterior powers are all included. No assumption on absence
or number of interior zeros of $g$ has entered these estimates.

There is an explicit constant for the claimed uniform order. Suppose
$D\le C_{\rm cut}a^2$ with fixed $C_{\rm cut}\ge0$, and put
\[
 \begin{gathered}
 K=C_{\rm cut}+1,\qquad G=\log^+(1/|g(0)|),\qquad
 H_* =\max\{1,H\sqrt{2v+1}\},\qquad
 \kappa_*=1+\sqrt{8/3},\\
 L_* =\max\{\log H_*,4\log\kappa_*\},\qquad
 J_*=KG+\tfrac v4(3+\log K).
 \end{gathered}                                                   \tag{LT21}
\]
Then $d\le Ka^2\le Kq$ and
$H_d^{\rm harm}\le1+\log d\le1+\log K+2\log a$.
Since $\log a\le a$, $a\le q$ and $a^2\le q$, LT20 implies
$J_{D,s}\le J_*q$. Moreover the expression in LRC8 is bounded by
$cdL_*$. Therefore for every retained rank, including zero,
\[
 0\le F_{D,s,r,c}\le cKL_*q,\qquad
 0\le I_{D,s,r,c}\le c(J_*+KL_*)q .                    \tag{LT22}
\]
All constants depend only on the actual conductor and the displayed
cutoff constant. These inequalities prove uniform $O(q)$ bounds at
all ranks and both original source orders, with no rank restriction
beyond $0\le r\le c(D+1)$.

\subsection{The literal lower covariance and every attained quotient}

Now impose the original endpoint domain $N\ge q-1$, $D=N-v$.
Let $\mathsf E_-$ and $\mathsf E_+$ be the full original observation
matrices in the unchanged root order:
\[
 (\mathsf E_-)_{n,(u,w)}=\phi_{n,s,c_1}(\beta_{a-8,u,w}),\ 0\le n\le D,
 \quad
 (\mathsf E_+)_{n,(u,w)}=\phi_{n,s,c_0}(\beta_{a,u,w}),\ 0\le n\le N.
                                                               \tag{LT23}
\]
The roots are distinct because their real and imaginary parts recover
the two indices. Since $q-Q=16a-48\ge96>v$, we have $D\ge Q-1$.
The leading coefficients in LT2 and the Vandermonde determinant prove
that $\mathsf E_-$ has rank $Q$ and $\mathsf E_+$ has rank $q$.
Define the original conductor matrix $C_a$ by collecting every term
in the identity
\[
 \sum_{u,w}(C_a^{\mathsf T}\eta)_{uw}e^{\beta_{a,u,w}z}
 =E_A(z)\sum_{u,w=0}^{a-8}\eta_{uw}e^{\beta_{a-8,u,w}z}.   \tag{LT24}
\]
This uses exactly
$\beta_{a-8,u,w}+\beta_{8,r,t}=\beta_{a,u+r,w+t}$.
All contributions producing the same root are retained.
If the left side is zero, the lower exponential sum vanishes on
an open set where $E_A\ne0$, and hence everywhere. Its derivative
Vandermonde forces $\eta=0$; thus $C_a^{\mathsf T}$ is injective.

For a lower functional $\Lambda$, complex linearity gives
$\Lambda(\mathcal T_A\phi_{v+j,s,c_0})
=\sum_i(A_{D,s})_{ij}\Lambda(\phi_{i,s,c_1})$.
This proves the ordinary-transpose square, including the vanishing
first $v$ rows,
\[
 \mathsf E_+C_a^{\mathsf T}
 =\begin{bmatrix}0_{v\times Q}\\A_{D,s}^{\mathsf T}\mathsf E_-\end{bmatrix},
 \quad K_{-,D,s}=\mathsf E_-^*\mathsf E_->0,\quad
 K_{C,N,s}=\mathsf E_-^*\overline{A_{D,s}}A_{D,s}^{\mathsf T}\mathsf E_->0.
                                                               \tag{LT25}
\]
This is the original WCF22 map with the unchanged orders $s=1,a$,
the lower centre $a/2-4$, and the shifted cutoff $N-v$.

For independent columns $e_1,\ldots,e_r$, the determinant of their
Gram matrix is $\|e_1\wedge\cdots\wedge e_r\|^2$. This follows
by expanding both sides as the determinant of their pairwise inner
products; it is the classical exterior Gram identity, for example
Lyons, \S4, (4.1) \cite{human:lyons2003}.
Apply $\bigwedge^r B$ to this wedge and $\bigwedge^r B^{-1}$
to the image. LRC8--9 give the two directed determinant bounds.
Taking $E=\mathsf E_-$, $B=A_{D,s}^{\mathsf T}$ and $r=Q$ proves
\[
 \boxed{-2I_{D,s,Q,1}\le e^C_{N,s}
 :=\log\det K_{C,N,s}-\log\det K_{-,N-v,s}
 \le2F_{D,s,Q,1}.}                                      \tag{LRC10}
\]

For completeness, the same argument controls every compatible attained
quotient, rather than only a restriction. Let $K_0\subset Y\subset H$
be subspaces of a finite-dimensional Hilbert space and $B:H\to H'$
an invertible linear map. Put
$Y_0=Y\cap K_0^\perp$ and $Y_1=BY\cap(BK_0)^\perp$.
Every affine fibre modulo $K_0$ has a unique minimum-norm element in
$Y_0$, by its orthogonal decomposition. The induced map and its inverse
in these minimum representatives are exactly
\[
 \overline B=P_{Y_1}B|_{Y_0},\qquad
 \overline B^{-1}=P_{Y_0}B^{-1}|_{Y_1}.                  \tag{LT26}
\]
Indeed subtraction in the first formula lies in $BK_0$ and subtraction
in the second lies in $K_0$, proving that the compositions are identities
on quotient classes and hence on their unique minimum representatives.
Projections and isometric inclusions contract every exterior norm,
as is immediate from their zero and unit singular values.
It follows that each exterior norm of $\overline B$ and its inverse
is bounded by the corresponding exterior norm of $B$ and $B^{-1}$.
For any common quotient coordinate frame of rank $r$, the quotient
Gram determinant ratio therefore lies in
\[
 -2I_{D,s,r,c}\ \le\ \log\frac{\det G_1}{\det G_0}
                    \ \le\ 2F_{D,s,r,c}                 \tag{LT27}
\]
when $B=(A_{D,s}^{\mathsf T})^{\oplus c}$.
A succession of restrictions and quotients reduces to the full
preimage pair $K_0\subset Y$ in the original ambient space: the
union of the affine fibres minimized at successive steps is exactly
the final affine fibre. Its minimum is attained by the same orthogonal
projection. Thus the complete denominators remain, and no extra factor
is paid for the dimension of a denominator or the number of steps.

Retain $N_i=(q-1,q,2q-1,2q)_i$, $D_i=N_i-v$ and
$\mathcal R f=-f(N_0)-f(N_1)+f(N_2)+f(N_3)$.
With $f_i^C=F_{D_i,s,Q,1}$ and $i_i^C=I_{D_i,s,Q,1}$, LRC10 gives
\[
 -2(f_0^C+f_1^C+i_2^C+i_3^C)
 \le\mathcal Re^C\le
 2(i_0^C+i_1^C+f_2^C+f_3^C).                            \tag{LT28}
\]
The actual cutoffs have $D_i\le2(a+1)^2\le3a^2$ for $a\ge9$;
LT22 proves that both finite allowances are $O_{\rm actual}(q)$.
The precise native $\Xi$ application of LT27, including its full
annihilators, coordinate section and original lower Gram denominators,
is given next. All earlier independently proved upper and lower
allowances may be retained by taking the minimum separately on each
directed side. The symbol, its zeros and their pole directions have
never been removed from the original map.
\endgroup


% Exact Xi receiver, reconstructed from NCT1--14,20 and TAC1b--i,5,7,9--11.
% The containing transport proof supplies the proved exterior bounds F and I.
\begingroup
\subsection{The exact analytic comparison and its four signed endpoints}
\label{lrc:sec:xi-exact-map}

Retain the actual simple quartet, original period and logarithm branch,
all coefficients of the original nonzero conductor, and the original root
order. Write
\[
\begin{gathered}
 a=4l+1\ge9,\quad q=(a+1)^2,\quad Q=(a-7)^2=q',\\
 m=q-Q-v,\quad N\ge q-1,\quad D=N-v,\quad
 d=D+1,\quad t=N-q+1,\\
 s\in\{1,a\},\quad c_0=a/2,\quad c_1=a/2-4,\quad
 \beta_{n,u,w}=n/2+(2u-n)\delta+i(2w-n)\gamma,\\
 E_A(z)=\sum_{r,h=0}^{8}a_{rh}e^{\beta_{8,r,h}z},\quad
 \mu_j=E_A^{(j)}(0),\quad v=\operatorname{ord}_0E_A\le80,\quad
 \mu_v\ne0,\quad 0<\delta<1/2,\quad\gamma>2.
\end{gathered}\tag{XIM1}
\]
Thus $t\ge0$, $q-Q=16a-48\ge96$, and $m>0$. All polynomial
spaces have complex coefficients. The original Hermitian norm and frame are
\[
\begin{gathered}
 dm_s(y)=\frac{(2\pi)^{s/2}2^{s/2}}{4\pi\Gamma(s/2)}
           |\Gamma(s/4+iy/2)|^2\,dy,\qquad M_s=(2\pi)^{s/2},\\
 p_{0,s}=1,\quad p_{1,s}=y,\quad
 p_{n+1,s}=yp_{n,s}-n(n-1+s/2)p_{n-1,s},\\
 h_{n,s}=M_sn!(s/2)_n,\qquad
 \phi_{n,s,c}(S)=p_{n,s}((S-c)/i)/\sqrt{h_{n,s}},\qquad
 \|P\|_{s,c}^2=\int_{\mathbb R}|P(c+iy)|^2dm_s(y).
\end{gathered}\tag{XIM2}
\]
The frame is the original orthonormal frame, with the unchanged
Meixner--Pollaczek dictionary and norm
\cite[eqs.~(18.19.8)--(18.19.9), (18.22.8)]{human:dlmf18}; its leading coefficient is
$i^{-n}/\sqrt{h_{n,s}}$. Neither the mass nor this phase is changed.
The duals below are complex-linear duals, with the dual Hilbert norms.

Put $\mathcal P_j=\mathbb C[S]_{\le j}$, with $\mathcal P_{-1}=0$,
and use $S'$ for the target variable. The unscaled conductor map and its
matrix in the displayed orthonormal frames are
\[
\begin{gathered}
 \mathcal T_AP(S')=\sum_{r,h=0}^{8}a_{rh}P(S'+\beta_{8,r,h}),\\
 A=A_{D,s}:\mathcal P_N/\mathcal P_{v-1}\longrightarrow
             \mathbb C[S']_{\le D},\\
 \mathcal T_A\phi_{v+j,s,c_0}
     =\sum_{i=0}^D A_{ij}\phi_{i,s,c_1},\qquad 0\le j\le D,\\
 \rho_n=\sqrt{n!/(s/2)_n},\quad
 w(z)=-i\arctan z,\quad
 g(z)=z^{-v}e^{-4w(z)}E_A(w(z))=\sum_{j\ge0}g_jz^j,\\
 g_0=(-i)^v\mu_v/v!,\qquad
 A_{ij}=\begin{cases}g_{j-i}\rho_{j+v}/\rho_i,&i\le j,\\
                         0,&i>j.
          \end{cases}
\end{gathered}\tag{XIM3}
\]
These are the literal NCT3--5 and WCF4--6 matrices. Indeed
$\mathcal T_AS^n=\sum_{j=v}^n\binom nj\mu_jS'^{n-j}$, with
nonzero leading coefficient $\binom nv\mu_v$ when $n\ge v$.
It follows by elimination in degree that the kernel is exactly
$\mathcal P_{v-1}$ and that $A$ is an isomorphism. The quotient
norm uses its original orthogonal representatives
$\phi_{v,s,c_0},\ldots,\phi_{N,s,c_0}$.
For the matrix entries, multiplying the recurrence in XIM2 by
$z^n/n!$ gives the differential equation
$(1+z^2)\partial_zG=(y-(s/2)z)G$, $G(y,0)=1$, whose solution is
$G(y,z)=(1+z^2)^{-s/4}\exp(y\arctan z)$.
Applying $\mathcal T_A$ to this generating identity centred at $c_0$
multiplies the same identity centred at $c_1$ by $z^vg(z)$, since
$c_0-c_1=4$. Coefficient comparison and
$n!/\sqrt{h_{n,s}}=\rho_n/\sqrt{M_s}$ give exactly XIM3.

Here are the full relation and annihilator spaces. Define
\[
 \chi_a(S)=\prod_{u,w=0}^{a}(S-\beta_{a,u,w}),\qquad
 \chi_-(S')=\prod_{u,w=0}^{a-8}(S'-\beta_{a-8,u,w}).
 \tag{XIM4}
\]
Use brackets $[P]_-$ for the coefficient column of $P$ in
$\phi_{0,s,c_1},\ldots,\phi_{D,s,c_1}$, and $[P]_+$ for the
coefficient column of its class modulo $\mathcal P_{v-1}$ in
$\phi_{v,s,c_0},\ldots,\phi_{N,s,c_0}$. Form the matrices
\[
\begin{gathered}
 \mathsf X_N=\bigl[[\chi_aS^j]_+\bigr]_{0\le j<t}
       \in\mathbb C^{d\times t},\qquad
 \mathsf R_N=\bigl[[\mathcal T_A(\chi_aS^j)]_-\bigr]_{0\le j<t}
       =A\mathsf X_N\in\mathbb C^{d\times t},\\
 \mathsf L_N=\bigl[[\chi_-S'^j]_-\bigr]_{0\le j<d-Q}
       \in\mathbb C^{d\times(d-Q)}.
\end{gathered}\tag{XIM5}
\]
At $t=0$, the first two matrices have no columns. The upper relation
columns are independent modulo $\mathcal P_{v-1}$ because their
nonzero linear combinations have degree at least $q>v-1$. The
native relation columns have distinct degrees $q+j-v$ and nonzero
leading coefficients $\binom{q+j}{v}\mu_v$. The ideal columns are
monic of distinct degrees $Q+j$. Thus their respective ranks are
$t,t,d-Q$.

For every lower root, the identity
$\beta_{a-8,u,w}+\beta_{8,r,h}=\beta_{a,u+r,w+h}$ shows that
$\mathcal T_A(\chi_aS^j)$ vanishes at that root. Distinctness of
the roots follows by comparing real and imaginary parts. Division
by the monic $\chi_-$ therefore has zero remainder. Consequently
$\operatorname{im}\mathsf R_N\subset\operatorname{im}\mathsf L_N$.
For a functional with coordinate column
$z=(\Lambda(\phi_{i,s,c_1}))_{i=0}^D$, evaluation on a polynomial
with coefficient column $p$ is $z^{\mathsf T}p$. Its norm is
$\|z\|_2$. The two exact annihilator spaces are therefore
\[
\begin{gathered}
 \mathscr Y_D=\ker\mathsf R_N^{\mathsf T}
    =(\operatorname{im}\overline{\mathsf R_N})^\perp,
       \qquad \dim\mathscr Y_D=d-t=q-v,\\
 \mathscr K_D=\ker\mathsf L_N^{\mathsf T}
    =(\operatorname{im}\overline{\mathsf L_N})^\perp,
       \qquad \dim\mathscr K_D=Q,\\
 \mathscr K_D\subset\mathscr Y_D,\qquad
 \dim(\mathscr Y_D/\mathscr K_D)=m.
\end{gathered}\tag{XIM6}
\]
Complex conjugation in the orthogonal-complement expressions is
necessary because the algebraic annihilator uses ordinary transpose.

We identify these spaces with the original analytic functionals and
lower evaluations, including their frames. Let
$\mathbb V[n,\beta]=\beta^n$, $0\le n<q$, be the upper
Vandermonde in the original root order, and set
\[
\begin{gathered}
 \mathscr N_v=\{\nu\in\mathbb C^q:(\mathbb V\nu)_n=0
                                     \ (0\le n<v)\},\\
 C_{(u,w),(i,j)}=
 \begin{cases}a_{i-u,j-w},&0\le i-u,j-w\le8,\\0,&\text{otherwise},\end{cases}
       \qquad C\in\mathbb C^{Q\times q},\\
 \mathcal N_\nu(z)=\sum_\beta\nu_\beta e^{\beta z},\qquad
 \mathcal N_{C^{\mathsf T}\eta}(z)
   =E_A(z)\sum_{u,w=0}^{a-8}\eta_{uw}e^{\beta_{a-8,u,w}z}.
\end{gathered}\tag{XIM7}
\]
The last identity collects all original shift coefficients. Its zero
left side forces the lower exponential sum to vanish wherever
$E_A\ne0$, hence everywhere by analyticity; the lower Vandermonde
then forces $\eta=0$. Thus $C^{\mathsf T}$ is injective and its
image lies in $\mathscr N_v$.

Let $P_v$ insert the last $q-v$ coordinates into $\mathbb C^q$ and
put $J_v=P_v^{\mathsf T}\mathbb VC^{\mathsf T}$. Retain the original
first nonzero $Q$-row minor $\mathcal I$ of $J_v$, its complementary
row set $\mathcal J$, and the coordinate insertion $E_{\mathcal J}$.
Then
\[
 Z=\mathbb V^{-1}P_vE_{\mathcal J}\in\mathbb C^{q\times m},
 \qquad [C^{\mathsf T},Z]\text{ is a frame of }\mathscr N_v,
 \qquad F_b=\mathcal N_{Ze_b}/E_A.
 \tag{XIM8}
\]
The frame assertion follows by applying $\mathbb V$, taking the
$\mathcal I$ rows first, and then taking the complementary rows.
Each $F_b$ is analytic at zero because the first $v$ numerator
derivatives vanish.

These are also exactly the two packet annihilators in TAC1g--h.
Put $E=\mathbb C[S]/(\chi_a)$, $E_-=\mathbb C[S']/(\chi_-)$,
$L_v=\mathbb C[S]_{<v}\subset E$, and
$V_A=\ker(\overline{\mathcal T}_A:E\longrightarrow E_-)$.
Root addition proves that this packet map is well defined, and in
root-value coordinates it is $C$. Since $C^{\mathsf T}$ is injective,
the map is onto and $\dim V_A=q-Q$. The degree calculation in XIM3
gives $L_v\subset V_A$. The exact complex-linear pairing is
\[
\begin{gathered}
 \Phi:\mathbb C^q\longrightarrow E^\vee,\qquad
 \Phi(\nu)([P])=\sum_\beta\nu_\beta P(\beta),\\
 \Phi(\mathscr N_v)=\operatorname{ann}(L_v),\qquad
 \Phi(\operatorname{im}C^{\mathsf T})=\operatorname{ann}(V_A),\\
 \Psi:\mathscr N_v/\operatorname{im}C^{\mathsf T}
       \longrightarrow (V_A/L_v)^\vee,\quad
 \Psi([\nu])([P]+L_v)=\sum_\beta\nu_\beta P(\beta).
\end{gathered}\tag{XIM8a}
\]
The upper Vandermonde makes $\Phi$ invertible. Its first $v$ moment
equations are precisely vanishing on $L_v$. The conductor dual
functionals vanish on $V_A$, and both that image and
$\operatorname{ann}(V_A)$ have dimension $Q$, proving their equality.
Restriction therefore has exactly this kernel. Every functional on
$V_A/L_v$ extends linearly from $V_A$ to $E$ while vanishing on
$L_v$, proving that $\Psi$ is onto and hence an isomorphism.
At $v=0$, $L_v=0$ and the corresponding moment conditions are empty.

For $\nu\in\mathscr N_v$, let
$\Lambda_{\mathcal N_\nu/E_A}(S'^r)=(\mathcal N_\nu/E_A)^{(r)}(0)$.
Leibniz's finite identity proves
\[
 \Lambda_{\mathcal N_\nu/E_A}(\mathcal T_AP)
       =\Lambda_{\mathcal N_\nu}(P),\qquad \deg P\le N.
 \tag{XIM9}
\]
It vanishes for $P=\chi_aS^j$, so the native functional belongs
to $\mathscr Y_D$. If its first $D+1$ derivatives vanish, XIM9
makes all numerator derivatives through $N$ vanish. The first $q$
such equations have invertible Vandermonde, so $\nu=0$.
The dimension $q-v$ now proves that these are all of $\mathscr Y_D$.
Every lower evaluation annihilates the ideal generated by $\chi_-$.
Their first $Q$ evaluation rows form an invertible triangular change
of the lower Vandermonde, since $D\ge Q-1$. They are independent,
and dimension $Q$ proves that they are all of $\mathscr K_D$.

For completeness, the finite analytic observation matrix can be
constructed without an unevaluated division of germs. Put
$n_{jb}=\sum_\beta Z_{\beta b}\beta^j$ and $f_{rb}=F_b^{(r)}(0)$.
For $0\le r\le D$, successively set
\[
 f_{rb}=
 \frac{n_{v+r,b}-\displaystyle\sum_{j=v+1}^{v+r}
          \binom{v+r}{j}\mu_j f_{v+r-j,b}}
      {\binom{v+r}{v}\mu_v}.
 \tag{XIM10}
\]
At $r=0$ the sum is empty. All denominators are nonzero. This is
exactly the derivative product identity, hence retains every original
coefficient. Writing $\phi_{n,s,c_1}(S')=\sum_r\kappa_{nr}S'^r$,
define
\[
\begin{gathered}
 H=\mathsf E_-\in\mathbb C^{d\times Q},\qquad
 H_{n,\beta'}=\phi_{n,s,c_1}(\beta'),\\
 J\in\mathbb C^{d\times m},\qquad
 J_{nb}=\sum_{r=0}^n\kappa_{nr}f_{rb},\qquad
 \mathscr K_D=\operatorname{im}H,\quad
 \mathscr Y_D=\operatorname{im}[H,J].
\end{gathered}\tag{XIM11}
\]
Here $[H,J]$ has full column rank $Q+m=q-v$.

Let $\mathsf E_+[n,\beta]=\phi_{n,s,c_0}(\beta)$, $0\le n\le N$,
and let $\mathsf E_+^{\rm red}$ consist of rows $v,\ldots,N$.
The upper observation matrix has full column rank because its first
$q$ rows are an invertible triangular change of the upper Vandermonde.
The first $v$ rows vanish on $\mathscr N_v$. The ordinary dual
matrix $B=A^{\mathsf T}$ therefore gives the exact equations
\[
\begin{gathered}
 BH=\mathsf E_+^{\rm red}C^{\mathsf T},\qquad
 BJ=\mathsf E_+^{\rm red}Z,\qquad
 \mathsf E_+C^{\mathsf T}=\begin{bmatrix}0_{v\times Q}\\BH\end{bmatrix},
 \qquad \mathsf E_+Z=\begin{bmatrix}0_{v\times m}\\BJ\end{bmatrix},\\
 B\mathscr Y_D=\ker\mathsf X_N^{\mathsf T},\qquad
 B\mathscr K_D=\ker(A^{-1}\mathsf L_N)^{\mathsf T}
                =\operatorname{im}(\mathsf E_+^{\rm red}C^{\mathsf T}).
\end{gathered}\tag{XIM12}
\]
For the first line, apply each functional to the basis equation XIM3
and use XIM9. For the second line, use $\mathsf R_N=A\mathsf X_N$
and $B=A^{\mathsf T}$; invertibility proves both inclusions are
equalities. This proves the full commutative square on actual finite
jets. No adjoint and no scalar $\alpha_A=v!/\mu_v$ occurs.

The polynomial source of the analytic quotient is also exact.
Restriction of functionals induces isometric isomorphisms
\[
 \frac{\operatorname{ann}(\operatorname{im}\mathsf R_N)}
      {\operatorname{ann}(\operatorname{im}\mathsf L_N)}
   \simeq
 \left(\frac{\operatorname{im}\mathsf L_N}
               {\operatorname{im}\mathsf R_N}\right)^\vee,
 \quad
 \frac{\operatorname{ann}(\operatorname{im}\mathsf X_N)}
      {\operatorname{ann}(\operatorname{im}A^{-1}\mathsf L_N)}
   \simeq
 \left(\frac{\operatorname{im}A^{-1}\mathsf L_N}
               {\operatorname{im}\mathsf X_N}\right)^\vee.
 \tag{XIM13}
\]
The primal metrics first restrict the original polynomial metric and
then minimize over the full relation space. To prove the dual metric
claim, extend a functional from the larger displayed polynomial
subspace by composition with its orthogonal projection. This is the
minimum extension, since every other extension differs on the
orthogonal complement and the squared norms add. It still vanishes
on the relation subspace. All such extensions differ by precisely
the displayed denominator annihilator. This proves isometry as well
as the algebraic quotient assertion. The primal map is $A$ on these
two polynomial subquotients, and XIM12 is its contravariant ordinary
dual. Each primal subquotient has dimension $(d-Q)-t=m$.
The original packet identification is
\[
\begin{gathered}
 V_A/L_v\longrightarrow
 \operatorname{im}(A^{-1}\mathsf L_N)/\operatorname{im}\mathsf X_N,
 \quad [P]+L_v\longmapsto[[P]_+],\\
 Y_N=\mathbb C[S']_{\le D-Q}\big/
   \{\mathcal T_A(\chi_a h)/\chi_-:h\in\mathcal P_{t-1}\},\\
 \theta_N:V_A/L_v\longrightarrow Y_N,\qquad
 [P]+L_v\longmapsto[\mathcal T_AP/\chi_-].
\end{gathered}
 \tag{XIM13a}
\]
for the upper polynomial quotient and its divided native counterpart,
respectively. A representative of degree at most $N$ exists since
$N\ge q-1$. Changing it by $\chi_a h$ gives exactly an XIM5 relation,
and changing it by a polynomial of degree below $v$ gives zero.
Conversely, every polynomial in $\operatorname{im}(A^{-1}\mathsf L_N)$
has a degree-at-most-$N$ representative whose conductor vanishes on
the lower roots, hence represents $V_A$. A zero quotient class differs
from a multiple of $\chi_a$ by a polynomial of degree below $v$.
This proves both surjectivity and injectivity. Dividing the native
relation by $\chi_-$ is legitimate by XIM5; its relation map is the
original $h\mapsto\mathcal T_A(\chi_a h)/\chi_-$. Thus XIM13a is
exactly the TAC1d map, including its full relation space and scalar one.

We now give both attained coefficient minima. Put
\[
\begin{gathered}
 K_0=H^*H>0,\quad K_1=H^*B^*BH>0,\qquad
 P_0=HK_0^{-1}H^*,\quad P_1=BHK_1^{-1}H^*B^*,\\
 W_0=(I-P_0)J,\qquad W_1=(I-P_1)BJ,\\
 D_{a,D}^{(s)}=W_0^*W_0,\qquad
 \widehat D_{a,D}^{(s)}=W_1^*W_1.
\end{gathered}\tag{XIM14}
\]
The first metric is the original NCT10 native analytic metric. The
second is precisely NCT8/TAC5: setting $K_+=\mathsf E_+^*\mathsf E_+$,
XIM12 gives
\[
 \widehat D_{a,D}^{(s)}=
 Z^*\left[K_+-K_+C^{\mathsf T}
    (\overline C K_+ C^{\mathsf T})^{-1}\overline C K_+\right]Z.
 \tag{XIM15}
\]
For every $x\in\mathbb C^m$, completing the two positive squares
in the Schur-complement calculation
\cite[Appendix A.5.5, p.~650; Appendix C.4.1, p.~673]{human:boyd-vandenberghe}
proves
\[
\begin{aligned}
 x^*D_{a,D}^{(s)}x
   &=\min_{\eta\in\mathbb C^Q}\|Jx+H\eta\|_2^2,
 &\eta_0(x)&=-K_0^{-1}H^*Jx,\\
 x^*\widehat D_{a,D}^{(s)}x
   &=\min_{\eta\in\mathbb C^Q}\|B(Jx+H\eta)\|_2^2,
 &\eta_1(x)&=-K_1^{-1}H^*B^*BJx.
\end{aligned}\tag{XIM16}
\]
Both minima retain every lower evaluation and all their cross terms.
Each minimizer is unique because $K_0,K_1$ are positive definite.
More explicitly, the difference between the first squared norm and
its displayed minimum is
$(\eta-\eta_0(x))^*K_0(\eta-\eta_0(x))$, and the difference for
the second is $(\eta-\eta_1(x))^*K_1(\eta-\eta_1(x))$.
These identities follow on expanding, using the two displayed
formulas for $\eta_0,\eta_1$; each cross term cancels exactly.
If either minimum were zero, invertibility of $B$ and independence
of $[H,J]$ would imply $x=0$. Thus both analytic metrics are positive
definite. The two minimizing coefficient columns need not coincide.

The induced map on the quotient has the following exact domain,
codomain, and inverse:
\[
\begin{gathered}
 Q_0=\mathscr Y_D\cap\mathscr K_D^\perp=\operatorname{im}W_0,
 \quad Q_1=B\mathscr Y_D\cap(B\mathscr K_D)^\perp
                         =\operatorname{im}W_1,\\
 \overline B:Q_0\longrightarrow Q_1,
       \quad \overline B=P_{Q_1}B|_{Q_0},\qquad
 \overline B^{-1}=P_{Q_0}B^{-1}|_{Q_1},\qquad
 \overline B(W_0x)=W_1x.
\end{gathered}\tag{XIM17}
\]
The quotient identifications are isometric because each representative
has a unique orthogonal minimum lift. Projection after $B$ subtracts
an element of $B\mathscr K_D$; applying $B^{-1}$ turns that into an
element of $\mathscr K_D$, which the second projection removes.
This proves both inverse identities. It also proves the last equation
using XIM14. In the orthonormal minimum-lift frames
$W_0(D_{a,D}^{(s)})^{-1/2}$ and
$W_1(\widehat D_{a,D}^{(s)})^{-1/2}$, the matrix is
\[
 [\overline B]=
 (\widehat D_{a,D}^{(s)})^{1/2}(D_{a,D}^{(s)})^{-1/2},\qquad
 [\overline B]^*[\overline B]=
 (D_{a,D}^{(s)})^{-1/2}\widehat D_{a,D}^{(s)}(D_{a,D}^{(s)})^{-1/2}.
 \tag{XIM18}
\]
Indeed $W_1^*BW_0=W_1^*W_1=\widehat D_{a,D}^{(s)}$, because
the discarded part lies in $B\mathscr K_D$. Thus the original TAC9
correction is exactly
\[
 \epsilon_{a,N,s}
 =\log\det\widehat D_{a,D}^{(s)}-\log\det D_{a,D}^{(s)}
 =2\log|\det[\overline B]|.
 \tag{XIM19}
\]

Write the established all-rank logarithmic exterior estimates for the
actual ambient matrix, including its original weights, as
\[
 \log\left\|\bigwedge^r\bigl((A_{D,s}^{\mathsf T})^{\oplus c}\bigr)\right\|
       \le F_{D,s,r,c},\qquad
 \log\left\|\bigwedge^r\big((A_{D,s}^{\mathsf T})^{\oplus c}\big)^{-1}
       \right\|\le I_{D,s,r,c},\qquad 0\le r\le cd.
 \tag{XIM20}
\]
The $F$ and $I$ in this display are logarithms of exterior norms,
without the Gram factor two. Orthogonal projections contract every
exterior power: in orthonormal bases their singular values are zero
or one. Restrictions and inclusions of Hilbert subspaces also
contract exterior powers. Applying XIM17 in both directions and
then XIM20 at $c=1$, $r=m$, proves the finite bound
\[
 \boxed{-2I_{D,s,m,1}\le\epsilon_{a,N,s}\le2F_{D,s,m,1}.}
 \tag{XIM21}
\]
There is no added $Q$-rank allowance: the complete denominator has
already been retained in the exact quotient map and in both minima.
For $c$ copies, the denominator has dimension $cQ$ and the full
analytic quotient dimension $cm$. Any further restriction followed
by an attained quotient has its full preimages $K'\subset Y'$ in
the ambient $c$ copies; necessarily $K'$ contains the original
$cQ$-dimensional denominator. Nested minimization over these fibres
equals minimization over $K'$, because the union of the original
affine fibres is the final affine fibre. XIM17 applies to $K'\subset Y'$
and proves the same determinant interval $[-2I_{D,s,r,c},2F_{D,s,r,c}]$
with its actual final rank $r$, including $r=0$ with empty determinant
one. No intermediate dimension is charged a second time.

Finally retain the actual TAC endpoints and dual sign orientation:
\[
 N_i=(q-1,q,2q-1,2q)_i,\quad D_i=N_i-v,\quad
 \Xi_{a,s}=-\epsilon_{a,N_0,s}-\epsilon_{a,N_1,s}
             +\epsilon_{a,N_2,s}+\epsilon_{a,N_3,s}.
 \tag{XIM22}
\]
For these four endpoints,
$d_i=(q-v,q+1-v,2q-v,2q+1-v)_i$ and
$t_i=(0,1,q,q+1)_i$. The annihilator dimensions remain exactly
$q-v$ and $Q$, and the analytic quotient rank remains $m$ at every
endpoint. Set $F_i=F_{D_i,s,m,1}$ and $I_i=I_{D_i,s,m,1}$.
Negating the two low-endpoint inequalities in XIM21 and adding the
two high-endpoint inequalities gives
\[
 \boxed{
 -2(F_0+F_1+I_2+I_3)
       \le\Xi_{a,s}\le
  2(I_0+I_1+F_2+F_3).}
 \tag{XIM23}
\]
These are bounds on the same four correlated matrices appearing in
TAC9. A separate radius infimum, minimum of established allowances,
or improved ambient estimate can be inserted in each $F_i$ or $I_i$
provided it bounds XIM20; no attainment of an infimum is used.
The source order remains $1$ or $a$ at the lower grid, its centre
remains $a/2-4$, and its cutoff remains $N_i-v$. In particular no
replacement by the order $a-8$ or by an unshifted cutoff occurs.
The explicit uniform bounds LT21--22 of the preceding transport proof
apply because $D_i\le3a^2$ for $a\ge9$. Inserting them in XIM23
proves $\Xi_{a,s}=O_{\rm actual}(q)$ for both original source orders.
The identical endpoint bound XIM21 holds at every cutoff
$q-1\le N\le2q$, including each growing-prefix endpoint $N=q+R$
with $0\le R<q$, with the same fixed constant.
\endgroup

% END PRESERVED INPUT provider_LRC_TRANSPORT_COMPLETE.tex
\endgroup

\clearpage
\begingroup


\setcounter{equation}{0}
\renewcommand{\theequation}{AKS.\arabic{equation}}
\section*{Complete absolute common-kernel and full-numerator estimates}
\addcontentsline{toc}{section}{Complete absolute common-kernel and full-numerator estimates}
% BEGIN PRESERVED INPUT provider_AKS.tex SHA256 bdb8a06e8b1f6562e3fbd3e78c3b33670d9784ddba82e56d6a6b42a86783a280
% AKS: new continuation, 14 September 2026.
% This body is self-contained at the level of its explicitly stated inherited inputs.
\section{The absolute common kernel: the new finite result}
\label{aks:sec:AKS}

This complete incoming calculation sharpens the current absolute common
kernel estimate. Its original comparison was against the preceding
logarithmic bound; the intervening KAF chapter already established
the $O(kq)$ scale. The new finite envelope reduces its uniform limsup
constant from $16(4+10\log2)$ to $32\log(25\log3/(4\pi))$.
AKJ below takes the pointwise minimum of both finite bounds. No
limiting kernel coefficient is inferred from these upper bounds.

Retain the actual hypothetical complete quartet and its entire quotient
\[
 \rho=\tfrac12+\delta+i\gamma,\qquad 0<\delta<\tfrac12,\quad \gamma>2,
 \qquad h(z)=\prod_{\alpha\in\{\rho,\bar\rho,1-\rho,1-\bar\rho\}}
 (z-\alpha)^m,\quad v_h=2\xi/h.
\]
The full primary unit is $U=M_{j_hv_h}$. The original tensor order, centre,
polynomial, and coefficient quotient are
\[
\begin{gathered}
 k=4l+1,\quad l\ge1,\quad e=1+k(m-1),\quad q=e(k+1)^2,\quad c=k/2,\\
 \lambda_{a,b}=c+(2a-k)\delta+i(2b-k)\gamma,\qquad
 \chi(S)=\prod_{a,b=0}^k(S-\lambda_{a,b})^e,
 \quad E=\mathbb C[S]/(\chi).
\end{gathered}\tag{AKS1}
\]
For the principal conclusion $m=1$. The source-operator estimates below
also hold with the stated repeated-root list for $m>1$, but no simple-quartet
rank is assigned to those multiplicities.

Let $\Lambda_k(u):E\twoheadrightarrow B_k(u)$ be the original observation,
$K(u)=\ker\Lambda_k(u)$, and $I:K(u)\hookrightarrow E$ its original
coefficient inclusion. For a source $s\in\{1,k\}$, let $G_N^{(s)}$ be the
original minimum-remainder Gram, in the frame $(1,S,\ldots,S^{q-1})$.
Define $H_{K,N}^{(s)}=I^*G_N^{(s)}I$ and retain exactly
\[
 F_K^{(s)}=
 \log\det H_{K,q-1}^{(s)}+\log\det H_{K,q}^{(s)}
 -\log\det H_{K,2q-1}^{(s)}-\log\det H_{K,2q}^{(s)}.
 \tag{AKS2}
\]
Every empty determinant is one. All conjugate transposes refer to the
specified coordinate matrices, not to a bilinear residue pairing.

Here are entirely explicit positive numbers used in the bound:
\[
\begin{gathered}
 R=k\sqrt{\delta^2+\gamma^2},\quad b_s=s/2,\quad
 t_q=\frac\pi2-\frac1q,\quad \mathfrak a_q=\frac{2\log3}{t_q},\quad
 \mathfrak s_q=\frac{\mathfrak a_q^q-1}{\mathfrak a_q-1},\\
 \mathfrak b_{s,q}=\left[\frac9{25}\sin\frac1q\right]^{-b_s}
       \exp\left(2R\left(\log3+\frac{t_q}{2}\right)\right),\\
 \mathcal Z_{s,j}=1+\mathfrak b_{s,q}\mathfrak s_q^2
        \sum_{d=q}^{q+j-1}\frac{d!}{(b_s)_d}
                          \left(\frac{25}{16}\right)^d,
 \quad \mathcal L_{s,j}=\log \mathcal Z_{s,j}\quad (0\le j\le q+1).
\end{gathered}\tag{AKS3}
\]
The empty sum gives $\mathcal Z_{s,0}=1$. In the admitted family $q\ge36$, so
$t_q\in(0,\pi/2)$ and $\mathfrak a_q>1$. These parameters are radii in the existing
Laplace and generating functions; neither is a new source order.

\paragraph{Main theorem (absolute finite bound).}
For the original simple quartet, every original logarithm branch, and
every $|u|\ge R_*$, where $R_*$ is precisely RPC21 or RPC23 (recorded in
Section~\ref{aks:sec:AKS-radius}), put $r=8k-16$. Then
\[
 \boxed{\quad
 0\le F_K^{(s)}(k,u)\le r\bigl(\mathcal L_{s,q}+\mathcal L_{s,q+1}\bigr),
 \qquad s\in\{1,k\}.\quad}
 \tag{AKS4}
\]
There is a stronger, fully period-dependent interval in AKS28 below.
In particular, with
\[
 \kappa_* =32\log\left(\frac{25\log3}{4\pi}\right),
 \qquad 25.02078097649<\kappa_*<25.02078097650,
 \tag{AKS5}
\]
one has the uniform statement
\[
 \boxed{\quad
 0\le\limsup_{k\to\infty\atop k\equiv1\ (4)}
       \ \sup_{\substack{|u|\ge R_*\\\text{original branches}}}
         \frac{F_K^{(s)}(k,u)}{kq}\le\kappa_*.
 \quad}\tag{AKS6}
\]
Here the source may be $s=1$ for every $k$, or $s=k$. The same bound holds
for the actual arithmetic kernel. Its exact finite allowances are in
AKS34--36. Thus the absolute kernel is $O_h(kq)$ on this domain, not only
$O_h(kq\log q)$. Neither AKS5 nor AKS6 calls $\kappa_*$ the actual leading
coefficient.

\section{An exact coefficient realization of the actual lost directions}
\label{aks:sec:AKS-frame}

\subsection{The period matrix and invariant restriction, with their full coefficients}
For $m=1$, order the one-factor roots as $(++),(+-),(-+),(--)$, and let
$\mathcal E_{\rm prim}$ have their original Lagrange idempotents as columns.
Set $\zeta=\exp(2\pi i/5)$ and $F_{jt}=\zeta^{jt}-1$, $1\le j,t\le4$.
The exact inverse is $(F^{-1})_{tj}=\zeta^{-tj}/5$. Retain
\[
\begin{gathered}
 \Pi[P]_j=\int_{\Gamma_j}e^{\Phi_h(z)/u}\operatorname{rem}_hP(z)\,dz,
 \qquad\Phi_h'=h,\quad\Phi_h(0)=0,\\
 \Gamma_j=-\ell_0+\ell_j,\quad
 \arg\ell_j=(\arg u+\pi+2\pi j)/5,\\
 \mathsf H=F^{-1}\Pi U\mathcal E_{\rm prim},\quad
 D=\operatorname{diag}(\zeta,\zeta^2,\zeta^3,\zeta^4),\\
 Q(x)=Q_0(\mathsf H^{-1}x),\qquad
 Q_0(t)=t_1t_4-t_2t_3,\quad Q_a(x)=Q(D^{-a}x).
\end{gathered}\tag{AKS7}
\]
The integral always uses the degree-less-than-four representative.
Invertibility of $\Pi$, all contour phases, and its full determinant are
the PDE/PQO results in the supplied proofs. The four unit values are
$(a,\bar a,\bar a,a)$, with the actual $a=v_h(\rho)\ne0$; no value of $a$
is selected. Hence $\mathsf H$ is the original invertible matrix.

The following formulas make the coefficient realization directly
computable. Write
\[
 \mathcal I_{k,0}=\{\nu\in\mathbb Z_{\ge0}^4:|\nu|=k,
                         \ \nu_1+2\nu_2+3\nu_3+4\nu_4\equiv0\pmod5\}.
\]
For $\nu\in\mathcal I_{k,0}$ define $B_{\nu;(a,b)}$ by summing over
$4\times4$ arrays $n_{t,\alpha}\ge0$ with row sums $\nu_t$, total plus-real
count $a$, and total plus-imaginary count $b$:
\[
 B_{\nu;(a,b)}=
 \sum_{\substack{\sum_\alpha n_{t,\alpha}=\nu_t\\
       \sum_{t,\alpha:\epsilon_\alpha=+}n_{t,\alpha}=a\\
       \sum_{t,\alpha:\eta_\alpha=+}n_{t,\alpha}=b}}
 \left(\prod_{t=1}^4\frac{\nu_t!}{\prod_\alpha n_{t,\alpha}!}\right)
                  \prod_{t,\alpha}\mathsf H_{t,\alpha}^{n_{t,\alpha}}.
 \tag{AKS8}
\]
Let $\mathsf V_{(a,b),d}=\lambda_{a,b}^d$ for $0\le d<q$ and set
$\mathsf A=B\mathsf V$. All entries of $\mathsf A$ are now finite
expressions in the actual periods and unit. There is no generic period
matrix in this definition.

For completeness, put $t_\alpha(w)=\exp((\eta_\alpha\gamma-
 i\epsilon_\alpha\delta)w)$ and $x(w)=\mathsf Ht(w)$. Expansion gives
\[
 x(w)^\nu=\sum_{a,b}B_{\nu;(a,b)}e^{\omega_{a,b}w},\qquad
 \omega_{a,b}=(2b-k)\gamma-i(2a-k)\delta.
\]
In the original primary algebra,
$e^{wY}[1]=\sum_{a,b}e^{\omega_{a,b}w}\varepsilon_{a,b}$, where
$Y=(M_S-cI)/i$. The exponentials are independent by the Vandermonde of
the distinct $\omega_{a,b}$. OCS9's unscaled symmetric word basis
therefore identifies the original observation with
\[
       \jmath\Lambda_k=\mathcal F_0\mathsf A,
       \qquad\ker\mathsf A=K(u)
       \quad\text{in the original coefficient frame.}
       \tag{AKS9}
\]
Every multinomial in AKS8 belongs to its actual ordered words; the
word multiplicity is not inserted a second time. The exact inherited
row Gram of $\mathcal F_0$ is the restriction to $\mathcal I_{k,0}$ of
\[
 \mathcal G_{\nu,\mu}=
 \sum_{\substack{n_{tu}\ge0\\\sum_u n_{tu}=\nu_t\\
                              \sum_t n_{tu}=\mu_u}}
 \frac{k!}{\prod_{t,u}n_{tu}!}\prod_{t,u}
          [5(\delta_{tu}+1)]^{n_{tu}}.
 \tag{AKS10}
\]
This follows by taking the coefficient of $\bar x^\nu y^\mu$ in
$(\bar x^TF^*Fy)^k$, since $F^*F=5(I+\mathbf1\mathbf1^*)$. It is positive
definite and includes every row cross term.

Here is the precise role of the orbit ideal. The map sending a homogeneous
polynomial $f$ to the covector specified by
$\Theta_k(f)(e^{wY}[1])=f(x(w))$ has kernel
$Q\mathbb C[x]_{k-2}$ and is onto $E^\vee$. Indeed equal exponential
indices correspond exactly to monomial differences divisible by
$t_1t_4-t_2t_3$; the displayed Vandermonde proves that there are no other
relations. Its invariant image is $\operatorname{im}\Lambda_k^\vee$.
On the five-member domain, the five $Q_a$ are distinct prime quadrics.
An invariant polynomial divisible by $Q$ is divisible by each $Q_a$,
and therefore by their full product; its quotient is invariant by
cancellation in the polynomial ring. Consequently
\[
\begin{gathered}
 \mathbb C[x]^{\langle D\rangle}\cap (Q)
       =\left(\prod_{a=0}^4Q_a\right)\mathbb C[x]^{\langle D\rangle},\\
 K(u)^\vee\simeq
 \mathbb C[x]_k\big/(\mathbb C[x]_k^{\langle D\rangle}
                              +Q\mathbb C[x]_{k-2}),\\
 b:=\operatorname{rank}\mathsf A
     =d_{5,k,0}-d_{5,k-10,0}=k^2-6k+17,\qquad r=q-b=8k-16.
\end{gathered}\tag{AKS11}
\]
Negative degrees are zero. The count is
$d_{5,d,0}=(\binom{d+3}{3}+4\epsilon_d)/5$,
$\epsilon_d=\mathbf1_{d\equiv0(5)}-\mathbf1_{d\equiv1(5)}$;
subtracting the cubics gives the stated expression for $k\ge10$, and
$k=5,9$ give $12,44$ directly. This is the inherited OQX calculation,
repeated to specify which kernel AKS9 realizes. The source metric is
still to be applied to this kernel, not to the invariant ring itself.

\subsection{A determinant formula requiring no assumed nonzero pivot}
Let $\mathcal G_0$ be the full positive row Gram in AKS10 and set
\[
 W=\mathsf A^*\mathcal G_0\mathsf A,\qquad
 \sigma_j=[z^j]\det(I_q+zW),\qquad
 P_K=\frac1{\sigma_b}\sum_{j=0}^b(-1)^j\sigma_{b-j}W^j.
 \tag{AKS12}
\]
Since $W$ has exactly $b$ positive eigenvalues, $\sigma_b>0$. The scalar
polynomial in AKS12 is one at zero and zero at every nonzero eigenvalue:
it is the normalized product of the factors $\lambda_i-z$ for the $b$
positive eigenvalues. Hence $P_K=P_K^*=P_K^2$ and
$\operatorname{im}P_K=K(u)$. This is the projector derived from the
literal observation, not a freely chosen projection. Replacing
$\mathcal G_0$ by any other positive coordinate Gram leaves this same
coefficient-orthogonal kernel projector, since its image and orthogonal
complement are unchanged. No such replacement is needed here.

At every original endpoint define
\[
 D_{K,N}^{(s)}=[z^r]\det(I_q+zP_KG_N^{(s)}P_K).
 \tag{AKS13}
\]
These numbers are strictly positive for $r>0$ and are one for $r=0$.
To prove their exact metric meaning, choose an orthonormal coefficient
basis $U_K$ of $\operatorname{im}P_K$ just for the proof. In a unitary
extension, $P_KG_NP_K$ has its positive block $U_K^*G_NU_K$ and a zero
block. Thus $D_{K,N}=\det(U_K^*G_NU_K)$. Write the original inclusion
as $I=U_KT$; then $|\det T|^2=\det(I^*I)$ and
\[
 \det H_{K,N}^{(s)}=\det(I^*I)D_{K,N}^{(s)},\qquad
 \boxed{\ F_K^{(s)}=
 \log\frac{D_{K,q-1}^{(s)}D_{K,q}^{(s)}}
               {D_{K,2q-1}^{(s)}D_{K,2q}^{(s)}}.\ }
 \tag{AKS14}
\]
This supplies an equivalent exact coefficient determinant formula for
all the $8k-16$ lost directions, including the entire period dependence.
The original frame factor is present at each endpoint and cancels only
with its four actual signs. It is unnecessary to suppose any particular
minor nonzero. Alternatively, selecting the first independent $r$ columns
of $P_K$ gives an explicit coefficient frame; the determinant formula
avoids even that chart choice.

For $m>1$ replace AKS8--9 by OCS12--18's full confluent matrix
$\mathfrak A_{\nu;(a,b,D)}=i^DD!c_{\nu;a,b,D}$ and its primary isomorphism.
AKS12--14 are then identical with the actual rank, not with the rank in
AKS11. All primary nilpotents and all unit derivatives remain in that
matrix.

\section{A bounded exact remainder operator on the original source}
\label{aks:sec:AKS-operator}

\subsection{The complete source norms and coefficient map}
For either existing source set $M_s=(2\pi)^{s/2}$ and retain
\[
 dm_{0,s}(y)=\frac{M_s2^{b_s}}{4\pi\Gamma(b_s)}
           |\Gamma(b_s/2+iy/2)|^2\,dy,\qquad
 \int e^{ty}dm_{0,s}=M_s(\cos t)^{-b_s},\quad |t|<\pi/2.
 \tag{AKS15}
\]
The supplied IGO1--3 proof derives this transform by the beta integral,
Fourier inversion and a contour shift within $|\operatorname{Im}x|<\pi/2$.
In particular it includes the exact masses $M_1=\sqrt{2\pi}$ and
$M_k=(2\pi)^{k/2}$ and exponential domination on every smaller strip.
The original real monic polynomials satisfy
\[
\begin{gathered}
 p_0=1,\quad p_1=y,\quad
 p_{d+1}=yp_d-d(b_s+d-1)p_{d-1},\\
 \sum_{d\ge0}\frac{p_d(y)}{d!}z^d
      =(1+z^2)^{-b_s/2}e^{y\arctan z},\qquad
 h_d^{(s)}=\|p_d\|^2=M_s d!(b_s)_d.
\end{gathered}\tag{AKS16}
\]
The generating formula follows by differentiating in $z$. To check the
normalization directly, integrate the product of two generating
functions at small real $z,w$. AKS15 and the cosine addition identity
give $M_s(1-zw)^{-b_s}$. Taking any finite pair of derivatives at zero,
justified by the stated exponential domination, gives orthogonality and
the norms in AKS16. Thus the leading phase in
$u_d(S)=p_d((S-c)/i)/\sqrt{h_d^{(s)}}$ is exactly $i^{-d}$.
The same generating function agrees with NIST DLMF 18.23.7 after the
original parameter rescaling; it is derived here from the supplied
recurrence, not substituted for it.

List the original roots
$\omega_\nu=(\lambda_{a,b}-c)/i$ in lexicographic order, repeating each
exactly $e$ times, and put
\[
 P(y)=i^{-q}\chi(c+iy)=\prod_{\nu=1}^q(y-\omega_\nu),\qquad
 \Phi:[f(y)]\mapsto[f((S-c)/i)].
 \tag{AKS17}
\]
The inverse is substitution $S=c+iy$, and
$P((S-c)/i)=i^{-q}\chi(S)$. This is an exact algebra isomorphism, including
the complete confluent orders; there is no monomial comparison divisor.
Let $C_s$ have columns $[u_0],\ldots,[u_{q-1}]$. Its nonzero determinant
is $i^{-q(q-1)/2}\prod_{d<q}(h_d^{(s)})^{-1/2}$. Define
\[
 v_j=C_s^{-1}[u_{q+j}],\qquad
 \mathsf K_j=I_q+\sum_{a=0}^{j-1}v_av_a^*,\qquad 0\le j\le q+1.
 \tag{AKS18}
\]
Parseval in the low-degree original orthonormal polynomial basis proves
\[
 \|v_j\|_2^2=
 \frac{\|\operatorname{rem}_P p_{q+j}\|_{m_{0,s}}^2}
      {M_s(q+j)!(b_s)_{q+j}}.
 \tag{AKS19}
\]
This is the full norm of the actual remainder, not a diagonal entry of
an approximated quotient Gram.

\subsection{Newton interpolation with all roots and their orders}
For $0\le a<q$ define $N_a(y)=\prod_{\nu=1}^a(y-\omega_\nu)$, with
$N_0=1$. For every complex $v$, the unique entire-function remainder is
\[
 \operatorname{rem}_P(e^{vy})
   =\sum_{a=0}^{q-1}[\omega_1,\ldots,\omega_{a+1}]e^{v\cdot}\ N_a(y),
 \qquad
 \left|[\omega_1,\ldots,\omega_{a+1}]e^{v\cdot}\right|
       \le\frac{|v|^a}{a!}e^{R|v|}.
 \tag{AKS20}
\]
Here divided differences at repeated nodes mean Hermite divided
differences, not evaluation with a zero denominator.

A proof covering the confluent case is useful. On the standard
$a$-simplex with barycentric coordinates $\theta_1,\ldots,\theta_{a+1}$,
\[
 [\omega_1,\ldots,\omega_{a+1}]e^{v\cdot}
       =v^a\int_{\Delta_a}e^{v\sum_\nu\theta_\nu\omega_\nu}\,d\theta.
\]
The integral of $\prod\theta_\nu^{n_\nu}$ is
$\prod n_\nu!/(a+\sum n_\nu)!$, as follows by iterated integration of
integer beta integrals. Expanding the exponential therefore gives
$\sum_{n\ge0}v^{a+n}h_n(\omega_1,\ldots,\omega_{a+1})/(a+n)!$.
For monomials, the divided difference of $y^{a+n}$ is exactly $h_n$:
this follows by multiplying the geometric series
$\prod(1-\omega_\nu z)^{-1}$ and the Newton division recurrence.
It is a polynomial identity in the nodes and so also holds at repeated
nodes. This proves the integral formula. Iterated Newton division for
polynomials, followed by the convergent exponential series, proves the
first equality in AKS20, including every required derivative.
Finally the simplex has volume $1/a!$ and
$|\sum\theta_\nu\omega_\nu|\le R$, proving its stated bound.

For any $0<t<\pi/2$, put $\tau=2/t$. Positivity of every term of the
cosh series in AKS15 gives
\[
 \|y^a\|_{m_{0,s}}
 \le\sqrt{M_s}(\cos t)^{-b_s/2}\tau^a a!.
\]
Indeed its square is at most
$M_s(\cos t)^{-b_s}(2a)!t^{-2a}$ and
$(2a)!\le4^a(a!)^2$. Expanding the \emph{whole} Newton polynomial, then
using $|e_j(\omega_1,\ldots,\omega_a)|\le\binom ajR^j$, proves
\[
\begin{split}
 \|N_a\|_{m_{0,s}}
 &\le\sqrt{M_s}(\cos t)^{-b_s/2}
        \sum_{j=0}^a\binom ajR^{a-j}\tau^j j!\\
 &=\sqrt{M_s}(\cos t)^{-b_s/2}\tau^a a!
           \sum_{n=0}^a\frac{(R/\tau)^n}{n!}\\
 &\le\sqrt{M_s}(\cos t)^{-b_s/2}\tau^a a!e^{R/\tau}.
\end{split}\tag{AKS21}
\]
There are no omitted polynomial coefficients. Combining AKS20--21 gives
for every complex $v$
\[
 \|\operatorname{rem}_P(e^{vy})\|_{m_{0,s}}
 \le\sqrt{M_s}(\cos t)^{-b_s/2}e^{R(|v|+t/2)}
                       \sum_{a=0}^{q-1}(2|v|/t)^a.
 \tag{AKS22}
\]

\subsection{Cauchy extraction in the actual Gamma generating function}
Fix $0<z<1$ and write $a_z=\operatorname{arctanh}z$.
For $|w|=z$ the defining power series gives
$|\arctan w|\le a_z$ and $|1+w^2|\ge1-z^2$.
Apply the literal finite-dimensional remainder map to AKS16 and use
Cauchy's formula for its $d$th coefficient. The norm triangle inequality
on that circle and AKS22 give
\[
\begin{split}
 \frac{\|\operatorname{rem}_P p_d\|_{m_{0,s}}}{d!}
 \le{}&z^{-d}(1-z^2)^{-b_s/2}\sqrt{M_s}(\cos t)^{-b_s/2}\\
 &\mathrel{}\times e^{R(a_z+t/2)}
                    \sum_{a=0}^{q-1}(2a_z/t)^a.
\end{split}\tag{AKS23}
\]
Cauchy's formula is legitimate for the finite remainder vector: its
entries are holomorphic for $|w|<1$ by the actual matrix generating
identity OPG7. The source norm is taken only after this finite-vector
identity; there is no unjustified integral of a nonintegrable generating
function on the large circle.

Substitution of AKS23 into AKS19 retains the factorial denominator and
cancels the two occurrences of the original mass, giving
\[
 \boxed{\ 
 \|v_j\|_2^2\le\frac{d!}{(b_s)_d}z^{-2d}
 [\cos t(1-z^2)]^{-b_s}e^{2R(a_z+t/2)}
 \left[\sum_{a=0}^{q-1}(2a_z/t)^a\right]^2,
 \quad d=q+j.\ }
 \tag{AKS24}
\]
Adding the \emph{complete} rank-one positive matrices in AKS18 now proves
$\mathsf K_j\preceq \mathcal Z_{s,j}I_q$ after setting $z=4/5$ and $t=t_q$.
Indeed $a_{4/5}=\log3$, $1-z^2=9/25$, and $\cos t_q=\sin(1/q)$, so
the exact scalar obtained is AKS3. The high index $j=q+1$ includes the
original polynomial of degree $2q$.

There is also a root-by-root finite refinement. Put
$\mu_{s,d}=\int y^{2d}dm_{0,s}$,
$R_a=\max_{\nu\le a+1}|\omega_\nu|$, and
\[
\begin{gathered}
 B_{s,a}^{\rm root}=\sum_{d=0}^a
     |e_{a-d}(\omega_1,\ldots,\omega_a)|\sqrt{\mu_{s,d}},\\
 J_s^{\rm root}(z)=\frac{(1-z^2)^{-b_s}}{M_s}
 \left[\sum_{a=0}^{q-1}\frac{a_z^a}{a!}e^{a_zR_a}
                      B_{s,a}^{\rm root}\right]^2.
\end{gathered}\tag{AKS25}
\]
Replace the scalar multiplying
$d!z^{-2d}/(b_s)_d$ in AKS24 by $J_s^{\rm root}(z)$ to get a valid
sharper bound, by the same proof before AKS21 is relaxed. Every moment
is exactly $M_s(2d)![w^{2d}](\cos w)^{-b_s}$. Thus this refinement keeps
every actual root coefficient and a finite original moment formula.
AKS21 bounds it by the explicit scalar in AKS24. The proof never replaces
the source determinant by a product of diagonal moments: AKS24 is a
full-vector operator bound, and the next step restricts the inverse of
the full covariance.

\section{The four original kernel Grams and the evaluated scale}
\label{aks:sec:AKS-volume}

The original orthonormal source map at degree $q+j-1$ is
$C_s[I_q,v_0,\ldots,v_{j-1}]$. Its minimum lift of $C_sx$ is
$[I_q,v_0,\ldots,v_{j-1}]^*\mathsf K_j^{-1}x$. It has that image and is
orthogonal to the full relation kernel, so Pythagoras proves
\[
 G_{q+j-1}^{(s)}=C_s^{-*}\mathsf K_j^{-1}C_s^{-1},\qquad
 H_j:=H_{K,q+j-1}^{(s)}=T_s^*\mathsf K_j^{-1}T_s,\quad T_s=C_s^{-1}I.
 \tag{AKS26}
\]
This verifies the source-minimum step in the coefficient determinant
AKS13--14. Neither a cross block nor any original relation is removed.

Since $I_q\preceq\mathsf K_j\preceq \mathcal Z_{s,j}I_q$, inversion and compression
give $\mathcal Z_{s,j}^{-1}H_0\preceq H_j\preceq H_0$, where $H_0=T_s^*T_s$.
Define $d_j=-\log\det(H_0^{-1/2}H_jH_0^{-1/2})$. Congruence and all
$r$ eigenvalues give
\[
 0=d_0\le d_1\le\cdots\le d_{q+1},\qquad
 0\le d_j\le r\mathcal L_{s,j},\qquad
 F_K^{(s)}=d_q+d_{q+1}-d_1.
 \tag{AKS27}
\]
The identity uses precisely the four endpoints in AKS2.
In particular it proves AKS4, as well as the sharper interval
\[
 \boxed{\quad d_1\le F_K^{(s)}\le
             r(\mathcal L_{s,q}+\mathcal L_{s,q+1})-d_1.\quad}
 \tag{AKS28}
\]
For a directly calculable expression for its lower endpoint set
$V=\|v_0\|^2$ and
$\theta=v_0^*T_s(T_s^*T_s)^{-1}T_s^*v_0$. Then
\[
 0\le\theta\le V,\qquad
 d_1=\log\frac{1+V}{1+V-\theta}.
 \tag{AKS29}
\]
Indeed the inverse of $I+v_0v_0^*$ is
$I-v_0v_0^*/(1+V)$; take its full compressed determinant. For $r=0$
set $\theta=d_1=0$. In GMB/OPR notation $V=t_0$ and $V-\theta=\xi_0$,
so this is exactly the first original section-corrected kernel increment,
not an independently chosen kernel observation.

Here is an evaluated finite remainder for AKS6. Put
\[
\begin{gathered}
 a_\infty=\frac{4\log3}{\pi},\qquad
 c_0=\log\frac{25\log3}{4\pi}>0,\\
 E_{s,q}(R)=\log2+\frac4{\pi-2/q}
  +\frac{s}{2}\log\frac{25\pi q}{18}
  +2R\left(\log3+\frac\pi4\right)\\
 \hspace{23mm}+\log((q+1)(4q+1))-2\log(a_\infty-1).
\end{gathered}\tag{AKS30}
\]
Then for both source orders, every $q\ge36$, and the same actual kernel,
\[
 \boxed{\quad
 \mathcal L_{s,q},\mathcal L_{s,q+1}\le2c_0q+E_{s,q}(R),\qquad
 0\le F_K^{(s)}\le4rc_0q+2rE_{s,q}(R).
 \quad}\tag{AKS31}
\]
To verify every factor, $b_s\ge1/2$ gives
\[
 \frac{d!}{(b_s)_d}\le\frac{4^d}{\binom{2d}{d}}\le2d+1.
\]
The last inequality uses that the central binomial is the largest of the
$2d+1$ positive coefficients whose sum is $4^d$. Thus the sum in AKS3
is at most $(q+1)(4q+1)(5/4)^{4q}$. Also
$\mathfrak s_q\le \mathfrak a_q^q/(a_\infty-1)$ and
\[
 2q\log \mathfrak a_q\le2q\log a_\infty+\frac4{\pi-2/q},\qquad
 \sin(1/q)\ge\frac2{\pi q}.
\]
The first follows from $-\log(1-x)\le x/(1-x)$ at $x=2/(\pi q)$;
the second is the chord bound for sine on $[0,\pi/2]$.
Use $t_q\le\pi/2$ in the remaining exponential. The order-$q$ terms
combine to
$4q\log(5/4)+2q\log a_\infty=2c_0q$.
The upper bound for the logarithm of the added positive scalar is
positive, so $\log(1+e^x)\le\log2+\max(0,x)$ proves the first bound in
AKS31. AKS28 proves the second.

On the original simple family $q=(k+1)^2$, $R=O_h(k)$ and
$r=8k-16$. Consequently $E_{1,q}(R)=O_h(k+\log q)$ and
$E_{k,q}(R)=O_h(k\log q)$. Dividing AKS31 by $kq$ proves AKS6, with
$32c_0=\kappa_*$. In fact the fixed-source upper bound is
\[
 F_K^{(1)}\le(32k-64)c_0q+(16k-32)E_{1,q}(R),
 \tag{AKS32}
\]
whose explicit second term is $O_h(k^2)$. The tensor-source version has
an $O_h(k^2\log q)$ upper remainder. These are remainders of upper
bounds, not assertions about the actual signed error or leading value.
The proof is uniform in all the period parameters in the stated domain:
no norm or condition number of a period matrix occurs in AKS30.

\paragraph{The precise earlier loss.}
OKV15 had already retained
$\|v_j\|^2\le q^22^{6q}(2D_*q)^{2d}/(d!(b_s)_d)$, $d=q+j$.
Replacing its denominator by $1/2$ caused the old $q\log q$ exponent.
Even without AKS20--24 one obtains the valid repair
\[
 \|\mathsf K_{q+1}\|\le
 1+(q+1)q^22^{6q}\frac{(4D_*q)^{4q}}{(4q)!},\qquad D_*\ge3.
 \tag{AKS33}
\]
Here $d!(b_s)_d\ge(2d)!/4^d$ and the remaining factorial ratio increases
up to $d=2q$ for $D_*\ge3$. The elementary inequality
$\log n!\ge n\log n-n+1$ already makes the logarithm of AKS33 $O_h(q)$.
The Newton calculation strengthens this repair to the packet-independent
limsup constant AKS5, while retaining $R$ in its finite remainder.


\paragraph{Exact enclosure of the displayed numerical ceiling.}
For $x>1$ and $z=(x-1)/(x+1)$ the integrated geometric series gives
\[
 0\le \log x-2\sum_{j=0}^{39}\frac{z^{2j+1}}{2j+1}
      \le\frac{2z^{81}}{81(1-z^2)}.
\]
For $0<x<1$, the alternating sum through $j=39$ for $\arctan x$
has error between zero and $x^{81}/81$. The tangent addition formula,
with the angle in $(0,\pi/2)$, proves
$\pi=16\arctan(1/5)-4\arctan(1/239)$: the tangent of
$4\arctan(1/5)-\arctan(1/239)$ is exactly one. Use the resulting
rational lower and upper bounds for $\pi$ and for $\log3$, then apply
the first displayed bound to the two endpoints of
$25\log3/(4\pi)$. Multiplication by $32$ gives, with outward decimal
rounding, the interval
\[
 \kappa_*\in[25.0207809764994871,\ 25.0207809764994872].
\]
The accompanying rational-arithmetic implementation checks these
inequalities by integer cross multiplication, rather than a floating-point
logarithm. This proves the less finely rounded strict interval AKS5.

\section{The complete arithmetic and mixed receivers}
\label{aks:sec:AKS-transfer}

\subsection{The unchanged finite arithmetic allowances}
Use exactly AMT4's constants. To prevent a loss hidden in order notation,
write them here:
\[
\begin{gathered}
 \alpha=\pi/2,\ p=8m-9/2,\ B_h=42+8m,\ M=21+4m,\\
 \vartheta_h=\int_{-1}^1w_h(y)dy,\quad
 M_\alpha=\int e^{\alpha y}w_h(y)dy,\quad
 w_h(y)=|v_h(1/2+iy)|^2/(2\pi),\\
 c_\Gamma=\sqrt2e^{-7/6},\quad C_\Gamma=\sqrt{2\sqrt5}e^{2/3},\\
 a_k=\frac{c_h\vartheta_h^{k-3}e^{-\alpha(k-3)}(k-2)^{-B_h}}
                    {C_\Gamma2^{B_h/2}},\quad
 A_k=\frac{C_h^{\rm tilt}}{c_\Gamma}k^{p+1}M_\alpha^{k-1},\\
 D_N=[1+4(N+M)^2]^M,\quad L_N^{\rm ar}=\log(A_kD_N/a_k),\\
 E_k^+=L_{q-1}^{\rm ar}+L_q^{\rm ar},\qquad
 E_k^-=L_{2q-1}^{\rm ar}+L_{2q}^{\rm ar}.
\end{gathered}\tag{AKS34}
\]
The constants $c_h,C_h^{\rm tilt}$ are the original TW/BT envelope
constants, not new choices. Their actual inequalities and complete
source construction are preserved in AMT4--6 and the current source
closure. Minimize AMT6 on the same remainder fibre and then restrict to
$I(K)$. AMT9--20 gives, with the original arithmetic source $w_h^{*k}$,
\[
 -rE_k^+\le F_K^{\rm ar}-F_K^{(1)}\le rE_k^-.
\]
Let $\underline U_s=d_1^{(s)}$ and
$\overline U_s=r(\mathcal L_{s,q}+\mathcal L_{s,q+1})-d_1^{(s)}$. The new finite conclusion is
\[
 \boxed{\quad
 \max\{0,\underline U_1-rE_k^+\}\le F_K^{\rm ar}
                  \le\overline U_1+rE_k^-.
 \quad}\tag{AKS35}
\]
This improves AMT37's absolute upper bound. For the sharper correlated
version retain $d_{E,N}=q\log A_k-
\log(\det G_N^{\rm ar}/\det G_N^{(1)})$ and
$c_{X,N}=\min\{r_XL_N^{\rm ar},d_{E,N}\}$, where $r_K=r$, $r_B=q-r$.
Put $C_X^-=c_{X,q-1}+c_{X,q}$,
$C_X^+=c_{X,2q-1}+c_{X,2q}$. AMT18 then gives
\[
\begin{split}
 \max\{0,\underline U_1-C_K^-,
                  \underline U_1+\delta_k^\sigma-C_B^+\}
 &\le F_K^{\rm ar}\\
 &\le\overline U_1+
            \min\{C_K^+,\delta_k^\sigma+C_B^-\}.
\end{split}\tag{AKS36}
\]
One may replace $d_{E,N}$ in the caps by the original full-source
$\Xi_{k,N}$ or its proved finite majorant. The scalar
$\delta_k^\sigma=\mathcal B_k^{\rm ar}-\mathcal B_k^\sigma$ retains
exactly AMT21:
\[
\begin{gathered}
 \max\{-qE_k^+,-\Xi_{k,q-1}-\Xi_{k,q}\}
       \le\delta_k^\sigma\\
       \le\min\{qE_k^-,\Xi_{k,2q-1}+\Xi_{k,2q}\}.
\end{gathered}\tag{AKS37}
\]
Since $E_k^\pm=O_h(k+\log q)$ by the exact AMT22 expression, AKS35 and
AKS32 prove the arithmetic version of AKS6, with an explicit
$O_h(k^2)$ upper remainder. Nonnegativity follows from the actual
nested source minima, as in AMT15, not from comparison with a possibly
negative lower endpoint.

\subsection{The original Gamma multiplier, not a replacement higher source}
Retain $t_j=2j+1/2$, $f_l(S)=\prod_{j=0}^{l-1}(S-c-t_j)$ and
\[
 dm_{0,k}=\beta_k|f_l(c+iy)|^2d\sigma,\quad
 \beta_k=\frac{(2\pi)^{k/2}}{\sqrt2\Gamma(k/2)},\quad
 T_l=\prod_{j=0}^{l-1}t_j^2,\quad
 U_{l,N}=\prod_{j=0}^{l-1}[4(N+j+1)^2+t_j^2].
 \tag{AKS38}
\]
The exact source map is $\mathcal P_N\to f_l\mathcal P_N\subset
\mathcal P_{N+l}$, $P\mapsto f_lP$, with inverse division on this
specified subspace. Let $\mathsf E_\chi$ be the complete divided-jet
map and $U_f$ its Leibniz multiplication matrix,
$(U_f)_{r,a}=f_l^{(r-a)}(\lambda)/(r-a)!$ for $r\ge a$.
AMT25--28 gives the precise transported observation and kernel
\[
\begin{gathered}
 \widetilde\Lambda=\Lambda\mathsf E_\chi^{-1}U_f^{-1},\quad
 \widetilde I=U_f\mathsf E_\chi I,\quad
 \ker\widetilde\Lambda=U_f\mathsf E_\chi K,\\
 C_N=\mathscr K_{\chi,N+l}
   -\mathscr K_{\chi,f;N+l}\mathscr K_{f,N+l}^{-1}
                               \mathscr K_{f,\chi;N+l},\\
 H_{K,N}^{(k)}=\beta_k\widetilde I^*C_N^{-1}\widetilde I.
\end{gathered}\tag{AKS39}
\]
Each covariance is the original finite sum of full polynomial jet outer
products divided by its original monic norm. These formulas retain the
extra $f$ constraints and their full cross covariance. In particular
$M_fK=K$ is not assumed. AKS26 calculates the same tensor metric directly
in its own orthonormal source; the two representations agree by their
identical source norm and identical constrained minimum.

Define
\[
 w_N=\log(U_{l,N}/T_l),\quad
 \Omega^{\rm low}=w_{q-1}+w_q,\quad
 \Omega^{\rm high}=w_{2q-1}+w_{2q}.
\]
The already proved same-coefficient comparison AMT42--44 retains all
mass factors before cancellation and gives
\[
 -r\Omega^{\rm low}\le\Delta_K^\Gamma:=F_K^{(1)}-F_K^{(k)}
                         \le r\Omega^{\rm high}.
 \tag{AKS40}
\]
It therefore sharpens the two independent absolute bounds by
\[
\begin{split}
 \max\{\underline U_k,\underline U_1-r\Omega^{\rm high},0\}
       &\le F_K^{(k)}\\
       &\le\min\{\overline U_k,\overline U_1+r\Omega^{\rm low}\}.
\end{split}\tag{AKS41}
\]
The original finite estimate
$0\le w_N\le l\log[1+16(N+l)^2]$ proves that AKS40 is $o(kq)$ on the
actual rank-$r$ domain. Along with AKS35 this preserves the completed
fact that all three source kernels differ by $o(kq)$. AKS4, not AKS40,
is the new absolute-scale calculation.

\subsection{Boundary, QGQ and arithmetic receiving identities}
At each endpoint retain the original minimum section
$S_N^{(s)}=(G_N^{(s)})^{-1}\Lambda^*
(\Lambda(G_N^{(s)})^{-1}\Lambda^*)^{-1}$. The full Schur identity is
\[
 \det(I^*G_N^{(s)}I)\det Q_N^{(s)}
       =|\det[I,S_*]|^2\det G_N^{(s)},\qquad
 Q_N^{(s)}=(\Lambda(G_N^{(s)})^{-1}\Lambda^*)^{-1}.
 \tag{AKS42}
\]
Thus the new finite absolute boundary interval is
\[
 \max\{0,\mathcal B_k^{(s)}-\overline U_s\}
       \le F_B^{(s)}\le\mathcal B_k^{(s)}-\underline U_s.
 \tag{AKS43}
\]
The baselines in this formula are the \emph{complete} original BSL/HBL
objects, with their full finite relation errors. No $q^2$ coefficient is
assigned to $K$; its proved coefficient there remains zero.

Write $\mathcal H_k=\mathcal B_k^\sigma-\mathcal B_k^0=-\mathfrak T_k$.
The exact full identities still read
\[
\begin{gathered}
 F_B^{(1)}-F_B^{(k)}=\mathcal H_k-\Delta_K^\Gamma,\\
 \boxed{\quad\mathcal B_k^{\rm ar}
   =F_K^{(k)}+F_B^{(k)}+\mathcal H_k+\delta_k^\sigma
   =F_K^{(1)}+F_B^{(1)}+\delta_k^\sigma.\quad}
\end{gathered}\tag{AKS44}
\]
The Gamma interval is not re-estimated in this argument. For either
complete QRI1 or QRI2 interval $[\omega^-,\omega^+]$ for its already
calculated centre $W_k$, retain its exact source errors
$e_0\in[-a_0,b_0]$ and $e_q+e_{q+1}\in[-b_k^+,b_k^-]$. Then
\[
 \underline{\mathcal H}_k=-\omega^+-b_k^--a_0,
 \qquad\overline{\mathcal H}_k=-\omega^-+b_k^++b_0,
 \qquad \mathcal H_k\in
       [\underline{\mathcal H}_k,\overline{\mathcal H}_k].
 \tag{AKS45}
\]
These are exactly QRI3, not new unspecified error choices. The complete
QRI1--9, QGT and QGQ definitions remain verbatim in each updated receiver;
the new interval is intersected with all their retained tighter finite
intervals. In particular the simple-quartet nonzero error radius stays
\[
 \limsup\left|
 \frac{\mathcal H_k+lq\mathcal C_\Gamma}{q}
       +\frac{\partial_\alpha\mathcal L(2,\pi)}8
       -\frac{\log2}{4}\right|\le\frac{A_{\rm QGT}}{\sqrt2}.
 \tag{AKS46}
\]
No zero-radius or $o(q)$ arithmetic assertion is introduced.

For example the complete finite absolute defect relative to the tensor
boundary, with the actual Gamma and arithmetic scalars retained, is
\[
 \underline U_k+\mathcal H_k+\delta_k^\sigma
 \le\mathcal B_k^{\rm ar}-F_B^{(k)}
 \le\overline U_k+\mathcal H_k+\delta_k^\sigma.
 \tag{AKS47}
\]
Inserting AKS41 makes this sharper. Replacing $\mathcal H_k$ by the
appropriate endpoints of AKS45 and $\delta_k^\sigma$ by AKS37 gives its
fully inherited finite scalar enclosure with the same signs.
Also $F_B^{\rm ar}=\mathcal B_k^\sigma+\delta_k^\sigma-F_K^{\rm ar}$;
subtract the interval AKS35 or AKS36 to get its new absolute boundary
interval. This retains the exact correlations rather than taking
independent hypothetical source errors.

The existing limit is
$\mathcal H_k/(kq)\to-\mathcal C_\Gamma/4$ and
$\delta_k^\sigma/(kq)\to0$. Therefore the new interval implies
\[
\begin{gathered}
 \liminf_k\inf_{|u|\ge R_*}
 \frac{\mathcal B_k^{\rm ar}-F_B^{(k)}}{kq}
          \ge-\mathcal C_\Gamma/4,\\
 \limsup_k\sup_{|u|\ge R_*}
 \frac{\mathcal B_k^{\rm ar}-F_B^{(k)}}{kq}
          \le\kappa_*-\mathcal C_\Gamma/4.
\end{gathered}\tag{AKS48}
\]
All original branches are included in the infima and suprema. This does
not decide its sign. The already proved boundary \emph{difference}
limits of AMT49 are unchanged. Every subsequential limit of the absolute
normalized kernel lies in $[0,\kappa_*]$ and is shared by its three source
versions when they are evaluated on the same sequence of period
parameters. Existence of a single limit is not proved by boundedness.

Finally the HC residual remains exactly
$\mathcal R_k=\mathcal B_k^{\rm ar}-4q\log(D_hk)$, with its original
$D_h$ and original nonnegative contraction, phase and trace terms.
AKS44 changes none of them. The supplied positive $q^2$ baseline and
its complete Gamma correction are not re-presented as new results here.

\section{Actions, support faces, and the actual arithmetic class}
\label{aks:sec:AKS-actions}

The coefficient formula AKS12 supplies all four blocks of the same sum
action: $P_KYP_K$, $P_KY(I-P_K)$, $(I-P_K)YP_K$ and
$(I-P_K)Y(I-P_K)$. There is no claim that either off-diagonal block
vanishes. In the actual minimum-section coordinates the corresponding
formula remains
\[
 [I,S_N]^{-1}Y[I,S_N]=
 \begin{pmatrix}
 H_{K,N}^{-1}I^*G_NYI&H_{K,N}^{-1}I^*G_NYS_N\\
 \Lambda YI&\Lambda YS_N
 \end{pmatrix}.
 \tag{AKS49}
\]
Its proof is multiplication by the inverse coordinate maps
$H_{K,N}^{-1}I^*G_N$ and $\Lambda$. In the invariant-ring model, the same
action is the derivation
$\mathcal Df=(\mathsf H\Omega\mathsf H^{-1}x)\cdot\nabla f$,
$\Omega=\operatorname{diag}(\gamma-i\delta,-\gamma-i\delta,
\gamma+i\delta,-\gamma+i\delta)$, with $\mathcal DQ=0$.
For the actual conductor $\mathcal A_Q=\prod_{a=1}^4Q_a$, its full defect
is $\mathcal D\mathcal A_Q=\sum_{a=1}^4(\mathcal DQ_a)
\prod_{b\ne a,\,1\le b\le4}Q_b$, exactly OQC22. Thus this coefficient
realization preserves the original action leaving the observed dual.

In particular the section-corrected vector is still OPR2's
\[
 C_s r_j^{\rm ker}=e_{q+j}-S_{q+j-1}^{(s)}\Lambda e_{q+j}
   =I\bigl(\zeta_{q+j}-\kappa S_{q+j-1}^{(s)}\Lambda e_{q+j}\bigr),
\]
with its full denominator $1+\xi_j$ in the kernel increment. The operator
proof estimates the complete covariance and its actual restriction, so
none of these coupled terms is replaced by a raw fixed-section vector.
The all-iterate kernel stays
$\bigcap_{a=0}^{q-1}\ker(\Lambda Y^a)$; nothing in AKS4 declares it zero.

On the existing arithmetic coefficient realization retain
$\mathcal V_{h,k}P=P(D_1+\cdots+D_k)F_h^{\otimes k}$ and
$J^{(k)}\mathcal V_{h,k}=\eta\operatorname{rem}_\chi$, where
$\eta=U^{\otimes k}\iota$ is the original injection. AKS17 and the
Newton coefficients compute precisely this same remainder. Any two
polynomial lifts differ by $\chi Q$ in the original relation space,
so their original arithmetic cohomology classes agree. The specified
primitive of the difference is EP.48, with the full differential and
cochain signs proved in AKP3--5 below; it depends on the lift difference.
The source comparison to $w_h^{*k}$ is the identity on original
polynomial coefficients, followed by the actual minima as in AMT8;
it is not a replacement arithmetic source.

The scalar object stays $G(R)=\{\tau\}\sqcup R^\bullet$, with supported
zero $0_R^\bullet$, injective amplitude/support pair, and infinite
integer target. Every coefficient map above is applied on each original
present support label and outer leg. It commutes with the specified
coefficient inclusions because both compositions apply the same linear
map to the same coefficient; a present zero is not converted into
absence. No new base object, purity condition, or identification of the
sum action with a Frobenius action is used. The actual distinguished
class and its higher-jet directions remain in the full fibre, even
when the observation kills them.

\section{Exact period domain and the interior alternative}
\label{aks:sec:AKS-radius}

For reference the inherited radius used above is completely specified
as follows; these are AMT33--34/RPC17--23, with no new constants. Put
\[
\begin{gathered}
 \beta=2(\gamma^2-\delta^2),\quad\eta=(\delta^2+\gamma^2)^2,\quad
 a=v_h(\rho)\ne0,\\
 c_3=(a^{-2}-\bar a^{-2})/(4i\delta\gamma),\quad
 c_1=(a^{-2}+\bar a^{-2})/2-(\delta^2-\gamma^2)c_3,\\
 g_j=5^{j/5-1}e^{\pi ij/5}\Gamma(j/5),\quad
 g_+=\max_j|g_j|,\quad g_-^{-1}=\max_j|g_j|^{-1},\quad
 \kappa=\max_{j<4}|g_{j+1}/g_j|,\\
 E_j^{\rm ray}=\frac85\int_0^\infty r^j(r^3/3+r)
       e^{-r^5/5+r^3/3+r}dr,\quad
 C_T=2\left(\sum_{j=0}^3(E_j^{\rm ray})^2\right)^{1/2},\\
 C_K=4C_Tg_-^{-1}+2g_+g_-^{-1}+2g_+g_-^{-2}C_T,\\
 C_{K,3}=C_K(3\kappa^2+3\kappa C_K+C_K^2),\quad
 \epsilon_* =\min\{1,(2C_Tg_-^{-1})^{-1}\},\quad B_* =\beta+\eta.
\end{gathered}\tag{AKS50}
\]
Then
\[
R_*=
\begin{cases}
\displaystyle\left[\max\left\{1,\frac{B_*}{\epsilon_*},
\frac{2(|c_3|C_{K,3}B_*+|c_1|(\kappa+C_K))}{|c_3g_4/g_1|}
\right\}\right]^{5/2},&c_3\ne0,\\[4pt]
\displaystyle\left[\max\left\{1,\frac{B_*}{\epsilon_*},
\frac{2C_KB_*}{|g_2/g_1|}\right\}\right]^{5/2},&c_3=0.
\end{cases}\tag{AKS51}
\]
The actual unit implies $(c_1,c_3)\ne(0,0)$ and all integrals converge.
RPC's full period comparison proves a nonzero $(4,1)$ or $(2,1)$
commutator entry, respectively, for $|u|\ge R_*$ on every branch. It
therefore proves the five-member orbit needed in AKS11.

Inside this radius the same coefficient determinant AKS12--14 and the
operator bound AKS24 remain valid. The orbit is decided by the original
noninvariant quadratic coefficients, not by the radius alone. Since
$Q$ is nondegenerate, its projective orbit has either one or five
members. A one-member orbit must have character one: its character has
both fourth and fifth powers equal to one, by the quadratic determinant
and $D^5=I$. In that case $Q=a'x_1x_4+b'x_2x_3$ with $a'b'\ne0$, and
\[
 b=d_{5,k,0}-d_{5,k-2,0}
   =\left\lfloor\frac{(k+1)^2}{5}\right\rfloor
           +\mathbf1_{k\equiv0\text{ or }3\ (5)},\qquad r=q-b.
 \tag{AKS52}
\]
If a noninvariant coefficient is nonzero, AKS11 applies instead. These
are the complete OQC alternatives. The bound AKS31 always uses that
actual $r$; it does not extend the $8k-16$ specialization to an
unevaluated exceptional point. This calculation does not assert that
the interior exceptional set is empty.

\section{Verification scope and the remaining absolute coefficient}
\label{aks:sec:AKS-scope}

The new mathematical content is AKS20--32 and its application to the
actual coefficient determinant AKS12--14, followed by the finite
receiving improvements AKS35--48. The complete source and orbit
identifications used in that application are the supplied OCS/OQC/OQX,
RPC, OPG/OPR and AMT results; their exact predecessors are preserved.
AKS33 pinpoints a repair obtainable immediately from the earlier proof,
whereas the Newton bound supplies the sharper explicit constant.

The bound \emph{at} scale $kq$ is now finite and uniform. What is not
proved is a matching lower coefficient, a unique limit for
$F_K^{(1)}/(kq)$, or a sign for the interval AKS48. Those questions involve
the orientation of the original period kernel within the inverse
remainder covariance; replacing that compression by a scalar bound
loses precisely that information. This is the finite-bounds outcome
allowed by the current request, not a claim to have evaluated an
asymptotic coefficient. The source differences are already completed
inputs and are not renamed as the remaining problem.

The accompanying exact checks distinguish finite algebraic diagnostics
from the analytic proof for the full family. They do not assert that a
diagnostic root is a zero of $\xi$, that a diagnostic matrix equals an
actual period matrix, or that a new Lean theorem has been checked.
No RH conclusion is drawn from this kernel bound.

% END PRESERVED INPUT provider_AKS.tex
\endgroup

\clearpage
\begingroup


\setcounter{equation}{0}
\renewcommand{\theequation}{PRT.\arabic{equation}}
\section*{Original rank-three primitive connecting maps and endpoint returns}
\addcontentsline{toc}{section}{Original rank-three primitive connecting maps and endpoint returns}
% BEGIN PRESERVED INPUT provider_PRT_COMPLETE_PREREQUISITE.tex SHA256 73a78a7c5dd1c6d4a04cc0c503a414cbc92a5a7ac5308676f1a90f805d559f1c
% BEGIN ADDITIVE MIXED PROFILE INTEGRATION PRT_CONNECTING_ENDPOINT_RETURN.tex
% EXACT INPUT SHA256 a0b5138b12764ffee5de29e9a579eb613abe544850e165a2f6624035878ff944
\clearpage
\label{prt:MPBUILD:PRT_CONNECTING_ENDPOINT_RETURN:start}
% Additive proof fragment. Parent source: accepted current14 SHA256
% 6b76e69e05a657c6d6675df2e5ad15345c0d0d5263da3bc0d40dc49b88319dc1.
\section{The original three-equation transfer of every connecting return}
\label{prt:sec:PRT}

Fix the simple quartet, its actual amplitude, period and logarithm branch
on the original five-orbit locus. Keep $k=4l+1\ge9$, $j=k+4$,
$q=(k+1)^2$, $Q=(j+1)^2$, $d=Q-q=8k+24$, and the full
nonreduced conductor
\[
 E_A(z)=\sum_{a,b=0}^8a_{ab}e^{\beta_{8,a,b}z},\qquad
 v=\operatorname{ord}_0E_A,\quad\mu_v=E_A^{(v)}(0)\ne0,
 \quad\alpha_A=v!/\mu_v.
 \tag{PRT1}
\]
No conductor factor is removed. Write $B=B_{\rm cof}$,
$\mathcal U=\mathbb C^{4m_j}$, $\mathbb J=\ker\mathbb B$,
$\mathscr B=\operatorname{im}B$ and $W=\mathscr B+\mathbb J$,
as in CPQ2. Define, in this same coefficient space,
\[
 S_e=\operatorname{im}[B,F_e],\qquad K_e=S_e\cap W,
 \qquad T_e=S_e+W\qquad(e=0,1).
 \tag{PRT2}
\]
These are the CPQ spaces, with $S_e$ abbreviated only to make the
two endpoint subscripts legible. The later norm cutoff $N$ is independent
of the defining endpoint $e$. At every $Q-1\le N\le2Q$ use exactly
\[
 D=N-v,\qquad \mathcal D_N^{(s)}=
 \operatorname{diag}(D_{j,D}^{(s)},4),\qquad
 c'_j=j/2-4,\quad h_{n,s}=(2\pi)^{s/2}n!(s/2)_n,
 \quad s=1\text{ or }j.
 \tag{PRT3}
\]
Thus both actual source pairs $(1,1)$ and $(k,j)$ are covered in
their target order. Their degree-$k$ proper source retains its separate
form; none of the identities below substitutes that form for $A_{\rm inv}$.

\subsection{The actual three rows in the original period coefficients}

Retain the original polynomials $L_i=(\mathsf H(ZW,Z,W,1)^{\mathsf T})_i$,
$b_{i,a,b}=[Z^aW^b]L_i^4$, the exponential $B_i(z)$ and the
translation maps $\mathcal T_i,\mathcal T_A$ of CPR2. In particular
\[
 R_i(S)=\frac{\sum_{a,b=0}^4b_{i,a,b}
                         \chi_j(S+\beta_{4,a,b})}{\chi_k(S)},
 \qquad b_i=B_i(0),\qquad
 \chi_n(S)=\prod_{a,b=0}^n(S-\beta_{n,a,b}).
 \tag{PRT4}
\]
The quotient is the complete polynomial quotient: root addition makes
every numerator vanish at every distinct root of $\chi_k$.
Choose the first original index $i_*$ with $b_{i_*}\ne0$, and put
\[
 C_i=R_i-(b_i/b_{i_*})R_{i_*}\quad(i\ne i_*),\qquad
 V=([S^a]C_i)_{0\le a<d,\ i\ne i_*}.
 \tag{PRT5}
\]
Existence of $i_*$ and independence of these three columns are the
complete PDS10--11/CPR17 calculation: the original Bezout identity at
zero gives a nonzero $b_i$; the four $R_i$ are independent; their leading
coefficients are $b_i$. Subtracting their exact leading multiples
therefore gives three independent columns of degree at most $d-1$.
No period coefficient is replaced in this operation.

Let $h_b$ be the original primitive germs used to define the $d$ columns
of $F_0$ in CIT7. They are specified modulo the entire raw analytic
source. Define their relation-value matrix by
\[
 (M_0)_{a,b}=\Lambda_{h_b}
                  \bigl(\mathcal T_A(\chi_k S^a)\bigr),
             \qquad0\le a<d,
 \qquad L=V^{\mathsf T}M_0.
 \tag{PRT6}
\]
Every entry uses the ordinary, complex-linear transpose. The complete
annihilator identity CPR9--10 proves that the displayed $d$ independent
relation polynomials are the full dual of the raw source quotient.
The $h_b$ are a basis of that quotient, so $M_0$ is invertible. Thus
$L$ has rank exactly three. The matrix is computable from the original
period arrays by finite polynomial division and pairing; it contains
neither a new numerical model nor a selected nonzero period minor.

Choose once a full column basis $Z$ of $\ker L$ and a three-column
complement $Y$ so that $[Z,Y]$ is invertible. Then the literal coefficient
spaces obey
\[
 \boxed{\quad S_1=\mathscr B+F_0\ker L\subset S_0,
       \qquad S_0/S_1\text{ has the three lifts }F_0Y.\quad}
 \tag{PRT7}
\]
Here equality is of subspaces of the original $\mathcal U$, not only
of abstract quotient dimensions. To prove it, use PDS11: truncation of
the endpoint-one primitive quotient identifies it with exactly
$\ker V^{\mathsf T}$ in endpoint-zero relation-value coordinates.
The missing top value is uniquely
$-u_*^{\mathsf T}x_0/b_{i_*}$ with the complete $u_*$ from PDS10.
For a representative $h$, set $f=E_Ah$. The original target numerator
in component $i$ is
\[
 N_i(f)=\mathbb V_j^{-1}
        \bigl(r![z^r](B_i f)\bigr)_{0\le r<Q},
 \qquad(\mathbb V_j)_{r,(a,b)}=\beta_{j,a,b}^{r}.
 \tag{PRT8}
\]
Those values use only the truncation through degree $Q-1$.
At endpoint one, the additional four equations assert exactly that
the same interpolants also have the required next jet; they do not
alter a coefficient in PRT8. Division by the complete $E_A$ gives
the same target analytic functions, modulo precisely the original lower
exponentials. An analytic source change has numerator $d$ and changes
the four target numerators by $(L_i^4d)_i$, hence changes their DNE
coordinates by $Bx$ with the full original scalar $\alpha_A$.
This is the coefficient square RRC5--9/CIT8. It proves both inclusions
in PRT7, including every analytic-source and lower-gauge change.
Finally $F_0$ induces an injection modulo $\mathscr B$ by CIT9;
the invertible matrix $[Z,Y]$ therefore proves the asserted dimension
and the stated three lifts.

The coefficient change can itself be written explicitly. For the original
endpoint-one primitive columns $h_{1,b}$ put
$ (M_1)_{a,b}=\Lambda_{h_{1,b}}(\mathcal T_A(\chi_kS^a))$
for $0\le a\le d$, and let $P_{<d}$ select the first $d$ rows.
Then
\[
 \Gamma=M_0^{-1}P_{<d}M_1,\qquad
 \operatorname{im}\Gamma=\ker L,\qquad
 F_1=F_0\Gamma+BH
 \tag{PRT7a}
\]
for the unique matrix $H$ with the displayed equality.
Indeed PDS11 proves the first two assertions, including that $\Gamma$
has rank $d-3$. The preceding numerator interpolation calculation
shows that each difference column lies in $\mathscr B$.
The original injectivity of $B$ gives its unique coordinate in $H$.
Equivalently $\Gamma=ZT$ for an invertible $(d-3)$-square matrix $T$,
so every original fixed frame change in PRT7 is specified by a literal
coefficient equation. No conjugation occurs in $M_0,M_1,L,\Gamma$.
Hermitian adjoints enter only when evaluating the actual metric below.

\subsection{An exact sequence of these three directions}

PRT7 gives $K_1\subset K_0$ and $T_1\subset T_0$. The identity on
coefficient representatives gives the exact sequence
\[
 0\longrightarrow K_0/K_1\xrightarrow{[x]\mapsto[x]}
 S_0/S_1\xrightarrow{[x]\mapsto[x]}T_0/T_1\longrightarrow0.
 \tag{PRT9}
\]
The first map is injective because $K_0\cap S_1=K_1$.
If $x\in S_0$ maps to zero on the right, write $x=s+w$ with
$s\in S_1$, $w\in W$; then $w=x-s\in K_0$, so its class is the
image of $K_0/K_1$. Conversely every such class maps to zero. Every
element of $T_0=S_0+W$ has the class of its $S_0$ summand modulo
$T_1=S_1+W$, proving surjectivity.
Let
\[
 a=\dim(T_0/T_1)=\dim(\mathcal G_k\cap\mathcal C_k)
                    \quad\text{from CPR17--18}.
 \quad
 \dim(K_0/K_1)=3-a,\qquad0\le a\le3.
 \tag{PRT10}
\]
The equality with CPR's actual polynomial intersection follows because
$\dim(T_e/W)=c_e$ and CPR18 proves $c_0-c_1=a$. It can also be
calculated directly in the three columns $F_0Y$: take their images
modulo $T_1$, retaining the full $\mathbb B$, $B$ and $L$ of
PRT2/PRT6. Their exact rank is $a$ and their kernel has dimension
$3-a$. This proves a complete original-data rank computation without
assuming a generic minor is nonzero. Rank strata, including $a=0$,
are retained.

Give the three terms of PRT9 their original attained metrics
\[
 A_N=\operatorname{Gram}_{\mathcal D_N}(S_0/S_1),\quad
 B_N=\operatorname{Gram}_{\mathcal D_N}(K_0/K_1),\quad
 C_N=\operatorname{Gram}_{\mathcal D_N}(T_0/T_1).
 \tag{PRT11}
\]
Every value frame is fixed in $N$ and in the source-order comparison.
For example if $X_1$ is any fixed independent frame of $S_1$, the first
one, in the exact three lifts $X_A=F_0Y$, is
\[
 A_N=X_A^*\mathcal D_NX_A-X_A^*\mathcal D_NX_1
 (X_1^*\mathcal D_NX_1)^{-1}X_1^*\mathcal D_NX_A.
 \tag{PRT12}
\]
Select the first independent images of $F_0Y$ modulo $T_1$ to obtain
the $a$ lifts $X_C$. A full basis of $\ker(F_0Y\bmod T_1)$, followed
by subtracting an $S_1$ representative as in PRT9, gives $3-a$
lifts $X_B$ in $K_0$. Applying PRT12 with $(X_1,X_A)$ replaced by
the frames of $(K_1,X_B)$ and $(T_1,X_C)$ gives $B_N$ and $C_N$.
All three are positive definite on their stated quotient values;
the determinant of an empty quotient is one.

For completeness every fixed frame factor can be retained before
taking a return. For a fixed inclusion $K\subset Y_0$, choose frames
$P_K,P_{Y_0}$ and quotient lifts $X$. Put
\[
 [P_K,X]=P_{Y_0}R_{Y_0/K},\qquad
 d_{Y_0/K}=|\det R_{Y_0/K}|^2,\qquad
 g_{Y_0}(N)=\det(P_{Y_0}^*\mathcal D_NP_{Y_0}).
 \tag{PRT13}
\]


% BEGIN HUMAN CITATION M05_PRT_QUOTIENT
This is the standard positive-block Schur calculation
\cite[\S2.1, (2.3)--(2.4)]{human:boyd1994}; the determinant
factorization follows from the same block elimination
\cite[p.~222, (6a)--(6b)]{human:hager1989}.
Here the block frame is exactly $[P_K,X]$, with its retained
factor $d_{Y_0/K}$; neither reference supplies the programme's
connecting sequence or its endpoint comparison.
% END HUMAN CITATION M05_PRT_QUOTIENT

The full block Schur determinant identity is
\[
 \det\operatorname{Gram}_{\mathcal D_N}(Y_0/K)
       =d_{Y_0/K}\,g_{Y_0}(N)/g_K(N).
 \tag{PRT14}
\]
It follows by expressing the restricted form in $[P_K,X]$,
taking its determinant, and eliminating its upper-left block.
This defines and evaluates every factor from the stated original
frames. It does not identify a varying minimizing representative
with a fixed lift. CPQ's frames for each quotient are used in
PRT13 when that quotient appears below.

\subsection{Evaluation of the difference of the actual CPQ returns}

Use the unchanged target return
$\mathcal Rf=-f(Q-1)-f(Q)+f(2Q-1)+f(2Q)$ and write
\[
 U_A=\mathcal R\log\det A_N,\qquad
 U_B=\mathcal R\log\det B_N,\qquad
 U_C=\mathcal R\log\det C_N.
 \tag{PRT15}
\]
These symbols are returns of precisely the three evaluated quotient
matrices PRT11--12. They are not freely assignable parameters.
The exact CPQ differences are
\[
 \begin{aligned}
 V_{{\rm conn},0}-V_{{\rm conn},1}&=U_A,&
 V_{h,0}-V_{h,1}&=U_B,\\
 V_{s,0}-V_{s,1}&=U_A-U_B,&
 V_{t,0}-V_{t,1}&=U_C,\\
 V_{{\rm rem},1}-V_{{\rm rem},0}&=U_C,&
 \mathcal L_0-\mathcal L_1&=U_C-U_A+U_B.
 \end{aligned}\tag{PRT16}
\]
Here $V_b,V_I,V_A$ are independent of the defining primitive endpoint
$e$: their spaces, coefficient maps and frames are the same CPQ3/9
objects. To prove every line, apply PRT14 to
$S_e/\mathscr B$, $K_e/\mathscr B$, $S_e/K_e$, $T_e/W$ and
$\mathcal U/T_e$, respectively. For example the logarithmic
determinant difference of the latter two endpoint quotients is
$\log g_{T_0}-\log g_{T_1}$ plus the explicitly defined fixed
frame-factor difference; PRT14 for $T_0/T_1$ identifies its
nonconstant term with $\log\det C_N$.
Each frame factor is independent of $N$, and
$\mathcal R1=-1-1+1+1=0$, so it cancels only after this full return.
The other four differences follow by the same stated determinant
identities. Finally substitute them into CPQ10's
$\mathcal L_e=V_I+V_{t,e}-V_{s,e}-V_A$.

This also gives the exact propagation, without deleting a complement,
\[
 \begin{aligned}
 Z_e&=V_b+V_{{\rm rem},e}+V_{s,e}+V_A,
 &Z_1-Z_0&=U_C-U_A+U_B,\\
 \mathcal P^{\rm rest}_0-\mathcal P^{\rm rest}_1
    &=U_C-U_A+U_B,
 &4S_{K,j}&=\mathcal L_0+Z_0=\mathcal L_1+Z_1.
 \end{aligned}\tag{PRT17}
\]
In particular a change of primitive endpoint changes the common loss
and complementary sum by exactly opposite amounts in their allocation
identity. The original total and the complete-action allocation
cancellation are preserved.

\subsection{Directed finite bounds with the original masses}

Use the proved original reconstruction/source constant of FRE1,
evaluated at degree $j$ and order $s$:
\[
 A_{j,s}^{\rm end}=
 \max\{1,Q\mathcal B_{2Q-v,s,j,A}^{2}\}
                    (M_{j,s}^{\rm rec})^2.
 \tag{PRT18}
\]
The two factors retain DNE13's full mass $(2\pi)^{s/2}$,
the full conductor scalar and division coefficients, and LJR16's
original fixed-disk reconstruction, phases and Vandermonde.
FRE2--4 proves on this exact metric
\[
 \mathcal D_{Q-1}^{(s)}\preceq\mathcal D_N^{(s)}
   \preceq A_{j,s}^{\rm end}\mathcal D_{Q-1}^{(s)},
 \qquad
 \log A_{j,s}^{\rm end}\le2Q\log j+O_{\rm actual}(Q).
 \tag{PRT19}
\]
Substitution in each full affine fibre defining PRT11 and subsequent
minimization preserve both inequalities. Original row extension also
proves that each of $A_N,B_N,C_N$ increases with $N$.
If $G_N$ is one of these forms of rank $r$, set
$f(N)=\log\det G_N-\log\det G_{Q-1}$.
The generalized eigenvalues in PRT19 show $0\le f(N)\le
r\log A_{j,s}^{\rm end}$ and $f$ is increasing. Hence
$\mathcal R\log\det G=f(2Q-1)+f(2Q)-f(Q)$ is nonnegative
and at most $2r\log A_{j,s}^{\rm end}$.
This proves the actual directed intervals
\[
 \begin{gathered}
 0\le U_A\le6\log A_{j,s}^{\rm end},\qquad
 0\le U_B\le2(3-a)\log A_{j,s}^{\rm end},\qquad
 0\le U_C\le2a\log A_{j,s}^{\rm end},\\
 \boxed{\quad0\le V_{{\rm rem},1}-V_{{\rm rem},0}
                   \le2a\log A_{j,s}^{\rm end}.\quad}\\
 -2(3-a)\log A_{j,s}^{\rm end}
 \le V_{s,0}-V_{s,1}\le6\log A_{j,s}^{\rm end},\\
 |\mathcal L_0-\mathcal L_1|\le6\log A_{j,s}^{\rm end},
 \qquad |Z_1-Z_0|\le6\log A_{j,s}^{\rm end}.
 \end{gathered}\tag{PRT20}
\]
The last bound uses $\dim B+\dim C=3$, keeping the exact
sequence rather than separately charging two unrelated rank-three
errors. These are finite bounds on actual CPQ expressions.
For each fixed original period they are
$O_{\rm actual}(Q\log j)=o(kq)$, for both source pairs.
At the actual rank stratum $a=0$, the remaining-target forms are
the same quotient in fixed value frames and their returns are exactly
equal. At $a=3$, the kernel spaces coincide and $U_B=0$ exactly.
All intermediate strata remain governed by their actual original
three-column matrix. No claim is made that either endpoint's common
leading return has now been evaluated; PRT16 evaluates their difference
and reduces its calculation to the stated three original directions.

The proof requires the five-orbit coefficient construction, its original
distinct quartet roots and $\mu_v\ne0$. It does not use the additional
pole-coordinate gate minors of RPI/RPK. Consequently a period excluded
by a pole-coordinate pivot remains covered when it belongs to the
original five-orbit locus: PRT6/8/12 construct all maps without that
pivot. Rank changes of the three-row image are handled by PRT10.
The value $v$ is retained throughout; neither $v\le32$ nor a generic
value of $v$ has been substituted for an actual coefficient.

\label{prt:MPBUILD:PRT_CONNECTING_ENDPOINT_RETURN:end}
% END ADDITIVE MIXED PROFILE INTEGRATION PRT_CONNECTING_ENDPOINT_RETURN.tex
% END PRESERVED INPUT provider_PRT_COMPLETE_PREREQUISITE.tex
\endgroup

\clearpage
\begingroup


\setcounter{equation}{0}
\renewcommand{\theequation}{BT.\arabic{equation}}
\section*{Complete boundary-tilt and Binet envelope proof}
\addcontentsline{toc}{section}{Complete boundary-tilt and Binet envelope proof}
% BEGIN PRESERVED INPUT provider_BOUNDARY_TILT_COMPLETE.tex SHA256 3c7785c28e7c713cddfacf405b261b62c8d1dfb5016764089a5b1aa7f8da49ca
\section{Actual boundary tilts and uniform convolution tails}\label{bt:actual-boundary-tilts-and-uniform-convolution-tails}

This calculation retains the original arithmetic density and the complete hypothetical actual quartet. It does not replace that density with a reference measure. The estimates do not assume RH or a pointwise lower bound for zeta.

\subsection{1. Objects and the exact retained expression}\label{bt:objects-and-the-exact-retained-expression}

Fix \(0<\delta<1/2\), \(\gamma>0\), and an integer \(m\ge1\), and suppose the selected full-order quartet is a quartet of zeros of the original \(g=2\xi\). Its monic polynomial is \[
h(s)=\prod_{\epsilon,\eta\in\{-1,1\}}
 (s-\tfrac12-\epsilon\delta-i\eta\gamma)^m,
\qquad \deg h=4m.
\tag{BT1}
\] The entire function \(g/h\), its original Taylor unit, source maps, and measure remain unchanged: \[
w_h(t)=\frac{|(g/h)(1/2+it)|^2}{2\pi},\qquad
\alpha=\frac\pi2,\qquad R_0=\sqrt{\delta^2+\gamma^2}.
\tag{BT2}
\] The original completion formula gives, exactly, \[
\begin{split}
w_h(t)
&=\frac{(t^2+1/4)^2}{\sqrt\pi}\,
 \frac{|\zeta(1/2+it)|^2}{|h(1/2+it)|^2}\,\sigma(t),\\
\sigma(t)
&=\frac{|\Gamma(1/4+it/2)|^2}{2\pi},\\
|h(1/2+it)|^2
&=\bigl[((t-\gamma)^2+\delta^2)
              ((t+\gamma)^2+\delta^2)\bigr]^{2m}.
\end{split}
\tag{BT3}
\] The two factors with each fixed imaginary sign contribute a minus sign before multiplication, so their product is the positive expression in the last line. Every denominator in (BT3) is positive. Conjugation of the original \(g\) and the real coefficients of \(h\) give \(w_h(-t)=w_h(t)\).

\subsection{2. A centered partial-sum bound from the original floor formula}\label{bt:a-centered-partial-sum-bound-from-the-original-floor-formula}

For every integer \(N\ge1\), initially for \(\Re s>1\), \[
\zeta(s)=\sum_{n=1}^{N}n^{-s}
 +\frac{N^{1-s}}{s-1}
 -s\int_N^\infty \{x\}x^{-s-1}\,dx.
\tag{BT4}
\] To verify (BT4), integrate the constant value \(\lfloor x\rfloor\) on each interval \([n,n+1)\), or telescope integration by parts of the absolutely convergent tail. The last integral is absolutely and locally uniformly convergent for \(\Re s>0\), so the identity continues to that half-plane with \(s\ne1\).

Let \(B_1(x)=\{x\}-1/2\), so \(|B_1(x)|\le1/2\). The identity \(\int_N^\infty x^{-s-1}dx=N^{-s}/s\) gives exactly \[
\zeta(s)=\sum_{n=1}^{N-1}n^{-s}+\frac12N^{-s}
 +\frac{N^{1-s}}{s-1}
 -s\int_N^\infty B_1(x)x^{-s-1}\,dx.
\tag{BT5}
\] No asymptotic remainder is substituted here.

On \(s=1/2+it\), put \(\varrho=\sqrt{t^2+1/4}\), which is both \(|s|\) and \(|s-1|\). For \(N\ge2\), \[
\sum_{n=1}^{N-1}n^{-1/2}+\tfrac12N^{-1/2}
\le 2\sqrt{N-1}-1+\tfrac1{2\sqrt N}
\le 2\sqrt N-1.
\] The first inequality follows from the integral bound for the decreasing function \(x^{-1/2}\). For the second, \(2(\sqrt N-\sqrt{N-1})=2/(\sqrt N+\sqrt{N-1})\ge1/\sqrt N\). At \(N=1\), the same asserted inequality holds directly, with left side \(1/2\). The absolute integral in (BT5), before its multiplier \(s\), is at most \(\tfrac12\int_N^\infty x^{-3/2}dx=N^{-1/2}\). Therefore \[
|\zeta(1/2+it)|
\le 2\sqrt N-1+\frac{\sqrt N}{\varrho}
                       +\frac{\varrho}{\sqrt N}
\quad(t\in\mathbb R,\ N\ge1).
\tag{BT6}
\] The infimum over these explicitly specified integer choices is also a bound.

For \(x=|t|\ge1\), choose \(N=\lfloor x\rfloor\). Then \[
x/2\le N\le x,\qquad
\varrho\le\frac{\sqrt5}{2}x,\qquad
\frac{\sqrt N}{\varrho}\le x^{-1/2}.
\] Since \(-1+x^{-1/2}\le0\), (BT6) yields \[
\boxed{|\zeta(1/2+it)|\le Z_*\sqrt{|t|}\quad(|t|\ge1),
\qquad Z_*=2+\sqrt{5/2}.}
\tag{BT7}
\] This improves the polynomial exponent in the uncentered estimate used in TG.15. It is not a claim that (BT7) is the strongest known zeta bound. On the whole real line the original \(N=1\) uncentered formula gives independently \[
|\zeta(1/2+it)|\le 1+2\sqrt{t^2+1/4}.
\tag{BT8}
\]

\subsection{3. Explicit two-sided Gamma bounds}\label{bt:explicit-two-sided-gamma-bounds}

The complete Binet calculation included below proves \[
\boxed{
c_\Gamma e^{-\alpha|t|}(1+|t|)^{-1/2}
\le \sigma(t)\le
C_\Gamma e^{-\alpha|t|}(1+|t|)^{-1/2}
\quad(t\in\mathbb R),}
\tag{BT9}
\] where \[
c_\Gamma=\sqrt2\,e^{-7/6},\qquad
C_\Gamma=\sqrt{2\sqrt5}\,e^{2/3}.
\tag{BT10}
\] For completeness, its exact starting identity for \(\Re z>0\) is \[
\log\Gamma z=(z-\tfrac12)\log z-z+\tfrac12\log(2\pi)+R(z),
\quad
R(z)=\int_0^\infty e^{-zu}
 \left(\frac1{e^u-1}-\frac1u+\frac12\right)\frac{du}{u}.
\] The logarithmic branch is the one real on the positive real axis. The kernel is \((\coth x-1/x)/(4x)\), with \(x=u/2>0\). Define \[
F(x)=x\cosh x-\sinh x,\qquad
G(x)=(x^2+3)\sinh x-3x\cosh x.
\] Both vanish at zero, and \(F'(x)=x\sinh x\ge0\), \(G'(x)=xF(x)\ge0\). Dividing their nonnegative values by the appropriate positive factors gives \(0\le\coth x-1/x\le x/3\). Thus the kernel lies between zero and \(1/12\). At \(z=1/4+it/2\), integration of \(e^{-u/4}/12\) proves \(|R(z)|\le1/3\).

Twice the real part of Binet's identity, with \(x=|t|\), gives exactly \[
\sigma(t)=|z|^{-1/2}
 \exp[-x\arctan(2x)-1/2+2\Re R(z)].
\] For \(x>0\), \[
-x\arctan(2x)-\tfrac12
=-\alpha x+x\arctan(1/(2x))-\tfrac12
\le-\alpha x,
\] because \(\arctan v\le v\) for \(v\ge0\). At \(x=0\), the same inequality is \(-1/2\le0\). The lower bound uses \(\arctan(2x)\le\pi/2\). Hence the full exponent is between \(-\alpha x-7/6\) and \(-\alpha x+2/3\), including zero.

Finally \(|z|^2=1/16+x^2/4\), and the exact squared differences \[
\frac{(1+x)^2}{4}-|z|^2=\frac{3+8x}{16},\qquad
|z|^2-\frac{(1+x)^2}{20}=\frac{(4x-1)^2}{80}
\] give \[
\sqrt2(1+x)^{-1/2}\le |z|^{-1/2}
\le\sqrt{2\sqrt5}(1+x)^{-1/2}.
\] These prove (BT9)--(BT10) with no unproved Stirling remainder and no unquantified compact Gamma minimum.

\subsection{4. A global explicit actual boundary-tilt bound}\label{bt:a-global-explicit-actual-boundary-tilt-bound}

Set \[
p=8m-\frac92,\qquad \beta=p-\frac12=8m-5,\qquad
T_h=\max(1,2R_0).
\tag{BT11}
\] Then \(p\ge7/2>1\) and \(\beta>0\). Define \[
\begin{split}
C_{\mathrm{comp}}
={}&\frac{C_\Gamma}{\sqrt\pi}
 (T_h^2+1/4)^2
 [1+2\sqrt{T_h^2+1/4}]^2
 \delta^{-8m}(1+T_h)^\beta,\\
C_{\mathrm{tail}}
={}&\frac{25C_\Gamma Z_*^2}{16\sqrt\pi}
 \left(\frac43\right)^{4m}
 \left(1+\frac1{T_h}\right)^\beta,\\
C_h^{\mathrm{tilt}}
={}&\max(C_{\mathrm{comp}},C_{\mathrm{tail}}).
\end{split}
\tag{BT12}
\] These constants are explicit in the fixed quartet and elementary constants. None is an unevaluated zeta supremum, and none changes the source mass.

For \(|t|\le T_h\), (BT3) has \(|h(1/2+it)|^2\ge\delta^{8m}\). Apply (BT8) and the upper inequality of (BT9). Multiplying by \(e^{\alpha|t|}(1+|t|)^p\) leaves the factor \((1+|t|)^\beta\); bounding \(|t|\) by \(T_h\) proves the constant \(C_{\mathrm{comp}}\).

For \(x=|t|\ge T_h\), retain the exact identity \[
((t-\gamma)^2+\delta^2)((t+\gamma)^2+\delta^2)
=(t^2-R_0^2)^2+4\delta^2t^2
\ge\frac9{16}t^4.
\tag{BT13}
\] The last inequality uses \(t^2\ge4R_0^2\). Thus \[
|h(1/2+it)|^2\ge(3/4)^{4m}x^{8m},
\qquad
(t^2+1/4)^2\le(25/16)x^4.
\] Applying (BT7) and (BT9) to (BT3) gives \[
e^{\alpha x}w_h(t)
\le\frac{25C_\Gamma Z_*^2}{16\sqrt\pi}
\left(\frac43\right)^{4m}
x^{5-8m}(1+x)^{-1/2}.
\] Multiplication by \((1+x)^p\) leaves \(((1+x)/x)^\beta\le(1+1/T_h)^\beta\). This proves the global bound \[
\boxed{
0\le e^{\alpha|t|}w_h(t)
\le C_h^{\mathrm{tilt}}(1+|t|)^{-p}
\quad(t\in\mathbb R).}
\tag{BT14}
\] The exact expression (BT3), including its actual zeta numerator and complete quartet denominator, is the object bounded throughout.

\subsection{5. Boundary moments and finite derivative orders}\label{bt:boundary-moments-and-finite-derivative-orders}

Define the original two boundary tilts and their actual masses: \[
f_\pm(t)=e^{\pm\alpha t}w_h(t),\qquad
M_\pm=\int_{\mathbb R}f_\pm(t)\,dt.
\tag{BT15}
\] The nonzero entire function \(g/h\) has only isolated zeros on the integration line, so both masses are positive. Evenness gives \(M_+=M_-=:M_\alpha\). Integration of (BT14) proves \[
\boxed{
0<M_\alpha=M_h(\alpha)=M_h(-\alpha)
\le\frac{2C_h^{\mathrm{tilt}}}{p-1}<\infty.}
\tag{BT16}
\] For every real \(j\ge0\) with \(j<p-1\), the inequality \(|t|^j\le(1+|t|)^j\) gives \[
\int_{\mathbb R}|t|^j e^{\pm\alpha t}w_h(t)\,dt
\le\frac{2C_h^{\mathrm{tilt}}}{p-j-1}.
\tag{BT17}
\] More precisely, for \(R\ge0\), the part \(|t|\ge R\) is at most \[
\frac{2C_h^{\mathrm{tilt}}}{p-j-1}
(1+R)^{j-p+1}.
\tag{BT18}
\] Thus every integer order \(0\le j\le8m-6\) is absolutely integrable at both boundary tilts. For a simple quartet, this includes orders zero, one and two.

For the actual Laplace function \(M_h(z)=\int e^{zt}w_h(t)\,dt\), differentiation in \(|\Re z|<\alpha\) gives \(\int t^j e^{zt}w_h(t)\,dt\). For each stated order, (BT14) dominates the absolute integrand uniformly on \(|\Re z|\le\alpha\) by the integrable function \(C_h^{\mathrm{tilt}}|t|^j(1+|t|)^{-p}\). Dominated convergence therefore extends those derivatives continuously to the closed strip with exactly the boundary moments above. This does not assert holomorphic continuation through the boundary or finiteness of every endpoint derivative.

\subsection{6. A uniform convolution envelope with the literal tilted mass}\label{bt:a-uniform-convolution-envelope-with-the-literal-tilted-mass}

If a nonnegative integrable function \(f\) has actual mass \(M>0\) and \(f(x)\le C(1+|x|)^{-p}\), then for every integer \(k\ge1\), \[
\boxed{
f^{*k}(u)\le kCM^{k-1}(1+|u|/k)^{-p}.}
\tag{BT19}
\] For \(k=1\) this is the hypothesis. In the convolution integral with \(x_1+\cdots+x_k=u\), at least one index has \(|x_i|\ge|u|/k\). Bound its integrand by the sum of the \(k\) restrictions to these events. On the event for index \(i\), the factor \(f(x_i)\) is at most \(C(1+|u|/k)^{-p}\). Parametrize the convolution hyperplane by the other \(k-1\) coordinates; the change from any such convolution parametrization to another has absolute determinant one. Dropping the event now leaves the full product integral \(M^{k-1}\). Summing proves (BT19), without a probability normalization.

Multiplication inside the original convolution gives exactly \[
f_\pm^{*k}(u)=e^{\pm\alpha u}m_{h,k}(u),
\qquad m_{h,k}=w_h^{*k}.
\tag{BT20}
\] Use the plus tilt for \(u\ge0\), the minus tilt for \(u<0\), and their same actual mass \(M_\alpha\). The result is \[
\boxed{
m_{h,k}(u)\le
kC_h^{\mathrm{tilt}}M_\alpha^{k-1}
e^{-\alpha|u|}(1+|u|/k)^{-p}}
\quad(k\ge1,\ u\in\mathbb R).
\tag{BT21}
\] All dependence on \(k\) is displayed. The constants \(C_h^{\mathrm{tilt}},p,M_\alpha\) belong to fixed \(h\).

For every integer \(n\ge0\) with \(2n>p-1\), use \((1+u/k)^{-p}\le k^pu^{-p}\) for \(u>0\). The Gamma integral gives \[
\boxed{
\int_{\mathbb R}u^{2n}m_{h,k}(u)\,du
\le 2C_h^{\mathrm{tilt}}k^{p+1}M_\alpha^{k-1}
\frac{\Gamma(2n-p+1)}{\alpha^{\,2n-p+1}}.}
\tag{BT22}
\] The condition \(2n>p-1\) is precisely what makes this comparison integrable at zero. The actual raw moments exist also when the condition fails, but (BT22) is not used there.

For the retained source comparison, one direct corollary is \[
m_{h,k}(u)\le b_k\sigma(u),\qquad
b_k=\frac{C_h^{\mathrm{tilt}}}{c_\Gamma}
k^{p+1}M_\alpha^{k-1}>0.
\tag{BT23}
\] Indeed \(1+|u|/k\ge(1+|u|)/k\); (BT21) and the lower Gamma bound leave the additional factor \((1+|u|)^{-p+1/2}\le1\). This is a proved identity-map norm comparison; it does not assert equality of the measures.

\section{Internal derivation of the exact Binet identity}\label{bt:internal-derivation-of-the-exact-binet-identity}

This note proves the exact identity used in the actual boundary-tilt calculation from the defining Gamma integral, elementary integration, locally uniform products, and the Beta substitution. No Stirling formula, Gamma asymptotic, or cited Binet identity is assumed.

All complex powers of positive real integration variables use their real logarithms. The notation \(\operatorname{Log}z\) on \(\Re z>0\) denotes the principal holomorphic logarithm. The holomorphic logarithm of Gamma will be constructed below and will be real on the positive real axis.

\subsection{1. Defining integrals and the finite Euler formula}\label{bt:defining-integrals-and-the-finite-euler-formula}

For \(\Re z>0\), define \[
\Gamma(z)=\int_0^\infty x^{z-1}e^{-x}\,dx.
\tag{BI1}
\] This is absolutely convergent. On a compact set with \(0<a\le\Re z\le b\), its integrand and each of its first finitely many derivatives in \(z\) are bounded in absolute value by \[
(x^{a-1}+x^{b-1})(1+|\log x|^j)e^{-x},
\] with the appropriate fixed integer \(j\). This function is integrable at zero and infinity. Differentiation under the integral is therefore valid, and \(\Gamma\) is holomorphic on the right half-plane. Integration by parts gives \[
\Gamma(z+1)=z\Gamma(z),
\qquad \Gamma(1)=1,
\tag{BI2}
\] because \(x^ze^{-x}\) tends to zero at both endpoints for \(\Re z>0\).

For \(\Re z,\Re w>0\), let \[
B(z,w)=\int_0^1 t^{z-1}(1-t)^{w-1}\,dt.
\tag{BI3}
\] For an integer \(N\ge1\), integration by parts gives \[
B(z,N+1)=\frac{N}{z}B(z+1,N).
\] Iterating, and using \(B(z+N,1)=1/(z+N)\), proves the exact finite formula \[
B(z,N+1)=\frac{N!}{z(z+1)\cdots(z+N)}.
\tag{BI4}
\] In particular none of these denominators vanishes in the right half-plane.

Set \[
E_N(z)=N^zB(z,N+1)
=\frac{N^zN!}{z(z+1)\cdots(z+N)},\qquad N\ge1.
\tag{BI5}
\] The substitution \(t=x/N\) gives \[
E_N(z)=\int_0^N x^{z-1}(1-x/N)^N\,dx.
\tag{BI6}
\] For \(0\le x<N\), \(\log(1-x/N)\le-x/N\), so \[
0\le(1-x/N)^N\le e^{-x}.
\] For each fixed \(x>0\), the expression on the left tends to \(e^{-x}\), with the indicator of \(x<N\) understood.

On a compact set as above, therefore \[
\begin{split}
\sup_z|E_N(z)-\Gamma(z)|
\le\int_0^\infty &(x^{a-1}+x^{b-1})\\
&\cdot\left|(1-x/N)^N\mathbf1_{\{x<N\}}-e^{-x}\right|dx
\longrightarrow0
\end{split}
\] by dominated convergence. The same argument with an additional factor \(|\log x|^j\), for \(j=1,2\), proves locally uniform convergence of those two derivatives as well. Thus the Euler limit and the derivative limits are established directly from their integrals.

\subsection{2. Nonvanishing and the holomorphic Gamma logarithm}\label{bt:nonvanishing-and-the-holomorphic-gamma-logarithm}

Write \(H_N=\sum_{j=1}^N1/j\). The real numbers \[
\gamma_N=H_N-\log N
\] decrease and are bounded below by zero. Indeed \[
\gamma_{N+1}-\gamma_N
=\frac1{N+1}-\log(1+1/N)<0
\] by integrating \(1/x>1/(N+1)\) on \([N,N+1)\), while \(H_N\ge\int_1^{N+1}dx/x=\log(N+1)\). Let \(\gamma_{\mathrm E}=\lim_N\gamma_N\).

For \(\Re z>0\), every \(1+z/j\) has positive real part. Its principal logarithm is holomorphic. Define \[
Q_N(z)=\operatorname{Log}z+z(H_N-\log N)
 +\sum_{j=1}^N\left(\operatorname{Log}(1+z/j)-z/j\right).
\tag{BI7}
\] Exponentiating this literal sum and using (BI5) gives exactly \[
e^{Q_N(z)}
=zN^{-z}\prod_{j=1}^N(1+z/j)=\frac1{E_N(z)}.
\tag{BI8}
\]

The logarithmic series in (BI7) converges locally uniformly. To prove this explicitly, for \(|w|\le1/2\), \[
\operatorname{Log}(1+w)-w
=-\int_0^1\frac{t w^2}{1+tw}\,dt,
\] whence \(|\operatorname{Log}(1+w)-w|\le|w|^2\). On a compact set with \(|z|\le R\), this bounds the \(j\)-th summand by \(R^2/j^2\) whenever \(j\ge2R\). The finite initial part is holomorphic and the summable bound proves uniform convergence of the tail. Hence \[
Q(z)=\operatorname{Log}z+\gamma_{\mathrm E}z
 +\sum_{j=1}^\infty
       \left(\operatorname{Log}(1+z/j)-z/j\right)
\tag{BI9}
\] is holomorphic and \(Q_N\to Q\) locally uniformly.

Taking limits in \(E_N e^{Q_N}=1\), using Section 1, proves \[
\Gamma(z)e^{Q(z)}=1.
\tag{BI10}
\] This proves that Gamma is nowhere zero on the entire right half-plane. It also constructs the required holomorphic logarithm: \[
\mathcal G(z):=-Q(z),\qquad e^{\mathcal G(z)}=\Gamma(z).
\tag{BI11}
\] For \(x>0\), every term in \(Q(x)\) is real and the defining Gamma integral is positive. Thus \(\mathcal G(x)\) is precisely the real logarithm of \(\Gamma(x)\); there is no unspecified additive \(2\pi i\) constant.

The second derivatives of the series in (BI9) converge locally uniformly, since \(\sum_j|z+j|^{-2}\) does. Termwise differentiation gives \[
\boxed{
\mathcal G''(z)=\sum_{j=0}^{\infty}\frac1{(z+j)^2}
\qquad(\Re z>0).}
\tag{BI12}
\] For justification without assuming a derivative-series theorem, on any slightly larger compact subset of the same half-plane the series converges uniformly; Cauchy's integral formula on small circles then gives uniform convergence of its derivatives. Direct estimates of the derivatives of its summands give the same conclusion.

\subsection{\texorpdfstring{3. Beta product, duplication, and the factor \(\sqrt\pi\)}{3. Beta product, duplication, and the factor \textbackslash sqrt\textbackslash pi}}\label{bt:beta-product-duplication-and-the-factor-sqrtpi}

The absolutely convergent product of the two defining Gamma integrals admits the change of variables \[
x=rt,\qquad y=r(1-t),\qquad r>0,\quad0<t<1.
\] Its Jacobian matrix has determinant \(-r\), so the positive integration Jacobian is \(r\). Consequently \[
\Gamma(z)\Gamma(w)=\Gamma(z+w)B(z,w)
\qquad(\Re z,\Re w>0).
\tag{BI13}
\] All powers here are powers of positive real bases. Fubini and the change of variables are justified by the same integrals with real parts of the exponents.

In \(B(z,z)\), first substitute \(t=(1+u)/2\), and then use symmetry and \(v=u^2\). Every factor is retained: \[
\begin{split}
B(z,z)
&=2^{1-2z}\int_{-1}^1(1-u^2)^{z-1}\,du\\
&=2^{2-2z}\int_0^1(1-u^2)^{z-1}\,du\\
&=2^{1-2z}B(1/2,z).
\end{split}
\tag{BI14}
\]

Finally \(x=t^2\) in (BI1) gives \(\Gamma(1/2)=2\int_0^\infty e^{-t^2}dt\). The latter is the positive integral over the whole real line. Its square is \(\int_{\mathbb R^2}e^{-(x^2+y^2)}dx\,dy\). Polar coordinates, whose positive Jacobian is \(r\), evaluate this square as \[
2\pi\int_0^\infty re^{-r^2}\,dr=\pi.
\] Thus \(\Gamma(1/2)=\sqrt\pi\). Inserting (BI13) into (BI14), and using the nonvanishing established in (BI10), proves \[
\boxed{
\Gamma(z)\Gamma(z+1/2)
=2^{1-2z}\sqrt\pi\,\Gamma(2z)
\qquad(\Re z>0).}
\tag{BI15}
\] On the positive real axis all these factors are positive, so their real logarithms obey the corresponding additive identity without any branch ambiguity.

\subsection{4. The bounded integral remainder}\label{bt:the-bounded-integral-remainder}

For \(u>0\), define \[
K(u)=\left(\frac1{e^u-1}-\frac1u+\frac12\right)\frac1u.
\tag{BI16}
\] We prove its exact bounds rather than assuming a Bernoulli expansion. With \(x=u/2>0\), \[
K(u)=\frac{\coth x-1/x}{4x}.
\] The two functions \[
F_0(x)=x\cosh x-\sinh x,\qquad
F_1(x)=(x^2+3)\sinh x-3x\cosh x
\] vanish at zero and satisfy \(F_0'(x)=x\sinh x\ge0\) and \(F_1'(x)=xF_0(x)\ge0\). Dividing their nonnegative values by positive factors gives \[
0\le\coth x-\frac1x\le\frac{x}{3}.
\] Therefore \[
\boxed{0\le K(u)\le\frac1{12}\quad(u>0).}
\tag{BI17}
\]

Define the actual remainder candidate \[
R(z)=\int_0^\infty e^{-zu}K(u)\,du,\qquad \Re z>0.
\tag{BI18}
\] It is absolutely convergent and holomorphic. Differentiating \(j\) times is justified on every compact set with \(\Re z\ge a>0\) by \(u^je^{-au}/12\). In particular \[
|R(z)|\le\frac1{12\Re z},
\qquad 0\le R(x)\le\frac1{12x}\quad(x>0).
\tag{BI19}
\]

\subsection{5. Equality of second derivatives}\label{bt:equality-of-second-derivatives}

Put \(c=\tfrac12\log(2\pi)\), and define \[
\mathcal F(z)=(z-\tfrac12)\operatorname{Log}z-z+c+R(z).
\tag{BI20}
\] Differentiation and the exact algebraic identity \[
u^2K(u)=\frac{u}{e^u-1}-1+\frac u2
\] give \[
\begin{split}
\mathcal F''(z)
&=\frac1z+\frac1{2z^2}
 +\int_0^\infty e^{-zu}
       \left(\frac{u}{e^u-1}-1+\frac u2\right)du\\
&=\frac1{z^2}
 +\int_0^\infty\frac{u e^{-zu}}{e^u-1}\,du.
\end{split}
\tag{BI21}
\] Here \(\int_0^\infty e^{-zu}du=1/z\) and \(\int_0^\infty ue^{-zu}du=1/z^2\) follow by elementary integration on \(\Re z>0\); all cancellation factors are displayed.

The geometric series \(1/(e^u-1)=\sum_{j=1}^\infty e^{-ju}\) is valid for \(u>0\). Its absolute integral after multiplication by \(ue^{-zu}\) is bounded by the same expression with \(\Re z\), which is integrable at zero and infinity. Hence termwise integration is valid, and \[
\mathcal F''(z)
=\frac1{z^2}+\sum_{j=1}^\infty
 \int_0^\infty ue^{-(z+j)u}du
=\sum_{j=0}^\infty\frac1{(z+j)^2}
=\mathcal G''(z).
\tag{BI22}
\] The right half-plane is connected, so the holomorphic function \(\mathcal F-\mathcal G\) has zero second derivative and consequently equals \(az+b\) there, for constants \(a,b\in\mathbb C\).

\subsection{6. The recurrence fixes the linear constant}\label{bt:the-recurrence-fixes-the-linear-constant}

For positive real \(x\), (BI2) and the real-axis branch in (BI11) give \[
\mathcal G(x+1)-\mathcal G(x)=\log x.
\] The corresponding exact difference of the explicit function (BI20) is \[
\mathcal F(x+1)-\mathcal F(x)-\log x
=(x+\tfrac12)\log(1+1/x)-1+R(x+1)-R(x).
\tag{BI23}
\] The left side equals \(a\). The elementary integral bounds \[
\frac{y}{1+y}\le\log(1+y)\le y\qquad(y>0)
\] show that \(x\log(1+1/x)\to1\) and \(\tfrac12\log(1+1/x)\to0\). The remainder bound (BI19) makes the remaining difference tend to zero. Thus \(a=0\), and \(\mathcal F-\mathcal G=b\) throughout the right half-plane.

\subsection{7. Duplication fixes the additive constant}\label{bt:duplication-fixes-the-additive-constant}

For \(x>0\), the positive-real logarithm of (BI15) is \[
\mathcal G(x)+\mathcal G(x+1/2)-\mathcal G(2x)
=(1-2x)\log2+\tfrac12\log\pi.
\tag{BI24}
\] Since \(\mathcal F=\mathcal G+b\), the difference between the same expression using \(\mathcal F\) and the right side of (BI24) equals exactly \(b+b-b=b\).

Insert (BI20) without dropping any constant. The terms containing \(\log x\) cancel, and the resulting difference is \[
\begin{split}
b={}&x\log(1+1/(2x))-\tfrac12\\
&+\left(c-\tfrac12\log(2\pi)\right)
 +R(x)+R(x+1/2)-R(2x).
\end{split}
\tag{BI25}
\] Our displayed \(c=\tfrac12\log(2\pi)\) cancels exactly. The same elementary logarithm bounds show \(x\log(1+1/(2x))\to1/2\), and each remainder tends to zero by (BI19). Thus \(b=0\).

Combining Sections 5--7 proves the complete exact identity \[
\boxed{
\log\Gamma(z)
=(z-\tfrac12)\operatorname{Log}z-z+\tfrac12\log(2\pi)
+\int_0^\infty e^{-zu}
\left(\frac1{e^u-1}-\frac1u+\frac12\right)\frac{du}{u},
\quad\Re z>0.}
\tag{BI26}
\] The logarithm on the left is the uniquely specified holomorphic branch from (BI11). The integral is absolutely convergent, and its bound and all differentiations have been justified. This is the identity used by the actual boundary-tilt proof.

% END PRESERVED INPUT provider_BOUNDARY_TILT_COMPLETE.tex
\endgroup

\clearpage
\begingroup


\setcounter{equation}{0}
\renewcommand{\theequation}{TW.\arabic{equation}}
\section*{Complete original tilted-convolution envelope proof}
\addcontentsline{toc}{section}{Complete original tilted-convolution envelope proof}
% BEGIN PRESERVED INPUT provider_THREE_FACTOR_ENVELOPE_COMPLETE.tex SHA256 3c600f651828fc2dcb637bd1f4d177d54a7f9e226da9fb975b93e448386f82e0
\section{Uniform original source comparison and the two growing norm windows}

Retain an actual complete off-line quartet
$1/2\pm\delta\pm i\gamma$ of common order $m$, with
$0<\delta<1/2$ and $\gamma>0$.  The polynomial $h$, the entire quotient
$v_h=(2\xi)/h$, its full Taylor unit, and the original source are unchanged.
Set
\[
 w_h(u)=\frac{|v_h(1/2+iu)|^2}{2\pi},\qquad
 m_{h,k}=w_h^{*k},\qquad \mu_h=\int_{\mathbb R}w_h(u)\,du,
 \qquad \alpha=\pi/2.
 \tag{TW1}
\]
The proved actual boundary-tilt calculation supplies
\[
 p=8m-\tfrac92>1,\qquad
 e^{\alpha|u|}w_h(u)\leq C_h(1+|u|)^{-p},\qquad
 M_\alpha=\int_{\mathbb R}e^{\alpha u}w_h(u)\,du\in(0,\infty).
 \tag{TW2}
\]
Here $C_h$ denotes exactly the explicit constant $C_h^{\mathrm{tilt}}$
of (BT12), with its original compact and tail expressions.
Its proof retains the exact factor
$-\pi^{-1/4}(u^2+1/4)e^{-iu\log\pi/2}\zeta(1/2+iu)$
and the full quartet denominator.  In particular (TW2) does not use a
positive lower bound for $\zeta$ or an assumption about its zero locations.

The comparison measure used below is the fixed original one-factor
Gamma measure
\[
 d\sigma(u)=\frac{|\Gamma(1/4+iu/2)|^2}{2\pi}\,du,
 \qquad \int d\sigma=\sqrt{2\pi}.
 \tag{TW3}
\]
Its proved two-sided boundary bounds are
\[
 c_\Gamma e^{-\alpha|u|}(1+|u|)^{-1/2}
 \leq\frac{d\sigma}{du}
 \leq C_\Gamma e^{-\alpha|u|}(1+|u|)^{-1/2},
 \quad
 c_\Gamma=\sqrt2e^{-7/6},\quad
 C_\Gamma=\sqrt{2\sqrt5}\,e^{2/3}.
 \tag{TW4}
\]
This fixed measure has mass $\sqrt{2\pi}$ for every $k$.  It is distinct
from the $k$-fold Gamma convolution of mass $(2\pi)^{k/2}$ in (AT4).
Both comparisons, with their exact connecting constants, will be retained.

\subsection{A uniform convolution bound at the actual exponential boundary}

For every integer $k\geq1$ and every $u\in\mathbb R$,
\[
 \boxed{m_{h,k}(u)\leq
 kC_hM_\alpha^{k-1}(1+|u|/k)^{-p}e^{-\alpha|u|}.}
 \tag{TW5}
\]
To prove this, put $f_+(x)=e^{\alpha x}w_h(x)$.  Its mass is exactly
$M_\alpha$, and $f_+(x)\leq C_h(1+|x|)^{-p}$.  The original product
integral gives $f_+^{*k}(u)=e^{\alpha u}m_{h,k}(u)$.
On the hyperplane $x_1+\cdots+x_k=u$, at least one coordinate satisfies
$|x_i|\geq|u|/k$.  Split this integral into the union of those $k$ sets.
On its $i$th set the $i$th factor is at most
$C_h(1+|u|/k)^{-p}$.  Eliminate $x_i=u-\sum_{j\ne i}x_j$; the change
from the original convolution coordinates has absolute Jacobian one.
Integrate the other $k-1$ nonnegative factors over all their original
domains; the resulting bound is
$C_h(1+|u|/k)^{-p}M_\alpha^{k-1}$.
Summing proves the bound with $e^{-\alpha u}$.
Conjugation of the original real-entire $2\xi$ and of the complete
quartet polynomial proves $w_h(-x)=w_h(x)$.  Thus
$f_-(x)=e^{-\alpha x}w_h(x)$ has the same mass, and the same argument
with $f_-$ proves the bound with $e^{\alpha u}$.  Taking the smaller
of these two proved expressions gives exactly (TW5).  This argument
also covers $k=1$ directly.  Tonelli's theorem applies because all
integrands are nonnegative and the displayed masses are finite.

For integers $n$ with $2n>p-1$, this yields the literal raw moment bound
\[
 \int_{\mathbb R}u^{2n}m_{h,k}(u)\,du
 \leq 2C_h k^{p+1}M_\alpha^{k-1}
 \frac{\Gamma(2n-p+1)}{\alpha^{2n-p+1}}.
 \tag{TW6}
\]
Indeed $(1+u/k)^{-p}\leq k^pu^{-p}$ for $u>0$, and substitution
$y=\alpha u$ evaluates the remaining positive half-line integral.
No arithmetic density has been assigned the value of a Gamma density.

\subsection{Pointwise comparison of the original convolution}

For $k\geq3$, retain the proved original convolution lower envelope
\[
 m_{h,k}(u)\geq c_h\vartheta_h^{k-3}
 e^{-\alpha(|u|+k-3)}(k-2+|u|)^{-B},\quad
 B=42+8m,\quad \vartheta_h=\int_{-1}^{1}w_h(u)\,du>0.
 \tag{TW7}
\]
Its complete proof is the original arithmetic envelope (AU.10), also
retained as (CJ.15); it concerns this same $g/h$ convolution.
Define the following explicit constants, without altering either measure:
\[
 \begin{gathered}
 M=B/2=21+4m,\\
 a_k=\frac{c_h\vartheta_h^{k-3}e^{-\alpha(k-3)}(k-2)^{-B}}
 {C_\Gamma 2^{B/2}},\qquad
 A_k=\frac{C_h}{c_\Gamma}k^{p+1}M_\alpha^{k-1}.
 \end{gathered}\tag{TW8}
\]
Then
\[
 \boxed{
 a_k(1+u^2)^{-M}\frac{d\sigma}{du}
 \leq m_{h,k}(u)\leq A_k\frac{d\sigma}{du}.}
 \tag{TW9}
\]
For the upper bound, $1+|u|/k\geq(1+|u|)/k$ in (TW5), and the
lower inequality of (TW4) leaves the factor
$(1+|u|)^{-(p-1/2)}\leq1$.  This gives precisely $A_k$.
For the lower bound, $k-2\geq1$ gives
$k-2+|u|\leq(k-2)(1+|u|)$.  The upper inequality of (TW4) implies
$e^{-\alpha|u|}\geq C_\Gamma^{-1}d\sigma/du$.
Finally $(1+|u|)^B\leq2^{B/2}(1+u^2)^{B/2}$.
Substitution of these three inequalities in (TW7) proves its lower
counterpart in (TW9), with its displayed constant and direction.

\subsection{Full finite polynomial forms, with the original coordinate maps}

Let $b_j$ be the fixed Gamma monic polynomials in $u$.  The original
generating-integral calculation (TG.7)--(TG.14) proves
\[
 \gamma_j=\|b_j\|_\sigma^2
 =\sqrt2\,j!\Gamma(j+1/2)
 =\sqrt{2\pi}\frac{(2j)!}{4^j},\qquad
 u b_j=b_{j+1}+j(j-1/2)b_{j-1}.
 \tag{TW10}
\]
The last term is absent at $j=0$.  The second equality follows by
induction from $\Gamma(1/2)=\sqrt\pi$ and the Gamma recurrence; no
asymptotic formula is used there.

For $N\geq0$ put
\[
 D_N=[1+4(N+M)^2]^M.
\]
Every polynomial $P$ of degree at most $N$ satisfies
\[
 \int |P(u)|^2(1+u^2)^M\,d\sigma(u)
 \leq D_N\int |P(u)|^2\,d\sigma(u).
 \tag{TW11}
\]
Here is the full multiplication estimate.  The orthonormal vectors
$e_j=b_j/\sqrt{\gamma_j}$ obey
\[
 u e_j=\sqrt{(j+1)(j+1/2)}\,e_{j+1}
       +\sqrt{j(j-1/2)}\,e_{j-1}.
\]
On degrees at most $N$, the first weighted shift has norm at most
$N+1$ and the second has norm at most $N$; this follows by taking
squared norms of their orthogonal images term by term.
The triangle inequality therefore gives
$\|uP\|_\sigma\leq2(N+1)\|P\|_\sigma$.
After each multiplication the degree increases by one.  Iteration gives
$\|u^jP\|_\sigma\leq[2(N+M)]^j\|P\|_\sigma$ for $0\leq j\leq M$.
Expand $(1+u^2)^M$, integrate its nonnegative summands, and use the
binomial theorem; the result is exactly (TW11).

Apply Cauchy--Schwarz with the two functions
$|P|(1+u^2)^{-M/2}$ and $|P|(1+u^2)^{M/2}$ in $L^2(\sigma)$:
\[
 \left(\int|P|^2\,d\sigma\right)^2
 \leq\left(\int|P|^2(1+u^2)^{-M}\,d\sigma\right)
       \left(\int|P|^2(1+u^2)^M\,d\sigma\right).
\]
For $P\ne0$, divide by its positive $\sigma$ norm and apply (TW11).
For $P=0$ the resulting inequality is immediate.
Combining with (TW9) proves the full finite-form comparison
\[
 \boxed{
 \frac{a_k}{D_N}\int|P|^2\,d\sigma
 \leq\int|P|^2m_{h,k}(u)\,du
 \leq A_k\int|P|^2\,d\sigma
 \qquad(\deg P\leq N).}
 \tag{TW12}
\]
It holds for every complex coefficient vector, not just monic polynomials.
In particular $A_kD_N/a_k\geq1$ by applying (TW12) to any nonzero $P$.

The original sum coordinate is $S=c+iu$, $c=k/2$.  Its exact algebra
map is $\Psi_k:\mathbb C[S]\to\mathbb C[u]$,
$\Psi_k P(u)=P(c+iu)$, with inverse $f\mapsto f((S-c)/i)$.
In degree at most $N$ its coefficient matrix has diagonal
$1,i,\ldots,i^N$.  Thus (TW12) is a congruence statement in the
original $S$ coordinates as well.  On the monic degree-$n$ affine
space, $P\mapsto i^{-n}\Psi_kP$ is a bijection to the monic $u$
polynomials and preserves the norm because $|i^{-n}|=1$.
No degree-dependent phase is asserted to be a single algebra map.
Consequently, writing $H_{k,N}$ for the actual original polynomial
Gram and $H_{N}^{\sigma,k}$ for the fixed $\sigma$ Gram evaluated at
$S=c+iu$, the fully typed matrix form is
\[
 \boxed{(a_k/D_N)H_N^{\sigma,k}\preceq H_{k,N}
 \preceq A_kH_N^{\sigma,k}.}\tag{TW13}
\]
The identity on the original polynomial source commutes with remainder
modulo the same $\chi_{h,k}$ and hence with its full unitful jet
observation.  Neither this comparison nor the coordinate congruence
changes $\chi_{h,k}$, $S$, its full jets, or the original source norm.

\subsection{Uniform arithmetic monic norms}

Let $\omega_{h,k,n}$ be the minimum original squared norm of a monic
degree-$n$ polynomial.  Minimizing (TW12) on that same affine space,
using (TW10) for its fixed comparison minimum, gives
\[
 \boxed{\frac{a_k\gamma_n}{D_n}\leq\omega_{h,k,n}
 \leq A_k\gamma_n,\qquad k\geq3,\ n\geq0.}\tag{TW14}
\]
The upper bound may equivalently be verified on the explicit trial
polynomial $b_n$; the lower bound applies to every monic polynomial.
Thus no minimizer was transferred between different metrics without
the stated inequality.

Here and below constants in $O_h$ may depend on the fixed complete
quartet and the fixed proven envelope constants, but not on $k,n$.
Taking logarithms of the exact constants in (TW8) gives
\[
 \log a_k=(k-3)(\log\vartheta_h-\alpha)-B\log(k-2)
            +\log c_h-\log C_\Gamma-(B/2)\log2,
\]
\[
 \log A_k=(k-1)\log M_\alpha+(p+1)\log k
                       +\log(C_h/c_\Gamma).
\]
Also $\log D_n=M\log(1+4(n+M)^2)=O_h(1+\log(n+1))$.
These displayed equalities prove, uniformly for $k\geq3,n\geq0$,
\[
 \boxed{\log\omega_{h,k,n}=
 \log\!\left(\sqrt{2\pi}\frac{(2n)!}{4^n}\right)
 +O_h(k+\log(n+1)).}\tag{TW15}
\]
This is a high-degree norm estimate for the actual arithmetic convolution.
It is not an assertion that consecutive norm ratios have an asymptotic;
the error in (TW15) does not justify that differentiation in $n$.

The same all-polynomial comparison also controls the original raw even
moments $I_{k,2n}=\int_{\mathbb R}u^{2n}m_{h,k}(u)\,du$.  For
$k\geq3$ and $n\geq1$ it gives the explicit two-sided bounds
\[
 \boxed{
 \frac{a_k\sqrt2c_\Gamma e^{-\alpha}}{D_n}
       \frac{(2n-1)!}{\alpha^{2n}}
 \leq I_{k,2n}
 \leq2A_k C_\Gamma\frac{(2n)!}{\alpha^{2n+1}}.}
 \tag{TW15a}
\]
For the lower comparison moment, restrict (TW4) to $u\geq1$ and
use $(1+u)^{-1/2}\geq1/(\sqrt2u)$.  Both half-lines give
$\int u^{2n}\,d\sigma\geq\sqrt2c_\Gamma
\alpha^{-2n}\Gamma(2n,\alpha)$.
Repeated integration by parts gives
$\Gamma(2n,\alpha)=(2n-1)!e^{-\alpha}
\sum_{j=0}^{2n-1}\alpha^j/j!\geq(2n-1)!e^{-\alpha}$.
For the upper comparison moment, drop the factor
$(1+|u|)^{-1/2}\leq1$ in (TW4) and evaluate the Gamma integral.
Applying (TW12) to the original degree-$n$ polynomial $u^n$ proves
(TW15a).  Its logarithms give, uniformly on these stated degrees,
\[
 \log I_{k,2n}=\log\bigl((2n)!\alpha^{-2n}\bigr)
                    +O_h(k+\log(n+1)).\tag{TW15b}
\]
The full one-factor zeta numerator and actual tilted mass in (TW2)
remain in $a_k,A_k$; (TW15a) is a bound on the original moments.

\subsection{The two original windows with full finite errors}

Put $q=[1+k(m-1)](k+1)^2$.  For $n=q-1$ or $n=q$, division of
(TW14) at the two exact degrees gives
\[
 \log\frac{\gamma_{n+q}}{\gamma_n}
       -\log\frac{A_kD_{n+q}}{a_k}
 \leq\log\frac{\omega_{h,k,n+q}}{\omega_{h,k,n}}
 \leq\log\frac{\gamma_{n+q}}{\gamma_n}
       +\log\frac{A_kD_n}{a_k}.
 \tag{TW16}
\]
All logarithmic error bounds here are nonnegative by (TW12).
The exact factorial formula (TW10) and monotonicity of $\log x$ give
\[
 \log\frac{\gamma_{2q}}{\gamma_q}
 =2q\log(4q/e)+\eta_q,\qquad 0\leq\eta_q\leq\log2.
 \tag{TW17}
\]
Indeed the logarithm of $(4q)!/(2q)!$ is
$\sum_{j=2q+1}^{4q}\log j$.  Its excess over
$\int_{2q}^{4q}\log x\,dx$ is nonnegative and at most
$\sum_{j=2q+1}^{4q}(\log j-\log(j-1))=\log2$.
Subtracting $q\log4$ gives precisely (TW17).
The other original window satisfies the exact identity
\[
 \frac{\gamma_{2q-1}}{\gamma_{q-1}}
 =\frac{\gamma_{2q}}{\gamma_q}\frac{q-1/2}{4q-1},
 \qquad \frac16\leq\frac{q-1/2}{4q-1}<\frac14.
\]
It follows in particular that
\[
 4q\log(4q/e)-\log6
 \leq\log\frac{\gamma_{2q-1}\gamma_{2q}}
                    {\gamma_{q-1}\gamma_q}
 \leq4q\log(4q/e).
 \tag{TW18}
\]
For fixed $h$, one has $k/q\to0$ and $\log q/q\to0$.
Thus each finite error in (TW16), divided by $2q$, tends to zero.
Equations (TW16)--(TW18) prove both limits on their stated windows:
\[
 \boxed{
 \lim_{k\to\infty}\frac1q
 \left(\frac{\omega_{h,k,2q}}{\omega_{h,k,q}}\right)^{1/(2q)}
 =\lim_{k\to\infty}\frac1q
 \left(\frac{\omega_{h,k,2q-1}}{\omega_{h,k,q-1}}\right)^{1/(2q)}
 =\frac4e.}\tag{TW19}
\]
No limit with an unspecified ratio $n/q$ is asserted.

\subsection{The resulting actual four-window lower estimate}

Retain the original relative controls $\varepsilon_N$, phases $\phi_N$,
and contraction penalties $P_-,P_+$ of (HC1)--(HC2).  Their exact first
boundary and regular products, inherited from (CJ.9)--(CJ.13), give
\[
 \mathcal B_k=\sum_{N=q-1}^{2q-1}w_N
 \log(\varepsilon_N^2+\phi_N^2)
 -\log\frac{\omega_{h,k,2q-1}\omega_{h,k,2q}}
                {\omega_{h,k,q-1}\omega_{h,k,q}}+P_-+P_+,
\]
where the endpoint weights are $1,2,\ldots,2,1$.  Set
\[
 E_k=2\log(A_k/a_k)+\log D_{q-1}+\log D_q\geq0,
 \qquad E_k=O_h(k+\log q).
\]
Adding the upper sides of (TW16) and applying (TW18) gives the
actual sum-norm upper bound $4q\log(4q/e)+E_k$.
Consequently, with the full original trace
$\mathcal L=2\delta\ell_k(k+1)\lfloor(k+1)^2/4\rfloor$,
\[
 \begin{split}
 \mathcal B_k\geq{}&4q\log(\delta e k/8)-E_k+P_-+P_+\\
 &+\sum_Nw_N\log(1+\phi_N^2/\mathcal L^2)
 +\sum_Nw_N\log\frac{\varepsilon_N^2+\phi_N^2}
                           {\mathcal L^2+\phi_N^2}
 +4q\log\frac{2\mathcal L}{\delta kq}.
 \end{split}\tag{TW20}
\]
All terms after $-E_k$ are the displayed original nonnegative costs.
The sums retain both endpoint degrees and all their multiplicities.
To verify the constant, $4q\log\mathcal L-4q\log(4q/e)$ equals
$4q\log(\delta e k/8)+4q\log(2\mathcal L/(\delta kq))$.
This proves the identity used in the bound with no rescaling of $S$.
Let $\mathcal R_k^{\rm TW}$ denote exactly the right side of
(TW20), with every displayed contraction, phase, control and trace term.
The original Gamma four-volume $\mathcal B_k^\Gamma$ uses the
same $\chi$ and the original tensor Gamma density. Substituting
(AT19), whose $\mathfrak C_*$ is computed on that precise pair of
source Grams, gives the additional current consequence
\[
\mathcal R_k^{\rm TW}\le\mathcal B_k
 \le\mathfrak b_*^+,\qquad
\mathfrak b_*^+\ge\mathcal R_k^{\rm TW},\qquad
\mathfrak C_*\ge J_*\ge\max\{0,\mathcal R_k^{\rm TW}-\mathcal B_k^\Gamma\}.
\tag{TW20a}
\]
Here $\mathfrak b_*^+$ is (AT21e) on the actual tensor Gamma
path, incorporating the nonlinear radius $I_*$ as well as the
additive control. All inequalities follow by substitution of the same scalar;
none of the nonnegative terms of (TW20) is discarded or assigned
an asymptotic cancellation. Since $E_k/q\to0$, the quantified refinement includes
\[
 \liminf_{k\to\infty}
 \frac{\mathcal B_k-4q\log k}{q}
 \geq4\log(\delta e/8).
 \tag{TW21}
\]
The leading threshold four is inherited.  The new content is the
actual boundary-tilt comparison, its uniform finite-degree errors,
and the sharper literal window constant.  These lower estimates
alone do not assign a sign to the signed arithmetic correction in
(AT14)--(AT17).

% END PRESERVED INPUT provider_THREE_FACTOR_ENVELOPE_COMPLETE.tex
\endgroup

\clearpage
\begingroup


\setcounter{equation}{0}
\renewcommand{\theequation}{ZLM.\arabic{equation}}
\section*{Original local zeta mass and complete three-factor lower envelope}
\addcontentsline{toc}{section}{Original local zeta mass and complete three-factor lower envelope}
% BEGIN PRESERVED INPUT provider_endpoint_arithmetic_original_source.tex SHA256 b83d8a31bcde06321afc31b011b47d53ab5d4d083809c71beafaf421eccbbce2
\textbf{Source metadata: author (MetaInlines)}

Owner-directed SplitZero analytic continuation

\textbf{Source metadata: date (MetaInlines)}

13 September 2026

\textbf{Source metadata: subtitle (MetaInlines)}

A uniform polynomial-norm estimate; integration with the four-volume criterion

\textbf{Source metadata: title (MetaInlines)}

Arithmetic endpoint bounds on the original \(\tau\)-base source

\subsection{1. Result and exact source boundary}\label{zlm:endpoint-arithmetic_original--result-and-exact-source-boundary}

The formalization session's endpoint criterion selects an actual member of the canonical arithmetic family. This continuation proves one of its previously unproved analytic inputs. For the \textbf{unchanged arithmetic measure}, a fixed finite packet \(h\), and integers \(n\ge k\ge3\), there is a finite constant \(C_h\), independent of \(n\) and \(k\), such that

\[
 \boxed{\left(\frac{\omega_{h,k,2n}}{\omega_{h,k,n}}\right)^{1/(2n)}
       \le C_h n.}                                                    \tag{1}
\]

The norms are the original monic polynomial norms for \(m_{h,k}=w_h^{*k}\); the total mass remains \(\mu_h^k\). No root-location hypothesis is used in this bound. The proof retains a compact-source mass, a lower convolution constant, and both exponential moments. These constants are defined explicitly below; they have not been numerically enclosed.

For the packet consisting exactly of a hypothetical nonreal off-line quartet, let

\[
 q_k=[1+k(m-1)](k+1)^2,
 \quad
 L_{h,k}=2\delta[1+k(m-1)](k+1)
                 \left\lfloor\frac{(k+1)^2}{4}\right\rfloor.             \tag{2}
\]

The existing exterior theorem gives \(L_{h,k}\le\epsilon_{h,k,N}\) at every admitted degree. Combining (1) with the other session's endpoint theorem gives

\[
 \boxed{\min_{q_k\le N<2q_k}\epsilon_{h,k,N}
 \le C_hq_k\sinh\!\left(\frac{\mathcal B_{h,k}}{2q_k}\right),\qquad
 \mathcal B_{h,k}=\log\frac{V_{q_k-1}V_{q_k}}
                              {V_{2q_k-1}V_{2q_k}}.}                    \tag{3}
\]

Consequently such a quartet forces

\[
 \boxed{\mathcal B_{h,k}\ge
 2q_k\operatorname{arsinh}\!\left(\frac{\delta k}{2C_h}\right),\qquad
 \liminf_{k\to\infty}\frac{\mathcal B_{h,k}}{q_k\log k}\ge2.}           \tag{4}
\]

The upper bound for this four-volume quantity is not proved here. In particular, (1) does not establish RH. It completes the polynomial-norm input in (3) on a specific admissible window and leaves the source/relation-volume comparison with its exact maps. Section 9 expresses that remaining quantity as two positive block log-determinants of actual new relation layers.

The finite quotient, canonical source, residue action and rank-two control are inputs from the preceding analytic notes. The new proofs in Sections 3--7 use only their explicitly defined arithmetic measure. The endpoint theorem is the formalizer's written result, not attributed to this continuation. Its finite telescopes are reported kernel-checked; the analytic estimates below have written proofs, not new Lean certificates.

\subsection{2. The base, source, and both kinds of quotient}\label{zlm:endpoint-arithmetic_original--the-base-source-and-both-kinds-of-quotient}

Retain the original marked coefficient square

\[
\begin{array}{ccc}
 G(\mathbb Z)&\xrightarrow{G(j)}&G(\mathbb C)\\
 p_{\mathbb Z}\downarrow&&\downarrow p_{\mathbb C}\\
 \mathbb Z&\xrightarrow{j}&\mathbb C.
\end{array}                                                          \tag{5}
\]

Here \(G(R)=\{\tau\}\sqcup\{r^\bullet:r\in R\}\), \(e_R=0_R^\bullet\), \(p_R(\tau)=0_R\), and \(p_R(r^\bullet)=r\). The target of \(p_{\mathbb Z}\) is infinite. The construction remains over the original pointed base \(\mathfrak b_\tau\); no new scalar zero is adjoined at a tensor degree or at a polynomial relation layer.

Keep \(C_+=[V\xrightarrow{\Theta}\mathscr B]\) in degrees \(0,1\), \(Q=\mathscr B/\Theta V\), and \(D=-x\partial_x\). The source spaces and the Fourier transform remain those of the earlier theta construction. In particular,

\[
 g(s)=2\xi(s)=s(s-1)\pi^{-s/2}\Gamma(s/2)\zeta(s),\quad
 \mathcal M\Theta\phi_*=g.                                            \tag{6}
\]

Let \(h\) be a fixed finite product of zeros of \(g\), with the full order of each selected zero retained, and write \(d=\deg h\). Define

\[
 v_h=g/h,\quad
 w_h(t)=\frac{|v_h(1/2+it)|^2}{2\pi},\quad
 \mu_h=\int_{\mathbb R}w_h(t)\,dt,\quad m_{h,k}=w_h^{*k}.               \tag{7}
\]

The quotient \(v_h\) is entire, including at the selected centres. It is not the meromorphic expression evaluated before removing its actual removable factors. For \(h=1\), the same analytic definitions are meaningful, but the finite arithmetic packet is zero. Only the nonempty case is used for the finite Gram inverses below.

The source map is

\[
 \mathcal V_{h,k}:\mathbb C[S]\longrightarrow\mathscr B^{\widehat\otimes k},
 \qquad P\longmapsto P(D_1+\cdots+D_k)(F_h^{\otimes k}),\qquad
 \mathcal MF_h=v_h.                                                    \tag{8}
\]

The original cyclic module is \(E_{h,k}=\mathbb C[S]/(\chi_{h,k})\). Its arithmetic inclusion keeps the Taylor unit:

\[
 \eta_{h,k}[P]=\upsilon_h^{\otimes k}P(A_k)1,\quad
 \upsilon_h=j_h(g/h),\quad A_k=\sum_iM_{s_i}.                           \tag{9}
\]

The supplied source identities are

\[
 J^{(k)}\mathcal V_{h,k}=\eta_{h,k}\pi_\chi,\qquad
 q^{(k)}\mathcal V_{h,k}=\sigma_h^{\otimes k}\eta_{h,k}\pi_\chi.          \tag{10}
\]

The actual norm of (8) is

\[
 \|\mathcal V_{h,k}P\|^2
 =\int_{\mathbb R}|P(k/2+iu)|^2m_{h,k}(u)\,du.                          \tag{11}
\]

For \(\mathcal P_n=\mathbb C[S]_{\le n}\), the monic norm is the quotient norm of the leading-coefficient map

\[
 0\longrightarrow\mathcal P_{n-1}\longrightarrow\mathcal P_n
 \xrightarrow{\operatorname{lc}_n}\mathbb C\longrightarrow0,
 \qquad
 \omega_{h,k,n}=\inf_{\operatorname{lc}_nP=1}\|\mathcal V_{h,k}P\|^2.
                                                                         \tag{12}
\]

The arithmetic quotient is the other specified map from the same source:

\[
 0\longrightarrow\mathcal P_{n-q}
 \xrightarrow{\times\chi}\mathcal P_n
 \xrightarrow{\pi_\chi}E_{h,k}\longrightarrow0,\qquad q=\deg\chi,
                                                                         \tag{13}
\]

for \(n\ge q-1\), with the preceding relation space zero at \(n=q-1\). Multiplication by the monic \(\chi\) carries \(\mathcal P_{n-q}/\mathcal P_{n-q-1}\) isomorphically to \(\mathcal P_n/\mathcal P_{n-1}\), with leading coefficient unchanged. The map \(\pi_\chi\) kills that very relation. Thus (12), (13), and multiplication by \(\chi\) specify the relation between the degree-line norm and the arithmetic quotient metric. One norm is not substituted for the other.

At each admitted original support label \(\lambda\), the lifted quotient is \((\lambda,P)\mapsto(\lambda,[P]_\chi)\). An original relation has image \((\lambda,0)\); external absence has image \(\tau\). Its norm before this projection remains available through (11).

\subsection{3. A quantitative local arithmetic mass lemma}\label{zlm:endpoint-arithmetic_original--a-quantitative-local-arithmetic-mass-lemma}

\subsubsection{3.1 Uniform polynomial growth on a fixed strip}\label{zlm:endpoint-arithmetic_original--uniform-polynomial-growth-on-a-fixed-strip}

There is a constant \(C_0\ge1\) such that, for \(|T|\ge10\), the function

\[
 f_T(z)=\zeta(1/2+iT+iz)                                                \tag{14}
\]

is holomorphic on and inside the ellipse

\[
 z=\frac12(w+w^{-1}),\qquad |w|=8,
                                                                         \tag{15}
\]

and has modulus at most \(M_T=C_0(1+|T|)^6\) there.

The ellipse has real and imaginary semiaxes \(65/16\) and \(63/16\). Hence its image in the \(s\)-plane lies in \(-55/16\le\Re s\le71/16\) and \(|\Im s-T|\le65/16\). The pole of \(\zeta\) is outside this set when \(|T|\ge10\).

For \(\Re s\ge1/4\), use the identity

\[
 \zeta(s)=\frac{s}{s-1}-s\int_1^\infty\{x\}x^{-s-1}\,dx.              \tag{16}
\]

It follows by summing the steps of \(\lfloor x\rfloor\) for \(\Re s>1\) and then by analytic continuation for \(\Re s>0\), \(s\ne1\). Its integral is bounded by \(1/\Re s\), giving \(O(1+|T|)\) on the indicated part of the ellipse. For \(\Re s\le1/4\), the functional equation

\[
 \zeta(s)=2^s\pi^{s-1}\sin(\pi s/2)\Gamma(1-s)\zeta(1-s)              \tag{17}
\]

and the vertical-strip gamma estimate give \(O((1+|T|)^{3/2-\Re s})\). On the stated strip this exponent is below \(5\). The common power \(6\) in (14) is therefore a valid upper bound with one fixed constant. The gamma estimate and functional equation are the classical inputs recorded in DLMF §§5.11 and 25.4; none assumes RH.

Every disc of radius one about a point of \([-1,1]\) lies inside the ellipse. The maximum principle and Cauchy's derivative formula therefore also give

\[
 \sup_{x\in[-1,1]}|f_T'(x)|\le M_T.                                    \tag{18}
\]

\subsubsection{3.2 Propagate a nonzero value to the critical-line interval}\label{zlm:endpoint-arithmetic_original--propagate-a-nonzero-value-to-the-critical-line-interval}

Put

\[
 r_*=(3+\sqrt{13})/2<4,\qquad
 z_*=-3i/2=\tfrac12(-ir_*+(-ir_*)^{-1}),\qquad
 c_* =\zeta(2)^{-1}.                                                    \tag{19}
\]

The absolutely convergent reciprocal Dirichlet series gives \(|1/\zeta(2+iT)|\le\zeta(2)\), so \(|f_T(z_*)|\ge c_*>0\). Let \(m_T=\max_{[-1,1]}|f_T|\).

Apply the maximum principle on \(1\le|w|\le8\) to \(f_T((w+w^{-1})/2)\). Equivalently subtract from its logarithmic modulus the harmonic function linear in \(\log|w|\) with boundary values \(\log m_T\) and \(\log M_T\). Zeros are handled by the usual subharmonic maximum principle. At \(w=-ir_*\) this gives

\[
 c_*\le m_T^{\vartheta}M_T^{1-\vartheta},\qquad
 \vartheta=1-\frac{\log r_*}{\log8}>\frac13.
                                                                         \tag{20}
\]

Here \(0<c_*<1\) and \(M_T\ge1\). Consequently

\[
 m_T\ge c_*^3M_T^{-2}.                                                 \tag{21}
\]

This is the three-circles argument with every coordinate map and radius fixed. No polynomial approximation to \(\zeta\) or assumption about its zeros is used.

Choose a point attaining \(m_T\). By (18), on a one-sided subinterval of length \(m_T/(2M_T)\le1/2\) inside \([-1,1]\), the modulus is at least \(m_T/2\). Thus

\[
 \int_{T-1}^{T+1}|\zeta(1/2+it)|^2\,dt
 \ge \frac{m_T^3}{8M_T}
 \ge \frac{c_*^9}{8C_0^7}(1+|T|)^{-42}.                                \tag{22}
\]

For \(T\in[-10,10]\), the integral in (22) is continuous and strictly positive: otherwise \(\zeta\) would vanish on an interval and hence identically on its holomorphic domain. Its minimum on that compact interval is positive. Combining it with (22) proves the all-real statement

\[
 \boxed{\exists c_\zeta>0\quad
 \int_{T-1}^{T+1}|\zeta(1/2+it)|^2\,dt
 \ge c_\zeta(1+|T|)^{-42}\quad(T\in\mathbb R).}                         \tag{23}
\]

The exponent \(42\) is deliberately an explicit sufficient exponent, not an optimal local mean-square estimate. The proof is a classical analytic propagation argument specialized here to the needed arithmetic interval.

\subsection{4. A lower envelope for the actual convolution, after three factors}\label{zlm:endpoint-arithmetic_original--a-lower-envelope-for-the-actual-convolution-after-three-factors}

\subsubsection{\texorpdfstring{4.1 Local mass of the unchanged \(g/h\)}{4.1 Local mass of the unchanged }}\label{zlm:endpoint-arithmetic_original--local-mass-of-the-unchanged-gh}

The gamma estimate, with a positive lower constant valid also on compact sets, implies

\[
 |\Gamma(1/4+it/2)|^2
 \ge c_\Gamma(1+|t|)^{-1/2}e^{-\pi|t|/2}.                              \tag{24}
\]

Let \(H_h=\prod_{\rho\in Z}(1+|\rho-1/2|)^{m_\rho}\), so \(|h(1/2+it)|\le H_h(1+|t|)^d\). The entire identity \(g=hv_h\) makes this bound valid at selected centres as well: \(|g|\le H_h(1+|t|)^d|v_h|\). The factors in (6), including \(|s(s-1)|^2=(t^2+1/4)^2\), then give a positive constant \(c_h^{(0)}\) with

\[
 w_h(t)\ge c_h^{(0)}e^{-\alpha_0|t|}(1+|t|)^{-2d}
                   |\zeta(1/2+it)|^2,\qquad \alpha_0=\pi/2.            \tag{25}
\]

Equation (25) is a weaker bound than the available factor \((1+|t|)^{7/2-2d}\); it follows by an inequality, not by changing the density. On \([T-1,T+1]\), its exponential and polynomial weights are bounded below by \(e^{-\alpha_0}e^{-\alpha_0|T|}\) and \(2^{-2d}(1+|T|)^{-2d}\), respectively. Hence (23) proves

\[
 \int_{T-1}^{T+1}w_h(t)\,dt
 \ge c_h^{(1)}e^{-\alpha_0|T|}(1+|T|)^{-B_h},\qquad B_h=42+2d.          \tag{26}
\]

\subsubsection{4.2 The positive compact convolution factor}\label{zlm:endpoint-arithmetic_original--the-positive-compact-convolution-factor}

The function \(w_h\) is continuous, bounded, integrable, and positive outside a discrete set. Thus \(m_{h,2}=w_h*w_h\) is continuous and strictly positive at every real point. Indeed its integrand is positive almost everywhere; continuity follows from translation continuity in \(L^1\) and boundedness of the other factor. Define the \textbf{actual} compact quantities

\[
 b_h=\min_{|v|\le1}m_{h,2}(v)>0,\qquad
 \vartheta_h=\int_{-1}^{1}w_h(v)\,dv>0.                                 \tag{27}
\]

Both are determined by the original density; neither is assigned the value one. Restricting a nonnegative convolution integral gives

\[
 m_{h,3}(u)\ge
 b_h\int_{u-1}^{u+1}w_h(t)\,dt
 \ge c_h e^{-\alpha_0|u|}(1+|u|)^{-B_h},\quad c_h=b_hc_h^{(1)}>0.
                                                                         \tag{28}
\]

The omitted part is the nonnegative integral over \(|v|>1\) and remains part of \(m_{h,3}\). We use a lower bound, not a new compactly supported density.

For \(k\ge3\), restrict the remaining \(k-3\) variables to \([-1,1]\) and use (28). Their sum has absolute value at most \(k-3\). This gives the central analytic estimate

\[
 \boxed{
 m_{h,k}(u)\ge
 c_h\vartheta_h^{k-3}
 e^{-\alpha_0(|u|+k-3)}
 (1+|u|+k-3)^{-B_h}
 \quad(k\ge3,\ u\in\mathbb R).}                                       \tag{29}
\]

All mass and tensor factors are explicit. For \(k=3\) the factor \(\vartheta_h^0\) is the empty product. No statement for \(k=1,2\) is inferred by using a negative exponent. The amplification below only needs \(k\ge3\).

\subsection{5. Upper moments and exact lower monic norms}\label{zlm:endpoint-arithmetic_original--upper-moments-and-exact-lower-monic-norms}

Fix \(0<b<\pi/2\) and retain both one-factor Laplace values

\[
 M_h(\pm b)=\int e^{\pm bt}w_h(t)\,dt>0.                               \tag{30}
\]

They are finite by the original gamma decay. Fubini and positivity give

\[
 \int e^{b|u|}m_{h,k}(u)\,du
 \le M_h(b)^k+M_h(-b)^k.                                                \tag{31}
\]

The elementary exponential-series inequality \(x^j\le j!b^{-j}e^{bx}\) for \(x\ge0\) yields

\[
 \omega_{h,k,j}
 \le (2j)!b^{-2j}\bigl(M_h(b)^k+M_h(-b)^k\bigr).                       \tag{32}
\]

The trial polynomial is \((S-k/2)^j\), whose value on \(S=k/2+iu\) is \(i^ju^j\). Its phase has modulus one and remains in the source map; the generator \(A=M_S\) has not been shifted to define a different spectral problem.

For every leading coefficient of modulus one, the exact least squared norm of a degree-\(j\) polynomial on \([-L,L]\) with Lebesgue measure is

\[
 \mathfrak l_j(L)=
 \frac{2L^{2j+1}}{2j+1}
       \left(\frac{2^j(j!)^2}{(2j)!}\right)^2.                           \tag{33}
\]

To check all constants, Rodrigues' formula gives the Legendre polynomial \(P_j(x)\) with leading coefficient \((2j)!/(2^j(j!)^2)\) and squared norm \(2/(2j+1)\). Its scaled monic version on \([-L,L]\) is the unique minimizer: subtract it from any competing monic polynomial, and the difference has lower degree and is orthogonal to it. For the source coordinate \(S=k/2+iu\), multiply this minimizer by the retained leading phase \(i^j\). This proves (33) at the original type. See also DLMF §§18.3 and 18.5.

Restricting (11) to \([-L,L]\) and using (29) proves

\[
 \boxed{
 \omega_{h,k,j}\ge
 c_h\vartheta_h^{k-3}
 e^{-\alpha_0(L+k-3)}(1+L+k-3)^{-B_h}\,\mathfrak l_j(L)
 \quad(k\ge3,L>0).}                                                    \tag{34}
\]

The observation has the explicit isometric splitting into the inside and outside interval components. Addition recovers the full source function, and their squared norms add. In the split lift each zero component is the supported zero at its label. The outside component has not been declared absent.

\subsection{6. The norm-window estimate is proved}\label{zlm:endpoint-arithmetic_original--the-norm-window-estimate-is-proved}

Set \(j=n\) and \(L=n\) in (34), and use \(j=2n\) in (32). The exact finite bound is

\[
 \boxed{
 \left(\frac{\omega_{h,k,2n}}{\omega_{h,k,n}}\right)^{1/(2n)}
 \le\left[
 \frac{(4n)!b^{-4n}(M_h(b)^k+M_h(-b)^k)
       e^{\alpha_0(n+k-3)}(1+n+k-3)^{B_h}}
      {c_h\vartheta_h^{k-3}\mathfrak l_n(n)}
       \right]^{1/(2n)}.}                                              \tag{35}
\]

For \(n\ge k\ge3\), this is at most \(C_hn\). Here is one explicit sufficient constant, with \(M_* =\max\{1,M_h(b),M_h(-b)\}\):

\[
 \boxed{
 C_h=64b^{-2}\exp(\alpha_0)\,2^{B_h/2}
       \exp\!\left(\frac{B_h}{2\exp(1)}\right)
       \sqrt{M_*\max(1,c_h^{-1})\max(1,\vartheta_h^{-1})}.
 }                                                                     \tag{36}
\]

Proof of the uniform estimate: \(\binom{2n}{n}\le4^n\) in (33) gives \(\mathfrak l_n(n)\ge2n^{2n+1}/((2n+1)4^n)\). Also \((4n)!\le(4n)^{4n}\), \(M_h(b)^k+M_h(-b)^k\le2M_*^k\), \(n+k-3\le2n\), and \(1+n+k-3\le2n\). After taking the \(2n\)-th root the factorial and monic factors contribute at most \(64n\). The mass factors contribute at most the square root in (36); the exponential contributes at most \(\exp(\alpha_0)\). Finally \(n^{1/n}\le\exp(1/\exp(1))\) gives \((2n)^{B_h/(2n)}\le2^{B_h/2}\exp(B_h/(2\exp(1)))\). This proves (1).

The constant is intentionally not optimized. Its dependence on the fixed packet is allowed; its independence of the polynomial and tensor degrees is what is needed. Defining it via positive compact minima does not claim to have provided a numerical interval enclosure for those minima.

For a general source window the same two inequalities give, without asymptotic notation,

\[
 \left(\frac{\omega_{h,k,n+r}}{\omega_{h,k,n}}\right)^{1/(2r)}
 \le\left[
 \frac{(2(n+r))!b^{-2(n+r)}(M_h(b)^k+M_h(-b)^k)
       e^{\alpha_0(L+k-3)}(1+L+k-3)^{B_h}}
      {c_h\vartheta_h^{k-3}\mathfrak l_n(L)}
       \right]^{1/(2r)}.                                               \tag{37}
\]

\subsection{7. Compose the proved bound with the formalizer's endpoint theorem}\label{zlm:endpoint-arithmetic_original--compose-the-proved-bound-with-the-formalizers-endpoint-theorem}

For a nonempty cyclic packet \(E_{h,k}\), let \(G_N\) be its original canonical quotient metric, \(V_N=\det G_N>0\), and \(\epsilon_{h,k,N}\) the original rank-two control allowance. The new endpoint theorem supplied by PR \#23 is

\[
 \min_{0\le j<r}\epsilon_{h,k,n+j}
 \le
 \left(\frac{\omega_{h,k,n+r}}{\omega_{h,k,n}}\right)^{1/(2r)}
 \sinh\!\left(\frac1{2r}
       \log\frac{V_{n-1}V_n}{V_{n+r-1}V_{n+r}}\right),\qquad n\ge q.
                                                                         \tag{38}
\]

Its proof multiplies the local Toda bounds, applies concavity of \(\log\sinh x\), and uses the two exact telescopes. The zero-contraction case is handled before taking logarithms. This result is credited to the parallel formalization/research session. The proof and the distinction between checked finite components and written Jensen assembly are retained in its source.

Choose \(n\ge\max(q,k)\) and \(r=n\). Equations (1) and (38) yield the actual canonical selection estimate

\[
 \boxed{
 \min_{n\le N<2n}\epsilon_{h,k,N}
 \le C_hn\sinh\!\left(\frac1{2n}
                 \log\frac{V_{n-1}V_n}{V_{2n-1}V_{2n}}\right).
 }                                                                     \tag{39}
\]

The chosen index is a member of the existing canonical family. The source representative, Taylor unit, kernel, action and trace are not chosen again after selecting this index.

For the packet consisting exactly of the quartet in (2), \(q_k\ge k\) and \(L_{h,k}/q_k\ge\delta k/2\). The earlier exterior bound holds at every admitted index, including the one selected by (39). Since \(\sinh\) is increasing on \([0,\infty)\), (4) follows. Moreover \(\operatorname{arsinh}(ak)=\log k+O_a(1)\) for every fixed \(a>0\), which proves the stated threshold \(2\).

A proved upper bound with \(\limsup\mathcal B_{h,k}/(q_k\log k)<2\) would contradict (4). Such an upper bound is \textbf{not} an output of this note. The previous sufficient target \(\mathcal B_{h,k}=o(q_k\log k)\) remains sufficient, but (4) specifies the larger quantitative threshold against which a future upper estimate can be compared. No repeated cost-free exterior operation is used: the multiplicity correction in the new formalization is retained.

\subsection{8. The remaining four volumes retain their finite jet maps}\label{zlm:endpoint-arithmetic_original--the-remaining-four-volumes-retain-their-finite-jet-maps}

The previous source/relation determinant identity remains

\[
 V_N=\frac{\mathfrak D_{N+1}}{\mathfrak B_{N-q+1}}.
                                                                         \tag{40}
\]

In the parallel full-jet transfer coordinates (PR \#24), let \(a=N-q+1\), let \(\mathsf F_a\) be its \(2q\) by \(2q\) raw confluent-jet matrix, and let \(\mathsf V\) be its raw-derivative Vandermonde matrix, with every factorial retained. Its established identity is

\[
 V_N=\left(\prod_{j=a}^{a+q-1}\omega_{h,k,j}\right)
                       \frac{\det\mathsf V}{\det\mathsf F_a}.           \tag{41}
\]

For the window \(n=r=q\), the four transfer indices are exactly \(0,1,q,q+1\). Put \(\Omega_j=\prod_{i=j}^{j+q-1}\omega_{h,k,i}\). Then

\[
 \boxed{
 \exp\mathcal B_{h,k}
 =\frac{\Omega_0\Omega_1}{\Omega_q\Omega_{q+1}}
   \frac{\det\mathsf F_q\det\mathsf F_{q+1}}
        {\det\mathsf F_0\det\mathsf F_1}.
 }                                                                     \tag{42}
\]

The common Vandermonde cancels in this explicit ratio, not inside any underlying map. The transfer matrices and their inverses still need their actual arithmetic control; the fixed width alone gives no such bound.

The empty packet has \(q=0\), \(V_N=1\) as an empty determinant, and no positive-rank inverse or quartet estimate. The analytic inequalities (23)--(37) still apply to its theta seed, but are not statements about a nonexistent finite packet.

\subsection{9. The volume window is also two explicit retained-boundary costs}\label{zlm:endpoint-arithmetic_original--the-volume-window-is-also-two-explicit-retained-boundary-costs}

There is a second useful representation of the same remaining quantity, with no reference-density change. Fix \(i<j\), both admissible. In the original monic polynomial basis put

\[
 F_{i,j}=[b_{i+1},\ldots,b_j],\qquad
 \Omega_{i,j}=\operatorname{diag}(\omega_{i+1},\ldots,\omega_j),
 \qquad b_a=[p_a]_\chi.                                                 \tag{43}
\]

The kernel update is the actual sum \(K_j=K_i+F_{i,j}\Omega_{i,j}^{-1}F_{i,j}^*\), where \(K_i=G_i^{-1}\). The matrix determinant lemma gives

\[
 \boxed{\frac{V_i}{V_j}
 =\det\!\left(I_{j-i}+\Omega_{i,j}^{-1}F_{i,j}^*G_iF_{i,j}\right).}       \tag{44}
\]

This determinant has positive real value. Its factors are \(1+\lambda_a\), where \(\lambda_a\ge0\) are the generalized eigenvalues of the Hermitian pair \((F_{i,j}^*G_iF_{i,j},\Omega_{i,j})\). No non-Hermitian matrix is asserted positive in an unspecified Euclidean metric.

Let \(T_{i,j}\) have columns \(\mathcal V_{h,k}p_{i+1},\ldots,
\mathcal V_{h,k}p_j\). The original boundary-layer isomorphism is

\[
 \mathfrak b_{i,j}=T_{i,j}-R_iF_{i,j}:
 \mathbb C^{j-i}\xrightarrow{\sim}
 \mathcal D_j\cap\mathcal D_i^\perp,
 \qquad
 \mathfrak b_{i,j}^*\mathfrak b_{i,j}
 =\Omega_{i,j}+F_{i,j}^*G_iF_{i,j}.                                     \tag{45}
\]

Its jets vanish; its highest polynomial coefficients give injectivity. Any new boundary can be reduced by these highest coefficients to an old boundary, and orthogonality then gives surjectivity. These statements prove the displayed map with the same original relation spaces, not only a dimension count. The representative update is

\[
 R_j=R_i+
 \mathfrak b_{i,j}
 (\Omega_{i,j}+F_{i,j}^*G_iF_{i,j})^{-1}F_{i,j}^*G_i.                    \tag{46}
\]

Every correction in (46) is an original theta boundary. Its split quotient image is \(e\) at the receiving fibre, and the input relation remains available in (45). Its cochain primitives are the same fixed-order divisions and tensor signs as in the previous cyclic source construction.

Apply (44) to \((i,j)=(q-1,2q-1)\) and \((q,2q)\). This proves

\[
 \boxed{\mathcal B_{h,k}
 =\sum_a\log(1+\lambda_a^{(0)})
   +\sum_a\log(1+\lambda_a^{(1)}).}                                    \tag{47}
\]

In particular

\[
 \mathcal B_{h,k}\le
 \operatorname{Tr}(\Omega_{q-1,2q-1}^{-1}F^*G_{q-1}F)
 +\operatorname{Tr}(\Omega_{q,2q}^{-1}F^*G_qF),                          \tag{48}
\]

where each \(F\) has its own indicated source and target from (43). When this trace bound is too large, (47) retains the full logarithmic cost rather than replacing it by its linear upper bound. The analytic norm estimate in (1) does not bound (47); it supplies the other factor in (39).

\subsection{10. Deligne reading and what has actually been transferred}\label{zlm:endpoint-arithmetic_original--deligne-reading-and-what-has-actually-been-transferred}

The six supplied S20 parts were reassembled in this turn. Their individual chunk hashes matched, and the assembled archive has 216,580,466 bytes and SHA-256 \nolinkurl{e2005bf31e1fcf362765f73c315c855e0f504f24f34d3a0ab188049da522b7c8}. This verifies the supplied bytes, not every transcription claim in them.

I read the TeX for Definition 1.3.5, Lemmas 3.2.10--3.2.12, the tensor argument 3.2.13, and Corollary 3.3.6 in the nested source files. I compared those exact passages with the supplied top-level French/English frozen page records at printed pages 158, 203, and 206. No PDF text extraction or OCR was used. The source review records a genuine inherited transcription fault: old \nolinkurl{fr_seg2.tex} writes a weight bound \(\le2\) at 3.2.10, whereas the controlling page record says \(<2\); only the latter, with integrality, implies \(\le1\). The old file also corrupts the explicit squared-eigenvalue step in 3.2.13. Those old files remain unmodified; they are not used as authority for these comparisons. No claim of a full independent audit of Weil II is made.

Deligne's determinant/rank operation and the upper/dual-lower architecture remain attached to the programme's specified exterior and residue maps. Here the additional estimate is a characteristic-zero analytic estimate for the actual source, (1), followed by the endpoint selection (39). This is not a statement that the new base satisfies Deligne's finite-field hypotheses. The remaining volume estimate has not been inferred from purity, positivity of a Gram, or the mere existence of the support.

The new result is therefore a completed analytic component of the current route: its two endpoint monic norms have the required \(O_h(q_k)\) scale on the actual quartet window. The other session's four-volume theorem now leaves the explicitly specified relation-volume cost as the remaining quantitative input on that window.

\subsection{11. Proof and computation status}\label{zlm:endpoint-arithmetic_original--proof-and-computation-status}

Sections 3--7 contain the new written analytic argument. Sections 8--9 contain explicit finite comparisons and a handoff for checking them using the existing quotient and matrix library. The standard complex-analysis, gamma and Legendre facts used in the argument are proved to the extent needed or referenced below. No general novelty claim is made for these classical ingredients.

The accompanying checker tests finite exact identities: coordinate factors, Legendre extremal norms, the local algebra in the interval argument, convolution moments on declared Laplace models, source quotients, block updates, raw confluent-jet endpoint factors, the volume telescope, and retained support cases. These tests do not certify the analytic inequalities by sampling them. There is no new Lean certificate and no interval-certified value of \(C_h\).

This continuation makes no remote changes. The integration files are add-only, and the parallel formalization branches and the reader's existing sources are left untouched.

\subsection{References and pinned inputs}\label{zlm:endpoint-arithmetic_original--references-and-pinned-inputs}

\begin{itemize}
\tightlist
\item
  Deligne, P. \emph{La conjecture de Weil. II}. Publications Mathématiques de l'IHÉS \textbf{52} (1980), 137--252. Supplied S20 source and current page-local records; the selected reading and discrepancies are itemized in \nolinkurl{SOURCE_REVIEW.md}.
\item
  NIST Digital Library of Mathematical Functions: §§25.4 (zeta functional equation), 5.11 (gamma estimates), 18.3 and 18.5 (Legendre polynomials). Public source pages consulted: https://dlmf.nist.gov/25.4, https://dlmf.nist.gov/5.11, https://dlmf.nist.gov/18.3, https://dlmf.nist.gov/18.5.
\item
  Original programme: the delivered \emph{Tau Gamma Convolution Descent}, \emph{Tau Toda Volume Control}, \emph{Tau Cyclic Sum Control}, and \emph{Tau Exterior Trace Amplification} notes. Their definitions of the source, action, unit and canonical quotient are retained; this is not a complete replay of all their upstream analytic claims.
\item
  Parallel endpoint criterion: PR \#23, \nolinkurl{c720f40530eed2f5969dbabe94dfd3fddc0f507f}, \nolinkurl{workbenches/tau-toda-volume-formal/ENDPOINT_CRITERION.md}.
\item
  Parallel confluent transfer: PR \#24, \nolinkurl{dfcbba5cbf7fec8c9301fe242c13e741d762013e}, \nolinkurl{workbenches/tau-confluent-transfer/RESEARCH_NOTE.md}.
\item
  Inspected main: \nolinkurl{b32ca2128e0deb0eec6eafb860776a5d6a28dcbb}, after its merge of the original-Gram exterior/trace contribution.
\end{itemize}

% END PRESERVED INPUT provider_endpoint_arithmetic_original_source.tex
\endgroup

\clearpage
\begingroup


\setcounter{equation}{0}
\renewcommand{\theequation}{TGR.\arabic{equation}}
\section*{Original Gamma mass, recurrence and exact monic norms}
\addcontentsline{toc}{section}{Original Gamma mass, recurrence and exact monic norms}
% BEGIN PRESERVED INPUT provider_theta_gamma_reference.tex SHA256 a379a1e19de4ce885680e14bda466a37b4c1eb25228c19b931d573f58be167f9
\section{The exact gamma reference and the arithmetic Gram determinants}
\label{tgr:sec:theta-gamma-reference}

The norm sequence in (KL.24) is attached to the actual measure
\[
 d\mu(t)=|g(\tfrac12+it)|^2\frac{dt}{2\pi},\qquad g=2\xi.
                                                               \tag{TG.1}
\]
We construct an explicit gamma reference for this measure, prove its
polynomial norms, and retain the full arithmetic multiplier in finite
Gram determinants.  This provides another exact formula for the
original boundary allowance; no asymptotic estimate for an arithmetic
determinant is assumed.

\subsection{The original measure, phase, and packet cancellation}

The defining identity for the completed function gives, on the original
line \(s=1/2+it\),
\[
 \begin{aligned}
 g(\tfrac12+it)&=\Gamma(\tfrac14+it/2)\,a(t),\\
 a(t)&=-\pi^{-1/4}(t^2+\tfrac14)
           e^{-it\log\pi/2}\zeta(\tfrac12+it).
 \end{aligned}                                                   \tag{TG.2}
\]
Indeed \(s(s-1)=-(t^2+1/4)\) and
\(\pi^{-s/2}=\pi^{-1/4}e^{-it\log\pi/2}\).  Both the minus
sign and this phase remain in the multiplier.

For \(\lambda>0\), define
\[
 d\sigma_\lambda(t)=|\Gamma(\lambda+it/2)|^2\frac{dt}{2\pi}.
                                                               \tag{TG.3}
\]
In particular
\[
 d\mu=|a|^2d\sigma_{1/4},\qquad
 d\nu_h=|a_h|^2d\sigma_{1/4},\qquad
 a_h(t)=\frac{g(1/2+it)/h(1/2+it)}{\Gamma(1/4+it/2)}.
                                                               \tag{TG.4}
\]
The last expression defines the values at every zero of the original
full packet \(h\); away from those points it equals
\(a(t)/h(1/2+it)\).  The function \(g/h\) is entire because
the complete order of each selected zero is removed.  Gamma has no
zero and has no pole on the real \(t\) axis.  Thus (TG.4) retains
the actual cancellation at each critical packet centre, without
introducing a pole into the measure.

Multiplication gives surjective Hilbert isometries
\[
 \begin{aligned}
 U_a:L^2(\mu)&\longrightarrow L^2(\sigma_{1/4}),
                  &F&\longmapsto aF,\\
 U_{a_h}:L^2(\nu_h)&\longrightarrow L^2(\sigma_{1/4}),
                  &F&\longmapsto a_hF.
 \end{aligned}                                                    \tag{TG.5}
\]
The integral identity proves the isometry.  The zero set of each
multiplier on the real axis is discrete, since it comes from a
nonzero analytic function.  Given \(H\in L^2(\sigma_{1/4})\),
define \(F=H/a\), respectively \(H/a_h\), off this null set
and give \(F\) arbitrary finite values on it.  Then
\(\int |F|^2d\mu=\int |H|^2d\sigma_{1/4}\), respectively
the packet identity, proving surjectivity and the stated inverse.
These are maps between the indicated Hilbert spaces.  Their images
of polynomial subspaces are exactly \(a\mathbb C[t]\) and
\(a_h\mathbb C[t]\), with the displayed multipliers retained.
Moreover the packet isometry from (KL.23) satisfies the exact diagram
\[
 U_{a_h}\bigl(h(1/2+it)P(t)\bigr)=U_aP(t).
                                                               \tag{TG.6}
\]

\subsection{Gamma norms from a generating integral}

We give the integral derivation to fix all constants independently of
different conventions for Meixner--Pollaczek polynomials.  The beta
integral, with \(r=e^x\), gives
\[
 \int_{\mathbb R}\frac{e^{\lambda x}}{(1+e^x)^{2\lambda}}
                     e^{iux}\,dx
       =\frac{\Gamma(\lambda+iu)\Gamma(\lambda-iu)}
                    {\Gamma(2\lambda)}.
                                                               \tag{TG.7}
\]
For completeness the beta integral follows by applying the change of
variables \(v=r_1+r_2\), \(y=r_1/(r_1+r_2)\) to the product of
the two absolutely convergent gamma integrals; the Jacobian is \(v\).
The conditions \(\operatorname{Re}(\lambda\pm iu)>0\) hold for
real \(u\).  The function on the left before Fourier transformation
and its second derivative are integrable: at each end they are
bounded by a constant times \(e^{-\lambda|x|}\).
Two integrations by parts therefore make its Fourier transform
integrable.  Fourier inversion applies.  One direct justification is
to multiply the transform by \(e^{-\varepsilon u^2}\), apply
Fubini, and use the Gaussian approximate identity; dominated
convergence removes that factor because the transform is integrable.
Using its evenness and then \(t=2u\), we obtain
\[
 \int_{\mathbb R}e^{iyt}\,d\sigma_\lambda(t)
   =\frac{2\Gamma(2\lambda)}{(2\cosh y)^{2\lambda}}
   =2^{1-2\lambda}\Gamma(2\lambda)(\cosh y)^{-2\lambda}.
                                                               \tag{TG.8}
\]
The factor \(2\) in the numerator is the Jacobian of the original
change \(t=2u\).

The following integral proves the required exponential integrability.

The function \((2\cosh(x/2))^{-2\lambda}\) extends holomorphically
to \(|\operatorname{Im}x|<\pi\), using the branch positive on
the real axis.  On each closed narrower strip it is bounded by
\(C e^{-\lambda|\operatorname{Re}x|}\).  In (TG.7), for
\(u>0\), shift the contour to \(\mathbb R+ic\), where
\(0<c<\pi\).  The vertical sides of the truncated rectangle
tend to zero by that bound.  For \(u<0\), shift down instead.
It follows that
\[
 |\Gamma(\lambda+iu)|^2\leq C_{\lambda,c}e^{-c|u|}.
\]
Consequently \(\int e^{v|t|}d\sigma_\lambda(t)<\infty\) for
every \(0<v<\pi/2\), by choosing \(2v<c<\pi\).
This proves holomorphic continuation of both sides of (TG.8) to
\(|\operatorname{Im}y|<\pi/2\), where they agree by the identity
theorem and use the branch inherited from the real axis.

For \(z\) near zero use the logarithms that vanish at \(z=0\) and set
\[
 F_\lambda(z,t)=(1-iz)^{-\lambda+it/2}
                        (1+iz)^{-\lambda-it/2}
      =\sum_{n\geq0}P_n^{(\lambda)}(t/2;\pi/2)z^n.
                                                               \tag{TG.9}
\]
For real \(z\) near zero this is
\((1+z^2)^{-\lambda}\exp(t\arctan z)\).  Apply (TG.8) at
\(y=-i(\arctan z+\arctan w)\), where \(z,w\) are real and
sufficiently small.  The identity
\[
 \cos(\arctan z+\arctan w)
        =\frac{1-zw}{\sqrt{1+z^2}\sqrt{1+w^2}}
\]
then proves
\[
 \int_{\mathbb R}F_\lambda(z,t)F_\lambda(w,t)
                            d\sigma_\lambda(t)
   =2^{1-2\lambda}\Gamma(2\lambda)(1-zw)^{-2\lambda}.
                                                               \tag{TG.10}
\]
All derivatives in \(z,w\) at zero can be passed through the
integral: on a smaller closed bidisc their integrands are bounded
by a polynomial in \(|t|\) times \(e^{c|t|}\), with
\(c<\pi/2\).  Thus comparison of the Taylor coefficients in
(TG.10) is justified, not formal.

Define the monic polynomial and its unchanged squared norm by
\[
 b_n^{(\lambda)}(t)=n!P_n^{(\lambda)}(t/2;\pi/2),\qquad
 \gamma_n^{(\lambda)}=2^{1-2\lambda}n!\Gamma(n+2\lambda).
                                                               \tag{TG.11}
\]
The coefficient of \(t^n\) in (TG.9) is \(1/n!\), since
\(\arctan z=z+O(z^3)\).  Consequently \(b_n^{(\lambda)}\)
is monic.  Equation (TG.10) gives the full orthogonality identity
\[
 \int_{\mathbb R}b_n^{(\lambda)}(t)b_j^{(\lambda)}(t)
                   d\sigma_\lambda(t)
       =\gamma_n^{(\lambda)}\delta_{nj}.
                                                               \tag{TG.12}
\]
All these polynomials have real coefficients, as is also clear from
the real form of the generating function.  Differentiating (TG.9)
gives \((1+z^2)\partial_zF_\lambda=(t-2\lambda z)F_\lambda\).
Its coefficient of \(z^n\) proves
\[
 \begin{aligned}
 b_0^{(\lambda)}&=1,\qquad b_1^{(\lambda)}=t,\\
 b_{n+1}^{(\lambda)}&=t b_n^{(\lambda)}
              -n(n+2\lambda-1)b_{n-1}^{(\lambda)},\quad n\geq1,\\
 \gamma_n^{(\lambda)}/\gamma_{n-1}^{(\lambda)}
              &=n(n+2\lambda-1),\quad n\geq1.
 \end{aligned}                                                    \tag{TG.13}
\]
For the actual reference \(\lambda=1/4\), abbreviate
\(b_n=b_n^{(1/4)}\) and \(\gamma_n=\gamma_n^{(1/4)}\).  Thus
\[
 \boxed{\gamma_n=\sqrt2\,n!\Gamma(n+\tfrac12),\qquad
           \gamma_{n+1}/\gamma_n=(n+1)(n+\tfrac12).}
                                                               \tag{TG.14}
\]
These formulas agree, with the displayed Jacobian and monic factors,
with NIST DLMF \S18.19, equations (18.19.8)--(18.19.9), and \S18.22,
equation (18.22.8).

\subsection{Finite arithmetic determinants in the retained metrics}

For \(N\geq0\) let \(\mathord{\text{\fontencoding{OT1}\fontfamily{cmss}\selectfont\char0}}_N=\operatorname{diag}
(\gamma_0,\ldots,\gamma_N)\), and keep both matrices
\[
 \begin{aligned}
 \mathsf B_N&=\left(\int_{\mathbb R}
              b_i(t)b_j(t)|a(t)|^2d\sigma_{1/4}(t)\right)_{0\leq i,j\leq N},\\
 \mathsf J_N&=\mathord{\text{\fontencoding{OT1}\fontfamily{cmss}\selectfont\char0}}_N^{-1/2}\mathsf B_N
                       \mathord{\text{\fontencoding{OT1}\fontfamily{cmss}\selectfont\char0}}_N^{-1/2},\qquad
 D_N=\det\mathsf J_N,\quad D_{-1}=1.
 \end{aligned}                                                    \tag{TG.15}
\]
The diagonal factors are explicit invertible coordinate maps;
\(\mathsf B_N=\mathord{\text{\fontencoding{OT1}\fontfamily{cmss}\selectfont\char0}}_N^{1/2}\mathsf J_N
\mathord{\text{\fontencoding{OT1}\fontfamily{cmss}\selectfont\char0}}_N^{1/2}\) recovers every original inner product.
The integrals converge with no unproved growth estimate.  Indeed,
integration of the step function \(\lfloor x\rfloor\) gives,
initially for \(\operatorname{Re}s>1\), and then by analytic
continuation for \(\operatorname{Re}s>0\), \(s\ne1\),
\[
 \zeta(s)=\frac{s}{s-1}
         -s\int_1^\infty\{x\}x^{-s-1}\,dx.
\]
The initial equality follows by integrating each constant step on
\([n,n+1)\), or by summation by parts of the absolutely convergent
Dirichlet series.  The last integral converges absolutely and
locally uniformly in the asserted half-plane since \(0\leq\{x\}<1\).
On \(s=1/2+it\), it proves
\( |\zeta(s)|\leq1+2\sqrt{t^2+1/4}\), because
\(|s/(s-1)|=1\).  Equation (TG.2) therefore bounds \(|a(t)|\)
by a polynomial of degree three.  The exponential integrability
just proved gives all its polynomial moments.  The packet
multiplier has the same property on the tail because \(h\) is
monic, and is bounded on each compact interval by its cancelled
analytic expression in (TG.4).
Each matrix is positive definite:
the integral of the squared modulus of a nonzero polynomial is
positive, because \(a\) vanishes only at isolated real points.

\begin{theorem}[The actual theta norms as arithmetic determinants]
For all \(n\geq0\), with the initial value included,
\[
 \boxed{\mathfrak h_n=\gamma_n\frac{D_n}{D_{n-1}},\qquad
 \frac{\mathfrak h_{n+1}}{\mathfrak h_n}
   =(n+1)(n+\tfrac12)\frac{D_{n+1}D_{n-1}}{D_n^2}.}
                                                               \tag{TG.16}
\]
\end{theorem}
\begin{proof}
Let \(Q_n(t)\) be the monic polynomial orthogonal to degrees below
\(n\) in \(L^2(\mu)\).  Its existence and uniqueness follow
from the positive definite Gram matrix.  Expressing
\(Q_0,\ldots,Q_N\) in \(b_0,\ldots,b_N\) gives an upper
triangular matrix with diagonal entries one.  The determinant of
the Gram matrix is therefore unchanged, and orthogonality gives
\(\det\mathsf B_N=\prod_{j=0}^N\|Q_j\|_\mu^2\).

The algebra isomorphism
\(\Psi:\mathbb C[s]\to\mathbb C[t]\),
\((\Psi p)(t)=p(1/2+it)\), preserves the original integral
inner product.  On the affine space of degree-\(n\) monic
polynomials the map \(p\mapsto i^{-n}\Psi p\) is its monic
version and preserves the norm because \(|i^{-n}|=1\).
It takes all lower-degree polynomial subspaces onto themselves.
Thus the original monic theta residual has squared norm
\(\mathfrak h_n=\|Q_n\|_\mu^2\).  The degree-dependent phase
is used on this affine space; it is not asserted to define a linear
map on the entire polynomial algebra.

Finally \(D_N=\det\mathsf B_N/\prod_{j=0}^N\gamma_j\).
Divide this equality at \(N=n\) by the equality at \(N=n-1\),
taking an empty product to be one.  This proves the first formula.
Dividing consecutive formulas and using (TG.14) proves the second.
\end{proof}

Define \(\mathsf B_N^{(h)},\mathsf J_N^{(h)},D_N^{(h)}\)
by (TG.15) with the entire-cancelled multiplier \(a_h\) of
(TG.4).  The same proof, with no restriction on whether the packet
centres are critical, gives
\[
 \kappa_n^{(h)}=\gamma_n\frac{D_n^{(h)}}{D_{n-1}^{(h)}}.
                                                               \tag{TG.17}
\]
Combining (TG.17) and (TG.16) with (KL.24) yields the exact
packet-to-reference determinant map, for every \(m\geq0\),
\[
 \boxed{r_m^{(h)}=
  \frac{(d+m)!\,\Gamma(d+m+1/2)}{m!\,\Gamma(m+1/2)}
  \frac{D_{d+m}^{(h)}}{D_{d+m-1}^{(h)}}
  \frac{D_{m-1}}{D_m}.}
                                                               \tag{TG.18}
\]
Every determinant is strictly positive and the full degree
\(d=\sum_\rho m_\rho\) is retained.

For a packet stable under both ordinary conjugation and
\(\rho\mapsto1-\overline\rho\), and for every \(m\geq0\),
(KL.32) becomes
\[
 \boxed{\epsilon_m^2=(m+1)(m+\tfrac12)
       \frac{D_{m+1}D_{m-1}}{D_m^2}
           (r_m^{-1}-1)(1-r_{m+1}).}
                                                               \tag{TG.19}
\]
This is the exact same finite-stage allowance.  Its factor
\((m+1)(m+1/2)\) is explicit; the arithmetic determinant
quotient remains in the formula.

\subsection{Full confluent kernel and the original coordinates}

Set \(\widehat b_j(s)=i^j b_j((s-1/2)/i)\).  These are monic
polynomials in the original variable \(s\).  With
\(\mathord{\textsf{Δ}}_N=\operatorname{diag}(1,i,\ldots,i^N)\),
their actual packet Gram matrix is
\[
 \mathsf B_{N,s}^{(h)}=\mathord{\textsf{Δ}}_N^*
                  \mathsf B_N^{(h)}\mathord{\textsf{Δ}}_N.
                                                               \tag{TG.20}
\]
Indeed \(\widehat b_j(1/2+it)=i^j b_j(t)\); taking the
conjugate in the first argument gives exactly the displayed
congruence, including every phase.

Let \(\mathsf C_N\) have columns
\([\widehat b_j]_h\in E_h=\mathbb C[s]/(h)\), in the original
monomial remainder coordinates.  For \(N\geq d-1\),
\[
 \boxed{K_N=\mathsf C_N(\mathsf B_{N,s}^{(h)})^{-1}
                           \mathsf C_N^*,\qquad
  \mathbf G_{N-d}=U_{\varepsilon_h}^*K_N^{-1}
                           U_{\varepsilon_h}.}
                                                               \tag{TG.21}
\]
Here \(\mathbf G_{-1}\) is the fixed initial-section metric,
and \(\varepsilon_h=(j_h(g/h))^{-1}\) is the same full local
unit as before.  To prove the first identity, choose the actual
monic orthogonal polynomial basis through degree \(N\).  If
\(T\) is its invertible coefficient matrix in the
\(\widehat b_j\) basis, then
\[
 T^*\mathsf B_{N,s}^{(h)}T
     =\operatorname{diag}(\kappa_0,\ldots,\kappa_N).
\]
Inverting gives
\((\mathsf B_{N,s}^{(h)})^{-1}
 =T\operatorname{diag}(\kappa_j^{-1})T^*\).
The columns of \(\mathsf C_NT\) are the full remainders
\([p_j]_h\).  Substitution proves exactly
\(K_N=\sum_{j=0}^N[p_j]_h[p_j]_h^*/\kappa_j\), the original
interpolation kernel.  The first \(d\) monic polynomials span all
remainders, so \(\mathsf C_N\) is onto and \(K_N\) is positive
definite.  The second identity is the fixed-unit minimum-lift
formula already proved in (KL.27), with the exact index
\(N=d+m\).  No single remainder column is required to be nonzero
or invertible.

For a centre \(\rho\), put \(t_\rho=(\rho-1/2)/i\) and
write \(f^{[\ell]}=f^{(\ell)}/\ell!\).  The full Taylor
coordinates of these columns are explicitly
\[
 \widehat b_j^{[\ell]}(\rho)
     =i^{j-\ell}b_j^{[\ell]}(t_\rho),
                 \qquad 0\leq\ell<m_\rho.
                                                               \tag{TG.22}
\]
This follows by differentiating the affine argument \(\ell\)
times.  CRT converts these full Taylor coordinates to the monomial
remainder coordinates, and multiplication by
\(U_{\upsilon_h}\) retains all derivatives of \(g/h\).  Thus
(TG.21) computes the entire confluent kernel, including nilpotent
orders and critical centres.

\subsection{The finite map to Romik's gamma parameter}

Romik's expansion uses \(P_n^{(3/4)}(t/2;\pi/2)\), with the
unchanged variable \(t/2\); see his \S3 and Appendix A.2.
The exact finite change to the parameter in (TG.3) is
\[
 \begin{aligned}
 b_n^{(1/4)}(t)
   &=\sum_{k=0}^{\lfloor n/2\rfloor}
       \frac{n!}{(n-2k)!}\binom{1/2}{k}
                       b_{n-2k}^{(3/4)}(t),\\
 b_n^{(3/4)}(t)
   &=\sum_{k=0}^{\lfloor n/2\rfloor}
       \frac{n!}{(n-2k)!}\binom{-1/2}{k}
                       b_{n-2k}^{(1/4)}(t).
 \end{aligned}                                                    \tag{TG.23}
\]
To prove these formulas, divide the two generating functions:
\(F_{1/4}(z,t)=(1+z^2)^{1/2}F_{3/4}(z,t)\).
The binomial series at zero and coefficient comparison prove the
first formula.  Multiplication by \((1+z^2)^{-1/2}\) proves the
inverse.  Each finite coefficient matrix is triangular with diagonal
one, preserves parity, and is invertible.  Applied to the same
measure, it transports Gram matrices by its full congruence and
transports every Taylor column by the same coefficient matrix.
Consequently the kernel expression (TG.21) is invariant under this
change of polynomial basis.

The two reference measures themselves have the exact Hilbert-space
map
\[
 L^2(\sigma_{1/4})\longrightarrow L^2(\sigma_{3/4}),
 \qquad F(t)\longmapsto
        \frac{\Gamma(1/4+it/2)}{\Gamma(3/4+it/2)}F(t).
                                                               \tag{TG.24}
\]
The gamma quotient has neither zeros nor poles for real \(t\).
Cancellation of its squared modulus proves the isometry and its
inverse; this isometry has a multiplier, whereas (TG.23) is the
finite polynomial coordinate map.  Their respective domains and
metrics have both been given explicitly.

Romik's coefficient asymptotics concern the expansion of \(\Xi\)
in a fixed gamma-orthogonal basis.  Inoue's \S2 similarly treats
expansion coefficients and convergence.  In the present construction
the arithmetic norms are the determinants (TG.16).  Equations
(TG.20)--(TG.24) are the exact maps connecting those polynomial
presentations to these determinants; no assertion that their
coefficient asymptotics already estimate \(D_n\) is used.

\subsection{Exact scope of the exponential-weight theorem}

The original theorem of Lubinsky--Mhaskar--Saff does include tail
order one.  In their 1988 paper, Definition 2.1 requires
\(W=e^{-Q}\), with \(Q\) even and continuous, \(Q'\) existing
on \((0,\infty)\), \(xQ'(x)\) bounded at zero, and, eventually,
\[
 Q'>0,\qquad x^2|Q'''|/Q'\leq C,\qquad
            1+xQ''/Q'\longrightarrow\alpha>0.
                                                               \tag{TG.25}
\]
Theorem 2.3 admits a finite generalized Jacobi factor, a lower-degree
polynomial exponential, and a bounded nonnegative factor
\(\Psi(t)\to1\).  Its Theorem 2.4 replaces the polynomial
exponential by \(V\), but requires
\(\log V(t)/Q(t)\to0\), as well as bounded reciprocal
approximation on the expanding interval by polynomials of degree
\(o(n)\).  The precise statements are on printed pages 67--70;
their proof is on pages 80--81.  These hypotheses concern the
full weight, not its gamma factor alone.

We prove a concrete obstruction to inserting the actual arithmetic
multiplier into those displayed hypotheses.  Hardy's theorem
supplies real numbers \(T_j\to\infty\) with
\(g(1/2+iT_j)=0\).  Removing a fixed full polynomial packet
leaves an infinite subsequence of these zeros.  At each remaining
\(T_j\), the multiplier \(a_h\) of (TG.4) has the local form
\[
 a_h(t)=(t-T_j)^{\ell_j}v_j(t),\qquad
 \ell_j\geq1,\quad v_j(T_j)\ne0,
                                                               \tag{TG.26}
\]
where \(v_j\) is analytic in a neighbourhood of \(T_j\).
This follows from the Taylor series at an isolated zero, with its
entire order \(\ell_j\) retained.

Let \(B(t)>0\) be any continuous comparison amplitude near all
sufficiently large \(T_j\), and suppose it has no zeros there.
For every \(j\) and every \(\eta>0\), continuity and (TG.26)
give a nonempty interval \(I_{j,\eta}\) about \(T_j\) on which
\[
                 |a_h(t)|/B(t)<\eta.
                                                               \tag{TG.27}
\]
Thus a decomposition with this ratio as the bounded multiplier
\(\Psi\) cannot satisfy \(\Psi\to1\).  The obstruction
holds on intervals of positive measure, so changing the value only
at each zero does not repair it.  A finite generalized Jacobi
factor has only finitely many real zeros or singularities; choosing
the subsequence beyond them leaves (TG.27) unchanged.

For the logarithmic hypothesis choose any eventually positive
continuous \(Q\).  At a remaining \(T_j\), shrink the
neighbourhood so that \(Q(t)\leq Q(T_j)+1\), and choose
\(\eta=\exp(-j(Q(T_j)+1))\) in (TG.27).  On the resulting
positive-measure interval,
\[
 \frac{\log(|a_h(t)|/B(t))}{Q(t)}<-j
                  \quad(t\ne T_j).
                                                               \tag{TG.28}
\]
Hence this ratio cannot be \(V\) with \(\log V/Q\to0\).
An additional factor \(\Psi\to1\) is bounded away from zero
and infinity on the tail and does not remove the contradiction.
This disproves the indicated application of Theorems 2.3--2.4,
including almost-everywhere modifications of the weight.
It makes no claim that other asymptotic methods are unavailable.
The exact positive Gram matrices (TG.15)--(TG.21) retain the
arithmetic information at all those zeros and are available for
such further analysis.

\subsection{Primary source locators}

The content-reading route and source hashes are recorded in
\path{work/theta_norm_literature_route_20260912.json} and
\path{work/theta_norm_literature_sources/download_receipt.json}.
The principal public locators are:
\begin{itemize}
 \item D.~S.~Lubinsky, H.~N.~Mhaskar and E.~B.~Saff,
 \emph{A proof of Freud's conjecture for exponential weights},
 Constructive Approximation 4 (1988), 65--83, Definition 2.1,
 Theorems 2.3--2.5 and \S4:
 \url{https://math.vanderbilt.edu/saffeb/texts/81.pdf}.
 \item Dan Romik, \emph{Orthogonal polynomial expansions for the
 Riemann xi function in the Hermite, Meixner--Pollaczek, and
 continuous Hahn bases}, \S3.1--3.3 and Appendix A.2:
 \url{https://www.math.ucdavis.edu/~romik/data/uploads/papers/riemannxi-final.pdf}.
 \item Hiroto Inoue, \emph{Expansion of Dirichlet L-function on the
 critical line in Meixner--Pollaczek polynomials}, arXiv:1412.1220,
 \S2.1--2.2:
 \url{https://arxiv.org/abs/1412.1220}.
 \item NIST DLMF, equations (18.19.8)--(18.19.9), (18.22.8):
 \url{https://dlmf.nist.gov/18.19},
 \url{https://dlmf.nist.gov/18.22}.
 \item G.~H.~Hardy, \emph{Sur les z\'eros de la fonction
 \(\zeta(s)\) de Riemann}, Comptes rendus 158 (1914),
 1012--1014, \S1--3.  The exact page images and transcriptions
 are linked in the machine-readable source route.
\end{itemize}

% END PRESERVED INPUT provider_theta_gamma_reference.tex
\endgroup

\endgroup

\clearpage
\begin{thebibliography}{999}

\bibitem{or:Ivic} A. Ivi\'c,
\emph{On the mean square of the Riemann zeta-function in short intervals},
arXiv:0803.0132v3, Section 2, (2.1) and the following \(E\)-analogue.
\url{https://arxiv.org/html/0803.0132}.
\bibitem{or:KVJ} H. T. Koelink and J. Van der Jeugt,
\emph{Bilinear generating functions for orthogonal polynomials},
arXiv:q-alg/9704016, Proposition 2.1, (2.9).
\url{https://arxiv.org/html/q-alg/9704016}.
\bibitem{or:DLMFhyper} NIST Digital Library of Mathematical Functions,
15.6.1, Euler integral for the regularized hypergeometric function.
\url{https://dlmf.nist.gov/15.6.E1}.
\bibitem{or:DLMFgamma} NIST Digital Library of Mathematical Functions,
5.11, Gamma asymptotic expansions.
\url{https://dlmf.nist.gov/5.11}.

\bibitem{rc:DLMF183}
NIST Digital Library of Mathematical Functions, Table 18.3.1,
\textit{Orthogonality properties for classical orthogonal polynomials},
\url{https://dlmf.nist.gov/18.3#T1}; consulted 19 September 2026.
\bibitem{rc:DLMF189}
NIST Digital Library of Mathematical Functions, Section 18.9(i),
\textit{Recurrence relations},
\url{https://dlmf.nist.gov/18.9#i}; consulted 19 September 2026.

\bibitem{peq:AKS} Split-Zero research programme (2026), \emph{Current kernel and multiplicity proofs}, supplied complete source edition; AKS3, AKS15--24, NCT, NPK and IGI6--14. The published 765 and supplied accepted 782 cuts are not identified with one another.
\bibitem{peq:LRC} Split-Zero mathematical continuation (2026, September 17), \emph{Lower-root covariance and all-rank conductor refinement}, supplied file \texttt{Pasted markdown(20260917-221927).md}, LRC6--10.
\bibitem{peq:PRT} Split-Zero research programme (2026), \emph{The original three-equation transfer of every connecting return}, PRT1--20; complete source in \texttt{ORIGINAL\_PROOF\_PROVIDERS(8).md}, DOC-0106, source SHA-256 \texttt{a0b5138b12764ffee5de29e9a579eb613abe544850e165a2f6624035878ff944}.
\bibitem{peq:EC} Coordinator continuation (2026, September 18), \emph{The retained eigenclass's marked pairing}, supplied complete EC proof, \texttt{Pasted markdown (2)(10).md}; EC37--41.
\bibitem{peq:human:boyd-vandenberghe}
Stephen Boyd and Lieven Vandenberghe,
\emph{Convex Optimization}, Cambridge University Press, 2004.
Appendix A.5.5, p.~650 (Schur minimum and determinant factorization),
and Appendix C.4.1, p.~673 (block elimination).
Author-hosted edition: \url{https://web.stanford.edu/~boyd/cvxbook/bv_cvxbook.pdf}.
Directly inspected PDF pages 664--665 and 687--688 on 17 September 2026.
The complex Hermitian form required here is proved explicitly in STF8--12.

\bibitem{peq:human:dlmf18}
T.~H. Koornwinder, R. Wong, R. Koekoek, R.~F. Swarttouw and
W.~P. Reinhardt,
``Orthogonal Polynomials,'' Chapter 18 of the
\emph{NIST Digital Library of Mathematical Functions},
version 1.2.8 (15 September 2026).
Christoffel--Darboux and its confluent form: equations (18.2.12)--(18.2.13),
\url{https://dlmf.nist.gov/18.2.E12},
\url{https://dlmf.nist.gov/18.2.E13}.
Classical norms and leading coefficients: Table 18.3.1,
\url{https://dlmf.nist.gov/18.3.T1}.
Rodrigues: equation (18.5.5) and Table 18.5.1,
\url{https://dlmf.nist.gov/18.5.E5},
\url{https://dlmf.nist.gov/18.5.T1}.
Laguerre series and zero value: equations (18.5.12) and (18.6.1),
\url{https://dlmf.nist.gov/18.5.E12},
\url{https://dlmf.nist.gov/18.6.E1}.
Classical recurrence: equation (18.9.1) and Table 18.9.1,
\url{https://dlmf.nist.gov/18.9.E1},
\url{https://dlmf.nist.gov/18.9.T1}.
Meixner--Pollaczek weight, norm and recurrence:
equations (18.19.8)--(18.19.9) and (18.22.8),
\url{https://dlmf.nist.gov/18.19},
\url{https://dlmf.nist.gov/18.22.E8}.
Chapter authors: \url{https://dlmf.nist.gov/18}.
Accessed 17 September 2026.


% BEGIN WCF_BIB_APPEND human:dlmf18
Generating function: equation (18.23.7),
\url{https://dlmf.nist.gov/18.23.E7}.
This locator was directly rechecked on 18 September 2026.
% END WCF_BIB_APPEND human:dlmf18

\bibitem{peq:human:globalization2026}
\emph{A Note on the Globalization Semiring G(Z)}, Zenodo,
12 March 2026. Deposited creators: ChatGPT, 5.4 and Gemini, 2.5 DeepThink.
Originating programme directed by KokunoYumeto with AI assistance,
as expressly identified by the programme author.
\url{https://doi.org/10.5281/zenodo.18976982}.
\url{https://zenodo.org/records/18976982}.
This attribution does not claim global novelty.

\bibitem{peq:human:splitzero2026}
The Clankers (deposited creator), \emph{A Secondary Note on Split-Zero
Globalization: Semimodule Classification, Multiplicative Reconstruction,
and All-Orders Zeta Deformation}, Zenodo, record publication date
20 March 2026. Originating programme directed by KokunoYumeto
with AI assistance. \url{https://doi.org/10.5281/zenodo.19120905}.
\url{https://zenodo.org/records/19120905}.
The deposited PDF prints 29 April 2026; record and printed dates
are retained separately. No global priority claim is made.

\bibitem{et:Incoming}
Split-Zero mathematical continuation,
\emph{Complete additive proofs}, 18 September 2026,
sections E, S and J, supplied as \texttt{COMPLETE\_PROOFS.md}.
The corresponding supplied narrative is
\emph{Bidirectional integration of the arithmetic centre, endpoint tensors,
and D032 calculations}, \texttt{INTEGRATION\_REPORT.md}.
\bibitem{et:Current}
Split-Zero research programme,
\emph{Current cumulative mathematical source},
\texttt{09\_UPDATED\_JOINT\_NOTE.tex}, accepted local successor,
OPG1--5; PQO5; RPC8; PCL1--6.
Exact local paths, byte hashes and extracted line intervals are
recorded in \texttt{SOURCE\_PINS.json}.
\bibitem{et:Envelopes}
Split-Zero research programme,
\emph{The full tilted mass in the original arithmetic tail}
and \emph{A global tilted-mass transfer with the central source retained},
RTM1--2 and GTM1, preserved provider
\texttt{RTM\_GTM\_original\_providers.tex}.
\bibitem{et:Deligne}
P. Deligne, \emph{La conjecture de Weil. II},
Publications Math\'ematiques de l'IH\'ES \textbf{52} (1980), 137--252,
\href{https://doi.org/10.1007/BF02684780}{doi:10.1007/BF02684780}.
Full supplied mathematical source retained in the source collection.
\bibitem{et:DLMF}
R. A. Askey and R. Roy, \emph{Gamma Function}, Chapter 5 of the
NIST Digital Library of Mathematical Functions,
\href{https://dlmf.nist.gov/5.2}{\S5.2(i)}, consulted 19 September 2026.

\bibitem{cpt:complete}
Split-Zero mathematical continuation,
\emph{Complete additive proofs}, 18 September 2026,
source file \texttt{COMPLETE\_PROOFS.md},
central arithmetic density note, equations C1--C22.
The new path transport in this note extends its C3--C13,
C17 and C21--C22; all original source bytes are preserved.

\bibitem{cpt:providers}
Split-Zero mathematical continuation,
\emph{The full tilted mass in the original arithmetic tail;
A global tilted-mass transfer with the central source retained},
source file \texttt{RTM\_GTM\_original\_providers.tex},
RTM1--RTM3 and GTM1, GTM7.
This is the original-source provider carried inside the
18 September 2026 signed arithmetic value package.

\bibitem{cpt:centralProvider}
Split-Zero mathematical continuation,
\emph{Central source window},
source file \texttt{CENTRAL\_SOURCE\_WINDOW.md},
ACW1--ACW13 and ACW20--ACW26,
carried inside the 18 September 2026 signed arithmetic
value package.  This note reconstructs its needed Fourier
estimate in the appendix.

\bibitem{cpt:sourceRelation}
Split-Zero mathematical continuation,
\emph{The signed return through the actual adjacent relation and quotient},
source file \texttt{SXR\_original\_source\_relation.tex},
SXR1--SXR10, carried inside the same package.
The finite Schur and signed-band identities are proved
again here with the two-parameter measure.

\bibitem{cpt:dlmfLegendreNorm}
T.\ H.\ Koornwinder, R.\ Wong, R.\ Koekoek, R.\ F.\ Swarttouw
and W.\ P.\ Reinhardt,
NIST Digital Library of Mathematical Functions,
Chapter 18, Table 18.3.1, Legendre orthogonality and leading
coefficient conventions.
\url{https://dlmf.nist.gov/18.3#T1}.
Consulted 19 September 2026.

\bibitem{cpt:dlmfLegendreRec}
T.\ H.\ Koornwinder, R.\ Wong, R.\ Koekoek, R.\ F.\ Swarttouw
and W.\ P.\ Reinhardt,
NIST Digital Library of Mathematical Functions,
Chapter 18, Table 18.9.1, Legendre recurrence.
\url{https://dlmf.nist.gov/18.9#T1}.
Consulted 19 September 2026.

\bibitem{cpt:dlmfMP}
T.\ H.\ Koornwinder, R.\ Wong, R.\ Koekoek, R.\ F.\ Swarttouw
and W.\ P.\ Reinhardt,
NIST Digital Library of Mathematical Functions,
Chapter 18, equation 18.22.8, Meixner--Pollaczek recurrence.
\url{https://dlmf.nist.gov/18.22.E8}.
Consulted 19 September 2026.

\bibitem{cpt:dlmfFourier}
NIST Digital Library of Mathematical Functions,
Section 1.14(i), especially (1.14.1), (1.14.4) and (1.14.8),
Fourier inversion and Parseval identities.
Its source notes cite E.\ C.\ Titchmarsh,
\emph{Introduction to the Theory of Fourier Integrals},
second edition, Chelsea, 1986, pp.\ 3--15, 42, 50--59.
\url{https://dlmf.nist.gov/1.14#i}.
Consulted 19 September 2026.

\bibitem{mii:deligne}
Pierre Deligne, \emph{La conjecture de Weil. II}, Publications
Math\'ematiques de l'IH\'ES \textbf{52} (1980), 137--252.
\href{https://doi.org/10.1007/BF02684780}{doi:10.1007/BF02684780}.
Primary record:
\url{https://www.numdam.org/item/PMIHES_1980__52__137_0/}.
Used: (1.4), (3.4.1)(iii), (3.7.2.3), and the lissity assertion
in (3.7.3). Local supplied full English mathematical source:
\texttt{D032\_FULL\_EN\_source\_checked\_corrected.tex}.

\bibitem{mii:dlmf}
NIST Digital Library of Mathematical Functions,
\emph{Gamma Function}, \S5.5(iii), Gauss's multiplication formula,
equations (5.5.6) and (5.5.7).
\url{https://dlmf.nist.gov/5.5.E6};
\url{https://dlmf.nist.gov/5.5.E7}.

\bibitem{mii:received}
Split-Zero programme, \emph{The original marked exponential family:
Deligne constituents and the retained-class jet divisor},
19 September 2026, nine-page received PDF
\texttt{DELIGNE\_MARKED\_IMAGE.pdf}, DMI1--DMI27.
The accompanying user-supplied transcript contains the same
calculation and its input/output record. Its source SHA256 and
the exact local paths are preserved in \texttt{ACCEPTANCE.json}.
This reference credits the received calculation separately from
the human results explicitly cited above.

\bibitem{ep:continuation}
The Clankers, \emph{Residue prices, constrained composition, and nonperiodic
optimal schedules}, 17 September 2026.  User-supplied
\texttt{EP817\_CONTINUATION\_20260917}; full source preserved at
\texttt{source/research/residue\_composition/}.
Baseline repository
\url{https://github.com/KokunoYumeto/erdos-problem-817-workbench},
revision \texttt{dbd0a93f2cf935103e88e5c2b2b71fe85ae7537b}.
\bibitem{ep:support}
KokunoYumeto, \emph{Reconstruction with changes of support index},
\texttt{workbenches/splitzero-tandem/tex/support\_diagrams.tex},
equations (D6)--(D8), revision
\texttt{58626a62cd5648e8ae7450fb54d4a4aab2330981},
\url{https://github.com/KokunoYumeto/zeta-function-research-reader}.
The pinned source and the retained local version with (D9)--(D14) were
read during this integration.
\bibitem{ep:devos}
Matt DeVos, \emph{A short proof of Kneser's addition theorem for abelian
groups}, 2013, \url{https://arxiv.org/abs/1303.3539v1}.
Role: the classical stabilizer inequality in the preserved residue-price
proof; the primary statement was checked on 19 September 2026.

\bibitem{human:boyd-vandenberghe}
Stephen Boyd and Lieven Vandenberghe,
\emph{Convex Optimization}, Cambridge University Press, 2004.
Appendix A.5.5, p.~650 (Schur minimum and determinant factorization),
and Appendix C.4.1, p.~673 (block elimination).
Author-hosted edition: \url{https://web.stanford.edu/~boyd/cvxbook/bv_cvxbook.pdf}.
Directly inspected PDF pages 664--665 and 687--688 on 17 September 2026.
The complex Hermitian form required here is proved explicitly in STF8--12.


\bibitem{human:boyd1994}
Stephen Boyd, Laurent El Ghaoui, Eric Feron and Venkataramanan Balakrishnan,
\emph{Linear Matrix Inequalities in System and Control Theory},
SIAM Studies in Applied Mathematics 15, SIAM, 1994.
\href{https://doi.org/10.1137/1.9781611970777}{doi:10.1137/1.9781611970777}.
Author-hosted edition: \url{https://web.stanford.edu/~boyd/lmibook/lmibook.pdf}.
The cited strict block formula is \S2.1, pp.~7--8, (2.3)--(2.4).


\bibitem{human:dlmf18}
T.~H. Koornwinder, R. Wong, R. Koekoek, R.~F. Swarttouw and
W.~P. Reinhardt,
``Orthogonal Polynomials,'' Chapter 18 of the
\emph{NIST Digital Library of Mathematical Functions},
version 1.2.8 (15 September 2026).
Christoffel--Darboux and its confluent form: equations (18.2.12)--(18.2.13),
\url{https://dlmf.nist.gov/18.2.E12},
\url{https://dlmf.nist.gov/18.2.E13}.
Classical norms and leading coefficients: Table 18.3.1,
\url{https://dlmf.nist.gov/18.3.T1}.
Rodrigues: equation (18.5.5) and Table 18.5.1,
\url{https://dlmf.nist.gov/18.5.E5},
\url{https://dlmf.nist.gov/18.5.T1}.
Laguerre series and zero value: equations (18.5.12) and (18.6.1),
\url{https://dlmf.nist.gov/18.5.E12},
\url{https://dlmf.nist.gov/18.6.E1}.
Classical recurrence: equation (18.9.1) and Table 18.9.1,
\url{https://dlmf.nist.gov/18.9.E1},
\url{https://dlmf.nist.gov/18.9.T1}.
Meixner--Pollaczek weight, norm and recurrence:
equations (18.19.8)--(18.19.9) and (18.22.8),
\url{https://dlmf.nist.gov/18.19},
\url{https://dlmf.nist.gov/18.22.E8}.
Chapter authors: \url{https://dlmf.nist.gov/18}.
Accessed 17 September 2026.


% BEGIN WCF_BIB_APPEND human:dlmf18
Generating function: equation (18.23.7),
\url{https://dlmf.nist.gov/18.23.E7}.
This locator was directly rechecked on 18 September 2026.
% END WCF_BIB_APPEND human:dlmf18

\bibitem{human:globalization2026}
\emph{A Note on the Globalization Semiring G(Z)}, Zenodo,
12 March 2026. Deposited creators: ChatGPT, 5.4 and Gemini, 2.5 DeepThink.
Originating programme directed by KokunoYumeto with AI assistance,
as expressly identified by the programme author.
\url{https://doi.org/10.5281/zenodo.18976982}.
\url{https://zenodo.org/records/18976982}.
This attribution does not claim global novelty.


\bibitem{human:hager1989}
William W.~Hager, ``Updating the inverse of a matrix,''
\emph{SIAM Review} \textbf{31}(2) (1989), 221--239.
\href{https://doi.org/10.1137/1031049}{doi:10.1137/1031049}.
\url{https://people.clas.ufl.edu/hager/files/update-1.pdf}.
Inverse formulas (1)--(2), p.~221; block elimination (6a)--(6b),
p.~222; historical discussion, pp.~221--223.


\bibitem{human:lyons2003}
Russell Lyons, ``Determinantal probability measures,''
\emph{Publications Math\'ematiques de l'IH\'ES} \textbf{98} (2003), 167--212.
\href{https://doi.org/10.1007/s10240-003-0016-0}{doi:10.1007/s10240-003-0016-0}.
Inspected author version: \href{https://arxiv.org/abs/math/0204325v4}{arXiv:math/0204325v4},
18 June 2003. Locators here refer to that version: \S4, (4.1), p.~11.


\bibitem{human:splitzero2026}
The Clankers (deposited creator), \emph{A Secondary Note on Split-Zero
Globalization: Semimodule Classification, Multiplicative Reconstruction,
and All-Orders Zeta Deformation}, Zenodo, record publication date
20 March 2026. Originating programme directed by KokunoYumeto
with AI assistance. \url{https://doi.org/10.5281/zenodo.19120905}.
\url{https://zenodo.org/records/19120905}.
The deposited PDF prints 29 April 2026; record and printed dates
are retained separately. No global priority claim is made.

\end{thebibliography}
\end{document}
